\documentclass[english]{book}
\usepackage{enumerate}
\usepackage{amssymb}
\usepackage{amsmath}
\usepackage{amsthm}
\usepackage{latexsym}
\usepackage[T1]{fontenc}
\usepackage[latin1]{inputenc}


\makeatletter
%% overlay.sty = overlay and align two symbols, respecting changing styles

\def\loverlay#1#2{\mathpalette\@overlay{{#1}{#2}{}{\hfil}}}
\def\overlay#1#2{\mathpalette\@overlay{{#1}{#2}{\hfil}{\hfil}}}
\def\roverlay#1#2{\mathpalette\@overlay{{#1}{#2}{\hfil}{}}}
   % calls to \@overlay look like
   %    \@overlay\textstyle{{x}{y}{\hfil}{\hfil}}
 
\def\@overlay#1#2{\@@overlay#1#2}
   % strip brackets from 2nd arg, to get
   %    \@@overlay\textstyle{x}{y}{\hfil}{\hfil}

\def\@@overlay#1#2#3#4#5{{
   \def\overlaystyle{#1}
   \setbox0=\hbox{\m@th$\overlaystyle#2$}
   \setbox1=\hbox{\m@th$\overlaystyle#3$}
   \ifdim \wd0<\wd1 \setbox2=\box1 \setbox1=\box0 \setbox0=\box2\fi
      % \box0 is now the wider box
   \rlap{\hbox to\wd0{#4\box1\relax#5}}\box0
}}
% end overlay command definitions
\makeatother


% angle measure symbol command, using overlay
\newcommand{\angm}{\overlay{<}{)} }

% arc command
\newcommand{\arc}[1]{\overset{\frown}{#1} }

% degree symbol command
\newcommand{\degrees}{\ensuremath{^{\circ}}}

%page settings
\evensidemargin 10pt
\topmargin 10pt
\makeatletter
\usepackage{babel}


% JDLM headers

\usepackage{fancyhdr}
\pagestyle{fancy}

\let\titlepageAuthor\author
\renewcommand{\author}[1]{\newcommand{\headerAuthor}{#1}\titlepageAuthor{#1}} 
\let\titlepageTitle\title
\renewcommand{\title}[1]{\newcommand{\headerTitle}{#1}\titlepageTitle{#1}} 

\lhead{\headerTitle: \leftmark} \chead{} \rhead{\thepage} 
\lfoot{\copyright \headerAuthor} \cfoot{} \rfoot{ - MathNerds Course Guide Collection - www.jiblm.org}
\renewcommand{\headrulewidth}{0.4pt}
\renewcommand{\footrulewidth}{0.4pt}


\begin{document}


\thispagestyle{empty}
\title{Trigonometry}
\author{Ed Reagh}
\date{}
\maketitle

\thispagestyle{empty}
\pagenumbering{roman}
\setcounter{tocdepth}{3} 
\tableofcontents

\chapter*{Chapter I}
\addcontentsline{toc}{chapter}{\numberline{}Chapter I}
\markboth{Chapter I}{Chapter I}

\pagenumbering{arabic}


We are going to try to solve the following problem: Suppose we are given a circle $K$ with center $C$ and radius $r > 0$. This means that, if $C: (c_{1},c_{2})$, then $K: (x-c_{1})^2 + (y-c_{2})^2 = r^{2}$.

Suppose also that we are given two points $P$ and $Q$ of $K$.

% Figure 1.1

\begin{center}
%TeXCAD Options
%\grade{\on}
%\emlines{\off}
%\epic{\off}
%\beziermacro{\on}
%\reduce{\on}
%\snapping{\off}
%\quality{8.00}
%\graddiff{0.01}
%\snapasp{1}
%\zoom{4.0000}
\unitlength 1mm % = 2.85pt
\linethickness{0.4pt}
\ifx\plotpoint\undefined\newsavebox{\plotpoint}\fi % GNUPLOT compatibility
\begin{picture}(110.5,77.49)(0,47)
%\circle(71.75,86.75){75.48}
\put(109.49,86.75){\line(0,1){1.383}}
\put(109.46,88.13){\line(0,1){1.381}}
\multiput(109.39,89.51)(-.0316,.3442){4}{\line(0,1){.3442}}
\multiput(109.26,90.89)(-.02947,.22857){6}{\line(0,1){.22857}}
\multiput(109.09,92.26)(-.03242,.19486){7}{\line(0,1){.19486}}
\multiput(108.86,93.63)(-.03075,.15053){9}{\line(0,1){.15053}}
\multiput(108.58,94.98)(-.03262,.13437){10}{\line(0,1){.13437}}
\multiput(108.26,96.32)(-.03127,.11091){12}{\line(0,1){.11091}}
\multiput(107.88,97.65)(-.032594,.101248){13}{\line(0,1){.101248}}
\multiput(107.46,98.97)(-.03369,.092844){14}{\line(0,1){.092844}}
\multiput(106.98,100.27)(-.032435,.080104){16}{\line(0,1){.080104}}
\multiput(106.47,101.55)(-.033268,.074224){17}{\line(0,1){.074224}}
\multiput(105.9,102.81)(-.032179,.065276){19}{\line(0,1){.065276}}
\multiput(105.29,104.05)(-.032821,.06085){20}{\line(0,1){.06085}}
\multiput(104.63,105.27)(-.03336,.056769){21}{\line(0,1){.056769}}
\multiput(103.93,106.46)(-.032338,.050682){23}{\line(0,1){.050682}}
\multiput(103.19,107.63)(-.032749,.047402){24}{\line(0,1){.047402}}
\multiput(102.4,108.77)(-.033085,.044324){25}{\line(0,1){.044324}}
\multiput(101.58,109.88)(-.033352,.041425){26}{\line(0,1){.041425}}
\multiput(100.71,110.95)(-.033557,.038687){27}{\line(0,1){.038687}}
\multiput(99.8,112)(-.033703,.036095){28}{\line(0,1){.036095}}
\multiput(98.86,113.01)(-.033796,.033635){29}{\line(-1,0){.033796}}
\multiput(97.88,113.98)(-.036255,.033531){28}{\line(-1,0){.036255}}
\multiput(96.86,114.92)(-.038847,.033372){27}{\line(-1,0){.038847}}
\multiput(95.81,115.82)(-.041583,.033155){26}{\line(-1,0){.041583}}
\multiput(94.73,116.68)(-.044481,.032873){25}{\line(-1,0){.044481}}
\multiput(93.62,117.51)(-.047558,.032523){24}{\line(-1,0){.047558}}
\multiput(92.48,118.29)(-.053146,.033555){22}{\line(-1,0){.053146}}
\multiput(91.31,119.03)(-.056927,.03309){21}{\line(-1,0){.056927}}
\multiput(90.11,119.72)(-.061006,.032531){20}{\line(-1,0){.061006}}
\multiput(88.89,120.37)(-.069063,.033638){18}{\line(-1,0){.069063}}
\multiput(87.65,120.98)(-.074381,.032914){17}{\line(-1,0){.074381}}
\multiput(86.39,121.54)(-.080258,.032053){16}{\line(-1,0){.080258}}
\multiput(85.1,122.05)(-.093004,.033247){14}{\line(-1,0){.093004}}
\multiput(83.8,122.51)(-.101402,.032112){13}{\line(-1,0){.101402}}
\multiput(82.48,122.93)(-.12115,.03353){11}{\line(-1,0){.12115}}
\multiput(81.15,123.3)(-.13453,.03198){10}{\line(-1,0){.13453}}
\multiput(79.8,123.62)(-.15067,.03004){9}{\line(-1,0){.15067}}
\multiput(78.45,123.89)(-.19501,.03149){7}{\line(-1,0){.19501}}
\multiput(77.08,124.11)(-.2287,.02838){6}{\line(-1,0){.2287}}
\put(75.71,124.28){\line(-1,0){1.378}}
\put(74.33,124.4){\line(-1,0){1.381}}
\put(72.95,124.47){\line(-1,0){1.383}}
\put(71.57,124.49){\line(-1,0){1.382}}
\put(70.19,124.46){\line(-1,0){1.38}}
\multiput(68.81,124.38)(-.3441,-.0333){4}{\line(-1,0){.3441}}
\multiput(67.43,124.24)(-.22842,-.03056){6}{\line(-1,0){.22842}}
\multiput(66.06,124.06)(-.1947,-.03335){7}{\line(-1,0){.1947}}
\multiput(64.7,123.83)(-.15038,-.03147){9}{\line(-1,0){.15038}}
\multiput(63.34,123.54)(-.13422,-.03326){10}{\line(-1,0){.13422}}
\multiput(62,123.21)(-.11076,-.0318){12}{\line(-1,0){.11076}}
\multiput(60.67,122.83)(-.101092,-.033076){13}{\line(-1,0){.101092}}
\multiput(59.36,122.4)(-.086504,-.031856){15}{\line(-1,0){.086504}}
\multiput(58.06,121.92)(-.079949,-.032816){16}{\line(-1,0){.079949}}
\multiput(56.78,121.39)(-.074064,-.033621){17}{\line(-1,0){.074064}}
\multiput(55.52,120.82)(-.065122,-.03249){19}{\line(-1,0){.065122}}
\multiput(54.29,120.21)(-.060693,-.033111){20}{\line(-1,0){.060693}}
\multiput(53.07,119.54)(-.056609,-.03363){21}{\line(-1,0){.056609}}
\multiput(51.88,118.84)(-.050527,-.032579){23}{\line(-1,0){.050527}}
\multiput(50.72,118.09)(-.047246,-.032974){24}{\line(-1,0){.047246}}
\multiput(49.59,117.3)(-.044166,-.033296){25}{\line(-1,0){.044166}}
\multiput(48.48,116.46)(-.041266,-.033549){26}{\line(-1,0){.041266}}
\multiput(47.41,115.59)(-.037151,-.032536){28}{\line(-1,0){.037151}}
\multiput(46.37,114.68)(-.034695,-.032707){29}{\line(-1,0){.034695}}
\multiput(45.36,113.73)(-.033474,-.033956){29}{\line(0,-1){.033956}}
\multiput(44.39,112.75)(-.033358,-.036415){28}{\line(0,-1){.036415}}
\multiput(43.46,111.73)(-.033187,-.039005){27}{\line(0,-1){.039005}}
\multiput(42.56,110.68)(-.032956,-.041741){26}{\line(0,-1){.041741}}
\multiput(41.71,109.59)(-.032661,-.044637){25}{\line(0,-1){.044637}}
\multiput(40.89,108.47)(-.0337,-.049786){23}{\line(0,-1){.049786}}
\multiput(40.11,107.33)(-.033301,-.053305){22}{\line(0,-1){.053305}}
\multiput(39.38,106.16)(-.032818,-.057084){21}{\line(0,-1){.057084}}
\multiput(38.69,104.96)(-.03224,-.06116){20}{\line(0,-1){.06116}}
\multiput(38.05,103.73)(-.033309,-.069222){18}{\line(0,-1){.069222}}
\multiput(37.45,102.49)(-.03256,-.074537){17}{\line(0,-1){.074537}}
\multiput(36.89,101.22)(-.03167,-.08041){16}{\line(0,-1){.08041}}
\multiput(36.39,99.93)(-.032804,-.093161){14}{\line(0,-1){.093161}}
\multiput(35.93,98.63)(-.031628,-.101554){13}{\line(0,-1){.101554}}
\multiput(35.52,97.31)(-.03296,-.12131){11}{\line(0,-1){.12131}}
\multiput(35.16,95.98)(-.03134,-.13468){10}{\line(0,-1){.13468}}
\multiput(34.84,94.63)(-.03298,-.16967){8}{\line(0,-1){.16967}}
\multiput(34.58,93.27)(-.03057,-.19516){7}{\line(0,-1){.19516}}
\multiput(34.36,91.91)(-.03275,-.2746){5}{\line(0,-1){.2746}}
\put(34.2,90.53){\line(0,-1){1.378}}
\put(34.09,89.15){\line(0,-1){1.381}}
\put(34.02,87.77){\line(0,-1){1.383}}
\put(34.01,86.39){\line(0,-1){1.382}}
\put(34.05,85.01){\line(0,-1){1.38}}
\multiput(34.14,83.63)(.02792,-.27514){5}{\line(0,-1){.27514}}
\multiput(34.28,82.25)(.03165,-.22827){6}{\line(0,-1){.22827}}
\multiput(34.47,80.88)(.02999,-.17022){8}{\line(0,-1){.17022}}
\multiput(34.71,79.52)(.03219,-.15023){9}{\line(0,-1){.15023}}
\multiput(35,78.17)(.03082,-.12187){11}{\line(0,-1){.12187}}
\multiput(35.34,76.83)(.03232,-.1106){12}{\line(0,-1){.1106}}
\multiput(35.73,75.5)(.033557,-.100933){13}{\line(0,-1){.100933}}
\multiput(36.16,74.19)(.032268,-.086351){15}{\line(0,-1){.086351}}
\multiput(36.65,72.89)(.033197,-.079792){16}{\line(0,-1){.079792}}
\multiput(37.18,71.62)(.032086,-.069798){18}{\line(0,-1){.069798}}
\multiput(37.75,70.36)(.0328,-.064966){19}{\line(0,-1){.064966}}
\multiput(38.38,69.13)(.0334,-.060535){20}{\line(0,-1){.060535}}
\multiput(39.05,67.92)(.032359,-.053882){22}{\line(0,-1){.053882}}
\multiput(39.76,66.73)(.032819,-.050371){23}{\line(0,-1){.050371}}
\multiput(40.51,65.57)(.033199,-.047088){24}{\line(0,-1){.047088}}
\multiput(41.31,64.44)(.033506,-.044007){25}{\line(0,-1){.044007}}
\multiput(42.15,63.34)(.032496,-.039583){27}{\line(0,-1){.039583}}
\multiput(43.02,62.27)(.032712,-.036996){28}{\line(0,-1){.036996}}
\multiput(43.94,61.24)(.032872,-.034539){29}{\line(0,-1){.034539}}
\multiput(44.89,60.24)(.034115,-.033312){29}{\line(1,0){.034115}}
\multiput(45.88,59.27)(.036573,-.033184){28}{\line(1,0){.036573}}
\multiput(46.91,58.34)(.039163,-.033){27}{\line(1,0){.039163}}
\multiput(47.96,57.45)(.041897,-.032757){26}{\line(1,0){.041897}}
\multiput(49.05,56.6)(.044792,-.032448){25}{\line(1,0){.044792}}
\multiput(50.17,55.79)(.049946,-.033462){23}{\line(1,0){.049946}}
\multiput(51.32,55.02)(.053463,-.033047){22}{\line(1,0){.053463}}
\multiput(52.5,54.29)(.05724,-.032546){21}{\line(1,0){.05724}}
\multiput(53.7,53.61)(.06454,-.03363){19}{\line(1,0){.06454}}
\multiput(54.93,52.97)(.06938,-.032979){18}{\line(1,0){.06938}}
\multiput(56.18,52.37)(.074691,-.032204){17}{\line(1,0){.074691}}
\multiput(57.44,51.83)(.08593,-.033373){15}{\line(1,0){.08593}}
\multiput(58.73,51.33)(.093316,-.03236){14}{\line(1,0){.093316}}
\multiput(60.04,50.87)(.11018,-.03374){12}{\line(1,0){.11018}}
\multiput(61.36,50.47)(.12146,-.03238){11}{\line(1,0){.12146}}
\multiput(62.7,50.11)(.13482,-.0307){10}{\line(1,0){.13482}}
\multiput(64.05,49.8)(.16982,-.03217){8}{\line(1,0){.16982}}
\multiput(65.41,49.55)(.1953,-.02964){7}{\line(1,0){.1953}}
\multiput(66.77,49.34)(.27476,-.03145){5}{\line(1,0){.27476}}
\put(68.15,49.18){\line(1,0){1.379}}
\put(69.52,49.08){\line(1,0){1.382}}
\put(70.91,49.02){\line(1,0){1.383}}
\put(72.29,49.01){\line(1,0){1.382}}
\put(73.67,49.06){\line(1,0){1.379}}
\multiput(75.05,49.15)(.275,.02923){5}{\line(1,0){.275}}
\multiput(76.43,49.3)(.22812,.03274){6}{\line(1,0){.22812}}
\multiput(77.79,49.5)(.17008,.0308){8}{\line(1,0){.17008}}
\multiput(79.16,49.74)(.15007,.0329){9}{\line(1,0){.15007}}
\multiput(80.51,50.04)(.12172,.0314){11}{\line(1,0){.12172}}
\multiput(81.84,50.39)(.11045,.03285){12}{\line(1,0){.11045}}
\multiput(83.17,50.78)(.093574,.031607){14}{\line(1,0){.093574}}
\multiput(84.48,51.22)(.086197,.032679){15}{\line(1,0){.086197}}
\multiput(85.77,51.71)(.079633,.033576){16}{\line(1,0){.079633}}
\multiput(87.05,52.25)(.069644,.032419){18}{\line(1,0){.069644}}
\multiput(88.3,52.83)(.064809,.033109){19}{\line(1,0){.064809}}
\multiput(89.53,53.46)(.060375,.033688){20}{\line(1,0){.060375}}
\multiput(90.74,54.14)(.053728,.032615){22}{\line(1,0){.053728}}
\multiput(91.92,54.85)(.050214,.033059){23}{\line(1,0){.050214}}
\multiput(93.08,55.61)(.046929,.033423){24}{\line(1,0){.046929}}
\multiput(94.2,56.42)(.043846,.033715){25}{\line(1,0){.043846}}
\multiput(95.3,57.26)(.039428,.032684){27}{\line(1,0){.039428}}
\multiput(96.36,58.14)(.03684,.032888){28}{\line(1,0){.03684}}
\multiput(97.4,59.06)(.034382,.033036){29}{\line(1,0){.034382}}
\multiput(98.39,60.02)(.033149,.034273){29}{\line(0,1){.034273}}
\multiput(99.35,61.01)(.033009,.036731){28}{\line(0,1){.036731}}
\multiput(100.28,62.04)(.032814,.03932){27}{\line(0,1){.03932}}
\multiput(101.16,63.1)(.032557,.042053){26}{\line(0,1){.042053}}
\multiput(102.01,64.2)(.033577,.046819){24}{\line(0,1){.046819}}
\multiput(102.82,65.32)(.033224,.050105){23}{\line(0,1){.050105}}
\multiput(103.58,66.47)(.032792,.05362){22}{\line(0,1){.05362}}
\multiput(104.3,67.65)(.032273,.057394){21}{\line(0,1){.057394}}
\multiput(104.98,68.86)(.033322,.0647){19}{\line(0,1){.0647}}
\multiput(105.61,70.09)(.032648,.069537){18}{\line(0,1){.069537}}
\multiput(106.2,71.34)(.031848,.074844){17}{\line(0,1){.074844}}
\multiput(106.74,72.61)(.032963,.086088){15}{\line(0,1){.086088}}
\multiput(107.24,73.9)(.031915,.093469){14}{\line(0,1){.093469}}
\multiput(107.68,75.21)(.03321,.11034){12}{\line(0,1){.11034}}
\multiput(108.08,76.54)(.0318,.12162){11}{\line(0,1){.12162}}
\multiput(108.43,77.87)(.0334,.14997){9}{\line(0,1){.14997}}
\multiput(108.73,79.22)(.03136,.16997){8}{\line(0,1){.16997}}
\multiput(108.98,80.58)(.03349,.22801){6}{\line(0,1){.22801}}
\multiput(109.18,81.95)(.03014,.2749){5}{\line(0,1){.2749}}
\put(109.33,83.33){\line(0,1){1.379}}
\put(109.43,84.7){\line(0,1){2.046}}
%\end
%\vector[both](71.75,87)(102.5,108)
\put(102.5,108){\vector(3,2){.07}}\put(71.75,87){\vector(-3,-2){.07}}
\multiput(71.75,87)(.049357945,.033707865){623}{\line(1,0){.049357945}}
%\end
%\qbezvec(108,118.25)(101.5,117.25)(100,113.25)
\put(100,113.25){\vector(-1,-2){.07}}\qbezier(108,118.25)(101.5,117.25)(100,113.25)
%\end
\put(33.25,106.25){\makebox(0,0)[cc]{P}}
\put(33.5,70.25){\makebox(0,0)[cc]{Q}}
\put(74.25,83.75){\makebox(0,0)[cc]{$C$}}
\put(86,94){\makebox(0,0)[cc]{$r$}}
\put(110.5,118.5){\makebox(0,0)[cc]{$K$}}
\put(39.75,106.25){\circle{.5}} 
\put(37,72){\circle{.5}}
\end{picture}

Figure 1.1
\end{center}


\begin{description}
\item[Notation:] $P:(x,y)$ means $P$ is a point with coordinates $(x,y)$. $M:y = ax+b$ means $M$ is a pointset with equation $y = ax+b$.
\end{description}

In this example, $M$ is a line with slope $a$ and y-intercept $b$. If we apply this notation to the information above, $C: (c_{1},c_{2})$ means $C$ is a point with coordinates $(c_{1},c_{2})$, and $K:(x-c_{1})^2 + (y-c_{2})^2 = r^2$ means $K$ is a pointset with equation $(x-c_{1})^{2} + (y-c_{2})^{2} = r^{2}$. You probably recognize this as an equation of a circle. From algebra, we can find the distance from $P$ to $Q$ by using the coordinates of $P$ and $Q$, if we know them. If $P:(p_{1},p_{2})$ and $Q:(q_{1},q_{2})$, then the distance from $P$ to $Q$, denoted $PQ$, is given by $PQ = \sqrt{(p_{1}-q_{1})^{2}+(p_{2}-q_{2})^{2}}$.

Note that this is a definition of distance between $P$ and $Q$.

% Figure 1.2

\begin{center}
%TeXCAD Options
%\grade{\on}
%\emlines{\off}
%\epic{\off}
%\beziermacro{\on}
%\reduce{\on}
%\snapping{\off}
%\quality{8.00}
%\graddiff{0.01}
%\snapasp{1}
%\zoom{4.0000}
\unitlength 1mm % = 2.85pt
\linethickness{0.4pt}
\ifx\plotpoint\undefined\newsavebox{\plotpoint}\fi % GNUPLOT compatibility
\begin{picture}(60.47,57.97)(0,51)
\leavevmode
%\circle(32.75,81.25){55.44}
\put(60.47,81.25){\line(0,1){1.15}}
\put(60.45,82.4){\line(0,1){1.148}}
\put(60.38,83.55){\line(0,1){1.144}}
\multiput(60.26,84.69)(-.03328,.22759){5}{\line(0,1){.22759}}
\multiput(60.09,85.83)(-.0305,.16144){7}{\line(0,1){.16144}}
\multiput(59.88,86.96)(-.03252,.14003){8}{\line(0,1){.14003}}
\multiput(59.62,88.08)(-.03064,.11085){10}{\line(0,1){.11085}}
\multiput(59.31,89.19)(-.03201,.09953){11}{\line(0,1){.09953}}
\multiput(58.96,90.28)(-.0331,.08994){12}{\line(0,1){.08994}}
\multiput(58.56,91.36)(-.031546,.075848){14}{\line(0,1){.075848}}
\multiput(58.12,92.42)(-.032354,.06951){15}{\line(0,1){.06951}}
\multiput(57.63,93.47)(-.033008,.063851){16}{\line(0,1){.063851}}
\multiput(57.11,94.49)(-.033532,.058755){17}{\line(0,1){.058755}}
\multiput(56.54,95.49)(-.032157,.051281){19}{\line(0,1){.051281}}
\multiput(55.93,96.46)(-.032544,.047408){20}{\line(0,1){.047408}}
\multiput(55.27,97.41)(-.03284,.043826){21}{\line(0,1){.043826}}
\multiput(54.58,98.33)(-.033055,.040498){22}{\line(0,1){.040498}}
\multiput(53.86,99.22)(-.033198,.037392){23}{\line(0,1){.037392}}
\multiput(53.09,100.08)(-.033273,.034484){24}{\line(0,1){.034484}}
\multiput(52.3,100.91)(-.034675,.033074){24}{\line(-1,0){.034675}}
\multiput(51.46,101.7)(-.037583,.032981){23}{\line(-1,0){.037583}}
\multiput(50.6,102.46)(-.040688,.032821){22}{\line(-1,0){.040688}}
\multiput(49.7,103.18)(-.044015,.032587){21}{\line(-1,0){.044015}}
\multiput(48.78,103.87)(-.047595,.03227){20}{\line(-1,0){.047595}}
\multiput(47.83,104.51)(-.054325,.033631){18}{\line(-1,0){.054325}}
\multiput(46.85,105.12)(-.058948,.033193){17}{\line(-1,0){.058948}}
\multiput(45.85,105.68)(-.064041,.032639){16}{\line(-1,0){.064041}}
\multiput(44.82,106.2)(-.069695,.031952){15}{\line(-1,0){.069695}}
\multiput(43.78,106.68)(-.081878,.0335){13}{\line(-1,0){.081878}}
\multiput(42.71,107.12)(-.09013,.03258){12}{\line(-1,0){.09013}}
\multiput(41.63,107.51)(-.09971,.03144){11}{\line(-1,0){.09971}}
\multiput(40.53,107.86)(-.12336,.03333){9}{\line(-1,0){.12336}}
\multiput(39.42,108.16)(-.14022,.03171){8}{\line(-1,0){.14022}}
\multiput(38.3,108.41)(-.16161,.02956){7}{\line(-1,0){.16161}}
\multiput(37.17,108.62)(-.22778,.03197){5}{\line(-1,0){.22778}}
\put(36.03,108.78){\line(-1,0){1.145}}
\put(34.89,108.89){\line(-1,0){1.148}}
\put(33.74,108.95){\line(-1,0){1.15}}
\put(32.59,108.97){\line(-1,0){1.15}}
\put(31.44,108.94){\line(-1,0){1.147}}
\multiput(30.29,108.86)(-.2858,-.0314){4}{\line(-1,0){.2858}}
\multiput(29.15,108.74)(-.1895,-.02883){6}{\line(-1,0){.1895}}
\multiput(28.01,108.56)(-.16126,-.03143){7}{\line(-1,0){.16126}}
\multiput(26.88,108.34)(-.13984,-.03333){8}{\line(-1,0){.13984}}
\multiput(25.77,108.08)(-.11067,-.03128){10}{\line(-1,0){.11067}}
\multiput(24.66,107.76)(-.09934,-.03258){11}{\line(-1,0){.09934}}
\multiput(23.57,107.41)(-.08975,-.03362){12}{\line(-1,0){.08975}}
\multiput(22.49,107)(-.075665,-.031983){14}{\line(-1,0){.075665}}
\multiput(21.43,106.56)(-.069322,-.032754){15}{\line(-1,0){.069322}}
\multiput(20.39,106.06)(-.06366,-.033376){16}{\line(-1,0){.06366}}
\multiput(19.37,105.53)(-.055307,-.031989){18}{\line(-1,0){.055307}}
\multiput(18.38,104.95)(-.051094,-.032453){19}{\line(-1,0){.051094}}
\multiput(17.4,104.34)(-.047219,-.032817){20}{\line(-1,0){.047219}}
\multiput(16.46,103.68)(-.043636,-.033092){21}{\line(-1,0){.043636}}
\multiput(15.54,102.99)(-.040306,-.033289){22}{\line(-1,0){.040306}}
\multiput(14.66,102.25)(-.0372,-.033413){23}{\line(-1,0){.0372}}
\multiput(13.8,101.49)(-.034291,-.033472){24}{\line(-1,0){.034291}}
\multiput(12.98,100.68)(-.032873,-.034865){24}{\line(0,-1){.034865}}
\multiput(12.19,99.85)(-.032764,-.037773){23}{\line(0,-1){.037773}}
\multiput(11.44,98.98)(-.032586,-.040876){22}{\line(0,-1){.040876}}
\multiput(10.72,98.08)(-.032332,-.044202){21}{\line(0,-1){.044202}}
\multiput(10.04,97.15)(-.033678,-.050295){19}{\line(0,-1){.050295}}
\multiput(9.4,96.19)(-.033317,-.054518){18}{\line(0,-1){.054518}}
\multiput(8.8,95.21)(-.032852,-.059138){17}{\line(0,-1){.059138}}
\multiput(8.24,94.21)(-.032269,-.064228){16}{\line(0,-1){.064228}}
\multiput(7.73,93.18)(-.031549,-.069879){15}{\line(0,-1){.069879}}
\multiput(7.25,92.13)(-.033027,-.08207){13}{\line(0,-1){.08207}}
\multiput(6.82,91.06)(-.03206,-.09032){12}{\line(0,-1){.09032}}
\multiput(6.44,89.98)(-.03086,-.09989){11}{\line(0,-1){.09989}}
\multiput(6.1,88.88)(-.03262,-.12355){9}{\line(0,-1){.12355}}
\multiput(5.81,87.77)(-.0309,-.1404){8}{\line(0,-1){.1404}}
\multiput(5.56,86.65)(-.0334,-.18874){6}{\line(0,-1){.18874}}
\multiput(5.36,85.51)(-.03065,-.22796){5}{\line(0,-1){.22796}}
\put(5.2,84.37){\line(0,-1){1.145}}
\put(5.1,83.23){\line(0,-1){1.149}}
\put(5.04,82.08){\line(0,-1){1.15}}
\put(5.03,80.93){\line(0,-1){1.149}}
\put(5.07,79.78){\line(0,-1){1.147}}
\multiput(5.15,78.63)(.0331,-.2856){4}{\line(0,-1){.2856}}
\multiput(5.28,77.49)(.02992,-.18933){6}{\line(0,-1){.18933}}
\multiput(5.46,76.36)(.03236,-.16108){7}{\line(0,-1){.16108}}
\multiput(5.69,75.23)(.03034,-.12413){9}{\line(0,-1){.12413}}
\multiput(5.96,74.11)(.03192,-.11049){10}{\line(0,-1){.11049}}
\multiput(6.28,73.01)(.03316,-.09915){11}{\line(0,-1){.09915}}
\multiput(6.65,71.92)(.031511,-.082664){13}{\line(0,-1){.082664}}
\multiput(7.06,70.84)(.032419,-.075479){14}{\line(0,-1){.075479}}
\multiput(7.51,69.78)(.033154,-.069132){15}{\line(0,-1){.069132}}
\multiput(8.01,68.75)(.031758,-.059733){17}{\line(0,-1){.059733}}
\multiput(8.55,67.73)(.032308,-.055122){18}{\line(0,-1){.055122}}
\multiput(9.13,66.74)(.032747,-.050906){19}{\line(0,-1){.050906}}
\multiput(9.75,65.77)(.033089,-.047029){20}{\line(0,-1){.047029}}
\multiput(10.41,64.83)(.033344,-.043444){21}{\line(0,-1){.043444}}
\multiput(11.11,63.92)(.033521,-.040113){22}{\line(0,-1){.040113}}
\multiput(11.85,63.04)(.033627,-.037006){23}{\line(0,-1){.037006}}
\multiput(12.62,62.19)(.033669,-.034097){24}{\line(0,-1){.034097}}
\multiput(13.43,61.37)(.035054,-.032671){24}{\line(1,0){.035054}}
\multiput(14.27,60.58)(.037961,-.032545){23}{\line(1,0){.037961}}
\multiput(15.15,59.83)(.041064,-.032349){22}{\line(1,0){.041064}}
\multiput(16.05,59.12)(.046607,-.03368){20}{\line(1,0){.046607}}
\multiput(16.98,58.45)(.050488,-.033388){19}{\line(1,0){.050488}}
\multiput(17.94,57.81)(.054709,-.033002){18}{\line(1,0){.054709}}
\multiput(18.93,57.22)(.059327,-.03251){17}{\line(1,0){.059327}}
\multiput(19.94,56.67)(.064413,-.031898){16}{\line(1,0){.064413}}
\multiput(20.97,56.16)(.075064,-.03337){14}{\line(1,0){.075064}}
\multiput(22.02,55.69)(.082259,-.032553){13}{\line(1,0){.082259}}
\multiput(23.09,55.27)(.0905,-.03154){12}{\line(1,0){.0905}}
\multiput(24.17,54.89)(.11008,-.03331){10}{\line(1,0){.11008}}
\multiput(25.27,54.56)(.12374,-.03191){9}{\line(1,0){.12374}}
\multiput(26.39,54.27)(.14057,-.03009){8}{\line(1,0){.14057}}
\multiput(27.51,54.03)(.18893,-.03231){6}{\line(1,0){.18893}}
\multiput(28.64,53.83)(.22813,-.02934){5}{\line(1,0){.22813}}
\put(29.79,53.69){\line(1,0){1.146}}
\put(30.93,53.59){\line(1,0){1.149}}
\put(32.08,53.54){\line(1,0){1.15}}
\put(33.23,53.53){\line(1,0){1.149}}
\put(34.38,53.58){\line(1,0){1.146}}
\multiput(35.53,53.67)(.22833,.02777){5}{\line(1,0){.22833}}
\multiput(36.67,53.81)(.18915,.03101){6}{\line(1,0){.18915}}
\multiput(37.8,53.99)(.16089,.03329){7}{\line(1,0){.16089}}
\multiput(38.93,54.23)(.12395,.03106){9}{\line(1,0){.12395}}
\multiput(40.04,54.5)(.1103,.03255){10}{\line(1,0){.1103}}
\multiput(41.15,54.83)(.09896,.03373){11}{\line(1,0){.09896}}
\multiput(42.24,55.2)(.08248,.031988){13}{\line(1,0){.08248}}
\multiput(43.31,55.62)(.075291,.032854){14}{\line(1,0){.075291}}
\multiput(44.36,56.08)(.068939,.033552){15}{\line(1,0){.068939}}
\multiput(45.4,56.58)(.059549,.032102){17}{\line(1,0){.059549}}
\multiput(46.41,57.13)(.054934,.032625){18}{\line(1,0){.054934}}
\multiput(47.4,57.71)(.050716,.03304){19}{\line(1,0){.050716}}
\multiput(48.36,58.34)(.046837,.033359){20}{\line(1,0){.046837}}
\multiput(49.3,59.01)(.043251,.033594){21}{\line(1,0){.043251}}
\multiput(50.21,59.71)(.038184,.032284){23}{\line(1,0){.038184}}
\multiput(51.08,60.46)(.035278,.03243){24}{\line(1,0){.035278}}
\multiput(51.93,61.23)(.032546,.032511){25}{\line(1,0){.032546}}
\multiput(52.74,62.05)(.032469,.035242){24}{\line(0,1){.035242}}
\multiput(53.52,62.89)(.032326,.038148){23}{\line(0,1){.038148}}
\multiput(54.27,63.77)(.033641,.043214){21}{\line(0,1){.043214}}
\multiput(54.97,64.68)(.033411,.046801){20}{\line(0,1){.046801}}
\multiput(55.64,65.61)(.033096,.05068){19}{\line(0,1){.05068}}
\multiput(56.27,66.58)(.032685,.054899){18}{\line(0,1){.054899}}
\multiput(56.86,67.57)(.032167,.059513){17}{\line(0,1){.059513}}
\multiput(57.41,68.58)(.033627,.068903){15}{\line(0,1){.068903}}
\multiput(57.91,69.61)(.032936,.075255){14}{\line(0,1){.075255}}
\multiput(58.37,70.66)(.032078,.082445){13}{\line(0,1){.082445}}
\multiput(58.79,71.74)(.03102,.09068){12}{\line(0,1){.09068}}
\multiput(59.16,72.82)(.03267,.11027){10}{\line(0,1){.11027}}
\multiput(59.49,73.93)(.03119,.12392){9}{\line(0,1){.12392}}
\multiput(59.77,75.04)(.03346,.16085){7}{\line(0,1){.16085}}
\multiput(60,76.17)(.03122,.18912){6}{\line(0,1){.18912}}
\multiput(60.19,77.3)(.02802,.2283){5}{\line(0,1){.2283}}
\put(60.33,78.44){\line(0,1){1.146}}
\put(60.42,79.59){\line(0,1){1.659}}
%\end
%\emline(51,101.75)(5.25,84.75)
\multiput(51,101.75)(-.09077381,-.033730159){504}{\line(-1,0){.09077381}}
%\end
\put(5.25,84.75){\line(0,1){0}} 
\put(5.25,84.75){\line(0,1){.25}}
\put(53.75,103.5){\makebox(0,0)[cc]{$Q$}}
\put(3.25,87.5){\makebox(0,0)[cc]{$P$}}
\put(34.75,79.25){\makebox(0,0)[cc]{$C$}}
\put(32.75,81.25){\circle{.71}}
\end{picture}
\\ Figure 1.2
\end{center}

From Figure 1.2, $PQ$ is the length of the linear path from $P$ to $Q$; i.e., the path from $P$ to $Q$ along the line containing these two points. This path is called the chord from $P$ to $Q$, or just the chord $PQ$. But note another path from $P$ to $Q$, namely along the circle. This path is called the circular arc from $P$ to $Q$, or just the arc from $P$ to $Q$. Can you see a difficulty in talking about ``the arc from $P$ to $Q$''? In any event, the problem we want to solve is this: If we can find the length of the chord from $P$ to $Q$, namely $PQ$, can we also find the length of the circular arc from $P$ to $Q$?

Before you read any further, see if you can solve this problem. If you can solve it satisfactorily, then maybe you can skip this course. And that suggests that spending a considerable amount of effort right now could be very profitable to you.

Now that you have spent some time wrestling with this problem, you probably realize that there does not appear to be any simple algebraic solution. How do you get a handle on such a problem? The question itself seems clear enough, but the answer is elusive.

One approach to seemingly intractable problems such as this is to try to find a special case of the problem that may be solvable, and then generalize the results to the original problem. Our original problem deals with any circle in the plane with any radius. Let's take a particular circle with a particular radius and see if this helps us gain some insight into our general problem. Consider the circle $K$ with radius $1$ and center $O: (0,0)$. Then $K: (x-0)^{2} + (y-0)^{2} =1^{2}$, or $K: x^{2} + y^{2} = 1$. $K$ is called the unit circle (radius $1$) with center at the origin ( $O: (0,0)$ ). Again, try to solve the problem before reading beyond this paragraph. We are \emph{not} in a hurry. There is no fast-forward to the learning process. I think that, today, we have a little difficulty acknowledging this idea, much more so than previous generations. Everything seems to get done much faster than before. We have computers that are so lightning fast that we come to expect results instantly. Being brought up in an atmosphere of instant results can lead us to think that there may be something wrong with us if we don't grasp everything instantly and without effort. Previous generations didn't have this problem. Things just didn't move as rapidly then, and they knew that things like understanding a difficult concept took time and effort. Although things move more quickly today, we are still the same as people of previous generations and we still need time and effort to understand things. Don't get frustrated because you don't get something immediately. It will come to you if you are patient with yourself and persevering.

Annoying, isn't it? There still doesn't seem to be any obvious handle that we can latch on to. We need some tools. We really ought to define the things we're talking about here. For example, what do we mean by an arc of a circle? It is almost always a crucial and very helpful step in solving a difficult problem to have everything defined precisely.

Consider the figure 1.3. Try to define the following statement: ``$A$ is the arc $PQ$ of the circle $K$.''

% Figure 1.3
\begin{center}
%TeXCAD Options
%\grade{\on}
%\emlines{\off}
%\epic{\off}
%\beziermacro{\on}
%\reduce{\on}
%\snapping{\off}
%\quality{8.00}
%\graddiff{0.01}
%\snapasp{1}
%\zoom{4.0000}
\unitlength 1mm % = 2.85pt
\linethickness{0.4pt}
\ifx\plotpoint\undefined\newsavebox{\plotpoint}\fi % GNUPLOT compatibility
\begin{picture}(72.25,68)(0,42)
\leavevmode
%\circle(32.75,81.25){55.44}
\put(60.47,81.25){\line(0,1){1.15}}
\put(60.45,82.4){\line(0,1){1.148}}
\put(60.38,83.55){\line(0,1){1.144}}
\multiput(60.26,84.69)(-.03328,.22759){5}{\line(0,1){.22759}}
\multiput(60.09,85.83)(-.0305,.16144){7}{\line(0,1){.16144}}
\multiput(59.88,86.96)(-.03252,.14003){8}{\line(0,1){.14003}}
\multiput(59.62,88.08)(-.03064,.11085){10}{\line(0,1){.11085}}
\multiput(59.31,89.19)(-.03201,.09953){11}{\line(0,1){.09953}}
\multiput(58.96,90.28)(-.0331,.08994){12}{\line(0,1){.08994}}
\multiput(58.56,91.36)(-.031546,.075848){14}{\line(0,1){.075848}}
\multiput(58.12,92.42)(-.032354,.06951){15}{\line(0,1){.06951}}
\multiput(57.63,93.47)(-.033008,.063851){16}{\line(0,1){.063851}}
\multiput(57.11,94.49)(-.033532,.058755){17}{\line(0,1){.058755}}
\multiput(56.54,95.49)(-.032157,.051281){19}{\line(0,1){.051281}}
\multiput(55.93,96.46)(-.032544,.047408){20}{\line(0,1){.047408}}
\multiput(55.27,97.41)(-.03284,.043826){21}{\line(0,1){.043826}}
\multiput(54.58,98.33)(-.033055,.040498){22}{\line(0,1){.040498}}
\multiput(53.86,99.22)(-.033198,.037392){23}{\line(0,1){.037392}}
\multiput(53.09,100.08)(-.033273,.034484){24}{\line(0,1){.034484}}
\multiput(52.3,100.91)(-.034675,.033074){24}{\line(-1,0){.034675}}
\multiput(51.46,101.7)(-.037583,.032981){23}{\line(-1,0){.037583}}
\multiput(50.6,102.46)(-.040688,.032821){22}{\line(-1,0){.040688}}
\multiput(49.7,103.18)(-.044015,.032587){21}{\line(-1,0){.044015}}
\multiput(48.78,103.87)(-.047595,.03227){20}{\line(-1,0){.047595}}
\multiput(47.83,104.51)(-.054325,.033631){18}{\line(-1,0){.054325}}
\multiput(46.85,105.12)(-.058948,.033193){17}{\line(-1,0){.058948}}
\multiput(45.85,105.68)(-.064041,.032639){16}{\line(-1,0){.064041}}
\multiput(44.82,106.2)(-.069695,.031952){15}{\line(-1,0){.069695}}
\multiput(43.78,106.68)(-.081878,.0335){13}{\line(-1,0){.081878}}
\multiput(42.71,107.12)(-.09013,.03258){12}{\line(-1,0){.09013}}
\multiput(41.63,107.51)(-.09971,.03144){11}{\line(-1,0){.09971}}
\multiput(40.53,107.86)(-.12336,.03333){9}{\line(-1,0){.12336}}
\multiput(39.42,108.16)(-.14022,.03171){8}{\line(-1,0){.14022}}
\multiput(38.3,108.41)(-.16161,.02956){7}{\line(-1,0){.16161}}
\multiput(37.17,108.62)(-.22778,.03197){5}{\line(-1,0){.22778}}
\put(36.03,108.78){\line(-1,0){1.145}}
\put(34.89,108.89){\line(-1,0){1.148}}
\put(33.74,108.95){\line(-1,0){1.15}}
\put(32.59,108.97){\line(-1,0){1.15}}
\put(31.44,108.94){\line(-1,0){1.147}}
\multiput(30.29,108.86)(-.2858,-.0314){4}{\line(-1,0){.2858}}
\multiput(29.15,108.74)(-.1895,-.02883){6}{\line(-1,0){.1895}}
\multiput(28.01,108.56)(-.16126,-.03143){7}{\line(-1,0){.16126}}
\multiput(26.88,108.34)(-.13984,-.03333){8}{\line(-1,0){.13984}}
\multiput(25.77,108.08)(-.11067,-.03128){10}{\line(-1,0){.11067}}
\multiput(24.66,107.76)(-.09934,-.03258){11}{\line(-1,0){.09934}}
\multiput(23.57,107.41)(-.08975,-.03362){12}{\line(-1,0){.08975}}
\multiput(22.49,107)(-.075665,-.031983){14}{\line(-1,0){.075665}}
\multiput(21.43,106.56)(-.069322,-.032754){15}{\line(-1,0){.069322}}
\multiput(20.39,106.06)(-.06366,-.033376){16}{\line(-1,0){.06366}}
\multiput(19.37,105.53)(-.055307,-.031989){18}{\line(-1,0){.055307}}
\multiput(18.38,104.95)(-.051094,-.032453){19}{\line(-1,0){.051094}}
\multiput(17.4,104.34)(-.047219,-.032817){20}{\line(-1,0){.047219}}
\multiput(16.46,103.68)(-.043636,-.033092){21}{\line(-1,0){.043636}}
\multiput(15.54,102.99)(-.040306,-.033289){22}{\line(-1,0){.040306}}
\multiput(14.66,102.25)(-.0372,-.033413){23}{\line(-1,0){.0372}}
\multiput(13.8,101.49)(-.034291,-.033472){24}{\line(-1,0){.034291}}
\multiput(12.98,100.68)(-.032873,-.034865){24}{\line(0,-1){.034865}}
\multiput(12.19,99.85)(-.032764,-.037773){23}{\line(0,-1){.037773}}
\multiput(11.44,98.98)(-.032586,-.040876){22}{\line(0,-1){.040876}}
\multiput(10.72,98.08)(-.032332,-.044202){21}{\line(0,-1){.044202}}
\multiput(10.04,97.15)(-.033678,-.050295){19}{\line(0,-1){.050295}}
\multiput(9.4,96.19)(-.033317,-.054518){18}{\line(0,-1){.054518}}
\multiput(8.8,95.21)(-.032852,-.059138){17}{\line(0,-1){.059138}}
\multiput(8.24,94.21)(-.032269,-.064228){16}{\line(0,-1){.064228}}
\multiput(7.73,93.18)(-.031549,-.069879){15}{\line(0,-1){.069879}}
\multiput(7.25,92.13)(-.033027,-.08207){13}{\line(0,-1){.08207}}
\multiput(6.82,91.06)(-.03206,-.09032){12}{\line(0,-1){.09032}}
\multiput(6.44,89.98)(-.03086,-.09989){11}{\line(0,-1){.09989}}
\multiput(6.1,88.88)(-.03262,-.12355){9}{\line(0,-1){.12355}}
\multiput(5.81,87.77)(-.0309,-.1404){8}{\line(0,-1){.1404}}
\multiput(5.56,86.65)(-.0334,-.18874){6}{\line(0,-1){.18874}}
\multiput(5.36,85.51)(-.03065,-.22796){5}{\line(0,-1){.22796}}
\put(5.2,84.37){\line(0,-1){1.145}}
\put(5.1,83.23){\line(0,-1){1.149}}
\put(5.04,82.08){\line(0,-1){1.15}}
\put(5.03,80.93){\line(0,-1){1.149}}
\put(5.07,79.78){\line(0,-1){1.147}}
\multiput(5.15,78.63)(.0331,-.2856){4}{\line(0,-1){.2856}}
\multiput(5.28,77.49)(.02992,-.18933){6}{\line(0,-1){.18933}}
\multiput(5.46,76.36)(.03236,-.16108){7}{\line(0,-1){.16108}}
\multiput(5.69,75.23)(.03034,-.12413){9}{\line(0,-1){.12413}}
\multiput(5.96,74.11)(.03192,-.11049){10}{\line(0,-1){.11049}}
\multiput(6.28,73.01)(.03316,-.09915){11}{\line(0,-1){.09915}}
\multiput(6.65,71.92)(.031511,-.082664){13}{\line(0,-1){.082664}}
\multiput(7.06,70.84)(.032419,-.075479){14}{\line(0,-1){.075479}}
\multiput(7.51,69.78)(.033154,-.069132){15}{\line(0,-1){.069132}}
\multiput(8.01,68.75)(.031758,-.059733){17}{\line(0,-1){.059733}}
\multiput(8.55,67.73)(.032308,-.055122){18}{\line(0,-1){.055122}}
\multiput(9.13,66.74)(.032747,-.050906){19}{\line(0,-1){.050906}}
\multiput(9.75,65.77)(.033089,-.047029){20}{\line(0,-1){.047029}}
\multiput(10.41,64.83)(.033344,-.043444){21}{\line(0,-1){.043444}}
\multiput(11.11,63.92)(.033521,-.040113){22}{\line(0,-1){.040113}}
\multiput(11.85,63.04)(.033627,-.037006){23}{\line(0,-1){.037006}}
\multiput(12.62,62.19)(.033669,-.034097){24}{\line(0,-1){.034097}}
\multiput(13.43,61.37)(.035054,-.032671){24}{\line(1,0){.035054}}
\multiput(14.27,60.58)(.037961,-.032545){23}{\line(1,0){.037961}}
\multiput(15.15,59.83)(.041064,-.032349){22}{\line(1,0){.041064}}
\multiput(16.05,59.12)(.046607,-.03368){20}{\line(1,0){.046607}}
\multiput(16.98,58.45)(.050488,-.033388){19}{\line(1,0){.050488}}
\multiput(17.94,57.81)(.054709,-.033002){18}{\line(1,0){.054709}}
\multiput(18.93,57.22)(.059327,-.03251){17}{\line(1,0){.059327}}
\multiput(19.94,56.67)(.064413,-.031898){16}{\line(1,0){.064413}}
\multiput(20.97,56.16)(.075064,-.03337){14}{\line(1,0){.075064}}
\multiput(22.02,55.69)(.082259,-.032553){13}{\line(1,0){.082259}}
\multiput(23.09,55.27)(.0905,-.03154){12}{\line(1,0){.0905}}
\multiput(24.17,54.89)(.11008,-.03331){10}{\line(1,0){.11008}}
\multiput(25.27,54.56)(.12374,-.03191){9}{\line(1,0){.12374}}
\multiput(26.39,54.27)(.14057,-.03009){8}{\line(1,0){.14057}}
\multiput(27.51,54.03)(.18893,-.03231){6}{\line(1,0){.18893}}
\multiput(28.64,53.83)(.22813,-.02934){5}{\line(1,0){.22813}}
\put(29.79,53.69){\line(1,0){1.146}}
\put(30.93,53.59){\line(1,0){1.149}}
\put(32.08,53.54){\line(1,0){1.15}}
\put(33.23,53.53){\line(1,0){1.149}}
\put(34.38,53.58){\line(1,0){1.146}}
\multiput(35.53,53.67)(.22833,.02777){5}{\line(1,0){.22833}}
\multiput(36.67,53.81)(.18915,.03101){6}{\line(1,0){.18915}}
\multiput(37.8,53.99)(.16089,.03329){7}{\line(1,0){.16089}}
\multiput(38.93,54.23)(.12395,.03106){9}{\line(1,0){.12395}}
\multiput(40.04,54.5)(.1103,.03255){10}{\line(1,0){.1103}}
\multiput(41.15,54.83)(.09896,.03373){11}{\line(1,0){.09896}}
\multiput(42.24,55.2)(.08248,.031988){13}{\line(1,0){.08248}}
\multiput(43.31,55.62)(.075291,.032854){14}{\line(1,0){.075291}}
\multiput(44.36,56.08)(.068939,.033552){15}{\line(1,0){.068939}}
\multiput(45.4,56.58)(.059549,.032102){17}{\line(1,0){.059549}}
\multiput(46.41,57.13)(.054934,.032625){18}{\line(1,0){.054934}}
\multiput(47.4,57.71)(.050716,.03304){19}{\line(1,0){.050716}}
\multiput(48.36,58.34)(.046837,.033359){20}{\line(1,0){.046837}}
\multiput(49.3,59.01)(.043251,.033594){21}{\line(1,0){.043251}}
\multiput(50.21,59.71)(.038184,.032284){23}{\line(1,0){.038184}}
\multiput(51.08,60.46)(.035278,.03243){24}{\line(1,0){.035278}}
\multiput(51.93,61.23)(.032546,.032511){25}{\line(1,0){.032546}}
\multiput(52.74,62.05)(.032469,.035242){24}{\line(0,1){.035242}}
\multiput(53.52,62.89)(.032326,.038148){23}{\line(0,1){.038148}}
\multiput(54.27,63.77)(.033641,.043214){21}{\line(0,1){.043214}}
\multiput(54.97,64.68)(.033411,.046801){20}{\line(0,1){.046801}}
\multiput(55.64,65.61)(.033096,.05068){19}{\line(0,1){.05068}}
\multiput(56.27,66.58)(.032685,.054899){18}{\line(0,1){.054899}}
\multiput(56.86,67.57)(.032167,.059513){17}{\line(0,1){.059513}}
\multiput(57.41,68.58)(.033627,.068903){15}{\line(0,1){.068903}}
\multiput(57.91,69.61)(.032936,.075255){14}{\line(0,1){.075255}}
\multiput(58.37,70.66)(.032078,.082445){13}{\line(0,1){.082445}}
\multiput(58.79,71.74)(.03102,.09068){12}{\line(0,1){.09068}}
\multiput(59.16,72.82)(.03267,.11027){10}{\line(0,1){.11027}}
\multiput(59.49,73.93)(.03119,.12392){9}{\line(0,1){.12392}}
\multiput(59.77,75.04)(.03346,.16085){7}{\line(0,1){.16085}}
\multiput(60,76.17)(.03122,.18912){6}{\line(0,1){.18912}}
\multiput(60.19,77.3)(.02802,.2283){5}{\line(0,1){.2283}}
\put(60.33,78.44){\line(0,1){1.146}}
\put(60.42,79.59){\line(0,1){1.659}}
%\end
\put(5.25,84.75){\line(0,1){0}} 
\put(5.25,84.75){\line(0,1){.25}}
\put(3.25,87.5){\makebox(0,0)[cc]{$P$}}
\put(34.75,79.25){\makebox(0,0)[cc]{$C$}}
\put(32.75,81.25){\circle{.71}}
%\emline(1.75,95.5)(47.25,49.75)
\multiput(1.75,95.5)(.0337286879,-.0339140104){1349}{\line(0,-1){.0339140104}}
%\end
\put(47.5,54.5){\makebox(0,0)[cc]{$Q$}}
%\qbezvec(69,103.5)(61.63,109)(55.75,101.5)
\put(55.75,101.5){\vector(-3,-4){.07}}\qbezier(69,103.5)(61.63,109)(55.75,101.5)
%\end
\put(72.25,101.25){\makebox(0,0)[cc]{$K$}}
\put(33.5,44){\makebox(0,0)[cc]{Figure 1.3}}
\end{picture}
\end{center}

Remember that, although you can use the picture as an aid in constructing your definition, the definition itself must stand alone; i.e., the definition must not be picture-dependent. Make a serious effort to define the statement above, \emph{before} you go further.

Definitions can often be tricky things. When you try to define a word, you must use other words in your definition. But what do these other words mean? If we try to define all of them, then we introduce even more words that must be defined. There seems to be no end to this process. What we end up doing is assuming that the person for whom we are defining things is familiar with a minimal set of common words such as ``if'', ``then'', ``such that'', ``or'', ``and'', and quite a few others in this general category. Then we use certain other special words without defining them, and we establish certain rules for their use.

Now, what am I talking about here? You did this very thing when you took geometry. In that course, you used the words ``point'', ``line'', ``set'', and ``between'', and no effort was made to define them. You could think of them as anything you liked as long as they conformed to some rules that were called postulates or axioms. The words ``rule'', ``axiom'', ``postulate'', and ``assumption'' can all be used interchangeably in mathematics. As an example of such an axiom in geometry, you had the following: If $L$ is a line, $L$ is a pointset. Hence, your visualization of what the words listed above mean would have to fit the axiom. The axiom above eliminates ridiculous things like thinking of a line as a tree and a point as a bran flake. Even if we use a sensible notion of set as a collection, the axiom ``if L is a line, L is a pointset'' is just not satisfied by the conclusion ``L is a tree. Therefore, L is a set of bran flakes''.  Admittedly, this is an absurd application of the idea of being able to think of our undefined terms as anything that we want. But it makes the point that the axiom puts restrictions on such thinking. 

However, consider the following. Suppose ``line'' means ``book'', and ``point'' means ``page''. If we also let ``set'' mean ``collection'', then ``$L$ is a book means $L$ is a collection of pages'' does not violate the axiom. In fact, if the only axiom that we had in geometry was ``if L is a line, L is a pointset'', then it would be perfectly legitimate to think of a line as a book and a point as a page. If your curiosity is aroused, you might look back at the axioms from your geometry course to see if there is at least one axiom there that eliminates books and pages from legitimacy as lines and points. Wouldn't it change your perspective on geometry if there were no such axiom? 

The main point here is that you go back so far in defining words until you are faced with having to use some words without definition. ``Number'' is such a word in mathematics.

Often in defining things in mathematics, we find that the task is to define a category or set of things distinguishing those elements that are in this category from those that are not. Hence, a good definition in mathematics is one which tells us whether or not something we are interested in is in the category that the definition created. For example, the statement ``an elephant is a large animal'' has some serious defects as a definition. If you are examining an elephant, then the definition tells us it is a large animal. But if we are examining a large animal, we do not know, on the basis of this definition, whether or not it is an elephant. In other words, if we are defining a set of objects with some property assigned to them, then we want the definition to enable us to accomplish the following: Given an object we should be able to use the definition to determine that the object \emph{does} belong to the set or that the object \emph{does not} belong to the set.

We may not always be able to construct our definitions of our categories or sets in such a useful way as described above, but if we can, then we should try to do so.

Another aspect of definitions is the difficulty in defining a word by itself. It is much easier to define a whole statement in which the particular word of interest is used. In other words, we define something in the context in which we use it. This is the technique that will be used throughout this course in defining terms. There are other ways of defining things, but it is useful to have one way very clearly in mind.

Note the following definition:

\begin{description}
\item[Definition 1.1:] The statement that $A$ is the arc $PQ$ of the circle $K$ means:
  \begin{enumerate}[\hspace{1cm}1)]
    \item $A$ is a pointset,
    \item each of $P$ and $Q$ is a point of $K$,
    \item a point $x$ belongs to $A$ if and only if
    \begin{enumerate}[\hspace{1cm}(a)]
      \item $x = P$,
      \item $x = Q$, or
      \item if $C$ is the center of $K$, then $x$ is a point of $K$ on the non-$C$-side of line $PQ$.
    \end{enumerate}
  \end{enumerate}
\end{description}

% Figure 1.4

\begin{center}
%TeXCAD Options
%\grade{\on}
%\emlines{\off}
%\epic{\off}
%\beziermacro{\on}
%\reduce{\on}
%\snapping{\off}
%\quality{8.00}
%\graddiff{0.01}
%\snapasp{1}
%\zoom{4.0000}
\unitlength 1mm % = 2.85pt
\linethickness{0.4pt}
\ifx\plotpoint\undefined\newsavebox{\plotpoint}\fi % GNUPLOT compatibility
\begin{picture}(78,62)(0,47)
\leavevmode
%\circle(32.75,81.25){55.44}
\put(60.47,81.25){\line(0,1){1.15}}
\put(60.45,82.4){\line(0,1){1.148}}
\put(60.38,83.55){\line(0,1){1.144}}
\multiput(60.26,84.69)(-.03328,.22759){5}{\line(0,1){.22759}}
\multiput(60.09,85.83)(-.0305,.16144){7}{\line(0,1){.16144}}
\multiput(59.88,86.96)(-.03252,.14003){8}{\line(0,1){.14003}}
\multiput(59.62,88.08)(-.03064,.11085){10}{\line(0,1){.11085}}
\multiput(59.31,89.19)(-.03201,.09953){11}{\line(0,1){.09953}}
\multiput(58.96,90.28)(-.0331,.08994){12}{\line(0,1){.08994}}
\multiput(58.56,91.36)(-.031546,.075848){14}{\line(0,1){.075848}}
\multiput(58.12,92.42)(-.032354,.06951){15}{\line(0,1){.06951}}
\multiput(57.63,93.47)(-.033008,.063851){16}{\line(0,1){.063851}}
\multiput(57.11,94.49)(-.033532,.058755){17}{\line(0,1){.058755}}
\multiput(56.54,95.49)(-.032157,.051281){19}{\line(0,1){.051281}}
\multiput(55.93,96.46)(-.032544,.047408){20}{\line(0,1){.047408}}
\multiput(55.27,97.41)(-.03284,.043826){21}{\line(0,1){.043826}}
\multiput(54.58,98.33)(-.033055,.040498){22}{\line(0,1){.040498}}
\multiput(53.86,99.22)(-.033198,.037392){23}{\line(0,1){.037392}}
\multiput(53.09,100.08)(-.033273,.034484){24}{\line(0,1){.034484}}
\multiput(52.3,100.91)(-.034675,.033074){24}{\line(-1,0){.034675}}
\multiput(51.46,101.7)(-.037583,.032981){23}{\line(-1,0){.037583}}
\multiput(50.6,102.46)(-.040688,.032821){22}{\line(-1,0){.040688}}
\multiput(49.7,103.18)(-.044015,.032587){21}{\line(-1,0){.044015}}
\multiput(48.78,103.87)(-.047595,.03227){20}{\line(-1,0){.047595}}
\multiput(47.83,104.51)(-.054325,.033631){18}{\line(-1,0){.054325}}
\multiput(46.85,105.12)(-.058948,.033193){17}{\line(-1,0){.058948}}
\multiput(45.85,105.68)(-.064041,.032639){16}{\line(-1,0){.064041}}
\multiput(44.82,106.2)(-.069695,.031952){15}{\line(-1,0){.069695}}
\multiput(43.78,106.68)(-.081878,.0335){13}{\line(-1,0){.081878}}
\multiput(42.71,107.12)(-.09013,.03258){12}{\line(-1,0){.09013}}
\multiput(41.63,107.51)(-.09971,.03144){11}{\line(-1,0){.09971}}
\multiput(40.53,107.86)(-.12336,.03333){9}{\line(-1,0){.12336}}
\multiput(39.42,108.16)(-.14022,.03171){8}{\line(-1,0){.14022}}
\multiput(38.3,108.41)(-.16161,.02956){7}{\line(-1,0){.16161}}
\multiput(37.17,108.62)(-.22778,.03197){5}{\line(-1,0){.22778}}
\put(36.03,108.78){\line(-1,0){1.145}}
\put(34.89,108.89){\line(-1,0){1.148}}
\put(33.74,108.95){\line(-1,0){1.15}}
\put(32.59,108.97){\line(-1,0){1.15}}
\put(31.44,108.94){\line(-1,0){1.147}}
\multiput(30.29,108.86)(-.2858,-.0314){4}{\line(-1,0){.2858}}
\multiput(29.15,108.74)(-.1895,-.02883){6}{\line(-1,0){.1895}}
\multiput(28.01,108.56)(-.16126,-.03143){7}{\line(-1,0){.16126}}
\multiput(26.88,108.34)(-.13984,-.03333){8}{\line(-1,0){.13984}}
\multiput(25.77,108.08)(-.11067,-.03128){10}{\line(-1,0){.11067}}
\multiput(24.66,107.76)(-.09934,-.03258){11}{\line(-1,0){.09934}}
\multiput(23.57,107.41)(-.08975,-.03362){12}{\line(-1,0){.08975}}
\multiput(22.49,107)(-.075665,-.031983){14}{\line(-1,0){.075665}}
\multiput(21.43,106.56)(-.069322,-.032754){15}{\line(-1,0){.069322}}
\multiput(20.39,106.06)(-.06366,-.033376){16}{\line(-1,0){.06366}}
\multiput(19.37,105.53)(-.055307,-.031989){18}{\line(-1,0){.055307}}
\multiput(18.38,104.95)(-.051094,-.032453){19}{\line(-1,0){.051094}}
\multiput(17.4,104.34)(-.047219,-.032817){20}{\line(-1,0){.047219}}
\multiput(16.46,103.68)(-.043636,-.033092){21}{\line(-1,0){.043636}}
\multiput(15.54,102.99)(-.040306,-.033289){22}{\line(-1,0){.040306}}
\multiput(14.66,102.25)(-.0372,-.033413){23}{\line(-1,0){.0372}}
\multiput(13.8,101.49)(-.034291,-.033472){24}{\line(-1,0){.034291}}
\multiput(12.98,100.68)(-.032873,-.034865){24}{\line(0,-1){.034865}}
\multiput(12.19,99.85)(-.032764,-.037773){23}{\line(0,-1){.037773}}
\multiput(11.44,98.98)(-.032586,-.040876){22}{\line(0,-1){.040876}}
\multiput(10.72,98.08)(-.032332,-.044202){21}{\line(0,-1){.044202}}
\multiput(10.04,97.15)(-.033678,-.050295){19}{\line(0,-1){.050295}}
\multiput(9.4,96.19)(-.033317,-.054518){18}{\line(0,-1){.054518}}
\multiput(8.8,95.21)(-.032852,-.059138){17}{\line(0,-1){.059138}}
\multiput(8.24,94.21)(-.032269,-.064228){16}{\line(0,-1){.064228}}
\multiput(7.73,93.18)(-.031549,-.069879){15}{\line(0,-1){.069879}}
\multiput(7.25,92.13)(-.033027,-.08207){13}{\line(0,-1){.08207}}
\multiput(6.82,91.06)(-.03206,-.09032){12}{\line(0,-1){.09032}}
\multiput(6.44,89.98)(-.03086,-.09989){11}{\line(0,-1){.09989}}
\multiput(6.1,88.88)(-.03262,-.12355){9}{\line(0,-1){.12355}}
\multiput(5.81,87.77)(-.0309,-.1404){8}{\line(0,-1){.1404}}
\multiput(5.56,86.65)(-.0334,-.18874){6}{\line(0,-1){.18874}}
\multiput(5.36,85.51)(-.03065,-.22796){5}{\line(0,-1){.22796}}
\put(5.2,84.37){\line(0,-1){1.145}}
\put(5.1,83.23){\line(0,-1){1.149}}
\put(5.04,82.08){\line(0,-1){1.15}}
\put(5.03,80.93){\line(0,-1){1.149}}
\put(5.07,79.78){\line(0,-1){1.147}}
\multiput(5.15,78.63)(.0331,-.2856){4}{\line(0,-1){.2856}}
\multiput(5.28,77.49)(.02992,-.18933){6}{\line(0,-1){.18933}}
\multiput(5.46,76.36)(.03236,-.16108){7}{\line(0,-1){.16108}}
\multiput(5.69,75.23)(.03034,-.12413){9}{\line(0,-1){.12413}}
\multiput(5.96,74.11)(.03192,-.11049){10}{\line(0,-1){.11049}}
\multiput(6.28,73.01)(.03316,-.09915){11}{\line(0,-1){.09915}}
\multiput(6.65,71.92)(.031511,-.082664){13}{\line(0,-1){.082664}}
\multiput(7.06,70.84)(.032419,-.075479){14}{\line(0,-1){.075479}}
\multiput(7.51,69.78)(.033154,-.069132){15}{\line(0,-1){.069132}}
\multiput(8.01,68.75)(.031758,-.059733){17}{\line(0,-1){.059733}}
\multiput(8.55,67.73)(.032308,-.055122){18}{\line(0,-1){.055122}}
\multiput(9.13,66.74)(.032747,-.050906){19}{\line(0,-1){.050906}}
\multiput(9.75,65.77)(.033089,-.047029){20}{\line(0,-1){.047029}}
\multiput(10.41,64.83)(.033344,-.043444){21}{\line(0,-1){.043444}}
\multiput(11.11,63.92)(.033521,-.040113){22}{\line(0,-1){.040113}}
\multiput(11.85,63.04)(.033627,-.037006){23}{\line(0,-1){.037006}}
\multiput(12.62,62.19)(.033669,-.034097){24}{\line(0,-1){.034097}}
\multiput(13.43,61.37)(.035054,-.032671){24}{\line(1,0){.035054}}
\multiput(14.27,60.58)(.037961,-.032545){23}{\line(1,0){.037961}}
\multiput(15.15,59.83)(.041064,-.032349){22}{\line(1,0){.041064}}
\multiput(16.05,59.12)(.046607,-.03368){20}{\line(1,0){.046607}}
\multiput(16.98,58.45)(.050488,-.033388){19}{\line(1,0){.050488}}
\multiput(17.94,57.81)(.054709,-.033002){18}{\line(1,0){.054709}}
\multiput(18.93,57.22)(.059327,-.03251){17}{\line(1,0){.059327}}
\multiput(19.94,56.67)(.064413,-.031898){16}{\line(1,0){.064413}}
\multiput(20.97,56.16)(.075064,-.03337){14}{\line(1,0){.075064}}
\multiput(22.02,55.69)(.082259,-.032553){13}{\line(1,0){.082259}}
\multiput(23.09,55.27)(.0905,-.03154){12}{\line(1,0){.0905}}
\multiput(24.17,54.89)(.11008,-.03331){10}{\line(1,0){.11008}}
\multiput(25.27,54.56)(.12374,-.03191){9}{\line(1,0){.12374}}
\multiput(26.39,54.27)(.14057,-.03009){8}{\line(1,0){.14057}}
\multiput(27.51,54.03)(.18893,-.03231){6}{\line(1,0){.18893}}
\multiput(28.64,53.83)(.22813,-.02934){5}{\line(1,0){.22813}}
\put(29.79,53.69){\line(1,0){1.146}}
\put(30.93,53.59){\line(1,0){1.149}}
\put(32.08,53.54){\line(1,0){1.15}}
\put(33.23,53.53){\line(1,0){1.149}}
\put(34.38,53.58){\line(1,0){1.146}}
\multiput(35.53,53.67)(.22833,.02777){5}{\line(1,0){.22833}}
\multiput(36.67,53.81)(.18915,.03101){6}{\line(1,0){.18915}}
\multiput(37.8,53.99)(.16089,.03329){7}{\line(1,0){.16089}}
\multiput(38.93,54.23)(.12395,.03106){9}{\line(1,0){.12395}}
\multiput(40.04,54.5)(.1103,.03255){10}{\line(1,0){.1103}}
\multiput(41.15,54.83)(.09896,.03373){11}{\line(1,0){.09896}}
\multiput(42.24,55.2)(.08248,.031988){13}{\line(1,0){.08248}}
\multiput(43.31,55.62)(.075291,.032854){14}{\line(1,0){.075291}}
\multiput(44.36,56.08)(.068939,.033552){15}{\line(1,0){.068939}}
\multiput(45.4,56.58)(.059549,.032102){17}{\line(1,0){.059549}}
\multiput(46.41,57.13)(.054934,.032625){18}{\line(1,0){.054934}}
\multiput(47.4,57.71)(.050716,.03304){19}{\line(1,0){.050716}}
\multiput(48.36,58.34)(.046837,.033359){20}{\line(1,0){.046837}}
\multiput(49.3,59.01)(.043251,.033594){21}{\line(1,0){.043251}}
\multiput(50.21,59.71)(.038184,.032284){23}{\line(1,0){.038184}}
\multiput(51.08,60.46)(.035278,.03243){24}{\line(1,0){.035278}}
\multiput(51.93,61.23)(.032546,.032511){25}{\line(1,0){.032546}}
\multiput(52.74,62.05)(.032469,.035242){24}{\line(0,1){.035242}}
\multiput(53.52,62.89)(.032326,.038148){23}{\line(0,1){.038148}}
\multiput(54.27,63.77)(.033641,.043214){21}{\line(0,1){.043214}}
\multiput(54.97,64.68)(.033411,.046801){20}{\line(0,1){.046801}}
\multiput(55.64,65.61)(.033096,.05068){19}{\line(0,1){.05068}}
\multiput(56.27,66.58)(.032685,.054899){18}{\line(0,1){.054899}}
\multiput(56.86,67.57)(.032167,.059513){17}{\line(0,1){.059513}}
\multiput(57.41,68.58)(.033627,.068903){15}{\line(0,1){.068903}}
\multiput(57.91,69.61)(.032936,.075255){14}{\line(0,1){.075255}}
\multiput(58.37,70.66)(.032078,.082445){13}{\line(0,1){.082445}}
\multiput(58.79,71.74)(.03102,.09068){12}{\line(0,1){.09068}}
\multiput(59.16,72.82)(.03267,.11027){10}{\line(0,1){.11027}}
\multiput(59.49,73.93)(.03119,.12392){9}{\line(0,1){.12392}}
\multiput(59.77,75.04)(.03346,.16085){7}{\line(0,1){.16085}}
\multiput(60,76.17)(.03122,.18912){6}{\line(0,1){.18912}}
\multiput(60.19,77.3)(.02802,.2283){5}{\line(0,1){.2283}}
\put(60.33,78.44){\line(0,1){1.146}}
\put(60.42,79.59){\line(0,1){1.659}}
%\end
\put(5.25,84.75){\line(0,1){0}} \put(5.25,84.75){\line(0,1){.25}}
\put(3.25,87.5){\makebox(0,0)[cc]{$P$}}
\put(34.75,79.25){\makebox(0,0)[cc]{$C$}}
\put(32.75,81.25){\circle{.71}}
%\emline(2.5,92.5)(78,84.5)
\multiput(2.5,92.5)(.31722689,-.03361345){238}{\line(1,0){.31722689}}
%\end
\put(63,89.5){\makebox(0,0)[cc]{$Q$}}
\put(52,101){\circle{.5}}
\put(55,102.5){\makebox(0,0)[cc]{$x$}}
%\qbezvec(69,103.5)(61.63,109)(55.75,101.5)
\put(50.75,106.5){\vector(-3,-4){.07}}\qbezier(64,108.5)(64.63,119)(50.75,106.5)
%\end
\put(67.25,106.25){\makebox(0,0)[cc]{$K$}}
\put(32,49.25){\makebox(0,0)[cc]{Figure 1.4}}
\end{picture}
\end{center}

It will pay you to study this definition carefully by asking yourself whether or not it accomplishes what we described previously as what we want a definition to accomplish. This method of definition will be used extensively throughout the course in order for you to get enough exposure to it so that you will begin to feel comfortable with it. You will be asked to define certain entities yourself, just as you were asked to define the arc of a circle. It is not so important that you reproduced exactly what was presented above as the definition of an arc. But it is important that you made an effort to define an arc. Comparing your definition with the one above can be very instructive to you, and, if you will continue to participate in this manner (attempting the definition before you look at the one presented in the text), you will find it very rewarding in terms of your understanding of the material.

Now, here is a concept from geometry that is very useful.

\begin{description}
\item[Definition 1.2:] The statement that the point $P$ is on the $A$-side of line $L$ means:
  \begin{enumerate}[\hspace{1cm}1)]
    \item $A$ is a point not on $L$,
    \item $P$ is a point not on $L$, and
    \item no point of $L$ is between $P$ and $A$.
  \end{enumerate}
\end{description}

Draw a picture to help see this, and then define what you would mean by the statement that the point $P$ is on the non-$A$-side of line $L$.

You may have noted a slight problem with our definition of arc of a circle. If $C$ is the center of the circle $K$, $P$ and $Q$ are points of $K$, and $P$, $C$, and $Q$ are colinear, then our definition won't work. Draw a picture of this condition.

You'll notice, if you've drawn a picture of the situation described above, that there is an arc that we could talk about, but it is a semi-circle. But which one? Before going on further, try to define the statement ``$A$ is the semi-circle $PRQ$ of the circle $K$.''

\begin{description}

\item[Definition 1.3:] The statement that $A$ is the semi-circle $PRQ$ of the circle $K$ means:
  \begin{enumerate}[\hspace{1cm}1)]
    \item $A$ is a pointset,
    \item $P$, $R$, and $Q$ are distinct points of $K$,
    \item if $C$ is the center of $K$, then $PCQ$ is true, and
    \item a point $x$ belongs to $A$ if and only if
    \begin{enumerate}[\hspace{1cm}(a)]
      \item $x$ is a point of arc $PR$, or
      \item $x$ is a point of arc $RQ$.
    \end{enumerate}
  \end{enumerate}

\item[Question:] Can you see how to use the notion of ``side of a line'' to define a semi-circle? Here is another definition from geometry.

\item[Definition 1.4:] The statement that $ABC$ is true, or $ABC$, means:
  \begin{enumerate}[\hspace{1cm}1)]
    \item $A$, $B$, and $C$ are distinct points,
    \item some line contains all of $A$, $B$, and $C$, and
    \item $B$ is between $A$ and $C$.
  \end{enumerate}

\end{description}

I cannot suggest too strongly that you draw pictures to see what the definitions mean.

Since we will be talking a great deal about lines and circles in this course, it will be useful to have some notation to use to distinguish between pointsets that are subsets of circles, and pointsets that are subsets of lines. Consequently, we adopt the following convention. When speaking of two points $P$ and $Q$ as points of a line, we represent the set which consists of $P$, $Q$, and the points between them, by the symbol $[P,Q]$, called the interval $PQ$, and the length of that interval by the symbol $PQ$.

When speaking of two points $P$ and $Q$ as points of a circle, we don't use the term ``between'' anymore, but we speak of the arc $PQ$, which we represent by the symbol $\arc{PQ}$, and we represent the length of $\arc{PQ}$ by the symbol $l(\arc{PQ})$.

Since $PQR$ means $P$, $Q$, and $R$ are distinct points of a line, and $Q$ is between $P$ and $R$, let's adopt the following convention: $\arc{PQR}$ means $P$, $Q$, and $R$ are distinct points of a circle, and $Q$ is a point of arc $\arc{PR}$. Since it is useful to refer to semi-circles as arcs, the semi-circle $PQR$ will be denoted $\arc{PQR}$.

\begin{description}
\item[Question:] Why not refer to the semi-circle $\arc{PQR}$ as just the semi-circle $\arc{PR}$?
\end{description}

Now, back to our main problem. We are concerned with the notion of length of an arc of a circle. Hence, we must establish that each arc of a circle does indeed have length. 

This may seem like a strange thing to have to do, but someone, or some thing, may have some mental image of an arc of a circle that is so weird that, as far as he is concerned, if an arc of a circle is really like he imagines it, then arcs of circles just don't have length. What we want to do is to require him to limit his notion of arc of a circle so that it is sensible to talk about arc length. 

We do this by simply imposing the notion of length on arcs of circles; that is, we make the assumption that circular arcs have length. You will notice that we have no definition of arc length so this is one of the terms that we use in this course without definition. In a later course, perhaps we will be better able to define it.

You'll notice also that we \emph{do} have a definition of length of an interval.

\begin{description}

\item[Axiom 1:] If each of $P$ and $Q$ is a point of the circle $K$, then there is one and only one non-negative number $x$ which is the length of the arc $\arc{PQ}$.

  Furthermore, if $\arc{PQR}$ is a semi-circle of $K$, then there is one and only one non-negative number $x$ which is the length of $\arc{PQR}$.

\item[Notation:] If each of $P$ and $Q$ is a point of the circle $K$, then the length of arc $\arc{PQ}$ is denoted by the symbol $l(\arc{PQ})$.

  Furthermore, if $\arc{PQR}$ is a semi-circle of $K$, then the length of $\arc{PQR}$ is denoted $l(\arc{PQR})$.

\end{description}

Note that this axiom and the notation formalize what was discussed rather informally earlier.

Note again that this does not define arc length, but it does tell us that arc length is a number.

\vspace{0.3cm}

There are some things from algebra and geometry that are very useful to us in suggesting tools for us to use in this course. For example, we used the notation $ABC$ from geometry to mean $A$, $B$, and $C$ are distinct points of a line, and $B$ is between $A$ and $C$. This suggested that we define $\arc{ABC}$ to mean $A$, $B$, and $C$ are distinct points of a circle, and $B$ is a point of arc $\arc{AC}$. This is not inconsistent with the way we defined a semi-circle, but it means that, if $\arc{ABC}$ is a semi-circle, then we must state that it is.

\begin{description}
\item[Definition 1.5:] $\arc{ABC}$ means:
  \begin{enumerate}[\hspace{1cm}1)]
    \item $A$, $B$, and $C$ are distinct points of a circle, and
    \item $B$ is a point of arc $\arc{AC}$ of that circle. You'll note that $\arc{ABC}$ means that $A$, $B$, and $C$ are contained in a semi-circle of the circle.
  \end{enumerate}
\end{description}

There is also an axiom from geometry that may be useful to us.

\begin{description}
\item[Axiom 2:] If $A$, $B$, and $C$ are distinct points of a line, then one and only one of the following is true:
  \begin{enumerate}[\hspace{1cm}1)]
    \item $ABC$ 
    \item $BAC$
    \item $ACB$
  \end{enumerate}
\end{description}

Your Geometry course may not have had such an axiom stated explicitly. But if so, I'll bet you a dollar to a donut hole that your course used this axiom tacitly. (There's another good word you might want to look up if you aren't comfortable with it.)

See if you can come up with something analogous to this for $A$, $B$, and $C$ being distinct points of a circle. Do this \emph{before} going on further. Draw some pictures because there may be a trap lurking here.

\begin{description}
\item[Axiom 3:] If $A$, $B$, and $C$ are distinct points of a semi-circle of the circle $K$, then one and only one of the following is true:
  \begin{enumerate}[\hspace{1cm}1)]
    \item $\arc{ABC}$ 
    \item $\arc{BAC}$
    \item $\arc{ACB}$
  \end{enumerate}
\end{description}

Do you see why the semi-circle is included in Axiom 3?

Recall from algebra the following, which you may have used as an axiom.

\begin{description}
\item[Axiom 4:] The following two statements are equivalent:
  \begin{enumerate}[\hspace{1cm}1)]
    \item $ABC$
    \item $AB + BC = AC$
  \end{enumerate}
\end{description}

Can you come up with a statement analogous to Axiom 4 from algebra that we could apply to our circles? Try to do so before you go on.

\begin{description}
\item[Axiom 5:] The following two statements are equivalent:
  \begin{enumerate}[\hspace{1cm}1)]
    \item $\arc{ABC}$
    \item $l(\arc{AB}) + l(\arc{BC}) = l(\arc{AC})$
  \end{enumerate}
\end{description}

Since $ABC$ means $A$, $B$, and $C$ are colinear and $B$ is between $A$ and $C$, what would $ABCD$ mean? Well, like you would expect, it means $ABC$, $ABD$, $ACD$, and $BCD$. 

Can you show that $ABCD$ also means that $AB + BC + CD = AD$?

Similarly, $\arc{PQR}$ means $Q$ is a point of arc $\arc{PR}$, so $\arc{PQRS}$ means $\arc{PQR}$, $\arc{PQS}$, $\arc{PRS}$, and $\arc{QRS}$. And I'll bet you can show that $\arc{PQRS}$ implies that $l(\arc{PQ}) + l(\arc{QR}) + l(\arc{RS}) = l(\arc{PS})$. 

What if we have $ABD$ and $ACD$? What can you conclude if $B \neq C$? Can you see that either $ABCD$ or $ACBD$? We'll take these ideas for granted in our struggles in this course. 

Now, let's examine Axiom 4 some more. Axiom 4 means that, if either of statement 1 or statement 2 is true, then the other is true. This axiom is strong enough to enable us to prove that the following two statements are equivalent:
\begin{enumerate}[\hspace{1cm}1)]
  \item $L: Ax + By + C = 0$, not both of $A$ and $B$ are $0$.
  \item $L$ is a line.
\end{enumerate}

This means that, if $L$ is a pointset with equation $Ax + By + C = 0$, where not both of $A$ and $B$ are $0$, then $L$ is a line. It also means that, if $L$ is a line, then there exists a number triple $A,B,C$ such that $L$ has equation $Ax + By + C = 0$. You may have either proved this or at least used it in your last algebra course.

Now, so far, we have established with our Axioms 1, 3, and 5 that circular arcs have length, and that arc length is additive. We already know how to find chord length if we have the coordinates of the end points. We simply use the definition of distance between two points. 

What we need now is some obvious or intuitive relationship between arcs and chords that we can use to solve our main problem; i.e., find arc length if we are given the chord length. Of course, if we could define distance along an arc of a circle, we would have solved our problem. But you have probably observed that this is not an easy task for us now. Hence, we are looking for some other route to solve our problem of relating arc length to chord length.

Consider the following picture:

\begin{center}
%TeXCAD Picture [fig1.5.pic]. Options:
%\grade{\on}
%\emlines{\off}
%\epic{\off}
%\beziermacro{\on}
%\reduce{\on}
%\snapping{\off}
%\quality{8.00}
%\graddiff{0.01}
%\snapasp{1}
%\zoom{4.0000}
\unitlength 0.8mm % = 2.85pt
\linethickness{0.4pt}
\ifx\plotpoint\undefined\newsavebox{\plotpoint}\fi % GNUPLOT compatibility
\begin{picture}(59.34,64.84)(0,0)
%\circle(32.25,37.75){54.19}
\put(59.34,37.75){\line(0,1){1.133}}
\put(59.32,38.88){\line(0,1){1.131}}
\put(59.25,40.01){\line(0,1){1.127}}
\multiput(59.13,41.14)(-.03307,.22426){5}{\line(0,1){.22426}}
\multiput(58.97,42.26)(-.0303,.15906){7}{\line(0,1){.15906}}
\multiput(58.75,43.38)(-.03231,.13795){8}{\line(0,1){.13795}}
\multiput(58.5,44.48)(-.03044,.10918){10}{\line(0,1){.10918}}
\multiput(58.19,45.57)(-.0318,.09801){11}{\line(0,1){.09801}}
\multiput(57.84,46.65)(-.03288,.08854){12}{\line(0,1){.08854}}
\multiput(57.45,47.71)(-.031336,.074649){14}{\line(0,1){.074649}}
\multiput(57.01,48.76)(-.032135,.068389){15}{\line(0,1){.068389}}
\multiput(56.53,49.78)(-.032781,.062798){16}{\line(0,1){.062798}}
\multiput(56,50.79)(-.033298,.057762){17}{\line(0,1){.057762}}
\multiput(55.44,51.77)(-.033702,.05319){18}{\line(0,1){.05319}}
\multiput(54.83,52.73)(-.032308,.046561){20}{\line(0,1){.046561}}
\multiput(54.18,53.66)(-.032597,.043018){21}{\line(0,1){.043018}}
\multiput(53.5,54.56)(-.032805,.039725){22}{\line(0,1){.039725}}
\multiput(52.78,55.44)(-.032941,.036653){23}{\line(0,1){.036653}}
\multiput(52.02,56.28)(-.03301,.033774){24}{\line(0,1){.033774}}
\multiput(51.23,57.09)(-.034393,.032364){24}{\line(-1,0){.034393}}
\multiput(50.4,57.87)(-.038964,.033706){22}{\line(-1,0){.038964}}
\multiput(49.54,58.61)(-.04226,.033573){21}{\line(-1,0){.04226}}
\multiput(48.66,59.31)(-.045809,.033365){20}{\line(-1,0){.045809}}
\multiput(47.74,59.98)(-.049646,.033074){19}{\line(-1,0){.049646}}
\multiput(46.8,60.61)(-.053819,.032689){18}{\line(-1,0){.053819}}
\multiput(45.83,61.2)(-.058382,.032199){17}{\line(-1,0){.058382}}
\multiput(44.84,61.74)(-.067635,.033693){15}{\line(-1,0){.067635}}
\multiput(43.82,62.25)(-.073912,.033037){14}{\line(-1,0){.073912}}
\multiput(42.79,62.71)(-.081016,.032218){13}{\line(-1,0){.081016}}
\multiput(41.73,63.13)(-.08915,.0312){12}{\line(-1,0){.08915}}
\multiput(40.66,63.51)(-.10845,.03294){10}{\line(-1,0){.10845}}
\multiput(39.58,63.83)(-.12193,.03152){9}{\line(-1,0){.12193}}
\multiput(38.48,64.12)(-.13853,.0297){8}{\line(-1,0){.13853}}
\multiput(37.37,64.36)(-.1862,.03183){6}{\line(-1,0){.1862}}
\multiput(36.26,64.55)(-.22485,.02882){5}{\line(-1,0){.22485}}
\put(35.13,64.69){\line(-1,0){1.129}}
\put(34,64.79){\line(-1,0){1.132}}
\put(32.87,64.84){\line(-1,0){1.133}}
\put(31.74,64.84){\line(-1,0){1.133}}
\put(30.6,64.79){\line(-1,0){1.13}}
\multiput(29.47,64.7)(-.22496,-.02793){5}{\line(-1,0){.22496}}
\multiput(28.35,64.56)(-.18633,-.03109){6}{\line(-1,0){.18633}}
\multiput(27.23,64.38)(-.15846,-.03331){7}{\line(-1,0){.15846}}
\multiput(26.12,64.14)(-.12205,-.03104){9}{\line(-1,0){.12205}}
\multiput(25.02,63.86)(-.10858,-.0325){10}{\line(-1,0){.10858}}
\multiput(23.94,63.54)(-.09739,-.03365){11}{\line(-1,0){.09739}}
\multiput(22.87,63.17)(-.081143,-.031896){13}{\line(-1,0){.081143}}
\multiput(21.81,62.75)(-.074043,-.032743){14}{\line(-1,0){.074043}}
\multiput(20.78,62.3)(-.067768,-.033424){15}{\line(-1,0){.067768}}
\multiput(19.76,61.79)(-.05851,-.031966){17}{\line(-1,0){.05851}}
\multiput(18.76,61.25)(-.053948,-.032475){18}{\line(-1,0){.053948}}
\multiput(17.79,60.67)(-.049777,-.032877){19}{\line(-1,0){.049777}}
\multiput(16.85,60.04)(-.045941,-.033183){20}{\line(-1,0){.045941}}
\multiput(15.93,59.38)(-.042393,-.033405){21}{\line(-1,0){.042393}}
\multiput(15.04,58.68)(-.039097,-.033551){22}{\line(-1,0){.039097}}
\multiput(14.18,57.94)(-.036022,-.033629){23}{\line(-1,0){.036022}}
\multiput(13.35,57.16)(-.033143,-.033643){24}{\line(0,-1){.033643}}
\multiput(12.55,56.36)(-.033086,-.036521){23}{\line(0,-1){.036521}}
\multiput(11.79,55.52)(-.032963,-.039595){22}{\line(0,-1){.039595}}
\multiput(11.07,54.65)(-.032767,-.042888){21}{\line(0,-1){.042888}}
\multiput(10.38,53.75)(-.032492,-.046432){20}{\line(0,-1){.046432}}
\multiput(9.73,52.82)(-.032128,-.050263){19}{\line(0,-1){.050263}}
\multiput(9.12,51.86)(-.033527,-.057629){17}{\line(0,-1){.057629}}
\multiput(8.55,50.88)(-.033031,-.062668){16}{\line(0,-1){.062668}}
\multiput(8.02,49.88)(-.032406,-.06826){15}{\line(0,-1){.06826}}
\multiput(7.54,48.86)(-.031632,-.074524){14}{\line(0,-1){.074524}}
\multiput(7.09,47.81)(-.03324,-.08841){12}{\line(0,-1){.08841}}
\multiput(6.69,46.75)(-.03219,-.09788){11}{\line(0,-1){.09788}}
\multiput(6.34,45.67)(-.03088,-.10906){10}{\line(0,-1){.10906}}
\multiput(6.03,44.58)(-.03286,-.13782){8}{\line(0,-1){.13782}}
\multiput(5.77,43.48)(-.03093,-.15894){7}{\line(0,-1){.15894}}
\multiput(5.55,42.37)(-.0283,-.18677){6}{\line(0,-1){.18677}}
\multiput(5.38,41.25)(-.0307,-.2817){4}{\line(0,-1){.2817}}
\put(5.26,40.12){\line(0,-1){1.131}}
\put(5.18,38.99){\line(0,-1){3.398}}
\put(5.24,35.59){\line(0,-1){1.128}}
\multiput(5.35,34.47)(.03218,-.22439){5}{\line(0,-1){.22439}}
\multiput(5.52,33.34)(.02967,-.15918){7}{\line(0,-1){.15918}}
\multiput(5.72,32.23)(.03176,-.13807){8}{\line(0,-1){.13807}}
\multiput(5.98,31.12)(.03334,-.12144){9}{\line(0,-1){.12144}}
\multiput(6.28,30.03)(.03141,-.09813){11}{\line(0,-1){.09813}}
\multiput(6.62,28.95)(.03253,-.08867){12}{\line(0,-1){.08867}}
\multiput(7.01,27.89)(.033426,-.080525){13}{\line(0,-1){.080525}}
\multiput(7.45,26.84)(.031863,-.068516){15}{\line(0,-1){.068516}}
\multiput(7.93,25.81)(.032532,-.062928){16}{\line(0,-1){.062928}}
\multiput(8.45,24.81)(.033068,-.057894){17}{\line(0,-1){.057894}}
\multiput(9.01,23.82)(.033491,-.053324){18}{\line(0,-1){.053324}}
\multiput(9.61,22.86)(.032122,-.046689){20}{\line(0,-1){.046689}}
\multiput(10.25,21.93)(.032426,-.043147){21}{\line(0,-1){.043147}}
\multiput(10.93,21.02)(.032647,-.039855){22}{\line(0,-1){.039855}}
\multiput(11.65,20.15)(.032795,-.036783){23}{\line(0,-1){.036783}}
\multiput(12.41,19.3)(.032875,-.033905){24}{\line(0,-1){.033905}}
\multiput(13.2,18.49)(.034264,-.032501){24}{\line(1,0){.034264}}
\multiput(14.02,17.71)(.037141,-.032389){23}{\line(1,0){.037141}}
\multiput(14.87,16.96)(.040212,-.032207){22}{\line(1,0){.040212}}
\multiput(15.76,16.25)(.045676,-.033547){20}{\line(1,0){.045676}}
\multiput(16.67,15.58)(.049514,-.033271){19}{\line(1,0){.049514}}
\multiput(17.61,14.95)(.053688,-.032903){18}{\line(1,0){.053688}}
\multiput(18.58,14.36)(.058254,-.03243){17}{\line(1,0){.058254}}
\multiput(19.57,13.81)(.063282,-.031839){16}{\line(1,0){.063282}}
\multiput(20.58,13.3)(.07378,-.03333){14}{\line(1,0){.07378}}
\multiput(21.61,12.83)(.080887,-.03254){13}{\line(1,0){.080887}}
\multiput(22.67,12.41)(.08903,-.03156){12}{\line(1,0){.08903}}
\multiput(23.73,12.03)(.10832,-.03337){10}{\line(1,0){.10832}}
\multiput(24.82,11.69)(.1218,-.03201){9}{\line(1,0){.1218}}
\multiput(25.91,11.41)(.13841,-.03025){8}{\line(1,0){.13841}}
\multiput(27.02,11.16)(.18608,-.03257){6}{\line(1,0){.18608}}
\multiput(28.14,10.97)(.22473,-.02972){5}{\line(1,0){.22473}}
\put(29.26,10.82){\line(1,0){1.129}}
\put(30.39,10.72){\line(1,0){1.132}}
\put(31.52,10.66){\line(1,0){1.133}}
\put(32.66,10.66){\line(1,0){1.133}}
\put(33.79,10.7){\line(1,0){1.13}}
\multiput(34.92,10.79)(.22507,.02703){5}{\line(1,0){.22507}}
\multiput(36.04,10.92)(.18645,.03035){6}{\line(1,0){.18645}}
\multiput(37.16,11.1)(.15859,.03268){7}{\line(1,0){.15859}}
\multiput(38.27,11.33)(.12217,.03055){9}{\line(1,0){.12217}}
\multiput(39.37,11.61)(.10871,.03207){10}{\line(1,0){.10871}}
\multiput(40.46,11.93)(.09752,.03326){11}{\line(1,0){.09752}}
\multiput(41.53,12.29)(.08127,.031573){13}{\line(1,0){.08127}}
\multiput(42.59,12.7)(.074172,.032448){14}{\line(1,0){.074172}}
\multiput(43.63,13.16)(.0679,.033154){15}{\line(1,0){.0679}}
\multiput(44.64,13.66)(.062301,.033717){16}{\line(1,0){.062301}}
\multiput(45.64,14.2)(.054077,.032261){18}{\line(1,0){.054077}}
\multiput(46.62,14.78)(.049908,.032678){19}{\line(1,0){.049908}}
\multiput(47.56,15.4)(.046072,.033){20}{\line(1,0){.046072}}
\multiput(48.48,16.06)(.042526,.033236){21}{\line(1,0){.042526}}
\multiput(49.38,16.76)(.03923,.033396){22}{\line(1,0){.03923}}
\multiput(50.24,17.49)(.036156,.033485){23}{\line(1,0){.036156}}
\multiput(51.07,18.26)(.033277,.033511){24}{\line(0,1){.033511}}
\multiput(51.87,19.06)(.033231,.03639){23}{\line(0,1){.03639}}
\multiput(52.64,19.9)(.03312,.039463){22}{\line(0,1){.039463}}
\multiput(53.36,20.77)(.032938,.042757){21}{\line(0,1){.042757}}
\multiput(54.06,21.67)(.032677,.046303){20}{\line(0,1){.046303}}
\multiput(54.71,22.59)(.032328,.050135){19}{\line(0,1){.050135}}
\multiput(55.32,23.55)(.031881,.054302){18}{\line(0,1){.054302}}
\multiput(55.9,24.52)(.03328,.062536){16}{\line(0,1){.062536}}
\multiput(56.43,25.52)(.032677,.068131){15}{\line(0,1){.068131}}
\multiput(56.92,26.55)(.031928,.074398){14}{\line(0,1){.074398}}
\multiput(57.37,27.59)(.03359,.08828){12}{\line(0,1){.08828}}
\multiput(57.77,28.65)(.03258,.09775){11}{\line(0,1){.09775}}
\multiput(58.13,29.72)(.03131,.10893){10}{\line(0,1){.10893}}
\multiput(58.44,30.81)(.03341,.13768){8}{\line(0,1){.13768}}
\multiput(58.71,31.91)(.03157,.15881){7}{\line(0,1){.15881}}
\multiput(58.93,33.02)(.02905,.18666){6}{\line(0,1){.18666}}
\multiput(59.1,34.14)(.0318,.2816){4}{\line(0,1){.2816}}
\put(59.23,35.27){\line(0,1){1.131}}
\put(59.31,36.4){\line(0,1){1.348}}
%\end
%\emline(14,18)(49.25,17)
\multiput(14,18)(1.175,-.033333){30}{\line(1,0){1.175}}
%\end
%\emline(10.75,54.5)(55.25,52.5)
\multiput(10.75,54.5)(.7416667,-.0333333){60}{\line(1,0){.7416667}}
%\end
\put(8.25,57.25){\makebox(0,0)[cc]{$P$}}
\put(58.25,53.5){\makebox(0,0)[cc]{$Q$}}
\put(11,16){\makebox(0,0)[cc]{$R$}}
\put(53.25,15.75){\makebox(0,0)[cc]{$S$}}
\put(32,37.75){\circle{.5}}
\put(31.5,5.25){\makebox(0,0)[cc]{Figure 1.5}}
\end{picture}
\end{center}

Suppose $l(\arc{PQ}) = l(\arc{RS})$. Do you believe it follows that $PQ = RS$?

Suppose $PQ = RS$. Do you believe it follows that $l(\arc{PQ}) = l(\arc{RS})$?

\begin{description}
\item[Axiom 6:] If each of $P$, $Q$, $R$, and $S$ is a point of the circle $K$, then the following two statements are equivalent:
  \begin{enumerate}[\hspace{1cm}1)]
    \item $PQ=RS$
    \item $l(\arc{PQ}) = l(\arc{RS})$
  \end{enumerate}
\end{description}

It will be useful to elaborate what it means to say that two statements are equivalent.

\begin{description}
\item[Definition 1.6:] The statement that Statement $A$ and Statement $B$ are equivalent means:
  \begin{enumerate}[\hspace{1cm}1)]
    \item if $A$ is true, then $B$ is true, and
    \item if $B$ is true, then $A$ is true.
  \end{enumerate}
\end{description}

Hence, Axiom 6 states that, if the chord lengths are the same, then the arc lengths are the same, and it also states that, if the arc lengths are the same, then the chord lengths are the same.

There is one more property from geometry that we can rely on and is probably familiar to you -- namely, the Pythagorean Theorem. Here is a statement of the Pythagorean Theorem in a form that we will find convenient, and we will use it as an axiom.

\begin{description}

\item[Axiom 7:] If $A$, $B$, and $C$ are non-colinear points and line $AB$ is perpendicular to line $BC$, then $(AB)^{2} + (BC)^{2} = (AC)^{2}$.

\item[Notation:] Sometimes it is convenient to denote ``perpendicular to'' by the symbol ``$\bot$''. 

\end{description}

These axioms may seem rather obvious to you. Certainly, looking at the pictures, you might conclude that they are obvious. If these things are obvious, why do we have to make such a big deal out of them? Why do we have to state them as assumptions? Let me try to give you a little insight into that question. 

Recall that, in your geometry course, you were not allowed to make your arguments picture-dependent. No proof that you gave was valid if it relied on the picture. Why was that, anyway? Well, imagine a race of beings that do not have eyes as we do. However, they are intelligent, can hear and understand our speech, and can reason as well as we. If we were challenged to explain our mathematics to them, everything we presented to them would have to be verbal. No pictures on the blackboard or on paper would be of any use. This would mean that every argument for every theorem would have to be crafted verbally, with no reference to pictures. What would be obvious to them? If we are going to communicate with our alien friends and we are going to use something from a picture that seems obvious to us, then we must state it as an axiom (assumption) or we must find some way of showing that it follows logically from what has been done previously.

I know what you're thinking. You know very few aliens and are not likely to discuss mathematics with them. So why not use the pictures since we're all among friends? If you are less than awestruck by the alien argument, how about this? One of the rules of the game of Mathematics is that no argument is valid that is picture-dependent. You don't like that either? Well then, how accurate are our pictures? Is it possible that a picture could actually mislead us into thinking that something is true, but an argument based on our assumptions shows that it's not true? 

Setting all the above aside, one of the main purposes of this course is to develop your innate ability to reason. We reason mostly with language and we use pictures as aids to our reasoning. Note that our pictures are aids, not crutches. The idea is to help you learn to develop your arguments verbally and use the pictures as aids to your arguments, not integral parts of your arguments.

In a later course, you will use more fundamental axioms, and show that our axioms here are theorems (logical consequences of them). For our purposes at this stage, however, Axioms 1-7 will be our fundamental assumptions. They, along with our definitions and your good sense, will be our tools to try to solve our main problem. Furthermore, by the end of this course, you will be able to explain all of this to your alien friends.

\begin{description}
\item[Problem 1.1:] Consider the following statement: 

  $A$: ``All elephants are large animals.''

  Which of the following statements absolutely, positively must be true if statement $A$ is false?
  \begin{enumerate}[\hspace{1cm}(a)]
    \item All elephants are not large animals.
    \item All elephants are small animals.
    \item Some large animal is not an elephant.
    \item Some elephant is not a large animal.
    \item Some elephant is not large.
    \item Some elephant is not an animal.
    \item Some elephant is either not large or is not an animal.
    \item Not all elephants are large animals.
  \end{enumerate}
\end{description}

What is the point of that problem about elephants? It is simply an exercise in language. Why would you do that in a mathematics course? Well, no matter what your chosen field in life, there will be a vocabulary and a certain use of the language that is peculiar to that field. Lawyers speak their own language, and if you are not conversant in lawyerese, it's difficult to follow their arguments. Doctors have their own peculiar language; and if you've ever listened to computer people conversing in geekese, it can be completely baffling if you're not one of them. Every field makes use of its own esoteric language. (There is another good word, ``esoteric''. Look it up to make sure I've used it properly.)

In mathematics, as opposed to our everyday use of language, we want to be able to state things with absolute preciseness. We want no ambiguity to creep into our statements, because we want it to be true that anyone conversant in the language of mathematics will understand exactly what the author means.

We don't normally speak with such preciseness of language. We use idioms and slang and various phrases that can be interpreted in several ways. But we communicate all right, most of the time. But here is the source of a difficulty. How many people do you know who have had difficulty in mathematics? Probably quite a few. There's an outside chance that you may even consider yourself one of them. 

My observation has been that a significant proportion of the people who have trouble in mathematics do so because they are not really aware that the language is being used in a slightly different manner in mathematics from their everyday use of it. They don't understand what they're asked to do, because the language just doesn't make sense to them. They're not stupid. They're just trying to interpret what's written in the text in a manner that is not appropriate. 

So, how do you learn to use the language in the manner of a mathematician? This technique of defining things that I'm trying to acquaint you with is one aspect of it. Another is to recognize that the purpose of saying things the way that we do in mathematics is to achieve as great a degree of preciseness as possible (i.e., with the least ambiguity as possible). This means we need to adopt some conventions about what certain things mean in mathematics that may have somewhat different meanings in our everyday conversations.

Let's consider the following statement. ``All of Pete's friends have I.Q.s greater than that of a fern.'' We may hope that this statement is true, but how do we go about demonstrating that it is? There are several devices that we can use to analyze statements to test them for truth or falsity. 

One device is to rewrite the statement in another form which means the same as our original statement, but somehow gives us another perspective on it. For example, an equivalent form of our fern statement is: ``If $x$ is a friend of Pete, then $x$ has an I.Q. greater than a fern.'' Does this help? Sometimes rewriting the statement just doesn't help. 

Another device is to consider a special case of our statement to see if this gives us a clue about it. We actually used such a device right in the beginning of this chapter, when we considered the special circle of radius $1$ and center at $O:(0,0)$. But how can we consider a special case of the fern statement? There's nothing obvious here. 

Another device is to generalize our statement, the opposite of considering a special case. Now, what am I talking about? Our statement is: ``All of Pete's friends have I.Q.s greater than a fern.'' A generalization of this: ``All elements in set $A$ are in set $B$.'' In this case, set $A$ is the set of Pete's friends and set $B$ is the set of things having I.Q.s greater than a fern. How can we generalize the other form of our statement? The other form is: ``If $x$ if a friend of Pete, then $x$ has an I.Q. greater than a fern.'' Try to find such a thing before you go on.

A generalization is: ``If statement $P$ is true, then statement $Q$ is true.'' In this case, $P$ is ``$x$ is a friend of Pete'' and $Q$ is ``$x$ has an I.Q. greater than a fern''. Since this form, ``if $P$ is true, then $Q$ is true'', comes up so often, it will benefit us to examine it carefully. Knowing how to interpret and handle such a statement will make your mathematics and science adventures much less painful, and may even make them fun. Let's let $S$ denote the statement ``If $P$ is true, then $Q$ is true.'' 

$S$: If $P$ is true, then $Q$ is true. 

Each of $S$, $P$, and $Q$ is a statement. As a result, each is either true or false. We use the words ``statement'', ``true'', and ``false'' in the following way. A statement is either true or false, and it cannot be both. Things that are not statements are neither true nor false. A question is not a statement, so it is not true. But, neither is it false. You're asking ``What about true love?'' Although true love can be one of the most important things in a person's life, the meaning in that case is different from the meaning intended when we speak of a statement. 

Now, how do we decide if statement $S$ is true?

We can prove that $S$ is true \emph{directly}. What does this mean? It means that we make the supposition that $P$ is true, and then by a series of justifiable steps, we deduce that $Q$ must follow. Note that we do \emph{not} assume that $S$ is true; we only assume $P$ is true. These two things are quite different. You want an example of proving something directly, don't you? Consider the following:

$S$: If each of $x$ and $y$ is a positive number, then $x+y$ is a positive number. If you've gotten this far in mathematics, then I'm sure you believe the statement $S$ above is true. The role of statement $P$ is played by ``each of $x$ and $y$ is a positive number'', and the role of statement $Q$ is played by ``$x+y$ is a positive number''. To show that $S$ is true, we suppose $P$ is true; i.e., suppose that each of $x$ and $y$ is a positive number. Now what do we do? Since each of $x$ and $y$ is a number, we have all the axioms and definitions about numbers that we can use as tools. 

To say $x$ is a positive number means $0<x$. Now that we have supposed that each of $x$ and $y$ is a positive number, we can make the following statements: $0<x$ and $0<y$. We now want to show that $x+y$ is a positive number, which means we have to show that $0<x+y$. Perusing our axioms about numbers, we find one that says ``If each of $a$, $b$, and $c$ is a number and $a<b$, then $a+c<b+c$. Since $0<x$, then this axiom allows us to state that $0+y<x+y$. Make sure that you see how I have used this axiom. Another axiom says $0+a=a$. Hence, using this axiom, we can write $y<x+y$.

We have another axiom, happily, that says ``If $a<b$ and $b<c$, then $a<c$. Using this axiom, we have $0<y$ and $y<x+y$, so that $0<x+y$. By definition, this means $x+y$ is a positive number. What we have done is to prove directly that the statement ``If each of $x$ and $y$ is a positive number, then $x+y$ is a positive number.'' is true. We supposed that $P$ is true, and showed that $Q$ had to follow by a sequence of steps, each of which we could justify by an axiom or definition. 

There is also an \emph{indirect} method that we can use to attack a statement of the form $S$: If $P$ is true, then $Q$ is true. By ``indirect'' is meant that we suppose that $S$ is \emph{not} true, and then show that this leads to a contradiction. Now, what does it mean to say that $S$ is not true? First of all, since $S$ is a statement, if it is not true, then it must be false. ``$S$ is false'' means that there is some situation or condition in which $P$ is true and $Q$ is false. Hence, if we suppose $S$ is false, there are two things that must be true.

The first is that there exists some condition or situation such that $P$ is true under this condition. The second thing that must be true under the supposition that $S$ is false is that $Q$ is false. If this leads to a contradiction, then our assumption that $S$ is false must itself be false. What we would have done is this:
\begin{enumerate}[\hspace{1cm}1)]
  \item We assume that $S$ is false.
  \item We conclude, on the basis of this assumption, that contradictory statements must both be true. This defies our sense of logic, so we conclude that our assumption about $S$ being false is not true. This means that $S$ must be true, because it cannot be false.
\end{enumerate}

Note the difference between the following two statements.
\begin{enumerate}[\hspace{1cm}1)]
  \item $S$: ``If $P$ is true, then $Q$ is true'' is false.
  \item If $P$ is true, then $Q$ is false.
\end{enumerate}
What we've been talking about here is that, if $S$: ``If $P$ is true, then $Q$ is true'' is false, then there exists a situation or condition under which $P$ is true and $Q$ is false.

Let's try this indirect method on our statement $S$: ``If each of $x$ and $y$ is a positive number, then $x+y$ is a positive number.'' Suppose $S$ is false. This means that $P$: ``each of $x$ and $y$ is a positive number'' is true for \emph{some} $x$ and $y$, and it also means that $Q$: ``$x+y$ is a positive number'' is false for that $x$ and that $y$.

Will you grant that there are positive numbers? You see, our assumption is that $S$: ``If each of $x$ and $y$ is a positive number, then $x+y$ is a positive number'' is false. This means that there exist a positive number $x$ and a positive number $y$ such that $x+y$ is not a positive number. Hence, we must show that there does indeed exist a positive number $x$ and a positive number $y$. Since you appear to be skeptical (which I encourage), then consider the following. 

You can show by the axioms for the number system that $0<1$. That makes $1$ a positive number. And, if we define $2$ as $1+1$, i.e. $2 = 1+1$, then we can show that $0<2$ and also that $1 \neq 2$. Hence, there are at least two positive numbers. Since there are positive numbers, let's suppose that each of $x$ and $y$ is a positive number and that, for this $x$ and this $y$, $x+y$ is \emph{not} a positive number (by our assumption that $S$ is false). We have, then, $0<x$, $0<y$, and $0 \neg < x+y$. 

One of our axioms states that, if each of $a$ and $b$ is a number, then one and only one of the following is true: 
\begin{enumerate}[\hspace{1cm}1)]
  \item $a = b$
  \item $a < b$
  \item $b < a$.
\end{enumerate}
From $0 \neg < x+y$ and this axiom, either $0 = x+y$ or $x+y < 0$. We can use the notation $x+y \leq 0$ to denote the condition $x+y = 0$ or $x+y < 0$. $x+y \leq 0$ is read ``$x+y$ is less than or equal to $0$.'' Hence, what we have is $0 < x$, $0 < y$, and $x+y \leq 0$.

Using $x+y \leq 0$ and $0 < x$, we conclude that $x+y < x$. (You can prove that if $a \leq b$ and $b < c$, then $a < c$.) From $x+y < x$, we can conclude that $y < 0$ by subtracting $x$ from both sides of the inequality. But up above we have $0 < y$. $0 < y$ and $y < 0$ lead to $0 < 0$, and this contradicts our axioms. Hence, $0 < x$, $0 < y$, and $x+y \leq 0$ leads to a contradiction of the axioms. 

What do we conclude? Well, if $S$ is false, then there must exist a number $0<x$ and a number $0<y$ such that $x+y \leq 0$. But we've seen that this leads to the statement $0<0$, which contradicts what we know to be true, namely $0=0$. Then $S$ cannot be false. If $S$ cannot be false, then it must be true, since every statement is either true or false. I don't know why anyone would use an indirect argument on this theorem, but it is included simply to illustrate how it works. 

Here is another example, that has some interesting subtleties. 

$S$: If $x$ is a number and $x=x+1$, then $x=-\frac{1}{2}$.

First, let's use the direct method.

$P$: $x$ is a number and $x=x+1$.

$Q$: $x=-\frac{1}{2}$.

Suppose $P$ is true; i.e., that $x$ is a number such that $x=x+1$. We need to show that $x= - \frac{1}{2}$.

\begin{enumerate}[\hspace{1cm}1)]
  \item $x=x+1$ (This is given to us in our assumption that $P$ is true.)
  \item $x^2=(x+1)^2$ (Squaring both sides.)
  \item $x^2=x^2+2x+1$ (Expanding the right side.)
  \item $0=2x+1$ (Adding $-x^2$ to both sides.)
  \item $x= - \frac{1}{2}$ (Solving for $x$.)
\end{enumerate}

So, we have shown $S$ is true. We supposed that $P$ was true, and then, by a series of steps, each justifiable within our system, we showed that $Q$ followed. Hence, we conclude that if $x = x+1$, then $x=-\frac{1}{2}$.

I can hear you saying ``That's ridiculous! There isn't any number $x$ such that $x=x+1$!'' What you say is absolutely true. But notice that what has been said is that $S$ is true. Nothing has been asserted that $P$ is true or that $Q$ is true. $S$ states that a relationship exists between $P$ and $Q$, and we have shown that that relationship is valid. 

Let's look at this problem from the perspective of an indirect method. Suppose that $S$: ``If $x$ is a number and $x = x+1$, then $x=-\frac{1}{2}$'' is false. This means that two things must be true.
\begin{enumerate}[\hspace{1cm}1)]
  \item There exists a number $x$ such that $x = x+1$.
  \item The $x$ in 1) above is not $-\frac{1}{2}$.
\end{enumerate}

But there is no number $x$ such that $x = x+1$. In other words, there is no condition such that $P$ is true. But, if $S$ is false, then such a condition must exist. Hence, it cannot be that $S$ is false. Therefore, $S$ is true. Does this seem strange? It gets stranger.

Consider this:
\begin{enumerate}[\hspace{1cm}1)]
  \item Suppose $x = x+1$
  \item Then $x-1 = x$
  \item $(x-1)^{2} = x^{2}$
  \item $x^{2} - 2x + 1 = x^2$
  \item $-2x+1 = 0$
  \item $x = \frac{1}{2}$ (not $-\frac{1}{2}$ as before)
\end{enumerate}

We've just shown that if $x = x+1$, then $x = \frac{1}{2}$. But before, we showed that if $x = x+1$, then $x=-\frac{1}{2}$. Isn't there some contradiction here since $-\frac{1}{2} \neq \frac{1}{2}$? Notice that we are saying that $S$ is true, and that does \emph{not} mean that $P$ is true, nor does it mean that $Q$ is true. It simply means that \emph{if} $P$ is true, \emph{then} $Q$ is true. See if you can show that if $x = x+1$, then $x=42$.

Showing $S$ to be true directly does not require us to demonstrate that there are any conditions under which $P$ is true. But if $S$ is false, then we must show that some condition exists such that $P$ is true. You'll note that if $P$ is, in fact, false, then no matter what $Q$ is, we cannot fulfill the condition of $S$ being false. Hence, $S$: ``If $P$ is true, then $Q$ is true'' is a true statement if $P$ is known to be false, and this is the case no matter what $Q$ states.

$S$: If the moon is made of green cheese, then fish are mammals. $S$ is true. You may protest that such statements are nonsense. They may be, and they may be ridiculous, but nevertheless, by the conventions that we've discussed here, they are true. By a statement being true, we shouldn't infer that it is profound or even useful. We find, though, that by adopting the conventions about language that we've discussed here, we can achieve great precision in our statements and also the ability to make very fine distinctions when necessary. These are necessary tools in both mathematics and science.

Now let's apply the principles discussed in this rather lengthy discourse, to the statement we started with. $S$: All of Pete's friends have I.Q.s greater than a fern. Converting this to the form ``if $P$, then $Q$'', we have ``If $x$ is a friend of Pete, then $x$ has an I.Q. greater than a fern.'' Let's try an indirect argument.

Suppose $S$ is false. What does that mean? 

$P$: $x$ is a friend of Pete.

$Q$: $x$ has an I.Q. greater than a fern. 

``$S$ is false'' means there is a friend of Pete who does not have an I.Q. greater than a fern. What all this means is that S is false if we can produce a friend of Pete whose I.Q. is not greater than a fern. If we cannot produce such a friend -- that is, if it is impossible to produce such a friend -- then we must conclude that $S$ is true. If there is such a friend, then $S$ is not true. 

But, what if Pete has no friends? Then it is impossible to produce a friend of Pete whose I.Q. is not greater than a fern. Hence, we would have to conclude that it is true that ``all of Pete's friends have an I.Q. greater than a fern''. That is similar to ``if $x = x+1$, then $x = -\frac{1}{2}''$ is true since there is no number $x$ such that $x = x+1$.

It turns out that this device of constructing a second statement which must be true if the original statement is false is so useful that it has a name; it is called a \emph{bare denial}.

\begin{description}
\item[Definition 1.7:] The statement that Statement $A$ is a bare denial of Statement $B$ means:
  \begin{enumerate}[\hspace{1cm}1)]
    \item if $A$ is not true, then $B$ is true, and
    \item if $B$ is not true, then $A$ is true.
  \end{enumerate}
\end{description}

If you are trying to prove that a certain statement is true, show that its bare denial is false. I'll bet you've already done this or at least witnessed it in some prior mathematics class. 

Here is another example. The statement to be proved is: If each of $x$ and $y$ is a positive number, then $\sqrt{x+y} \leq \sqrt{x} + \sqrt{y}$. I'm sure you're familiar with this. Now, what must be true if the statement above is false? 

Try to construct the bare denial before you go on further. What is set $A$? What is set $B$? Or, if you look at it as an example of $S$: ``If $P$ is true, then $Q$ is true'', what is Statement $P$ and what is Statement $Q$?

Statement: If each of $x$ and $y$ is a positive number, then $\sqrt{x+y} \leq \sqrt{x} + \sqrt{y}$.

Bare denial: There exist a positive number $x$ and a positive number $y$ such that $\sqrt{x+y} \neg \leq \sqrt{x} + \sqrt{y}$. 

Suppose the statement is false. Then the bare denial is true.

We know that there exist positive numbers. Hence, let's suppose that each of $x$ and $y$ is a positive number, and that $\sqrt{x+y} \neg \leq \sqrt{x} + \sqrt{y}$. We have $0<x$, $0<y$, and $\sqrt{x+y} \neg \leq \sqrt{x} + \sqrt{y}$. 

Recall that $a \neg \leq b$ means $a>b$, which means that $b<a$. Hence, $\sqrt{x+y} \neg \leq \sqrt{x} + \sqrt{y}$ means $\sqrt{x+y} < \sqrt{x} + \sqrt{y}$. 

Then, since $0< \sqrt{x} \sqrt{y}$, $x+2 \sqrt{x} \sqrt{y} + y < x+y$, squaring both sides. (We have to remember, or show, that $0<a$ implies $0 < \sqrt{a}$.) 

Then $2 \sqrt{x} \sqrt{y} < 0$. But $0<2$, $0 < \sqrt{x}$, and $0 < \sqrt{y}$. Hence, $0 < 2 \sqrt{x} \sqrt{y}$. This is a contradiction.

What all this means is that if the bare denial is true, then $0<0$. Since that's false, the bare denial cannot be true, and that means the statement must be true. We have as a theorem then that if $0<x$ and $0<y$, then $\sqrt{x+y} \leq \sqrt{x} + \sqrt{y}$.

It is probably not an exaggeration to say that if you cannot construct a bare denial to a statement, then you don't understand the statement. It is certainly true that constructing a bare denial to a statement can sometimes help immensely in understanding the original statement.

\begin{description}
\item[Problem 1.2:] Construct a bare denial to each of the following statements.
  \begin{enumerate}[\hspace{1cm}1)]
    \item All of my friends like brussels sprouts.
    \item All of the teachers in my school are bald.
    \item All numbers are positive.
    \item If $x$ is a number, then $x<x+1$.
    \item Some politicians tell the truth, occasionally. (Some means ``at least one''.)
    \item All that glitters is not gold.
    \item All great golfers play dominoes and eat peanut butter and spaghetti sandwiches.
    \item If each of $x$ and $y$ is a number, then $\vert x+y \vert  \leq  \vert x \vert + \vert y \vert$. Can you prove this by showing that the bare denial is false?
    \item All of the elements in set $A$ are in set $B$. 
  \end{enumerate}
\end{description}

Now to some more mundane problems, many of which rely on stuff you got in algebra.

\begin{description}

\item[Problem 1.3:] In each of the following, find the distance between the two points.
  \begin{enumerate}[\hspace{1cm}(a)]
    \item $I:(1,0)$ and $i:(0,1)$
    \item $I:(1,0)$ and $J:(-1,0)$
    \item $I:(1,0)$ and $j:(0,-1)$
    \item $J:(-1,0)$ and $j:(0,-1)$
    \item $P:(1,2)$ and $Q:(3,4)$
    \item $P:(1,2)$ and $Q:(-1,-2)$
    \item $P:(x,y)$ and $Q:(u,v)$
  \end{enumerate}

\item[Problem 1.4:] Find an equation for each of the following pointsets.
  \begin{enumerate}[\hspace{1cm}(a)]
    \item The circle $K$ with center $O:(0,0)$ and radius $1$.
    \item The circle $K$ with center $O:(0,0)$ and radius $3$.
    \item The circle $K$ with center $C:(1,1)$ and radius $2$.
    \item The circle $K$ with center $C:(-2,-3)$ and radius $1$.
    \item The circle $K$ with center $C:(c_{1},c_{2})$ and radius $r>0$.
    \item The line $L$ with slope $2$ and $y$-intercept $3$.
    \item The line $L$ with slope $-1$ and $y$-intercept $1$.
    \item The line $L$ with slope $0$ and $y$-intercept $2$.
    \item The line $L$ with $x$-intercept $1$ and $y$-intercept $2$.
    \item The line $L$ with $x$-intercept $a$ and $y$-intercept $b$.
    \item The line $L$ with slope $0$ and $y$-intercept $b$.
    \item The vertical line $L$ containing $I:(1,0)$.
    \item The vertical line $L$ containing $P:(1,7)$.
    \item The vertical line $L$ containing $P:(a,b)$.
  \end{enumerate}

\item[Problem 1.5:] For each of the following sets of points, determine whether the points are colinear.
  \begin{enumerate}[\hspace{1cm}(a)]
    \item $O: (0,0)$, $I: (1,0)$, $P: (11,0)$
    \item $O: (0,0)$, $P: (1,1)$, $Q: (3,4)$
    \item $P: (1,3)$, $Q: (4,1)$, $R: (7,-1)$
    \item $P: (-2,-1)$, $Q: (3,3)$, $R: (7,6)$
  \end{enumerate}

\item[Problem 1.6:] Suppose $P:(0,5)$, $Q:(2,v)$, and $R:(5,15)$. Suppose also that $PQR$. Solve for $v$ two ways.
  \begin{enumerate}[\hspace{1cm}(a)]
    \item Use ``$PQR$ means $PQ+QR = PR$'' and grind through the algebra to find $v$. This will give you a pretty good algebra workout.
    \item Next, use the coordinates of $P$ and $R$ to get an equation for the line containing them, and use that equation to find $v$.
  \end{enumerate}

\end{description}

After doing (a) and (b) above, you may wonder why anyone would use the first method for anything. Well, consider the following problems.

\begin{description}

\item[Problem 1.7:] Suppose $i: (0,1)$, $Q: (r,s)$, $I: (1,0)$ and $iQI$. Show that $r>0$ and $s>0$.

\item[Problem 1.8:] Suppose $O:(0,0)$, $I:(1,0)$, $i:(0,l)$, $J:(-l,0)$, and $j:(0,-l)$. Show that: (a) $IOJ$, and (b) $iOj$.

\item[Problem 1.9:] Suppose $P:(u,v)$, $P$ is a point of $K:x^{2} + y^{2} = 1$, and $POI$. Show that $P=J$.

\end{description}



\newpage
\begin{center}
\section*{Chapter I Hints and Solutions}
\end{center}

\textbf{Problem 1.1}:
\begin{enumerate}[\hspace{1cm}(a)]
  \item does not have to be true if $A$ is false. It means that no elephant is a large animal.
  \item does not have to be true if $A$ is false. What is the difference in meaning of (a) and (b)?
  \item does not have to be true if $A$ is false. Can you think of an example to prove this?
  \item must be true if $A$ is false. What if there are no elephants? Will it still follow that (d) is true if $A$ is false?
  \item does not have to be true if $A$ is false. Give an example to show this.
  \item does not have to be true if $A$ is false. Give an example. Will the example in (e) work here?
  \item must be true if $A$ is false. Does this mean the same thing as (d)?
  \item must be true if $A$ is false.
\end{enumerate}

\textbf{Problem 1.2}:
\begin{enumerate}[\hspace{1cm}1)]
  \item One of my friends does not like brussels sprouts. An alternative form is ``There exists a friend of mine who does not like brussels sprouts.'' If five of my friends do not like brussels sprouts, is it true that one of my friends does not like brussels sprouts?
  \item Some teacher in my school is not bald. Note: ``Some'' means ``at least one''.
  \item Some number is not positive. All of the following are equivalent to ``Some number is not positive.'':
    \begin{enumerate}[\hspace{1cm}(a)]
      \item There exists a number which is not positive.
      \item There exists a number $x$ such that $x \neg \geq 0$.
      \item There exists a number $x$ such that $x \leq 0$.
    \end{enumerate}
  \item There exists a number $x$ such that $x \geq x+1$.
  \item No politician tells the truth, ever.
  \item There is something that glitters that is gold.

        Note that 6) (original) is an idiom. What it means in everyday language is that just because something glitters doesn't mean that it is gold. But in the context in which we use the language in mathematics, it means that anything that glitters must \emph{not} be gold.

  \item At least one great golfer either does not play dominoes or does not eat peanut butter and spaghetti sandwiches.
  \item There exist a number $x$ and a number $y$ such that $\vert x+y \vert > \vert x \vert + \vert y \vert$.
  \item There exists an element in set $A$ that is not in set $B$.

        Note that if $S$: ``All elements in set $A$ are in set $B$'' is false, then the bare denial ``There exists an element in set $A$ that is not in set $B$'' must be true. To test the bare denial, you must show that there exists an element in set $A$. If you cannot do that because there aren't any elements in set $A$, then the bare denial is false so that the original statement is true. Strange world, isn't it?
\end{enumerate}

\textbf{Problem 1.3}:
\begin{enumerate}[\hspace{1cm}(a)]
  \item $\sqrt{2}$ 
  \item $2$
  \item $\sqrt{2}$
  \item $2$
  \item $2 \sqrt{2}$
  \item $2 \sqrt{5}$ 
  \item $\sqrt{ (x - u)^{2} + (y - v)^{2} }$
\end{enumerate}

\textbf{Problem 1.4}:
\begin{enumerate}[\hspace{1cm}(a)]
  \item x$^{2}$+y$^{2}$ = 1 
  \item x$^{2}$+y$^{2}$ = 9. 
  \item (x-1)$^{2 }$+ (y-l)$^{2}$ = 4 
  \item (x+2)$^{2}$ + (y+3)$^{2 }$= 1
  \item (x-c$_{1})^{2}$ + (y-c$_{2})^{2}$ = r

    In (f) through (n), the answers are given in three forms of an equation of a line, namely $y=Ax + B$, $(\frac{x}{A}) + (\frac{y}{B}) = 1$, and $Ax + By + C = 0$.

  \item $y = 2x+3$, or $\frac{x}{(-3/2)} + \frac{y}{3} = 1$, or $2x - y + 3 = 0$
  \item $y = -x + 1$, $\frac{x}{1} + \frac{y}{2} = 1$, $x + y - 1 = 0$
  \item $y = 2$, ---, $y - 2 = 0$
  \item $y = -2x + 2$, $\frac{x}{1} + \frac{y}{2} = 1$, $x + (\frac{1}{2})y - 0 = 0$
  \item $y = (-b/a)x + b$, $\frac{x}{a} + \frac{y}{b} = 1$, $bx + ay - ab = 0$
  \item $y = b$, ---, $y - b = 0$
  \item ---, ---, $x = 1$
  \item ---, ---, $x = 1$
  \item ---, ---, $x = a$
\end{enumerate}

\textbf{Problem 1.5}:
\begin{enumerate}[\hspace{1cm}(a)]
  \item yes
  \item no
  \item yes
  \item no
\end{enumerate}

\textbf{Problem 1.6}:
\begin{enumerate}[\hspace{1cm}(a)]
  \item $PQ+QR = PR$, which means 
        \begin{align*}
        \sqrt{ (0-2)^2 + (5-v)^2 } & + \sqrt{ (2-5)^2 + (v-15)^2 } \\
                                   & = \sqrt{ (0-5)^2 + (5-15)^2 }
        \end{align*}
        or $\sqrt{ 4 + (5-v)^2 } = 5 \sqrt{5} - \sqrt{ 9 + (v-15)^2 }$.

        Squaring both sides, collecting terms, and reducing, yields 
        \begin{equation*}
        \sqrt{5} \cdot \sqrt{ 9 + (v-15)^2 } = 33-2v
        \end{equation*}

        Squaring again and collecting terms, we have $v^2-18v+81=0$. Hence $(v-9)^2=0$, so that $v=9$.
  \item Using $P$ and $R$, the slope of line $PR$ is $\frac{(15-5)}{(5-0)}=2$. From $P:(0,5)$, the $y$-intercept is $5$.

        Hence, line $PR:y = 2x+5$. If $x=2$, then $y=9$.
\end{enumerate}

\textbf{Problem 1.7}:
\begin{enumerate}[\hspace{1cm}1)]
  \item $i:(0,1)$, $Q:(r,s)$, $I:(1,0)$, and $iQI$.
  \item $iQI$ means $iQ+QI=iI$.
  \item Hence, 
    \begin{equation*}
    \sqrt{(0-r)^{2} + (1-s)^{2} } + \sqrt{(r-1)^{2} + (s-0)^{2} } = \sqrt{ (0-1)^{2}+(1-0)^{2} }
    \end{equation*}
    or $\sqrt{r^{2} + (1-s)^{2} } + \sqrt{(r-1)^{2} + s^{2} } = \sqrt{2}$.
  \item Since $iQI$, $Q$ is on line $iI: y=-x+1$.
  \item Hence, $s = -r+1$.
  \item Then $\sqrt{r^{2} + r^{2}} + \sqrt{ (r-1)^{2} + (1-r)^{2} } = \sqrt{2}$, which reduces to $\vert r \vert  +  \vert r-1 \vert  = 1$
  \item Recall that $\sqrt{x^{2}} = \vert x \vert$. 
  \item $\vert r-1 \vert = 1 - \vert r \vert$, and squaring both sides and collecting terms yields $\vert r \vert = r \geq 0$. 
  \item $r = 0$ implies $s = 1$, and $Q:(0,l)$ implies $Q = i$. But $Q \neq i$ since $iQI$. Hence, $r > 0$. 
  \item $\vert r \vert = r$ implies $\vert r-1 \vert = 1 - r \geq 0$, so $r \leq 1$. This is from $\vert r \vert + \vert r-1 \vert = 1$ and $\vert r \vert = r$.
  \item $r = 1$ implies $s = 0$, which means $Q:(1,0)$, which means $Q=I$. But $Q \neq I$ since $iQI$.
  \item Then $r \neq 1$. Hence, $r<1$. 
  \item Then $s = 1-r > 0$. Hence, $r>0$ and $s>0$.
\end{enumerate}

\textbf{Problem 1.8}:
\begin{enumerate}[\hspace{1cm}1)]
  \item $IO + OJ = \sqrt{(1 - 0)^{2} + (0 - 0)^{2}} +$ $\sqrt{(0 + 1)^{2} + (0 - 0)^{2 }} =$ $\sqrt{1} +$ $\sqrt{1} =$ $1 + 1 = 2$
  \item $IJ = \sqrt{(1+1)^{2} + (0-0)^{2}} = \sqrt{2^{2}} = 2$
  \item Hence, $IO + OJ = IJ$, which means $IOJ$.
  \item Similarly, $iOj$.
\end{enumerate}

\textbf{Problem 1.9}:
\begin{enumerate}[\hspace{1cm}1)]
  \item $P:(u,v)$ is a point of $K:x^{2} + y^{2} = 1$.
  \item Hence $u^{2} + v^{2} = 1$. $POI$ implies $PO+OI = PI$.
  \item Therefore 
    \begin{equation*}
    \sqrt{(u-0)^{2} + (v-0)^{2}} + \sqrt{(0-1)^{2}+(0-0)^{2}} = \sqrt{(u - 1)^{2}+ (v - 0)^{2}}
    \end{equation*}
  \item Then $\sqrt{u^{2}+ v^{2}} +$ $\sqrt{1} =$ $\sqrt{ u^{2}-2u + 1 + v^{2} }$ or $1+1 =$ $\sqrt{1+1-2u} =$ $\sqrt{2-2u}$.
  \item Then $4=2-2u$, which means that $u=-l$. Hence, $v=0$ so that $P:(u,v)$ means $P:(-1,0)$.
  \item That means $P=J$.
\end{enumerate}


\vspace{2cm}


\begin{center}
News Headlines:
\vspace{0.5cm}

1) Stolen painting found by tree.
\vspace{0.5cm}

2) War dims hope for peace.
\end{center}



\chapter*{Chapter II}
\addcontentsline{toc}{chapter}{\numberline{}Chapter II}
\markboth{Chapter II}{Chapter II}


Now, on to our main problem. Recall that our first attempt at simplification of the problem was to limit the discussion to the unit circle with center at the origin. Let's simplify the problem even further, by limiting the discussion to what we might call the ``upper semi-circle of the unit circle $K: x^{2} + y^{2} = 1$''; that is, the semi-circle $\arc{JiI}$. We know from Axiom 1 that there is a non-negative number which is the length of $\arc{JiI}$. Let the symbol ``$\pi$'' denote that number.

\begin{description}

\item[Definition 2.1:] $\pi = l(\arc{JiI})$.

\end{description}

This doesn't tell us anything about the size of $\pi$, but perhaps we can get some idea of how big it is.

\begin{center}
%TeXCAD Options
%\grade{\on}
%\emlines{\off}
%\epic{\off}
%\beziermacro{\on}
%\reduce{\on}
%\snapping{\off}
%\quality{8.00}
%\graddiff{0.01}
%\snapasp{1}
%\zoom{4.0000}
\unitlength 1mm % = 2.85pt
\linethickness{0.4pt}
\ifx\plotpoint\undefined\newsavebox{\plotpoint}\fi % GNUPLOT compatibility
\begin{picture}(115.75,59)(0,70)
\put(19,82.5){\line(1,0){82.5}} \put(59.5,82.5){\line(0,1){46.5}}
\qbezier(59.5,122.75)(91.13,122.75)(105.25,82.75)
\qbezier(59.5,122.75)(22.25,120.38)(15,82.5)
\put(101.5,82.5){\line(1,0){14.25}}
\put(19,82.5){\line(-1,0){11.75}}
%\emline(7.25,82.5)(7,82.25)
\multiput(7.25,82.5)(-.03125,-.03125){8}{\line(0,-1){.03125}}
%\end
\put(58.75,124.75){\makebox(0,0)[cc]{$i:(0,1)$}}
\put(59.25,79){\makebox(0,0)[cc]{$O:(0,0)$}}
\put(104.5,79.75){\makebox(0,0)[cc]{$I:(1,0)$}}
\put(15,79.75){\makebox(0,0)[cc]{$J:(-1,0)$}}
\put(60,71.5){\makebox(0,0)[cc]{Figure 2.1}}
\end{picture}
\end{center}

\begin{description}

\item[Question 2.1:] Can you see a way to convince yourself that $2<\pi<4$? I'm not looking for a proof here, just an argument that is convincing on the surface.

  What you have done to get this inequality probably won't convince our alien friend, but I'll bet you're reasonably sure it's true.

\item[Question 2.2:] Is it true that no matter what semi-circle you look at in $K:x^{2} + y^{2}=1$, its length is $\pi$? Wouldn't this be awkward if it were not true?

\end{description}

Here is a problem set to get you started in earnest on our main problem. Draw a picture of $K:x^{2} + y^{2} = 1$, and label $I:(1,0)$, $i:(0,1)$, $J:(-1,0)$, $j:(0,-1)$, and $O:(0,0)$. Put your axioms and definitions in front of you, since these are your major tools, and then use the picture as an aid in the following:

\begin{description}

\item[Problem 2.1:] Show that $iI = Ij = jJ = Ji = \sqrt{2}$.

\item[Problem 2.2:] Show that $l(\arc{iI}) = l(\arc{Ij}) = l(\arc{jJ}) = l(\arc{Ji}) = \frac{\pi}{2}$.

\item[Problem 2.3:] Suppose $P:(u,v)$ is a point of $K:x^{2} + y^{2} = 1$, $Q$ is a point of $K$, and $POQ$ is true. Show that $Q:(-u,-v)$.

\item[Problem 2.4:] Show that if $S$ is a semi-circle of $K$, then $S$ has length $\pi$.

\item[Problem 2.5:] Show that if $\arc{ipI}$ and $Pi = PI$, then $P:(\frac{\sqrt{2}}{2}, \frac{\sqrt{2}}{2})$ and $l(\arc{Pi}) = l(\arc{PI}) = \frac{\pi}{4}$.

\end{description}

Our alien student seems to be squirming over this last problem. What is his problem? Did you have great difficulty with this one as he seems to have had? He claims that there are two possible answers to the problem, not just one. He states that $P:(\frac{\sqrt{2}}{2}, \frac{\sqrt{2}}{2})$ is a solution but $P:(-\frac{\sqrt{2}}{2}, -\frac{\sqrt{2}}{2})$ is also a solution. Why did we pick the solution with the positive coordinates? Well, our picture makes this seem perfectly obvious. But our alien student can't see our picture.

\begin{description}

\item[Question 2.3:] Can you figure out a way of convincing our alien student that the only valid solution to this problem is $P:(\frac{\sqrt{2}}{2}, \frac{\sqrt{2}}{2})$ because $P:(-\frac{\sqrt{2}}{2}, -\frac{\sqrt{2}}{2})$ is not a point of $\arc{iI}$? Remember, you can't convince our student of anything by reference to the picture. It must rest entirely on our axioms and definitions.

  You will note that a very similar problem will annoy us in problem 2.6 and in problem 2.7 .

\item[Problem 2.6:] Show that, if $\arc{JQi}$ and $JQ = Qi$, then $Q:(-\frac{\sqrt{2}}{2}, \frac{\sqrt{2}}{2})$ and $l(\arc{QI}) = \frac{3\pi}{4}$.

\item[Problem 2.7:] Suppose $\arc{iPQI}$ and $l(\arc{iP}) = l(\arc{PQ}) = l(\arc{QI})$. Show that $P:(\frac{1}{2}, \frac{\sqrt{3}}{2})$ and $Q:(\frac{\sqrt{3}}{2},\frac{1}{2})$, and furthermore, that $l(\arc{PQ}) = \frac{\pi}{6}$.

\item[Problem 2.8:] Convince yourself that $3 < \pi < 3.5$.

\end{description}

What good is all this going to do for us in solving our problem? It's certainly not clear how this relates to finding a relationship between arc length and chord length. But, let's look at what we've got.

We have the arc length $\pi$ associated with the chord length of $2$; i.e., $l(\arc{JI}) = \pi$ and $JI=2$.

We have the arc length $\frac{\pi}{2}$ associated with the chord length of $\sqrt{2}$; i.e., $l(\arc{iI}) = \frac{\pi}{2}$ and $iI=\sqrt{2}$. 

We have $l(\arc{II}) = 0$ and $II = 0$.

We a1so have the arc 1ength $\frac{\pi}{4}$ associated with the chord length-- oops, we haven't calculated that yet, have we? 

We have the point $P$ of arc $\arc{iI}$ such that $l(\arc{PI}) = \frac{\pi}{4}$, and we found that the coordinates of $P$ are $(\frac{\sqrt{2}}{2}, \frac{\sqrt{2}}{2})$; i.e., $P: (\frac{\sqrt{2}}{2}, \frac{\sqrt{2}}{2})$. Let's assume these coordinates are correct at present, even though our student (and maybe even you) is not convinced of this. Knowing these coordinates, calculate $PI$. You should get $PI = \sqrt{2 - \sqrt{2}}$. We can, by the same reasoning, find the chord length associated with the arc length $\frac{3\pi}{4}$. We can even get the chord lengths associated with the arc 1engths $\frac{\pi}{6}$ and $\frac{\pi}{3}$. Can you get these chord lengths associated with the arc lengths $\frac{2\pi}{3}$ and $\frac{5\pi}{6}$?

Let's put this information in a more convenient form. We'll construct a table that will make it easier for us to see all the information at once.

\begin{center}
Table 2.1
\end{center}
\begin{table}[htbp]
\begin{center}
\begin{tabular}{|l|l|}
  \hline 
  arc length    &  chord length \\
  \hline 
  $l(\arc{PI})$ & $PI$ \\
  
  \hline
  0             &  0 \\
  $\pi/6$       &    \\
  $\pi/4$       &    \\
  $\pi/3$       &    \\
  $\pi/2$       & $\sqrt{2}$ \\
  $2\pi/3$      &    \\
  $3\pi/4$      &    \\
  $5\pi/6$      &    \\
  $\pi$         &  2 \\
  \hline
\end{tabular}
\label{tbl2.1}
\end{center}
\end{table}

Let's read the table. The first pair is $(0,0)$. It says that if $P$ is a point of $\arc{JiI}$ and $l(\arc{PI}) = 0$, then $PI = 0$. Note that this means $P=I$. The second pair that is entered in the table so far is $(\frac{\pi}{2}, \sqrt{2})$. It says that if $P$ is a point of $\arc{JiI}$ and $l(\arc{PI}) = \frac{\pi}{2}$, then $PI = \sqrt{2}$. This means that $P = i$.

Now fill in the remaining entries, pairing arc length with chord length. Use a picture of semi-circle $\arc{JiI}$ to aid in your thinking. (You're worried that maybe we don't have the correct coordinates for some of these points, aren't you? I certainly hope you are worried, unless of course you have proved that they are correct.) 

I hope it seems to you that listing these number pairs, arc length paired with its corresponding chord length, in a table like Table 2.1 makes it easier to see all of the information at once rather than just enumerating them in the text. An even more illuminating procedure is to plot the number pairs in a graph so that we actually get a picture of how the pairs look in relation to one another. How do we do this? Well, here is one way.

\begin{center}
%TeXCAD Options
%\grade{\on}
%\emlines{\off}
%\epic{\off}
%\beziermacro{\on}
%\reduce{\on}
%\snapping{\off}
%\quality{8.00}
%\graddiff{0.01}
%\snapasp{1}
%\zoom{4.0000}
\unitlength 1mm % = 2.85pt
\linethickness{0.4pt}
\ifx\plotpoint\undefined\newsavebox{\plotpoint}\fi % GNUPLOT compatibility
\begin{picture}(91.75,49.75)(0,90)
\put(15.75,139.5){\line(0,-1){36.75}}
\put(15.75,102.75){\line(1,0){62.75}}
%\dashline{1}(16,139.5)(77.25,139.75)
\put(15.93,139.43){\line(1,0){.988}}
\put(17.91,139.44){\line(1,0){.988}}
\put(19.88,139.45){\line(1,0){.988}}
\put(21.86,139.45){\line(1,0){.988}}
\put(23.83,139.46){\line(1,0){.988}}
\put(25.81,139.47){\line(1,0){.988}}
\put(27.78,139.48){\line(1,0){.988}}
\put(29.76,139.49){\line(1,0){.988}}
\put(31.74,139.49){\line(1,0){.988}}
\put(33.71,139.5){\line(1,0){.988}}
\put(35.69,139.51){\line(1,0){.988}}
\put(37.66,139.52){\line(1,0){.988}}
\put(39.64,139.53){\line(1,0){.988}}
\put(41.62,139.53){\line(1,0){.988}}
\put(43.59,139.54){\line(1,0){.988}}
\put(45.57,139.55){\line(1,0){.988}}
\put(47.54,139.56){\line(1,0){.988}}
\put(49.52,139.57){\line(1,0){.988}}
\put(51.49,139.57){\line(1,0){.988}}
\put(53.47,139.58){\line(1,0){.988}}
\put(55.45,139.59){\line(1,0){.988}}
\put(57.42,139.6){\line(1,0){.988}}
\put(59.4,139.61){\line(1,0){.988}}
\put(61.37,139.62){\line(1,0){.988}}
\put(63.35,139.62){\line(1,0){.988}}
\put(65.32,139.63){\line(1,0){.988}}
\put(67.3,139.64){\line(1,0){.988}}
\put(69.28,139.65){\line(1,0){.988}}
\put(71.25,139.66){\line(1,0){.988}}
\put(73.23,139.66){\line(1,0){.988}}
\put(75.2,139.67){\line(1,0){.988}}
%\end
%\dashline{1}(77.25,139.75)(77.25,103.25)
\put(77.18,139.68){\line(0,-1){.986}}
\put(77.18,137.71){\line(0,-1){.986}}
\put(77.18,135.73){\line(0,-1){.986}}
\put(77.18,133.76){\line(0,-1){.986}}
\put(77.18,131.79){\line(0,-1){.986}}
\put(77.18,129.81){\line(0,-1){.986}}
\put(77.18,127.84){\line(0,-1){.986}}
\put(77.18,125.87){\line(0,-1){.986}}
\put(77.18,123.9){\line(0,-1){.986}}
\put(77.18,121.92){\line(0,-1){.986}}
\put(77.18,119.95){\line(0,-1){.986}}
\put(77.18,117.98){\line(0,-1){.986}}
\put(77.18,116){\line(0,-1){.986}}
\put(77.18,114.03){\line(0,-1){.986}}
\put(77.18,112.06){\line(0,-1){.986}}
\put(77.18,110.09){\line(0,-1){.986}}
\put(77.18,108.11){\line(0,-1){.986}}
\put(77.18,106.14){\line(0,-1){.986}}
\put(77.18,104.17){\line(0,-1){.986}}
%\end
%\dashline{1}(16,126.5)(45.5,126.75)
\put(15.93,126.43){\line(1,0){.983}}
\put(17.9,126.45){\line(1,0){.983}}
\put(19.86,126.46){\line(1,0){.983}}
\put(21.83,126.48){\line(1,0){.983}}
\put(23.8,126.5){\line(1,0){.983}}
\put(25.76,126.51){\line(1,0){.983}}
\put(27.73,126.53){\line(1,0){.983}}
\put(29.7,126.55){\line(1,0){.983}}
\put(31.66,126.56){\line(1,0){.983}}
\put(33.63,126.58){\line(1,0){.983}}
\put(35.6,126.6){\line(1,0){.983}}
\put(37.56,126.61){\line(1,0){.983}}
\put(39.53,126.63){\line(1,0){.983}}
\put(41.5,126.65){\line(1,0){.983}}
\put(43.46,126.66){\line(1,0){.983}}
%\end
%\dashline{1}(45.5,126.75)(45.5,103.25)
\put(45.43,126.68){\line(0,-1){.979}}
\put(45.43,124.72){\line(0,-1){.979}}
\put(45.43,122.76){\line(0,-1){.979}}
\put(45.43,120.8){\line(0,-1){.979}}
\put(45.43,118.85){\line(0,-1){.979}}
\put(45.43,116.89){\line(0,-1){.979}}
\put(45.43,114.93){\line(0,-1){.979}}
\put(45.43,112.97){\line(0,-1){.979}}
\put(45.43,111.01){\line(0,-1){.979}}
\put(45.43,109.05){\line(0,-1){.979}}
\put(45.43,107.1){\line(0,-1){.979}}
\put(45.43,105.14){\line(0,-1){.979}}
%\end
\put(15.75,120.75){\circle{.5}} 
\put(30.25,102.75){\circle{.5}}
\put(61,102.75){\circle{.5}}
\put(80.25,139.25){\makebox(0,0)[cc]{$x$}}
\put(49.25,127.25){\makebox(0,0)[cc]{$x$}}
\put(18,105){\makebox(0,0)[cc]{$x$}}
\put(12.75,101.5){\makebox(0,0)[cc]{$0$}}
\put(12.75,139){\makebox(0,0)[cc]{$2$}}
\put(12.75,119){\makebox(0,0)[cc]{$1$}}
\put(30,100.25){\makebox(0,0)[cc]{$\frac{\pi}{4}$}}
\put(45.5,100.75){\makebox(0,0)[cc]{$\frac{\pi}{2}$}}
\put(60.75,100){\makebox(0,0)[cc]{$\frac{3\pi}{4}$}}
\put(77,100.5){\makebox(0,0)[cc]{$\pi$}}
\put(91.75,119.25){\makebox(0,0)[cc]{Figure 2.2}}
%\qbezvec(78.25,92)(64.5,88.75)(52.75,101.5)
\put(52.75,101.5){\vector(-1,1){.07}}\qbezier(78.25,92)(64.5,88.75)(52.75,101.5)
%\end
%\qbezvec(4.5,118.75)(6.88,113.25)(14.75,112.75)
\put(14.75,112.75){\vector(1,0){.07}}\qbezier(4.5,118.75)(6.88,113.25)(14.75,112.75)
%\end
\put(4,121.75){\makebox(0,0)[cc]{$PI$}}
\put(82.75,92){\makebox(0,0)[cc]{$l(\arc{PI})$}}
\end{picture}
\end{center}

I have plotted the three number pairs that were already in Table 2.1 before you completed it. Each of these plotted points is marked with an $x$. Plot the remaining entries from Table 2.1. Obviously you will have to use approximations to your entries in plotting the points, but this is alright. We are merely sketching the graph to see what it might suggest to us. 

Well, what do you think? Does the sketch you've just made suggest anything to you? Let me ask some very specific questions about it.

\begin{description}

\item[Question 2.4:] Is the relationship between arc length and chord length linear?

\item[Question 2.5:] Is the relationship circular?

\item[Question 2.6:] Is it parabolic?

\end{description}

And what do these questions mean, anyway? This may be unfamiliar territory to you, so let me elaborate. You were asked to generate some number pairs (arc length paired with its corresponding chord length), and then you were asked to plot these number pairs as points on a graph. Now you are being asked whether some rather familiar geometric figure will fit your graph.

Let's start with the first familiar geometric figure, and I'll give you a hint or two to get you started. Here is the question: Is the relationship between arc length and chord length linear? This means, does there exist a line that contains all of the points on your graph? Is it obvious from your graph that there is no such line? Would it be obvious to our alien friend who cannot see your graph? How would we convince him that the relationship is not linear?

If there were such a line, then it would have an equation of the form $Ax + By + C = 0$. Do you recall this from your algebra? Hence, we can rephrase the question in algebraic terms. Does there exist a number triple $(A,B,C)$ such that every number pair in your list (arc length, chord length) will satisfy this equation? That is, if you pick any number pair $(l(\arc{PI}),PI)$ in your list (table), then $A(l(\arc{PI})) + B(PI) + C = 0$? If no such number triple exists, then the relationship between arc length and chord length is not linear.

Do not read beyond this point before you give a very serious effort to this problem of showing that the relationship between arc length and chord length is not linear.

Well, how did you do? Does the question make sense to you? Let's examine this problem in some detail, and then I'll leave the harder ones for you. (Why do teachers always work the easy ones and leave the difficult ones for the students?) 

We have the following table of number pairs:

\begin{center}
Table 2.1 (abbreviated)
\end{center}
\begin{table}[htbp]
\begin{center}
\begin{tabular}{|l|l|}
  \hline 
  $l(\arc{PI})$ & $PI$ \\
  \hline 
  0             & 0 \\
  $\pi/4$       & $\sqrt{2-\sqrt{2}}$ \\
  $\pi/2$       & $\sqrt{2}$ \\
  $3\pi/4$      & $\sqrt{2+\sqrt{2}}$ \\
  \hline
\end{tabular}
\label{tbl-q2.6}
\end{center}
\end{table}

You have more in your table, but these will be sufficient to illustrate what we're doing.

You read the table in the following way: The numbers in the left column under the symbol $l(\arc{PI})$ are arc lengths, and the numbers in the right column under the symbol $PI$ are chord lengths. For example, if $P$ is a point of our upper semi-circle $\arc{JiI}$ such that the arc length from $P$ to $I$ -- $l(\arc{PI})$ -- is $\pi/4$, then paired with that arc length is the chord length $PI = \sqrt{2-\sqrt{2}}$.

Let's examine in some detail how to get this chord length. First, let's look at the picture and put the given information in it. We are given that $P$ is a point of arc $\arc{iI}$ and that $l(\arc{PI}) = \frac{\pi}{4}$. Can you see that this means $P \neq i$ and $P \neq I$?

\begin{center}
%TeXCAD Options
%\grade{\on}
%\emlines{\off}
%\epic{\off}
%\beziermacro{\on}
%\reduce{\on}
%\snapping{\off}
%\quality{8.00}
%\graddiff{0.01}
%\snapasp{1}
%\zoom{4.0000}
\unitlength 0.8mm % = 2.85pt
\linethickness{0.4pt}
\ifx\plotpoint\undefined\newsavebox{\plotpoint}\fi % GNUPLOT compatibility
\begin{picture}(115.75,56.5)(0,77)
\put(19,82.5){\line(1,0){82.5}} \put(59.5,82.5){\line(0,1){46.5}}
\qbezier(59.5,122.75)(91.13,122.75)(105.25,82.75)
\qbezier(59.5,122.75)(22.25,120.38)(15,82.5)
\put(101.5,82.5){\line(1,0){14.25}}
\put(19,82.5){\line(-1,0){11.75}}
%\emline(7.25,82.5)(7,82.25)
\multiput(7.25,82.5)(-.03125,-.03125){8}{\line(0,-1){.03125}}
%\end
\put(58.75,124.75){\makebox(0,0)[cc]{$i:(0,1)$}}
\put(59.25,79){\makebox(0,0)[cc]{$O:(0,0)$}}
\put(104.5,79.75){\makebox(0,0)[cc]{$I:(1,0)$}}
\put(15,79.75){\makebox(0,0)[cc]{$J:(-1,0)$}}
\put(59.25,133.5){\makebox(0,0)[cc]{Figure 2.3}}
\end{picture}
\end{center}

To find the chord length $PI$, we need the coordinates of $P$. Suppose $P:(u,v)$. Our problem, then, is to calculate $u$ and $v$. By definition, $l(\arc{JiI}) = \pi$, or $l(\arc{JI}) = \pi$. We also have that $l(\arc{Ji}) + l(\arc{iI}) = l(\arc{JI}) = \pi$, from Axiom 1. But it is also true that $Ji=iI$. Hence from Axiom 3, $l(\arc{Ji}) = l(\arc{iI})$. Follow this on your picture. All this combines to yield $l(\arc{iI}) = \frac{\pi}{2}$. Since $\arc{iPI}$ and $P \neq i$ and $P \neq I$, $l(\arc{iP}) + l(\arc{PI}) = l(\arc{iI})$, from Axiom 1. But $l(\arc{PI}) = \frac{\pi}{4}$ and $l(\arc{iI}) = \frac{\pi}{2}$, so that $l(\arc{iP}) = \frac{\pi}{4}$; i.e., $l(\arc{iP}) = l(\arc{PI})$. Therefore, from Axiom 3, $iP = PI$. Make sure you follow this line of reasoning, because it will come up again.

Now, if you have not already calculated the coordinates of $P$ (i.e., $u$ and $v$), use the results above ($iP = PI$) to do so now before you go on.

Calculating $u$ and $v$ of $P:(u,v)$ using $iP = PI$. $P:(u,v)$, $i:(0,1)$, $I:(1,0)$, and we have $iP = PI$. Hence $\sqrt{(0-u)^{2} + (1-v)^{2}} =$ $\sqrt{(u-1)^{2} + (v-0)^{2}}$. Then $u^{2} + (1-v)^{2} = (u-1)^{2} + v^{2}$. This reduces to $-2v = -2u$. (Remember $P$ is a point of $K:x + y^{2} = 1$. Hence $u^{2} + v^{2} = 1$.) So $v = u$. This means that, since $u^{2} + v^{2} = 1$, $u^{2} + u^{2} = 1$, or $u = \pm \sqrt{\frac{1}{2}} = \pm \frac{\sqrt{2}}{2}$. How do we eliminate the possibility that $u = - \frac{\sqrt{2}}{2}$? This is our alien student's objection. Well, assuming that we can handle this, we have $u = \frac{\sqrt{2}}{2}$, and since $v = u$, $v = \frac{\sqrt{2}}{2}$ so that $P:(\frac{\sqrt{2}}{2}, \frac{\sqrt{2}}{2})$. Then $PI = \sqrt{2-\sqrt{2}}$.

Our alien friend is jumping up and down about to grow feathers. What is his problem now? Well, he claims that we fudged back there. We claimed that, if $P$ is a point of the upper semi-circle $\arc{JiI}$ such that $l(\arc{PI}) = \frac{\pi}{4}$, then the chord length $PI = \sqrt{2-\sqrt{2}}$. But, he says, what we showed was that, if $P$ is a point of arc $\arc{iI}$ such that $l(\arc{PI}) = \frac{\pi}{4}$, then $PI = \sqrt{2-\sqrt{2}}$. Do you see the difference? Are we in trouble? Can you resolve this? We found that, if $\arc{iPI}$ and $l(\arc{PI}) = \frac{\pi}{4}$, then $P:(\frac{\sqrt{2}}{2},\frac{\sqrt{2}}{2})$, which led to $PI = \sqrt{2-\sqrt{2}}$.

\vspace{0.3cm}

\textbf{Question 2.7}: Is it possible that there is some point $Q$ of $\arc{JiI}$ different from $P:(\frac{\sqrt{2}}{2},\frac{\sqrt{2}}{2})$ such that $l(\arc{QI}) = \frac{\pi}{4}$? This would be very awkward if it's true. Of course, our picture indicates that there cannot be any such point $Q$, but our alien friend, sadly, cannot see our picture. Can you convince him that there is only one point of $\arc{JiI}$ such that the arc length from that point to $I$ is $\frac{\pi}{4}$?

A note about the chord lengths that you calculate. It's better to leave them in square root form rather than decimal form. Decimals are only approximations usually, and we may lose something important in our calculations if we use, for example, $1.41$ instead of $\sqrt{2}$. You will have to use approximations when sketching graphs, but other than that, we should usually avoid approximations.

Getting back to the reading of the table, you can see that if the arc length from a point of $\arc{JiI}$ to $I$ is $\sqrt{2}$, then the chord length is $\sqrt{2}$. Furthermore, the only such point is $i:(0,1)$. We have the following number pairs in our collection in the abbreviated Table 2.1: $(0,0)$, $(\frac{\pi}{4}, \sqrt{2-\sqrt{2}})$, $(\frac{\pi}{2}, \sqrt{2})$, $(\frac{3\pi}{4}, \sqrt{2+\sqrt{2}})$ and $(\pi, 2)$. It is just more convenient for us if we put them in a table like Table 2.1. 

Now, to the question that we want to answer. If you plotted these number pairs, would some line contain all these points? Sketching a graph of these is not too difficult.

\begin{center}
%TeXCAD Options
%\grade{\on}
%\emlines{\off}
%\epic{\off}
%\beziermacro{\on}
%\reduce{\on}
%\snapping{\off}
%\quality{8.00}
%\graddiff{0.01}
%\snapasp{1}
%\zoom{4.0000}
\unitlength 1mm % = 2.85pt
\linethickness{0.4pt}
\ifx\plotpoint\undefined\newsavebox{\plotpoint}\fi % GNUPLOT compatibility
\begin{picture}(73.75,55.5)(0,87)
\put(18.5,142.25){\line(0,-1){48.75}}
\put(12.75,99.5){\line(1,0){61}} 
\put(18.5,142.25){\circle{.5}}
\put(73.25,99.5){\circle{.5}} 
\put(10,124.5){\makebox(0,0)[cc]{$PI$}}
\put(45,96.25){\makebox(0,0)[cc]{$l(\arc{PI})$}}
\put(73.5,96.5){\makebox(0,0)[cc]{$\pi$}}
\put(15.75,97){\makebox(0,0)[cc]{$0$}}
\put(15,142.5){\makebox(0,0)[cc]{$2$}}
\put(44.75,88.25){\makebox(0,0)[cc]{Graph 2.4}}
\end{picture}
\end{center}

Note that the arc lengths go from $0$ to $\pi$, and the chord lengths go from $0$ to $2$. The scales don't have to be the same on the two axes to sketch this, so we just pick some convenient points to represent $\pi$ the $l(\arc{PI})$ axis and $2$ on the $PI$ axis. These determine the rest of the points on the two axes. Filling some of those in on both axes, we have the following after plotting the number pairs.

\begin{center}
%TeXCAD Picture [fig2.5.pic]. Options:
%\grade{\on}
%\emlines{\off}
%\epic{\off}
%\beziermacro{\on}
%\reduce{\on}
%\snapping{\off}
%\quality{8.00}
%\graddiff{0.01}
%\snapasp{1}
%\zoom{4.0000}
\unitlength 1mm % = 2.85pt
\linethickness{0.4pt}
\ifx\plotpoint\undefined\newsavebox{\plotpoint}\fi % GNUPLOT compatibility
\begin{picture}(77,57.5)(0,87)
\put(18.5,142.25){\line(0,-1){48.75}}
\put(12.75,99.5){\line(1,0){61}}
\put(18.5,142.25){\circle{.5}}
\put(73.25,99.5){\circle{.5}}
\put(73,96.5){\makebox(0,0)[cc]{$\pi$}}
\put(15.75,96.5){\makebox(0,0)[cc]{$0$}}
\put(44.25,96.5){\makebox(0,0)[cc]{$\pi/2$}}
\put(31.5,96.5){\makebox(0,0)[cc]{$\pi/4$}}
\put(57.75,96.5){\makebox(0,0)[cc]{$3\pi/4$}}
\put(66,138.5){\makebox(0,0)[cc]{$(\frac{3\pi}{4},\sqrt{2+\sqrt{2}})$}}
\put(51,128.75){\makebox(0,0)[cc]{$(\frac{\pi}{2},\sqrt{2})$}}
\put(33.75,117.25){\makebox(0,0)[cc]{$(\frac{\pi}{4},\sqrt{2-\sqrt{2}})$}}
\put(77,144.5){\makebox(0,0)[cc]{$(\pi,2)$}}
\put(24,102.25){\makebox(0,0)[cc]{$(0,0)$}}
\put(15,142.5){\makebox(0,0)[cc]{$2$}}
\put(15,120.25){\makebox(0,0)[cc]{$1$}}
\put(15,131.25){\makebox(0,0)[cc]{$1.5$}}
\put(15,109){\makebox(0,0)[cc]{$0.5$}}
\put(44.75,88.25){\makebox(0,0)[cc]{Graph 2.5}}
%\dottedline(18.5,142.5)(73.75,142.5)
\multiput(18.43,142.43)(.98661,0){57}{{\rule{.4pt}{.4pt}}}
%\end
%\dottedline(73.75,142.5)(73.75,99.75)
\multiput(73.68,142.43)(0,-.99419){44}{{\rule{.4pt}{.4pt}}}
%\end
%\dottedline(57.75,99.5)(57.75,139.75)
\multiput(57.68,99.43)(0,.98171){42}{{\rule{.4pt}{.4pt}}}
%\end
%\dottedline(57.75,139.75)(18.75,139.75)
\multiput(57.68,139.68)(-.975,0){41}{{\rule{.4pt}{.4pt}}}
%\end
%\dottedline(44.25,99.5)(44.25,128.75)
\multiput(44.18,99.43)(0,.975){31}{{\rule{.4pt}{.4pt}}}
%\end
%\dottedline(44.25,128.75)(18.75,129)
\multiput(44.18,128.68)(-.9808,.0096){27}{{\rule{.4pt}{.4pt}}}
%\end
%\dottedline(31.5,99.5)(31.5,115)
\multiput(31.43,99.43)(0,.9688){17}{{\rule{.4pt}{.4pt}}}
%\end
%\dottedline(31.5,115)(18.5,115)
\multiput(31.43,114.93)(-.9286,0){15}{{\rule{.4pt}{.4pt}}}
%\end
\end{picture}
\end{center}

Approximations were necessary in order to plot these number pairs as points. It looks fairly obvious that no line contains all five of our points. But now we have to prove that's true by algebraic methods. You have already worked on this problem, and if you solved it, that's great. But if you got stuck, then please follow this carefully.

I am going to give you a series of steps. If you have already done the first step, then go on to the next step. Continue this until you get to the first step that you didn't do. Read that step, but read no further until you have given yourself the chance to complete the problem from there on. The object is for you to do as much of the problem on your own as you can.

\vspace{0.3cm}

\emph{Step 1}: Suppose $L: Ax + By + C = 0$, and $L$ contains all five of the points listed in our table. You will recall from algebra that every line in the plane has an equation of the form above, and furthermore, every pointset in the plane with such an equation is a line as long as not both of $A$ and $B$ are $0$.

  Note that what you are doing here is assuming that there \emph{is} a line that contains all five of the points of your graph. Why would we do that, if our objective is to show that no such line exists? Well, what if, on the basis of that assumption, we arrive at something that contradicts something that we know is true? What would that mean? For example, what if, by assuming that there is a line that contains all five of our points, we arrive at the conclusion that $0 = 1$? This means that, if such a line exists, then $0 = 1$. But, we know that $0 \neq 1$. What does this tell us about the existence of such a line? Think about this very carefully. Such thinking will be rewarded.

\vspace{0.3cm}

\emph{Step 2}: Since $(0,0)$ is in our collection, this should yield you the result that $C = 0$. Hence, $L: Ax + By = 0$.

\vspace{0.3cm}

\emph{Step 3}: Put two other number pairs from the collection in Table 2.1 in the equation $Ax + By = 0$. This yields two equations involving $A$ and $B$.

\vspace{0.3cm}

\emph{Step 4}: Solve for $A$ and $B$ from the two equations in Step 3.

\vspace{0.3cm}

\emph{Step 5}: You have now used three of the number pairs in our collection -- $(0,0)$ and the two you selected in step 3. Furthermore, you know what $A$, $B$, and $C$ are. Put one of the other two number pairs into your equation, for which you now have values for $A$ and $B$. Does it fit the equation? If it does, try the other number pair. If one of these two number pairs does not satisfy the equation for your line, what does that mean?

\vspace{0.3cm}

\emph{Step 6}: You assumed that your line $L$ contained all five points of your graph. But from step 5, you saw that was false. You have a contradiction, so something is wrong. The contradiction is a result of assuming that some line contains all five points of the graph. Hence, that assumption must be false. That means that no line contains all five points of the graph. And this is what we were trying to get. Did you find a simpler way of doing this? Like most problems, there is more than one way to do it. 

  You can use a similar technique to show that no circle contains all of the points. You will find that the arithmetic is much easier if you add at least one more number pair to your table. Try finding the number pair in your table whose first term is $\frac{\pi}{3}$. If you don't see how to do this yet, you can still show that there is no circle containing all five of your points, but the arithmetic is a little tedious. For a parabola, first see if a parabola with a vertical axis of symmetry contains all the points; i.e., one with equation $y = Ax^{2} + Bx + C$. This one is easier than the circle.

  If you solve this problem for both the circle and the parabola with a vertical axis and you are very ambitious, you might try your hand at something more general.

\begin{description}

\item[Question 2.8:] $M: Ax^{2} + Bxy + Cy^{2} + Dx + Ey + F = 0$ can be used to describe all of the conic sections in the plane. See if you can show that no pointset with such an equation contains all of the points in our graph. You will note that there are six unknowns in this equation: $A$, $B$, $C$, $D$, $E$, and $F$. Hence, you will need at least one more number pair for your table. The algebra and the arithmetic involved in this should give you an excellent workout, and quite possibly a brain cramp.

  One of the observations that I have made over the years is that students really develop facility with their algebra while taking Trigonometry. If you are going to go further in mathematics, then being able to manipulate algebraically with confidence is going to be very important to you. You don't want to get bogged down in the algebra and miss the point of the discussion at hand.

\end{description}

Recall that our alien friend was unconvinced that we had solved the problem of calculating the coordinates of the point $P$ of arc $\arc{iI}$ such that $l(\arc{PI}) = \frac{\pi}{4}$. Here is a sequence of problems that you can pose to him that will convince him of the validity of our claim that $P:(\frac{\sqrt{2}}{2}, \frac{\sqrt{2}}{2})$. Of course, you had better work these yourself before you give them to him. You don't want to be embarrassed if he should ask for help. Study the statements of the problems themselves and draw pictures to see what they mean. If you'll do that, then you'll see how they fit together to solve the problem of the coordinates of $P$, above, very nicely. We state the problem as a theorem, which is a consequence of our axioms and definitions, and we list three helping theorems (lemmas) to prove the main theorem.

\begin{description}

\item[Theorem 2.1:] Suppose $P:(u,v)$ and $\arc{iPI}$. Then $u>0$ and $v>0$. You see, if this theorem is invoked, it settles the question that our alien friend has. He had two solutions for the coordinates of the point in question: $P:(\frac{\sqrt{2}}{2}, \frac{\sqrt{2}}{2})$, and the alternative $P:(-\frac{\sqrt{2}}{2}, -\frac{\sqrt{2}}{2})$, but only one can be valid according to the theorem.

\end{description}

The following problems, ``helping theorems'' called \emph{lemmas}, are designed to help you prove the preceding theorem, so work on these first.

\begin{description}

\item[Lemma 2.1.1:] Suppose $P:(u,v)$ and $\arc{iPI}$. Show that there is a point $Q$ on line $iI$ such that $PQO$ and $iQI$.

\item[Lemma 2.1.2:] Suppose $Q:(r,s)$ and $\arc{iQI}$. Show that $r>0$ and $s>0$. (You've already done this).

\item[Lemma 2.1.3:] Suppose $Q:(r,s)$, $r>0$, $s>0$, $P:(u,v)$, and $PQO$ is true. Show that $u>0$ and $v>0$.

\end{description}

A word about the problems, theorems, and questions that are sprinkled throughout this text. If you get stuck on one, see if there are some hints at the end of the chapter. If so, read the hints one at a time, going back to try the problem after each hint. Hopefully, it won't be necessary for you to use all of the hints on a problem. If, however, there are no hints for your problem, don't just become bogged down with that problem. Go on with the material, and every now and then, go back and try your sticky problem again. 

The point is, don't let yourself become stymied by a problem or set of problems. Assume the results of the problem and move on. I have gone back to some problems years after abandoning all hope, and have had a lot of success, probably because of increased mathematical maturity, which comes only with time and sweat and perseverance and sweat and frustration and sweat. Mostly sweat.

\vspace{0.3cm}

After you have digested the preceding lemmas and the theorem, see if you can construct three analogous lemmas to help you prove the following theorem.

\begin{description}

\item[Theorem 2.2:] Suppose $P:(u,v)$ and $\arc{JPi}$. Then $u<0$ and $v>0$.

\end{description}

Can you see that our picture of the upper semi-circle $\arc{JiI}$ suggests this result?

\vspace{0.3cm}

Here are some more problems that you might pose to our alien friend. They will all seem perfectly obvious to us, but he may challenge us to demonstrate them. These problems have to do with circles, lines, interior points of circles, and exterior points of circles. Recall that we assume the Pythagorean Theorem in the following form: If $A$, $B$, and $C$ are non-colinear points and line $AB$ is perpendicular to line $BC$, then $(AB)^{2} + (BC)^{2} = (AC)^{2}$. This was stated as Axiom 7.

\begin{description}

\item[Problem 2.9:] Suppose $ABC$, and line $DA$ is perpendicular to line $AC$. Show that $DA<DB<DC$.

\item[Problem 2.10:] Suppose $A$, $B$, and $C$ are non-colinear points. Show that there exists a point $N$ of line $AB$ which is nearer to $C$ than any other point of line $AB$, and that line $CN \bot$ line $AB$. This will test your ability to analyze quadratics. We'll call $N$ the \emph{near point} of line $AB$ to $C$.

\item[Problem 2.11:] Suppose $P$ is an interior point of $K:x^{2}+y^{2} = 1$, $T$ is a point of $K$, and $Q$ is a point such that $PTQ$. Show that $Q$ is an exterior point of $K$. Note: To say that $P$ is an interior point of $K$ means $PO<1$, and to say that $Q$ is an exterior point of $K$ means $QO>1$.

\item[Problem 2.12:] Suppose $P$ and $Q$ are points of $K:x^{2}+y^{2} = 1$, and $N$ is the near point of line $PQ$ to $O$. Show that $PNQ$, and that $N$ is an interior point of $K$.

\item[Problem 2.13:] Suppose $P$ and $Q$ are points of $K:x^{2}+y^{2} = 1$, and $PTQ$ is true. Show that $T$ is an interior point of $K$.

\item[Problem 2.14:] Theorem 2.1 states that, if $P:(u,v)$ and $\arc{iPI}$, then $u>0$ and $v>0$. Can you show that, if $P:(u,v)$ is a point of $\arc{JiI}$ and $u>0$ and $v>0$, then $\arc{iPI}$?

\item[Problem 2.15:] Show that, if $P:(u,v)$ is a point of $\arc{JiI}$ and $u<0$ and $v>0$, then $\arc{JPi}$. 

\end{description}

All of these will seem very obvious to you if you'll draw pictures of them. The point is that I want you to get used to the idea of questioning everything that is not an axiom or a statement that has already been proven; i.e., a theorem. 

This course is, above all else, a mathematics course, and, as such, the idea is to acquaint you with mathematical methods. The fact that we are talking about very familiar things, lines and circles, is not important. We might just as well be talking about boonfarkles and heebie-jeebies, as long as we had axioms that referred to these. I suspect that there would not be anything any more obvious to us about boonfarkles and heebie-jeebies than is obvious to our alien friend about lines and circles.

You'll notice that some of the problems stated here and elsewhere in the text are outlined as steps or hints at the ends of the chapters in which they occur. Please use these hints as a last resort only. Furthermore, as you read the hints for a problem that you got stuck on, read the hints one at a time. That is, after you read a hint, go back and try to solve the problem again. If you are still stuck, then go back and read the next hint. The object is for you to get as little help as necessary in the way of hints to work the problem. The only way to learn mathematics is to do mathematics. You may need the hints in the beginning, but as time wears on, perhaps you can wean yourself from them. Once you work some problems all on your own, the sense of accomplishment that you will feel will make you very reluctant to resort to using hints.

If you are comfortable with what's been done up to this point, then you are well prepared for the remainder. Establishing the ground rules and the foundation is always the most difficult, but also the most important, thing that we have to do.

\vspace{0.3cm}

Here's a puzzle for you which is a change of pace. This puzzle has nothing to do with the course, but you may find it interesting and challenging. Here is the puzzle:

Suppose that you start from some point on the Earth and walk 1 mile south, then 1 mile east, then 1 mile north, and you find that you are right back where you started from. The question is, where on earth did you start from? You can suppose a perfectly smooth Earth (borrowed from some physics department where they keep such things), and you can even walk on water, if you're inclined to do that sort of thing.

Now, here is the challenging part. Once you've found one such place, can you find another? In fact, can you find infinitely many more? One hint: There is only one such place north of the Equator, and you probably already found it. 

Don't spend so much time on this that you neglect your vitamins.



\newpage
\begin{center}
\section*{Chapter II Hints and Solutions}
\end{center}


\textbf{Question 2.1}:

\begin{center}
%TeXCAD Options
%\grade{\on}
%\emlines{\off}
%\epic{\off}
%\beziermacro{\on}
%\reduce{\on}
%\snapping{\off}
%\quality{8.00}
%\graddiff{0.01}
%\snapasp{1}
%\zoom{4.0000}
\unitlength 1mm % = 2.85pt
\linethickness{0.4pt}
\ifx\plotpoint\undefined\newsavebox{\plotpoint}\fi % GNUPLOT compatibility
\begin{picture}(115.75,54.5)(0,71)
\put(19,82.5){\line(1,0){82.5}}
\qbezier(59.5,122.75)(91.13,122.75)(105.25,82.75)
\qbezier(59.5,122.75)(22.25,120.38)(15,82.5)
\put(101.5,82.5){\line(1,0){14.25}}
\put(19,82.5){\line(-1,0){11.75}}
%\emline(7.25,82.5)(7,82.25)
\multiput(7.25,82.5)(-.03125,-.03125){8}{\line(0,-1){.03125}}
%\end
%\dashline{1}(15,82.75)(15,123)
\put(14.93,82.68){\line(0,1){.982}}
\put(14.93,84.64){\line(0,1){.982}}
\put(14.93,86.61){\line(0,1){.982}}
\put(14.93,88.57){\line(0,1){.982}}
\put(14.93,90.53){\line(0,1){.982}}
\put(14.93,92.5){\line(0,1){.982}}
\put(14.93,94.46){\line(0,1){.982}}
\put(14.93,96.42){\line(0,1){.982}}
\put(14.93,98.39){\line(0,1){.982}}
\put(14.93,100.35){\line(0,1){.982}}
\put(14.93,102.31){\line(0,1){.982}}
\put(14.93,104.28){\line(0,1){.982}}
\put(14.93,106.24){\line(0,1){.982}}
\put(14.93,108.2){\line(0,1){.982}}
\put(14.93,110.17){\line(0,1){.982}}
\put(14.93,112.13){\line(0,1){.982}}
\put(14.93,114.09){\line(0,1){.982}}
\put(14.93,116.06){\line(0,1){.982}}
\put(14.93,118.02){\line(0,1){.982}}
\put(14.93,119.98){\line(0,1){.982}}
\put(14.93,121.95){\line(0,1){.982}}
%\end
%\dashline{1}(58.75,122.75)(15,122.25)
\put(58.68,122.68){\line(-1,0){.994}}
\put(56.69,122.66){\line(-1,0){.994}}
\put(54.7,122.63){\line(-1,0){.994}}
\put(52.71,122.61){\line(-1,0){.994}}
\put(50.73,122.59){\line(-1,0){.994}}
\put(48.74,122.57){\line(-1,0){.994}}
\put(46.75,122.54){\line(-1,0){.994}}
\put(44.76,122.52){\line(-1,0){.994}}
\put(42.77,122.5){\line(-1,0){.994}}
\put(40.78,122.48){\line(-1,0){.994}}
\put(38.79,122.45){\line(-1,0){.994}}
\put(36.8,122.43){\line(-1,0){.994}}
\put(34.82,122.41){\line(-1,0){.994}}
\put(32.83,122.38){\line(-1,0){.994}}
\put(30.84,122.36){\line(-1,0){.994}}
\put(28.85,122.34){\line(-1,0){.994}}
\put(26.86,122.32){\line(-1,0){.994}}
\put(24.87,122.29){\line(-1,0){.994}}
\put(22.88,122.27){\line(-1,0){.994}}
\put(20.9,122.25){\line(-1,0){.994}}
\put(18.91,122.23){\line(-1,0){.994}}
\put(16.92,122.2){\line(-1,0){.994}}
%\end
%\dashline{1}(58.75,122.75)(103.75,123)
\put(58.68,122.68){\line(1,0){.978}}
\put(60.64,122.69){\line(1,0){.978}}
\put(62.59,122.7){\line(1,0){.978}}
\put(64.55,122.71){\line(1,0){.978}}
\put(66.51,122.72){\line(1,0){.978}}
\put(68.46,122.73){\line(1,0){.978}}
\put(70.42,122.74){\line(1,0){.978}}
\put(72.38,122.76){\line(1,0){.978}}
\put(74.33,122.77){\line(1,0){.978}}
\put(76.29,122.78){\line(1,0){.978}}
\put(78.24,122.79){\line(1,0){.978}}
\put(80.2,122.8){\line(1,0){.978}}
\put(82.16,122.81){\line(1,0){.978}}
\put(84.11,122.82){\line(1,0){.978}}
\put(86.07,122.83){\line(1,0){.978}}
\put(88.03,122.84){\line(1,0){.978}}
\put(89.98,122.85){\line(1,0){.978}}
\put(91.94,122.86){\line(1,0){.978}}
\put(93.9,122.88){\line(1,0){.978}}
\put(95.85,122.89){\line(1,0){.978}}
\put(97.81,122.9){\line(1,0){.978}}
\put(99.77,122.91){\line(1,0){.978}}
\put(101.72,122.92){\line(1,0){.978}}
%\end
%\dashline{1}(103.75,123)(105.25,83)
\put(103.68,122.93){\line(0,-1){.976}}
\put(103.75,120.98){\line(0,-1){.976}}
\put(103.83,119.03){\line(0,-1){.976}}
\put(103.9,117.08){\line(0,-1){.976}}
\put(103.97,115.12){\line(0,-1){.976}}
\put(104.05,113.17){\line(0,-1){.976}}
\put(104.12,111.22){\line(0,-1){.976}}
\put(104.19,109.27){\line(0,-1){.976}}
\put(104.27,107.32){\line(0,-1){.976}}
\put(104.34,105.37){\line(0,-1){.976}}
\put(104.41,103.42){\line(0,-1){.976}}
\put(104.48,101.47){\line(0,-1){.976}}
\put(104.56,99.52){\line(0,-1){.976}}
\put(104.63,97.56){\line(0,-1){.976}}
\put(104.7,95.61){\line(0,-1){.976}}
\put(104.78,93.66){\line(0,-1){.976}}
\put(104.85,91.71){\line(0,-1){.976}}
\put(104.92,89.76){\line(0,-1){.976}}
\put(105,87.81){\line(0,-1){.976}}
\put(105.07,85.86){\line(0,-1){.976}}
\put(105.14,83.91){\line(0,-1){.976}}
%\end
\put(58.75,124.75){\makebox(0,0)[cc]{$i:(0,1)$}}
\put(59.25,79){\makebox(0,0)[cc]{$O:(0,0)$}}
\put(104.5,79.75){\makebox(0,0)[cc]{$I:(1,0)$}}
\put(15,79.75){\makebox(0,0)[cc]{$J:(-1,0)$}}
\put(60,72.75){\makebox(0,0)[cc]{Figure 2.1}}
\put(107,124){\makebox(0,0)[cc]{$(1,1)$}}
\put(12,124){\makebox(0,0)[cc]{$(-1,1)$}}
\put(58.75,122.5){\circle{.5}}
\end{picture}
\end{center}

$JI = 2$ and $2<\pi$, at least it looks that way in the picture. Can you see a path of length $4$ that is greater than $\pi$?

\vspace{0.3cm}

\textbf{Problem 2.3}: $P:(u,v)$ and $u^{2} + v^{2} = 1$. Let $Q:(r,s)$.
\begin{enumerate}[\hspace{1cm}1)]
  \item $POQ$ implies $PO+OQ = PQ$.
  \item Translate Step 1 into an algebraic statement using the coordinates of $P$, $Q$, and $O$.
  \item Solve for $r$ and $s$. Recall that $Q:(r,s)$ is a point of $K$, so $r^{2} + s^{2} = 1$.
\end{enumerate}

\textbf{Problem 2.4}: Suppose $S$ is a semi-circle of $K$.
\begin{enumerate}[\hspace{1cm}1)]
  \item Suppose $S = \arc{PQ}$, where $P$ and $Q$ are points of $K$. Let $P:(u,v)$. Draw the picture.
  \item What are the coordinates of $Q$?
  \item Suppose line $PQ$ is not vertical. Get an equation for line $PQ$. Note that $POQ$.
  \item Suppose $R:(r,s)$ is a point of $K$ and that line $RO \bot$ line $PQ$.
  \item Find $r$ and $s$ in terms of $u$ and $v$.
  \item Show that $PR = RQ = iI$.
  \item Conclude that $l(\arc{PR}) = l(\arc{RQ}) = l(\arc{iI})$.
  \item Note that $l(S) = l(\arc{PRQ}) = l(\arc{PR}) + l(\arc{RQ}) = \pi$
\end{enumerate}

\textbf{Problems 2.5, 2.6, and 2.7}: These are very much like Problem 2.4.

\vspace{0.3cm}

\textbf{Problem 2.8}:
\begin{enumerate}
  \item[\hspace{1cm}1)] Illustration:
\end{enumerate}
\begin{center}
%TeXCAD Picture [prob2.8hint.pic]. Options:
%\grade{\on}
%\emlines{\off}
%\epic{\off}
%\beziermacro{\on}
%\reduce{\on}
%\snapping{\off}
%\quality{8.00}
%\graddiff{0.01}
%\snapasp{1}
%\zoom{4.0000}
\unitlength 1mm % = 2.85pt
\linethickness{0.4pt}
\ifx\plotpoint\undefined\newsavebox{\plotpoint}\fi % GNUPLOT compatibility
\begin{picture}(112.5,32.5)(0,0)
\put(62.5,7.5){\line(1,0){50}}
\put(1.64,7.5){\line(1,0){50}}
\qbezier(4.25,7.5)(4.38,31.63)(27,29.25)
\qbezier(65,7.5)(65.13,31.63)(87.75,29.25)
\qbezier(48.75,7.5)(48.63,31.63)(26,29.25)
\qbezier(109.5,7.5)(109.38,31.63)(86.75,29.25)
\put(4.25,4.5){\makebox(0,0)[cc]{$J:(-1,0)$}}
\put(65,4.5){\makebox(0,0)[cc]{$J:(-1,0)$}}
\put(48.25,4.5){\makebox(0,0)[cc]{$I:(1,0)$}}
\put(109,4.5){\makebox(0,0)[cc]{$I:(1,0)$}}
\put(26.5,4.5){\makebox(0,0)[cc]{$O:(0,0)$}}
\put(87.25,4.5){\makebox(0,0)[cc]{$O:(0,0)$}}
\put(26.5,32.5){\makebox(0,0)[cc]{$i:(0,1)$}}
\put(87.25,32.5){\makebox(0,0)[cc]{$i:(0,1)$}}
\put(53,26){\makebox(0,0)[cc]{$R:(\frac{\sqrt{2}}{2},\frac{\sqrt{2}}{2})$}}
\put(105.75,25.75){\makebox(0,0)[cc]{$R$}}
\put(111,20.5){\makebox(0,0)[cc]{$P$}}
\put(99.75,31.5){\makebox(0,0)[cc]{$Q$}}
%\emline(4.25,7.5)(9.5,24.5)
\multiput(4.25,7.5)(.03365385,.10897436){156}{\line(0,1){.10897436}}
%\end
\put(9.5,24.5){\line(0,1){0}}
%\emline(9.5,24.5)(26.5,29.25)
\multiput(9.5,24.5)(.12056738,.03368794){141}{\line(1,0){.12056738}}
%\end
%\emline(26.5,29.25)(43.75,24)
\multiput(26.5,29.25)(.11057692,-.03365385){156}{\line(1,0){.11057692}}
%\end
%\emline(43.75,24)(48.25,8)
\multiput(43.75,24)(.03358209,-.11940299){134}{\line(0,-1){.11940299}}
%\end
\put(65,7.5){\line(0,1){6.75}}
\put(65,14.25){\line(0,1){5.5}}
\put(75.5,29.5){\line(1,0){23.75}}
\put(99.25,30){\line(0,1){0}}
\put(109.25,19.5){\line(0,1){0}}
%\emline(109.25,19.5)(109.5,7.25)
\multiput(109.25,19.5)(.03125,-1.53125){8}{\line(0,-1){1.53125}}
%\end
\put(109.5,7.25){\line(0,1){0}}
%\emline(99,29.5)(109.25,19.5)
\multiput(99,29.5)(.034511785,-.033670034){297}{\line(1,0){.034511785}}
%\end
\put(109.25,19.5){\line(-1,0){.25}}
\put(109,19.5){\line(0,1){0}}
%\emline(75.5,29.5)(65.25,19.5)
\multiput(75.5,29.5)(-.034511785,-.033670034){297}{\line(-1,0){.034511785}}
%\end
\put(26.25,29.25){\circle*{1.12}}
\put(4.25,7.5){\circle*{1.12}}
\put(26.64,7.5){\circle*{1.12}}
\put(48.92,7.61){\circle*{1.12}}
\put(65.11,7.61){\circle*{1.12}}
\put(43.9,24.05){\circle*{1.12}}
\put(87.25,29.5){\circle*{1.41}}
\put(99,29.5){\circle*{1}}
\put(103.75,25){\circle*{1.12}}
\put(109.25,19.5){\circle*{1}}
\put(109.5,7.5){\circle*{1}}
\put(87.25,7.5){\circle*{1.12}}
\end{picture}
\end{center}

\begin{enumerate}
  \item[\hspace{1cm}2)] Note $4(RI) < \pi < 8(PI)$.
  \item[\hspace{1cm}3)] You know the coordinates of $R$. Calculate the coordinates of $P$ and $Q$ by getting the equation of line $PQ$.
  \item[\hspace{1cm}4)] Note that 3 $<$ 4(RI) and 8(PI) $<$ 3.5.
\end{enumerate}

\textbf{Question 2.5}:
\begin{enumerate}[\hspace{1cm}1)]
  \item Suppose the relationship between arc length and chord length is circular. Then there is a circle $K: (x-c_{1})^{2} + (y-c_{2})^{2} = r^{2}$ which contains all of the points associated with the number pairs in our table.
  \item Fill out the Table 2.1 in its entirety; i.e., fill in all nine number pairs.
  \item Note that $(0,0)$ is a number pair in the table, and that it means that $c_{1}^{2} + c_{2}^{2} = r^{2}$.
  \item Put $(\frac{\pi}{3}, 1)$ into your equation for $K$, and put $(\pi,2)$ into your equation for $K$. This should yield you two equations involving $c_{1}$ and $c_{2}$. Solve for one of them.
  \item Now take two other pairs, say $(\frac{\pi}{3},1)$ and $(\frac{\pi}{2}, \sqrt{2})$ and put them in your equation for $K$ yielding two other equations involving $c_{1}$ and $c_{2}$. Solve these for the same one you solved for in step 4. You probably get conflicting results. What does it mean if you did get, for example, $c_{1} = x$ in one case and $c_{1} = y$ in the other, where $x \neq y$?
\end{enumerate}

\textbf{Question 2.6}:
\begin{enumerate}[\hspace{1cm}1)]
  \item Suppose the relationship between arc length and chord length is $\ $ parabolic, and the points associated with the number pairs in Table 2.1 fit the equation $y = Ax^{2} + Bx + C$.
  \item Proceed as in the case of the circle. You should conclude that no such number triple $(A,B,C)$ exists such that \emph{all} the number pairs satisfy the equation.
\end{enumerate}

\textbf{Question 2.7}:
\begin{enumerate}[\hspace{1cm}1)]
  \item Suppose $Q$ is a point of semi-circle $\arc{JiI}$, $Q \neq P:(\frac{\sqrt{2}}{2}, \frac{\sqrt{2}}{2})$, and $l(\arc{QI}) = \frac{\pi}{4}$. Note that we already know that $l(\arc{PI}) = \frac{\pi}{4}$.
  \item Semi-circle $\arc{JiI}$ is semi-circle $\arc{JPI}$. Hence, $Q$ is a point of semi-circle $\arc{JPI}$.
  \item From Step 2, either $\arc{JQPI}$ or $\arc{JPQI}$.
  \item Show that both possibilities in Step 3 result in a contradiction.
  \item Your conclusion is that the assumption that $l(\arc{QI}) = \frac{\pi}{4}$ and $Q \neq P:(\frac{\sqrt{2}}{2}, \frac{\sqrt{2}}{2})$ leads to a contradiction. Hence, your assumption is false. This means there is only one point of semi-circle $\arc{JiI}$ whose arc length to $I$ is $\frac{\pi}{4}$.
\end{enumerate}

\textbf{Lemma 2.1.1}: (This is Lemma 1 for Theorem 2.1. Hence, the name is Lemma 2.1.1).
\begin{enumerate}[\hspace{1cm}1)]
  \item Note the definition of $P$ being a point of arc $\arc{iI}$.
  \item You now have a point $Q$ of line $iI$ such that $PQO$. But where is $Q$ on line $iI$?
  \item Show $Q \neq i$ and $Q \neq I$. Where else can $Q$ be on line $iI$?
  \item Show $QiI$ is false. Note that $QiI$ implies $Qi + iI = QI$ and you can get an algebraic statement out of this.
  \item Show $iIQ$ is false also. What's left?
\end{enumerate}

\textbf{Lemma 2.1.3}:
\begin{enumerate}[\hspace{1cm}1)]
  \item $PQO$ implies $PQ + QO = PO$.
  \item From Step 1, $\sqrt{(u-r)^{2} + (v-s)^{2}} + \sqrt{r^{2} + s^{2}} = \sqrt{u^{2} + v^{2}}$.
  \item Since $PQO$, line $PQ: y = mx$. Hence, $v = mu$ and $s = mr$.
  \item Incorporating Step 3 in Step 2, you should get 

      $\sqrt{(u-r)^{2}(1+m^{2})} + \sqrt{r^{2}(1+m^{2})} = \sqrt{u^{2}(1+m^{2})}$
  \item Step 4 reduces to $\vert u-r \vert  +  \vert r \vert  =  \vert u \vert$.
  \item With appropriate algebraic chicanery, this will reduce to $ur = \vert u \vert  \vert r \vert  \geq  0$. From this, since $r>0$, you can conclude that $u>0$.
\end{enumerate}

\textbf{Theorem 2.1}:
\begin{enumerate}[\hspace{1cm}1)]
  \item Use Lemma 2.1.1.
  \item Use Lemma 2.1.2.
  \item Use Lemma 2.1.3.
\end{enumerate}

\textbf{Problem 2.9}:
\begin{enumerate}[\hspace{1cm}1)]
  \item Draw the picture.
  \item The picture should look something like the following:
\end{enumerate}

\begin{center}
%TeXCAD Options
%\grade{\on}
%\emlines{\off}
%\epic{\off}
%\beziermacro{\on}
%\reduce{\on}
%\snapping{\off}
%\quality{8.00}
%\graddiff{0.01}
%\snapasp{1}
%\zoom{4.0000}
\unitlength 1mm % = 2.85pt
\linethickness{0.4pt}
\ifx\plotpoint\undefined\newsavebox{\plotpoint}\fi % GNUPLOT compatibility
\begin{picture}(63,37.75)(0,103)
\put(15,137.5){\line(0,-1){30.75}} \put(15,106.75){\line(1,0){46}}
%\emline(61,106.75)(15.25,137.75)
\multiput(61,106.75)(-.0497823721,.0337323177){919}{\line(-1,0){.0497823721}}
%\end
%\emline(15.25,137.75)(40,107.25)
\multiput(15.25,137.75)(.033719346,-.0415531335){734}{\line(0,-1){.0415531335}}
%\end
%\emline(40,107.25)(40.25,106.75)
\multiput(40,107.25)(.03125,-.0625){8}{\line(0,-1){.0625}}
%\end
\put(15.25,110.5){\line(1,0){2.75}} \put(18,110.5){\line(0,-1){3.5}}
\put(13,136){\makebox(0,0)[cc]{$D$}}
\put(63,104){\makebox(0,0)[cc]{$C$}}
\put(39.5,104){\makebox(0,0)[cc]{$B$}}
\put(13,104){\makebox(0,0)[cc]{$A$}}
\end{picture}
\end{center}

The little $\square$ near point $A$ means a right angle; i.e., angle $DAB$ is a right angle. Use Axiom 7 (the Pythagorean Theorem) to get your result. Note $AB + BC = AC$ since $ABC$.

\textbf{Problem 2.10}:
\begin{enumerate}[\hspace{1cm}1)]
  \item Draw the picture.
  \item Your picture could look something like the following:
\end{enumerate}

\begin{center}
%TeXCAD Options
%\grade{\on}
%\emlines{\off}
%\epic{\off}
%\beziermacro{\on}
%\reduce{\on}
%\snapping{\off}
%\quality{8.00}
%\graddiff{0.01}
%\snapasp{1}
%\zoom{4.0000}
\unitlength 1mm % = 2.85pt
\linethickness{0.4pt}
\ifx\plotpoint\undefined\newsavebox{\plotpoint}\fi % GNUPLOT compatibility
\begin{picture}(81.5,43.75)(0,100)
%\emline(15.5,106.5)(61.5,141.25)
\multiput(15.5,106.5)(.0446601942,.0337378641){1030}{\line(1,0){.0446601942}}
%\end
\put(61.5,141.25){\line(1,-1){18}}
%\emline(58.75,139.25)(61.5,136.5)
\multiput(58.75,139.25)(.0335366,-.0335366){82}{\line(0,-1){.0335366}}
%\end
%\emline(61.5,136.5)(64,138.75)
\multiput(61.5,136.5)(.0373134,.0335821){67}{\line(1,0){.0373134}}
%\end
\put(77.5,125.25){\circle{.71}} 
\put(22.75,111.75){\circle{.5}}
\put(43.75,127.75){\circle{.5}}
\put(20.75,113.5){\makebox(0,0)[cc]{$A$}}
\put(40.75,130.5){\makebox(0,0)[cc]{$B$}}
\put(81.5,127.25){\makebox(0,0)[cc]{$C$}}
\put(61.5,143.75){\makebox(0,0)[cc]{$N$}}
\end{picture}
\end{center}

Now suppose line $AB$ is not vertical. Then line $AB: y = mx + b$. Let $C:(c_{1}, c_{2})$.

\begin{enumerate}
  \item[\hspace{1cm}3)] If $P:(x,y)$ is a point of line $AB$, then $PC = \sqrt{(x-c_{1})^{2} + (y-c_{2})^{2}}$.
  \item[\hspace{1cm}4)] $(PC)^{2} = (x-c_{1})^{2} + (y-c_{2})^{2}$, which is a quadratic in $x$ if you replace $y$ by $mx + b$.
  \item[\hspace{1cm}5)] Recall that the minimum of the quadratic $y = Ax^{2} + Bx + C$ if $A>0$ occurs where $x = -\frac{B}{2A}$.
  \item[\hspace{1cm}6)] Find $y$ where $x = -\frac{B}{2A}$, and let $N:(x,y)$. $N$ is the nearest point of line $AB$ to $C$.
  \item[\hspace{1cm}7)] You now have the coordinates of $N$ in terms of $m$, $b$, $c_{1}$, and $c_{2}$. Find the slope of line $NC$.
  \item[\hspace{1cm}8)] The slope of line $NC$ should be $-\frac{1}{m}$, which means line $NC$ is perpendicular to line $AB$ (which has slope $m$).
\end{enumerate}

\textbf{Problem 2.11}:
\begin{enumerate}[\hspace{1cm}1)]
  \item Draw the picture.
  \item Your picture could look something like the following:
\end{enumerate}

\begin{center}
%TeXCAD Options
%\grade{\on}
%\emlines{\off}
%\epic{\off}
%\beziermacro{\on}
%\reduce{\on}
%\snapping{\off}
%\quality{8.00}
%\graddiff{0.01}
%\snapasp{1}
%\zoom{4.0000}
\unitlength 1mm % = 2.85pt
\linethickness{0.4pt}
\ifx\plotpoint\undefined\newsavebox{\plotpoint}\fi % GNUPLOT compatibility
\begin{picture}(96.5,73.75)(0,55)
%\circle(53,90.25){68.48}
\put(87.24,90.25){\line(0,1){1.308}}
\put(87.21,91.56){\line(0,1){1.306}}
\multiput(87.14,92.86)(-.0312,.3255){4}{\line(0,1){.3255}}
\multiput(87.01,94.17)(-.02906,.21603){6}{\line(0,1){.21603}}
\multiput(86.84,95.46)(-.03196,.18409){7}{\line(0,1){.18409}}
\multiput(86.62,96.75)(-.03031,.14212){9}{\line(0,1){.14212}}
\multiput(86.34,98.03)(-.03214,.12678){10}{\line(0,1){.12678}}
\multiput(86.02,99.3)(-.0336,.11405){11}{\line(0,1){.11405}}
\multiput(85.65,100.55)(-.032096,.095349){13}{\line(0,1){.095349}}
\multiput(85.24,101.79)(-.033163,.087335){14}{\line(0,1){.087335}}
\multiput(84.77,103.01)(-.031915,.075254){16}{\line(0,1){.075254}}
\multiput(84.26,104.22)(-.032721,.069629){17}{\line(0,1){.069629}}
\multiput(83.7,105.4)(-.033392,.064532){18}{\line(0,1){.064532}}
\multiput(83.1,106.56)(-.032249,.056889){20}{\line(0,1){.056889}}
\multiput(82.46,107.7)(-.03276,.052967){21}{\line(0,1){.052967}}
\multiput(81.77,108.81)(-.033179,.049329){22}{\line(0,1){.049329}}
\multiput(81.04,109.9)(-.033516,.045937){23}{\line(0,1){.045937}}
\multiput(80.27,110.95)(-.032426,.041054){25}{\line(0,1){.041054}}
\multiput(79.46,111.98)(-.032664,.038255){26}{\line(0,1){.038255}}
\multiput(78.61,112.98)(-.032838,.03561){27}{\line(0,1){.03561}}
\multiput(77.72,113.94)(-.032953,.033104){28}{\line(0,1){.033104}}
\multiput(76.8,114.86)(-.03546,.033){27}{\line(-1,0){.03546}}
\multiput(75.84,115.76)(-.038106,.032838){26}{\line(-1,0){.038106}}
\multiput(74.85,116.61)(-.040905,.032613){25}{\line(-1,0){.040905}}
\multiput(73.83,117.42)(-.045784,.033725){23}{\line(-1,0){.045784}}
\multiput(72.78,118.2)(-.049177,.033404){22}{\line(-1,0){.049177}}
\multiput(71.69,118.94)(-.052817,.033002){21}{\line(-1,0){.052817}}
\multiput(70.59,119.63)(-.056741,.032509){20}{\line(-1,0){.056741}}
\multiput(69.45,120.28)(-.064379,.033686){18}{\line(-1,0){.064379}}
\multiput(68.29,120.88)(-.069479,.033039){17}{\line(-1,0){.069479}}
\multiput(67.11,121.45)(-.075108,.032259){16}{\line(-1,0){.075108}}
\multiput(65.91,121.96)(-.087183,.033562){14}{\line(-1,0){.087183}}
\multiput(64.69,122.43)(-.095201,.032531){13}{\line(-1,0){.095201}}
\multiput(63.45,122.86)(-.10441,.03128){12}{\line(-1,0){.10441}}
\multiput(62.2,123.23)(-.12663,.03272){10}{\line(-1,0){.12663}}
\multiput(60.93,123.56)(-.14198,.03096){9}{\line(-1,0){.14198}}
\multiput(59.65,123.84)(-.18394,.0328){7}{\line(-1,0){.18394}}
\multiput(58.37,124.07)(-.2159,.03004){6}{\line(-1,0){.2159}}
\multiput(57.07,124.25)(-.3253,.0327){4}{\line(-1,0){.3253}}
\put(55.77,124.38){\line(-1,0){1.305}}
\put(54.46,124.46){\line(-1,0){1.308}}
\put(53.16,124.49){\line(-1,0){1.308}}
\put(51.85,124.47){\line(-1,0){1.306}}
\put(50.54,124.4){\line(-1,0){1.302}}
\multiput(49.24,124.28)(-.2594,-.03369){5}{\line(-1,0){.2594}}
\multiput(47.94,124.11)(-.18423,-.03112){7}{\line(-1,0){.18423}}
\multiput(46.65,123.9)(-.16004,-.03337){8}{\line(-1,0){.16004}}
\multiput(45.37,123.63)(-.12692,-.03156){10}{\line(-1,0){.12692}}
\multiput(44.1,123.31)(-.1142,-.03308){11}{\line(-1,0){.1142}}
\multiput(42.85,122.95)(-.095494,-.031661){13}{\line(-1,0){.095494}}
\multiput(41.61,122.54)(-.087486,-.032764){14}{\line(-1,0){.087486}}
\multiput(40.38,122.08)(-.080426,-.033676){15}{\line(-1,0){.080426}}
\multiput(39.18,121.57)(-.069777,-.032403){17}{\line(-1,0){.069777}}
\multiput(37.99,121.02)(-.064684,-.033097){18}{\line(-1,0){.064684}}
\multiput(36.82,120.43)(-.060037,-.033673){19}{\line(-1,0){.060037}}
\multiput(35.68,119.79)(-.053116,-.032518){21}{\line(-1,0){.053116}}
\multiput(34.57,119.1)(-.04948,-.032954){22}{\line(-1,0){.04948}}
\multiput(33.48,118.38)(-.04609,-.033306){23}{\line(-1,0){.04609}}
\multiput(32.42,117.61)(-.042918,-.033581){24}{\line(-1,0){.042918}}
\multiput(31.39,116.81)(-.038404,-.032489){26}{\line(-1,0){.038404}}
\multiput(30.39,115.96)(-.03576,-.032675){27}{\line(-1,0){.03576}}
\multiput(29.43,115.08)(-.033254,-.032802){28}{\line(-1,0){.033254}}
\multiput(28.49,114.16)(-.033162,-.035309){27}{\line(0,-1){.035309}}
\multiput(27.6,113.21)(-.033012,-.037956){26}{\line(0,-1){.037956}}
\multiput(26.74,112.22)(-.032799,-.040756){25}{\line(0,-1){.040756}}
\multiput(25.92,111.2)(-.03252,-.043728){24}{\line(0,-1){.043728}}
\multiput(25.14,110.15)(-.033628,-.049024){22}{\line(0,-1){.049024}}
\multiput(24.4,109.08)(-.033243,-.052666){21}{\line(0,-1){.052666}}
\multiput(23.7,107.97)(-.032767,-.056592){20}{\line(0,-1){.056592}}
\multiput(23.05,106.84)(-.032192,-.060844){19}{\line(0,-1){.060844}}
\multiput(22.44,105.68)(-.033356,-.069327){17}{\line(0,-1){.069327}}
\multiput(21.87,104.5)(-.032601,-.07496){16}{\line(0,-1){.07496}}
\multiput(21.35,103.3)(-.031696,-.081227){15}{\line(0,-1){.081227}}
\multiput(20.87,102.08)(-.032966,-.095052){13}{\line(0,-1){.095052}}
\multiput(20.44,100.85)(-.03175,-.10426){12}{\line(0,-1){.10426}}
\multiput(20.06,99.6)(-.0333,-.12648){10}{\line(0,-1){.12648}}
\multiput(19.73,98.33)(-.0316,-.14184){9}{\line(0,-1){.14184}}
\multiput(19.44,97.06)(-.03364,-.18379){7}{\line(0,-1){.18379}}
\multiput(19.21,95.77)(-.03103,-.21576){6}{\line(0,-1){.21576}}
\multiput(19.02,94.48)(-.02732,-.26014){5}{\line(0,-1){.26014}}
\put(18.89,93.18){\line(0,-1){1.305}}
\put(18.8,91.87){\line(0,-1){1.307}}
\put(18.76,90.56){\line(0,-1){1.308}}
\put(18.78,89.25){\line(0,-1){1.306}}
\put(18.84,87.95){\line(0,-1){1.303}}
\multiput(18.95,86.65)(.0325,-.25955){5}{\line(0,-1){.25955}}
\multiput(19.11,85.35)(.03028,-.18437){7}{\line(0,-1){.18437}}
\multiput(19.33,84.06)(.03264,-.16019){8}{\line(0,-1){.16019}}
\multiput(19.59,82.78)(.03098,-.12706){10}{\line(0,-1){.12706}}
\multiput(19.9,81.5)(.03256,-.11435){11}{\line(0,-1){.11435}}
\multiput(20.25,80.25)(.031224,-.095638){13}{\line(0,-1){.095638}}
\multiput(20.66,79)(.032365,-.087634){14}{\line(0,-1){.087634}}
\multiput(21.11,77.78)(.033309,-.080579){15}{\line(0,-1){.080579}}
\multiput(21.61,76.57)(.032084,-.069925){17}{\line(0,-1){.069925}}
\multiput(22.16,75.38)(.032801,-.064835){18}{\line(0,-1){.064835}}
\multiput(22.75,74.21)(.033398,-.060191){19}{\line(0,-1){.060191}}
\multiput(23.38,73.07)(.032275,-.053264){21}{\line(0,-1){.053264}}
\multiput(24.06,71.95)(.032727,-.04963){22}{\line(0,-1){.04963}}
\multiput(24.78,70.86)(.033095,-.046242){23}{\line(0,-1){.046242}}
\multiput(25.54,69.79)(.033385,-.043071){24}{\line(0,-1){.043071}}
\multiput(26.34,68.76)(.033605,-.040094){25}{\line(0,-1){.040094}}
\multiput(27.18,67.76)(.032511,-.035909){27}{\line(0,-1){.035909}}
\multiput(28.06,66.79)(.03265,-.033404){28}{\line(0,-1){.033404}}
\multiput(28.98,65.85)(.035157,-.033323){27}{\line(1,0){.035157}}
\multiput(29.93,64.95)(.037804,-.033185){26}{\line(1,0){.037804}}
\multiput(30.91,64.09)(.040606,-.032985){25}{\line(1,0){.040606}}
\multiput(31.92,63.27)(.043579,-.032719){24}{\line(1,0){.043579}}
\multiput(32.97,62.48)(.046745,-.03238){23}{\line(1,0){.046745}}
\multiput(34.04,61.74)(.052514,-.033483){21}{\line(1,0){.052514}}
\multiput(35.15,61.03)(.056442,-.033025){20}{\line(1,0){.056442}}
\multiput(36.28,60.37)(.060697,-.032469){19}{\line(1,0){.060697}}
\multiput(37.43,59.76)(.069174,-.033672){17}{\line(1,0){.069174}}
\multiput(38.61,59.18)(.07481,-.032943){16}{\line(1,0){.07481}}
\multiput(39.8,58.66)(.081081,-.032066){15}{\line(1,0){.081081}}
\multiput(41.02,58.18)(.0949,-.033399){13}{\line(1,0){.0949}}
\multiput(42.25,57.74)(.10412,-.03223){12}{\line(1,0){.10412}}
\multiput(43.5,57.35)(.11484,-.0308){11}{\line(1,0){.11484}}
\multiput(44.76,57.02)(.1417,-.03225){9}{\line(1,0){.1417}}
\multiput(46.04,56.73)(.16068,-.03017){8}{\line(1,0){.16068}}
\multiput(47.33,56.48)(.21562,-.03201){6}{\line(1,0){.21562}}
\multiput(48.62,56.29)(.26002,-.02851){5}{\line(1,0){.26002}}
\put(49.92,56.15){\line(1,0){1.305}}
\put(51.22,56.06){\line(1,0){1.307}}
\put(52.53,56.01){\line(1,0){1.308}}
\put(53.84,56.02){\line(1,0){1.307}}
\put(55.15,56.08){\line(1,0){1.304}}
\multiput(56.45,56.19)(.25969,.03132){5}{\line(1,0){.25969}}
\multiput(57.75,56.34)(.18451,.02944){7}{\line(1,0){.18451}}
\multiput(59.04,56.55)(.16034,.0319){8}{\line(1,0){.16034}}
\multiput(60.32,56.8)(.1272,.0304){10}{\line(1,0){.1272}}
\multiput(61.59,57.11)(.1145,.03204){11}{\line(1,0){.1145}}
\multiput(62.85,57.46)(.10376,.03335){12}{\line(1,0){.10376}}
\multiput(64.1,57.86)(.087781,.031964){14}{\line(1,0){.087781}}
\multiput(65.33,58.31)(.08073,.03294){15}{\line(1,0){.08073}}
\multiput(66.54,58.8)(.07007,.031764){17}{\line(1,0){.07007}}
\multiput(67.73,59.34)(.064984,.032505){18}{\line(1,0){.064984}}
\multiput(68.9,59.93)(.060342,.033123){19}{\line(1,0){.060342}}
\multiput(70.05,60.56)(.056082,.033633){20}{\line(1,0){.056082}}
\multiput(71.17,61.23)(.049778,.032501){22}{\line(1,0){.049778}}
\multiput(72.26,61.94)(.046392,.032883){23}{\line(1,0){.046392}}
\multiput(73.33,62.7)(.043223,.033188){24}{\line(1,0){.043223}}
\multiput(74.37,63.5)(.040247,.033422){25}{\line(1,0){.040247}}
\multiput(75.37,64.33)(.037444,.033591){26}{\line(1,0){.037444}}
\multiput(76.35,65.21)(.034795,.0337){27}{\line(1,0){.034795}}
\multiput(77.29,66.11)(.033483,.035005){27}{\line(0,1){.035005}}
\multiput(78.19,67.06)(.033357,.037653){26}{\line(0,1){.037653}}
\multiput(79.06,68.04)(.03317,.040455){25}{\line(0,1){.040455}}
\multiput(79.89,69.05)(.032918,.043429){24}{\line(0,1){.043429}}
\multiput(80.68,70.09)(.032593,.046596){23}{\line(0,1){.046596}}
\multiput(81.43,71.16)(.033722,.05236){21}{\line(0,1){.05236}}
\multiput(82.13,72.26)(.033283,.05629){20}{\line(0,1){.05629}}
\multiput(82.8,73.39)(.032746,.060548){19}{\line(0,1){.060548}}
\multiput(83.42,74.54)(.032099,.065185){18}{\line(0,1){.065185}}
\multiput(84,75.71)(.033284,.074659){16}{\line(0,1){.074659}}
\multiput(84.53,76.91)(.032436,.080934){15}{\line(0,1){.080934}}
\multiput(85.02,78.12)(.031416,.087979){14}{\line(0,1){.087979}}
\multiput(85.46,79.35)(.0327,.10397){12}{\line(0,1){.10397}}
\multiput(85.85,80.6)(.03132,.1147){11}{\line(0,1){.1147}}
\multiput(86.2,81.86)(.0329,.14155){9}{\line(0,1){.14155}}
\multiput(86.49,83.14)(.0309,.16054){8}{\line(0,1){.16054}}
\multiput(86.74,84.42)(.033,.21547){6}{\line(0,1){.21547}}
\multiput(86.94,85.71)(.02969,.25988){5}{\line(0,1){.25988}}
\put(87.09,87.01){\line(0,1){1.304}}
\put(87.18,88.32){\line(0,1){1.932}}
%\end
%\emline(4.75,82.75)(84.5,128.75)
\multiput(4.75,82.75)(.0584677419,.0337243402){1364}{\line(1,0){.0584677419}}
%\end
%\qbezvec(96.5,97)(94.5,91.25)(88.5,90.5)
\put(88.5,90.5){\vector(-4,-1){.07}}\qbezier(96.5,97)(94.5,91.25)(88.5,90.5)
%\end
\put(61.25,115.25){\circle{.5}} 
\put(11,86.25){\circle{.5}}
\put(63.25,113){\makebox(0,0)[cc]{$P$}}
\put(10.75,83.75){\makebox(0,0)[cc]{$Q$}}
\put(16.5,92){\makebox(0,0)[cc]{$T$}}
\put(53.75,86.5){\makebox(0,0)[cc]{$O$}}
\put(99,99.25){\makebox(0,0)[cc]{$K:x^{2}+y^{2}=1$}}
\end{picture}
\end{center}

\begin{enumerate}
  \item[\hspace{1cm}3)] Suppose $N$ is the near point of line $PT$ to $O$.
  \item[\hspace{1cm}4)] Recognize that one of the following must be true: $N = P$, $N = T$, $N = Q$, $NPTQ$, $PNTQ$, $PTNQ$, or $PTQN$.
  \item[\hspace{1cm}5)] In each of the cases above, draw the picture and use the results of Problem 2.9.
  \item[\hspace{1cm}6)] You should have found that some of the cases above are impossible, in that they contradict something that you know to be true. In all the other cases, you should have the result that $OT < OQ$, implying that $Q$ is an exterior point of $K$.
\end{enumerate}

\textbf{Problem 2.12}: Use the same technique as in Problem 2.11.

\vspace{0.3cm}

\textbf{Problem 2.13}: Use the same technique as in Problem 2.11.

\vspace{0.3cm}

\textbf{Problem 2.14}:
\begin{enumerate}[\hspace{1cm}a)]
  \item Suppose $P:(u,v)$, $P$ is a point of $\arc{JiI}$, and $u>0$ and $v>0$.
  \item Since $P$ is a point of $\arc{JiI}$, $P$ is a point of $\arc{Ji}$ or $P$ is a point of $\arc{iI}$.
  \item If $P$ is a point of $\arc{Ji}$, then $P \neq J$ since $v>0$, and $P \neq i$ since $u>0$. Hence, $\arc{JPi}$ is true. By Theorem 2.2, $P:(u,v)$ and $\arc{JPi}$ mean $u<0$ and $v>0$. But $u>0$. Hence, $\arc{JPi}$ is false.
  \item Then $\arc{iPI}$. Note $P \neq I$ since $v>0$.
\end{enumerate}

\textbf{Problem 2.15}: Similar argument as in Problem 2.14.

\vspace{0.3cm}

\vspace{2cm}

\begin{center}
Science breakthrough: Cold wave linked to temperatures.
\vspace{0.5cm}

Revelation: Typhoon rips through cemetery. Hundreds dead.
\end{center}



\chapter*{Chapter III}
\addcontentsline{toc}{chapter}{\numberline{}Chapter III}
\markboth{Chapter III}{Chapter III}


The relationship between $l(\arc{PI})$ and $PI$ (arc length and chord length) for points $P$ of $\arc{JiI}$ is apparently not something that we are familiar with. So what can we do now?

Let's try something that is not at all obvious. So far, we have been dealing with a collection of number pairs $(l(\arc{PI}), PI)$ for various points $P$ of $\arc{JiI}$. These number pairs are ordered in the sense that the arc length, $l(\arc{PI})$, comes first, and the chord length, $PI$, comes second. Arc length is the first term and chord length is the second term. It's arbitrary which one comes first, but it is helpful to use one of them consistently as the first term. This allows us to graph the number pairs to see if the picture of them -- the graph -- suggests anything useful to us. 

Unfortunately, it didn't seem to help us in this case, although it did suggest some questions to ask about it; e.g., does a line fit the data?

But let's generate another collection of ordered number pairs and see if it will help us. For each point $P$ of $\arc{JiI}$, consider the number pair where:
\begin{enumerate}[\hspace{1cm}a)]
  \item the first term of the number pair is the arc length $l(\arc{PI})$, and
  \item the second term of the number pair is the first coordinate of $P$.
\end{enumerate}
I know what you're thinking. You are thinking either
\begin{enumerate}[\hspace{1cm}1)]
  \item What's the point of doing that?, or
  \item If you're going to pair the arc length $l(\arc{PI})$ with the first coordinate of $P$, what about pairing $l(\arc{PI})$ with the second coordinate of $P$?
\end{enumerate}
The answer to the first question is: experience.

I don't want to discourage you from trying some other approach to the problem at this point. As the old saying goes, ``there are lots of ways to skin a cat'', if you're inclined to do such a thing. In any event, the approach I've made here is what occurred to me because of my background. Your background is different from mine, so you might think of something else to do, and your way may lead to a solution to our problem very nicely. 

I might mention here that sometimes it is very difficult to answer the question ``Why did you do that?''. The only satisfactory answer sometimes is ``Because it leads to the result I'm after.'' But other times it is because you have spent so much time on the problem and become so immersed in it, that you have finally managed to get straight in your mind, all at the same time, enough of the information available that you can ``see'' what you have to do. Have you ever experienced something like that? 

\vfill

You've heard of people claiming that they have ``seen the light''. Well, perhaps what that amounts to is a situation where the person has finally gotten a rather large collection of information together, has understood each piece of the information, and has even been able to see how these various bits of information fit together. 

\vfill

You see, understanding something is more than just being able to recall the information surrounding it. It also means seeing how all those bits of information relate to one another, seeing links among the various categories of the information. 

\vfill

This takes time, sometimes lots and lots of time, to be able to absorb each of the bits of information and then see interconnections among them all. During the time that you are trying to reach this level of understanding, it seems that you are not making much progress. 

\vfill

But when you finally get it all straight in your mind at the same time, lo and behold, the solution is almost obvious! You have seen the light! This idea of ``seeing the light'' is not usually associated with mathematics, but the underlying idea is the same no matter what the subject is.

\vfill

So, what's this got to do with what we were discussing a moment or two ago? Well, how do you explain to someone your solution to a problem when you have gone through what was just described, particularly if it took you months of concentrated study to get to the point where everything suddenly cleared up? It's not easy. 

\vfill

Hence, at this stage, my explanation of why we might pursue our problem by considering the ordered number pairs ($l(\arc{PI})$, first coordinate of $P$) is ``experience''. You will find that as you mature in your chosen field, the approaches you make to problems in your field become less and less transparent to people out of your field. Experience is a great teacher.

\vfill

The answer to the second question is: Let's do that also. This generates two new collections of ordered number pairs to study. Let's list the number pairs that we know that belong to these collections, and graph them just as we did with our first collection $(l(\arc{PI}), PI)$.

\vspace{3.5cm}

Please fill in the blanks.

\begin{center}

\underline{Table 3.1}: First coordinate collection

\begin{table}[htbp]
\begin{center}
\begin{tabular}{|l|l|}
  \hline 
  arc length  &  first coordinate \\
  \hline
  0           &  1 \\
  $\pi/6$     &    \\
  $\pi/4$     &    \\
  $\pi/3$     &    \\
  $\pi/2$     &  0 \\
  $2\pi/3$    &    \\
  $3\pi/4$    &    \\
  $5\pi/6$    &    \\
  $\pi$       & -1 \\
  \hline
\end{tabular}
\label{tbl3.1}
\end{center}
\end{table}

\underline{Table 3.2}: Second coordinate collection

\begin{table}[htbp]
\begin{center}
\begin{tabular}{|l|l|}
  \hline 
  arc length  &  second coordinate \\
  \hline
  0           & 0 \\
  $\pi/6$     &   \\
  $\pi/4$     &   \\
  $\pi/3$     &   \\
  $\pi/2$     & 1 \\
  $2\pi/3$    &   \\
  $3\pi/4$    &   \\
  $5\pi/6$    &   \\
  $\pi$       & 0 \\
  \hline
\end{tabular}
\label{tbl3.2}
\end{center}
\end{table}

\end{center}


Now sketch the graphs of these two collections separately.

\vspace{0.3cm}

This is probably a good place to point out that this device of considering collections of ordered number pairs like we have been discussing is one of the most useful inventions in mathematics. It may not appear to be so important to you right now, but it turns out to be so significant in mathematics that such collections have been given a name. They are called ``functions''.

\vspace{0.3cm}

\textbf{Definition 3.1}: The statement that $f$ is a \emph{function} means:
\begin{enumerate}[\hspace{1cm}1)]
  \item $f$ is a collection, each element of which is an ordered number pair; and
  \item no two number pairs in $f$ have the same first term.
\end{enumerate}
Note that each of the collections that we have been discussing is such that no two ordered number pairs in the collection have the same first term. So, all three are functions.

\vspace{0.3cm}

\textbf{Definition 3.2}: The statement that $g$ is a simple graph means:
\begin{enumerate}[\hspace{1cm}1)]
  \item $g$ is a pointset; and
  \item no vertical line contains two points of $g$.
\end{enumerate}
Find a connection between functions and simple graphs.

\vspace{0.3cm}

Certain conventions pertaining to terminology have grown up around the notion of a function. For example, the collection of first terms of the number pairs in a function $f$ is called the \emph{initial set} of $f$, or sometimes the \emph{domain} of $f$. Furthermore, the collection of second terms of the number pairs in a function $f$ is called the \emph{final set} of $f$, or sometimes the \emph{range} of $f$ or the range of values of $f$.

In our very first collection, the initial set is a collection of arc lengths, and the final set is a set of chord lengths. In fact, the initial set is the interval $[0,\pi]$, and the final set is the interval $[0,2]$.

There is also some notation associated with functions that we need to become accustomed to. Oftentimes it is useful to have notation for the second term of a number pair in a function which not only tells us what function is under discussion, but also what the first term of that number pair is. Hence, if $f$ is a function and $(x,y)$ is a number pair in $f$, by convention, the second term, $y$, is denoted $f(x)$ and read ``$f$ of $x$''. Hence, $(x,y) = (x,f(x))$ or $y = f(x)$. $f(x)$ is the second term of the number pair in the function $f$ whose first term is $x$.

In our problem, let's call the first-coordinate function $C$ and the second-coordinate function $S$.

\vspace{0.3cm}

\textbf{Definition 3.3}: $C$ is the collection to which the number pair $(a,x)$ belongs if and only if:
\begin{enumerate}[\hspace{1cm}1)]
  \item $a$ is the arc length from some point $P$ of $\arc{JiI}$ to $I$; i.e., $a = l(\arc{PI})$; and
  \item $x$ is the first coordinate of $P$.
\end{enumerate}

\textbf{Definition 3.4}: $S$ is the collection to which the number pair $(a,y)$ belongs if and only if:
\begin{enumerate}[\hspace{1cm}1)]
  \item $a$ is the arc length from some point $P$ of $\arc{JiI}$ to $I$; i.e., $a = l(\arc{PI})$; and
  \item $y$ is the second coordinate of $P$.
\end{enumerate}

Draw a picture of the semi-circle $\arc{JiI}$, pick a point $P:(x,y)$ of $\arc{JiI}$, and label everything you can, to see what the definitions mean.

Note: If $P:(x,y)$ is a point of $\arc{JiI}$, and $l(\arc{PI}) = a$, then $(a,x) = (a,C(a))$ implies that $C(a) = x$, and $(a,y) = (a,S(a))$ implies that $S(a) = y$. 

\vspace{0.3cm}

It is going to take considerable effort on your part to get used to this idea of ``function'' and the notation surrounding it. So, let's give you some practice with it. Complete the following:

\vspace{0.3cm}

\textbf{Problem 3.1}: (The first four are intervals.)
\begin{enumerate}[\hspace{1cm}a)]
  \item The initial set of $C$ is ?
  \item The initial set of $S$ is ?
  \item The final set of $C$ is ?
  \item The final set of $S$ is ?
  \item $C(0) =$ ? , and $S(0) =$ ?
  \item $C(\frac{\pi}{2}) =$ ?, and $S(\frac{\pi}{2}) =$ ?
  \item $C(\pi) =$ ? , and $S(\pi) =$ ?
\end{enumerate}

\textbf{Problem 3.2}: Suppose $P$ is a point of $\arc{JiI}$, and $l(\arc{PI}) = a$. What are the coordinates of $P$ in terms of $a$, $C$, and $S$?

\begin{center}
%TeXCAD Options
%\grade{\on}
%\emlines{\off}
%\epic{\off}
%\beziermacro{\on}
%\reduce{\on}
%\snapping{\off}
%\quality{8.00}
%\graddiff{0.01}
%\snapasp{1}
%\zoom{4.0000}
\unitlength 1mm % = 2.85pt
\linethickness{0.4pt}
\ifx\plotpoint\undefined\newsavebox{\plotpoint}\fi % GNUPLOT compatibility
\begin{picture}(115.75,52.75)(0,72)
\put(19,82.5){\line(1,0){82.5}}
\qbezier(59.5,122.75)(91.13,122.75)(105.25,82.75)
\qbezier(59.5,122.75)(22.25,120.38)(15,82.5)
\put(101.5,82.5){\line(1,0){14.25}}
\put(19,82.5){\line(-1,0){11.75}}
%\emline(7.25,82.5)(7,82.25)
\multiput(7.25,82.5)(-.03125,-.03125){8}{\line(0,-1){.03125}}
%\end
\put(58.75,124.75){\makebox(0,0)[cc]{$i:(0,1)$}}
\put(59.25,79){\makebox(0,0)[cc]{$O:(0,0)$}}
\put(104.5,79.75){\makebox(0,0)[cc]{$I:(1,0)$}}
\put(15,79.75){\makebox(0,0)[cc]{$J:(-1,0)$}}
\put(105,108.75){\makebox(0,0)[cc]{$P:(?,?)$}}
\put(96,97){\makebox(0,0)[cc]{$a$}}
%\qbezvec(94.5,99.5)(94.25,101.88)(90,106.75)
\put(90,106.75){\vector(-3,4){.07}}\qbezier(94.5,99.5)(94.25,101.88)(90,106.75)
%\end
%\qbezvec(97.5,95.25)(99.63,90.5)(101.25,83.75)
\put(101.25,83.75){\vector(1,-4){.07}}\qbezier(97.5,95.25)(99.63,90.5)(101.25,83.75)
%\end
\put(90.75,108.75){\circle{.71}} 
\put(58.75,82.5){\circle{.71}}
\put(59.25,73.5){\makebox(0,0)[cc]{Figure 3.1}}
\end{picture}
\end{center}

Complete the following table:

\begin{center}
\underline{Table 3.3}
\end{center}

\begin{table}[htbp]
\begin{center}
\begin{tabular}{|p{30pt}|p{70pt}|p{70pt}|}
\hline 
$\theta$ & $C(\theta)$ & $S(\theta)$ \\
\hline 
$0$ \par $\pi/6$ \par $\pi/4$ \par $\pi/3$ \par $\pi/2$ \par $2\pi/3$ \par $3\pi/4$ \par $5\pi/6$ \par $\pi$ 
& $C(0)=$ \par $C(\pi/6)=$ \par $C(\pi/4)=$ \par $C(\pi/3)=$ \par $C(\pi/2)=$ \par $C(2\pi/3)=$ \par $C(3\pi/4)=$ \par $C(5\pi/6)=$ \par $C(\pi)=$
& $S(0)=$ \par $S(\pi/6)=$ \par $S(\pi/4)=$ \par $S(\pi/3)=$ \par $S(\pi/2)=$ \par $S(2\pi/3)=$ \par $S(3\pi/4)=$ \par $S(5\pi/6)=$ \par $S(\pi)=$ \\
\hline
\end{tabular}
\label{tbl3.3}
\end{center}
\end{table}

When we sketch the graphs of $C$ and $S$, we realize that we need some more number pairs that belong to these functions. There are too many gaps in it. How can we get some more? 

\vfill

Essentially, what is being asked is, how can we find the coordinates of a point of $\arc{JiI}$ if we know the arc length from that point to $I$? Does this seem to relate to our main problem? 

\vfill

Can we find the coordinates of a point of $\arc{JiI}$ which is in some sense halfway between two other points of $\arc{JiI}$ whose coordinates we know?

\vspace{1cm}

\vfill

\textbf{Problem 3.3}: For example, we know $(C(\frac{\pi}{6}), S(\frac{\pi}{6}))$ and $(C(0)), S(0))$. Can we find $(C(\frac{\pi}{12}), S(\frac{\pi}{12}))$?

\vfill

\begin{center}
%TeXCAD Options
%\grade{\on}
%\emlines{\off}
%\epic{\off}
%\beziermacro{\on}
%\reduce{\on}
%\snapping{\off}
%\quality{8.00}
%\graddiff{0.01}
%\snapasp{1}
%\zoom{4.0000}
\unitlength 0.75mm % = 2.85pt
\linethickness{0.4pt}
\ifx\plotpoint\undefined\newsavebox{\plotpoint}\fi % GNUPLOT compatibility
\begin{picture}(160,46)(0,78)
\put(19,82.5){\line(1,0){82.5}}
\qbezier(59.5,122.75)(91.13,122.75)(105.25,82.75)
\qbezier(59.5,122.75)(22.25,120.38)(15,82.5)
\put(101.5,82.5){\line(1,0){14.25}}
\put(19,82.5){\line(-1,0){11.75}}
%\emline(7.25,82.5)(7,82.25)
\multiput(7.25,82.5)(-.03125,-.03125){8}{\line(0,-1){.03125}}
%\end
\put(90.75,108.75){\circle{.71}} 
\put(58.75,82.5){\circle{.71}}
\put(59.25,79){\makebox(0,0)[cc]{$O:(0,0)$}}
\put(120.75,108.75){\makebox(0,0)[cc]{$(C(\frac{\pi}{6}),S(\frac{\pi}{6})) = (\frac{\sqrt{3}}{2},\frac{1}{2})$}}
\put(125.75,94.75){\makebox(0,0)[cc]{$(C(\frac{\pi}{12}),S(\frac{\pi}{12})) = (?,?)$}} 
\put(105.25,79){\makebox(0,0)[cc]{$I:(1,0) = (C(0),S(0))$}}
\put(15.25,120){\makebox(0,0)[cc]{Figure 3.2}}
\end{picture}
\end{center}

\vfill

Show that $C(\frac{\pi}{12}) = (\frac{1}{2}) \sqrt{2+\sqrt{3}}$, and $S(\frac{\pi}{12}) = (\frac{1}{2}) \sqrt{2-\sqrt{3}}$.

\vfill

Note that if $0 \leq \theta \leq \pi$, then $(C(\theta),S(\theta))$ are the coordinates of a point of the unit circle with center at the origin $O:(0,0)$. Since the distance from a point on this unit circle to the origin is $1$, this means that $\sqrt{[C(\theta) - 0]^{2} + [S(\theta) - 0]^{2}} = 1$, or, more conveniently, $[C(\theta)]^{2} + [S(\theta)]^{2} = 1$. 

\vfill

What does $\theta$ represent in all this? Is this true of the results that we have above for $C(\frac{\pi}{12})$ and $S(\frac{\pi}{12})$? That is, is it true that $[C(\frac{\pi}{12})]^{2} + [S(\frac{\pi}{12})]^{2} = 1$ if $C(\frac{\pi}{12}) = (\frac{1}{2}) \sqrt{2+\sqrt{3}}$ and $S(\frac{\pi}{12}) = (\frac{1}{2}) \sqrt{2-\sqrt{3}}$?

\vfill

Use this idea to cross-check your calculations of the other entries in your table of $C$ and $S$. What if your result doesn't check out? What would this mean if you didn't make any errors? Scary question!

\vspace{0.3cm}

\vfill

Construct the following table.

\vspace{4cm}

\begin{center}
\underline{Table 3.4}
\end{center}
\begin{table}[htbp]
\begin{center}
\begin{tabular}{|l|l|l|}
  \hline 
  $\theta$ & $C(\theta)$ & $S(\theta)$ \\
  \hline 
  $0$        & & \\ 
  $\pi/12$   & & \\ 
  $2\pi/12$  & & \\  
  $3\pi/12$  & & \\  
  $4\pi/12$  & & \\  
  $5\pi/12$  & & \\  
  $6\pi/12$  & & \\  
  $7\pi/12$  & & \\  
  $8\pi/12$  & & \\ 
  $9\pi/12$  & & \\ 
  $10\pi/12$ & & \\ 
  $11\pi/12$ & & \\ 
  $12\pi/12$ & & \\ 
  \hline
\end{tabular}
\label{tbl3.4}
\end{center}
\end{table}

\vfill

It gets a little tedious, doesn't it? Perhaps there are some better ways of getting some of these number pairs than grinding out all this arithmetic. Let's look at the picture of the semi-circle $\arc{JiI}$ and see if we can make some guesses about some relationships among various points of $\arc{JiI}$ that may help us.

\vfill

\begin{center}
%TeXCAD Options
%\grade{\on}
%\emlines{\off}
%\epic{\off}
%\beziermacro{\on}
%\reduce{\on}
%\snapping{\off}
%\quality{8.00}
%\graddiff{0.01}
%\snapasp{1}
%\zoom{4.0000}
\unitlength 1mm % = 2.85pt
\linethickness{0.4pt}
\ifx\plotpoint\undefined\newsavebox{\plotpoint}\fi % GNUPLOT compatibility
\begin{picture}(116.75,55.75)(0,70)
\put(19,82.5){\line(1,0){82.5}}
\qbezier(59.5,122.75)(91.13,122.75)(105.25,82.75)
\qbezier(59.5,122.75)(22.25,120.38)(15,82.5)
\put(101.5,82.5){\line(1,0){14.25}}
\put(19,82.5){\line(-1,0){11.75}}
%\emline(7.25,82.5)(7,82.25)
\multiput(7.25,82.5)(-.03125,-.03125){8}{\line(0,-1){.03125}}
%\end
\put(90.75,108.75){\circle{.71}} 
\put(58.75,82.5){\circle{.71}}
\put(4,108.75){\line(1,0){112.75}} 
\put(59,79){\makebox(0,0)[cc]{$O:(0,0)$}}
\put(59,125.75){\makebox(0,0)[cc]{$i:(0,1)$}}
\put(105.25,79){\makebox(0,0)[cc]{$I:(1,0)$}}
\put(14.25,79){\makebox(0,0)[cc]{$J:(-1,0)$}}
\put(6,111){\makebox(0,0)[cc]{$h$}}
\put(102,111.5){\makebox(0,0)[cc]{$P:(C(\theta),S(\theta))$}}
\put(59,73){\makebox(0,0)[cc]{Figure 3.3}}
\end{picture}
\end{center}

\vspace{0.3cm}

\vfill

Suppose $\arc{iPI}$, $P:(C(\theta),S(\theta))$, and $h$ is the horizontal line containing $P$.

\vspace{0.3cm}

\begin{center}
%TeXCAD Options
%\grade{\on}
%\emlines{\off}
%\epic{\off}
%\beziermacro{\on}
%\reduce{\on}
%\snapping{\off}
%\quality{8.00}
%\graddiff{0.01}
%\snapasp{1}
%\zoom{4.0000}
\unitlength 1mm % = 2.85pt
\linethickness{0.4pt}
\ifx\plotpoint\undefined\newsavebox{\plotpoint}\fi % GNUPLOT compatibility
\begin{picture}(116.75,55.75)(0,70)
\put(19,82.5){\line(1,0){82.5}}
\qbezier(59.5,122.75)(91.13,122.75)(105.25,82.75)
\qbezier(59.5,122.75)(22.25,120.38)(15,82.5)
\put(101.5,82.5){\line(1,0){14.25}}
\put(19,82.5){\line(-1,0){11.75}}
%\emline(7.25,82.5)(7,82.25)
\multiput(7.25,82.5)(-.03125,-.03125){8}{\line(0,-1){.03125}}
%\end
\put(90.75,108.75){\circle{.71}} 
\put(58.75,82.5){\circle{.71}}
\put(4,108.75){\line(1,0){112.75}} 
\put(59,82.75){\line(0,1){40}}
\put(59,79){\makebox(0,0)[cc]{$O:(0,0)$}}
\put(109.25,79){\makebox(0,0)[cc]{$I:(1,0)$}}
\put(14.25,79){\makebox(0,0)[cc]{$J:(-1,0)$}}
\put(21.75,111.5){\makebox(0,0)[cc]{$Q:(u,v)$}}
\put(6,111){\makebox(0,0)[cc]{$h$}}
\put(102,111.5){\makebox(0,0)[cc]{$P:(C(\theta),S(\theta))$}}
\put(59,125.75){\makebox(0,0)[cc]{$i:(0,1)$}}
\put(59,73){\makebox(0,0)[cc]{Figure 3.4}}
\end{picture}
\end{center}

Note that $h$ intersects the semi-circle $\arc{JiI}$ at another point $Q:(u,v)$. Make some guesses about the coordinates of $Q$ based on the geometry. See if, by knowing the coordinates of $P$, i.e., $(C(\theta),S(\theta))$, you can determine the coordinates of $Q$. Don't go on until you've given this a good try. Note that you're not asked to prove anything at this point, but if you want to try to prove that your conjecture is true, please do so.

\vspace{0.3cm}

\textbf{Problem 3.4}: What I hope you conjectured (guessed) is that the coordinates of $Q$ are $(C(\pi-\theta), S(\pi-\theta))$. Let's see if we can actually prove this.
\begin{enumerate}[\hspace{1cm}a)]
  \item Note that $Q:(u,v)$ is a point of $K:x^{2} + y^{2} = 1$. Hence $u^{2} + v^{2} = 1$.
  \item $h$ is a horizontal line, so $h: y = b$ for some number $b$. Since both $P:(C(\theta),S(\theta))$ and $Q:(u,v)$ belong to $h$, $v = S(\theta)$.
  \item Show that $u = \pm C(\theta)$, and since $Q \neq P$, $u = -C(\theta)$.
  \item Show that $JQ = PI$.
  \item Use the results of (d) to conclude that $l(\arc{JQ}) = \theta$.
  \item Since $\arc{JQI}$ is true, show that $l(\arc{QI}) = \pi - \theta$.
  \item Hence, $Q: (C(\pi-\theta), S(\pi-\theta))$. But $Q: (-C(\theta), S(\theta))$ so that $C(\pi-\theta) = -C(\theta)$ and $S(\pi-\theta) = S(\theta)$.
\end{enumerate}
If you have successfully answered the questions above, then you know that the coordinates of $Q$ are $(C(\pi-\theta), S(\pi-\theta))$, and that $C(\pi-\theta) = -C(\theta)$ and $S(\pi-\theta) = S(\theta)$.

Can you see how to use these results to fill in Table 3.4 on a previous page with a lot less effort?

\vspace{0.3cm}

Here is another brain teaser for you.

\vspace{0.6cm}

\textbf{Problem 3.5}: Suppose $0 \leq \theta \leq \phi \leq \pi$.
\begin{enumerate}[\hspace{1cm}a)]
  \item Construct a picture of $\arc{JiI}$ and label the coordinates of two points of $\arc{JiI}$ in terms of $C$ and $S$, one with arc length $\theta$ to $I$, and the other with arc length $\phi$ to $I$.
  \item What is the arc length on $\arc{JiI}$ between these two points?
  \item Find a point on $\arc{JiI}$ such that the arc length from that point to $I$ is the arc length you found in (b) above.
  \item What are the coordinates of the point in (c) above in terms of $C$ and $S$?
  \item It you know $C(\theta)$, $S(\theta)$, $C(\phi)$, and $S(\phi)$, can you find the coordinates of the point in (d) above?
\end{enumerate}

Here is what you may have gotten from the preceding set of questions.

\begin{center}
%TeXCAD Options
%\grade{\on}
%\emlines{\off}
%\epic{\off}
%\beziermacro{\on}
%\reduce{\on}
%\snapping{\off}
%\quality{8.00}
%\graddiff{0.01}
%\snapasp{1}
%\zoom{4.0000}
\unitlength 0.8mm % = 2.85pt
\linethickness{0.4pt}
\ifx\plotpoint\undefined\newsavebox{\plotpoint}\fi % GNUPLOT compatibility
\begin{picture}(115.75,54.75)(0,71)
\put(19,82.5){\line(1,0){82.5}}
\qbezier(59.5,122.75)(91.13,122.75)(105.25,82.75)
\qbezier(59.5,122.75)(22.25,120.38)(15,82.5)
\put(101.5,82.5){\line(1,0){14.25}}
\put(19,82.5){\line(-1,0){11.75}}
%\emline(7.25,82.5)(7,82.25)
\multiput(7.25,82.5)(-.03125,-.03125){8}{\line(0,-1){.03125}}
%\end
\put(59,82.5){\circle{.71}}
\put(59,82.75){\line(0,1){40}}
%\qbezvec(98.5,93)(98.25,95.63)(95,99.75)
\put(95,99.75){\vector(-3,4){.07}}\qbezier(98.5,93)(98.25,95.63)(95,99.75)
%\end
%\qbezvec(100.75,87.75)(101.75,86.63)(102.75,84)
\put(102.75,84){\vector(1,-3){.07}}\qbezier(100.75,87.75)(101.75,86.63)(102.75,84)
%\end
\put(30.25,111.75){\circle{.5}} 
\put(19.25,96){\circle{.5}}
\put(96,101.75){\circle{.5}} 
\put(59,79){\makebox(0,0)[cc]{$O:(0,0)$}}
\put(105.25,79){\makebox(0,0)[cc]{$I:(1,0)$}}
\put(14.25,79){\makebox(0,0)[cc]{$J:(-1,0)$}}
\put(59,125.75){\makebox(0,0)[cc]{$i:(0,1)$}}
\put(95.75,90){\makebox(0,0)[cc]{$\phi - \theta$}}
\put(114.75,100.75){\makebox(0,0)[cc]{$(C(\phi - \theta),S(\phi - \theta))$}}
\put(16.5,110.25){\makebox(0,0)[cc]{$(C(\theta),S(\theta))$}}
\put(8,95){\makebox(0,0)[cc]{$(C(\phi),S(\phi))$}}
\put(59,73){\makebox(0,0)[cc]{Figure 3.5}}
\end{picture}
\end{center}

With a rather straightforward application of Axiom 3, the following result can be obtained: $C(\phi-\theta) = C(\phi) C(\theta) + S(\phi) S(\theta)$. With a little imagination and perseverance, and of course Axiom 3, the following result can be obtained:

\vspace{0.3cm}

\textbf{Problem 3.6}: $S(\phi-\theta) = S(\phi) C(\theta) - S(\theta) C(\phi)$

Don't ignore the condition that $0 \leq \theta \leq \phi \leq \pi$.

We now have these relationships regarding $C$ and $S$:
\begin{enumerate}[\hspace{1cm}1)]
  \item If $0 \leq \theta \leq \pi$, then $[C(\theta)]^{2} + [S(\theta)]^{2} = 1$.
  \item If $0 \leq \theta \leq \phi \leq \pi$, then $C(\phi-\theta) = C(\phi) C(\theta) + S(\phi) S(\theta)$, and,
  \item If $0 \leq \theta \leq \phi \leq \pi$, then $S(\phi-\theta) = S(\phi) C(\theta) - S(\theta) C(\phi)$.
\end{enumerate}

There are two other relationships that we can get concerning $C$ and $S$.

\vspace{0.3cm}

\textbf{Problem 3.7}:
\begin{enumerate}
  \item[\hspace{1cm}4)] If $0 \leq \theta \leq \pi$, $0 \leq \phi \leq \pi$, and $0 \leq \theta + \phi \leq \pi$, then $C(\phi + \theta) = C(\phi) C(\theta) - S(\phi) S(\theta)$.
\end{enumerate}

\textbf{Problem 3.8}:
\begin{enumerate}
  \item[\hspace{1cm}5)] If $0 \leq \theta \leq \pi$, $0 \leq \phi \leq \pi$, and $0 \leq \theta+\phi \leq \pi$, then $S(\phi + \theta) = ?$ (try to guess the formula and then derive it by the use of the axioms and previous relationships).
\end{enumerate}

It doesn't appear to be as easy as the first two relationships, does it? Here's a little hint. It is always tempting to try to use something that we already have to get new results. What I'm getting at is this: Can you use the first two relationships to get these last two; i.e., can you use
\begin{enumerate}[\hspace{1cm}1)]
  \item $C(\phi-\theta) = C(\phi) C(\theta) + S(\phi) S(\theta)$ and
  \item $S(\phi-\theta) = S(\phi) C(\theta) - S(\theta) C(\phi)$ to derive
  \item $C(\phi+\theta) = C(\phi) C(\theta) - S(\phi) S(\theta)$ and
  \item $S(\phi+\theta) = ?$
\end{enumerate}

Write $\phi + \theta = \pi - \pi + \phi + \theta$ and see if you can find a way to use what you already know to get 3) and 4) above. The trick introduced above really amounts to just adding zero. The trick is in knowing (or guessing) how to write $0$. It's sometimes a hard trick to learn. Even more obscure sometimes is multiplying by $1$. Again, the trick is in how to write the $1$. Adding $0$ and multiplying by $1$ turn out to be tantamount to magic sometimes. If you learn how to do it, you'll have a very powerful weapon in your arsenal.

Don't go any further until you've given this a real try.

\vspace{0.3cm}

It is convenient to introduce some rather minor notation here. We have written $[C(\theta)]^{2}$ for $C(\theta) C(\theta)$. This is certainly all right, but we will find it written $C^{2}(\theta)$ in the literature.

\vspace{0.3cm}

\textbf{Definition 3.5}: If $f$ is a function and $x$ is a number in the initial set of $f$, then $f^{2}(x) = [f(x)]^{2}$, and is read ``$f$ squared of $x$''.

\vspace{0.3cm}

Hence, our basic formulas, often referred to as ``expansion formulas'', will be written as follows:
\begin{enumerate}[\hspace{1cm}1)]
  \item If $0 \leq \theta \leq \pi$, then $C^{2}(\theta) + S^{2}(\theta) = 1$. 
  \item If $0 \leq \theta \leq \phi \leq \pi$, then $C(\phi-\theta) = C(\phi) C(\theta) + S(\phi) S(\theta)$.
  \item If $0 \leq \theta \leq \phi \leq \pi$, then $S(\phi-\theta) = S(\phi) C(\theta) - S(\theta) C(\phi)$.
  \item If $0 \leq \theta \leq \phi \leq \pi$ and $0 \leq \phi + \theta \leq \pi$, then $C(\phi+\theta) = C(\phi) C(\theta) - S(\phi) S(\theta)$.
  \item If $0 \leq \theta \leq \phi \leq \pi$ and $0 \leq \phi + \theta \leq \pi$, then $S(\phi+\theta) = S(\phi) C(\theta) + S(\theta) C(\phi)$.
\end{enumerate}

By the use of the five basic formulas for $C$ and $S$, show that each of the following is true:

\vspace{0.3cm}

\textbf{Problem 3.10}:
\begin{enumerate}[\hspace{1cm}a)]
  \item if $0 \leq x \leq \pi$, $C(\pi-x) = -C(x)$
  \item if $0 \leq x \leq \pi$, $S(\pi-x) = S(x)$
  \item if $0 \leq x \leq \frac{\pi}{2}$, $C(\frac{\pi}{2}-x) = S(x)$
  \item if $0 \leq x \leq \frac{\pi}{2}$, $S(\frac{\pi}{2}-x) = C(x)$
  \item if $0 \leq x \leq \frac{\pi}{2}$, $C(\frac{\pi}{2}+x) = -S(x)$
  \item if $0 \leq x \leq \frac{\pi}{2}$, $S(\frac{\pi}{2}+x) = C(x)$
  \item if $0 \leq x \leq \frac{\pi}{2}$, $S(2x) = 2 S(x) C(x)$
  \item if $0 \leq x \leq \frac{\pi}{2}$, 
        \begin{enumerate}[\hspace{1cm}i)]
          \item $C(2x) = C^{2}(x) - S^{2}(x)$
          \item $C(2x) = 1 - 2S^{2}(x)$
          \item $C(2x) = 2C^{2}(x) - 1$
        \end{enumerate}
  \item if $0 \leq x \leq \pi$, $C(\frac{x}{2}) = \pm \sqrt{\frac{(1+C(x))}{2}}$
  \item if $0 \leq x \leq \pi$, $S(\frac{x}{2}) = \sqrt{\frac{(1-C(x))}{2}}$

        Why does the formula for $C(\frac{x}{2})$ have a $\pm$, but the one for $S(\frac{x}{2})$ does not?
  \item if $0 \leq x \leq \frac{\pi}{3}$, $S(3x) = 3S(x) -4S^{3}(x)$
  \item if $0 \leq x \leq \pi$, $S(x) = 3S(\frac{x}{3})- 4S^{3}(\frac{x}{3})$
  \item if $0 \leq x \leq \frac{\pi}{3}$, $C(3x) = 4C^{3}(x) - 3C(x)$
  \item if $0 \leq x \leq \pi$, $C(x) = 4C^{3}(\frac{x}{3}) - 3C(\frac{x}{3})$
\end{enumerate}

\textbf{Problem 3.11}: You know that $S(x) \geq 0$ for $0 \leq x \leq \pi$. How do you know this? Suppose you know that $C(\pi) = -l$ and you know no other number pairs in the collection of $C$ and $S$. But suppose you do know the five basic formulas for $C$ and $S$. Do you believe that you have sufficient information (with just what is listed above) to fill in the following table?

\begin{center}
\underline{Table 3.5}
\end{center}
\begin{table}[htbp]
\begin{center}
\begin{tabular}{|l|l|l|}
  \hline 
  $x$        & $C(x)$ & $S(x)$ \\
  \hline 
  $0$        & \qquad & \qquad \\
  $\pi/12$   &        &        \\
  $\pi/6$    &        &        \\
  $\pi/4$    &        &        \\
  $\pi/3$    &        &        \\
  $5\pi/12$  &        &        \\
  $\pi/2$    &        &        \\
  $7\pi/12$  &        &        \\
  $2\pi/3$   &        &        \\
  $3\pi/4$   &        &        \\
  $5\pi/6$   &        &        \\
  $11\pi/12$ &        &        \\
  $\pi$      &  $-1$  &        \\
  \hline
\end{tabular}
\label{tbl3.5}
\end{center}
\end{table}

Now, how do you get started on this kind of problem? Let's see what we can do. What is it that we know in order to start? If $0 \leq x \leq \pi$, $S(x)$ is the second coordinate of a point of the semi-circle $\arc{JiI}$. We know from some previous somewhat agonizing travail that the second coordinate of each point of $\arc{JiI}$ is positive or zero; i.e., non-negative. Hence, if $0 \leq x \leq \pi$, then $S(x) \geq 0$. Now what? Well, we have the five basic formulas and all those properties that you derived from them. It's probably a good idea to put those in front of you rather prominently. Staring at what you know when you have some questions can sometimes cause a neuron to fire. 

Consider the first pair in your table, $(C(0),S(0))$. Note that $S(0) = S(0-0)$. Now use your expansion formula for $S(x-y)$ and you'll see that $S(0) = 0$. With persistence and sometimes a little cleverness, you can get the rest of your entries.

Show that you do believe you can do it by deriving all of these entries without having to use the arc-chord technique directly. Of course, you can use the formulas in Problem 3.10, since they come directly from the five basic formulas. You'll find this much easier to do now as opposed to the last time you constructed this table.

\vspace{0.3cm}

\textbf{Problem 3.12}: Sketch the graphs of $C$ and $S$ separately, based on Table 3.5.

Would the graphs of $C$ and $S$ look better if your table were in increments of $\frac{\pi}{24}$ instead of $\frac{\pi}{12}$?

Can you see that we could fill in our table as densely as we like, if we had the time? Suppose we could set up some computer program that would grind out these entries for us? You could do this yourself if you're a computer wizard. You might wonder for a moment whether it would be worthwhile to print out such a table in increments of, say, $\frac{\pi}{10}^{6}$. This raises another interesting, even practical, question -- assuming that you were charged with creating a table of values for $C$ and $S$. Note that $C(\pi-\theta) = -C(\theta)$ and that $S(\pi-\theta) = S(\theta)$. Does this suggest anything to you about how to lessen your workload? This assumes, of course, that anyone interested in using your table would be aware of your formulas like the ones above. Can you see how to cut your workload in half again? Do you think that you could carry this workload-cutting to ridiculous extremes also?

Let's assume, for the moment, that you could calculate $C(\theta)$ and $S(\theta)$ for any value of $\theta$ in $[0,\pi]$ in which you have an interest. Can we solve our main problem? Remember, we specialized our main problem by considering a special circle, the unit circle with center at the origin. We then specialized even further by limited our discussion to the semi-circle $\arc{JiI}$ of the unit circle with center at the origin. Can we solve our problem of finding a relationship between arc length and chord length if we are confined to the semi-circle $\arc{JiI}$? I believe that you can do it in at least one direction; that is; it you are given the arc length, you can find the associated chord length. Going the other way -- i.e., given the chord length to find the associated arc length -- is a little more interesting. Before you go on, work on both ways. I think you'll find it rewarding to realize that we've made some progress in solving our problem.

\vspace{0.3cm}

Here is another irrelevant puzzle for you to chew on.

Suppose that there are ten people who owe you money. Suppose also that they are all going to pay you on the same day and in the same manner; i.e., each will pay you with a bag of gold coins where each gold coin is to weigh exactly one pound. They owe you a \emph{lot} of money. The problem is that you have heard through your spy network that one of your debtors is going to cheat you by shorting each of his gold coins by exactly one ounce. This is not sufficient for you to tell by looking at or feeling the coins. Your problem is to determine who the cheater is. The only tool you have is a weighing machine that will work exactly, but only once. In other words, you can weigh as many coins as you like and get an exact weighing, but the machine will work only once. Sort of a government issue weighing machine. One other thing that you know is that there are at least ten coins in each bag. How do you catch the cheater by making only one weighing?



\newpage
\begin{center}
\section*{Chapter III Hints and Solutions}
\end{center}


\textbf{Problem 3.1}:
\begin{enumerate}[\hspace{1cm}a)]
  \item $[0, \pi]$
  \item $[0, \pi]$
  \item $[-1, 1]$
  \item $[0,1]$
  \item $C(0) = 1$, $S(0) = 0$
  \item $C(\frac{\pi}{2}) = 0$, $S(\frac{\pi}{2}) = 1$
  \item $C(\pi) = -1$, $S(\pi) = 0$
\end{enumerate}

\textbf{Problem 3.2}: $P$ is a point of $\arc{JiI}$, and $l(\arc{PI}) = x$ implies that $P:(C(x),S(x))$.

\vspace{0.3cm}

\begin{center}
\underline{Table 3.3}
\end{center}
\begin{table}[htbp]
\begin{center}
\begin{tabular}{|p{40pt}|p{40pt}|p{40pt}|}
\hline $\theta$ & $C(\theta)$ & $S(\theta)$ \\
\hline $0$ \par $\pi/6$ \par $\pi/4$ \par $\pi/3$ \par $\pi/2$ \par $2\pi/3$ \par $3\pi/4$ \par $5\pi/6$ \par $\pi$
& $1$ \par $\sqrt{3}/2$ \par $\sqrt{2}/2$ \par $1/2$ \par $0$ \par $-1/2$ \par $-\sqrt{2}/2$ \par $-\sqrt{3}/2$ \par $-1$
& $0$ \par $1/2$ \par $\sqrt{2}/2$ \par $\sqrt{3}/2$ \par $1$ \par $\sqrt{3}/2$ \par $\sqrt{2}/2$ \par $1/2$ \par $0$ \\
\hline
\end{tabular}
\label{tbl3.3hints}
\end{center}
\end{table}

\textbf{Problem 3.3}: Let $P:(C(\frac{\pi}{6}),S(\frac{\pi}{6}))$, and $Q:(C(\frac{\pi}{12}),S(\frac{\pi}{2}))$. Note that $l(\arc{PI}) = \frac{\pi}{6}$, and $l(\arc{QI}) = \frac{\pi}{12}$, and $\arc{PQI}$. Hence, $l(\arc{PQ}) = l(\arc{QI})$ implies that $PQ = QI$. Now use your algebraic prowess to solve for $C(\frac{\pi}{12})$ and $S(\frac{\pi}{12})$.

\vspace{0.3cm}

\textbf{Problem 3.4}:
\begin{enumerate}[\hspace{1cm}a)]
  \item No hints.
  \item No hints.
  \item Since $v = S(\theta)$, $\ $ $u^{2} + v^{2} = 1$ $\ $ means $\ $ $u^{2} + [S(\theta)]^{2} = 1$ $\ $ so that $u^{2} = 1 - [S(\theta)]^{2} = [C(\theta)]^{2}$ , and that means $u = \pm C(\theta)$. But $u = C(\theta)$ means $Q:(C(\theta)$, $S(\theta))$, which means that $Q = P$. Since $Q \neq P$, $Q:(-C(\theta)$, $S(\theta))$.
  \item $\ $
      \begin{align*}
      JQ & = \sqrt{(u+1)^{2}+v^{2}} \\
         & = \sqrt{u^{2}+2u+1+v^{2}} \\
         & = \sqrt{2+2u} \\
         & = \sqrt{2-2C(\theta)}
      \end{align*}

      \begin{align*}
      PI & = \sqrt{[C(\theta)-1]^{2} + [S(\theta)-0]^{2}} \\
         & = \sqrt{[C(\theta)]^{2} -2C(\theta)+1 + [S(\theta)]^{2}} \\
         & = \sqrt{2 - 2C(\theta)}
      \end{align*}

      Hence, $JQ=PI$.
  \item From (d), $l(\arc{JQ}) = l(\arc{PI})$. But $l(\arc{PI}) = \theta$. Hence $l(\arc{JQ}) = \theta$.
  \item $\arc{JQI}$ means $l(\arc{JQ}) + l(\arc{QI}) = l(\arc{JI})$ or $\theta + l(\arc{QI}) = \pi$ which means $l(\arc{QI}) = \pi-\theta$.
  \item Hence, $Q:(C(\pi-\theta),S(\pi-\theta))$. But $Q:(-C(\theta), S(\theta))$. Therefore, $C(\pi-\theta) = -C(\theta)$ and $S(\pi-\theta) = S(\theta)$.
\end{enumerate}

\textbf{Problem 3.5}: Let $P:(C(\phi),S(\phi))$, $Q:(C(\theta),S(\theta))$, and $R:(C(\phi-\theta),S(\phi-\theta))$, and note that $l(\arc{PQ}) = l(\arc{RI})$. Hence $PQ = RI$; i.e., 

$\sqrt{[C(\phi)-C(\theta)]^{2}+[S(\phi)-S(\theta)]^{2}} = \sqrt{[C(\phi-\theta)-1]^{2}+ [S(\phi-\theta)-0]^{2}}$. 

Simplification by squaring both sides and collecting terms yields $C(\phi-\theta) = C(\phi) C(\theta) + S(\phi) S(\theta)$.

\vspace{0.3cm}

\textbf{Problem 3.6}: Note that $l(\arc{PR}) = l(\arc{QI})$ since $\phi - (\phi - \theta) = \theta$. Then $PR = QI$. When you convert $PR = QI$ into an algebraic statement and simplify, you will have to take care of a special case where $S(\phi) = 0$. Don't ignore this. You can handle it.

\vspace{0.3cm}

\textbf{Problem 3.7}: Write $\phi + \theta = \pi - \pi + (\phi+\theta)$

\vspace{0.3cm}

\textbf{Problem 3.8}: Write $\phi + \theta = \pi - [\pi - (\phi+\theta)]$

Write $\phi + \theta = \pi - [(\pi - \phi) - \theta]$.

Then $C(\phi + \theta) = C\{\pi-[(\pi-\phi)-\theta]\}$. Do you remember what $C(\pi-x)$ is?

\vspace{0.3cm}

\textbf{Problem 3.10}: These are straightforward applications of the expansion formulas and your knowledge of $C(0)$, $C(\frac{\pi}{2})$, $C(\pi)$, $S(0)$, $S(\frac{\pi}{2})$, and $S(\pi)$.
\begin{enumerate}
  \item[\hspace{1cm}g)] Use $S(x+y)$, letting $y=x$.
  \item[\hspace{1cm}h)] Use $C(x+y)$, letting $y=x$. Also note that $C^{2}(x)+S^{2}(x) = 1$.
  \item[\hspace{1cm}{i})] Use the results of Problem (h)(iii).
  \item[\hspace{1cm}j)] Use the results of Problem (h)(ii).
  \item[\hspace{1cm}k)] Expand $S(3x) = S(x+2x)$. Why is the condition $0 \leq x \leq \frac{\pi}{3}$ imposed?
  \item[\hspace{1cm}l)] This follows directly from Problem (k) above.
  \item[\hspace{1cm}m)] Expand $C(3x) = C(x+2x)$.
  \item[\hspace{1cm}n)] This follows directly from Problem (m) above.
\end{enumerate}

How much trouble would it be with the tools you have now to construct such a table in increments of $\frac{\pi}{24}$? You don't suppose you'd be asked to do that, do you?

\begin{center}
\underline{Table 3.5}
\end{center}
\begin{table}[htbp]
\begin{center}
\begin{tabular}{|p{40pt}|p{80pt}|p{80pt}|}
\hline $x$ & $C(x)$ & $S(x)$ \\
\hline $0$ \par $\pi/12$ \par $\pi/6$ \par $\pi/4$ \par $\pi/3$ \par $5\pi/12$ \par $\pi/2$
\par $7\pi/12$ \par $2\pi/3$ \par $3\pi/4$ \par $5\pi/6$ \par $11\pi/12$ \par $\pi$
& $1$ \par $(1/2)\sqrt{2+\sqrt{3}}$ \par $\sqrt{3}/2$ \par $\sqrt{2}/2$ \par $1/2$ \par $(1/2)\sqrt{2-\sqrt{3}}$ \par $0$ \par $-(1/2)\sqrt{2-\sqrt{3}}$ \par $-1/2$ \par $-\sqrt{2}/2$ \par $-\sqrt{3}/2$ \par $-(1/2)\sqrt{2+\sqrt{3}}$ \par $-1$
& $0$ \par $(1/2)\sqrt{2-\sqrt{3}}$ \par $1/2$ \par $\sqrt{2}/2$ \par $\sqrt{3}/2$ \par $(1/2)\sqrt{2+\sqrt{3}}$ \par $1$ \par $(1/2)\sqrt{2+\sqrt{3}}$ \par $\sqrt{3}/2$ \par $\sqrt{2}/2$ \par $1/2$ \par $(1/2)\sqrt{2-\sqrt{3}}$ \par $0$ \\
\hline
\end{tabular}
\label{tbl3.5hints}
\end{center}
\end{table}

\textbf{Problem 3.11}: You should already have $S(0) = 0$.
\begin{enumerate}[\hspace{1cm}1)]
  \item $C(0) = C(0 - 0) = C^{2}(0) + S^{2}(0)$. But $C^{2}(x) + S^{2}(x) = l$ for each $x$ in $[0,\pi]$. Hence $C^{2}(0) + S^{2}(0) = 1$. Recall what $S(0)$ is.
  \item Note $C^{2}(\pi) + S^{2}(\pi) = 1$, and we are given $C(\pi) = -l$. This yields $S(\pi) = 0$.
  \item Write $-1= C(\pi) = C(\frac{\pi}{2} + \frac{\pi}{2})$, and expand. This should yield $C(\frac{\pi}{2}) = 0$ and $S(\frac{\pi}{2}) = 1$. Remember $S(x) \geq 0$ for any $x$ in $[0, \pi]$.
  \item Use $S(\frac{x}{2}) = \sqrt{\frac{[1 - C(x)]}{2}}$ to find $S(\frac{\pi}{4})$ and use this result with $S(2x) = 2S(x)C(x)$ to find $C(\frac{\pi}{4})$.
  \item For $C(\frac{3\pi}{4})$ and $S(\frac{3\pi}{4})$, use $C(\pi-x)$ and $S(\pi-x)$ and some previous results.
  \item For $C(\frac{\pi}{6})$ and $S(\frac{\pi}{6})$, use $S(3x)$ and $C(3x)$, and some ingenuity.
  \item For $C(\frac{\pi}{3})$ and $S(\frac{\pi}{3})$, use $C(2x)$ and $S(2x)$ and previous results.
  \item For $C(\frac{2\pi}{3})$ and $S(\frac{2\pi}{3})$, use $C(\pi-x)$ and $S(\pi-x)$. You should be catching on to the technique by now.
  \item For $C(\frac{5\pi}{6})$ and $S(\frac{5\pi}{6})$, use $C(\pi-x)$ and $S(\pi-x)$ and previous results.
  \item $\frac{\pi}{12} = \frac{\pi}{3} - \frac{\pi}{4}$.
  \item $\frac{5\pi}{12} = \frac{\pi}{2} - \frac{\pi}{12}$.
  \item $\frac{7\pi}{12} = \frac{\pi}{2} + \frac{\pi}{12}$.
  \item $\frac{11\pi}{12} = \pi - \frac{\pi}{12}$.
\end{enumerate}


\vspace{2cm}

\begin{center}
Cooking school: Include your children when baking cookies.
\vspace{0.5cm}

International news: Iraqi head seeks arms.
\end{center}



\chapter*{Chapter IV}
\addcontentsline{toc}{chapter}{\numberline{}Chapter IV}
\markboth{Chapter IV}{Chapter IV}


The last problem in the last chapter that was relevant to our course can be stated as follows:

Suppose $P$ is a point of semi-circle $\arc{JiI}$, and $l(\arc{PI} = \theta$. Find $PI$. See if you can show that $PI = \sqrt{2-2C(\theta)}$.

Let's recap what we have done so far. We have two functions, $C$ and $S$, which satisfy the following five conditions:
\begin{enumerate}[\hspace{1cm}1)]
  \item if $0 \leq x \leq \pi$, then $C^{2}(x) + S^{2}(x) = 1$ 
  \item if $0 \leq y \leq x \leq \pi$, then $C(x-y) = C(x)C(y) + S(x)S(y)$
  \item if $0 \leq y \leq x \leq \pi$, then $S(x-y) = S(x)C(y) - S(y)C(x)$
  \item if $0 \leq x \leq \pi$, $0 \leq y \leq \pi$, and $0 \leq x + y \leq \pi$, then $C(x + y) = C(x)C(y) - S(x)S(y)$
  \item if $0 \leq x \leq \pi$, $0 \leq y \leq \pi$, and $0 \leq x + y \leq \pi$, then $S(x+y) = S(x)C(y) + S(y)C(x)$
\end{enumerate}

Note that these five properties are true about $C$ and $S$, whether or not we have our geometric interpretation of them on the upper semi-circle $\arc{JiI}$ of the unit circle with center at the origin.

\begin{center}
%TeXCAD Options
%\grade{\on}
%\emlines{\off}
%\epic{\off}
%\beziermacro{\on}
%\reduce{\on}
%\snapping{\off}
%\quality{8.00}
%\graddiff{0.01}
%\snapasp{1}
%\zoom{4.0000}
\unitlength 0.9mm % = 2.85pt
\linethickness{0.4pt}
\ifx\plotpoint\undefined\newsavebox{\plotpoint}\fi % GNUPLOT compatibility
\begin{picture}(115.75,54.75)(0,71)
\put(19,82.5){\line(1,0){82.5}}
\qbezier(59.5,122.75)(91.13,122.75)(105.25,82.75)
\qbezier(59.5,122.75)(22.25,120.38)(15,82.5)
\put(101.5,82.5){\line(1,0){14.25}}
\put(19,82.5){\line(-1,0){11.75}}
%\emline(7.25,82.5)(7,82.25)
\multiput(7.25,82.5)(-.03125,-.03125){8}{\line(0,-1){.03125}}
%\end
\put(59,82.5){\circle{.71}}
\put(59,82.75){\line(0,1){40}}
%\qbezvec(98.5,93)(98.25,95.63)(95,99.75)
\put(95,99.75){\vector(-3,4){.07}}\qbezier(98.5,93)(98.25,95.63)(95,99.75)
%\end
%\qbezvec(100.75,87.75)(101.75,86.63)(102.75,84)
\put(102.75,84){\vector(1,-3){.07}}\qbezier(100.75,87.75)(101.75,86.63)(102.75,84)
%\end
\put(30.25,111.75){\circle{.5}} 
\put(19.25,96){\circle{.5}}
\put(96,101.75){\circle{.5}} 
\put(59,79){\makebox(0,0)[cc]{$O$}}
\put(105.25,79){\makebox(0,0)[cc]{$I$}}
\put(14.25,79){\makebox(0,0)[cc]{$J$}}
\put(59,125.75){\makebox(0,0)[cc]{$i$}}
\put(95.75,90){\makebox(0,0)[cc]{$x$}}
\put(114.75,100.75){\makebox(0,0)[cc]{$(C(x),S(x))$}}
\put(16.5,110.25){\makebox(0,0)[cc]{$\arc{JiI}$}}
\put(59,73){\makebox(0,0)[cc]{Figure 4.1}}
\end{picture}
\end{center}

Our geometric interpretation of them is that if $0 \leq x \leq \pi$, then the point of $\arc{JiI}$ whose arc length to $I$ is $x$ has coordinates $(C(x), S(x))$. 

How can we generalize this, and what in the world does it mean to generalize it? Mathematicians love this type of thing. They will get some little obscure property to be true in some extremely limited and specialized circumstance and then try to find some way to generalize it to some universal profundity. Sometimes, they come close. Here is one perspective on generalizing what we have done.

$C$ and $S$ are confined to the interval $[0,\pi]$ in the sense that the initial set of each is $[0,\pi]$. It makes no sense to write $C(-l)$ or $C(2\pi)$, because neither $-1$ nor $2\pi$ is in the initial set of $C$. Let us try to define two \emph{new} functions $C_{1}$ and $S_{1}$ with the following properties:
\begin{enumerate}[\hspace{1cm}1)]
  \item The initial set of each is the interval $[0, 2\pi]$.
  \item if $0 \leq x \leq 2\pi$, then $C_{1}^{2}(x) + S_{1}^{2}(x) = 1$.
  \item if $0 \leq y \leq x \leq 2 \pi$, then $C_{1}(x-y) = C_{1}(x) C_{1}(y) + S_{1}(x) S_{1}(y)$
  \item if $0 \leq y \leq x \leq 2 \pi$, then $S_{1}(x-y) = S_{1}(x) C_{1}(y) - S_{1}(y) C_{1}(x)$
  \item if $0 \leq x \leq 2\pi$, $0 \leq y \leq 2\pi$, and $0 \leq x+y \leq 2 \pi$, then $C_{1}(x+y) = C_{1}(x) C_{1}(y) - S_{1}(x) S_{1}(y)$ 
  \item if $0 \leq x \leq 2\pi$, $0 \leq y \leq 2\pi$, and $0 \leq x+y \leq 2 \pi$, then $S_{1}(x+y) = S_{1}(x) C_{1}(y) + S_{1}(y) C_{1}(x)$
\end{enumerate}

Note that the only difference between $(C,S)$ and $(C_{1},S_{1})$ is that $C_{1}$ and $S_{1}$ are defined on $[0,2\pi]$ instead of just $[0,\pi]$.

We would further like to have a geometric interpretation of $C_{1}$ and $S_{1}$ on the whole unit circle, not just the upper half as are $C$ and $S$.

\begin{center}
%TeXCAD Options
%\grade{\on}
%\emlines{\off}
%\epic{\off}
%\beziermacro{\on}
%\reduce{\on}
%\snapping{\off}
%\quality{8.00}
%\graddiff{0.01}
%\snapasp{1}
%\zoom{4.0000}
\unitlength 1mm % = 2.85pt
\linethickness{0.4pt}
\ifx\plotpoint\undefined\newsavebox{\plotpoint}\fi % GNUPLOT compatibility
\begin{picture}(77.5,57.75)(0,60)
%\circle(41,91.5){49.24}
\put(65.62,91.5){\line(0,1){1.065}}
\put(65.6,92.56){\line(0,1){1.063}}
\put(65.53,93.63){\line(0,1){1.059}}
\multiput(65.42,94.69)(-.03213,.21057){5}{\line(0,1){.21057}}
\multiput(65.25,95.74)(-.02943,.14928){7}{\line(0,1){.14928}}
\multiput(65.05,96.78)(-.03138,.12938){8}{\line(0,1){.12938}}
\multiput(64.8,97.82)(-.03284,.11369){9}{\line(0,1){.11369}}
\multiput(64.5,98.84)(-.03087,.09177){11}{\line(0,1){.09177}}
\multiput(64.16,99.85)(-.03191,.08282){12}{\line(0,1){.08282}}
\multiput(63.78,100.85)(-.03273,.075106){13}{\line(0,1){.075106}}
\multiput(63.35,101.82)(-.033379,.068362){14}{\line(0,1){.068362}}
\multiput(62.89,102.78)(-.031766,.058497){16}{\line(0,1){.058497}}
\multiput(62.38,103.72)(-.032251,.053712){17}{\line(0,1){.053712}}
\multiput(61.83,104.63)(-.032624,.049363){18}{\line(0,1){.049363}}
\multiput(61.24,105.52)(-.032901,.045385){19}{\line(0,1){.045385}}
\multiput(60.62,106.38)(-.033091,.041723){20}{\line(0,1){.041723}}
\multiput(59.96,107.21)(-.033204,.038337){21}{\line(0,1){.038337}}
\multiput(59.26,108.02)(-.033248,.035189){22}{\line(0,1){.035189}}
\multiput(58.53,108.79)(-.034739,.033718){22}{\line(-1,0){.034739}}
\multiput(57.76,109.53)(-.037886,.033717){21}{\line(-1,0){.037886}}
\multiput(56.97,110.24)(-.041274,.033649){20}{\line(-1,0){.041274}}
\multiput(56.14,110.92)(-.044938,.033508){19}{\line(-1,0){.044938}}
\multiput(55.29,111.55)(-.04892,.033286){18}{\line(-1,0){.04892}}
\multiput(54.41,112.15)(-.053273,.032971){17}{\line(-1,0){.053273}}
\multiput(53.5,112.71)(-.058065,.032551){16}{\line(-1,0){.058065}}
\multiput(52.57,113.23)(-.063379,.03201){15}{\line(-1,0){.063379}}
\multiput(51.62,113.71)(-.074659,.033737){13}{\line(-1,0){.074659}}
\multiput(50.65,114.15)(-.08239,.03302){12}{\line(-1,0){.08239}}
\multiput(49.66,114.55)(-.09135,.0321){11}{\line(-1,0){.09135}}
\multiput(48.66,114.9)(-.10192,.03093){10}{\line(-1,0){.10192}}
\multiput(47.64,115.21)(-.12895,.03312){8}{\line(-1,0){.12895}}
\multiput(46.61,115.48)(-.14887,.03144){7}{\line(-1,0){.14887}}
\multiput(45.57,115.7)(-.1751,.02913){6}{\line(-1,0){.1751}}
\multiput(44.51,115.87)(-.2643,.0323){4}{\line(-1,0){.2643}}
\put(43.46,116){\line(-1,0){1.062}}
\put(42.4,116.08){\line(-1,0){1.064}}
\put(41.33,116.12){\line(-1,0){1.065}}
\put(40.27,116.11){\line(-1,0){1.064}}
\put(39.2,116.06){\line(-1,0){1.06}}
\multiput(38.14,115.96)(-.21099,-.02929){5}{\line(-1,0){.21099}}
\multiput(37.09,115.81)(-.1746,-.03199){6}{\line(-1,0){.1746}}
\multiput(36.04,115.62)(-.12979,-.02963){8}{\line(-1,0){.12979}}
\multiput(35,115.38)(-.11412,-.03131){9}{\line(-1,0){.11412}}
\multiput(33.97,115.1)(-.1014,-.03259){10}{\line(-1,0){.1014}}
\multiput(32.96,114.77)(-.09081,-.03359){11}{\line(-1,0){.09081}}
\multiput(31.96,114.4)(-.075539,-.031716){13}{\line(-1,0){.075539}}
\multiput(30.98,113.99)(-.068805,-.032456){14}{\line(-1,0){.068805}}
\multiput(30.02,113.54)(-.062847,-.033041){15}{\line(-1,0){.062847}}
\multiput(29.07,113.04)(-.057525,-.033495){16}{\line(-1,0){.057525}}
\multiput(28.15,112.5)(-.049798,-.031957){18}{\line(-1,0){.049798}}
\multiput(27.26,111.93)(-.045823,-.032287){19}{\line(-1,0){.045823}}
\multiput(26.39,111.32)(-.042165,-.032527){20}{\line(-1,0){.042165}}
\multiput(25.54,110.67)(-.03878,-.032685){21}{\line(-1,0){.03878}}
\multiput(24.73,109.98)(-.035633,-.032771){22}{\line(-1,0){.035633}}
\multiput(23.94,109.26)(-.032697,-.032791){23}{\line(0,-1){.032791}}
\multiput(23.19,108.5)(-.032668,-.035728){22}{\line(0,-1){.035728}}
\multiput(22.47,107.72)(-.032573,-.038874){21}{\line(0,-1){.038874}}
\multiput(21.79,106.9)(-.032405,-.042259){20}{\line(0,-1){.042259}}
\multiput(21.14,106.06)(-.032155,-.045916){19}{\line(0,-1){.045916}}
\multiput(20.53,105.18)(-.033684,-.052825){17}{\line(0,-1){.052825}}
\multiput(19.96,104.29)(-.033329,-.057621){16}{\line(0,-1){.057621}}
\multiput(19.42,103.36)(-.03286,-.062943){15}{\line(0,-1){.062943}}
\multiput(18.93,102.42)(-.032257,-.068898){14}{\line(0,-1){.068898}}
\multiput(18.48,101.46)(-.031498,-.075631){13}{\line(0,-1){.075631}}
\multiput(18.07,100.47)(-.03332,-.09091){11}{\line(0,-1){.09091}}
\multiput(17.7,99.47)(-.0323,-.10149){10}{\line(0,-1){.10149}}
\multiput(17.38,98.46)(-.03098,-.11421){9}{\line(0,-1){.11421}}
\multiput(17.1,97.43)(-.03344,-.14843){7}{\line(0,-1){.14843}}
\multiput(16.87,96.39)(-.03149,-.17469){6}{\line(0,-1){.17469}}
\multiput(16.68,95.34)(-.02868,-.21107){5}{\line(0,-1){.21107}}
\put(16.54,94.29){\line(0,-1){1.061}}
\put(16.44,93.23){\line(0,-1){2.129}}
\put(16.38,91.1){\line(0,-1){1.064}}
\put(16.42,90.03){\line(0,-1){1.062}}
\multiput(16.51,88.97)(.0331,-.2642){4}{\line(0,-1){.2642}}
\multiput(16.64,87.91)(.02964,-.17502){6}{\line(0,-1){.17502}}
\multiput(16.82,86.86)(.03187,-.14878){7}{\line(0,-1){.14878}}
\multiput(17.04,85.82)(.03349,-.12885){8}{\line(0,-1){.12885}}
\multiput(17.31,84.79)(.03122,-.10183){10}{\line(0,-1){.10183}}
\multiput(17.62,83.77)(.03236,-.09126){11}{\line(0,-1){.09126}}
\multiput(17.98,82.77)(.03325,-.08229){12}{\line(0,-1){.08229}}
\multiput(18.38,81.78)(.031528,-.069235){14}{\line(0,-1){.069235}}
\multiput(18.82,80.81)(.032193,-.063286){15}{\line(0,-1){.063286}}
\multiput(19.3,79.86)(.032718,-.05797){16}{\line(0,-1){.05797}}
\multiput(19.82,78.94)(.033124,-.053178){17}{\line(0,-1){.053178}}
\multiput(20.39,78.03)(.033427,-.048823){18}{\line(0,-1){.048823}}
\multiput(20.99,77.15)(.033638,-.044841){19}{\line(0,-1){.044841}}
\multiput(21.63,76.3)(.032161,-.039216){21}{\line(0,-1){.039216}}
\multiput(22.3,75.48)(.032289,-.036071){22}{\line(0,-1){.036071}}
\multiput(23.01,74.68)(.032348,-.033135){23}{\line(0,-1){.033135}}
\multiput(23.76,73.92)(.035285,-.033146){22}{\line(1,0){.035285}}
\multiput(24.53,73.19)(.038432,-.033093){21}{\line(1,0){.038432}}
\multiput(25.34,72.5)(.041819,-.03297){20}{\line(1,0){.041819}}
\multiput(26.18,71.84)(.04548,-.032769){19}{\line(1,0){.04548}}
\multiput(27.04,71.22)(.049457,-.032482){18}{\line(1,0){.049457}}
\multiput(27.93,70.63)(.053805,-.032095){17}{\line(1,0){.053805}}
\multiput(28.85,70.09)(.062495,-.033704){15}{\line(1,0){.062495}}
\multiput(29.78,69.58)(.068458,-.033182){14}{\line(1,0){.068458}}
\multiput(30.74,69.12)(.0752,-.032512){13}{\line(1,0){.0752}}
\multiput(31.72,68.69)(.08291,-.03167){12}{\line(1,0){.08291}}
\multiput(32.71,68.31)(.10105,-.03366){10}{\line(1,0){.10105}}
\multiput(33.73,67.98)(.11379,-.03251){9}{\line(1,0){.11379}}
\multiput(34.75,67.68)(.12947,-.031){8}{\line(1,0){.12947}}
\multiput(35.79,67.44)(.14936,-.029){7}{\line(1,0){.14936}}
\multiput(36.83,67.23)(.21067,-.03152){5}{\line(1,0){.21067}}
\put(37.88,67.08){\line(1,0){1.059}}
\put(38.94,66.96){\line(1,0){1.063}}
\put(40.01,66.9){\line(1,0){2.13}}
\put(42.14,66.9){\line(1,0){1.063}}
\put(43.2,66.98){\line(1,0){1.058}}
\multiput(44.26,67.09)(.21048,.03274){5}{\line(1,0){.21048}}
\multiput(45.31,67.26)(.14919,.02986){7}{\line(1,0){.14919}}
\multiput(46.35,67.47)(.12929,.03175){8}{\line(1,0){.12929}}
\multiput(47.39,67.72)(.1136,.03317){9}{\line(1,0){.1136}}
\multiput(48.41,68.02)(.09168,.03113){11}{\line(1,0){.09168}}
\multiput(49.42,68.36)(.08273,.03214){12}{\line(1,0){.08273}}
\multiput(50.41,68.75)(.075011,.032946){13}{\line(1,0){.075011}}
\multiput(51.39,69.18)(.068265,.033577){14}{\line(1,0){.068265}}
\multiput(52.34,69.65)(.058405,.031935){16}{\line(1,0){.058405}}
\multiput(53.28,70.16)(.053618,.032406){17}{\line(1,0){.053618}}
\multiput(54.19,70.71)(.049269,.032767){18}{\line(1,0){.049269}}
\multiput(55.08,71.3)(.04529,.033032){19}{\line(1,0){.04529}}
\multiput(55.94,71.93)(.041628,.033211){20}{\line(1,0){.041628}}
\multiput(56.77,72.59)(.03824,.033315){21}{\line(1,0){.03824}}
\multiput(57.57,73.29)(.035093,.033349){22}{\line(1,0){.035093}}
\multiput(58.34,74.02)(.033618,.034836){22}{\line(0,1){.034836}}
\multiput(59.08,74.79)(.033607,.037984){21}{\line(0,1){.037984}}
\multiput(59.79,75.59)(.03353,.041371){20}{\line(0,1){.041371}}
\multiput(60.46,76.41)(.033378,.045035){19}{\line(0,1){.045035}}
\multiput(61.09,77.27)(.033144,.049016){18}{\line(0,1){.049016}}
\multiput(61.69,78.15)(.032817,.053368){17}{\line(0,1){.053368}}
\multiput(62.25,79.06)(.032383,.058158){16}{\line(0,1){.058158}}
\multiput(62.77,79.99)(.031827,.063471){15}{\line(0,1){.063471}}
\multiput(63.24,80.94)(.033521,.074756){13}{\line(0,1){.074756}}
\multiput(63.68,81.91)(.03278,.08248){12}{\line(0,1){.08248}}
\multiput(64.07,82.9)(.03183,.09144){11}{\line(0,1){.09144}}
\multiput(64.42,83.91)(.03063,.102){10}{\line(0,1){.102}}
\multiput(64.73,84.93)(.03274,.12904){8}{\line(0,1){.12904}}
\multiput(64.99,85.96)(.03101,.14896){7}{\line(0,1){.14896}}
\multiput(65.21,87)(.02863,.17519){6}{\line(0,1){.17519}}
\multiput(65.38,88.06)(.0315,.2644){4}{\line(0,1){.2644}}
\put(65.51,89.11){\line(0,1){1.062}}
\put(65.59,90.18){\line(0,1){1.325}}
%\end
\put(41,91.5){\line(1,0){36.5}} \put(41,91.5){\line(-1,0){34.75}}
%\qbezvec(60.25,99.25)(63.25,97.5)(62.25,91.75)
\put(62.25,91.75){\vector(-1,-4){.07}}\qbezier(60.25,99.25)(63.25,97.5)(62.25,91.75)
%\end
%\qbezvec(58,103.75)(58.5,106.5)(55,109.25)
\put(55,109.25){\vector(-4,3){.07}}\qbezier(58,103.75)(58.5,106.5)(55,109.25)
%\end
%\qbezvec(26.75,107)(55.13,117.75)(58,92.5)
\put(58,92.5){\vector(1,-4){.07}}\qbezier(26.75,107)(55.13,117.75)(58,92.5)
%\end
%\qbezvec(22.25,102.5)(15.88,89.88)(24,81.75)
\put(24,81.75){\vector(1,-1){.07}}\qbezier(22.25,102.5)(15.88,89.88)(24,81.75)
%\end
\put(55.75,111){\circle{.8}} 
\put(20.75,77.5){\circle{.8}}
\put(67.75,111.75){\makebox(0,0)[cc]{$(C(x),S(x))$}}
\put(59,101.5){\makebox(0,0)[cc]{$x$}}
\put(69,88.75){\makebox(0,0)[cc]{$I$}}
\put(24.25,104.25){\makebox(0,0)[cc]{$y$}}
\put(13.25,89){\makebox(0,0)[cc]{$J$}}
\put(10.25,77.25){\makebox(0,0)[cc]{$(C(y),S(y))$}}
\put(41.25,61.75){\makebox(0,0)[cc]{Figure 4.2}}
\end{picture}
\end{center}

In the preceding picture, $x$ is legitimately illustrated as an arc length from a point on the upper semi-circle to the point $I$. However, $Y$ is not, by our definition, an arc length. But we can interpret $Y$ as the length of $\arc{JiI}$, $\pi$, plus the arc length from $J$ to some point on the lower semi-circle, $\arc{JiI}$, of the unit circle. How can we define $C_{1}$ and $S_{1}$ so that they satisfy all six properties listed on the previous page and so that, if $0 \leq x \leq 2\pi$, then $(C_{1}(x), S_{1}(x))$ are the coordinates of a point on the unit circle with center $O$? See if you can come up with something that would work before you go on. Draw some pictures.

Consider the following picture and partial definitions.

\begin{center}
%TeXCAD Options
%\grade{\on}
%\emlines{\off}
%\epic{\off}
%\beziermacro{\on}
%\reduce{\on}
%\snapping{\off}
%\quality{8.00}
%\graddiff{0.01}
%\snapasp{1}
%\zoom{4.0000}
\unitlength 1mm % = 2.85pt
\linethickness{0.4pt}
\ifx\plotpoint\undefined\newsavebox{\plotpoint}\fi % GNUPLOT compatibility
\begin{picture}(89,77)(0,52)
%\circle(49.5,95.5){60.92}
\put(79.96,95.5){\line(0,1){1.219}}
\put(79.94,96.72){\line(0,1){1.217}}
\multiput(79.86,97.94)(-.0305,.3034){4}{\line(0,1){.3034}}
\multiput(79.74,99.15)(-.02839,.20127){6}{\line(0,1){.20127}}
\multiput(79.57,100.36)(-.03122,.1714){7}{\line(0,1){.1714}}
\multiput(79.35,101.56)(-.0333,.14876){8}{\line(0,1){.14876}}
\multiput(79.09,102.75)(-.03138,.11785){10}{\line(0,1){.11785}}
\multiput(78.77,103.93)(-.0328,.10591){11}{\line(0,1){.10591}}
\multiput(78.41,105.09)(-.031316,.088431){13}{\line(0,1){.088431}}
\multiput(78,106.24)(-.032343,.080885){14}{\line(0,1){.080885}}
\multiput(77.55,107.37)(-.033184,.074224){15}{\line(0,1){.074224}}
\multiput(77.05,108.49)(-.031878,.064267){17}{\line(0,1){.064267}}
\multiput(76.51,109.58)(-.032513,.059443){18}{\line(0,1){.059443}}
\multiput(75.93,110.65)(-.033031,.055036){19}{\line(0,1){.055036}}
\multiput(75.3,111.69)(-.033447,.050986){20}{\line(0,1){.050986}}
\multiput(74.63,112.71)(-.032238,.045097){22}{\line(0,1){.045097}}
\multiput(73.92,113.71)(-.032538,.041867){23}{\line(0,1){.041867}}
\multiput(73.17,114.67)(-.032764,.038842){24}{\line(0,1){.038842}}
\multiput(72.39,115.6)(-.03292,.035999){25}{\line(0,1){.035999}}
\multiput(71.56,116.5)(-.033015,.03332){26}{\line(0,1){.03332}}
\multiput(70.7,117.37)(-.035695,.033251){25}{\line(-1,0){.035695}}
\multiput(69.81,118.2)(-.038539,.03312){24}{\line(-1,0){.038539}}
\multiput(68.89,118.99)(-.041565,.032922){23}{\line(-1,0){.041565}}
\multiput(67.93,119.75)(-.044798,.032652){22}{\line(-1,0){.044798}}
\multiput(66.95,120.47)(-.048263,.032301){21}{\line(-1,0){.048263}}
\multiput(65.93,121.15)(-.054729,.033537){19}{\line(-1,0){.054729}}
\multiput(64.89,121.79)(-.059141,.033059){18}{\line(-1,0){.059141}}
\multiput(63.83,122.38)(-.06397,.032469){17}{\line(-1,0){.06397}}
\multiput(62.74,122.93)(-.069295,.03175){16}{\line(-1,0){.069295}}
\multiput(61.63,123.44)(-.080583,.033086){14}{\line(-1,0){.080583}}
\multiput(60.5,123.9)(-.088139,.032129){13}{\line(-1,0){.088139}}
\multiput(59.36,124.32)(-.0968,.03096){12}{\line(-1,0){.0968}}
\multiput(58.2,124.69)(-.11755,.03247){10}{\line(-1,0){.11755}}
\multiput(57.02,125.02)(-.13196,.03082){9}{\line(-1,0){.13196}}
\multiput(55.83,125.29)(-.17111,.0328){7}{\line(-1,0){.17111}}
\multiput(54.63,125.52)(-.201,.03024){6}{\line(-1,0){.201}}
\multiput(53.43,125.71)(-.3031,.0333){4}{\line(-1,0){.3031}}
\put(52.22,125.84){\line(-1,0){1.217}}
\put(51,125.92){\line(-1,0){1.219}}
\put(49.78,125.96){\line(-1,0){1.219}}
\put(48.56,125.95){\line(-1,0){1.218}}
\put(47.34,125.88){\line(-1,0){1.215}}
\multiput(46.13,125.77)(-.24182,-.03184){5}{\line(-1,0){.24182}}
\multiput(44.92,125.61)(-.17168,-.02964){7}{\line(-1,0){.17168}}
\multiput(43.72,125.41)(-.14906,-.03193){8}{\line(-1,0){.14906}}
\multiput(42.53,125.15)(-.13126,-.03366){9}{\line(-1,0){.13126}}
\multiput(41.34,124.85)(-.1062,-.03182){11}{\line(-1,0){.1062}}
\multiput(40.18,124.5)(-.09611,-.03304){12}{\line(-1,0){.09611}}
\multiput(39.02,124.1)(-.081179,-.031596){14}{\line(-1,0){.081179}}
\multiput(37.89,123.66)(-.074526,-.032499){15}{\line(-1,0){.074526}}
\multiput(36.77,123.17)(-.068592,-.03324){16}{\line(-1,0){.068592}}
\multiput(35.67,122.64)(-.059739,-.031964){18}{\line(-1,0){.059739}}
\multiput(34.6,122.06)(-.055338,-.032523){19}{\line(-1,0){.055338}}
\multiput(33.54,121.45)(-.051292,-.032976){20}{\line(-1,0){.051292}}
\multiput(32.52,120.79)(-.047553,-.033336){21}{\line(-1,0){.047553}}
\multiput(31.52,120.09)(-.044081,-.033613){22}{\line(-1,0){.044081}}
\multiput(30.55,119.35)(-.039142,-.032404){24}{\line(-1,0){.039142}}
\multiput(29.61,118.57)(-.036301,-.032588){25}{\line(-1,0){.036301}}
\multiput(28.7,117.75)(-.033623,-.032706){26}{\line(-1,0){.033623}}
\multiput(27.83,116.9)(-.033578,-.035387){25}{\line(0,-1){.035387}}
\multiput(26.99,116.02)(-.033473,-.038232){24}{\line(0,-1){.038232}}
\multiput(26.19,115.1)(-.033304,-.041261){23}{\line(0,-1){.041261}}
\multiput(25.42,114.15)(-.033063,-.044495){22}{\line(0,-1){.044495}}
\multiput(24.69,113.17)(-.032744,-.047963){21}{\line(0,-1){.047963}}
\multiput(24,112.17)(-.032337,-.051697){20}{\line(0,-1){.051697}}
\multiput(23.36,111.13)(-.033602,-.058834){18}{\line(0,-1){.058834}}
\multiput(22.75,110.07)(-.033057,-.063669){17}{\line(0,-1){.063669}}
\multiput(22.19,108.99)(-.032387,-.069){16}{\line(0,-1){.069}}
\multiput(21.67,107.89)(-.031572,-.074924){15}{\line(0,-1){.074924}}
\multiput(21.2,106.76)(-.032939,-.087839){13}{\line(0,-1){.087839}}
\multiput(20.77,105.62)(-.03185,-.09651){12}{\line(0,-1){.09651}}
\multiput(20.39,104.46)(-.03355,-.11725){10}{\line(0,-1){.11725}}
\multiput(20.05,103.29)(-.03203,-.13167){9}{\line(0,-1){.13167}}
\multiput(19.77,102.11)(-.03008,-.14945){8}{\line(0,-1){.14945}}
\multiput(19.52,100.91)(-.03209,-.20071){6}{\line(0,-1){.20071}}
\multiput(19.33,99.71)(-.02884,-.2422){5}{\line(0,-1){.2422}}
\put(19.19,98.5){\line(0,-1){1.216}}
\put(19.09,97.28){\line(0,-1){4.872}}
\multiput(19.2,92.41)(.02962,-.24211){5}{\line(0,-1){.24211}}
\multiput(19.35,91.2)(.03274,-.20061){6}{\line(0,-1){.20061}}
\multiput(19.54,89.99)(.03056,-.14935){8}{\line(0,-1){.14935}}
\multiput(19.79,88.8)(.03245,-.13156){9}{\line(0,-1){.13156}}
\multiput(20.08,87.61)(.03084,-.10649){11}{\line(0,-1){.10649}}
\multiput(20.42,86.44)(.03216,-.09641){12}{\line(0,-1){.09641}}
\multiput(20.8,85.29)(.03322,-.087733){13}{\line(0,-1){.087733}}
\multiput(21.24,84.15)(.031812,-.074822){15}{\line(0,-1){.074822}}
\multiput(21.71,83.02)(.032607,-.068896){16}{\line(0,-1){.068896}}
\multiput(22.23,81.92)(.03326,-.063563){17}{\line(0,-1){.063563}}
\multiput(22.8,80.84)(.032012,-.055635){19}{\line(0,-1){.055635}}
\multiput(23.41,79.78)(.032503,-.051593){20}{\line(0,-1){.051593}}
\multiput(24.06,78.75)(.032897,-.047858){21}{\line(0,-1){.047858}}
\multiput(24.75,77.75)(.033205,-.044389){22}{\line(0,-1){.044389}}
\multiput(25.48,76.77)(.033436,-.041154){23}{\line(0,-1){.041154}}
\multiput(26.25,75.82)(.033596,-.038125){24}{\line(0,-1){.038125}}
\multiput(27.05,74.91)(.033691,-.035279){25}{\line(0,-1){.035279}}
\multiput(27.9,74.03)(.033727,-.032598){26}{\line(1,0){.033727}}
\multiput(28.77,73.18)(.036405,-.032471){25}{\line(1,0){.036405}}
\multiput(29.68,72.37)(.040952,-.033682){23}{\line(1,0){.040952}}
\multiput(30.63,71.59)(.044189,-.033471){22}{\line(1,0){.044189}}
\multiput(31.6,70.86)(.04766,-.033184){21}{\line(1,0){.04766}}
\multiput(32.6,70.16)(.051397,-.032812){20}{\line(1,0){.051397}}
\multiput(33.63,69.5)(.055442,-.032346){19}{\line(1,0){.055442}}
\multiput(34.68,68.89)(.063362,-.033642){17}{\line(1,0){.063362}}
\multiput(35.76,68.32)(.068699,-.033021){16}{\line(1,0){.068699}}
\multiput(36.86,67.79)(.07463,-.03226){15}{\line(1,0){.07463}}
\multiput(37.98,67.3)(.08128,-.031336){14}{\line(1,0){.08128}}
\multiput(39.11,66.87)(.09621,-.03273){12}{\line(1,0){.09621}}
\multiput(40.27,66.47)(.10631,-.03148){11}{\line(1,0){.10631}}
\multiput(41.44,66.13)(.13137,-.03324){9}{\line(1,0){.13137}}
\multiput(42.62,65.83)(.14916,-.03145){8}{\line(1,0){.14916}}
\multiput(43.81,65.58)(.17178,-.02909){7}{\line(1,0){.17178}}
\multiput(45.02,65.37)(.24192,-.03107){5}{\line(1,0){.24192}}
\put(46.23,65.22){\line(1,0){1.215}}
\put(47.44,65.11){\line(1,0){1.218}}
\put(48.66,65.05){\line(1,0){1.22}}
\put(49.88,65.04){\line(1,0){1.219}}
\put(51.1,65.08){\line(1,0){1.216}}
\multiput(52.31,65.17)(.24237,.02739){5}{\line(1,0){.24237}}
\multiput(53.53,65.31)(.2009,.03089){6}{\line(1,0){.2009}}
\multiput(54.73,65.49)(.171,.03335){7}{\line(1,0){.171}}
\multiput(55.93,65.73)(.13186,.03124){9}{\line(1,0){.13186}}
\multiput(57.11,66.01)(.11745,.03284){10}{\line(1,0){.11745}}
\multiput(58.29,66.34)(.0967,.03127){12}{\line(1,0){.0967}}
\multiput(59.45,66.71)(.088036,.032411){13}{\line(1,0){.088036}}
\multiput(60.59,67.13)(.080477,.033344){14}{\line(1,0){.080477}}
\multiput(61.72,67.6)(.069193,.031971){16}{\line(1,0){.069193}}
\multiput(62.83,68.11)(.063866,.032674){17}{\line(1,0){.063866}}
\multiput(63.91,68.67)(.059034,.033248){18}{\line(1,0){.059034}}
\multiput(64.98,69.26)(.054622,.033712){19}{\line(1,0){.054622}}
\multiput(66.01,69.9)(.048159,.032455){21}{\line(1,0){.048159}}
\multiput(67.02,70.59)(.044693,.032795){22}{\line(1,0){.044693}}
\multiput(68.01,71.31)(.04146,.033055){23}{\line(1,0){.04146}}
\multiput(68.96,72.07)(.038432,.033243){24}{\line(1,0){.038432}}
\multiput(69.88,72.87)(.035588,.033365){25}{\line(1,0){.035588}}
\multiput(70.77,73.7)(.032908,.033426){26}{\line(0,1){.033426}}
\multiput(71.63,74.57)(.032805,.036105){25}{\line(0,1){.036105}}
\multiput(72.45,75.47)(.032639,.038947){24}{\line(0,1){.038947}}
\multiput(73.23,76.41)(.032404,.041971){23}{\line(0,1){.041971}}
\multiput(73.98,77.37)(.033622,.047352){21}{\line(0,1){.047352}}
\multiput(74.68,78.37)(.033284,.051093){20}{\line(0,1){.051093}}
\multiput(75.35,79.39)(.032855,.055141){19}{\line(0,1){.055141}}
\multiput(75.97,80.44)(.032322,.059546){18}{\line(0,1){.059546}}
\multiput(76.56,81.51)(.033652,.068392){16}{\line(0,1){.068392}}
\multiput(77.09,82.6)(.032946,.074329){15}{\line(0,1){.074329}}
\multiput(77.59,83.72)(.032083,.080988){14}{\line(0,1){.080988}}
\multiput(78.04,84.85)(.03362,.09591){12}{\line(0,1){.09591}}
\multiput(78.44,86)(.03246,.10601){11}{\line(0,1){.10601}}
\multiput(78.8,87.17)(.03101,.11795){10}{\line(0,1){.11795}}
\multiput(79.11,88.35)(.03282,.14887){8}{\line(0,1){.14887}}
\multiput(79.37,89.54)(.03067,.1715){7}{\line(0,1){.1715}}
\multiput(79.59,90.74)(.0333,.24163){5}{\line(0,1){.24163}}
\put(79.75,91.95){\line(0,1){1.214}}
\put(79.87,93.16){\line(0,1){2.339}}
%\end
\put(49.25,135.75){\line(0,-1){78.25}}
\put(8.5,95.25){\line(1,0){80.5}}
%\emline(27,75.25)(72,115.75)
\multiput(27,75.25)(.037468776,.0337218984){1201}{\line(1,0){.037468776}}
%\end
\put(26.75,75.5){\circle{.5}}
\put(49.5,128.75){\makebox(0,0)[cc]{$i:(0,1)$}}
\put(80,91){\makebox(0,0)[cc]{$I:(1,0)$}}
\put(49.5,90.75){\makebox(0,0)[cc]{$O:(0,0)$}}
\put(19,91){\makebox(0,0)[cc]{$J:(-1,0)$}}
\put(49,60.5){\makebox(0,0)[cc]{$j:(0,-1)$}}
\put(17,74){\makebox(0,0)[cc]{$(C(x),S(x))$}}
\put(49.25,53.25){\makebox(0,0)[cc]{Figure 4.3}}
\end{picture}
\end{center}

\noindent if $0 \leq x \leq \pi$, then $C_{1}(x) = C(x)$ 

\noindent if $\pi \leq x \leq \pi$, then $C_{1}(x) = ?$

\noindent if $0 \leq x \leq \pi$, then $S_{1}(x) = S(x)$ 

\noindent if $\pi \leq x \leq \pi$, then $S_{1}(x) = ?$

Use the picture to get reasonable entries to test in the definitions above before going on.

\vspace{0.3cm}

Let's add a bit to the previous picture.

\begin{center}
%TeXCAD Options
%\grade{\on}
%\emlines{\off}
%\epic{\off}
%\beziermacro{\on}
%\reduce{\on}
%\snapping{\off}
%\quality{8.00}
%\graddiff{0.01}
%\snapasp{1}
%\zoom{4.0000}
\unitlength 1mm % = 2.85pt
\linethickness{0.4pt}
\ifx\plotpoint\undefined\newsavebox{\plotpoint}\fi % GNUPLOT compatibility
\begin{picture}(110,83.75)(0,52)
%\circle(49.5,95.5){60.92}
\put(79.96,95.5){\line(0,1){1.219}}
\put(79.94,96.72){\line(0,1){1.217}}
\multiput(79.86,97.94)(-.0305,.3034){4}{\line(0,1){.3034}}
\multiput(79.74,99.15)(-.02839,.20127){6}{\line(0,1){.20127}}
\multiput(79.57,100.36)(-.03122,.1714){7}{\line(0,1){.1714}}
\multiput(79.35,101.56)(-.0333,.14876){8}{\line(0,1){.14876}}
\multiput(79.09,102.75)(-.03138,.11785){10}{\line(0,1){.11785}}
\multiput(78.77,103.93)(-.0328,.10591){11}{\line(0,1){.10591}}
\multiput(78.41,105.09)(-.031316,.088431){13}{\line(0,1){.088431}}
\multiput(78,106.24)(-.032343,.080885){14}{\line(0,1){.080885}}
\multiput(77.55,107.37)(-.033184,.074223){15}{\line(0,1){.074223}}
\multiput(77.05,108.49)(-.031878,.064267){17}{\line(0,1){.064267}}
\multiput(76.51,109.58)(-.032513,.059443){18}{\line(0,1){.059443}}
\multiput(75.93,110.65)(-.033031,.055036){19}{\line(0,1){.055036}}
\multiput(75.3,111.69)(-.033447,.050986){20}{\line(0,1){.050986}}
\multiput(74.63,112.71)(-.032238,.045097){22}{\line(0,1){.045097}}
\multiput(73.92,113.71)(-.032538,.041867){23}{\line(0,1){.041867}}
\multiput(73.17,114.67)(-.032764,.038842){24}{\line(0,1){.038842}}
\multiput(72.39,115.6)(-.03292,.035999){25}{\line(0,1){.035999}}
\multiput(71.56,116.5)(-.033015,.03332){26}{\line(0,1){.03332}}
\multiput(70.7,117.37)(-.035695,.033251){25}{\line(-1,0){.035695}}
\multiput(69.81,118.2)(-.038539,.03312){24}{\line(-1,0){.038539}}
\multiput(68.89,118.99)(-.041565,.032922){23}{\line(-1,0){.041565}}
\multiput(67.93,119.75)(-.044798,.032652){22}{\line(-1,0){.044798}}
\multiput(66.95,120.47)(-.048263,.032301){21}{\line(-1,0){.048263}}
\multiput(65.93,121.15)(-.054729,.033537){19}{\line(-1,0){.054729}}
\multiput(64.89,121.79)(-.059141,.033059){18}{\line(-1,0){.059141}}
\multiput(63.83,122.38)(-.06397,.032469){17}{\line(-1,0){.06397}}
\multiput(62.74,122.93)(-.069295,.03175){16}{\line(-1,0){.069295}}
\multiput(61.63,123.44)(-.080583,.033086){14}{\line(-1,0){.080583}}
\multiput(60.5,123.9)(-.088139,.032129){13}{\line(-1,0){.088139}}
\multiput(59.36,124.32)(-.0968,.03096){12}{\line(-1,0){.0968}}
\multiput(58.2,124.69)(-.11755,.03247){10}{\line(-1,0){.11755}}
\multiput(57.02,125.02)(-.13196,.03082){9}{\line(-1,0){.13196}}
\multiput(55.83,125.29)(-.17111,.0328){7}{\line(-1,0){.17111}}
\multiput(54.63,125.52)(-.201,.03024){6}{\line(-1,0){.201}}
\multiput(53.43,125.71)(-.3031,.0333){4}{\line(-1,0){.3031}}
\put(52.22,125.84){\line(-1,0){1.217}}
\put(51,125.92){\line(-1,0){1.219}}
\put(49.78,125.96){\line(-1,0){1.219}}
\put(48.56,125.95){\line(-1,0){1.218}}
\put(47.34,125.88){\line(-1,0){1.215}}
\multiput(46.13,125.77)(-.24182,-.03184){5}{\line(-1,0){.24182}}
\multiput(44.92,125.61)(-.17168,-.02964){7}{\line(-1,0){.17168}}
\multiput(43.72,125.41)(-.14906,-.03193){8}{\line(-1,0){.14906}}
\multiput(42.53,125.15)(-.13126,-.03366){9}{\line(-1,0){.13126}}
\multiput(41.34,124.85)(-.1062,-.03182){11}{\line(-1,0){.1062}}
\multiput(40.18,124.5)(-.09611,-.03304){12}{\line(-1,0){.09611}}
\multiput(39.02,124.1)(-.081179,-.031596){14}{\line(-1,0){.081179}}
\multiput(37.89,123.66)(-.074526,-.032499){15}{\line(-1,0){.074526}}
\multiput(36.77,123.17)(-.068592,-.03324){16}{\line(-1,0){.068592}}
\multiput(35.67,122.64)(-.059739,-.031964){18}{\line(-1,0){.059739}}
\multiput(34.6,122.06)(-.055338,-.032523){19}{\line(-1,0){.055338}}
\multiput(33.54,121.45)(-.051292,-.032976){20}{\line(-1,0){.051292}}
\multiput(32.52,120.79)(-.047553,-.033336){21}{\line(-1,0){.047553}}
\multiput(31.52,120.09)(-.044081,-.033613){22}{\line(-1,0){.044081}}
\multiput(30.55,119.35)(-.039142,-.032404){24}{\line(-1,0){.039142}}
\multiput(29.61,118.57)(-.036301,-.032588){25}{\line(-1,0){.036301}}
\multiput(28.7,117.75)(-.033623,-.032706){26}{\line(-1,0){.033623}}
\multiput(27.83,116.9)(-.033578,-.035387){25}{\line(0,-1){.035387}}
\multiput(26.99,116.02)(-.033473,-.038232){24}{\line(0,-1){.038232}}
\multiput(26.19,115.1)(-.033304,-.041261){23}{\line(0,-1){.041261}}
\multiput(25.42,114.15)(-.033063,-.044495){22}{\line(0,-1){.044495}}
\multiput(24.69,113.17)(-.032744,-.047963){21}{\line(0,-1){.047963}}
\multiput(24,112.17)(-.032337,-.051697){20}{\line(0,-1){.051697}}
\multiput(23.36,111.13)(-.033602,-.058834){18}{\line(0,-1){.058834}}
\multiput(22.75,110.07)(-.033057,-.063669){17}{\line(0,-1){.063669}}
\multiput(22.19,108.99)(-.032387,-.069){16}{\line(0,-1){.069}}
\multiput(21.67,107.89)(-.031572,-.074924){15}{\line(0,-1){.074924}}
\multiput(21.2,106.76)(-.032939,-.087839){13}{\line(0,-1){.087839}}
\multiput(20.77,105.62)(-.03185,-.09651){12}{\line(0,-1){.09651}}
\multiput(20.39,104.46)(-.03355,-.11725){10}{\line(0,-1){.11725}}
\multiput(20.05,103.29)(-.03203,-.13167){9}{\line(0,-1){.13167}}
\multiput(19.77,102.11)(-.03008,-.14945){8}{\line(0,-1){.14945}}
\multiput(19.52,100.91)(-.03209,-.20071){6}{\line(0,-1){.20071}}
\multiput(19.33,99.71)(-.02884,-.2422){5}{\line(0,-1){.2422}}
\put(19.19,98.5){\line(0,-1){1.216}}
\put(19.09,97.28){\line(0,-1){4.872}}
\multiput(19.2,92.41)(.02962,-.24211){5}{\line(0,-1){.24211}}
\multiput(19.35,91.2)(.03274,-.20061){6}{\line(0,-1){.20061}}
\multiput(19.54,89.99)(.03056,-.14935){8}{\line(0,-1){.14935}}
\multiput(19.79,88.8)(.03245,-.13156){9}{\line(0,-1){.13156}}
\multiput(20.08,87.61)(.03084,-.10649){11}{\line(0,-1){.10649}}
\multiput(20.42,86.44)(.03216,-.09641){12}{\line(0,-1){.09641}}
\multiput(20.8,85.29)(.03322,-.087733){13}{\line(0,-1){.087733}}
\multiput(21.24,84.15)(.031812,-.074822){15}{\line(0,-1){.074822}}
\multiput(21.71,83.02)(.032607,-.068896){16}{\line(0,-1){.068896}}
\multiput(22.23,81.92)(.03326,-.063562){17}{\line(0,-1){.063562}}
\multiput(22.8,80.84)(.032012,-.055635){19}{\line(0,-1){.055635}}
\multiput(23.41,79.78)(.032503,-.051593){20}{\line(0,-1){.051593}}
\multiput(24.06,78.75)(.032897,-.047858){21}{\line(0,-1){.047858}}
\multiput(24.75,77.75)(.033205,-.044389){22}{\line(0,-1){.044389}}
\multiput(25.48,76.77)(.033436,-.041154){23}{\line(0,-1){.041154}}
\multiput(26.25,75.82)(.033596,-.038125){24}{\line(0,-1){.038125}}
\multiput(27.05,74.91)(.033691,-.035279){25}{\line(0,-1){.035279}}
\multiput(27.9,74.03)(.033727,-.032598){26}{\line(1,0){.033727}}
\multiput(28.77,73.18)(.036405,-.032471){25}{\line(1,0){.036405}}
\multiput(29.68,72.37)(.040952,-.033682){23}{\line(1,0){.040952}}
\multiput(30.63,71.59)(.044189,-.033471){22}{\line(1,0){.044189}}
\multiput(31.6,70.86)(.04766,-.033184){21}{\line(1,0){.04766}}
\multiput(32.6,70.16)(.051397,-.032812){20}{\line(1,0){.051397}}
\multiput(33.63,69.5)(.055442,-.032346){19}{\line(1,0){.055442}}
\multiput(34.68,68.89)(.063362,-.033642){17}{\line(1,0){.063362}}
\multiput(35.76,68.32)(.068699,-.033021){16}{\line(1,0){.068699}}
\multiput(36.86,67.79)(.07463,-.03226){15}{\line(1,0){.07463}}
\multiput(37.98,67.3)(.08128,-.031336){14}{\line(1,0){.08128}}
\multiput(39.11,66.87)(.09621,-.03273){12}{\line(1,0){.09621}}
\multiput(40.27,66.47)(.10631,-.03148){11}{\line(1,0){.10631}}
\multiput(41.44,66.13)(.13137,-.03324){9}{\line(1,0){.13137}}
\multiput(42.62,65.83)(.14916,-.03145){8}{\line(1,0){.14916}}
\multiput(43.81,65.58)(.17178,-.02909){7}{\line(1,0){.17178}}
\multiput(45.02,65.37)(.24192,-.03107){5}{\line(1,0){.24192}}
\put(46.23,65.22){\line(1,0){1.215}}
\put(47.44,65.11){\line(1,0){1.218}}
\put(48.66,65.05){\line(1,0){1.22}}
\put(49.88,65.04){\line(1,0){1.219}}
\put(51.1,65.08){\line(1,0){1.216}}
\multiput(52.31,65.17)(.24237,.02739){5}{\line(1,0){.24237}}
\multiput(53.53,65.31)(.2009,.03089){6}{\line(1,0){.2009}}
\multiput(54.73,65.49)(.171,.03335){7}{\line(1,0){.171}}
\multiput(55.93,65.73)(.13186,.03124){9}{\line(1,0){.13186}}
\multiput(57.11,66.01)(.11745,.03284){10}{\line(1,0){.11745}}
\multiput(58.29,66.34)(.0967,.03127){12}{\line(1,0){.0967}}
\multiput(59.45,66.71)(.088036,.032411){13}{\line(1,0){.088036}}
\multiput(60.59,67.13)(.080477,.033344){14}{\line(1,0){.080477}}
\multiput(61.72,67.6)(.069193,.031972){16}{\line(1,0){.069193}}
\multiput(62.83,68.11)(.063866,.032674){17}{\line(1,0){.063866}}
\multiput(63.91,68.67)(.059034,.033248){18}{\line(1,0){.059034}}
\multiput(64.98,69.26)(.054622,.033712){19}{\line(1,0){.054622}}
\multiput(66.01,69.9)(.048159,.032455){21}{\line(1,0){.048159}}
\multiput(67.02,70.59)(.044693,.032795){22}{\line(1,0){.044693}}
\multiput(68.01,71.31)(.04146,.033055){23}{\line(1,0){.04146}}
\multiput(68.96,72.07)(.038432,.033243){24}{\line(1,0){.038432}}
\multiput(69.88,72.87)(.035588,.033365){25}{\line(1,0){.035588}}
\multiput(70.77,73.7)(.032908,.033426){26}{\line(0,1){.033426}}
\multiput(71.63,74.57)(.032805,.036105){25}{\line(0,1){.036105}}
\multiput(72.45,75.47)(.032639,.038947){24}{\line(0,1){.038947}}
\multiput(73.23,76.41)(.032404,.041971){23}{\line(0,1){.041971}}
\multiput(73.98,77.37)(.033621,.047352){21}{\line(0,1){.047352}}
\multiput(74.68,78.37)(.033284,.051093){20}{\line(0,1){.051093}}
\multiput(75.35,79.39)(.032855,.055141){19}{\line(0,1){.055141}}
\multiput(75.97,80.44)(.032322,.059546){18}{\line(0,1){.059546}}
\multiput(76.56,81.51)(.033652,.068392){16}{\line(0,1){.068392}}
\multiput(77.09,82.6)(.032946,.074329){15}{\line(0,1){.074329}}
\multiput(77.59,83.72)(.032083,.080988){14}{\line(0,1){.080988}}
\multiput(78.04,84.85)(.03362,.09591){12}{\line(0,1){.09591}}
\multiput(78.44,86)(.03246,.10601){11}{\line(0,1){.10601}}
\multiput(78.8,87.17)(.03101,.11795){10}{\line(0,1){.11795}}
\multiput(79.11,88.35)(.03282,.14887){8}{\line(0,1){.14887}}
\multiput(79.37,89.54)(.03067,.1715){7}{\line(0,1){.1715}}
\multiput(79.59,90.74)(.0333,.24163){5}{\line(0,1){.24163}}
\put(79.75,91.95){\line(0,1){1.214}}
\put(79.87,93.16){\line(0,1){2.339}}
%\end
\put(49.25,135.75){\line(0,-1){78.25}}
\put(8.5,95.25){\line(1,0){80.5}}
%\emline(21.25,84.5)(77.75,106.5)
\multiput(21.25,84.5)(.086523737,.033690658){653}{\line(1,0){.086523737}}
%\end
\put(95,107.75){\makebox(0,0)[cc]{$(C(x - \pi),S(x - \pi))$}}
\put(10,83.5){\makebox(0,0)[cc]{$(C(x),S(x))$}}
\put(19,61.75){\makebox(0,0)[cc]{Figure 4.4}}
\end{picture}
\end{center}

Now try again, if you need to, before going on. We adopt the following definition.

\vspace{0.3cm}

\textbf{Definition 4.1}:

\noindent if $0 \leq x \leq \pi$, then $C_{1}(x) = C(x)$ 

\noindent if $\pi \leq x \leq \pi$, then $C_{1}(x) = -C(x-\pi)$

\noindent if $0 \leq x \leq \pi$, then $S_{1}(x) = S(x)$ 

\noindent if $\pi \leq x \leq \pi$, then $S_{1}(x) = -S(x-\pi)$

Do you see how the picture suggests these definitions?

\vspace{0.3cm}

\textbf{Problem 4.1}: Now, can you show that on the basis of Definition 4.1, $C_{1}$ and $S_{1}$ satisfy the five algebraic properties listed previously; i.e., $C_{1}^{2}(x) + S_{1}^{2}(x) = 1$, etc?

\vspace{0.3cm}

\textbf{Problem 4.2}: After you have done the above, construct a table for $C_{1}$ and $S_{1}$ in intervals of length $\frac{\pi}{24}$ from $0$ to $2\pi$.

\vspace{0.3cm}

\textbf{Problem 4.3}: Then use the table to sketch the graphs of $C_{1}$ and $S_{1}$ on the interval $[0, 2\pi]$. How much trouble would it be to construct your table for $C_{1}$ and $S_{1}$ in intervals of length $\frac{\pi}{48}$ from $0$ to $2\pi$?

\vspace{0.3cm}

\textbf{Problem 4.4}:
\begin{enumerate}[\hspace{1cm}1)]
  \item Suppose $P:(C_{1}(u),S_{1}(u))$ is a point of $K:x^{2} + y^{2}= 1$. What is $l(\arc{PI})$?
  \item Suppose $P:(C_{1}(u),S_{1}(u))$ and $Q:(C_{1}(v),S_{1}(v))$ are points of $K:x^{2}+ y^{2} = 1$. What is $l(\arc{PQ})$? Draw pictures to aid your thinking.
\end{enumerate}

Here is a little algebra teaser for you algebra wizards. I am going to prove to you that an elephant and a flea weigh the same. You don't believe that I can do it? Then your mission, if you should choose to accept it, is to find a mistake in the following argument. 
\begin{enumerate}[\hspace{1cm}1)]
  \item Suppose $E$ = weight of an elephant, and $F$ = weight of a flea.
  \item Suppose $S = E + F$.
  \item Then $S(E - F) = (E + F) (E - F)$.
  \item Hence, $SE - SF = E^{2} - F^{2}$.
  \item Therefore, $F^{2} - SF = E^{2} - SE$.
  \item Then $F^{2} - SF + (\frac{S}{2})^{2} = E^{2} - SE + (\frac{S}{2})^{2}$, which means $(F-\frac{S}{2})^{2} = (E-\frac{S}{2})^{2}$.
  \item Hence $F - \frac{S}{2} = E - \frac{S}{2}$, so that $F = E$.
\end{enumerate}

If there isn't anything wrong with this argument, then we might just as well abandon any study of mathematics because it leads to absurd results. History, anyone?


\newpage
\begin{center}
\section*{Chapter IV Hints and Solutions}
\end{center}


\textbf{Problem 4.1}: 

\emph{Part 1.} Show that if $0 \leq x \leq 2\pi$, then $C_{1}^{2}(x)+S_{1}^{2}(x)=1$.
\begin{enumerate}[\hspace{1cm}1)]
  \item Suppose $0 \leq x \leq 2\pi$. Either $0 \leq x \leq \pi$, or $\pi \leq x \leq 2\pi$.
  \item If $0 \leq x \leq \pi$, show that $C_{1}^{2}(x) + S_{1}^{2}(x) = 1$.
  \item If $\pi \leq x \leq 2\pi$, show that $C_{1}^{2}(x) + S_{1}^{2}(x) = 1$.

        You must refer to the definitions of $C_{1}$ and $S_{1}$. Of course you can use all of the properties that you know about $C$ and $S$, but you must be sure that when you write $C(x)$, then $0 \leq x \leq \pi$. $C(x)$ is not meaningful if $\pi < x \leq 2 \pi$.
\end{enumerate}

\emph{Part 2.} Show that if $0 \leq y \leq x \leq 2\pi$, then $C_{1}(x-y) = C_{1}(x) C_{1}(y) + S_{1}(x) S_{1}(y)$.
\begin{enumerate}[\hspace{1cm}1)]
  \item There are three cases.
    \begin{enumerate}[\hspace{1cm}a)]
      \item $0 \leq y \leq x \leq \pi$
      \item $0 \leq y \leq \pi < x \leq 2 \pi$
      \item $\pi < y \leq x \leq 2 \pi$
    \end{enumerate}
  \item Consider each case separately from the others, and by using the definitions of $C_{1}$ and $S_{1}$, show that $C_{1}(x-y) = C_{1}(x) C_{1}(y) + S_{1}(x) S_{1}(y)$ in each case.
\end{enumerate}

\emph{Part 3.} Show that if $0 \leq y \leq x \leq 2\pi$, then $S_{1}(x-y) = S_{1}(x) C_{1}(y) - S_{1}(y) C_{1}(x)$.
\begin{enumerate}[\hspace{1cm}1)]
  \item Use the same procedure as in 2) above.
\end{enumerate}

\emph{Part 4.} Show that if $0 \leq y \leq x \leq 2\pi$ and $0 \leq x+y \leq 2\pi$, then $C_{1}(x+y) = C_{1}(x) C_{1}(y) - S_{1}(x) S_{1}(y)$ and $S_{1}(x+y) = S_{1}(x) C_{1}(y) +
S_{1}(y) C_{1}(x)$.
\begin{enumerate}[\hspace{1cm}1)]
  \item List all of the cases that apply here, and then take care of each case individually.
  \item If $0 \leq y \leq x \leq \pi$ and $\pi < x+y \leq 2\pi$, write $C_{1}(x+y) = -C[(x+y) - \pi]$ by definition.
  \item Now apply a little algebraic witchcraft. $(x+y) - \pi = x + (y-\pi) = x - (\pi-y)$. You might muse over this to see how it does what you want. Once you see how to use this little device above, you should be able to handle all of the cases.
\end{enumerate}

\textbf{Problem 4.2}:
\begin{enumerate}[\hspace{1cm}1)]
  \item Now that you have the five basic properties for $C_{1}$ and $S_{1}$ that were true for $C$ and $S$, all of the properties about $C$ and $S$ that you derived from these five basic properties are derivable for $C_{1}$ and $S_{1}$. Do this now.
  \item Use the properties derived in Step 1 above to generate the table for $C_{1}$ and $S_{1}$ in increments of $\frac{\pi}{24}$.
\end{enumerate}

\textbf{Problem 4.3}: The graphs should look something like the following:

\begin{center}
%TeXCAD Picture [prob4.3hint1.pic]. Options:
%\grade{\on}
%\emlines{\off}
%\epic{\off}
%\beziermacro{\on}
%\reduce{\on}
%\snapping{\off}
%\quality{8.00}
%\graddiff{0.01}
%\snapasp{1}
%\zoom{4.0000}
\unitlength 1mm % = 2.85pt
\linethickness{0.4pt}
\ifx\plotpoint\undefined\newsavebox{\plotpoint}\fi % GNUPLOT compatibility
\begin{picture}(114,61.25)(0,0)
\put(11.5,61){\line(0,-1){59.5}}
\put(11.75,32.75){\line(1,0){102.25}}
\put(112.75,33.25){\line(0,-1){1}}
\put(80,33.5){\line(0,-1){1.25}}
\put(66.25,33.5){\line(0,-1){1.25}}
\put(53.75,33.5){\line(0,-1){1.5}}
\put(25.5,33.5){\line(0,-1){1}}
\put(25.5,32.5){\line(0,-1){.25}}
\put(39.5,32.75){\circle*{.71}}
\put(12.25,47.25){\line(-1,0){1.5}}
\put(12.25,15.75){\line(-1,0){1.25}}
\put(11.75,60.75){\circle*{.71}}
\put(67.25,2){\circle*{.5}}
\put(112.5,60){\circle*{.5}}
\put(7.75,60.25){\makebox(0,0)[cc]{$1$}}
\put(7.75,32.75){\makebox(0,0)[cc]{$0$}}
\put(8,1.75){\makebox(0,0)[cc]{$-1$}}
\put(38.75,50.25){\makebox(0,0)[cc]{$C$}}
\put(36.25,29){\makebox(0,0)[cc]{$\pi/2$}}
\put(66.25,28.75){\makebox(0,0)[cc]{$\pi$}}
\put(94.5,29){\makebox(0,0)[cc]{$3\pi/2$}}
\put(112.5,29.75){\makebox(0,0)[cc]{$2\pi$}}
\put(10.75,1.75){\line(1,0){1.5}}
%\qbezvec(37.5,47.5)(37.5,44.88)(34.5,42.75)
\put(34.5,42.75){\vector(-4,-3){.07}}\qbezier(37.5,47.5)(37.5,44.88)(34.5,42.75)
%\end
\qbezier(39.5,33)(68.75,-28.75)(91,32.5)
\qbezier(39.5,33)(28.88,54.88)(11.75,61.25)
\qbezier(91,32.75)(97.88,52.5)(112.25,60.25)
\end{picture}
\end{center}

\begin{center}
%TeXCAD Picture [prob4.3hint2.pic]. Options:
%\grade{\on}
%\emlines{\off}
%\epic{\off}
%\beziermacro{\on}
%\reduce{\on}
%\snapping{\off}
%\quality{8.00}
%\graddiff{0.01}
%\snapasp{1}
%\zoom{4.0000}
\unitlength 1mm % = 2.85pt
\linethickness{0.4pt}
\ifx\plotpoint\undefined\newsavebox{\plotpoint}\fi % GNUPLOT compatibility
\begin{picture}(114,61.25)(0,0)
\put(12.25,62.5){\line(0,-1){59.5}}
\put(12.5,34.25){\line(1,0){102.25}}
\put(80.75,35){\line(0,-1){1.25}}
\put(67,35){\line(0,-1){1.25}}
\put(54.5,35){\line(0,-1){1.5}}
\put(26.25,35){\line(0,-1){1}}
\put(26.25,34){\line(0,-1){.25}}
\put(40.25,34.25){\circle{.71}}
\put(13,48.75){\line(-1,0){1.5}}
\put(13,17.25){\line(-1,0){1.25}}
\put(12.5,62.25){\circle{.71}}
\put(11.5,3.25){\line(1,0){1.5}}
\put(38.5,63.5){\circle{.71}}
\put(91.75,33.5){\line(0,1){1.5}}
\put(103.25,33.75){\line(0,1){1.75}}
\put(62.75,55.5){\makebox(0,0)[cc]{$S$}}
\put(8.5,61.75){\makebox(0,0)[cc]{$1$}}
\put(8.5,34.25){\makebox(0,0)[cc]{$0$}}
\put(8.75,3.25){\makebox(0,0)[cc]{$-1$}}
\put(39,30.25){\makebox(0,0)[cc]{$\pi/2$}}
\put(64.5,30.5){\makebox(0,0)[cc]{$\pi$}}
\put(91.75,30.5){\makebox(0,0)[cc]{$3\pi/2$}}
\put(116.75,31.25){\makebox(0,0)[cc]{$2\pi$}}
%\qbezvec(61.75,54.25)(60.5,52.63)(56.25,51.5)
\put(56.25,51.5){\vector(-4,-1){.07}}\qbezier(61.75,54.25)(60.5,52.63)(56.25,51.5)
%\end
\qbezier(12.5,34.25)(36.75,93.25)(67,34.25)
\qbezier(67,34.25)(162,34.25)(67,34.25)
\qbezier(67,34.25)(92.75,-28.75)(114.5,34.25)
\end{picture}
\end{center}
 
\textbf{Problem 4.4}:

\emph{Part 1.} 
\begin{enumerate}[\hspace{1cm}1)]
  \item Note that there are two cases:
    \begin{enumerate}[\hspace{1cm}a)]
      \item $P$ is a point of semi-circle $\arc{JiI}$.
      \item $P$ is a point of semi-circle $\arc{JjI}$.
    \end{enumerate}
  \item For case (a), $C_{1}(u) = C(u)$ and $S_{1}(u) = S(u)$.
  \item Still in case (a), use the definitions of $C$ and $S$ to determine $l(\arc{PI})$.
  \item Still in case (a), you should have determined that $l(\arc{PI}) = u$.
  \item Case (b): $C_{1}(u) = -C(u-\pi)$ and $S_{1}(u) = -S(u-\pi)$.
  \item Let $Q:(C(u-\pi), S(u-\pi))$ and note that $Q$ is a point of semi-circle $\arc{JiI}$.
  \item What relationship exists between $PI$ and $JQ$?
  \item $PI=JQ$ implies that $l(\arc{PI}) = l(\arc{JQ})$.
  \item $Q$ is a point of $\arc{JiI}$, which means $l(\arc{JQ}) + l(\arc{QI}) = l(\arc{JiI}) = \pi$.
  \item $l(\arc{JQ}) = 2\pi - u$ since $l(\arc{QI}) = u - \pi$.
  \item Therefore, $l(\arc{PI}) = 2\pi - u$.
\end{enumerate}

\emph{Part 2.}
\begin{enumerate}[\hspace{1cm}1)]
  \item Suppose $u<v$. Then $0 \leq u < v \leq 2\pi$.
  \item There are two cases.
    \begin{enumerate}[\hspace{1cm}a)]
      \item $0 \leq v-u \leq \pi$,
      \item $\pi < v-u \leq 2\pi$.
    \end{enumerate}
  \item In case (a), let $R:(C(v-u),S(u-v))$ and consider $PQ$ and $RI$.
  \item $PQ = RI$ implies that $l(\arc{PQ}) = l(\arc{RI})$.
  \item Note that $R$ is a point of semi-circle $\arc{JiI}$ so that $l(\arc{RI}) = v - u$.
  \item In case (b), let $R:(C_{1}(v-u),S_{1}(v-u))$. Note that since $\pi \leq v-u \leq 2\pi$, $R$ is in the bottom semi-circle, $\arc{JjI}$. Hence, by case (b) of part 1 of the problem, $l(\arc{RI}) = 2\pi - (v-u)$.
  \item Show $RI = PQ$.
  \item $l(\arc{PQ}) = 2\pi - (v-u)$.
\end{enumerate}

\vspace{2cm}

\begin{center}
Government Probes:
\vspace{0.5cm}

1) Plane too close to ground, crash probe reveals.
\vspace{0.5cm}

2) Two Soviet ships collide. One dies.
\end{center}



\chapter*{Chapter V}
\addcontentsline{toc}{chapter}{\numberline{}Chapter V}
\markboth{Chapter V}{Chapter V}


We have accomplished something that is very interesting from a mathematical point of view. We started with the functions $C$ and $S$, defined on the interval $[0,\pi]$, which satisfied the five basic properties. We then extended this idea by the invention of the functions $C_{1}$ and $S_{1}$, defined on the interval $[0,2\pi]$, which also had the same five basic properties. 

How much further can we carry the idea? Is it possible that we could invent two functions defined over the whole x-axis that had these five basic properties? This is one of the things that mathematicians like to do -- start with some concept that is true in a limited way, and then find some way to generalize the notion as far as they can. Let's see if we can do this in our case.

Here is the problem: Can we find two functions $f$ and $g$, each defined on the whole $x$-axis, such that the following are all true for each number $x$ and each number $y$?

\begin{enumerate}[\hspace{1cm}1)]
  \item $f^{2}(x)+g^{2}(x)=1$
  \item $f(x-y)=f(x)f(y)+g(x)g(y)$
  \item $g(x-y)=g(x)f(y)-g(y)f(x)$
  \item $f(x+y)=f(x)f(y)-g(x)g(y)$
  \item $g(x+y)=g(x)f(y)+g(y)f(x)$
\end{enumerate}

There are probably a lot of ways of going about this, but consider the following approach. Let's use $C_{1}$ and $S_{1}$ as a basis for generating our new functions. Sketches of the graphs of these are as follows:

\begin{center}
%TeXCAD Picture [fig5.1.pic]. Options:
%\grade{\on}
%\emlines{\off}
%\epic{\off}
%\beziermacro{\on}
%\reduce{\on}
%\snapping{\off}
%\quality{8.00}
%\graddiff{0.01}
%\snapasp{1}
%\zoom{4.0000}
\unitlength 0.85mm % = 2.85pt
\linethickness{0.4pt}
\ifx\plotpoint\undefined\newsavebox{\plotpoint}\fi % GNUPLOT compatibility
\begin{picture}(140,72.5)(9,71)
\put(19.75,124.25){\line(0,-1){43}}
\put(19.75,103.75){\line(1,0){50.25}}
\put(87.5,124){\line(0,-1){41.25}}
\put(87.75,104.5){\line(1,0){45.25}}
\put(45.75,103.25){\line(0,1){1}}
\put(69.75,103.25){\line(0,1){1}}
\put(58.5,103.25){\line(0,1){1}}
\put(32,103){\line(0,1){1.25}}
%\emline(110.5,105)(110.25,104.25)
\multiput(110.5,105)(-.03125,-.09375){8}{\line(0,-1){.09375}}
%\end
\put(132.75,104){\line(0,1){1}}
\put(122,105){\line(0,-1){1}}
\put(98.75,105.25){\line(0,-1){1.25}}
%\dashline{1}(19.75,124)(70,124)
\put(19.68,123.93){\line(1,0){.985}}
\put(21.65,123.93){\line(1,0){.985}}
\put(23.62,123.93){\line(1,0){.985}}
\put(25.59,123.93){\line(1,0){.985}}
\put(27.56,123.93){\line(1,0){.985}}
\put(29.53,123.93){\line(1,0){.985}}
\put(31.5,123.93){\line(1,0){.985}}
\put(33.47,123.93){\line(1,0){.985}}
\put(35.44,123.93){\line(1,0){.985}}
\put(37.42,123.93){\line(1,0){.985}}
\put(39.39,123.93){\line(1,0){.985}}
\put(41.36,123.93){\line(1,0){.985}}
\put(43.33,123.93){\line(1,0){.985}}
\put(45.3,123.93){\line(1,0){.985}}
\put(47.27,123.93){\line(1,0){.985}}
\put(49.24,123.93){\line(1,0){.985}}
\put(51.21,123.93){\line(1,0){.985}}
\put(53.18,123.93){\line(1,0){.985}}
\put(55.15,123.93){\line(1,0){.985}}
\put(57.12,123.93){\line(1,0){.985}}
\put(59.09,123.93){\line(1,0){.985}}
\put(61.06,123.93){\line(1,0){.985}}
\put(63.03,123.93){\line(1,0){.985}}
\put(65,123.93){\line(1,0){.985}}
\put(66.97,123.93){\line(1,0){.985}}
\put(68.94,123.93){\line(1,0){.985}}
%\end
%\dashline{1}(19.75,81.5)(70,81.25)
\put(19.68,81.43){\line(1,0){.985}}
\put(21.65,81.42){\line(1,0){.985}}
\put(23.62,81.41){\line(1,0){.985}}
\put(25.59,81.4){\line(1,0){.985}}
\put(27.56,81.39){\line(1,0){.985}}
\put(29.53,81.38){\line(1,0){.985}}
\put(31.5,81.37){\line(1,0){.985}}
\put(33.47,81.36){\line(1,0){.985}}
\put(35.44,81.35){\line(1,0){.985}}
\put(37.42,81.34){\line(1,0){.985}}
\put(39.39,81.33){\line(1,0){.985}}
\put(41.36,81.32){\line(1,0){.985}}
\put(43.33,81.31){\line(1,0){.985}}
\put(45.3,81.3){\line(1,0){.985}}
\put(47.27,81.29){\line(1,0){.985}}
\put(49.24,81.28){\line(1,0){.985}}
\put(51.21,81.27){\line(1,0){.985}}
\put(53.18,81.26){\line(1,0){.985}}
\put(55.15,81.25){\line(1,0){.985}}
\put(57.12,81.24){\line(1,0){.985}}
\put(59.09,81.23){\line(1,0){.985}}
\put(61.06,81.22){\line(1,0){.985}}
\put(63.03,81.21){\line(1,0){.985}}
\put(65,81.2){\line(1,0){.985}}
\put(66.97,81.19){\line(1,0){.985}}
\put(68.94,81.18){\line(1,0){.985}}
%\end
%\dashline{1}(87.5,82.75)(87.5,81)
\put(87.43,82.68){\line(0,-1){.875}}
%\end
%\dashline{1}(87.5,81)(133.25,81.25)
\put(87.43,80.93){\line(1,0){.995}}
\put(89.42,80.94){\line(1,0){.995}}
\put(91.41,80.95){\line(1,0){.995}}
\put(93.4,80.96){\line(1,0){.995}}
\put(95.39,80.97){\line(1,0){.995}}
\put(97.38,80.98){\line(1,0){.995}}
\put(99.36,80.99){\line(1,0){.995}}
\put(101.35,81.01){\line(1,0){.995}}
\put(103.34,81.02){\line(1,0){.995}}
\put(105.33,81.03){\line(1,0){.995}}
\put(107.32,81.04){\line(1,0){.995}}
\put(109.31,81.05){\line(1,0){.995}}
\put(111.3,81.06){\line(1,0){.995}}
\put(113.29,81.07){\line(1,0){.995}}
\put(115.28,81.08){\line(1,0){.995}}
\put(117.27,81.09){\line(1,0){.995}}
\put(119.26,81.1){\line(1,0){.995}}
\put(121.24,81.11){\line(1,0){.995}}
\put(123.23,81.13){\line(1,0){.995}}
\put(125.22,81.14){\line(1,0){.995}}
\put(127.21,81.15){\line(1,0){.995}}
\put(129.2,81.16){\line(1,0){.995}}
\put(131.19,81.17){\line(1,0){.995}}
%\end
%\dashline{1}(87.5,124.25)(133.5,124)
\put(87.43,124.18){\line(1,0){.979}}
\put(89.39,124.17){\line(1,0){.979}}
\put(91.34,124.16){\line(1,0){.979}}
\put(93.3,124.15){\line(1,0){.979}}
\put(95.26,124.14){\line(1,0){.979}}
\put(97.22,124.13){\line(1,0){.979}}
\put(99.17,124.12){\line(1,0){.979}}
\put(101.13,124.11){\line(1,0){.979}}
\put(103.09,124.09){\line(1,0){.979}}
\put(105.05,124.08){\line(1,0){.979}}
\put(107,124.07){\line(1,0){.979}}
\put(108.96,124.06){\line(1,0){.979}}
\put(110.92,124.05){\line(1,0){.979}}
\put(112.88,124.04){\line(1,0){.979}}
\put(114.83,124.03){\line(1,0){.979}}
\put(116.79,124.02){\line(1,0){.979}}
\put(118.75,124.01){\line(1,0){.979}}
\put(120.71,124){\line(1,0){.979}}
\put(122.66,123.99){\line(1,0){.979}}
\put(124.62,123.98){\line(1,0){.979}}
\put(126.58,123.97){\line(1,0){.979}}
\put(128.54,123.96){\line(1,0){.979}}
\put(130.49,123.95){\line(1,0){.979}}
\put(132.45,123.94){\line(1,0){.979}}
%\end
%\dashline{1}(132,81)(134,81)
\put(131.93,80.93){\line(1,0){.667}}
\put(133.26,80.93){\line(1,0){.667}}
%\end
%\qbezvec(84.5,115.75)(87,117.75)(92.5,118.75)
\put(92.5,118.75){\vector(4,1){.07}}\qbezier(84.5,115.75)(87,117.75)(92.5,118.75)
%\end
%\qbezvec(14.5,116)(17.38,118.63)(25.75,116.75)
\put(25.75,116.75){\vector(4,-1){.07}}\qbezier(14.5,116)(17.38,118.63)(25.75,116.75)
%\end
\qbezier(32,103.5)(46.75,59.63)(58.5,103.25)
\qbezier(87.5,104.5)(98.5,143.5)(110.5,104.5)
\qbezier(110.5,104.25)(123.88,58.13)(132.75,104.5)
\qbezier(58.5,103.5)(61.25,114.25)(70,124)
\qbezier(32,104)(30.13,113)(19.75,124)
\put(14.5,111.5){\makebox(0,0)[cc]{$C_{1}$}}
\put(17,124){\makebox(0,0)[cc]{$1$}}
\put(17,103.75){\makebox(0,0)[cc]{$0$}}
\put(17,81.75){\makebox(0,0)[cc]{$-1$}}
\put(33.25,101){\makebox(0,0)[cc]{$\pi/2$}}
\put(45.75,101){\makebox(0,0)[cc]{$\pi$}}
\put(60.25,101){\makebox(0,0)[cc]{$3\pi/2$}}
\put(69.5,101){\makebox(0,0)[cc]{$2\pi$}}
\put(83,111.5){\makebox(0,0)[cc]{$S_{1}$}}
\put(85,124){\makebox(0,0)[cc]{$1$}}
\put(85,103.75){\makebox(0,0)[cc]{$0$}}
\put(85,81.75){\makebox(0,0)[cc]{$-1$}}
\put(99,101){\makebox(0,0)[cc]{$\pi/2$}}
\put(110.75,101){\makebox(0,0)[cc]{$\pi$}}
\put(122,101){\makebox(0,0)[cc]{$3\pi/2$}}
\put(132.75,101){\makebox(0,0)[cc]{$2\pi$}}
\put(77,73.75){\makebox(0,0)[cc]{Figure 5.1}}
\end{picture}
\end{center}

\vfill

Recall that, when we defined $C_{1}$ and $S_{1}$, we made them identical to $C$ and $S$ respectively on the interval from $0$ to $\pi$.

\vfill

Let's define our new functions $f$ and $g$ so that they are identical to $C_{1}$ and $S_{1}$ respectively on $[0,2\pi]$. Now what?

\vfill

Before you go on, see if you can come up with some scheme to define $f$ and $g$ over the rest of the $x$-axis so that they will satisfy the five basic properties.

\vspace{0.3cm}

\vfill

Consider the following scheme. Let's define $f$ such that on $[2\pi,4\pi]$ it is identical to $C_{1}$ on $[0,2\pi]$, and then on $[4\pi,6\pi]$ so that it also repeats what it is on $[0,2\pi]$. To carry this further, we define $f$ so that it is periodic with period $2\pi$; i.e., on any interval of length $2\pi$ whose left end is a multiple of $2\pi$, $f$ is identical to $C_{1}$ on $[0,2\pi]$. How do we say this algebraically?

\vfill

This is a nice puzzle, and you should work on it before you go on further.

\vspace{0.3cm}

\vfill

\textbf{Definition 5.1}: Suppose $f$ is the function to which the number pair $(x,y)$ belongs if and only if $y=C_{1}(x-n_{x}2\pi)$ where $n_{x}$ is the greatest integer $n$ such that $n(2\pi) \leq x$.

\vspace{0.3cm}

\vfill

This certainly sounds complicated enough. But study this picture. Note that $0 \leq x - n_{x}(2\pi) \leq 2\pi$, so that $C_{1}(x-n_{x}(2\pi))$ makes sense.

\vfill

Essentially what has happened is that we've divided the whole $x$-axis up into intervals of length $2\pi$ where the ends of the intervals are at multiples of $2\pi$. Now define $g$ analogously, relating it to $S_{1}$.

\begin{center}
%TeXCAD Options
%\grade{\on}
%\emlines{\off}
%\epic{\off}
%\beziermacro{\on}
%\reduce{\on}
%\snapping{\off}
%\quality{8.00}
%\graddiff{0.01}
%\snapasp{1}
%\zoom{4.0000}
\unitlength 1mm % = 2.85pt
\linethickness{0.4pt}
\ifx\plotpoint\undefined\newsavebox{\plotpoint}\fi % GNUPLOT compatibility
\begin{picture}(78,97.5)(0,26)
\put(14.5,102.75){\line(1,0){62.25}} 
\put(46.5,103){\line(0,-1){.5}}
\put(29.75,103){\line(0,-1){.5}} 
\put(63.25,103){\line(0,-1){.25}}
%\dashline{1}(13.5,123.25)(13.5,83.25)
\put(13.43,123.18){\line(0,-1){.976}}
\put(13.43,121.23){\line(0,-1){.976}}
\put(13.43,119.28){\line(0,-1){.976}}
\put(13.43,117.33){\line(0,-1){.976}}
\put(13.43,115.37){\line(0,-1){.976}}
\put(13.43,113.42){\line(0,-1){.976}}
\put(13.43,111.47){\line(0,-1){.976}}
\put(13.43,109.52){\line(0,-1){.976}}
\put(13.43,107.57){\line(0,-1){.976}}
\put(13.43,105.62){\line(0,-1){.976}}
\put(13.43,103.67){\line(0,-1){.976}}
\put(13.43,101.72){\line(0,-1){.976}}
\put(13.43,99.77){\line(0,-1){.976}}
\put(13.43,97.81){\line(0,-1){.976}}
\put(13.43,95.86){\line(0,-1){.976}}
\put(13.43,93.91){\line(0,-1){.976}}
\put(13.43,91.96){\line(0,-1){.976}}
\put(13.43,90.01){\line(0,-1){.976}}
\put(13.43,88.06){\line(0,-1){.976}}
\put(13.43,86.11){\line(0,-1){.976}}
\put(13.43,84.16){\line(0,-1){.976}}
%\end
%\dashline{1}(13.5,83.25)(77.25,83.25)
\put(13.43,83.18){\line(1,0){.996}}
\put(15.42,83.18){\line(1,0){.996}}
\put(17.41,83.18){\line(1,0){.996}}
\put(19.41,83.18){\line(1,0){.996}}
\put(21.4,83.18){\line(1,0){.996}}
\put(23.39,83.18){\line(1,0){.996}}
\put(25.38,83.18){\line(1,0){.996}}
\put(27.38,83.18){\line(1,0){.996}}
\put(29.37,83.18){\line(1,0){.996}}
\put(31.36,83.18){\line(1,0){.996}}
\put(33.35,83.18){\line(1,0){.996}}
\put(35.34,83.18){\line(1,0){.996}}
\put(37.34,83.18){\line(1,0){.996}}
\put(39.33,83.18){\line(1,0){.996}}
\put(41.32,83.18){\line(1,0){.996}}
\put(43.31,83.18){\line(1,0){.996}}
\put(45.3,83.18){\line(1,0){.996}}
\put(47.3,83.18){\line(1,0){.996}}
\put(49.29,83.18){\line(1,0){.996}}
\put(51.28,83.18){\line(1,0){.996}}
\put(53.27,83.18){\line(1,0){.996}}
\put(55.27,83.18){\line(1,0){.996}}
\put(57.26,83.18){\line(1,0){.996}}
\put(59.25,83.18){\line(1,0){.996}}
\put(61.24,83.18){\line(1,0){.996}}
\put(63.23,83.18){\line(1,0){.996}}
\put(65.23,83.18){\line(1,0){.996}}
\put(67.22,83.18){\line(1,0){.996}}
\put(69.21,83.18){\line(1,0){.996}}
\put(71.2,83.18){\line(1,0){.996}}
\put(73.2,83.18){\line(1,0){.996}}
\put(75.19,83.18){\line(1,0){.996}}
%\end
%\dashline{1}(77.25,83.25)(77.25,122.75)
\put(77.18,83.18){\line(0,1){.987}}
\put(77.18,85.15){\line(0,1){.988}}
\put(77.18,87.13){\line(0,1){.987}}
\put(77.18,89.1){\line(0,1){.987}}
\put(77.18,91.08){\line(0,1){.987}}
\put(77.18,93.05){\line(0,1){.987}}
\put(77.18,95.03){\line(0,1){.988}} 
\put(77.18,97){\line(0,1){.987}}
\put(77.18,98.98){\line(0,1){.987}}
\put(77.18,100.95){\line(0,1){.987}}
\put(77.18,102.93){\line(0,1){.987}}
\put(77.18,104.9){\line(0,1){.988}}
\put(77.18,106.88){\line(0,1){.987}}
\put(77.18,108.85){\line(0,1){.987}}
\put(77.18,110.83){\line(0,1){.987}}
\put(77.18,112.8){\line(0,1){.987}}
\put(77.18,114.78){\line(0,1){.988}}
\put(77.18,116.75){\line(0,1){.987}}
\put(77.18,118.73){\line(0,1){.987}}
\put(77.18,120.7){\line(0,1){.987}}
%\end
%\dashline{1}(77.25,122.75)(13.5,123.5)
\put(77.18,122.68){\line(-1,0){.996}}
\put(75.19,122.7){\line(-1,0){.996}}
\put(73.2,122.73){\line(-1,0){.996}}
\put(71.2,122.75){\line(-1,0){.996}}
\put(69.21,122.77){\line(-1,0){.996}}
\put(67.22,122.8){\line(-1,0){.996}}
\put(65.23,122.82){\line(-1,0){.996}}
\put(63.23,122.84){\line(-1,0){.996}}
\put(61.24,122.87){\line(-1,0){.996}}
\put(59.25,122.89){\line(-1,0){.996}}
\put(57.26,122.91){\line(-1,0){.996}}
\put(55.27,122.94){\line(-1,0){.996}}
\put(53.27,122.96){\line(-1,0){.996}}
\put(51.28,122.98){\line(-1,0){.996}}
\put(49.29,123.01){\line(-1,0){.996}}
\put(47.3,123.03){\line(-1,0){.996}}
\put(45.3,123.05){\line(-1,0){.996}}
\put(43.31,123.08){\line(-1,0){.996}}
\put(41.32,123.1){\line(-1,0){.996}}
\put(39.33,123.13){\line(-1,0){.996}}
\put(37.34,123.15){\line(-1,0){.996}}
\put(35.34,123.17){\line(-1,0){.996}}
\put(33.35,123.2){\line(-1,0){.996}}
\put(31.36,123.22){\line(-1,0){.996}}
\put(29.37,123.24){\line(-1,0){.996}}
\put(27.38,123.27){\line(-1,0){.996}}
\put(25.38,123.29){\line(-1,0){.996}}
\put(23.39,123.31){\line(-1,0){.996}}
\put(21.4,123.34){\line(-1,0){.996}}
\put(19.41,123.36){\line(-1,0){.996}}
\put(17.41,123.38){\line(-1,0){.996}}
\put(15.42,123.41){\line(-1,0){.996}}
%\end
%\dashline{1}(13.5,83)(13.5,36)
\put(13.43,82.93){\line(0,-1){.979}}
\put(13.43,80.97){\line(0,-1){.979}}
\put(13.43,79.01){\line(0,-1){.979}}
\put(13.43,77.05){\line(0,-1){.979}}
\put(13.43,75.1){\line(0,-1){.979}}
\put(13.43,73.14){\line(0,-1){.979}}
\put(13.43,71.18){\line(0,-1){.979}}
\put(13.43,69.22){\line(0,-1){.979}}
\put(13.43,67.26){\line(0,-1){.979}}
\put(13.43,65.3){\line(0,-1){.979}}
\put(13.43,63.35){\line(0,-1){.979}}
\put(13.43,61.39){\line(0,-1){.979}}
\put(13.43,59.43){\line(0,-1){.979}}
\put(13.43,57.47){\line(0,-1){.979}}
\put(13.43,55.51){\line(0,-1){.979}}
\put(13.43,53.55){\line(0,-1){.979}}
\put(13.43,51.6){\line(0,-1){.979}}
\put(13.43,49.64){\line(0,-1){.979}}
\put(13.43,47.68){\line(0,-1){.979}}
\put(13.43,45.72){\line(0,-1){.979}}
\put(13.43,43.76){\line(0,-1){.979}}
\put(13.43,41.8){\line(0,-1){.979}}
\put(13.43,39.85){\line(0,-1){.979}}
\put(13.43,37.89){\line(0,-1){.979}}
%\end
%\dashline{1}(13.5,36)(78,36)
\put(13.43,35.93){\line(1,0){.992}}
\put(15.41,35.93){\line(1,0){.992}}
\put(17.4,35.93){\line(1,0){.992}}
\put(19.38,35.93){\line(1,0){.992}}
\put(21.37,35.93){\line(1,0){.992}}
\put(23.35,35.93){\line(1,0){.992}}
\put(25.34,35.93){\line(1,0){.992}}
\put(27.32,35.93){\line(1,0){.992}}
\put(29.31,35.93){\line(1,0){.992}}
\put(31.29,35.93){\line(1,0){.992}}
\put(33.28,35.93){\line(1,0){.992}}
\put(35.26,35.93){\line(1,0){.992}}
\put(37.25,35.93){\line(1,0){.992}}
\put(39.23,35.93){\line(1,0){.992}}
\put(41.21,35.93){\line(1,0){.992}}
\put(43.2,35.93){\line(1,0){.992}}
\put(45.18,35.93){\line(1,0){.992}}
\put(47.17,35.93){\line(1,0){.992}}
\put(49.15,35.93){\line(1,0){.992}}
\put(51.14,35.93){\line(1,0){.992}}
\put(53.12,35.93){\line(1,0){.992}}
\put(55.11,35.93){\line(1,0){.992}}
\put(57.09,35.93){\line(1,0){.992}}
\put(59.08,35.93){\line(1,0){.992}}
\put(61.06,35.93){\line(1,0){.992}}
\put(63.05,35.93){\line(1,0){.992}}
\put(65.03,35.93){\line(1,0){.992}}
\put(67.01,35.93){\line(1,0){.992}} 
\put(69,35.93){\line(1,0){.992}}
\put(70.98,35.93){\line(1,0){.992}}
\put(72.97,35.93){\line(1,0){.992}}
\put(74.95,35.93){\line(1,0){.992}}
\put(76.94,35.93){\line(1,0){.992}}
%\end
%\dashline{1}(77.25,82.75)(77.75,36.25)
\put(77.18,82.68){\line(0,-1){.989}}
\put(77.2,80.7){\line(0,-1){.989}}
\put(77.22,78.72){\line(0,-1){.989}}
\put(77.24,76.74){\line(0,-1){.989}}
\put(77.26,74.76){\line(0,-1){.989}}
\put(77.29,72.79){\line(0,-1){.989}}
\put(77.31,70.81){\line(0,-1){.989}}
\put(77.33,68.83){\line(0,-1){.989}}
\put(77.35,66.85){\line(0,-1){.989}}
\put(77.37,64.87){\line(0,-1){.989}}
\put(77.39,62.89){\line(0,-1){.989}}
\put(77.41,60.91){\line(0,-1){.989}}
\put(77.44,58.94){\line(0,-1){.989}}
\put(77.46,56.96){\line(0,-1){.989}}
\put(77.48,54.98){\line(0,-1){.989}}
\put(77.5,53){\line(0,-1){.989}}
\put(77.52,51.02){\line(0,-1){.989}}
\put(77.54,49.04){\line(0,-1){.989}}
\put(77.56,47.06){\line(0,-1){.989}}
\put(77.58,45.08){\line(0,-1){.989}}
\put(77.61,43.11){\line(0,-1){.989}}
\put(77.63,41.13){\line(0,-1){.989}}
\put(77.65,39.15){\line(0,-1){.989}}
\put(77.67,37.17){\line(0,-1){.989}}
%\end
%\dashline{1}(13.5,71.25)(77.5,70.75)
\put(13.43,71.18){\line(1,0){.985}}
\put(15.4,71.16){\line(1,0){.985}}
\put(17.37,71.15){\line(1,0){.985}}
\put(19.34,71.13){\line(1,0){.985}}
\put(21.31,71.12){\line(1,0){.985}}
\put(23.28,71.1){\line(1,0){.985}}
\put(25.25,71.09){\line(1,0){.985}}
\put(27.21,71.07){\line(1,0){.985}}
\put(29.18,71.06){\line(1,0){.985}}
\put(31.15,71.04){\line(1,0){.985}}
\put(33.12,71.03){\line(1,0){.985}}
\put(35.09,71.01){\line(1,0){.985}} 
\put(37.06,71){\line(1,0){.985}}
\put(39.03,70.98){\line(1,0){.985}} 
\put(41,70.96){\line(1,0){.985}}
\put(42.97,70.95){\line(1,0){.985}}
\put(44.94,70.93){\line(1,0){.985}}
\put(46.91,70.92){\line(1,0){.985}}
\put(48.88,70.9){\line(1,0){.985}}
\put(50.85,70.89){\line(1,0){.985}}
\put(52.81,70.87){\line(1,0){.985}}
\put(54.78,70.86){\line(1,0){.985}}
\put(56.75,70.84){\line(1,0){.985}}
\put(58.72,70.83){\line(1,0){.985}}
\put(60.69,70.81){\line(1,0){.985}}
\put(62.66,70.8){\line(1,0){.985}}
\put(64.63,70.78){\line(1,0){.985}}
\put(66.6,70.76){\line(1,0){.985}}
\put(68.57,70.75){\line(1,0){.985}}
\put(70.54,70.73){\line(1,0){.985}}
\put(72.51,70.72){\line(1,0){.985}}
\put(74.48,70.7){\line(1,0){.985}}
\put(76.45,70.69){\line(1,0){.985}}
%\end
\put(13.5,53.5){\line(1,0){64}} 
\put(46.5,53.75){\line(0,-1){.75}}
\put(64.5,54){\line(0,-1){.75}} 
\put(28.75,53.75){\line(0,-1){.5}}
\qbezier(13.75,123.25)(28.25,116.88)(29.75,103)
\put(46.25,83.5){\line(0,-1){.5}} 
\put(46.5,36.5){\line(0,-1){.5}}
\qbezier(29.75,102.5)(34.88,87.13)(46.5,83.25)
\qbezier(46.5,83.25)(60.63,89.63)(63.25,102.5)
\qbezier(63.25,103)(65.13,118.13)(77.5,122.75)
\qbezier(13.5,71.25)(27.75,66.13)(29,53.5)
\qbezier(29,53.25)(33.75,39.88)(46.5,36)
\qbezier(46.75,36)(61.38,41.25)(64.5,53.5)
\qbezier(64.5,53.5)(67.5,66.88)(77.5,70.75)
\put(58.75,103.25){\line(0,-1){.5}} 
\put(59,93.75){\line(0,-1){1}}
\put(60,54){\line(0,-1){.75}} 
\put(60,45.25){\line(0,-1){.75}}
%\dashline{1}(58.75,102.75)(60,44.5)
\put(58.68,102.68){\line(0,-1){.987}}
\put(58.72,100.71){\line(0,-1){.987}}
\put(58.76,98.73){\line(0,-1){.987}}
\put(58.81,96.76){\line(0,-1){.987}}
\put(58.85,94.78){\line(0,-1){.987}}
\put(58.89,92.81){\line(0,-1){.987}}
\put(58.93,90.83){\line(0,-1){.987}}
\put(58.98,88.86){\line(0,-1){.987}}
\put(59.02,86.88){\line(0,-1){.987}}
\put(59.06,84.91){\line(0,-1){.987}}
\put(59.1,82.93){\line(0,-1){.987}}
\put(59.15,80.96){\line(0,-1){.987}}
\put(59.19,78.98){\line(0,-1){.987}}
\put(59.23,77.01){\line(0,-1){.987}}
\put(59.27,75.04){\line(0,-1){.987}}
\put(59.32,73.06){\line(0,-1){.987}}
\put(59.36,71.09){\line(0,-1){.987}}
\put(59.4,69.11){\line(0,-1){.987}}
\put(59.44,67.14){\line(0,-1){.987}}
\put(59.48,65.16){\line(0,-1){.987}}
\put(59.53,63.19){\line(0,-1){.987}}
\put(59.57,61.21){\line(0,-1){.987}}
\put(59.61,59.24){\line(0,-1){.987}}
\put(59.65,57.26){\line(0,-1){.987}}
\put(59.7,55.29){\line(0,-1){.987}}
\put(59.74,53.32){\line(0,-1){.987}}
\put(59.78,51.34){\line(0,-1){.987}}
\put(59.82,49.37){\line(0,-1){.987}}
\put(59.87,47.39){\line(0,-1){.987}}
\put(59.91,45.42){\line(0,-1){.987}}
%\end
%\qbezvec(53.5,56.75)(54.5,54.88)(58.5,54.5)
\put(58.5,54.5){\vector(4,-1){.07}}\qbezier(53.5,56.75)(54.5,54.88)(58.5,54.5)
%\end
%\qbezvec(27,91.75)(30.25,91.88)(32.5,94.5)
\put(32.5,94.5){\vector(1,1){.07}}\qbezier(27,91.75)(30.25,91.88)(32.5,94.5)
%\end
%\qbezvec(25.5,45.5)(27.13,47.88)(30.25,46.75)
\put(30.25,46.75){\vector(4,-1){.07}}\qbezier(25.5,45.5)(27.13,47.88)(30.25,46.75)
%\end
\put(13.5,99.75){\makebox(0,0)[cc]{$m_{x}(2\pi)$}}
\put(26,90.5){\makebox(0,0)[cc]{$f$}}
\put(58.75,105.25){\makebox(0,0)[cc]{$x$}}
\put(77.25,99.75){\makebox(0,0)[cc]{$(m_{x}+1)2\pi$}}
\put(62,91.75){\makebox(0,0)[cc]{$(x,f(x))$}}
\put(54.25,58.75){\makebox(0,0)[cc]{$x-m_{x}(2\pi)$}}
\put(46.5,51){\makebox(0,0)[cc]{$\pi$}}
\put(66.5,51.25){\makebox(0,0)[cc]{$3\pi/2$}}
\put(77.5,51){\makebox(0,0)[cc]{$2\pi$}}
\put(28,51.25){\makebox(0,0)[cc]{$\pi/2$}}
\put(13.25,50.75){\makebox(0,0)[cc]{$0$}}
\put(26.25,42.25){\makebox(0,0)[cc]{$C_{1}$}}
\put(67.5,44.25){\makebox(0,0)[cc]{$( x-m_{2}(2\pi), C_{1}(x-m_{2}(2\pi)) )$}}
\put(43.75,27.75){\makebox(0,0)[cc]{Figure 5.2}}
\end{picture}
\end{center}

\vfill

Of course, the big question now is, ``Is it true that $f$ and $g$ have the five basic properties?'' Do they? This will take you some time to prove, and you should know that $f$ and $g$ are very famous functions and have special names.

\vfill

$f$ is commonly known as the cosine function, and $g$ is known as the sine function. So instead of writing $f(x)=C_{1}(x-n_{x}(2\pi))$, we write $\cos x=C_{1}(x-n_{x}(2\pi))$, and instead of writing $g(x)=S_{1}(x-n_{x}(2\pi))$, we write $\sin x=S_{1}(x-n_{x}(2\pi))$. We read $\cos x$ as ``cosine x'', and we read $\sin x$ as ``sine x''.

\vfill

We already know what the graphs of these functions look like. The cosine graph is identical to $C_{1}$ on $[0,2\pi]$, and it repeats over the whole $x$-axis in every interval of length $2\pi$.

\begin{center}
%TeXCAD Options
%\grade{\on}
%\emlines{\off}
%\epic{\off}
%\beziermacro{\on}
%\reduce{\on}
%\snapping{\off}
%\quality{8.00}
%\graddiff{0.01}
%\snapasp{1}
%\zoom{4.0000}
\unitlength 0.8mm % = 2.85pt
\linethickness{0.4pt}
\ifx\plotpoint\undefined\newsavebox{\plotpoint}\fi % GNUPLOT compatibility
\begin{picture}(134.5,122.25)(0,16)
%\dashline{1}(14.75,119.75)(110,120)
\put(14.68,119.68){\line(1,0){.992}}
\put(16.66,119.68){\line(1,0){.992}}
\put(18.65,119.69){\line(1,0){.992}}
\put(20.63,119.7){\line(1,0){.992}}
\put(22.62,119.7){\line(1,0){.992}}
\put(24.6,119.71){\line(1,0){.992}}
\put(26.59,119.71){\line(1,0){.992}}
\put(28.57,119.72){\line(1,0){.992}}
\put(30.55,119.72){\line(1,0){.992}}
\put(32.54,119.73){\line(1,0){.992}}
\put(34.52,119.73){\line(1,0){.992}}
\put(36.51,119.74){\line(1,0){.992}}
\put(38.49,119.74){\line(1,0){.992}}
\put(40.48,119.75){\line(1,0){.992}}
\put(42.46,119.75){\line(1,0){.992}}
\put(44.45,119.76){\line(1,0){.992}}
\put(46.43,119.76){\line(1,0){.992}}
\put(48.41,119.77){\line(1,0){.992}}
\put(50.4,119.77){\line(1,0){.992}}
\put(52.38,119.78){\line(1,0){.992}}
\put(54.37,119.78){\line(1,0){.992}}
\put(56.35,119.79){\line(1,0){.992}}
\put(58.34,119.79){\line(1,0){.992}}
\put(60.32,119.8){\line(1,0){.992}}
\put(62.3,119.8){\line(1,0){.992}}
\put(64.29,119.81){\line(1,0){.992}}
\put(66.27,119.82){\line(1,0){.992}}
\put(68.26,119.82){\line(1,0){.992}}
\put(70.24,119.83){\line(1,0){.992}}
\put(72.23,119.83){\line(1,0){.992}}
\put(74.21,119.84){\line(1,0){.992}}
\put(76.2,119.84){\line(1,0){.992}}
\put(78.18,119.85){\line(1,0){.992}}
\put(80.16,119.85){\line(1,0){.992}}
\put(82.15,119.86){\line(1,0){.992}}
\put(84.13,119.86){\line(1,0){.992}}
\put(86.12,119.87){\line(1,0){.992}}
\put(88.1,119.87){\line(1,0){.992}}
\put(90.09,119.88){\line(1,0){.992}}
\put(92.07,119.88){\line(1,0){.992}}
\put(94.05,119.89){\line(1,0){.992}}
\put(96.04,119.89){\line(1,0){.992}}
\put(98.02,119.9){\line(1,0){.992}}
\put(100.01,119.9){\line(1,0){.992}}
\put(101.99,119.91){\line(1,0){.992}}
\put(103.98,119.91){\line(1,0){.992}}
\put(105.96,119.92){\line(1,0){.992}}
\put(107.95,119.92){\line(1,0){.992}}
%\end
%\dashline{1}(109.5,83.25)(14.25,83.5)
\put(109.43,83.18){\line(-1,0){.992}}
\put(107.45,83.18){\line(-1,0){.992}}
\put(105.46,83.19){\line(-1,0){.992}}
\put(103.48,83.2){\line(-1,0){.992}}
\put(101.49,83.2){\line(-1,0){.992}}
\put(99.51,83.21){\line(-1,0){.992}}
\put(97.52,83.21){\line(-1,0){.992}}
\put(95.54,83.22){\line(-1,0){.992}}
\put(93.55,83.22){\line(-1,0){.992}}
\put(91.57,83.23){\line(-1,0){.992}}
\put(89.59,83.23){\line(-1,0){.992}}
\put(87.6,83.24){\line(-1,0){.992}}
\put(85.62,83.24){\line(-1,0){.992}}
\put(83.63,83.25){\line(-1,0){.992}}
\put(81.65,83.25){\line(-1,0){.992}}
\put(79.66,83.26){\line(-1,0){.992}}
\put(77.68,83.26){\line(-1,0){.992}}
\put(75.7,83.27){\line(-1,0){.992}}
\put(73.71,83.27){\line(-1,0){.992}}
\put(71.73,83.28){\line(-1,0){.992}}
\put(69.74,83.28){\line(-1,0){.992}}
\put(67.76,83.29){\line(-1,0){.992}}
\put(65.77,83.29){\line(-1,0){.992}}
\put(63.79,83.3){\line(-1,0){.992}}
\put(61.8,83.3){\line(-1,0){.992}}
\put(59.82,83.31){\line(-1,0){.992}}
\put(57.84,83.32){\line(-1,0){.992}}
\put(55.85,83.32){\line(-1,0){.992}}
\put(53.87,83.33){\line(-1,0){.992}}
\put(51.88,83.33){\line(-1,0){.992}}
\put(49.9,83.34){\line(-1,0){.992}}
\put(47.91,83.34){\line(-1,0){.992}}
\put(45.93,83.35){\line(-1,0){.992}}
\put(43.95,83.35){\line(-1,0){.992}}
\put(41.96,83.36){\line(-1,0){.992}}
\put(39.98,83.36){\line(-1,0){.992}}
\put(37.99,83.37){\line(-1,0){.992}}
\put(36.01,83.37){\line(-1,0){.992}}
\put(34.02,83.38){\line(-1,0){.992}}
\put(32.04,83.38){\line(-1,0){.992}}
\put(30.05,83.39){\line(-1,0){.992}}
\put(28.07,83.39){\line(-1,0){.992}}
\put(26.09,83.4){\line(-1,0){.992}}
\put(24.1,83.4){\line(-1,0){.992}}
\put(22.12,83.41){\line(-1,0){.992}}
\put(20.13,83.41){\line(-1,0){.992}}
\put(18.15,83.42){\line(-1,0){.992}}
\put(16.16,83.42){\line(-1,0){.992}}
%\end
\put(3.5,101.25){\line(1,0){113.25}}
\put(14.25,101.75){\line(0,-1){.75}}
\put(21.75,101.75){\line(0,-1){.75}}
\put(28.5,101.75){\line(0,-1){1}} 
\put(35.25,102){\line(0,-1){1}}
\put(43,101.75){\line(0,-1){.75}}
\put(50.25,101.5){\line(0,-1){.75}}
\put(57.5,101.5){\line(0,-1){.75}}
\put(64.75,101.75){\line(0,-1){.75}}
\put(71.75,101.75){\line(0,-1){.75}}
\put(79,101.75){\line(0,-1){.75}} 
\put(92,102){\line(0,-1){1}}
\put(98.5,102){\line(0,-1){.5}} 
\put(98.5,101.5){\line(0,-1){.5}}
\put(104.75,102){\line(0,-1){1.25}} 
\put(110.25,102){\line(0,-1){1}}
\put(85.75,102){\line(0,-1){1.25}}
\qbezier(14.5,119.75)(21.38,113.5)(21.75,101.25)
\qbezier(21.75,101.25)(28.5,65.38)(35.25,101)
\qbezier(35.25,101.25)(43.25,138.25)(50.25,101.25)
\qbezier(50.25,101.25)(58.5,65.88)(64.75,101)
\qbezier(64.75,101.5)(71.63,138.13)(79,101.25)
\qbezier(79,101.25)(87,65.75)(92,101.25)
\qbezier(92,101.25)(97.13,138.25)(104.75,101.25)
\qbezier(104.75,101.25)(111.25,65.75)(116.75,101.25)
%\dashline{1}(109,119.75)(127.5,119.75)
\put(108.93,119.68){\line(1,0){.974}}
\put(110.88,119.68){\line(1,0){.974}}
\put(112.82,119.68){\line(1,0){.974}}
\put(114.77,119.68){\line(1,0){.974}}
\put(116.72,119.68){\line(1,0){.974}}
\put(118.67,119.68){\line(1,0){.974}}
\put(120.61,119.68){\line(1,0){.974}}
\put(122.56,119.68){\line(1,0){.974}}
\put(124.51,119.68){\line(1,0){.974}}
\put(126.46,119.68){\line(1,0){.974}}
%\end
%\dashline{1}(111.75,83.25)(128.75,83.25)
\put(111.68,83.18){\line(1,0){.944}}
\put(113.57,83.18){\line(1,0){.944}}
\put(115.46,83.18){\line(1,0){.944}}
\put(117.35,83.18){\line(1,0){.944}}
\put(119.24,83.18){\line(1,0){.944}}
\put(121.12,83.18){\line(1,0){.944}}
\put(123.01,83.18){\line(1,0){.944}}
\put(124.9,83.18){\line(1,0){.944}}
\put(126.79,83.18){\line(1,0){.944}}
%\end
\qbezier(116.75,101.25)(118.88,115.5)(126.5,119.75)
\put(126,101.75){\line(0,-1){.75}}
%\dashline{1}(14.5,64.13)(109.75,64.38)
\put(14.43,64.06){\line(1,0){.992}}
\put(16.41,64.06){\line(1,0){.992}}
\put(18.4,64.07){\line(1,0){.992}}
\put(20.38,64.08){\line(1,0){.992}}
\put(22.37,64.08){\line(1,0){.992}}
\put(24.35,64.09){\line(1,0){.992}}
\put(26.34,64.09){\line(1,0){.992}}
\put(28.32,64.1){\line(1,0){.992}} 
\put(30.3,64.1){\line(1,0){.992}}
\put(32.29,64.11){\line(1,0){.992}}
\put(34.27,64.11){\line(1,0){.992}}
\put(36.26,64.12){\line(1,0){.992}}
\put(38.24,64.12){\line(1,0){.992}}
\put(40.23,64.13){\line(1,0){.992}}
\put(42.21,64.13){\line(1,0){.992}}
\put(44.2,64.14){\line(1,0){.992}}
\put(46.18,64.14){\line(1,0){.992}}
\put(48.16,64.15){\line(1,0){.992}}
\put(50.15,64.15){\line(1,0){.992}}
\put(52.13,64.16){\line(1,0){.992}}
\put(54.12,64.16){\line(1,0){.992}}
\put(56.1,64.17){\line(1,0){.992}}
\put(58.09,64.17){\line(1,0){.992}}
\put(60.07,64.18){\line(1,0){.992}}
\put(62.05,64.18){\line(1,0){.992}}
\put(64.04,64.19){\line(1,0){.992}}
\put(66.02,64.2){\line(1,0){.992}}
\put(68.01,64.2){\line(1,0){.992}}
\put(69.99,64.21){\line(1,0){.992}}
\put(71.98,64.21){\line(1,0){.992}}
\put(73.96,64.22){\line(1,0){.992}}
\put(75.95,64.22){\line(1,0){.992}}
\put(77.93,64.23){\line(1,0){.992}}
\put(79.91,64.23){\line(1,0){.992}}
\put(81.9,64.24){\line(1,0){.992}}
\put(83.88,64.24){\line(1,0){.992}}
\put(85.87,64.25){\line(1,0){.992}}
\put(87.85,64.25){\line(1,0){.992}}
\put(89.84,64.26){\line(1,0){.992}}
\put(91.82,64.26){\line(1,0){.992}}
\put(93.8,64.27){\line(1,0){.992}}
\put(95.79,64.27){\line(1,0){.992}}
\put(97.77,64.28){\line(1,0){.992}}
\put(99.76,64.28){\line(1,0){.992}}
\put(101.74,64.29){\line(1,0){.992}}
\put(103.73,64.29){\line(1,0){.992}}
\put(105.71,64.3){\line(1,0){.992}}
\put(107.7,64.3){\line(1,0){.992}}
%\end
%\dashline{1}(109.25,27.63)(14,27.88)
\put(109.18,27.56){\line(-1,0){.992}}
\put(107.2,27.56){\line(-1,0){.992}}
\put(105.21,27.57){\line(-1,0){.992}}
\put(103.23,27.58){\line(-1,0){.992}}
\put(101.24,27.58){\line(-1,0){.992}}
\put(99.26,27.59){\line(-1,0){.992}}
\put(97.27,27.59){\line(-1,0){.992}}
\put(95.29,27.6){\line(-1,0){.992}}
\put(93.3,27.6){\line(-1,0){.992}}
\put(91.32,27.61){\line(-1,0){.992}}
\put(89.34,27.61){\line(-1,0){.992}}
\put(87.35,27.62){\line(-1,0){.992}}
\put(85.37,27.62){\line(-1,0){.992}}
\put(83.38,27.63){\line(-1,0){.992}}
\put(81.4,27.63){\line(-1,0){.992}}
\put(79.41,27.64){\line(-1,0){.992}}
\put(77.43,27.64){\line(-1,0){.992}}
\put(75.45,27.65){\line(-1,0){.992}}
\put(73.46,27.65){\line(-1,0){.992}}
\put(71.48,27.66){\line(-1,0){.992}}
\put(69.49,27.66){\line(-1,0){.992}}
\put(67.51,27.67){\line(-1,0){.992}}
\put(65.52,27.67){\line(-1,0){.992}}
\put(63.54,27.68){\line(-1,0){.992}}
\put(61.55,27.68){\line(-1,0){.992}}
\put(59.57,27.69){\line(-1,0){.992}}
\put(57.59,27.7){\line(-1,0){.992}}
\put(55.6,27.7){\line(-1,0){.992}}
\put(53.62,27.71){\line(-1,0){.992}}
\put(51.63,27.71){\line(-1,0){.992}}
\put(49.65,27.72){\line(-1,0){.992}}
\put(47.66,27.72){\line(-1,0){.992}}
\put(45.68,27.73){\line(-1,0){.992}}
\put(43.7,27.73){\line(-1,0){.992}}
\put(41.71,27.74){\line(-1,0){.992}}
\put(39.73,27.74){\line(-1,0){.992}}
\put(37.74,27.75){\line(-1,0){.992}}
\put(35.76,27.75){\line(-1,0){.992}}
\put(33.77,27.76){\line(-1,0){.992}}
\put(31.79,27.76){\line(-1,0){.992}}
\put(29.8,27.77){\line(-1,0){.992}}
\put(27.82,27.77){\line(-1,0){.992}}
\put(25.84,27.78){\line(-1,0){.992}}
\put(23.85,27.78){\line(-1,0){.992}}
\put(21.87,27.79){\line(-1,0){.992}}
\put(19.88,27.79){\line(-1,0){.992}}
\put(17.9,27.8){\line(-1,0){.992}}
\put(15.91,27.8){\line(-1,0){.992}}
%\end
\put(3.25,45.63){\line(1,0){113.25}}
\put(14,46.13){\line(0,-1){.75}} 
\put(21.5,46.13){\line(0,-1){.75}}
\put(28.25,46.13){\line(0,-1){1}} 
\put(35,46.38){\line(0,-1){1}}
\put(42.75,46.13){\line(0,-1){.75}} 
\put(50,45.88){\line(0,-1){.75}}
\put(57.25,45.88){\line(0,-1){.75}}
\put(64.5,46.13){\line(0,-1){.75}}
\put(71.5,46.13){\line(0,-1){.75}}
\put(78.75,46.13){\line(0,-1){.75}}
\put(91.75,46.38){\line(0,-1){1}} 
\put(98.25,46.38){\line(0,-1){.5}}
\put(98.25,45.88){\line(0,-1){.5}}
\put(104.5,46.38){\line(0,-1){1.25}} 
\put(110,46.38){\line(0,-1){1}}
\put(85.5,46.38){\line(0,-1){1.25}}
%\dashline{1}(108.75,64.13)(127.25,64.13)
\put(108.68,64.06){\line(1,0){.974}}
\put(110.63,64.06){\line(1,0){.974}}
\put(112.57,64.06){\line(1,0){.974}}
\put(114.52,64.06){\line(1,0){.974}}
\put(116.47,64.06){\line(1,0){.974}}
\put(118.42,64.06){\line(1,0){.974}}
\put(120.36,64.06){\line(1,0){.974}}
\put(122.31,64.06){\line(1,0){.974}}
\put(124.26,64.06){\line(1,0){.974}}
\put(126.21,64.06){\line(1,0){.974}}
%\end
%\dashline{1}(111.5,27.63)(128.5,27.63)
\put(111.43,27.56){\line(1,0){.944}}
\put(113.32,27.56){\line(1,0){.944}}
\put(115.21,27.56){\line(1,0){.944}}
\put(117.1,27.56){\line(1,0){.944}}
\put(118.99,27.56){\line(1,0){.944}}
\put(120.87,27.56){\line(1,0){.944}}
\put(122.76,27.56){\line(1,0){.944}}
\put(124.65,27.56){\line(1,0){.944}}
\put(126.54,27.56){\line(1,0){.944}}
%\end
\put(125.75,46.13){\line(0,-1){.75}}
\put(116.5,45.5){\line(1,0){18}} 
\put(117,101.25){\line(1,0){17.5}}
\qbezier(14,45.5)(21.88,82.5)(28.25,45.5)
\qbezier(28.25,45.5)(35.5,10.5)(42.75,45.5)
\qbezier(42.75,45.75)(50.5,82.38)(57.25,45.5)
\qbezier(57.25,45.5)(64.63,9.5)(71.5,45.5)
\qbezier(71.5,45.75)(78.5,82.88)(85.5,45.5)
\qbezier(85.5,45.5)(93.63,9.38)(98.25,45.75)
\qbezier(98.25,45.75)(104.88,82.75)(110,45.75)
\qbezier(110,45.75)(118.75,9.88)(125.5,45.5)
\put(118.25,46){\line(0,-1){1}}
\put(116.75,101.75){\line(0,-1){.75}}
\put(70,124.25){\makebox(0,0)[cc]{cosine}}
\put(14.25,98.5){\makebox(0,0)[cc]{$-4\pi$}}
\put(42.75,98.75){\makebox(0,0)[cc]{$-2\pi$}}
\put(72,98.25){\makebox(0,0)[cc]{$0$}}
\put(98.5,98.25){\makebox(0,0)[cc]{$4\pi$}}
\put(126,98.25){\makebox(0,0)[cc]{$6\pi$}}
\put(14,42.88){\makebox(0,0)[cc]{-4$\pi$}}
\put(42.5,43.13){\makebox(0,0)[cc]{-2$\pi$}}
\put(71.75,42.63){\makebox(0,0)[cc]{$0$}}
\put(98.25,42.63){\makebox(0,0)[cc]{4$\pi$}}
\put(125.75,42.63){\makebox(0,0)[cc]{6$\pi$}}
\put(70,68.5){\makebox(0,0)[cc]{sine}}
\put(70,19){\makebox(0,0)[cc]{Figure 5.3}}
\end{picture}
\end{center}

We rephrase the idea now, using the common names of our new functions.

\vspace{0.3cm}

\textbf{Theorem 5.1}:
\begin{enumerate}[\hspace{1cm}1)]
  \item $\cos^{2}x+\sin^{2}x=1$
  \item $\cos(x-y) = \cos x\cos y + \sin x\sin y$
  \item $\sin(x-y) = \sin x\cos y - \sin y\cos x$
  \item $\cos(x+y) = \cos x\cos y - \sin x\sin y$
  \item $\sin(x+y) = \sin x\cos y + \sin y\cos x$
\end{enumerate}
for each number $x$ and each number $y$.

Recall now that if $x$ is a number and $n_{x}$ is the greatest integer $n$ such that $n(2\pi)\leq x$, then $\cos x=C_{1}(x-n_{x}(2\pi))$ and $\sin x=S_{1}(x-n_{x}(2\pi))$.

Furthermore, we already know that the five basic properties are true about $C_{1}$ and $S_{1}$.

Do not go on further until you have worked very hard to prove the properties above are true about the cosine and sine functions. Don't get discouraged if this gives you some trouble. Help will arrive in the nick of time, but you must be ready for the help. Working on these and enduring moderate to mind-numbing frustration is excellent preparation for the help.

\vspace{0.3cm}

The following problems are designed to help you establish that the basic properties are true about the cosine and sine functions. If you have worked hard on this problem already (like you were supposed to do), then you will see how helpful the following ideas are.

\vspace{0.3cm}

Take the following statement for granted: If $x$ is a number, then there exists one and only one integer $n$ such that $n(2\pi) \leq x < (n+1)(2\pi)$. Do you see what this is saying? It says that, if $x$ is a number (any number what soever), then $x$ is in an interval of length $2\pi$ whose left end is a multiple of $2\pi$ and whose right end is the next multiple of $2\pi$. Note also that $x$ is not at the right end.

\begin{center}
%TeXCAD Options
%\grade{\on}
%\emlines{\off}
%\epic{\off}
%\beziermacro{\on}
%\reduce{\on}
%\snapping{\off}
%\quality{8.00}
%\graddiff{0.01}
%\snapasp{1}
%\zoom{4.0000}
\unitlength 1mm % = 2.85pt
\linethickness{0.4pt}
\ifx\plotpoint\undefined\newsavebox{\plotpoint}\fi % GNUPLOT compatibility
\begin{picture}(88.75,6.5)(0,123)
\put(11.5,128.75){\line(1,0){77.25}}
\put(18.75,129.25){\line(0,-1){1.25}}
\put(41.75,129.5){\line(0,-1){1.25}} 
\put(74,129.25){\line(0,-1){1}}
\put(18.5,125){\makebox(0,0)[cc]{$n(2\pi)$}}
\put(41.75,125){\makebox(0,0)[cc]{$x$}}
\put(73.75,125){\makebox(0,0)[cc]{$(n+1)(2\pi)$}}
\end{picture}
\end{center}

It states, furthermore, that there are not two such intervals; i.e., $n$ is unique. The multiple $(n)$ of $2\pi$ associated with the number $x$ is unique, and is denoted $n_{x}$. Hence, we have

\begin{center}
%TeXCAD Options
%\grade{\on}
%\emlines{\off}
%\epic{\off}
%\beziermacro{\on}
%\reduce{\on}
%\snapping{\off}
%\quality{8.00}
%\graddiff{0.01}
%\snapasp{1}
%\zoom{4.0000}
\unitlength 1mm % = 2.85pt
\linethickness{0.4pt}
\ifx\plotpoint\undefined\newsavebox{\plotpoint}\fi % GNUPLOT compatibility
\begin{picture}(101,6.75)(0,129)
\put(10.5,135){\line(1,0){90.5}}
\put(23.75,135.75){\line(0,-1){1.25}}
\put(44.25,135.5){\line(0,-1){1}}
\put(84.75,135.5){\line(0,-1){1.25}}
\put(23.75,131.5){\makebox(0,0)[cc]{$n_{x}(2\pi)$}}
\put(44.25,131.5){\makebox(0,0)[cc]{$x$}}
\put(85,130.75){\makebox(0,0)[cc]{$(n_{x}+1)(2\pi)$}}
\end{picture}
\end{center}

or $n_{x}(2\pi) \leq x < (n_{x}+1)(2\pi)$, which leads to $0 \leq x-n_{x}(2\pi) < 2\pi$.

How on earth do you go about proving something like this?

In view of the statement above, let $n_{x}$ denote the integer $n$ such that $0 \leq x -n_{x}(2\pi) < 2\pi$.

\vspace{0.3cm}

\textbf{Lemma 5.1.1}: Show that if $x$ is a number and $n$ is an integer, then $\cos(x+n(2\pi)) = \cos x$ and $\sin(x+n(2\pi)) = \sin x$.

\vspace{0.3cm}

\textbf{Lemma 5.1.2}: Show that if $x$ is a number, then $\cos(-x)=\cos x$ and $\sin(-x)=-\sin x$.

\vspace{0.3cm}

Note that the two properties in Lemma 5.1.2 would be trivial to show if we know that the five basic properties were true of the cosine and sine functions.

Now go back and work on Theorem 5.1.

\vspace{0.3cm}

\textbf{Problem 5.1}: Suppose that $x$ is a number and $n$ is an integer. By the use of Theorem 5.1, show that each of the following is true:
\begin{enumerate}[\hspace{1cm}1)]
  \item $\sin 2x = 2\sin x\cos x$
  \item $\cos 2x = \cos^{2}x - \sin^{2}x = 2\cos^{2}x - 1 = 1 - 2\sin^{2}x$
  \item $\sin (\frac{x}{2}) = \pm \sqrt{\frac{(1-\cos x)}{2}}$
  \item $\cos (\frac{x}{2}) = \pm \sqrt{\frac{(1+\cos x)}{2}}$
  \item $\sin 3x = \sin x(3-4\sin^{2}x)$
  \item $\cos 3x = \cos x(4\cos^{2}x-3)$
  \item $\cos (2\pi-x) = \cos x$
  \item $\sin (2\pi-x) = -\sin x$
  \item $\cos (\pi-x) = -\cos x$
  \item $\sin (\pi-x) = \sin x$
  \item $\cos (\frac{\pi}{2}-x) = \sin x$
  \item $\sin (\frac{\pi}{2}-x) = \cos x$
  \item $\cos (2n\pi) = 1$
  \item $\sin (2n\pi) = 0$
\end{enumerate}

\textbf{Problem 5.2}: Sketch the graph of each of the following on the interval $[-2\pi, 2\pi]$.
\begin{enumerate}[\hspace{1cm}1)]
  \item $y=\cos x$
  \item $y=\sin x$
  \item $y=-\cos x$
  \item $y=-\sin x$
  \item $y=2\cos x$
  \item $y=3\sin x$
  \item $y=-2\sin x$
  \item $y=(\frac{1}{2})\cos x$
  \item $y=\cos 2x$
  \item $y=\sin 3x$
  \item $y=\cos 6x$
  \item $y=2\cos 3x$
  \item $y=-3\sin 3x$
  \item $y=\sin (\frac{1}{2})x$
  \item $y=-2\cos (\frac{1}{3})x$
  \item $y=1+\cos x$
  \item $y=2-\sin x$
  \item $y=1+2\sin 3x$
  \item $y=2-3\sin (\frac{1}{4})x$
  \item $y=\sin (x+\frac{\pi}{2})$
  \item $y=\cos (x-\frac{\pi}{3})$
  \item $y=\sin (2x+\frac{\pi}{2})$
  \item $y=\cos (3x-\frac{\pi}{4})$
  \item $y=2\sin ((\frac{1}{2})x+\frac{\pi}{3})$
  \item $y=1-(\frac{3}{4}) \sin (2x-\frac{\pi}{4})$
  \item $y=a+\sin x$
  \item $y=b\sin x$
  \item $y=a+b\sin x$
\end{enumerate}

\textbf{Definition 5.2}:
\begin{enumerate}[\hspace{1cm}1)]
  \item tangent function: $\tan x = \frac{\sin x}{\cos x}$ if $\cos x \neq 0$
  \item cotangent function: $\cot x = \frac{\cos x}{\sin x}$ if $\sin x \neq 0$
  \item secant function: $\sec x = \frac{1}{\cos x}$ if $\cos x \neq 0$
  \item cosecant function: $\csc x = \frac{1}{\sin x}$ if $\sin x \neq 0$
\end{enumerate}

Sketch the graph of each of the functions above in the interval $[-2\pi,2\pi]$.

\vspace{0.3cm}

\textbf{Problem 5.3}: Find an expansion formula for each of the following in terms of the given function. For example, find an expansion formula for $\tan(x-y)$ in terms of $\tan x$ and $\tan y$.
\begin{enumerate}[\hspace{1cm}1)]
  \item $\tan (x-y)$
  \item $\tan (x+y)$
  \item $\cot (x-y)$
  \item $\cot (x+y)$
\end{enumerate}
Can you do the same for the secant and cosecant functions?

\vspace{0.3cm}

\textbf{Problem 5.4}: Show that each of the following is true:
\begin{enumerate}[\hspace{1cm}1)]
  \item if $\cos x \neq 0$, then $1+\tan^{2}x = \sec^{2}x$
  \item if $\sin x \neq 0$, then $1+\cot^{2}x = \csc^{2}x$
\end{enumerate}

\textbf{Problem 5.5}: Find some geometric interpretations of the sine and cosine functions on the unit circle with center at the origin.

\vspace{0.3cm}

\textbf{Problem 5.6}: Suppose $P$ and $Q$ are points of $K:x^{2}+y^{2}=1$. Consider each of the following.
\begin{enumerate}[\hspace{1cm}a)]
  \item If $l(\arc{PQ})=\frac{\pi}{3}$, what is $PQ$?
  \item If $PQ=\sqrt{2-\sqrt{2}}$, what is $l(\arc{PQ})$? 
  \item How would you go about finding $PQ$ if $l(\arc{PQ})=\frac{19\pi}{24}$? 
  \item How close can you get to $PQ$ if $l(\arc{PQ})=3$? 
  \item If $l(\arc{PQ})=5$, what is $PQ$?
\end{enumerate}

It appears that the accuracy of our solutions to some of the problems above is dependent on the table of values that we have for the cosine and sine functions. It would be nice if, given a number $x$, we had an algebraic formula for calculating $\cos x$ and $\sin x$. Of course, we wouldn't have to calculate these numbers except for $x$ in the interval $[0,2\pi]$, since we have the formulas $\cos(x+n(2\pi))=\cos x$ and $\sin(x+n(2\pi))=\sin x$ for each number $x$ and each integer $n$. 

Can you see how to narrow the interval of required calculation down from $[0,2\pi]$ even further? What is the smallest such interval needed in order to have all values of the sine and cosine functions just by using the five basic properties or those that we derived from them? The object is to make a table of values for cosine and sine that does not have more entries than necessary.

\vspace{0.3cm}

\textbf{Problem 5.7}: See if you can solve our main problem restricted to the circle $K:x^{2}+y^{2}=1$. That is, given two points $P$ and $Q$ of $K$, find a relationship between $PQ$ and $l(\arc{PQ})$.

\vspace{0.3cm}

\textbf{Problem 5.8}: Suppose $P:(\frac{3}{4},\frac{\sqrt{7}}{4})$.
\begin{enumerate}[\hspace{1cm}a)]
  \item Show $P$ is a point of $K:x^{2}+y^{2}=1$. 
  \item Calculate $PI$. 
  \item What is $l(\arc{PI})$?
\end{enumerate}

\textbf{Theorem 5.2}: If $A:(a,b)$ is a point of $K:x^{2}+y^{2}=1$, then $a = \cos \theta$ and $b = \sin \theta$ for some $0 \leq \theta \leq 2\pi$.

\vspace{0.3cm}

\textbf{Theorem 5.3}: Suppose $P:(u,v)$, $OP=r>0$, and $A:(a,b)$ is a point of $K:x^{2}+y^{2}=1$ such that $OAP$ or $OPA$. Show that $u=ra$ and $v=rb$.

\vspace{0.3cm}

\textbf{Corollary 5.3}: If $P:(u,v)$ and $OP=r>0$, then $u=r\cos\theta$ and $v = r \sin \theta$ for some $0 \leq \theta \leq 2\pi$.

\vspace{0.3cm}

Here is a particularly challenging puzzle. Suppose you have a piece of carpet that is rectangular in shape with dimensions $9 \times 12$. The units don't matter. The problem is the following: You are to cut the carpet into three pieces. One piece is $1 \times 8$ and the other two are such that, when fitted together, they form a square $10 \times 10$. No tricks, no subterfuge. This is a real and straight-forward problem that has a legitimate solution.


\newpage
\begin{center}
\section*{Chapter V Hints and Solutions}
\end{center}


\textbf{Lemma 5.1.1}:
\begin{enumerate}[\hspace{1cm}1)]
  \item We defined $n_{x}$ to be the least integer $n$ such that $n(2\pi) \leq x$. Hence, $n_{x}(2\pi) \leq x$, but $(n_{x}+1)2\pi > x$. What is the least integer $k$ such that $k(2\pi) \leq x+n(2\pi)$ where $n$ is an integer? Draw some pictures. Draw a picture of the interval that contains $x$; i.e.,
\end{enumerate}

\begin{center}
%TeXCAD Options
%\grade{\on}
%\emlines{\off}
%\epic{\off}
%\beziermacro{\on}
%\reduce{\on}
%\snapping{\off}
%\quality{8.00}
%\graddiff{0.01}
%\snapasp{1}
%\zoom{4.0000}
\unitlength 1mm % = 2.85pt
\linethickness{0.4pt}
\ifx\plotpoint\undefined\newsavebox{\plotpoint}\fi % GNUPLOT compatibility
\begin{picture}(101,6.75)(0,129)
\put(10.5,135){\line(1,0){90.5}}
\put(23.75,135.75){\line(0,-1){1.25}}
\put(44.25,135.5){\line(0,-1){1}}
\put(84.75,135.5){\line(0,-1){1.25}}
\put(23.75,131.5){\makebox(0,0)[cc]{$n_{x}(2\pi)$}}
\put(44.25,131.5){\makebox(0,0)[cc]{$x$}}
\put(85,130.75){\makebox(0,0)[cc]{$(n_{x}+1)(2\pi)$}}
\end{picture}
\end{center}

Now draw the same type interval for $x+n(2\pi)$. What is the left end of this interval?

\begin{enumerate}
  \item[\hspace{1cm}2)] $n_{x+n(2\pi)}$ is the least integer $k$ such that $k(2\pi) \leq x+n(2\pi)$, and you can show that $n_{x+n(2\pi)} = n_{x}+n$.
  \item[\hspace{1cm}3)] $\cos(x+n(2\pi)) = C_{1}[x+n(2\pi)-n_{x+n(2\pi)}(2\pi)]$ by definition.
  \item[\hspace{1cm}4)] $C_{1}[x+n(2\pi)-n_{x+n(2\pi)}(2\pi)] = C_{1}[x+n(2\pi)-(n_{x}+n)(2\pi)]$.
  \item[\hspace{1cm}5)] $\cos(x+n(2\pi)) = C_{1}[x-n_{x}(2\pi)] = \cos x$.
  \item[\hspace{1cm}6)] $\sin(x+n(2\pi)) = \sin x$ can be shown with the same procedure as with $\cos(x+n(2\pi))$.
\end{enumerate}

\textbf{Lemma 5.1.2}:
\begin{enumerate}[\hspace{1cm}1)]
  \item Suppose $n_{x}2\pi \leq x < (n_{x}+1)2\pi$. Either
    \begin{enumerate}[\hspace{1cm}a)]
      \item $n_{x}2\pi = x$, or
      \item $n_{x}2\pi < x$.
    \end{enumerate}
  \item Suppose (a); i.e. $n_{x}2\pi = x$. Then $-n_{x}2\pi = -x$, which you should be able to show means $n_{-x} = -n_{x}$.
  \item $\cos(-x) = C_{1}[-x-n_{-x}2\pi]$ by definition.
  \item Show that $-x-n_{-x}2\pi = 0$.
  \item Hence, $\cos(-x) = C_{1}(0) = 1$. Then show that $\cos x = C_{1}(0) = 1$ so that $\cos(-x) = \cos x$.
  \item Suppose (b); i.e. $n_{x}2\pi < x$. You can show that this means that $n_{-x} = -(n_{x}+1)$.
  \item Then $\cos(-x) = C_{1}[-x-n_{-x}2\pi] = C_{1}[-x+(n_{x}+1)2\pi]$.
  \item Show that $-x+(n_{x}+1)2\pi = 2\pi-(x-n_{x}2\pi)$.
  \item Show that $x-n_{x}2\pi$ is in $[0,2\pi]$ and that $2\pi-(x-n_{x}2\pi)$ is in $[0,2\pi]$.
  \item Recall (or show) that $C_{1}(2\pi-\theta) = C_{1}(\theta)$ so that $C_{1}[-x+(n_{x}+1)2\pi] = C_{1}[2\pi-(x-n_{x}2\pi)] = C_{1}(x-n_{x}2\pi)$, which is $\cos x$.
  \item $\sin(-x) = -\sin x$can be shown by the same procedure.
\end{enumerate}

\textbf{Theorem 5.1}:

\emph{Part 1.} Show that $\cos^{2}x + \sin^{2}x = 1$.
\begin{enumerate}[\hspace{1cm}1)]
  \item Suppose $x$ is a number and $0 \leq x-n_{x}(2\pi) < 2\pi$. Then $\cos x = C_{1}[x-n_{x}(2\pi)]$ and $\sin x = S_{1}[x-n_{x}(2\pi)]$.
  \item $\cos^{2}x + \sin^{2}x = C_{1}^{2}[x-n_{x}(2\pi)] + S_{1}^{2}[x-n_{x}(2\pi)] = 1$.
\end{enumerate}

\emph{Part 2.} Show $\cos(x-y) = \cos x\cos y + \sin x\sin y$.
\begin{enumerate}[\hspace{1cm}1)]
  \item Suppose that each of $x$ and $y$ is a number. Write the inequality involving $x$ and $n_{x}$, and the inequality involving $y$ and $n_{y}$. 
  \item Let $x-y = x-n_{x}(2\pi) + n_{x}(2\pi)-y + n_{y}(2\pi)-n_{y}(2\pi)$.
  \item $\cos(x-y) = \cos{[x-n_{x}(2\pi)] - [y-n_{y}(2\pi)] + (n_{x}-n_{y})2\pi}$.
  \item $\cos(x-y) = \cos{[x-n_{x}(2\pi)] - [y-n_{y}(2\pi)]}$, by Lemma 5.1.1.
  \item There are two cases.
    \begin{enumerate}
      \item $[x-n_{x}(2\pi)] - [y-n_{y}(2\pi)] \geq 0$.
      \item $[x-n_{x}(2\pi)] - [y-n_{y}(2\pi)] < 0$.
    \end{enumerate}
  \item Note that, if $0 \leq a < T$, $0 \leq b < T$, and $a-b \geq 0$, then $a-b < T$.
  \item If (a) in step 5, then by step 6, $0 \leq [x-n_{x}(2\pi)] - [y-n_{y}(2\pi)] < 2\pi$.
  \item By step 7, $\cos(x-y) = C_{1}{[x-n_{x}(2\pi)] - [y-n_{y}(2\pi)]}$.
  \item Expand $C_{1}$ in step 8, noting that $x-n_{x}(2\pi)$, $y-n_{y}(2\pi)$, and their difference $[x-n_{x}(2\pi)] - [y-n_{y}(2\pi)]$ are all in $[0,2\pi]$.
  \item $\cos(x-y) = \cos x\cos y + \sin x\sin y$ if (a) in step 5 is true.
  \item Suppose (b) in step 5 is true. Recall that $\cos(-x) = \cos x$.
  \item $\cos(x-y) = \cos(y-x)$.
  \item $[x-n_{x}(2\pi)] - [y-n_{y}(2\pi)] < 0$ implies that $[y-n_{y}(2\pi)] - [x-n_{x}(2\pi)] > 0$.
  \item Use the results just obtained and show that $\cos(x-y) = \cos x\cos y + \sin x\sin y$.
\end{enumerate}

\emph{Part 3.} Show that $\sin(x-y) = \sin x\cos y - \sin y\cos x$.
\begin{enumerate}[\hspace{1cm}1)]
  \item Show that $\sin(x-y) = \sin \{ [x-n_{x}(2\pi)]-[y-n_{y}(2\pi)] \}$.
  \item Consider two cases for $[x-n_{x}(2\pi)] - [y-n_{y}(2\pi)]$. Either it is non-negative or it is negative.
  \item Proceed as in showing that $\cos(x-y) = \cos x\cos y + \sin x\sin y$.
  \item Recall that $\sin(-x) = -\sin x$ when you get to the case where $[x-n_{x}(2\pi)] - [y-n_{y}(2\pi)]$ is negative.
\end{enumerate}

\emph{Part 4.} Show that $\cos(x+y) = \cos x\cos y - \sin x\sin y$.
\begin{enumerate}[\hspace{1cm}1)]
  \item Note that $x+y = x-(-y)$, and use previous results.
\end{enumerate}

\emph{Part 5.} Show that $\sin(x+y) = \sin x\cos y + \sin y\cos x$. Use the hint in 4).

\vspace{0.3cm}

\textbf{Problem 5.1}: You will recognize these properties. All of them follow from the five basic properties; i.e., Theorem 5.1. You might find Lemma 5.1.1 useful in some of the problems.

\vspace{0.3cm}

\textbf{Problem 5.2}: Sketching graphs of functions like $a \cos x$ or $a \sin x$ is fairly easy to do. But what do you do with something like $\cos bx$? If $b=2$, we have the example in Problem 5.2 to graph $y = \cos 2x$. You'll note that our alien friend is at a loss regarding graphing. But for us, pictures are extremely helpful aids.

\begin{enumerate}[\hspace{1cm}1)]
  \item Sketch the following basic framework.
\end{enumerate}

\begin{center}
%TeXCAD Options
%\grade{\on}
%\emlines{\off}
%\epic{\off}
%\beziermacro{\on}
%\reduce{\on}
%\snapping{\off}
%\quality{8.00}
%\graddiff{0.01}
%\snapasp{1}
%\zoom{4.0000}
\unitlength 1mm % = 2.85pt
\linethickness{0.4pt}
\ifx\plotpoint\undefined\newsavebox{\plotpoint}\fi % GNUPLOT compatibility
\begin{picture}(106.75,58.25)(0,80)
%\dashline{1}(14,136.5)(95.5,136.75)
\put(13.93,136.43){\line(1,0){.994}}
\put(15.92,136.44){\line(1,0){.994}}
\put(17.91,136.44){\line(1,0){.994}}
\put(19.89,136.45){\line(1,0){.994}}
\put(21.88,136.45){\line(1,0){.994}}
\put(23.87,136.46){\line(1,0){.994}}
\put(25.86,136.47){\line(1,0){.994}}
\put(27.84,136.47){\line(1,0){.994}}
\put(29.83,136.48){\line(1,0){.994}}
\put(31.82,136.48){\line(1,0){.994}}
\put(33.81,136.49){\line(1,0){.994}}
\put(35.8,136.5){\line(1,0){.994}}
\put(37.78,136.5){\line(1,0){.994}}
\put(39.77,136.51){\line(1,0){.994}}
\put(41.76,136.52){\line(1,0){.994}}
\put(43.75,136.52){\line(1,0){.994}}
\put(45.73,136.53){\line(1,0){.994}}
\put(47.72,136.53){\line(1,0){.994}}
\put(49.71,136.54){\line(1,0){.994}}
\put(51.7,136.55){\line(1,0){.994}}
\put(53.69,136.55){\line(1,0){.994}}
\put(55.67,136.56){\line(1,0){.994}}
\put(57.66,136.56){\line(1,0){.994}}
\put(59.65,136.57){\line(1,0){.994}}
\put(61.64,136.58){\line(1,0){.994}}
\put(63.62,136.58){\line(1,0){.994}}
\put(65.61,136.59){\line(1,0){.994}}
\put(67.6,136.59){\line(1,0){.994}}
\put(69.59,136.6){\line(1,0){.994}}
\put(71.58,136.61){\line(1,0){.994}}
\put(73.56,136.61){\line(1,0){.994}}
\put(75.55,136.62){\line(1,0){.994}}
\put(77.54,136.62){\line(1,0){.994}}
\put(79.53,136.63){\line(1,0){.994}}
\put(81.52,136.64){\line(1,0){.994}}
\put(83.5,136.64){\line(1,0){.994}}
\put(85.49,136.65){\line(1,0){.994}}
\put(87.48,136.66){\line(1,0){.994}}
\put(89.47,136.66){\line(1,0){.994}}
\put(91.45,136.67){\line(1,0){.994}}
\put(93.44,136.67){\line(1,0){.994}}
%\end
%\dashline{1}(14.25,84.5)(94,84.5)
\put(14.18,84.43){\line(1,0){.997}}
\put(16.17,84.43){\line(1,0){.997}}
\put(18.17,84.43){\line(1,0){.997}}
\put(20.16,84.43){\line(1,0){.997}}
\put(22.15,84.43){\line(1,0){.997}}
\put(24.15,84.43){\line(1,0){.997}}
\put(26.14,84.43){\line(1,0){.997}}
\put(28.14,84.43){\line(1,0){.997}}
\put(30.13,84.43){\line(1,0){.997}}
\put(32.12,84.43){\line(1,0){.997}}
\put(34.12,84.43){\line(1,0){.997}}
\put(36.11,84.43){\line(1,0){.997}}
\put(38.1,84.43){\line(1,0){.997}}
\put(40.1,84.43){\line(1,0){.997}}
\put(42.09,84.43){\line(1,0){.997}}
\put(44.09,84.43){\line(1,0){.997}}
\put(46.08,84.43){\line(1,0){.997}}
\put(48.07,84.43){\line(1,0){.997}}
\put(50.07,84.43){\line(1,0){.997}}
\put(52.06,84.43){\line(1,0){.997}}
\put(54.05,84.43){\line(1,0){.997}}
\put(56.05,84.43){\line(1,0){.997}}
\put(58.04,84.43){\line(1,0){.997}}
\put(60.04,84.43){\line(1,0){.997}}
\put(62.03,84.43){\line(1,0){.997}}
\put(64.02,84.43){\line(1,0){.997}}
\put(66.02,84.43){\line(1,0){.997}}
\put(68.01,84.43){\line(1,0){.997}} 
\put(70,84.43){\line(1,0){.997}}
\put(72,84.43){\line(1,0){.997}} 
\put(73.99,84.43){\line(1,0){.997}}
\put(75.99,84.43){\line(1,0){.997}}
\put(77.98,84.43){\line(1,0){.997}}
\put(79.97,84.43){\line(1,0){.997}}
\put(81.97,84.43){\line(1,0){.997}}
\put(83.96,84.43){\line(1,0){.997}}
\put(85.95,84.43){\line(1,0){.997}}
\put(87.95,84.43){\line(1,0){.997}}
\put(89.94,84.43){\line(1,0){.997}}
\put(91.94,84.43){\line(1,0){.997}}
%\end
%\dashline{1}(93.25,84.5)(95.25,84.5)
\put(93.18,84.43){\line(1,0){.667}}
\put(94.51,84.43){\line(1,0){.667}}
%\end
\put(57.75,83.75){\line(0,1){1}} 
\put(57.75,84.5){\line(0,1){52}}
\put(9.75,111.5){\line(1,0){91.25}}
%\dashline{1}(94.75,136.5)(99.5,136.5)
\put(94.68,136.43){\line(1,0){.95}}
\put(96.58,136.43){\line(1,0){.95}}
\put(98.48,136.43){\line(1,0){.95}}
%\end
%\dashline{1}(95.25,84.5)(100.25,84.5)
\put(95.18,84.43){\line(1,0){.833}}
\put(96.85,84.43){\line(1,0){.833}}
\put(98.51,84.43){\line(1,0){.833}}
%\end
\put(101,111.5){\line(1,0){1.75}}
\put(102.5,111.5){\line(1,0){4.25}} 
\put(8.5,111.5){\line(1,0){2}}
\put(57.5,124.5){\line(1,0){1}} 
\put(57.25,98.25){\line(1,0){1.25}}
\put(12,112){\line(0,-1){1.25}} 
\put(23,112){\line(0,-1){1.25}}
\put(35.25,112.25){\line(0,-1){1.25}}
%\emline(46,112.25)(46.25,111.25)
\multiput(46,112.25)(.03125,-.125){8}{\line(0,-1){.125}}
%\end
\put(68,112.25){\line(0,-1){1.25}}
\put(78.75,112.25){\line(0,-1){1.25}}
\put(89.25,112.25){\line(0,-1){1}}
\put(99.75,112.25){\line(0,-1){1.5}}
\put(58,139){\makebox(0,0)[cc]{$1$}}
\put(60.5,124.25){\makebox(0,0)[cc]{$+\frac{1}{2}$}}
\put(60.5,98){\makebox(0,0)[cc]{$-\frac{1}{2}$}}
\put(58,82){\makebox(0,0)[cc]{$-1$}}
\put(11.75,108.75){\makebox(0,0)[cc]{$-2\pi$}}
\put(22.75,109){\makebox(0,0)[cc]{$-3\pi/2$}}
\put(35,108.75){\makebox(0,0)[cc]{$-\pi$}}
\put(46,108.75){\makebox(0,0)[cc]{$-\pi/2$}}
\put(58,108.75){\makebox(0,0)[cc]{$0$}}
\put(68.25,108.5){\makebox(0,0)[cc]{$\pi/2$}}
\put(78.75,108.5){\makebox(0,0)[cc]{$\pi$}}
\put(89.25,108.25){\makebox(0,0)[cc]{$3\pi/2$}}
\put(99.5,108.75){\makebox(0,0)[cc]{$2\pi$}}
\end{picture}
\end{center}

\begin{enumerate}
  \item[\hspace{1cm}2)] How many cycles does $\cos x$ have in $[-2\pi,2\pi]$? How many cycles does $\cos 2x$ have in $[-2\pi,2\pi]$?
  \item[\hspace{1cm}3)] Sketch $\cos x$ in your picture from step 1. Then construct another basic framework near your first one. You are going to use this second one to sketch $\cos 2x$.
  \item[\hspace{1cm}4)] You realize from step 2 that for every cycle of $\cos x$, you will have 2 cycles for $\cos 2x$. $\cos x$ has one cycle on $[0,2\pi]$. Hence, $\cos 2x$ will have 2 cycles in $[0,2\pi]$. How many cycles would $\cos 3x$ have in $[0,2\pi]$? How many would $\cos (\frac{1}{2})x$ have? Have you got the idea?
  \item[\hspace{1cm}5)] $\cos x$ has amplitude $1$; i.e., the graph never gets more than $1$ away from the $x$-axis, and it gets to $1$ and $-1$ several times. How many times does $\cos x$ reach its maximum ($1$), and how many times does it reach its minimum ($-1$), in $[-2\pi,2\pi]$?
  \item[\hspace{1cm}6)] What is the amplitude of $\cos 2x$? How many times does it reach its maximum, and how many times does it reach its minimum, in $[-2\pi,2\pi]$?
  \item[\hspace{1cm}7)] Your graph of the maximum points ($x$), minimum points ($x$), and $0$ points ($0$) of $\cos 2x$ should look something like the following:
\end{enumerate}

%TeXCAD Picture [prob5.2hint2.pic]. Options:
%\grade{\on}
%\emlines{\off}
%\epic{\off}
%\beziermacro{\on}
%\reduce{\on}
%\snapping{\off}
%\quality{8.00}
%\graddiff{0.01}
%\snapasp{1}
%\zoom{4.0000}
\unitlength 0.85mm % = 2.85pt
\linethickness{0.4pt}
\ifx\plotpoint\undefined\newsavebox{\plotpoint}\fi % GNUPLOT compatibility
\begin{picture}(135,42.17)(8,96)
\put(10.25,117.25){\line(1,0){121.25}}
%\dashline{1}(15.75,137.75)(128,137.75)
\put(15.68,137.68){\line(1,0){.993}}
\put(17.67,137.68){\line(1,0){.993}}
\put(19.65,137.68){\line(1,0){.993}}
\put(21.64,137.68){\line(1,0){.993}}
\put(23.63,137.68){\line(1,0){.993}}
\put(25.61,137.68){\line(1,0){.993}}
\put(27.6,137.68){\line(1,0){.993}}
\put(29.59,137.68){\line(1,0){.993}}
\put(31.57,137.68){\line(1,0){.993}}
\put(33.56,137.68){\line(1,0){.993}}
\put(35.55,137.68){\line(1,0){.993}}
\put(37.53,137.68){\line(1,0){.993}}
\put(39.52,137.68){\line(1,0){.993}}
\put(41.51,137.68){\line(1,0){.993}}
\put(43.49,137.68){\line(1,0){.993}}
\put(45.48,137.68){\line(1,0){.993}}
\put(47.47,137.68){\line(1,0){.993}}
\put(49.45,137.68){\line(1,0){.993}}
\put(51.44,137.68){\line(1,0){.993}}
\put(53.43,137.68){\line(1,0){.993}}
\put(55.41,137.68){\line(1,0){.993}}
\put(57.4,137.68){\line(1,0){.993}}
\put(59.39,137.68){\line(1,0){.993}}
\put(61.37,137.68){\line(1,0){.993}}
\put(63.36,137.68){\line(1,0){.993}}
\put(65.35,137.68){\line(1,0){.993}}
\put(67.33,137.68){\line(1,0){.993}}
\put(69.32,137.68){\line(1,0){.993}}
\put(71.31,137.68){\line(1,0){.993}}
\put(73.29,137.68){\line(1,0){.993}}
\put(75.28,137.68){\line(1,0){.993}}
\put(77.27,137.68){\line(1,0){.993}}
\put(79.25,137.68){\line(1,0){.993}}
\put(81.24,137.68){\line(1,0){.993}}
\put(83.23,137.68){\line(1,0){.993}}
\put(85.22,137.68){\line(1,0){.993}}
\put(87.2,137.68){\line(1,0){.993}}
\put(89.19,137.68){\line(1,0){.993}}
\put(91.18,137.68){\line(1,0){.993}}
\put(93.16,137.68){\line(1,0){.993}}
\put(95.15,137.68){\line(1,0){.993}}
\put(97.14,137.68){\line(1,0){.993}}
\put(99.12,137.68){\line(1,0){.993}}
\put(101.11,137.68){\line(1,0){.993}}
\put(103.1,137.68){\line(1,0){.993}}
\put(105.08,137.68){\line(1,0){.993}}
\put(107.07,137.68){\line(1,0){.993}}
\put(109.06,137.68){\line(1,0){.993}}
\put(111.04,137.68){\line(1,0){.993}}
\put(113.03,137.68){\line(1,0){.993}}
\put(115.02,137.68){\line(1,0){.993}}
\put(117,137.68){\line(1,0){.993}}
\put(118.99,137.68){\line(1,0){.993}}
\put(120.98,137.68){\line(1,0){.993}}
\put(122.96,137.68){\line(1,0){.993}}
\put(124.95,137.68){\line(1,0){.993}}
\put(126.94,137.68){\line(1,0){.993}}
%\end
%\dashline{1}(15.75,97.5)(128,97.5)
\put(15.68,97.43){\line(1,0){.993}}
\put(17.67,97.43){\line(1,0){.993}}
\put(19.65,97.43){\line(1,0){.993}}
\put(21.64,97.43){\line(1,0){.993}}
\put(23.63,97.43){\line(1,0){.993}}
\put(25.61,97.43){\line(1,0){.993}}
\put(27.6,97.43){\line(1,0){.993}}
\put(29.59,97.43){\line(1,0){.993}}
\put(31.57,97.43){\line(1,0){.993}}
\put(33.56,97.43){\line(1,0){.993}}
\put(35.55,97.43){\line(1,0){.993}}
\put(37.53,97.43){\line(1,0){.993}}
\put(39.52,97.43){\line(1,0){.993}}
\put(41.51,97.43){\line(1,0){.993}}
\put(43.49,97.43){\line(1,0){.993}}
\put(45.48,97.43){\line(1,0){.993}}
\put(47.47,97.43){\line(1,0){.993}}
\put(49.45,97.43){\line(1,0){.993}}
\put(51.44,97.43){\line(1,0){.993}}
\put(53.43,97.43){\line(1,0){.993}}
\put(55.41,97.43){\line(1,0){.993}}
\put(57.4,97.43){\line(1,0){.993}}
\put(59.39,97.43){\line(1,0){.993}}
\put(61.37,97.43){\line(1,0){.993}}
\put(63.36,97.43){\line(1,0){.993}}
\put(65.35,97.43){\line(1,0){.993}}
\put(67.33,97.43){\line(1,0){.993}}
\put(69.32,97.43){\line(1,0){.993}}
\put(71.31,97.43){\line(1,0){.993}}
\put(73.29,97.43){\line(1,0){.993}}
\put(75.28,97.43){\line(1,0){.993}}
\put(77.27,97.43){\line(1,0){.993}}
\put(79.25,97.43){\line(1,0){.993}}
\put(81.24,97.43){\line(1,0){.993}}
\put(83.23,97.43){\line(1,0){.993}}
\put(85.22,97.43){\line(1,0){.993}}
\put(87.2,97.43){\line(1,0){.993}}
\put(89.19,97.43){\line(1,0){.993}}
\put(91.18,97.43){\line(1,0){.993}}
\put(93.16,97.43){\line(1,0){.993}}
\put(95.15,97.43){\line(1,0){.993}}
\put(97.14,97.43){\line(1,0){.993}}
\put(99.12,97.43){\line(1,0){.993}}
\put(101.11,97.43){\line(1,0){.993}}
\put(103.1,97.43){\line(1,0){.993}}
\put(105.08,97.43){\line(1,0){.993}}
\put(107.07,97.43){\line(1,0){.993}}
\put(109.06,97.43){\line(1,0){.993}}
\put(111.04,97.43){\line(1,0){.993}}
\put(113.03,97.43){\line(1,0){.993}}
\put(115.02,97.43){\line(1,0){.993}}
\put(117,97.43){\line(1,0){.993}}
\put(118.99,97.43){\line(1,0){.993}}
\put(120.98,97.43){\line(1,0){.993}}
\put(122.96,97.43){\line(1,0){.993}}
\put(124.95,97.43){\line(1,0){.993}}
\put(126.94,97.43){\line(1,0){.993}}
%\end
\put(19,118){\line(0,-1){1.5}}
\put(122.25,118){\line(0,-1){1.5}}
\put(71.5,118){\line(0,-1){1.5}}
\put(83.5,118){\line(0,-1){1.5}}
\put(96.75,118){\line(0,-1){1.5}}
\put(109.5,118){\line(0,-1){1.5}}
\put(32.25,118){\line(0,-1){1.5}}
\put(45.75,118){\line(0,-1){1.5}}
\put(58.75,118){\line(0,-1){1.5}}
\put(19,114){\makebox(0,0)[cc]{$-2\pi$}}
\put(32.25,114){\makebox(0,0)[cc]{$-3\pi/2$}}
\put(45.75,114){\makebox(0,0)[cc]{$-\pi$}}
\put(71.5,114){\makebox(0,0)[cc]{$0$}}
\put(25.75,117.25){\makebox(0,0)[cc]{$0$}}
\put(39,117.25){\makebox(0,0)[cc]{$0$}}
\put(52.5,117.25){\makebox(0,0)[cc]{$0$}}
\put(65.25,117.25){\makebox(0,0)[cc]{$0$}}
\put(77.5,117.25){\makebox(0,0)[cc]{$0$}}
\put(90.25,117.25){\makebox(0,0)[cc]{$0$}}
\put(103,117.25){\makebox(0,0)[cc]{$0$}}
\put(116,117.25){\makebox(0,0)[cc]{$0$}}
\put(83.5,114){\makebox(0,0)[cc]{$\pi/2$}}
\put(96.75,114){\makebox(0,0)[cc]{$\pi$}}
\put(109.25,114){\makebox(0,0)[cc]{$3\pi/2$}}
\put(122.25,114){\makebox(0,0)[cc]{$2\pi$}}
\put(59,114){\makebox(0,0)[cc]{$-\pi/2$}}
\put(10.5,138){\makebox(0,0)[cc]{$1$}}
\put(133,138){\makebox(0,0)[cc]{$1$}}
\put(9.75,97.5){\makebox(0,0)[cc]{$-1$}}
\put(132.25,97.5){\makebox(0,0)[cc]{$-1$}}
\put(19.08,138.17){\makebox(0,0)[cc]{$x$}}
\put(46.08,138.17){\makebox(0,0)[cc]{$x$}}
\put(31.83,97.92){\makebox(0,0)[cc]{$x$}}
\put(71.58,138.17){\makebox(0,0)[cc]{$x$}}
\put(59.33,97.67){\makebox(0,0)[cc]{$x$}}
\put(96.58,138.17){\makebox(0,0)[cc]{$x$}}
\put(84.33,97.67){\makebox(0,0)[cc]{$x$}}
\put(122.58,138.17){\makebox(0,0)[cc]{$x$}}
\put(110.33,97.67){\makebox(0,0)[cc]{$x$}}
\end{picture}


Since you know the shape of the graph, you need only find the maximum points, the minimum points, and the points where it crosses the $x$-axis ($0$ points). Sketching it from there is fairly easy. What do you do with something of the form $a \sin bx$? Let's consider an example. Let's sketch $y = 3 \sin 3x$ in $[-2\pi,2\pi]$.

\begin{enumerate}[\hspace{1cm}1)]
  \item How many cycles does it have in $[0,2\pi]$? 
  \item What is the amplitude? Note that the amplitude of $\sin 3x$ is $1$. Hence, the amplitude of $3 \sin 3x$ is $3$. This means that your basic framework will have to have upper limit $3$ and lower limit $-3$. 
  \item How many maximum points (at $3$) will the graph have in $[-2\pi,2\pi]$, how many minimum points (at $-3$), and how many $0$ points? Once you've located these, the rest is easy.
  \item Here is a sketch of the maximum points ($x$).
\end{enumerate}

%TeXCAD Picture [prob5.2hint3.pic]. Options:
%\grade{\on}
%\emlines{\off}
%\epic{\off}
%\beziermacro{\on}
%\reduce{\on}
%\snapping{\off}
%\quality{8.00}
%\graddiff{0.01}
%\snapasp{1}
%\zoom{4.0000}
\unitlength 0.75mm % = 2.85pt
\linethickness{0.4pt}
\ifx\plotpoint\undefined\newsavebox{\plotpoint}\fi % GNUPLOT compatibility
\begin{picture}(150,32.5)(7,104)
%\dashline{1}(11.5,136)(107,136)
\put(11.43,135.93){\line(1,0){.995}}
\put(13.42,135.93){\line(1,0){.995}}
\put(15.41,135.93){\line(1,0){.995}}
\put(17.4,135.93){\line(1,0){.995}}
\put(19.39,135.93){\line(1,0){.995}}
\put(21.38,135.93){\line(1,0){.995}}
\put(23.37,135.93){\line(1,0){.995}}
\put(25.36,135.93){\line(1,0){.995}}
\put(27.35,135.93){\line(1,0){.995}}
\put(29.34,135.93){\line(1,0){.995}}
\put(31.33,135.93){\line(1,0){.995}}
\put(33.32,135.93){\line(1,0){.995}}
\put(35.3,135.93){\line(1,0){.995}}
\put(37.29,135.93){\line(1,0){.995}}
\put(39.28,135.93){\line(1,0){.995}}
\put(41.27,135.93){\line(1,0){.995}}
\put(43.26,135.93){\line(1,0){.995}}
\put(45.25,135.93){\line(1,0){.995}}
\put(47.24,135.93){\line(1,0){.995}}
\put(49.23,135.93){\line(1,0){.995}}
\put(51.22,135.93){\line(1,0){.995}}
\put(53.21,135.93){\line(1,0){.995}}
\put(55.2,135.93){\line(1,0){.995}}
\put(57.19,135.93){\line(1,0){.995}}
\put(59.18,135.93){\line(1,0){.995}}
\put(61.17,135.93){\line(1,0){.995}}
\put(63.16,135.93){\line(1,0){.995}}
\put(65.15,135.93){\line(1,0){.995}}
\put(67.14,135.93){\line(1,0){.995}}
\put(69.13,135.93){\line(1,0){.995}}
\put(71.12,135.93){\line(1,0){.995}}
\put(73.11,135.93){\line(1,0){.995}}
\put(75.1,135.93){\line(1,0){.995}}
\put(77.09,135.93){\line(1,0){.995}}
\put(79.08,135.93){\line(1,0){.995}}
\put(81.07,135.93){\line(1,0){.995}}
\put(83.05,135.93){\line(1,0){.995}}
\put(85.04,135.93){\line(1,0){.995}}
\put(87.03,135.93){\line(1,0){.995}}
\put(89.02,135.93){\line(1,0){.995}}
\put(91.01,135.93){\line(1,0){.995}}
\put(93,135.93){\line(1,0){.995}}
\put(94.99,135.93){\line(1,0){.995}}
\put(96.98,135.93){\line(1,0){.995}}
\put(98.97,135.93){\line(1,0){.995}}
\put(100.96,135.93){\line(1,0){.995}}
\put(102.95,135.93){\line(1,0){.995}}
\put(104.94,135.93){\line(1,0){.995}}
%\end
%\dashline{1}(11.5,113.5)(107,113.5)
\put(11.43,113.43){\line(1,0){.995}}
\put(13.42,113.43){\line(1,0){.995}}
\put(15.41,113.43){\line(1,0){.995}}
\put(17.4,113.43){\line(1,0){.995}}
\put(19.39,113.43){\line(1,0){.995}}
\put(21.38,113.43){\line(1,0){.995}}
\put(23.37,113.43){\line(1,0){.995}}
\put(25.36,113.43){\line(1,0){.995}}
\put(27.35,113.43){\line(1,0){.995}}
\put(29.34,113.43){\line(1,0){.995}}
\put(31.33,113.43){\line(1,0){.995}}
\put(33.32,113.43){\line(1,0){.995}}
\put(35.3,113.43){\line(1,0){.995}}
\put(37.29,113.43){\line(1,0){.995}}
\put(39.28,113.43){\line(1,0){.995}}
\put(41.27,113.43){\line(1,0){.995}}
\put(43.26,113.43){\line(1,0){.995}}
\put(45.25,113.43){\line(1,0){.995}}
\put(47.24,113.43){\line(1,0){.995}}
\put(49.23,113.43){\line(1,0){.995}}
\put(51.22,113.43){\line(1,0){.995}}
\put(53.21,113.43){\line(1,0){.995}}
\put(55.2,113.43){\line(1,0){.995}}
\put(57.19,113.43){\line(1,0){.995}}
\put(59.18,113.43){\line(1,0){.995}}
\put(61.17,113.43){\line(1,0){.995}}
\put(63.16,113.43){\line(1,0){.995}}
\put(65.15,113.43){\line(1,0){.995}}
\put(67.14,113.43){\line(1,0){.995}}
\put(69.13,113.43){\line(1,0){.995}}
\put(71.12,113.43){\line(1,0){.995}}
\put(73.11,113.43){\line(1,0){.995}}
\put(75.1,113.43){\line(1,0){.995}}
\put(77.09,113.43){\line(1,0){.995}}
\put(79.08,113.43){\line(1,0){.995}}
\put(81.07,113.43){\line(1,0){.995}}
\put(83.05,113.43){\line(1,0){.995}}
\put(85.04,113.43){\line(1,0){.995}}
\put(87.03,113.43){\line(1,0){.995}}
\put(89.02,113.43){\line(1,0){.995}}
\put(91.01,113.43){\line(1,0){.995}}
\put(93,113.43){\line(1,0){.995}}
\put(94.99,113.43){\line(1,0){.995}}
\put(96.98,113.43){\line(1,0){.995}}
\put(98.97,113.43){\line(1,0){.995}}
\put(100.96,113.43){\line(1,0){.995}}
\put(102.95,113.43){\line(1,0){.995}}
\put(104.94,113.43){\line(1,0){.995}}
%\end
%\dashline{1}(106,136)(143.5,136)
\put(105.93,135.93){\line(1,0){.987}}
\put(107.9,135.93){\line(1,0){.987}}
\put(109.88,135.93){\line(1,0){.987}}
\put(111.85,135.93){\line(1,0){.987}}
\put(113.82,135.93){\line(1,0){.987}}
\put(115.8,135.93){\line(1,0){.987}}
\put(117.77,135.93){\line(1,0){.987}}
\put(119.75,135.93){\line(1,0){.987}}
\put(121.72,135.93){\line(1,0){.987}}
\put(123.69,135.93){\line(1,0){.987}}
\put(125.67,135.93){\line(1,0){.987}}
\put(127.64,135.93){\line(1,0){.987}}
\put(129.61,135.93){\line(1,0){.987}}
\put(131.59,135.93){\line(1,0){.987}}
\put(133.56,135.93){\line(1,0){.987}}
\put(135.53,135.93){\line(1,0){.987}}
\put(137.51,135.93){\line(1,0){.987}}
\put(139.48,135.93){\line(1,0){.987}}
\put(141.46,135.93){\line(1,0){.987}}
%\end
%\dashline{1}(106,113.5)(143.5,113.5)
\put(105.93,113.43){\line(1,0){.987}}
\put(107.9,113.43){\line(1,0){.987}}
\put(109.88,113.43){\line(1,0){.987}}
\put(111.85,113.43){\line(1,0){.987}}
\put(113.82,113.43){\line(1,0){.987}}
\put(115.8,113.43){\line(1,0){.987}}
\put(117.77,113.43){\line(1,0){.987}}
\put(119.75,113.43){\line(1,0){.987}}
\put(121.72,113.43){\line(1,0){.987}}
\put(123.69,113.43){\line(1,0){.987}}
\put(125.67,113.43){\line(1,0){.987}}
\put(127.64,113.43){\line(1,0){.987}}
\put(129.61,113.43){\line(1,0){.987}}
\put(131.59,113.43){\line(1,0){.987}}
\put(133.56,113.43){\line(1,0){.987}}
\put(135.53,113.43){\line(1,0){.987}}
\put(137.51,113.43){\line(1,0){.987}}
\put(139.48,113.43){\line(1,0){.987}}
\put(141.46,113.43){\line(1,0){.987}}
%\end
%\dashline{1}(142.75,136)(148,136)
\put(142.68,135.93){\line(1,0){.875}}
\put(144.43,135.93){\line(1,0){.875}}
\put(146.18,135.93){\line(1,0){.875}}
%\end
%\dashline{1}(142.75,113.5)(148,113.5)
\put(142.68,113.43){\line(1,0){.875}}
\put(144.43,113.43){\line(1,0){.875}}
\put(146.18,113.43){\line(1,0){.875}}
%\end
\put(9.25,136){\makebox(0,0)[cc]{$3$}}
\put(9.25,113.5){\makebox(0,0)[cc]{$-3$}}
\put(149.5,136){\makebox(0,0)[cc]{$3$}}
\put(149.5,113.5){\makebox(0,0)[cc]{$-3$}}
\put(84.25,122){\makebox(0,0)[cc]{$0$}}
\put(144.5,122){\makebox(0,0)[cc]{$2\pi$}}
\put(125.5,122){\makebox(0,0)[cc]{$4\pi/3$}}
\put(104.5,122){\makebox(0,0)[cc]{$2\pi/3$}}
\put(62,122){\makebox(0,0)[cc]{$-2\pi/3$}}
\put(38.75,122){\makebox(0,0)[cc]{$-4\pi/3$}}
\put(16,122){\makebox(0,0)[cc]{$-2\pi$}}
\put(16.25,136.5){\makebox(0,0)[cc]{$x$}}
\put(84.25,136.5){\makebox(0,0)[cc]{$x$}}
\put(104.75,136.5){\makebox(0,0)[cc]{$x$}}
\put(125.5,136.5){\makebox(0,0)[cc]{$x$}}
\put(144.75,136.5){\makebox(0,0)[cc]{$x$}}
\put(38,136.5){\makebox(0,0)[cc]{$x$}}
\put(62.25,136.5){\makebox(0,0)[cc]{$x$}}
\put(9.25,125.25){\line(1,0){140.75}}
\put(150,125.25){\line(0,1){0}}
%\dottedline(16.25,135.75)(16.25,125.25)
\multiput(16.18,135.68)(0,-.9545){12}{{\rule{.4pt}{.4pt}}}
%\end
%\dottedline(38.5,135.75)(38.5,125.25)
\multiput(38.43,135.68)(0,-.9545){12}{{\rule{.4pt}{.4pt}}}
%\end
%\dottedline(62,135.75)(62,125.25)
\multiput(61.93,135.68)(0,-.9545){12}{{\rule{.4pt}{.4pt}}}
%\end
%\dottedline(84.25,135.75)(84.25,125.25)
\multiput(84.18,135.68)(0,-.9545){12}{{\rule{.4pt}{.4pt}}}
%\end
%\dottedline(104.5,135.75)(104.5,125.25)
\multiput(104.43,135.68)(0,-.9545){12}{{\rule{.4pt}{.4pt}}}
%\end
%\dottedline(125.5,135.75)(125.5,125.25)
\multiput(125.43,135.68)(0,-.9545){12}{{\rule{.4pt}{.4pt}}}
%\end
%\dottedline(144.5,135.75)(144.5,125.25)
\multiput(144.43,135.68)(0,-.9545){12}{{\rule{.4pt}{.4pt}}}
%\end
%\dottedline(16.25,125.25)(16.25,125.25)
\multiput(16.18,125.18)(0,0){3}{{\rule{.4pt}{.4pt}}}
%\end
\end{picture}

Note that we have three cycles in $[0,2\pi]$. Hence we need to divide $[0,2\pi]$ into thirds. $[0,\frac{2\pi}{3}]$, $[\frac{2\pi}{3},\frac{4\pi}{3}]$, $[\frac{4\pi}{3},2\pi]$. Similarly with $[-2\pi,0]$. This gives us maximum points at $0$, $\frac{2\pi}{3}$, $\frac{4\pi}{3}$, and $\frac{6\pi}{3} = 2\pi$. We also get maximum points at $-\frac{2\pi}{3}$, $-\frac{4\pi}{3}$, and $-\frac{6\pi}{3} = -2\pi$. Finding the rest of the significant points is relatively easy. One minimum between each two successive maxima, and one $0$ point between each successive maximum and minimum point. 

\begin{enumerate}
  \item[\hspace{1cm}5)] Here is a sketch of the maxima ($x$), minima ($x$) and the $0$ points ($0$).
\end{enumerate}

%TeXCAD Picture [prob5.2hint4.pic]. Options:
%\grade{\on}
%\emlines{\off}
%\epic{\off}
%\beziermacro{\on}
%\reduce{\on}
%\snapping{\off}
%\quality{8.00}
%\graddiff{0.01}
%\snapasp{1}
%\zoom{4.0000}
\unitlength 0.75mm % = 2.85pt
\linethickness{0.4pt}
\ifx\plotpoint\undefined\newsavebox{\plotpoint}\fi % GNUPLOT compatibility
\begin{picture}(150,32.5)(7,104)
%\dashline{1}(11.5,136)(107,136)
\put(11.43,135.93){\line(1,0){.995}}
\put(13.42,135.93){\line(1,0){.995}}
\put(15.41,135.93){\line(1,0){.995}}
\put(17.4,135.93){\line(1,0){.995}}
\put(19.39,135.93){\line(1,0){.995}}
\put(21.38,135.93){\line(1,0){.995}}
\put(23.37,135.93){\line(1,0){.995}}
\put(25.36,135.93){\line(1,0){.995}}
\put(27.35,135.93){\line(1,0){.995}}
\put(29.34,135.93){\line(1,0){.995}}
\put(31.33,135.93){\line(1,0){.995}}
\put(33.32,135.93){\line(1,0){.995}}
\put(35.3,135.93){\line(1,0){.995}}
\put(37.29,135.93){\line(1,0){.995}}
\put(39.28,135.93){\line(1,0){.995}}
\put(41.27,135.93){\line(1,0){.995}}
\put(43.26,135.93){\line(1,0){.995}}
\put(45.25,135.93){\line(1,0){.995}}
\put(47.24,135.93){\line(1,0){.995}}
\put(49.23,135.93){\line(1,0){.995}}
\put(51.22,135.93){\line(1,0){.995}}
\put(53.21,135.93){\line(1,0){.995}}
\put(55.2,135.93){\line(1,0){.995}}
\put(57.19,135.93){\line(1,0){.995}}
\put(59.18,135.93){\line(1,0){.995}}
\put(61.17,135.93){\line(1,0){.995}}
\put(63.16,135.93){\line(1,0){.995}}
\put(65.15,135.93){\line(1,0){.995}}
\put(67.14,135.93){\line(1,0){.995}}
\put(69.13,135.93){\line(1,0){.995}}
\put(71.12,135.93){\line(1,0){.995}}
\put(73.11,135.93){\line(1,0){.995}}
\put(75.1,135.93){\line(1,0){.995}}
\put(77.09,135.93){\line(1,0){.995}}
\put(79.08,135.93){\line(1,0){.995}}
\put(81.07,135.93){\line(1,0){.995}}
\put(83.05,135.93){\line(1,0){.995}}
\put(85.04,135.93){\line(1,0){.995}}
\put(87.03,135.93){\line(1,0){.995}}
\put(89.02,135.93){\line(1,0){.995}}
\put(91.01,135.93){\line(1,0){.995}}
\put(93,135.93){\line(1,0){.995}}
\put(94.99,135.93){\line(1,0){.995}}
\put(96.98,135.93){\line(1,0){.995}}
\put(98.97,135.93){\line(1,0){.995}}
\put(100.96,135.93){\line(1,0){.995}}
\put(102.95,135.93){\line(1,0){.995}}
\put(104.94,135.93){\line(1,0){.995}}
%\end
%\dashline{1}(11.5,113.5)(107,113.5)
\put(11.43,113.43){\line(1,0){.995}}
\put(13.42,113.43){\line(1,0){.995}}
\put(15.41,113.43){\line(1,0){.995}}
\put(17.4,113.43){\line(1,0){.995}}
\put(19.39,113.43){\line(1,0){.995}}
\put(21.38,113.43){\line(1,0){.995}}
\put(23.37,113.43){\line(1,0){.995}}
\put(25.36,113.43){\line(1,0){.995}}
\put(27.35,113.43){\line(1,0){.995}}
\put(29.34,113.43){\line(1,0){.995}}
\put(31.33,113.43){\line(1,0){.995}}
\put(33.32,113.43){\line(1,0){.995}}
\put(35.3,113.43){\line(1,0){.995}}
\put(37.29,113.43){\line(1,0){.995}}
\put(39.28,113.43){\line(1,0){.995}}
\put(41.27,113.43){\line(1,0){.995}}
\put(43.26,113.43){\line(1,0){.995}}
\put(45.25,113.43){\line(1,0){.995}}
\put(47.24,113.43){\line(1,0){.995}}
\put(49.23,113.43){\line(1,0){.995}}
\put(51.22,113.43){\line(1,0){.995}}
\put(53.21,113.43){\line(1,0){.995}}
\put(55.2,113.43){\line(1,0){.995}}
\put(57.19,113.43){\line(1,0){.995}}
\put(59.18,113.43){\line(1,0){.995}}
\put(61.17,113.43){\line(1,0){.995}}
\put(63.16,113.43){\line(1,0){.995}}
\put(65.15,113.43){\line(1,0){.995}}
\put(67.14,113.43){\line(1,0){.995}}
\put(69.13,113.43){\line(1,0){.995}}
\put(71.12,113.43){\line(1,0){.995}}
\put(73.11,113.43){\line(1,0){.995}}
\put(75.1,113.43){\line(1,0){.995}}
\put(77.09,113.43){\line(1,0){.995}}
\put(79.08,113.43){\line(1,0){.995}}
\put(81.07,113.43){\line(1,0){.995}}
\put(83.05,113.43){\line(1,0){.995}}
\put(85.04,113.43){\line(1,0){.995}}
\put(87.03,113.43){\line(1,0){.995}}
\put(89.02,113.43){\line(1,0){.995}}
\put(91.01,113.43){\line(1,0){.995}}
\put(93,113.43){\line(1,0){.995}}
\put(94.99,113.43){\line(1,0){.995}}
\put(96.98,113.43){\line(1,0){.995}}
\put(98.97,113.43){\line(1,0){.995}}
\put(100.96,113.43){\line(1,0){.995}}
\put(102.95,113.43){\line(1,0){.995}}
\put(104.94,113.43){\line(1,0){.995}}
%\end
%\dashline{1}(106,136)(143.5,136)
\put(105.93,135.93){\line(1,0){.987}}
\put(107.9,135.93){\line(1,0){.987}}
\put(109.88,135.93){\line(1,0){.987}}
\put(111.85,135.93){\line(1,0){.987}}
\put(113.82,135.93){\line(1,0){.987}}
\put(115.8,135.93){\line(1,0){.987}}
\put(117.77,135.93){\line(1,0){.987}}
\put(119.75,135.93){\line(1,0){.987}}
\put(121.72,135.93){\line(1,0){.987}}
\put(123.69,135.93){\line(1,0){.987}}
\put(125.67,135.93){\line(1,0){.987}}
\put(127.64,135.93){\line(1,0){.987}}
\put(129.61,135.93){\line(1,0){.987}}
\put(131.59,135.93){\line(1,0){.987}}
\put(133.56,135.93){\line(1,0){.987}}
\put(135.53,135.93){\line(1,0){.987}}
\put(137.51,135.93){\line(1,0){.987}}
\put(139.48,135.93){\line(1,0){.987}}
\put(141.46,135.93){\line(1,0){.987}}
%\end
%\dashline{1}(106,113.5)(143.5,113.5)
\put(105.93,113.43){\line(1,0){.987}}
\put(107.9,113.43){\line(1,0){.987}}
\put(109.88,113.43){\line(1,0){.987}}
\put(111.85,113.43){\line(1,0){.987}}
\put(113.82,113.43){\line(1,0){.987}}
\put(115.8,113.43){\line(1,0){.987}}
\put(117.77,113.43){\line(1,0){.987}}
\put(119.75,113.43){\line(1,0){.987}}
\put(121.72,113.43){\line(1,0){.987}}
\put(123.69,113.43){\line(1,0){.987}}
\put(125.67,113.43){\line(1,0){.987}}
\put(127.64,113.43){\line(1,0){.987}}
\put(129.61,113.43){\line(1,0){.987}}
\put(131.59,113.43){\line(1,0){.987}}
\put(133.56,113.43){\line(1,0){.987}}
\put(135.53,113.43){\line(1,0){.987}}
\put(137.51,113.43){\line(1,0){.987}}
\put(139.48,113.43){\line(1,0){.987}}
\put(141.46,113.43){\line(1,0){.987}}
%\end
%\dashline{1}(142.75,136)(148,136)
\put(142.68,135.93){\line(1,0){.875}}
\put(144.43,135.93){\line(1,0){.875}}
\put(146.18,135.93){\line(1,0){.875}}
%\end
%\dashline{1}(142.75,113.5)(148,113.5)
\put(142.68,113.43){\line(1,0){.875}}
\put(144.43,113.43){\line(1,0){.875}}
\put(146.18,113.43){\line(1,0){.875}}
%\end
\put(9.25,136){\makebox(0,0)[cc]{$3$}}
\put(9.25,113.5){\makebox(0,0)[cc]{$-3$}}
\put(149.5,136){\makebox(0,0)[cc]{$3$}}
\put(149.5,113.5){\makebox(0,0)[cc]{$-3$}}
\put(84.25,122){\makebox(0,0)[cc]{$0$}}
\put(22,126){\makebox(0,0)[cc]{$0$}}
\put(33.5,126){\makebox(0,0)[cc]{$0$}}
\put(44,126){\makebox(0,0)[cc]{$0$}}
\put(55.5,126){\makebox(0,0)[cc]{$0$}}
\put(67.75,126){\makebox(0,0)[cc]{$0$}}
\put(79,126){\makebox(0,0)[cc]{$0$}}
\put(89.75,126){\makebox(0,0)[cc]{$0$}}
\put(100.5,126){\makebox(0,0)[cc]{$0$}}
\put(110.25,126){\makebox(0,0)[cc]{$0$}}
\put(120.5,126){\makebox(0,0)[cc]{$0$}}
\put(131,126){\makebox(0,0)[cc]{$0$}}
\put(140.5,126){\makebox(0,0)[cc]{$0$}}
\put(144.5,122){\makebox(0,0)[cc]{$2\pi$}}
\put(125.5,122){\makebox(0,0)[cc]{$4\pi/3$}}
\put(104.5,122){\makebox(0,0)[cc]{$2\pi/3$}}
\put(62,122){\makebox(0,0)[cc]{$-2\pi/3$}}
\put(38.75,122){\makebox(0,0)[cc]{$-4\pi/3$}}
\put(16,122){\makebox(0,0)[cc]{$-2\pi$}}
\put(16.25,136.5){\makebox(0,0)[cc]{$x$}}
\put(27.25,114){\makebox(0,0)[cc]{$x$}}
\put(84.25,136.5){\makebox(0,0)[cc]{$x$}}
\put(95.25,114){\makebox(0,0)[cc]{$x$}}
\put(104.75,136.5){\makebox(0,0)[cc]{$x$}}
\put(115.75,114){\makebox(0,0)[cc]{$x$}}
\put(125.5,136.5){\makebox(0,0)[cc]{$x$}}
\put(136.5,114){\makebox(0,0)[cc]{$x$}}
\put(144.75,136.5){\makebox(0,0)[cc]{$x$}}
\put(38,136.5){\makebox(0,0)[cc]{$x$}}
\put(49,114){\makebox(0,0)[cc]{$x$}}
\put(62.25,136.5){\makebox(0,0)[cc]{$x$}}
\put(73.25,114){\makebox(0,0)[cc]{$x$}}
\put(9.25,125.25){\line(1,0){140.75}}
\put(150,125.25){\line(0,1){0}}
%\dottedline(16.25,125.25)(16.25,125.25)
\multiput(16.18,125.18)(0,0){3}{{\rule{.4pt}{.4pt}}}
%\end
\put(16,124.5){\line(0,1){2.5}}
\put(27.25,124.5){\line(0,1){2.5}}
\put(38.75,124.5){\line(0,1){2.5}}
\put(49,124.5){\line(0,1){2.5}}
\put(62,124.5){\line(0,1){2.5}}
\put(73.25,124.5){\line(0,1){2.5}}
\put(84.25,124.5){\line(0,1){2.5}}
\put(95.25,124.5){\line(0,1){2.5}}
\put(104.75,124.5){\line(0,1){2.5}}
\put(115.75,124.5){\line(0,1){2.5}}
\put(125.5,124.5){\line(0,1){2.5}}
\put(136.5,124.5){\line(0,1){2.5}}
\put(144.75,124.5){\line(0,1){2.5}}
\end{picture}

Where are the minima? Where are the $0$ points?

\begin{enumerate}
  \item[\hspace{1cm}6)] Maxima: $-2\pi$, $-\frac{4\pi}{3}$, $-\frac{2\pi}{3}$, $0$, $\frac{2\pi}{3}$, $\frac{4\pi}{3}$, $2\pi$. 

      Minima:

      $\frac{[-2\pi+(-\frac{4\pi}{3})]}{2} = -\frac{10\pi}{6} = -\frac{5\pi}{3}$

      $\frac{[-\frac{4\pi}{3}+(-\frac{2\pi}{3})]}{2} = -\frac{6\pi}{6} = -\pi$

      $\frac{[-\frac{2\pi}{3}+0]}{2} = -\frac{\pi}{3}$

      Do you see how to get them? 

      $0$s:

      $\frac{[-2\pi+(-\frac{5\pi}{3})]}{2} = -\frac{11\pi}{6}$

      $\frac{[-\frac{5\pi}{3}+(-\frac{4\pi}{3})]}{2} = -\frac{9\pi}{6} = -\frac{3\pi}{2}$.

      Have you got it?
  \item[\hspace{1cm}7)] We really ought to check to see if our results are consistent with how the cosine function really works. We have claimed that $3 \sin 3x$ has a maximum at $-2\pi$, $-\frac{4\pi}{3}$, etc. This means that $3 \cos 3x = 3$ if $x$ is any of those values that we've listed. Is it true that $3 \cos 3x = 3$ if $x = -\frac{4\pi}{3}$; i.e., is it true that $3 \cos 3(-\frac{4\pi}{3}) = 3$? $3 \cos 3(-\frac{4\pi}{3}) = 3 \cos (-4\pi) = 3(1)$, by Problem 5.1 item 12. You might try some of the others for practice, including the minimum values ($-3$) and the $0$s.

       What do you do with something of the form $y = a+b \sin cx$? We've just examined the $b \sin cx$ part. What do we do with the $a$? When you look at the graph of $\sin x$, it oscillates about the $x$-axis. Adding $a$ to $b \sin cx$ merely has the graph of $b \sin cx$ oscillating about the horizontal line $y=a$ instead of about the $x$-axis ($y=0$).
\end{enumerate}

\textbf{Problem 5.3}: This just involves a little weird algebra. Consider the first one, $\tan (x-y)$.
\begin{enumerate}[\hspace{1cm}1)]
  \item $\tan (x-y) = \frac{\sin (x-y)}{\cos (x-y)}$.
  \item Expand $\sin (x-y)$ and $\cos (x-y)$.
  \item You should now have $\tan (x-y) = \frac{(\sin x \cos y - \sin y \cos x)}{(\cos x \cos y + \sin x \sin y)}$.
  \item Here is where the weird algebra comes in. We want to end up with everything in terms of $\tan x$ and $\tan y$. Divide both the numerator and the denominator of the expression in step 3 by $\cos x\cos y$.
  \item You should end up with $\tan (x-y)= \frac{(\tan x - \tan y)}{(1 + \tan x \tan y)}$. 
\end{enumerate}
The same witchcraft will work for the other expansion problems in Problem 5.3. You might have some difficulty when you try to get expansion formulas for $\sec (x+y)$ in terms of $\sec x$ and $\sec y$ and for $\csc(x+y)$ in terms of $\csc x$ and $\csc y$. 

You might try the following strategy to get a hint about how to get what you want. See if you can express all six of these functions (sine, cosine, tangent, cotangent, cosecant, and secant) in terms of one another. For example, find all six in terms of the sine function. You have five relationships to establish, since $\sin x = \sin x$. Then find $\cos x$ in terms of $\sin x$. Recall $\sin^{2} x + \cos^{2} x = 1$. This will yield $\cos x = \pm \sqrt{1 - \sin^{2} x}$. Now find $\tan x$ in terms of $\sin x$. Do you see the idea? 

Now you have to find $\cot x$, $\sec x$, and $\csc x$, each in terms of $\sin x$. Then you have five other problems like this. Finding each of the six functions in terms of $\cos x$, finding each of the six functions in terms of $\tan x$, etc. This should keep you quiet and busy long enough that your mother will check to see what you're doing.

\vspace{0.3cm}

\textbf{Problem 5.4}: These relationships are derivable from the first basic relationship, $\sin^{2} x + \cos^{2} x = 1$. Dividing through by the proper term will yield what you want.

\vspace{0.3cm}

\textbf{Problem 5.5}:
\begin{enumerate}[\hspace{1cm}1)]
  \item Use the definitions of $\cos x$ and $\sin x$. 
  \item Note that the definitions of $\sin x$ and $\cos x$ are in terms of $C_{1}$ and $S_{1}$, for which you already have geometric interpretations.
\end{enumerate}

\textbf{Problem 5.6}:

\begin{enumerate}[\hspace{1cm}a)]
  \item Suppose $l(\arc{PQ})=\frac{\pi}{3}$. Find $PQ$.
    \begin{enumerate}[\hspace{1cm}1)]
      \item Find the coordinates of the point $A$ of $\arc{JiI}$ such that $l(\arc{AI}) = \frac{\pi}{3}$. 
      \item Since $l(\arc{AI})=l(\arc{PQ})$, one of your axioms gives you $AI=PQ$. 
      \item Use the coordinates of $A$, which you can find in one of your tables, to find $AI$. You should get $AI=PQ=1$.
    \end{enumerate}
  \item $PQ=\sqrt{2-\sqrt{2}}$. Find $l(\arc{PQ})$.
    \begin{enumerate}[\hspace{1cm}1)]
      \item This chord length is in one of your early tables.
      \item You should find that $l(\arc{PQ})=\frac{\pi}{4}$.
    \end{enumerate}
  \item $l(\arc{PQ}) = \frac{19\pi}{24}$. Find $PQ$.
    \begin{enumerate}[\hspace{1cm}1)]
      \item One of your tables is in increments of $\frac{\pi}{24}$. 
      \item Find the coordinates of the point $A$ of $K$ such that 

          $A:(\cos \frac{19\pi}{24}, \sin \frac{19\pi}{24})$.
      \item A picture can help.
     \end{enumerate}
\end{enumerate}

\begin{center}
%TeXCAD Options
%\grade{\on}
%\emlines{\off}
%\epic{\off}
%\beziermacro{\on}
%\reduce{\on}
%\snapping{\off}
%\quality{8.00}
%\graddiff{0.01}
%\snapasp{1}
%\zoom{4.0000}
\unitlength 1mm % = 2.85pt
\linethickness{0.4pt}
\ifx\plotpoint\undefined\newsavebox{\plotpoint}\fi % GNUPLOT compatibility
\begin{picture}(100,45.5)(0,88)
\put(12.75,92.75){\line(1,0){87.25}}
\qbezier(58.5,131)(24.13,130.13)(20.25,92.75)
\qbezier(58.75,130.75)(89.25,127.25)(90.75,92.75)
\put(56.75,131.25){\line(0,-1){.75}}
%\emline(56.75,130.5)(57,130)
\multiput(56.75,130.5)(.03125,-.0625){8}{\line(0,-1){.0625}}
%\end
%\emline(23,108)(24.25,107.5)
\multiput(23,108)(.083333,-.033333){15}{\line(1,0){.083333}}
%\end
\put(20.5,108.25){\makebox(0,0)[cc]{$A$}}
\put(20,89.5){\makebox(0,0)[cc]{$J:(\cos 24\pi / 24, \sin 24\pi / 24)$}}
\put(90.75,89.5){\makebox(0,0)[cc]{$I:(\cos 0, \sin 0)$}}
\put(56.75,133.5){\makebox(0,0)[cc]{$i:(\cos 12\pi / 24, \sin 12\pi / 24)$}}
\end{picture}
\end{center}

Wherever $P$ and $Q$ are on $K$, the arc length from one to the other is the same as the arc length from $A$ to $I$; i.e., $l(\arc{PQ})=l(\arc{AI})$. Hence, from our main axiom, $PQ=AI$.

\begin{enumerate}
  \item[\hspace{2cm}4)] $\cos \frac{19\pi}{24} = \left( -\frac{\sqrt{2}}{4} \right) \left( \sqrt{2+\sqrt{2+\sqrt{3}}} + \sqrt{2-\sqrt{2+\sqrt{3}}} \right)$ and 
      
      $\sin \frac{19\pi}{24} = \left( \frac{\sqrt{2}}{4} \right) \left(\sqrt{2+\sqrt{2+\sqrt{3}}} - \sqrt{2-\sqrt{2+\sqrt{3}}} \right)$.
  \item[\hspace{2cm}5)] Where did I get $\cos \frac{19\pi}{24}$ and $\sin \frac{19\pi}{24}$? You may have these in your table from a previous chapter. If not, consider the following: $\cos \frac{19\pi}{24} = \cos (\frac{18\pi}{24} + \frac{\pi}{24}) = \cos (\frac{3\pi}{4} + \frac{\pi}{24})$. Expand this using your expansion formula. Now the problem is to find $\cos \frac{\pi}{24}$ and $\sin \frac{\pi}{24}$. Recall $\cos \frac{x}{2} = \pm \sqrt{\frac{(1+\cos x)}{2}}$ and $\sin \frac{x}{2} = \pm \sqrt{\frac{(1-\cos x)}{2}}$. If $0 \leq x \leq \pi$ (and note that $0 \leq \frac{\pi}{12} \leq \pi$), then $\cos \frac{x}{2} = \sqrt{\frac{(1+\cos x)}{2}}$ and $\sin \frac{x}{2} = \sqrt{\frac{(1-\cos x)}{2}}$. Since you have $\cos \frac{\pi}{12}$, you can get $\cos \frac{\pi}{24}$ and $\sin \frac{\pi}{24}$. $\cos \frac{\pi}{24} = (\frac{1}{2}) \sqrt{2+\sqrt{2+\sqrt{3}}}$ and $\sin \frac{\pi}{24} = (\frac{1}{2}) \sqrt{2-\sqrt{2+\sqrt{3}}}$.
  \item[\hspace{1cm}6)] You have $l(\arc{PQ}) = l(\arc{AI})$, so that $PQ=AI$. You also have $A:(\cos \frac{19\pi}{24},$ $\sin \frac{19\pi}{24})$, and you have numerical values for $\cos \frac{19\pi}{24}$ and $\sin \frac{19\pi}{24}$. All you need to do now is calculate $AI$.
  \item[\hspace{1cm}7)] $AI = \sqrt{(\cos \frac{19\pi}{24}-1)^{2} + (\sin \frac{19\pi}{24}-0)^{2}}$

      $\cos \frac{19\pi}{24} = -\frac{\sqrt{2}}{4} \left( \sqrt{2+\sqrt{2+\sqrt{3}}} + \sqrt{2-\sqrt{2+\sqrt{3}}} \right)$

      $\sin \frac{19\pi}{24} = \frac{\sqrt{2}}{4} \left( \sqrt{2+\sqrt{2+\sqrt{3}}} - \sqrt{2-\sqrt{2+\sqrt{3}}} \right)$

      You should get 
      \begin{equation*}
      AI = \sqrt{2+\frac{\sqrt{2}}{2} \left( \sqrt{2+\sqrt{2+\sqrt{3}}} + \sqrt{2-\sqrt{2+\sqrt{3}}} \right)}
      \end{equation*}
      which is $1.8939$ to four decimals.
\end{enumerate}

\begin{enumerate}
  \item[\hspace{1cm}d)] If $l(\arc{PQ})=3$, what is $PQ$? Note that the accuracy to which we can calculate $PQ$ is dependent on how dense or how fine our table of values for $\cos x$ and $\sin x$ is.
    \begin{enumerate}[\hspace{1cm}1)]
      \item Note that $\frac{22\pi}{24} < 3 < \frac{23\pi}{24}$. If $A:(\cos \frac{22\pi}{24}, \sin \frac{22\pi}{24})$ 

          and $B:(\cos \frac{23\pi}{24}, \sin \frac{23\pi}{24})$, then $AI < PQ < BI$.
      \item Draw a picture placing $A$ and $B$ on $\arc{JiI}$, and $P$ and $Q$ somewhere on $K:x^{2}+y^{2}=1$. This should help you see the logic of step 1.
      \item Show that $\cos \frac{22\pi}{24} = -(\frac{1}{2}) \sqrt{2+\sqrt{3}}$

          and $\sin \frac{22\pi}{24} = (\frac{1}{2}) \sqrt{2-\sqrt{3}}$. You probably have already done this in a previous chapter.
      \item Show that $\cos \frac{23\pi}{24} = -(\frac{1}{2}) \sqrt{2+\sqrt{2+\sqrt{3}}}$

          and $\sin \frac{23\pi}{24} = (\frac{1}{2}) \sqrt{2-\sqrt{2+\sqrt{3}}}$.
      \item If you don't see how to get $\cos \frac{23\pi}{24}$ and $\sin \frac{23\pi}{24}$, note that $\frac{23\pi}{24} = \pi - \frac{\pi}{24}$.
      \item $AI = \sqrt{\cos[\frac{22\pi}{24}-1]^{2} + \sin[\frac{22\pi}{24}-0]^{2}}$, 

          which turns out to be $\sqrt{2+\sqrt{2+\sqrt{3}}}$.
      \item $BI$ turns out to be $\sqrt{2+\sqrt{2+\sqrt{2+\sqrt{3}}}}$.
      \item Hence, $1.982 < PQ < 1.996$, and we accomplished this with a table of cosines and sines in increments of $\frac{\pi}{24}$. Are you curious how much closer you could get to $PQ$ if you used increments of $\frac{\pi}{48}$?
      \item Can you see how to use a scientific calculator to get your results? Recall that if $T:(\cos x,\sin x)$ is a point of $\arc{JiI}$, then

          $TI = \sqrt{(\cos x-1)^{2} + (\sin x-0)^{2}} = \sqrt{2-2 \cos x}$. 

          If $x=3$, then we have our problem. Now, to solve this problem, put your calculator in the radian mode. Normally your display will show ``deg''. Find out how to make it show ``rad'' for radian measure. Once you're in radian mode, do the following:
        \begin{enumerate}[\hspace{1cm}1.]
          \item press $3$
          \item press $\cos$
          \item press $\times$ (times)
          \item press $2$
          \item press $=$
          \item press $\pm$ to change signs (taking care of the $-$ sign in $2-2 \cos x$)
          \item press $+$
          \item press $2$
          \item press $=$
          \item press $\sqrt{\ }$
        \end{enumerate}
        You should have $1.994989973$ displayed. Calculators sure save a lot of time. But what are you going to do if you're stranded on a desert island with no calculator, and you need to know $\cos \frac{23\pi}{24}$? See how useful all this will be to you?
    \end{enumerate}

  \item[\hspace{1cm}e)] This is sort of a trick question. Draw a picture of $K:x^{2}+y^{2}=1$, and try to place $P$ and $Q$ on $K$. You will see that there are no two points on $K$ where the arc length between them is $5$. The greatest the are length can be between two such points is $\pi$.
\end{enumerate}


\textbf{Problem 5.7}:
\begin{enumerate}[\hspace{1cm}1)]
  \item Suppose $P:(\cos u, \sin u)$ and $Q:(\cos v, \sin v)$ are points of $K:x^{2}+y^{2}=1$, and suppose $u < v$.
  \item What is $l(\arc{PQ})$?
  \item Recall Problem 4.4 to solve step 2.
  \item Now that you know $l(\arc{PQ})$, calculate $PQ$. Have we really solved the main problem as it pertains to the unit circle with center at the origin?
\end{enumerate}

\textbf{Problem 5.8}: 

$P$ is a point of $K:x^{2}+y^{2}=1$, which means $P:(\cos u, \sin u)$ for some number $u$ in $[0,2\pi)$. Note how the left end is square, ``$[$'', and the right end is round, ``$)$''. What this means is that the set $[0,2\pi)$ contains $0$, all the numbers between $0$ and $2\pi$, but not $2\pi$. This is called a ``left closed set''. (Aren't you glad you know that?) $P:(\cos u,\sin u)$ implies that $l(\arc{PI})=u$ if $0 \leq u \leq \pi$, or $l(\arc{PI})=2\pi-u$ if $\pi < u < 2\pi$.

Furthermore, $P:(\frac{3}{4}, \frac{\sqrt{7}}{4})$ so that $\cos u = \frac{3}{4}$ and $\sin u = \frac{\sqrt{7}}{4}$. Does this tell you whether $0 \leq u \leq \pi$ or $\pi < u < 2\pi$? Unfortunately, our table doesn't contain these entries, so all we can do is approximate what $u$ is. Can you see how to do this? Can you show that $\frac{5\pi}{24} < u < \frac{6\pi}{24}$? Note that $\cos (\frac{5\pi}{24}) > \cos u > \cos (\frac{6\pi}{24})$. Now, where did that come from?

Note $\frac{5\pi}{24}=\frac{(10\pi/24)}{2}$ or $\frac{(5\pi/12)}{2}$. One of our formulas tells us that $\cos(\frac{x}{2})=\sqrt{\frac{(1+\cos x)}{2}}$ if $0\leq x\leq\pi$. Well, $0\leq\frac{5\pi}{12}\leq\pi$. Hence,

$\cos (\frac{(5\pi/12)}{2}) = \sqrt{\frac{1+\cos5\pi/12}{2}} = \sqrt{\frac{[1+\frac{\sqrt{2}}{4} (\sqrt{3}-1)]}{2}} \doteq 0.79^{+}$.

(Note: $\doteq$ means ``is approximately'').

Then $\cos (\frac{5\pi}{24}) > \cos u$ since $0.79^{+} > 0.75$ . Similarly, $\cos u > \cos (\frac{6\pi}{24})$. Hence, $\cos (\frac{5\pi}{24}) > \cos u > \cos (\frac{6\pi}{24})$. So, $(\frac{5\pi}{24}) < u < (\frac{6\pi}{24})$. Do you see why the inequality is reversed? Now, see if you can show that $(\frac{11\pi}{48}) < u < (\frac{12\pi}{48})$, and then that $(\frac{22\pi}{96}) < u < (\frac{23\pi}{96})$. 

Gets tedious, doesn't it? I'll bet someone who is pretty good at computers could set up a program that would narrow this down so that you could get $u$ to as many decimal places as you wanted. Grinding it out by hand merely shows that you can do it, but it's not very practical if you have a lot of them to do. But note the operative word ``can''. You \emph{can} do it.

\vspace{0.3cm}

\textbf{Theorem 5.2}:
\begin{enumerate}[\hspace{1cm}1)]
  \item Suppose $A:(a,b)$ is a point of $K:x^{2}+y^{2}=1$. Draw a picture.
  \item Suppose $A:(a,b)$ is a point of semi-circle $\arc{JiI}$.
\end{enumerate}

\begin{center}
%TeXCAD Options
%\grade{\on}
%\emlines{\off}
%\epic{\off}
%\beziermacro{\on}
%\reduce{\on}
%\snapping{\off}
%\quality{8.00}
%\graddiff{0.01}
%\snapasp{1}
%\zoom{4.0000}
\unitlength 1mm % = 2.85pt
\linethickness{0.4pt}
\ifx\plotpoint\undefined\newsavebox{\plotpoint}\fi % GNUPLOT compatibility
\begin{picture}(100,41)(0,90)
\put(12,92.75){\line(1,0){87.25}}
\qbezier(58.5,131)(24.13,130.13)(20.25,92.75)
\qbezier(58.75,130.75)(89.25,127.25)(90.75,92.75)
%\emline(23,108)(24.25,107.5)
\multiput(23,108)(.083333,-.033333){15}{\line(1,0){.083333}}
%\end
\put(16,108.5){\makebox(0,0)[cc]{$A:(a,b)$}}
\end{picture}
\end{center}

\begin{enumerate}
  \item[\hspace{1cm}3)] $\theta = l(\arc{AI})$ and $0 \leq \theta \leq \pi$.

      $a = C(\theta) = C_{1}(\theta) = \cos \theta$ and

      $b = S(\theta) = S_{1}(\theta) = \sin\theta$.
  \item[\hspace{1cm}4)] Suppose $A:(a,b)$ is a point of semi-circle $\arc{JiI}$.
\end{enumerate}

\begin{center}
%TeXCAD Options
%\grade{\on}
%\emlines{\off}
%\epic{\off}
%\beziermacro{\on}
%\reduce{\on}
%\snapping{\off}
%\quality{8.00}
%\graddiff{0.01}
%\snapasp{1}
%\zoom{4.0000}
\unitlength 1mm % = 2.85pt
\linethickness{0.4pt}
\ifx\plotpoint\undefined\newsavebox{\plotpoint}\fi % GNUPLOT compatibility
\begin{picture}(72.25,56.47)(0,84)
%\circle(38.25,114.75){51.44}
\put(63.97,114.75){\line(0,1){1.096}}
\put(63.95,115.85){\line(0,1){1.094}}
\put(63.88,116.94){\line(0,1){1.09}}
\multiput(63.76,118.03)(-.03257,.21677){5}{\line(0,1){.21677}}
\multiput(63.6,119.11)(-.02984,.1537){7}{\line(0,1){.1537}}
\multiput(63.39,120.19)(-.03182,.13325){8}{\line(0,1){.13325}}
\multiput(63.13,121.26)(-.0333,.11714){9}{\line(0,1){.11714}}
\multiput(62.83,122.31)(-.03131,.09459){11}{\line(0,1){.09459}}
\multiput(62.49,123.35)(-.03236,.08541){12}{\line(0,1){.08541}}
\multiput(62.1,124.37)(-.033207,.077493){13}{\line(0,1){.077493}}
\multiput(61.67,125.38)(-.031614,.065874){15}{\line(0,1){.065874}}
\multiput(61.2,126.37)(-.032242,.060438){16}{\line(0,1){.060438}}
\multiput(60.68,127.34)(-.032741,.055538){17}{\line(0,1){.055538}}
\multiput(60.12,128.28)(-.033129,.051088){18}{\line(0,1){.051088}}
\multiput(59.53,129.2)(-.033419,.047018){19}{\line(0,1){.047018}}
\multiput(58.89,130.09)(-.033622,.043274){20}{\line(0,1){.043274}}
\multiput(58.22,130.96)(-.032213,.038002){22}{\line(0,1){.038002}}
\multiput(57.51,131.8)(-.032333,.035004){23}{\line(0,1){.035004}}
\multiput(56.77,132.6)(-.033795,.033595){23}{\line(-1,0){.033795}}
\multiput(55.99,133.37)(-.036796,.033585){22}{\line(-1,0){.036796}}
\multiput(55.18,134.11)(-.040012,.03351){21}{\line(-1,0){.040012}}
\multiput(54.34,134.82)(-.043473,.033364){20}{\line(-1,0){.043473}}
\multiput(53.47,135.48)(-.047216,.033138){19}{\line(-1,0){.047216}}
\multiput(52.57,136.11)(-.051284,.032824){18}{\line(-1,0){.051284}}
\multiput(51.65,136.7)(-.055732,.03241){17}{\line(-1,0){.055732}}
\multiput(50.7,137.25)(-.060628,.031882){16}{\line(-1,0){.060628}}
\multiput(49.73,137.77)(-.070779,.033452){14}{\line(-1,0){.070779}}
\multiput(48.74,138.23)(-.077689,.032745){13}{\line(-1,0){.077689}}
\multiput(47.73,138.66)(-.0856,.03186){12}{\line(-1,0){.0856}}
\multiput(46.71,139.04)(-.09478,.03074){11}{\line(-1,0){.09478}}
\multiput(45.66,139.38)(-.11733,.03261){9}{\line(-1,0){.11733}}
\multiput(44.61,139.67)(-.13344,.03102){8}{\line(-1,0){.13344}}
\multiput(43.54,139.92)(-.15388,.02893){7}{\line(-1,0){.15388}}
\multiput(42.46,140.12)(-.21696,.03128){5}{\line(-1,0){.21696}}
\put(41.38,140.28){\line(-1,0){1.09}}
\put(40.29,140.39){\line(-1,0){1.094}}
\put(39.19,140.45){\line(-1,0){1.096}}
\put(38.1,140.47){\line(-1,0){1.096}}
\put(37,140.44){\line(-1,0){1.093}}
\multiput(35.91,140.36)(-.2723,-.0308){4}{\line(-1,0){.2723}}
\multiput(34.82,140.24)(-.18047,-.02822){6}{\line(-1,0){.18047}}
\multiput(33.74,140.07)(-.15352,-.03076){7}{\line(-1,0){.15352}}
\multiput(32.66,139.86)(-.13306,-.03261){8}{\line(-1,0){.13306}}
\multiput(31.6,139.6)(-.10524,-.0306){10}{\line(-1,0){.10524}}
\multiput(30.54,139.29)(-.0944,-.03187){11}{\line(-1,0){.0944}}
\multiput(29.51,138.94)(-.08521,-.03287){12}{\line(-1,0){.08521}}
\multiput(28.48,138.54)(-.077294,-.033667){13}{\line(-1,0){.077294}}
\multiput(27.48,138.11)(-.065684,-.032005){15}{\line(-1,0){.065684}}
\multiput(26.49,137.63)(-.060245,-.032601){16}{\line(-1,0){.060245}}
\multiput(25.53,137.11)(-.055342,-.033071){17}{\line(-1,0){.055342}}
\multiput(24.59,136.54)(-.05089,-.033432){18}{\line(-1,0){.05089}}
\multiput(23.67,135.94)(-.046818,-.033698){19}{\line(-1,0){.046818}}
\multiput(22.78,135.3)(-.041022,-.032265){21}{\line(-1,0){.041022}}
\multiput(21.92,134.62)(-.03781,-.032439){22}{\line(-1,0){.03781}}
\multiput(21.09,133.91)(-.034811,-.032541){23}{\line(-1,0){.034811}}
\multiput(20.29,133.16)(-.033393,-.033995){23}{\line(0,-1){.033995}}
\multiput(19.52,132.38)(-.033365,-.036995){22}{\line(0,-1){.036995}}
\multiput(18.79,131.57)(-.033271,-.04021){21}{\line(0,-1){.04021}}
\multiput(18.09,130.72)(-.033105,-.043671){20}{\line(0,-1){.043671}}
\multiput(17.43,129.85)(-.032857,-.047412){19}{\line(0,-1){.047412}}
\multiput(16.8,128.95)(-.032519,-.051478){18}{\line(0,-1){.051478}}
\multiput(16.22,128.02)(-.032078,-.055924){17}{\line(0,-1){.055924}}
\multiput(15.67,127.07)(-.033622,-.064872){15}{\line(0,-1){.064872}}
\multiput(15.17,126.1)(-.03303,-.070977){14}{\line(0,-1){.070977}}
\multiput(14.7,125.1)(-.032282,-.077883){13}{\line(0,-1){.077883}}
\multiput(14.28,124.09)(-.03135,-.08579){12}{\line(0,-1){.08579}}
\multiput(13.91,123.06)(-.0332,-.10445){10}{\line(0,-1){.10445}}
\multiput(13.58,122.02)(-.03191,-.11752){9}{\line(0,-1){.11752}}
\multiput(13.29,120.96)(-.03023,-.13362){8}{\line(0,-1){.13362}}
\multiput(13.05,119.89)(-.03268,-.17972){6}{\line(0,-1){.17972}}
\multiput(12.85,118.81)(-.02999,-.21714){5}{\line(0,-1){.21714}}
\put(12.7,117.73){\line(0,-1){1.091}}
\put(12.6,116.63){\line(0,-1){1.095}}
\put(12.54,115.54){\line(0,-1){1.096}}
\put(12.53,114.44){\line(0,-1){1.095}}
\put(12.57,113.35){\line(0,-1){1.093}}
\multiput(12.65,112.26)(.0324,-.2721){4}{\line(0,-1){.2721}}
\multiput(12.78,111.17)(.02929,-.1803){6}{\line(0,-1){.1803}}
\multiput(12.96,110.09)(.03167,-.15334){7}{\line(0,-1){.15334}}
\multiput(13.18,109.01)(.0334,-.13287){8}{\line(0,-1){.13287}}
\multiput(13.44,107.95)(.03122,-.10506){10}{\line(0,-1){.10506}}
\multiput(13.76,106.9)(.03243,-.09421){11}{\line(0,-1){.09421}}
\multiput(14.11,105.86)(.03338,-.08502){12}{\line(0,-1){.08502}}
\multiput(14.51,104.84)(.031688,-.071586){14}{\line(0,-1){.071586}}
\multiput(14.96,103.84)(.032395,-.065493){15}{\line(0,-1){.065493}}
\multiput(15.44,102.86)(.032959,-.06005){16}{\line(0,-1){.06005}}
\multiput(15.97,101.9)(.0334,-.055145){17}{\line(0,-1){.055145}}
\multiput(16.54,100.96)(.033734,-.05069){18}{\line(0,-1){.05069}}
\multiput(17.15,100.05)(.032277,-.044286){20}{\line(0,-1){.044286}}
\multiput(17.79,99.16)(.032509,-.04083){21}{\line(0,-1){.04083}}
\multiput(18.47,98.3)(.032663,-.037616){22}{\line(0,-1){.037616}}
\multiput(19.19,97.48)(.032748,-.034617){23}{\line(0,-1){.034617}}
\multiput(19.95,96.68)(.034193,-.033191){23}{\line(1,0){.034193}}
\multiput(20.73,95.92)(.037193,-.033145){22}{\line(1,0){.037193}}
\multiput(21.55,95.19)(.040408,-.033032){21}{\line(1,0){.040408}}
\multiput(22.4,94.49)(.043867,-.032844){20}{\line(1,0){.043867}}
\multiput(23.28,93.84)(.047607,-.032574){19}{\line(1,0){.047607}}
\multiput(24.18,93.22)(.051671,-.032212){18}{\line(1,0){.051671}}
\multiput(25.11,92.64)(.059621,-.033729){16}{\line(1,0){.059621}}
\multiput(26.06,92.1)(.06507,-.033236){15}{\line(1,0){.06507}}
\multiput(27.04,91.6)(.071172,-.032607){14}{\line(1,0){.071172}}
\multiput(28.04,91.14)(.078073,-.031819){13}{\line(1,0){.078073}}
\multiput(29.05,90.73)(.09379,-.03364){11}{\line(1,0){.09379}}
\multiput(30.08,90.36)(.10465,-.03257){10}{\line(1,0){.10465}}
\multiput(31.13,90.03)(.11771,-.03121){9}{\line(1,0){.11771}}
\multiput(32.19,89.75)(.15292,-.03364){7}{\line(1,0){.15292}}
\multiput(33.26,89.52)(.17991,-.03161){6}{\line(1,0){.17991}}
\multiput(34.34,89.33)(.21731,-.0287){5}{\line(1,0){.21731}}
\put(35.43,89.18){\line(1,0){1.092}}
\put(36.52,89.09){\line(1,0){1.095}}
\put(37.61,89.04){\line(1,0){1.096}}
\put(38.71,89.03){\line(1,0){1.095}}
\put(39.8,89.08){\line(1,0){1.092}}
\multiput(40.9,89.17)(.21751,.02719){5}{\line(1,0){.21751}}
\multiput(41.98,89.3)(.18013,.03036){6}{\line(1,0){.18013}}
\multiput(43.06,89.48)(.15314,.03258){7}{\line(1,0){.15314}}
\multiput(44.14,89.71)(.11793,.03039){9}{\line(1,0){.11793}}
\multiput(45.2,89.99)(.10487,.03185){10}{\line(1,0){.10487}}
\multiput(46.25,90.3)(.09402,.03299){11}{\line(1,0){.09402}}
\multiput(47.28,90.67)(.078292,.031277){13}{\line(1,0){.078292}}
\multiput(48.3,91.07)(.071396,.032114){14}{\line(1,0){.071396}}
\multiput(49.3,91.52)(.065299,.032784){15}{\line(1,0){.065299}}
\multiput(50.28,92.01)(.059853,.033315){16}{\line(1,0){.059853}}
\multiput(51.24,92.55)(.054945,.033727){17}{\line(1,0){.054945}}
\multiput(52.17,93.12)(.047831,.032244){19}{\line(1,0){.047831}}
\multiput(53.08,93.73)(.044093,.03254){20}{\line(1,0){.044093}}
\multiput(53.96,94.38)(.040635,.032751){21}{\line(1,0){.040635}}
\multiput(54.81,95.07)(.037421,.032886){22}{\line(1,0){.037421}}
\multiput(55.64,95.8)(.034422,.032953){23}{\line(1,0){.034422}}
\multiput(56.43,96.55)(.032987,.034389){23}{\line(0,1){.034389}}
\multiput(57.19,97.34)(.032923,.037389){22}{\line(0,1){.037389}}
\multiput(57.91,98.17)(.032791,.040603){21}{\line(0,1){.040603}}
\multiput(58.6,99.02)(.032583,.044062){20}{\line(0,1){.044062}}
\multiput(59.25,99.9)(.032291,.0478){19}{\line(0,1){.0478}}
\multiput(59.87,100.81)(.031904,.051861){18}{\line(0,1){.051861}}
\multiput(60.44,101.74)(.033374,.05982){16}{\line(0,1){.05982}}
\multiput(60.97,102.7)(.032848,.065267){15}{\line(0,1){.065267}}
\multiput(61.47,103.68)(.032184,.071365){14}{\line(0,1){.071365}}
\multiput(61.92,104.68)(.031354,.078261){13}{\line(0,1){.078261}}
\multiput(62.32,105.7)(.03308,.09398){11}{\line(0,1){.09398}}
\multiput(62.69,106.73)(.03195,.10484){10}{\line(0,1){.10484}}
\multiput(63.01,107.78)(.03051,.1179){9}{\line(0,1){.1179}}
\multiput(63.28,108.84)(.03273,.15311){7}{\line(0,1){.15311}}
\multiput(63.51,109.91)(.03054,.1801){6}{\line(0,1){.1801}}
\multiput(63.69,110.99)(.02741,.21748){5}{\line(0,1){.21748}}
\put(63.83,112.08){\line(0,1){1.092}}
\put(63.92,113.17){\line(0,1){1.579}}
%\end
\put(9.25,114.75){\line(1,0){63}}
%\emline(56,97.25)(56.75,96.25)
\multiput(56,97.25)(.032609,-.043478){23}{\line(0,-1){.043478}}
%\end
\put(37.75,89.75){\line(0,-1){1.25}}
\put(8.75,112.5){\makebox(0,0)[cc]{$J$}}
\put(66,112.75){\makebox(0,0)[cc]{$I$}}
\put(64.25,95.25){\makebox(0,0)[cc]{$A:(a,b)$}}
\put(37.5,86){\makebox(0,0)[cc]{$j$}}
\end{picture}
\end{center}

\begin{enumerate}
  \item[\hspace{1cm}5)] Let $\theta = 2\pi - l(\arc{AI}) = \pi + l(\arc{JA})$.
  \item[\hspace{1cm}6)] For $\pi < \theta < 2\pi$: 
     \\ $a = -C(\theta-\pi) = C_{1}(\theta) = \cos \theta$
     \\ $b = -S(\theta-\pi) = S_{1}(\theta) = \sin \theta$
     \\ From step 3 and step 5, regardless of where $A:(a,b)$ is on $K:x^{2}+y^{2}=1$, $a = \cos \theta$ and $b = \sin \theta$, for some $0 \leq \theta < 2\pi$.
\end{enumerate}

\textbf{Theorem 5.3}:
\begin{enumerate}[\hspace{1cm}1)]
  \item Suppose $P:(u,v)$, $OP=r>0$, and $A:(a,b)$ is a point of $K:x^{2}+y^{2}=1$ such that $OPA$ or $OAP$. Draw a picture or two.
  \item Either $OPA$ or $OAP$. Hence, either $OP+PA=OA$ or $OA+AP=OP$.
  \item Either $r+PA=1$ or $1+AP=r$. Hence $PA=1-r$ or $AP=r-1$.
  \item Hence $(a-u)^{2} + (b-v)^{2} = (1-r)^{2}$.
  \item Reducing this, noting that $u^{2}+v^{2}=r^{2}$ and $a^{2}+b^{2}=1$, we have $au+bv=r$.
  \item Now note that $P$ and $A$ are colinear with $0$.
  \item Note that $\frac{b}{a} = \frac{v}{u}$. Hence, $av-bu=0$.
  \item Combine steps 5 and 7.
  \item Solving $au+bv=r$ and $-bu+av=0$, you should get $u=ar$ and $v=br$.
\end{enumerate}

\textbf{Corollary 5.3}: Combining Theorems 5.2 and 5.3 yields $u = r \cos \theta$ and $v = r \sin \theta$ for some $\theta$ such that $0 \leq \theta < 2\pi$.


\vspace{2cm}

\begin{center}
Advertisement: Violin for sale by young man in good condition except for a loose peg in the head.
\vspace{0.5cm}

Justice: Juvenile court to try shooting defendant.
\end{center}


\chapter*{Chapter VI}
\addcontentsline{toc}{chapter}{\numberline{}Chapter VI}
\markboth{Chapter VI}{Chapter VI}

The problem of finding $\theta$ when we are given $\sin \theta$ or $\cos \theta$ keeps cropping up. Of course, that's because I keep bringing it up. However, this is not an uncommon problem. It turns out that we are faced with this idea when dealing with a variety of functions, not just our newly invented ones that we call cosine and sine. Hence, the following definition:

\vspace{0.3cm}

\textbf{Definition 6.1}: If $f$ is a function, then $f^{-1}$ denotes the collection of ordered pairs to which $(x,y)$ belongs if and only if $(y,x)$ belongs to $f$. $f^{-1}$ is called the \emph{inverse} of $f$.

\vspace{0.3cm}

\textbf{Problem 6.1}:
\begin{enumerate}[\hspace{1cm}1)]
  \item Find a function $f$ such that $f^{-1}$ is \emph{not} a function.
  \item Find a function $f$ such that $f^{-1}$ \emph{is} a function.
  \item Suppose $f^{-1}$ is the inverse of the function $f$, and $(1,3)$ is the only number pair belonging to $f$. By definition, $(3,1)$ belongs to $f^{-1}$.
    \begin{enumerate}[\hspace{1cm}a)]
      \item What is $f(1)$?
      \item What is $f^{-1}(3)$?
      \item Does $f[f^{-1}(3)]$ make sense, and if so, what is it?
      \item Does $f^{-1}[f(1)]$ make sense, and if so, what is it?
    \end{enumerate}
  \item The graph of a function is a simple graph; i.e., no vertical line contains two points of it. Can you find a similar statement to make about the graph of a function whose inverse is also a function?
  \item Is $\cos^{-1}$ a function?
  \item How many solutions are there for $x$ in the following equations?
    \begin{enumerate}[\hspace{1cm}a)]
      \item $\cos x = 1$
      \item $\cos x = 1.5$
      \item $\cos x = -\frac{1}{2}$
    \end{enumerate}
  \item How many solutions are there of the equations in 6) if we restrict $x$ to the interval $[0,2\pi]$?
  \item How many solutions are there of the equations in 6) if we restrict $x$ to the interval $[0,\pi]$?
  \item Note that no horizontal line contains two points of the graph of the cosine function in the interval $[0,11]$. Can you find an interval over which this will be true for the sine function? How many such intervals are there?
\end{enumerate}

\textbf{Definition 6.2}: Solutions to $y = \sin x$ and $y = \cos x$ for $x$ in the intervals $[-\frac{\pi}{2},\frac{\pi}{2}]$ and $[0,\pi]$ respectively are called \emph{principal values} of $\sin^{-1}$ and $\cos^{-1}$.

\vspace{0.3cm}

Consider the problem of finding the, or a, solution for $x$ in the equation $y = \sin x$. A specific example will help. Consider $\frac{1}{2} = \sin x$, and we are asked to solve for $x$. You will note that $\frac{\pi}{6}$ is a solution, but so is $\frac{5\pi}{6}$.

Furthermore, if $n$ is an integer, then $\frac{\pi}{6}+n(2\pi)$ and $\frac{5\pi}{6}+n(2\pi)$ are also solutions. Hence, there are infinitely many solutions to the equation $\frac{1}{2} = \sin x$ for the number $x$. We write these solutions as $x = \sin^{-1}\frac{1}{2}$. \emph{The} solution in the interval $[-\frac{\pi}{2},\frac{\pi}{2}]$ is called the principal value of $\sin^{-1} \frac{1}{2}$, and you'll note that the solution is $\frac{\pi}{6}$. You'll also note that the graph of the sine function in the interval $[-\frac{\pi}{2},\frac{\pi}{2}]$ is such that no horizontal line contains two points of it. Can you see the connection between this geometric property of the sine function on $[-\frac{\pi}{2},\frac{\pi}{2}]$ and the idea that $\sin^{-1} \frac{1}{2}$ is unique; i.e., there is only one number $x$ in $[-\frac{\pi}{2},\frac{\pi}{2}]$ such that $x = \sin^{-1} \frac{1}{2}$, and we know that number is $\frac{\pi}{6}$; i.e., $\sin^{-1} \frac{1}{2} = \frac{\pi}{6}$.

Will the inverse of the cosine function be a function on the same interval, $[-\frac{\pi}{2},\frac{\pi}{2}]$? A quick look at the graph of the cosine function will convince you that its inverse is \emph{not} a function on $[-\frac{\pi}{2},\frac{\pi}{2}]$  but that it \emph{is} a function on $[0,\pi]$. Hence, we say that the principal values of $\cos^{-1}$ are in $[0,\pi]$. You'll note that the inverse of the cosine function is also a function on $[2\pi,3\pi]$, but the interval containing $0$ is reserved for the name principal values.

\vspace{0.3cm}

\textbf{Problem 6.2}: Suppose that $f$ is a function, and on the interval $[a,b]$, $f^{-1}$ is also a function. Then, on the interval $[a,b]$, no vertical line contains two points of the graph of $f$ (since $f$ is a function), and no horizontal line contains two points of $f$ (since $f^{-1}$ is a function).
\begin{enumerate}[\hspace{1cm}1)]
  \item Is it true that on $[a,b]$, $f$ is either increasing at each point or decreasing at each point?
  \item What would you mean by saying that a simple graph \emph{is} increasing at a point?
  \item What would you mean by saying that a simple graph is \emph{not} increasing at a point?
\end{enumerate}

Still in reference to $f$ above, if $(x,y)$ belongs to $f$, we write $(x,y) = (x,f(x))$, or $y = f(x)$. Also, if $(x,y)$ belongs to $i$, then $(y,x)$ belongs to $f^{-1}$. We write $(y,x) = (y,f^{-1}(y))$, or $x = f^{-1}(y)$. 

Solutions for $x$ in $[a,b]$ of the equation $y = f(x)$ are written $x = f^{-1}(y)$. Only for special functions is reference usually made to principal values.

\vspace{0.3cm}

I have stated that if we start with the equation $y = \sin x$ (for example, $\frac{3}{4} = \sin x$) and we want to solve for $x$, we write $x = \sin^{-1} y$. Of course, this doesn't tell you how to calculate $x$, but the problem at this time is to decide how to read ``$x = \sin^{-1} y$.'' The ``$=$'' is read ``is'', but how do you read $\sin^{-1} y$? We could just read it ``the inverse sine of $y$''. But we can do a little better than that to give some insight into these inverse functions. We have $x = \sin^{-1} y$. What is $x$ anyway? We started with $y = \sin x$. Hence, $x$ is a number whose sine is $y$. Let's read ``$x = \sin^{-1} y$'' as ``$x$ is a number whose sine is $y$.'' If we restrict ourselves to principal values -- i.e., $x$ in $[-\frac{\pi}{2},\frac{\pi}{2}]$ -- then we can read it ``$x$ is \emph{the} number in $[-\frac{\pi}{2},\frac{\pi}{2}]$ whose sine is $y$''.

\vspace{0.3cm}

We have $y = \sin x$ and $x = \sin^{-1}y$, so that $\sin x = \sin(\sin^{-1}y) = y$. Is it true that $\sin^{-1}(\sin x) = x$? Do we have to restrict ourselves to principal values in order for these manipulative shenanigans to make sense?

Our development of the sine and cosine functions could lead us to interpret $x$ in $y = \sin x$ as an arc length in some cases. It used to be that $x = \sin^{-1} y$ was written $x =$ arcsin $y$, and read ``$x$ is arcsine $y$'' or ``$x$ is an arc whose sine is $y$''. This symbolism is not used much in today's literature. Almost always, the symbol $\sin^{-1}$ is used. Perhaps one reason you don't see ``arcsin $y$'' much anymore is that it is interpreted as ``an arc whose sine is $y$'' when really it is not the arc at all, but the length of an arc whose sine is $y$. So, if you write $x =$ arcsin $y$, $x$ is (the length of) an arc whose sine is $y$, then we automatically think of $x$ as a non-negative number. We have difficulty dealing with the notion of negative length. But as we have seen, the principal values of $\sin^{-1}$ or arcsin are in $[-\frac{\pi}{2},\frac{\pi}{2}]$. So $x$ in $x =$ arcsin $y$ can be negative. This is awkward. 

\vspace{0.3cm}

For example, consider $-\frac{\pi}{4} =$ arcsin $(-\frac{\sqrt{2}}{2})$. We read this as ``$\frac{\pi}{4}$ is the length of an arc whose sine is $-\sqrt{2}$''. This sounds silly to us. Hence, we read this as ``$-\frac{\pi}{4}$ is a number whose sine is $-\sqrt{2}$''. In general, we read $x = \sin^{-1} y$ as ``$x$ is a number whose sine is $y$'', and if we are restricting the discussion to principal values, we read it as ``$x$ is the number in $[-\frac{\pi}{2},\frac{\pi}{2}]$ whose sine is $y$''.

Suppose we are asked ``What is $\sin(\sin^{-1}x)$?'' We read $\sin^{-1}x$ as ``a number whose sine is $x$''. Hence the question, ``What is $\sin(\sin^{-1}x)$?'' is ``What is the sine of a number whose sine is $x$?'' Isn't this like asking ``What is the name of a person whose name is George?'' Is it obvious that the name of a person whose name is George is George? Hence, the sine of a number whose sine is $x$ is $x$; i.e., $\sin(\sin^{-1}x) = x$. Nothing but $x$ will satisfy this.

Now what is $\sin^{-1}(\sin x)$? In general, we read $\sin^{-1} y$ as ``a number whose sine is $y$''. Hence, $\sin^{-1}(\sin x)$ is read as ``a number whose sine is $\sin x$''. Well, $x$ is certainly one such number, but can you see that there are a lot of numbers that qualify? Let's consider an example. 

\vspace{0.3cm}

Suppose $\sin x = \frac{1}{2}$. We ask, what is $\sin^{-1}(\sin x)$? Well, $\sin^{-1}(\sin x) = \sin^{-1}(\frac{1}{2})$, which is ``a number whose sine is $\frac{1}{2}$''. $x$ is one such number, since we are given that $\sin x = \frac{1}{2}$. But can you see that $x+2\pi$ is also such a number? There's no unique answer to ``what is $\sin(\sin x)$?''

What if we restrict ourselves to principal values of $\sin^{-1}$? Then we read 

$\sin^{-1} (\sin x)$ as ``$\pi$ number in $[-\frac{\pi}{2},\frac{\pi}{2}]$ whose sine is $\sin x$''. There is only one such number, but is it $x$? What if $x$ is not in $[-\frac{\pi}{2},\frac{\pi}{2}]$? Alas, the conjecture $\sin^{-1}(\sin x) = x$ doesn't seem to work even though $\sin(\sin^{-1}x) = x$ is valid. Doesn't this seem strange? We seem to be caught in the WorldWide Weird Web. Do you think it is true that, given a function $f$, then $f[f^{-1}(x)] = x$ in all cases, but it is not always true that $f^{-1}[f(x)] = x$? This is a nice little puzzle for you that could be very instructive.

\vspace{0.3cm}

When speaking of the graph of a function, it is clear what is meant when we say ``Consider the graph of the function $f$ on the interval $[a,b]$.'' This implies that $f$ is defined on $[a,b]$ and the graph is a simple graph. When speaking of the graph of the inverse of a function, we have to be a little careful how we phrase things so that what we say makes sense.

There are two cases to consider. The first case is the one in which the inverse of the function under consideration is a function itself. Suppose $f$ is the function and $f$ is defined on the interval $[a,b]$. This mean that the initial set of $f$ is $[a,b]$. Suppose, furthermore, that the final set of $f$ is the interval $[c,d]$. Then, since $f^{-1}$ is a function in this case, the initial set of $f^{-1}$ is $[c,d]$, and the final set of $f^{-1}$ is $[a,b]$. We graph $f^{-1}$ on $[c,d]$, its initial set. There doesn't appear to be any confusion here. In fact, even if we say the initial set off is the set $M$ (perhaps an interval or the whole $x$-axis or whatever) and the final set of $f$ is the set $K$, then the initial set of $f^{-1}$ is $K$ and the final set of $f^{-1}$ is $M$, and we say that we graph $f^{-1}$ on the set $K$, the final set of $f$.

But what if the inverse of $f$, $f^{-1}$, is \emph{not} a function? Then its graph is \emph{not} a simple graph. Does this present a difficulty when we say ``graph $f^{-1}$''? Let's consider an example. 

Suppose $f(x) = \cos x$. In this case, the initial set of $f$ is the set of all real numbers, and the final set of $f$ is the interval $[-1,1]$. Hence, the initial set of $f^{-1}$ (namely $\cos^{-1}$) is the interval $[-1,1]$, and the final set of $f^{-1}$ is the set of all real numbers. But the graph of $f = \cos^{-1}$ is not a simple graph. The graph is defined on $[-1,1]$, and every vertical line intersecting $[-1,1]$ intersects the graph of $f^{-1} = \cos^{-1}$ infinitely many times. Can you see how this comes about? How many times does a horizontal line that intersects the $y$-axis in $[-1,1]$ intersect the graph of the cosine function? You're right! Infinitely many times. There's a connection between the horizontal line intersecting the graph of the cosine function infinitely many times, and vertical line intersecting the graph of $\cos^{-1}$ infinitely many times. 

Now, what if we wanted to look at that portion of the graph of $\cos^{-1}$ that is derived from that portion of the graph of the cosine function defined on $[0,\pi]$? How shall we describe this without having to write a long discourse about it? Let's do this by saying ``graph $\cos^{-1}$ over $[0,\pi]$''. Since $[0,\pi]$ is what we have called the principal values of $\cos^{-1}$, we can say ``graph $\cos^{-1}$ over its principal values''. This will distinguish that portion of the graph of $\cos^{-1}$ such that if $P:(x,y)$ is a point of this portion of the graph, then $-1<x<1$ and $0<y<\pi$. The value of $\cos^{-1}$ is $\sin[0,\pi]$.

Similarly, we can say ``graph $\sin^{-1}$ over its principal values'' and mean that if $P:(x,y)$ is a point of this portion of the graph of $\sin^{-1}$, then $-1<x<1$ and $-\frac{\pi}{2} < y < \frac{\pi}{2}$. Make sure you see what I'm talking about here. It's not difficult, but it's important to know what is being referred to by the terminology.

\vspace{0.3cm}

\textbf{Problem 6.3}:
\begin{enumerate}[\hspace{1cm}a)]
  \item Sketch the graphs of $\cos^{-1}$ and $\sin^{-1}$ over their respective principal values.
  \item It will be instructive for you to sketch the graphs of $\cos{-1}$ and $\sin{-1}$ over the intervals $[2\pi,3\pi]$ and $[2\pi-\frac{\pi}{2}, 2\pi+\frac{\pi}{2}]$, respectively.
  \item If you could sketch the graphs of $\cos^{-1}$ in its entirety, how many points of this graph would be contained in the vertical line $x=0$?
\end{enumerate}

\textbf{Problem 6.4}: Suppose $P$ is a point of $K:x^2+y^2=1$. Find a relationship between $PI$ and $l(\arc{PI})$.

\vspace{0.3cm}

\textbf{Problem 6.5}: Suppose $P$ and $Q$ are points of $K:x^2+y^2=1$. Find a relationship between $PQ$ and $l(\arc{PQ})$.

\vspace{0.3cm}

\textbf{Problem 6.6}: 
\begin{enumerate}[\hspace{1cm}a)]
  \item Sketch the graph of the tangent function on $[-2\pi,2\pi]$.
  \item Can you find a sub-interval of $[-2\pi,2\pi]$ that would be appropriate to call the principal values of $\tan^{-1}$?
  \item Sketch the graph of $\tan^{-1}$ over $[-2\pi,2\pi]$.
  \item You'll need to know what the initial set of $\tan^{-1}$ is.
  \item Also, don't lose sight of the fact that the tangent function is not defined at $\frac{\pi}{2}$. Where else in $[-2\pi,2\pi]$ is the tangent function not defined?
\end{enumerate}



\newpage
\begin{center}
\section*{Chapter VI Hints and Solutions}
\end{center}


\textbf{Problem 6.1}:
\begin{enumerate}[\hspace{1cm}1)]
  \item Consider $y = x^{2}$ on $[-1,1]$.
  \item Consider $y = x^{2}$ on $[0,1]$.
  \item 
    \begin{enumerate}[\hspace{1cm}a)]
      \item ``$(1,3)$ belongs to $f$'' implies that $(1,3) = (1,f(1))$. Hence, $f(1)= 3$.
      \item ``$(3,1)$ belongs to $f^{-1}$'' implies that $(3,1) = (3,f^{-1}(3))$. Hence, $f^{-1}(3) = 1$.
      \item $f[f^{-1}(3)] = f(1)$, since $f^{-1}(3)=1$; and since $f(1) = 3$, it follows that $f[f^{-1}(3)] = 3$.
      \item $f^{-1}[f(1)] = f^{-1}(3)$ since $f(1) = 3$, and since $f^{-1}(3) = 1$, it follows that $f^{-1}[f(1)] = 1$.
    \end{enumerate}
  \item If $f$ is a function and $f^{-1}$ is a function, then no horizontal line contains two points of the graph of $f$.
  \item Since it is not true that no horizontal line contains two points of the graph of the cosine function, $\cos^{-1}$ is \emph{not} a function.
  \item
    \begin{enumerate}[\hspace{1cm}a)]
      \item infinitely many
      \item none
      \item infinitely many
    \end{enumerate}
  \item
    \begin{enumerate}[\hspace{1cm}a)]
      \item two
      \item none
      \item two
    \end{enumerate}
  \item
    \begin{enumerate}[\hspace{1cm}a)]
      \item one
      \item none
      \item one
    \end{enumerate}
  \item $[-\frac{\pi}{2}, \frac{\pi}{2}]$ for one. For each integer $n$, $[n(2\pi)-\frac{\pi}{2}, n(2\pi)+\frac{\pi}{2}]$. Hence, there are infinitely many.
\end{enumerate}

\textbf{Problem 6.2}:
\begin{enumerate}[\hspace{1cm}1)]
  \item yes
  \item The statement that the simple graph $g$ is increasing at the point $P$ means:
    \begin{enumerate}[\hspace{1cm}a)]
      \item $P$ is a point of $g$.
      \item If $h$ and $k$ are vertical lines with $P$ between them, then there exists a point of $g$ between $h$ and $k$ and to the right of $P$.
      \item There exist vertical lines $h$ and $k$ with $P$ between them such that every point of $g$ between $h$ and $k$ and to the right of $P$ is above $P$.
    \end{enumerate}
    Draw pictures of this and ponder it carefully. If you can see how this works for you, then you are ready for some really cool ideas in your next course.
  \item If you worked on this, you probably appreciated the definition just given. If you construct the bare denial to the definition, you will have the idea.

        The bare denial goes like this: The statement that the simple graph $g$ is \emph{not} increasing at a point means: either 
    \begin{enumerate}[\hspace{1cm}a)]
      \item $P$ is not a point of $g$, or
      \item there exist vertical lines $h$ and $k$ with $P$ between them such that there is no point of $g$ between $h$ and $k$ and to the right of $P$, or 
      \item if $h$ and $k$ are vertical lines with $P$ between them, then there is some point of $g$ between $h$ and $k$ and to the right of $P$ which is not above $P$.
    \end{enumerate}
        Note that for a graph to be increasing at $P$, it must satisfy \emph{all} three conditions in the definition -- (a), (b), and (c). Hence, if it fails any \emph{one} of these three, it is not increasing at $P$.
\end{enumerate}

\textbf{Problem 6.3}:
\begin{enumerate}[\hspace{1cm}1)]
  \item sketch the cosine function on $[0,\pi]$.
  \item Note that $\cos^{-1}$ has initial set $[-1,1]$ and final set $[0,\pi]$ over its principal values.
  \item Set up the following:
\end{enumerate}

\begin{center}
%TeXCAD Options
%\grade{\on}
%\emlines{\off}
%\epic{\off}
%\beziermacro{\on}
%\reduce{\on}
%\snapping{\off}
%\quality{8.00}
%\graddiff{0.01}
%\snapasp{1}
%\zoom{4.0000}
\unitlength 1mm % = 2.85pt
\linethickness{0.4pt}
\ifx\plotpoint\undefined\newsavebox{\plotpoint}\fi % GNUPLOT compatibility
\begin{picture}(78,41)(0,94)
\put(45.25,100.75){\line(0,1){34.25}}
\put(44.5,134.75){\line(1,0){1.5}}
\put(9.75,100.25){\line(1,0){68.25}} 
\put(9.5,101){\line(0,-1){1}}
\put(25.5,101.25){\line(0,-1){1.75}}
\put(77.5,101.5){\line(0,-1){1.5}} 
\put(62,101.5){\line(0,-1){1.25}}
\put(44.5,119){\line(1,0){1.5}}
\put(44.5,108.75){\line(1,0){1.5}} 
\put(44.5,127){\line(1,0){1.5}}
\put(61.5,96){\makebox(0,0)[cc]{$1/2$}}
\put(77.75,96){\makebox(0,0)[cc]{$1$}}
\put(25.5,96){\makebox(0,0)[cc]{$-1/2$}}
\put(9.75,96){\makebox(0,0)[cc]{$-1$}}
\put(45.75,96){\makebox(0,0)[cc]{$0$}}
\put(51,134.75){\makebox(0,0)[cc]{$\pi$}}
\put(51,126.75){\makebox(0,0)[cc]{$3\pi/4$}}
\put(51,118.75){\makebox(0,0)[cc]{$\pi/2$}}
\put(51,108.75){\makebox(0,0)[cc]{$\pi/4$}}
\end{picture}
\end{center}

\begin{enumerate}
  \item[\hspace{1cm}4)] Find $\cos^{-1}(-1)$, $\cos^{-1}(0)$, and $\cos^{-1}(1)$, and sketch these in the figure above.
  \item[\hspace{1cm}5)] Find some more values of cos$^{-1}$ to fill in enough points to sketch the graph reasonably.
  \item[\hspace{1cm}6)] Your final result should look something like the following: (a)
\end{enumerate}

\begin{center}
%TeXCAD Options
%\grade{\on}
%\emlines{\off}
%\epic{\off}
%\beziermacro{\on}
%\reduce{\on}
%\snapping{\off}
%\quality{8.00}
%\graddiff{0.01}
%\snapasp{1}
%\zoom{4.0000}
\unitlength 1mm % = 2.85pt
\linethickness{0.4pt}
\ifx\plotpoint\undefined\newsavebox{\plotpoint}\fi % GNUPLOT compatibility
\begin{picture}(80,40)(0,95)
\put(45.25,100.75){\line(0,1){34.25}}
\put(44.5,134.75){\line(1,0){1.5}}
\put(45.75,109.5){\line(-1,0){1.25}}
\put(44.75,126.25){\line(1,0){1.25}} 
\put(44.75,129){\line(1,0){1}}
\put(44.75,124){\line(1,0){1}} 
\put(44.75,106.75){\line(1,0){1.25}}
\put(44.5,112){\line(1,0){1.5}} 
\put(9.75,100.25){\line(1,0){68.25}}
\put(9.5,101){\line(0,-1){1}} 
\put(25.5,101.25){\line(0,-1){1.75}}
\put(20.5,101.25){\line(0,-1){1.5}}
\put(14.25,101.25){\line(0,-1){1.25}}
\put(77.5,101.5){\line(0,-1){1.5}} 
\put(62,101.5){\line(0,-1){1.25}}
\put(66.75,101.5){\line(0,-1){1}}
\put(71.75,101.5){\line(0,-1){1.5}}
\put(41.5,135){\makebox(0,0)[cc]{$\pi$}}
\put(61.5,97){\makebox(0,0)[cc]{$\frac{1}{2}$}}
\put(66.75,103){\makebox(0,0)[cc]{$\frac{\sqrt{2}}{2}$}}
\put(71,97){\makebox(0,0)[cc]{$\frac{\sqrt{3}}{3}$}}
\put(77.75,97){\makebox(0,0)[cc]{$1$}}
\put(25,103){\makebox(0,0)[cc]{$-\frac{1}{2}$}}
\put(19.5,97){\makebox(0,0)[cc]{$-\frac{\sqrt{2}}{2}$}}
\put(14.25,103){\makebox(0,0)[cc]{$-\frac{\sqrt{3}}{2}$}}
\put(9.75,97){\makebox(0,0)[cc]{$-1$}}
\put(45.75,97){\makebox(0,0)[cc]{$0$}}
\put(40,129.25){\makebox(0,0)[cc]{$5\pi/6$}}
\put(50.5,126.25){\makebox(0,0)[cc]{$3\pi/4$}}
\put(39.25,123.5){\makebox(0,0)[cc]{$2\pi/3$}}
\put(50.5,112){\makebox(0,0)[cc]{$\pi/3$}}
\put(41,109.25){\makebox(0,0)[cc]{$\pi/4$}}
\put(50.25,106.5){\makebox(0,0)[cc]{$\pi/6$}}
\put(44.5,119){\line(1,0){1.5}}
\put(50.25,120.25){\makebox(0,0)[cc]{$\pi/2$}}
\put(71.25,108){\circle{.71}} 
\put(66.5,110.75){\circle{.5}}
\put(60.75,113.5){\circle{.5}} 
\put(21.5,124.75){\circle{.71}}
\put(16,127.25){\circle{.5}} 
\put(10.25,134.25){\circle{.5}}
\put(13,130.25){\circle{.71}}
\qbezier(10.25,134.5)(10.5,132.13)(12.75,130.25)
\qbezier(13,130.25)(14.13,128.63)(15.75,127.5)
\qbezier(16,127.25)(17.88,125.5)(21.25,124.75)
\qbezier(21.25,125)(34.13,119.5)(45.5,119)
\qbezier(45.25,119)(47.13,120.75)(60.5,113.5)
\qbezier(60.5,113.5)(63.75,112)(66,110.5)
\qbezier(66.25,111)(69.38,108.75)(71,108.5)
\qbezier(71.25,108.5)(80,103)(77.75,100.5)
%\qbezvec(65,124)(58.63,124.13)(57.75,116.75)
\put(57.75,116.75){\vector(-1,-4){.07}}\qbezier(65,124)(58.63,124.13)(57.75,116.75)
%\end
\put(70,124.25){\makebox(0,0)[cc]{$\cos^{-1}$}}
\end{picture}
\end{center}

\begin{enumerate}
  \item[\hspace{1cm}7)] Repeat these procedures for the sine function to get a sketch of the graph of $\sin^{-1}$ over the interval $[-\frac{\pi}{2},\frac{\pi}{2}]$.
  \item[\hspace{1cm}8)] Your final result should look something like the following:
\end{enumerate}

\begin{center}
%TeXCAD Options
%\grade{\on}
%\emlines{\off}
%\epic{\off}
%\beziermacro{\on}
%\reduce{\on}
%\snapping{\off}
%\quality{8.00}
%\graddiff{0.01}
%\snapasp{1}
%\zoom{4.0000}
\unitlength 1mm % = 2.85pt
\linethickness{0.4pt}
\ifx\plotpoint\undefined\newsavebox{\plotpoint}\fi % GNUPLOT compatibility
\begin{picture}(82.5,37.25)(0,98)
\put(45.25,100.75){\line(0,1){34.25}}
\put(44.5,134.75){\line(1,0){1.5}} 
\put(46,117.75){\line(-1,0){1.5}}
\put(45.75,109.5){\line(-1,0){1.25}}
\put(44.75,126.25){\line(1,0){1.25}} 
\put(44.75,129){\line(1,0){1}}
\put(44.75,124){\line(1,0){1}} 
\put(44.75,106.75){\line(1,0){1.25}}
\put(44.5,112){\line(1,0){1.5}} 
\put(9.75,117.75){\line(1,0){68.25}}
\put(9.5,118.5){\line(0,-1){1}} 
\put(25.5,118.75){\line(0,-1){1.75}}
\put(20.5,118.75){\line(0,-1){1.5}}
\put(14.25,118.75){\line(0,-1){1.25}}
\put(77.5,119){\line(0,-1){1.5}} 
\put(62,119){\line(0,-1){1.25}}
\put(66.75,119){\line(0,-1){1}} 
\put(71.75,119){\line(0,-1){1.5}}
\put(61,124.5){\circle{.71}} 
\put(65.75,127){\circle{.71}}
\put(70,129.25){\circle{.5}} 
\put(77.5,135){\circle{.5}}
\put(25.75,111.5){\circle{.71}} 
\put(20.5,110){\circle{.5}}
\put(15,107){\circle{.5}} 
\put(10.5,100.75){\circle{.71}}
\put(44.5,100.5){\line(1,0){1.5}}
\qbezier(45.25,117.75)(45.38,117.25)(61,124.75)
\qbezier(61,125)(63.25,125.38)(65.5,127.25)
\qbezier(66,127.25)(69.13,128.63)(69.75,129.5)
\qbezier(69.75,129.5)(70.63,129.88)(70,129.75)
\qbezier(70.25,129)(73.63,131.38)(77.5,135.25)
\qbezier(45.5,117.75)(43,116.75)(25.5,111.75)
\qbezier(25.75,111.75)(22.38,111)(20.5,110.25)
\qbezier(20.5,110)(16.88,109.38)(14.75,107.25)
\qbezier(15,107)(10.75,105.5)(10.5,101)
\put(41.5,135){\makebox(0,0)[cc]{$\pi$}}
\put(48.75,129.25){\makebox(0,0)[cc]{$\pi/3$}}
\put(39.25,126.25){\makebox(0,0)[cc]{$\pi/4$}}
\put(48.75,124){\makebox(0,0)[cc]{$\pi/6$}}
\put(49,115){\makebox(0,0)[cc]{$0$}}
\put(61.5,115.5){\makebox(0,0)[cc]{$\frac{1}{2}$}}
\put(66.75,113.75){\makebox(0,0)[cc]{$\frac{\sqrt{2}}{2}$}}
\put(72,112.75){\makebox(0,0)[cc]{$\frac{\sqrt{3}}{3}$}}
\put(77.75,116.25){\makebox(0,0)[cc]{$1$}}
\put(25.5,115.5){\makebox(0,0)[cc]{$-\frac{1}{2}$}}
\put(20.5,114.75){\makebox(0,0)[cc]{$-\frac{\sqrt{2}}{2}$}}
\put(14.25,120.75){\makebox(0,0)[cc]{$-\frac{\sqrt{3}}{2}$}}
\put(9.75,115.75){\makebox(0,0)[cc]{$-1$}}
\put(47.75,111.75){\makebox(0,0)[cc]{$-\pi/6$}}
\put(40.75,109.5){\makebox(0,0)[cc]{$-\pi/4$}}
\put(49.5,106.25){\makebox(0,0)[cc]{$-\pi/3$}}
\put(40.75,100.25){\makebox(0,0)[cc]{$-\pi$}}
%\qbezvec(81.25,127.25)(74,121)(64.75,124.75)
\put(64.75,124.75){\vector(-3,1){.07}}\qbezier(81.25,127.25)(74,121)(64.75,124.75)
%\end
\put(82.5,129.5){\makebox(0,0)[cc]{$\sin^{-1}$}}
\end{picture}
\end{center}

\begin{enumerate}
  \item[\hspace{1cm}9)] To get a sketch of the graph of $\cos^{-1}$ on the interval $[2\pi,3\pi]$, first set up the following basic framework:
\end{enumerate}

\begin{center}
%TeXCAD Options
%\grade{\on}
%\emlines{\off}
%\epic{\off}
%\beziermacro{\on}
%\reduce{\on}
%\snapping{\off}
%\quality{8.00}
%\graddiff{0.01}
%\snapasp{1}
%\zoom{4.0000}
\unitlength 1mm % = 2.85pt
\linethickness{0.4pt}
\ifx\plotpoint\undefined\newsavebox{\plotpoint}\fi % GNUPLOT compatibility
\begin{picture}(78,41)(0,94)
\put(45.25,100.75){\line(0,1){34.25}}
\put(9.75,100.25){\line(1,0){68.25}} 
\put(9.5,101){\line(0,-1){1}}
\put(25.5,101.25){\line(0,-1){1.75}}
\put(20.5,101.25){\line(0,-1){1.5}}
\put(14.25,101.25){\line(0,-1){1.25}}
\put(77.5,101.5){\line(0,-1){1.5}} 
\put(62,101.5){\line(0,-1){1.25}}
\put(66.75,101.5){\line(0,-1){1}}
\put(71.75,101.5){\line(0,-1){1.5}}
\put(61.5,103.5){\makebox(0,0)[cc]{$\frac{1}{2}$}}
\put(66.75,96){\makebox(0,0)[cc]{$\frac{\sqrt{2}}{2}$}}
\put(72,103.5){\makebox(0,0)[cc]{$\frac{\sqrt{3}}{3}$}}
\put(77.75,96){\makebox(0,0)[cc]{$1$}}
\put(26.5,103.5){\makebox(0,0)[cc]{$-\frac{1}{2}$}}
\put(20.5,96){\makebox(0,0)[cc]{$-\frac{\sqrt{2}}{2}$}}
\put(14.25,103.5){\makebox(0,0)[cc]{$-\frac{\sqrt{3}}{2}$}}
\put(9.75,96){\makebox(0,0)[cc]{$-1$}}
\put(45.75,96){\makebox(0,0)[cc]{$0$}} 
\put(44.5,109){\line(1,0){1.5}}
\put(44.5,121){\line(1,0){1.5}}
\put(44.5,133){\line(1,0){1.5}}
\put(48.75,109){\makebox(0,0)[cc]{$\pi$}}
\put(48.75,121){\makebox(0,0)[cc]{$2\pi$}}
\put(49.5,133){\makebox(0,0)[cc]{$3\pi$}}
\end{picture}
\end{center}

\begin{enumerate}
  \item[\hspace{1cm}10)] Fill in a sketch of $\cos^{-1}$ on $[0,\pi]$ as you did in step 6.
  \item[\hspace{1cm}11)] The sketch of $\cos^{-1}$ on $[2\pi,3\pi]$ is congruent to that on $[0,\pi]$.
  \item[\hspace{1cm}12)] Your sketch should now look like the following: (b)
\end{enumerate}

\begin{center}
%TeXCAD Options
%\grade{\on}
%\emlines{\off}
%\epic{\off}
%\beziermacro{\on}
%\reduce{\on}
%\snapping{\off}
%\quality{8.00}
%\graddiff{0.01}
%\snapasp{1}
%\zoom{4.0000}
\unitlength 1mm % = 2.85pt
\linethickness{0.4pt}
\ifx\plotpoint\undefined\newsavebox{\plotpoint}\fi % GNUPLOT compatibility
\begin{picture}(78,41.25)(0,94)
\put(45.25,100){\line(0,1){34.25}}
\put(44.5,134.75){\line(1,0){1.5}} 
\put(9.75,100){\line(1,0){68.25}}
\put(9.5,100){\line(0,-1){1.5}} 
\put(25.5,100){\line(0,-1){1.75}}
\put(20.5,100){\line(0,-1){1.5}} 
\put(14.25,100){\line(0,-1){1.25}}
\put(77.5,100){\line(0,-1){1.5}} 
\put(62,100){\line(0,-1){1.25}}
\put(66.75,100){\line(0,-1){3}} 
\put(71.75,100){\line(0,-1){7.5}}
\put(44.5,108.75){\line(1,0){1.5}}
\put(44.75,128){\line(1,0){1.25}}
\put(44.75,104.5){\line(1,0){1.25}}
\put(44.5,119){\line(1,0){1.5}}
\qbezier(9.75,109)(11,107.63)(45.25,104.75)
\qbezier(45.25,104.5)(67.38,105.13)(78,100.25)
\qbezier(77.25,120)(64.75,128)(45.25,128)
\qbezier(45.25,127.75)(23.13,127.5)(9.5,135.25)
\put(61.5,96){\makebox(0,0)[cc]{$\frac{1}{2}$}}
\put(66.75,94){\makebox(0,0)[cc]{$\frac{\sqrt{2}}{2}$}}
\put(72,90){\makebox(0,0)[cc]{$\frac{\sqrt{3}}{3}$}}
\put(77.75,96){\makebox(0,0)[cc]{$1$}}
\put(25.5,103){\makebox(0,0)[cc]{$-\frac{1}{2}$}}
\put(20.5,96){\makebox(0,0)[cc]{$-\frac{\sqrt{2}}{2}$}}
\put(14.25,103){\makebox(0,0)[cc]{$-\frac{\sqrt{3}}{2}$}}
\put(9.75,96){\makebox(0,0)[cc]{$-1$}}
\put(45.75,96){\makebox(0,0)[cc]{$0$}} 
\put(48.75,118.75){\makebox(0,0)[cc]{$2\pi$}}
\put(48.75,108.5){\makebox(0,0)[cc]{$\pi$}}
\put(49.5,133.75){\makebox(0,0)[cc]{$3\pi$}}
\end{picture}
\end{center}

Are you catching on?

\vspace{0.3cm}

(c) infinitely many.

\vspace{0.3cm}

Does it seem to you that given the sketch of cosine in $[0,\pi]$, we ought to be able to get $\cos^{-1}$ very easily? Given the point $(1,3)$, where is $(3,1)$? Given the point $(x,y)$, where is $(y,x)$? I'll bet you can find some slick way to get a sketch of the graph of $\cos^{-1}$ without a lot of painful point plotting.

\vspace{0.3cm}

\textbf{Problem 6.4}:
\begin{enumerate}[\hspace{1cm}1)]
  \item suppose $P:(u,v)$ is a point of $K:x^{2}+y^{2} = 1$. Draw a picture.
  \item Note that $u = \cos \theta$ and $v = \sin \theta$ for some $\theta$ such that $0 \leq \theta \leq 2\pi$.
  \item What is $l(\arc{PI})$?
  \item $l(\arc{PI})$ depends on whether $0 \leq \theta \leq \pi$ or $\pi < \theta < 2$.
  \item $l(\arc{PI}) = \theta$ or, $l(\arc{PI}) = 2\pi - \theta$.
  \item Find $PI$.
  \item $PI = \sqrt{(\cos \theta - 1)^{2}+(\sin \theta - 0)^{2}} = \sqrt{\cos^{2} \theta - 2 \cos \theta + 1 + \sin^{2} \theta} = \sqrt{2 - 2 \cos \theta}$

        Hence, $PI = \sqrt{2 - 2 \cos \theta}$, and $l(\arc{PI}) = \theta$ or $2\pi - \theta$, depending on where $\theta$ is in $[0,2\pi]$. Note that $\cos(2\pi - \theta) = \cos \theta$.

  \item If we are given $PI$, then PI = $\sqrt{2 - \cos\theta}$ , $(PI)^{2} = 2 -2 \cos \theta$ and $\cos \theta = \frac{ (2-(PI)^{2}) }{2}$ and $\theta = \frac{ \cos((2-(PI)^{2}) }{2})$

        If we are given $l(\arc{PI})$, then $l(\arc{PI}) = \theta$ or $2\pi - \theta$, so that $PI = \sqrt{2 - 2 \cos \theta}$.

        Since $\cos (2\pi - \theta) = \cos \theta$, $PI = \sqrt{2 - 2 \cos (l(\arc{PI}))}$.
\end{enumerate}

\textbf{Problem 6.5}: Now look at the more general problem of two points $P$ and $Q$ on $K: x^{2}+y^{2} = 1$.
\begin{enumerate}[\hspace{1cm}1)]
  \item Draw a picture.
  \item Write the coordinates of $P$ and $Q$ in terms of the functions cosine and sine.
  \item You should have $P:(\cos \phi, \sin \phi)$ for some number $\phi$ such that $0 \leq \phi < 2\pi$ and $Q:(\cos \theta, \sin \theta)$ for some number $\theta$ such that $0 \leq \theta < 2\pi$. (You could have used any symbols you want here, not just $\phi$ and $\theta$.) Suppose $0 < \theta < \phi < 2\pi$. One of them must be bigger than the other, or else we don't have two points in $P$ and $Q$.
  \item Note that $l(\arc{PQ}) = \phi - \theta$, or $l(\arc{PQ}) = 2\pi - (\phi - \theta)$.
  \item Show that $PQ = \sqrt{2} \sqrt{1 - \cos (\phi - \theta)}$.
  \item Show that on the basis on step 5, $\phi - \theta = \cos^{-1} (1- \frac{(PQ)^{2}}{2})$.
  \item Conclude that $l(\arc{PQ}) = \cos^{-1} (1- \frac{(PQ)^{2}}{2})$, or $l(\arc{PQ}) = 2\pi - \cos^{-1} (1- \frac{(PQ)^{2}}{2})$. How do you decide which one of these is the solution?
\end{enumerate}

If we know one of $PQ$ and $l(\arc{PQ})$, we can find the other. At least we can write down what we call a solution to the problem. As we have seen, however, getting numerical results means, in most cases, getting some approximation to what we want. But then, most things that we encounter in life are approximations, aren't they? Can you really cut a board exactly one foot long? Can you really weigh out exactly one pound of gold? Not really. As we progress in our science and technology, we develop better tools that allow us to make more accurate measurements. We always seem to be searching for one more decimal place of accuracy. But we're always going to have trouble measuring anything exactly. Note the distinction between measuring and counting. We can count five beans, but when we weigh five beans, one thing we can be absolutely certain of is that our weighing is wrong.

We have solved our main problem on the unit circle, at least in terms of our new functions sine and cosine. We have seen, however, that numerical solutions usually involve approximations, and although I tried to convince you that approximations are not abnormal, you might find this a tad unsatisfactory. So, let me make another stab at it. 

Consider this: suppose we are working on some cutting edge technology, and, for some reason, we need to know the length of the arc on the unit circle from the point on the circle with coordinates $P:( \cos 2.5, \sin 2.5)$ to the point with coordinates $Q:( \cos 1.26, \sin 1.26)$; i.e., we need $l(\arc{PQ})$. Suppose also that we need to know the distance between them; i.e. $PQ$. What are $\cos 2.5$, $\sin 2.5$, $\cos 1.26$, and $\sin 1.26$? We might find a table of values for our functions sine and cosine, and the numbers $2.5$ and $1.26$ might even be in the table. Even if they are, the entries listed above will be given as decimal numbers and will only be approximations. Your table may go out to ten decimal places. How much accuracy do you need? You may have some method of determining how close you need for your results to be to the correct result. Furthermore, I'll bet that if you need your results to be within one ten-trillionth of the correct result, then you can see how to get such accuracy as long as you had the time and the patience to do it. 

Of course, a computer would be a big help. Since the personal computer and personal calculator have become so popular, tables of values of our new functions aren't nearly as useful as they used to be. You want $\cos 2.5$? Simply put your hand-held calculator on the radian mode (we'll talk about what this means later), punch in 2.5, and hit the $\cos$ button. My calculator gives $-0.801143615$. What does yours say? 

Even after you get numerical approximations to these numbers $(\cos 2.5,$ $\sin 2.5)$, $(\cos 1.26, \sin 1.26)$, what do you do with them? Look back at the formulas we got just a few pages ago. $PQ = \sqrt{2} \sqrt{1- \cos (\phi - \theta)}$. Our $\phi$ and $\theta$ in this case are $2.5$ and $1.26$, respectively. Hence, $PQ =$ $\sqrt{2} \sqrt{1- \cos (2.5 - 1.26)} =$ $\sqrt{2} \sqrt{1- \cos (1.24)} \doteq$ $\sqrt{2} \sqrt{1 - 0.324796284}$ (which I got from my Casio FX-115 scientific calculator which, if they still made them today, would probably be given away with the purchase of a full tank of gasoline). 

Note the introduction of the symbol ``$\doteq$''. Since I have introduced the decimal approximation of $\cos l.26$, I no longer have equality, only an approximation to it. Hence, I must switch from $=$ to $\doteq$. Hence, $PQ \doteq \sqrt{2} \sqrt{0.675203715} \doteq  1.162070321$. Note the $\doteq$ here also. Do you see why? 

But, what about $l(\arc{PQ})$? Our formula that we got earlier is $l(\arc{PQ}) = \cos^{-1} (1- \frac{(PQ)^{2}}{2})$, or $l(\arc{PQ}) = 2\pi - \cos^{-1} (1- \frac{(PQ)^{2}}{2})$. Which is it? $PQ \doteq \sqrt{2} \sqrt{1.675203715}$, so $\frac{(PQ)^{2}}{2} \doteq 0.675203715$, and $1- \frac{(PQ)^{2}}{2} \doteq 0.324796284$.

That's a familiar number, isn't it? It's $\cos 1.24 = \cos (2.5 - 1.26) = \cos (\phi - \theta)$. Hence, $l(\arc{PQ}) = \cos^{-1} (\cos (\phi - \theta)) = \cos^{-1} (\cos 1.24) = 1.24$. Recall that we can't always say that $\cos^{-1} (\cos x) = x$. Why can we do it here? We knew this all along, didn't we? In this case, we used $l(\arc{PQ}) = \cos^{-1} (1-\frac{(PQ)^{2}}{2})$. Where would we use $l(\arc{PQ}) = 2\pi - \cos^{-1} (1- \frac{(PQ)^{2}}{2})$? 

And here is the {\$}64,000 question: How do we solve such problems as these if we're on a circle which is not the unit circle? Anyway, the point that I'm trying to make is that even though we are doomed to make approximations, we can usually get our results as close as we want or need them to be. Our limitations are usually in the area of physical measurements, not numerical calculations.

\vspace{0.3cm}

\textbf{Problem 6.6}:
\begin{enumerate}
  \item[\hspace{1cm}a)] A sketch of the tangent function is the following:
\end{enumerate}

\begin{center}
%TeXCAD Options
%\grade{\on}
%\emlines{\off}
%\epic{\off}
%\beziermacro{\on}
%\reduce{\on}
%\snapping{\off}
%\quality{8.00}
%\graddiff{0.01}
%\snapasp{1}
%\zoom{4.0000}
\unitlength 0.75mm % = 2.85pt
\linethickness{0.4pt}
\ifx\plotpoint\undefined\newsavebox{\plotpoint}\fi % GNUPLOT compatibility
\begin{picture}(150,79.25)(0,60)
\put(27.25,98){\line(1,0){89.5}}
%\dashline{1}(120.5,138.75)(120.5,66.5)
\put(120.43,138.68){\line(0,-1){.99}}
\put(120.43,136.7){\line(0,-1){.99}}
\put(120.43,134.72){\line(0,-1){.99}}
\put(120.43,132.74){\line(0,-1){.99}}
\put(120.43,130.76){\line(0,-1){.99}}
\put(120.43,128.78){\line(0,-1){.99}}
\put(120.43,126.8){\line(0,-1){.99}}
\put(120.43,124.82){\line(0,-1){.99}}
\put(120.43,122.84){\line(0,-1){.99}}
\put(120.43,120.86){\line(0,-1){.99}}
\put(120.43,118.89){\line(0,-1){.99}}
\put(120.43,116.91){\line(0,-1){.99}}
\put(120.43,114.93){\line(0,-1){.99}}
\put(120.43,112.95){\line(0,-1){.99}}
\put(120.43,110.97){\line(0,-1){.99}}
\put(120.43,108.99){\line(0,-1){.99}}
\put(120.43,107.01){\line(0,-1){.99}}
\put(120.43,105.03){\line(0,-1){.99}}
\put(120.43,103.05){\line(0,-1){.99}}
\put(120.43,101.07){\line(0,-1){.99}}
\put(120.43,99.09){\line(0,-1){.99}}
\put(120.43,97.11){\line(0,-1){.99}}
\put(120.43,95.13){\line(0,-1){.99}}
\put(120.43,93.15){\line(0,-1){.99}}
\put(120.43,91.17){\line(0,-1){.99}}
\put(120.43,89.19){\line(0,-1){.99}}
\put(120.43,87.21){\line(0,-1){.99}}
\put(120.43,85.23){\line(0,-1){.99}}
\put(120.43,83.26){\line(0,-1){.99}}
\put(120.43,81.28){\line(0,-1){.99}}
\put(120.43,79.3){\line(0,-1){.99}}
\put(120.43,77.32){\line(0,-1){.99}}
\put(120.43,75.34){\line(0,-1){.99}}
\put(120.43,73.36){\line(0,-1){.99}}
\put(120.43,71.38){\line(0,-1){.99}}
\put(120.43,69.4){\line(0,-1){.99}}
\put(120.43,67.42){\line(0,-1){.99}}
%\end
%\dashline{1}(93.25,139)(93.25,67.25)
\put(93.18,138.93){\line(0,-1){.997}}
\put(93.18,136.94){\line(0,-1){.997}}
\put(93.18,134.94){\line(0,-1){.997}}
\put(93.18,132.95){\line(0,-1){.997}}
\put(93.18,130.96){\line(0,-1){.997}}
\put(93.18,128.96){\line(0,-1){.997}}
\put(93.18,126.97){\line(0,-1){.997}}
\put(93.18,124.98){\line(0,-1){.997}}
\put(93.18,122.99){\line(0,-1){.997}}
\put(93.18,120.99){\line(0,-1){.997}}
\put(93.18,119){\line(0,-1){.997}}
\put(93.18,117.01){\line(0,-1){.997}}
\put(93.18,115.01){\line(0,-1){.997}}
\put(93.18,113.02){\line(0,-1){.997}}
\put(93.18,111.03){\line(0,-1){.997}}
\put(93.18,109.03){\line(0,-1){.997}}
\put(93.18,107.04){\line(0,-1){.997}}
\put(93.18,105.05){\line(0,-1){.997}}
\put(93.18,103.05){\line(0,-1){.997}}
\put(93.18,101.06){\line(0,-1){.997}}
\put(93.18,99.07){\line(0,-1){.997}}
\put(93.18,97.08){\line(0,-1){.997}}
\put(93.18,95.08){\line(0,-1){.997}}
\put(93.18,93.09){\line(0,-1){.997}}
\put(93.18,91.1){\line(0,-1){.997}}
\put(93.18,89.1){\line(0,-1){.997}}
\put(93.18,87.11){\line(0,-1){.997}}
\put(93.18,85.12){\line(0,-1){.997}}
\put(93.18,83.12){\line(0,-1){.997}}
\put(93.18,81.13){\line(0,-1){.997}}
\put(93.18,79.14){\line(0,-1){.997}}
\put(93.18,77.14){\line(0,-1){.997}}
\put(93.18,75.15){\line(0,-1){.997}}
\put(93.18,73.16){\line(0,-1){.997}}
\put(93.18,71.17){\line(0,-1){.997}}
\put(93.18,69.17){\line(0,-1){.997}}
%\end
%\dashline{1}(64,67.25)(64,138.75)
\put(63.93,67.18){\line(0,1){.993}}
\put(63.93,69.17){\line(0,1){.993}}
\put(63.93,71.15){\line(0,1){.993}}
\put(63.93,73.14){\line(0,1){.993}}
\put(63.93,75.12){\line(0,1){.993}}
\put(63.93,77.11){\line(0,1){.993}}
\put(63.93,79.1){\line(0,1){.993}}
\put(63.93,81.08){\line(0,1){.993}}
\put(63.93,83.07){\line(0,1){.993}}
\put(63.93,85.05){\line(0,1){.993}}
\put(63.93,87.04){\line(0,1){.993}}
\put(63.93,89.03){\line(0,1){.993}}
\put(63.93,91.01){\line(0,1){.993}} 
\put(63.93,93){\line(0,1){.993}}
\put(63.93,94.99){\line(0,1){.993}}
\put(63.93,96.97){\line(0,1){.993}}
\put(63.93,98.96){\line(0,1){.993}}
\put(63.93,100.94){\line(0,1){.993}}
\put(63.93,102.93){\line(0,1){.993}}
\put(63.93,104.92){\line(0,1){.993}}
\put(63.93,106.9){\line(0,1){.993}}
\put(63.93,108.89){\line(0,1){.993}}
\put(63.93,110.87){\line(0,1){.993}}
\put(63.93,112.86){\line(0,1){.993}}
\put(63.93,114.85){\line(0,1){.993}}
\put(63.93,116.83){\line(0,1){.993}}
\put(63.93,118.82){\line(0,1){.993}}
\put(63.93,120.8){\line(0,1){.993}}
\put(63.93,122.79){\line(0,1){.993}}
\put(63.93,124.78){\line(0,1){.993}}
\put(63.93,126.76){\line(0,1){.993}}
\put(63.93,128.75){\line(0,1){.993}}
\put(63.93,130.74){\line(0,1){.993}}
\put(63.93,132.72){\line(0,1){.993}}
\put(63.93,134.71){\line(0,1){.993}}
\put(63.93,136.69){\line(0,1){.993}}
%\end
%\dashline{1}(34.25,66.5)(34.25,138.25)
\put(34.18,66.43){\line(0,1){.997}}
\put(34.18,68.42){\line(0,1){.997}}
\put(34.18,70.42){\line(0,1){.997}}
\put(34.18,72.41){\line(0,1){.997}}
\put(34.18,74.4){\line(0,1){.997}}
\put(34.18,76.39){\line(0,1){.997}}
\put(34.18,78.39){\line(0,1){.997}}
\put(34.18,80.38){\line(0,1){.997}}
\put(34.18,82.37){\line(0,1){.997}}
\put(34.18,84.37){\line(0,1){.997}}
\put(34.18,86.36){\line(0,1){.997}}
\put(34.18,88.35){\line(0,1){.997}}
\put(34.18,90.35){\line(0,1){.997}}
\put(34.18,92.34){\line(0,1){.997}}
\put(34.18,94.33){\line(0,1){.997}}
\put(34.18,96.33){\line(0,1){.997}}
\put(34.18,98.32){\line(0,1){.997}}
\put(34.18,100.31){\line(0,1){.997}}
\put(34.18,102.3){\line(0,1){.997}}
\put(34.18,104.3){\line(0,1){.997}}
\put(34.18,106.29){\line(0,1){.997}}
\put(34.18,108.28){\line(0,1){.997}}
\put(34.18,110.28){\line(0,1){.997}}
\put(34.18,112.27){\line(0,1){.997}}
\put(34.18,114.26){\line(0,1){.997}}
\put(34.18,116.26){\line(0,1){.997}}
\put(34.18,118.25){\line(0,1){.997}}
\put(34.18,120.24){\line(0,1){.997}}
\put(34.18,122.24){\line(0,1){.997}}
\put(34.18,124.23){\line(0,1){.997}}
\put(34.18,126.22){\line(0,1){.997}}
\put(34.18,128.21){\line(0,1){.997}}
\put(34.18,130.21){\line(0,1){.997}}
\put(34.18,132.2){\line(0,1){.997}}
\put(34.18,134.19){\line(0,1){.997}}
\put(34.18,136.19){\line(0,1){.997}}
%\end
\put(116.75,98){\line(1,0){10.75}} 
\put(127.5,98){\line(1,0){10}}
\put(21.75,98){\line(1,0){6.5}} 
\put(137,98){\line(1,0){13}}
\put(22.25,98){\line(-1,0){13}} 
\put(20,99){\line(0,-1){1.5}}
\put(49.25,99){\line(0,-1){1.5}} 
\put(79.25,99.25){\line(0,-1){2}}
\put(108.25,99.25){\line(0,-1){1.75}}
\put(136.5,99){\line(0,-1){1.5}}
\qbezier(34,136.75)(29.5,115.5)(20,98.25)
\qbezier(34.75,66.25)(40,83.38)(49.25,98)
\qbezier(49.5,98)(59.38,111.75)(63.75,137.5)
\qbezier(64.75,67)(70.5,83.5)(79.25,98)
\qbezier(79.5,98)(86.13,105.88)(93.25,139.25)
\qbezier(120.25,138.5)(114.25,112.13)(108.25,98.25)
\qbezier(108.5,98)(97.63,82.88)(93.25,68.25)
\qbezier(121,66.25)(124.75,80.25)(136.5,98.25)
\put(80.5,94.75){\makebox(0,0)[cc]{$0$}}
\put(110.25,95.25){\makebox(0,0)[cc]{$\pi$}}
\put(95.25,100.75){\makebox(0,0)[cc]{$\frac{\pi}{2}$}}
\put(137.75,94.75){\makebox(0,0)[cc]{$2\pi$}}
\put(123.25,100.5){\makebox(0,0)[cc]{$\frac{3\pi}{2}$}}
\put(66.5,100.5){\makebox(0,0)[cc]{$-\frac{\pi}{2}$}}
\put(52.5,94.5){\makebox(0,0)[cc]{$-\pi$}}
\put(18.25,95.25){\makebox(0,0)[cc]{$-2\pi$}}
\put(38,100.5){\makebox(0,0)[cc]{$-\frac{3\pi}{2}$}}
\end{picture}
\end{center}

Some comments are in order.
\begin{enumerate}[\hspace{1cm}1)]
  \item For those places where the function is undefined, it is useful to sketch in a vertical dotted line.
  \item Find all the places where the function is $0$. The graph intersects the $x$-axis at these numbers.
  \item Consider what happens to the graph as it gets closer to one of the dotted vertical lines. For example, $\tan x = \frac{\sin x}{\cos x}$, and $\tan x$ is not defined at $\frac{\pi}{2}$. Why is it not defined at $\frac{\pi}{2}$? It's because $\cos (\frac{\pi}{2}) = 0$. What happens to $\frac{\sin x}{\cos x}$ as $x$ gets closer and closer to $\frac{\pi}{2}$ from below $\frac{\pi}{2}$? Well, $\sin x$ gets closer to $1$ and $\cos x$ gets closer to $0$. At any step along the way, this quotient is positive, and as you get closer to $\frac{\pi}{2}$, the quotient gets larger and larger. In fact, it increases without any bound as you get closer to $\frac{\pi}{2}$ from below. Furthermore, the graph of $\tan x$ gets closer to the vertical line through $\frac{\pi}{2}$.
  \item Now examine what happens to the tangent graph on the other side (the right side) of the vertical line through $\frac{\pi}{2}$. As you get closer and closer to $\frac{\pi}{2}$ from above, $\tan x = \frac{\sin x}{\cos x}$ becomes negative since $\sin x > 0$ but $\cos x < 0$ in that area. Furthermore, $\cos x$ gets closer to $0$ and $\sin x$ gets closer to $1$, so that in this case, the quotient decreases without any bound.
  \item The same analysis applies to each of the vertical lines through an odd multiple of $\frac{\pi}{2}$. As a result, you can get a good idea of what the graph must look like without having to calculate and plot a whole bunch of points. 
\end{enumerate}

Here is a nice language puzzle for you, pertaining to the graph of the tangent function. Put the sketch of this graph in front of you, and keep in mind that it repeats over every interval of length $\pi$ whose left end is an odd multiple of $\frac{\pi}{2}$. I am going to give you four statements about the tangent graph that look somewhat alike. Your job is to analyze each statement to see if it is true. Here we go.

\vfill

Suppose $L$ is the vertical line containing the number $\frac{\pi}{2}$ on the $x$-axis, and $T$ is the graph of the tangent function.

\vfill

\begin{enumerate}[\hspace{1cm}1)]
  \item If $V$ is a vertical line to the left of $L$, then there exists a horizontal line $H$ such that every point of $T$ above $H$ is between $V$ and $L$.
  \item If $V$ is a vertical line to the left of $L$, then there exists a horizontal line $H$ such that every point of $T$ between $V$ and $L$ is above $H$.
  \item If $H$ is a horizontal line, then there exists a vertical line $V$ to the left of $L$ such that every point of $T$ above $H$ is between $V$ and $L$.
  \item If $H$ is a horizontal line, then there exists a vertical line $V$ to the left of $L$ such that every point of $T$ between $V$ and $L$ is above $H$.
\end{enumerate}

\vfill

Are any of these true? Are they all true? Draw some pictures. You will note that subtle differences in language can sometimes make a huge difference in meaning. Incidentally, the line $L$ containing $\frac{\pi}{2}$ on the $x$-axis is called an \emph{asymptote} to the tangent graph, and the tangent graph is said to be \emph{asymptotic} to that line. It looks like the tangent graph has a plethora of asymptotes. (How many is a plethora?)

\vfill

\begin{enumerate}
  \item[\hspace{1cm}b)] Looking at the graph of $\tan x$, it appears that the interval $[-\frac{\pi}{2},\frac{\pi}{2}]$ is a good candidate. However, $\frac{\pi}{2}$ and $\frac{\pi}{2}$ are not even in the initial set of the tangent function. It is defined at every number in between $-\frac{\pi}{2}$ and $\frac{\pi}{2}$, and that set, known as a \emph{segment}, is denoted $(-\frac{\pi}{2},\frac{\pi}{2})$. Sometimes $(-\frac{\pi}{2},\frac{\pi}{2})$ is referred to as an \emph{open interval}. (Isn't that fascinating?) So $\tan^{-1}$ is a function over the segment $(-\frac{\pi}{2},\frac{\pi}{2})$, and its principal values are the numbers in that segment.
  \item[\hspace{1cm}c)] 
    \begin{enumerate}[\hspace{1cm}1)]
      \item Here is a sketch of $\tan^{-1}$.
    \end{enumerate}
\end{enumerate}

\begin{center}
%TeXCAD Options
%\grade{\on}
%\emlines{\off}
%\epic{\off}
%\beziermacro{\on}
%\reduce{\on}
%\snapping{\off}
%\quality{8.00}
%\graddiff{0.01}
%\snapasp{1}
%\zoom{4.0000}
\unitlength 1mm % = 2.85pt
\linethickness{0.4pt}
\ifx\plotpoint\undefined\newsavebox{\plotpoint}\fi % GNUPLOT compatibility
\begin{picture}(117,80.25)(0,47)
\put(63.5,127.25){\line(0,-1){68.25}}
%\dashline{1}(15.75,114)(114.75,114.5)
\put(15.68,113.93){\line(1,0){.99}}
\put(17.66,113.94){\line(1,0){.99}}
\put(19.64,113.95){\line(1,0){.99}}
\put(21.62,113.96){\line(1,0){.99}}
\put(23.6,113.97){\line(1,0){.99}}
\put(25.58,113.98){\line(1,0){.99}}
\put(27.56,113.99){\line(1,0){.99}} 
\put(29.54,114){\line(1,0){.99}}
\put(31.52,114.01){\line(1,0){.99}}
\put(33.5,114.02){\line(1,0){.99}}
\put(35.48,114.03){\line(1,0){.99}}
\put(37.46,114.04){\line(1,0){.99}}
\put(39.44,114.05){\line(1,0){.99}}
\put(41.42,114.06){\line(1,0){.99}}
\put(43.4,114.07){\line(1,0){.99}}
\put(45.38,114.08){\line(1,0){.99}}
\put(47.36,114.09){\line(1,0){.99}}
\put(49.34,114.1){\line(1,0){.99}}
\put(51.32,114.11){\line(1,0){.99}}
\put(53.3,114.12){\line(1,0){.99}}
\put(55.28,114.13){\line(1,0){.99}}
\put(57.26,114.14){\line(1,0){.99}}
\put(59.24,114.15){\line(1,0){.99}}
\put(61.22,114.16){\line(1,0){.99}}
\put(63.2,114.17){\line(1,0){.99}}
\put(65.18,114.18){\line(1,0){.99}}
\put(67.16,114.19){\line(1,0){.99}}
\put(69.14,114.2){\line(1,0){.99}}
\put(71.12,114.21){\line(1,0){.99}}
\put(73.1,114.22){\line(1,0){.99}}
\put(75.08,114.23){\line(1,0){.99}}
\put(77.06,114.24){\line(1,0){.99}}
\put(79.04,114.25){\line(1,0){.99}}
\put(81.02,114.26){\line(1,0){.99}} 
\put(83,114.27){\line(1,0){.99}}
\put(84.98,114.28){\line(1,0){.99}}
\put(86.96,114.29){\line(1,0){.99}}
\put(88.94,114.3){\line(1,0){.99}}
\put(90.92,114.31){\line(1,0){.99}}
\put(92.9,114.32){\line(1,0){.99}}
\put(94.88,114.33){\line(1,0){.99}}
\put(96.86,114.34){\line(1,0){.99}}
\put(98.84,114.35){\line(1,0){.99}}
\put(100.82,114.36){\line(1,0){.99}}
\put(102.8,114.37){\line(1,0){.99}}
\put(104.78,114.38){\line(1,0){.99}}
\put(106.76,114.39){\line(1,0){.99}}
\put(108.74,114.4){\line(1,0){.99}}
\put(110.72,114.41){\line(1,0){.99}}
\put(112.7,114.42){\line(1,0){.99}}
%\end
%\dashline{1}(115,96.75)(13.75,96.75)
\put(114.93,96.68){\line(-1,0){.993}}
\put(112.94,96.68){\line(-1,0){.993}}
\put(110.96,96.68){\line(-1,0){.993}}
\put(108.97,96.68){\line(-1,0){.993}}
\put(106.99,96.68){\line(-1,0){.993}}
\put(105,96.68){\line(-1,0){.993}}
\put(103.02,96.68){\line(-1,0){.993}}
\put(101.03,96.68){\line(-1,0){.993}}
\put(99.05,96.68){\line(-1,0){.993}}
\put(97.06,96.68){\line(-1,0){.993}}
\put(95.08,96.68){\line(-1,0){.993}}
\put(93.09,96.68){\line(-1,0){.993}}
\put(91.11,96.68){\line(-1,0){.993}}
\put(89.12,96.68){\line(-1,0){.993}}
\put(87.14,96.68){\line(-1,0){.993}}
\put(85.15,96.68){\line(-1,0){.993}}
\put(83.17,96.68){\line(-1,0){.993}}
\put(81.18,96.68){\line(-1,0){.993}}
\put(79.19,96.68){\line(-1,0){.993}}
\put(77.21,96.68){\line(-1,0){.993}}
\put(75.22,96.68){\line(-1,0){.993}}
\put(73.24,96.68){\line(-1,0){.993}}
\put(71.25,96.68){\line(-1,0){.993}}
\put(69.27,96.68){\line(-1,0){.993}}
\put(67.28,96.68){\line(-1,0){.993}}
\put(65.3,96.68){\line(-1,0){.993}}
\put(63.31,96.68){\line(-1,0){.993}}
\put(61.33,96.68){\line(-1,0){.993}}
\put(59.34,96.68){\line(-1,0){.993}}
\put(57.36,96.68){\line(-1,0){.993}}
\put(55.37,96.68){\line(-1,0){.993}}
\put(53.39,96.68){\line(-1,0){.993}}
\put(51.4,96.68){\line(-1,0){.993}}
\put(49.42,96.68){\line(-1,0){.993}}
\put(47.43,96.68){\line(-1,0){.993}}
\put(45.44,96.68){\line(-1,0){.993}}
\put(43.46,96.68){\line(-1,0){.993}}
\put(41.47,96.68){\line(-1,0){.993}}
\put(39.49,96.68){\line(-1,0){.993}}
\put(37.5,96.68){\line(-1,0){.993}}
\put(35.52,96.68){\line(-1,0){.993}}
\put(33.53,96.68){\line(-1,0){.993}}
\put(31.55,96.68){\line(-1,0){.993}}
\put(29.56,96.68){\line(-1,0){.993}}
\put(27.58,96.68){\line(-1,0){.993}}
\put(25.59,96.68){\line(-1,0){.993}}
\put(23.61,96.68){\line(-1,0){.993}}
\put(21.62,96.68){\line(-1,0){.993}}
\put(19.64,96.68){\line(-1,0){.993}}
\put(17.65,96.68){\line(-1,0){.993}}
\put(15.67,96.68){\line(-1,0){.993}}
%\end
%\dashline{1}(115.5,79.75)(14.75,80.25)
\put(115.43,79.68){\line(-1,0){.998}}
\put(113.43,79.69){\line(-1,0){.998}}
\put(111.44,79.7){\line(-1,0){.998}}
\put(109.44,79.71){\line(-1,0){.998}}
\put(107.45,79.72){\line(-1,0){.998}}
\put(105.45,79.73){\line(-1,0){.998}}
\put(103.46,79.74){\line(-1,0){.998}}
\put(101.46,79.75){\line(-1,0){.998}}
\put(99.47,79.76){\line(-1,0){.998}}
\put(97.47,79.77){\line(-1,0){.998}}
\put(95.48,79.78){\line(-1,0){.998}}
\put(93.48,79.79){\line(-1,0){.998}}
\put(91.49,79.8){\line(-1,0){.998}}
\put(89.49,79.81){\line(-1,0){.998}}
\put(87.5,79.82){\line(-1,0){.998}}
\put(85.5,79.83){\line(-1,0){.998}}
\put(83.51,79.84){\line(-1,0){.998}}
\put(81.51,79.85){\line(-1,0){.998}}
\put(79.52,79.86){\line(-1,0){.998}}
\put(77.52,79.87){\line(-1,0){.998}}
\put(75.53,79.88){\line(-1,0){.998}}
\put(73.53,79.89){\line(-1,0){.998}}
\put(71.54,79.9){\line(-1,0){.998}}
\put(69.54,79.91){\line(-1,0){.998}}
\put(67.55,79.92){\line(-1,0){.998}}
\put(65.55,79.93){\line(-1,0){.998}}
\put(63.56,79.94){\line(-1,0){.998}}
\put(61.56,79.95){\line(-1,0){.998}}
\put(59.57,79.96){\line(-1,0){.998}}
\put(57.57,79.97){\line(-1,0){.998}}
\put(55.58,79.98){\line(-1,0){.998}}
\put(53.58,79.99){\line(-1,0){.998}}
\put(51.59,80){\line(-1,0){.998}}
\put(49.59,80.01){\line(-1,0){.998}}
\put(47.6,80.02){\line(-1,0){.998}}
\put(45.6,80.03){\line(-1,0){.998}}
\put(43.61,80.04){\line(-1,0){.998}}
\put(41.61,80.05){\line(-1,0){.998}}
\put(39.62,80.06){\line(-1,0){.998}}
\put(37.62,80.07){\line(-1,0){.998}}
\put(35.63,80.08){\line(-1,0){.998}}
\put(33.63,80.09){\line(-1,0){.998}}
\put(31.64,80.1){\line(-1,0){.998}}
\put(29.64,80.11){\line(-1,0){.998}}
\put(27.65,80.12){\line(-1,0){.998}}
\put(25.65,80.13){\line(-1,0){.998}}
\put(23.66,80.14){\line(-1,0){.998}}
\put(21.66,80.15){\line(-1,0){.998}}
\put(19.67,80.15){\line(-1,0){.998}}
\put(17.67,80.16){\line(-1,0){.998}}
\put(15.68,80.17){\line(-1,0){.998}}
%\end
%\dashline{1}(117,61)(14.5,62)
\put(116.93,60.93){\line(-1,0){.995}}
\put(114.94,60.95){\line(-1,0){.995}}
\put(112.95,60.97){\line(-1,0){.995}}
\put(110.96,60.99){\line(-1,0){.995}}
\put(108.97,61.01){\line(-1,0){.995}}
\put(106.98,61.03){\line(-1,0){.995}}
\put(104.99,61.05){\line(-1,0){.995}}
\put(103,61.07){\line(-1,0){.995}}
\put(101.01,61.09){\line(-1,0){.995}}
\put(99.02,61.1){\line(-1,0){.995}}
\put(97.03,61.12){\line(-1,0){.995}}
\put(95.04,61.14){\line(-1,0){.995}}
\put(93.05,61.16){\line(-1,0){.995}}
\put(91.06,61.18){\line(-1,0){.995}}
\put(89.07,61.2){\line(-1,0){.995}}
\put(87.08,61.22){\line(-1,0){.995}}
\put(85.09,61.24){\line(-1,0){.995}}
\put(83.09,61.26){\line(-1,0){.995}}
\put(81.1,61.28){\line(-1,0){.995}}
\put(79.11,61.3){\line(-1,0){.995}}
\put(77.12,61.32){\line(-1,0){.995}}
\put(75.13,61.34){\line(-1,0){.995}}
\put(73.14,61.36){\line(-1,0){.995}}
\put(71.15,61.38){\line(-1,0){.995}}
\put(69.16,61.4){\line(-1,0){.995}}
\put(67.17,61.42){\line(-1,0){.995}}
\put(65.18,61.43){\line(-1,0){.995}}
\put(63.19,61.45){\line(-1,0){.995}}
\put(61.2,61.47){\line(-1,0){.995}}
\put(59.21,61.49){\line(-1,0){.995}}
\put(57.22,61.51){\line(-1,0){.995}}
\put(55.23,61.53){\line(-1,0){.995}}
\put(53.24,61.55){\line(-1,0){.995}}
\put(51.25,61.57){\line(-1,0){.995}}
\put(49.26,61.59){\line(-1,0){.995}}
\put(47.27,61.61){\line(-1,0){.995}}
\put(45.28,61.63){\line(-1,0){.995}}
\put(43.29,61.65){\line(-1,0){.995}}
\put(41.3,61.67){\line(-1,0){.995}}
\put(39.31,61.69){\line(-1,0){.995}}
\put(37.32,61.71){\line(-1,0){.995}}
\put(35.33,61.73){\line(-1,0){.995}}
\put(33.34,61.75){\line(-1,0){.995}}
\put(31.35,61.76){\line(-1,0){.995}}
\put(29.36,61.78){\line(-1,0){.995}}
\put(27.37,61.8){\line(-1,0){.995}}
\put(25.38,61.82){\line(-1,0){.995}}
\put(23.39,61.84){\line(-1,0){.995}}
\put(21.4,61.86){\line(-1,0){.995}}
\put(19.41,61.88){\line(-1,0){.995}}
\put(17.42,61.9){\line(-1,0){.995}}
\put(15.42,61.92){\line(-1,0){.995}}
%\end
\put(63.5,59.25){\line(0,-1){8.25}} 
\put(63.5,51){\line(0,-1){3.25}}
\put(63,126.5){\line(1,0){1}} 
\put(63,104.5){\line(1,0){1}}
\put(63.25,88.25){\line(1,0){.5}} 
\put(62.75,88.25){\line(1,0){1.5}}
%\emline(63,48.75)(64,49)
\multiput(63,48.75)(.125,.03125){8}{\line(1,0){.125}}
%\end
%\vector[both](14.75,88)(115.75,87.75)
\put(115.75,87.75){\vector(1,0){.07}}\put(14.75,88){\vector(-1,0){.07}}\multiput(14.75,88)(12.625,-.03125){8}{\line(1,0){12.625}}
%\end
\qbezier(15.5,115.25)(53,118)(63.5,126.75)
\qbezier(15,97.5)(46.38,98)(63.25,104.5)
\qbezier(63.5,104.25)(91.88,113.38)(113.75,114)
\qbezier(15,80.75)(44.25,81)(63.5,88.25)
\qbezier(63.5,88.25)(73.63,93.63)(115.25,96.5)
\qbezier(14.75,62.5)(44.88,63.38)(63.5,71.75)
\qbezier(63.5,71.75)(80,78)(115.5,79.25)
\qbezier(63.5,49)(79.75,60.88)(117,60.25)
%\emline(62.75,71.75)(64.25,72)
\multiput(62.75,71.75)(.1875,.03125){8}{\line(1,0){.1875}}
%\end
\put(66.5,127){\makebox(0,0)[cc]{$2\pi$}}
\put(65.75,115.75){\makebox(0,0)[cc]{$3\pi/2$}}
\put(67.5,104){\makebox(0,0)[cc]{$\pi$}}
\put(65.5,95.5){\makebox(0,0)[cc]{$\pi/2$}}
\put(67.75,88){\makebox(0,0)[cc]{$0$}}
\put(65.75,81.75){\makebox(0,0)[cc]{$-\pi/2$}}
\put(65.25,63){\makebox(0,0)[cc]{$-3\pi/2$}}
\put(67.75,48.75){\makebox(0,0)[cc]{$-2\pi$}}
\put(66,70.5){\makebox(0,0)[cc]{$-\pi$}}
\end{picture}
\end{center}

\vfill

\begin{enumerate}
  \item[\hspace{2cm}] The scale on the horizontal axis must show, or attempt to show, the final set of the tangent function, since the final set of the tangent function is the initial set of $\tan^{-1}$. The final set of the tangent function is the set of all real numbers. Do you see that?
  \item[\hspace{2cm}2)] Since the tangent function has vertical asymptotes at $-\frac{3\pi}{2}$, $-\frac{\pi}{2}$, $\frac{\pi}{2}$, and $\frac{3\pi}{2}$, the inverse will have horizontal asymptotes at these numbers. So, sketch in horizontal dotted lines at these numbers.
  \item[\hspace{2cm}3)] Do an analysis similar to what you did for the tangent function around its asymptotes and where it crossed the $x$-axis.
  \item[\hspace{2cm}4)] If you hold the graph of $\tan^{-1}$ up to the light and look through from the back side of the paper properly, you'll see a sketch of the tangent function.
  \item[\hspace{2cm}5)] You might contemplate the following picture to see if it relates in any way to the problem of plotting the inverse of a function.
\end{enumerate}

\begin{center}
%TeXCAD Picture [prob6.6hintc.pic]. Options:
%\grade{\on}
%\emlines{\off}
%\epic{\off}
%\beziermacro{\on}
%\reduce{\on}
%\snapping{\off}
%\quality{8.00}
%\graddiff{0.01}
%\snapasp{1}
%\zoom{4.0000}
\unitlength 1mm % = 2.85pt
\linethickness{0.4pt}
\ifx\plotpoint\undefined\newsavebox{\plotpoint}\fi % GNUPLOT compatibility
\begin{picture}(96,66)(0,0)
\put(7,13.75){\line(1,0){89}}
\put(96,13.75){\line(0,1){0}}
\put(28,3.25){\line(0,1){62.75}}
\put(28,66){\line(0,1){0}}
\put(28,13.75){\line(1,1){47}}
\put(75,60.75){\line(0,1){0}}
\put(28,13.75){\line(-1,-1){9.75}}
\put(18.25,4){\line(0,1){0}}
%\dashline{1}(41,26.75)(58.25,27)
\put(40.93,26.68){\line(1,0){.958}}
\put(42.85,26.71){\line(1,0){.958}}
\put(44.76,26.74){\line(1,0){.958}}
\put(46.68,26.76){\line(1,0){.958}}
\put(48.6,26.79){\line(1,0){.958}}
\put(50.51,26.82){\line(1,0){.958}}
\put(52.43,26.85){\line(1,0){.958}}
\put(54.35,26.87){\line(1,0){.958}}
\put(56.26,26.9){\line(1,0){.958}}
%\end
%\dashline{1}(58.25,27)(58.25,27)
\put(58.18,26.93){\line(0,1){0}}
%\end
%\dashline{1}(57.25,26.75)(57.25,43.5)
\put(57.18,26.68){\line(0,1){.985}}
\put(57.18,28.65){\line(0,1){.985}}
\put(57.18,30.62){\line(0,1){.985}}
\put(57.18,32.59){\line(0,1){.985}}
\put(57.18,34.56){\line(0,1){.985}}
\put(57.18,36.53){\line(0,1){.985}}
\put(57.18,38.5){\line(0,1){.985}}
\put(57.18,40.47){\line(0,1){.985}}
\put(57.18,42.44){\line(0,1){.985}}
%\end
%\dashline{1}(57.25,43.5)(41,43.5)
\put(57.18,43.43){\line(-1,0){.956}}
\put(55.27,43.43){\line(-1,0){.956}}
\put(53.36,43.43){\line(-1,0){.956}}
\put(51.44,43.43){\line(-1,0){.956}}
\put(49.53,43.43){\line(-1,0){.956}}
\put(47.62,43.43){\line(-1,0){.956}}
\put(45.71,43.43){\line(-1,0){.956}}
\put(43.8,43.43){\line(-1,0){.956}}
\put(41.89,43.43){\line(-1,0){.956}}
%\end
%\dashline{1}(41,43.5)(41,27.25)
\put(40.93,43.43){\line(0,-1){.956}}
\put(40.93,41.52){\line(0,-1){.956}}
\put(40.93,39.61){\line(0,-1){.956}}
\put(40.93,37.69){\line(0,-1){.956}}
\put(40.93,35.78){\line(0,-1){.956}}
\put(40.93,33.87){\line(0,-1){.956}}
\put(40.93,31.96){\line(0,-1){.956}}
\put(40.93,30.05){\line(0,-1){.956}}
\put(40.93,28.14){\line(0,-1){.956}}
%\end
%\vector(62.5,58.5)(67.75,54.5)
\put(67.75,54.5){\vector(4,-3){.07}}\multiput(62.5,58.5)(.04411765,-.03361345){119}{\line(1,0){.04411765}}
%\end
\put(67.75,54.5){\vector(0,1){0}}
\put(61.75,60){\makebox(0,0)[cc]{$L: y=x$}}
\put(37,42){\makebox(0,0)[cc]{$(u,v)$}}
\put(37,26.5){\makebox(0,0)[cc]{$(u,u)$}}
\put(61,26.5){\makebox(0,0)[cc]{$(v,u)$}}
\put(61,42){\makebox(0,0)[cc]{$(v,v)$}}
\end{picture}
\end{center}

Does anything just leap out at you about plotting inverses?


\vspace{2cm}

\begin{center}
Sound decision: Coal miners refuse to work after death.
\vspace{0.5cm}

Stay in school: Local high school dropouts cut in half.
\end{center}



\chapter*{Chapter VII}
\addcontentsline{toc}{chapter}{\numberline{}Chapter VII}
\markboth{Chapter VII}{Chapter VII}


Often, when the topic of trigonometry comes up, the discussion is about the functions cosine and sine and their applications to solving triangles. We should investigate what this all means.

First of all, let's dredge up some things from geometry to make the investigation more precise. Recall that, in geometry, we used the words ``set'', ``point'', ``line'', and ``between'' as undefined terms. All other terms had to be traceable back to one or more of these. Here is an example.

\begin{description}

\item[Definition 7.1:] The statement that $S$ is the segment $AB$ means:
  \begin{enumerate}[\hspace{1cm}1)]
    \item $S$ is a pointset,
    \item $A$ and $B$ are points,
    \item a point $P$ belongs to $S$ if and only if $APB$; i.e., $P$ is between $A$ and $B$.
  \end{enumerate}

\end{description}

Note that a segment is defined in terms of ``point'', ``set'', and ``between''.

\begin{description}

\item[Definition 7.2:] The statement that $I$ is the interval $AB$ means:
  \begin{enumerate}[\hspace{1cm}1)]
    \item $I$ is a pointset,
    \item $A$ and $B$ are points,
    \item a point $P$ belongs to $I$ if and only if
    \begin{enumerate}[\hspace{1cm}a)]
      \item $P$ is $A$ or $B$, or
      \item $P$ is a point of segment $AB$.
    \end{enumerate}
  \end{enumerate}

\end{description}

An interval is defined in terms of ``point'' and ``segment'', where ``segment'' is traceable back to the basic terms ``set'', ``point'', ``line'', and ``between''.

\begin{description}

\item[Problem 7.1:] Using the two definitions above as models, try to define the statement ``$T$ is the triangle $ABC$''.

\end{description}

Do you recall the notion of a ray from geometry? How would you define the statement ``$R$ is the ray $AB$''? Try this before going on, even if you've never seen the term.

\begin{description}

\item[Definition 7.3:] The statement that $R$ is the ray $AB$ means:
  \begin{enumerate}[\hspace{1cm}1)]
    \item $R$ is a pointset,
    \item $A$ and $B$ are points,
    \item a point $P$ belongs to $R$ if and only if
    \begin{enumerate}[\hspace{1cm}a)]
      \item $P$ is a point of line $AB$, and
      \item $PAB$ is not true.
    \end{enumerate}
  \end{enumerate}

\end{description}

Study this by drawing some pictures. You may have defined it somewhat differently, but that's all right as long as your definition defines the same pointset that the one above defines.

We used the notion of ``side of a line'' early in the course so you have undoubtedly ruminated about what this means. Notice how the definition of that concept also illustrates the idea that everything traces back to the terms ``set'', ``point'', ``line'', and ``between''.

\begin{description}

\item[Problem 7.2:] Can you use the notion of ``ray'' to define the statement ``$A$ is the angle $PQR$''?

\item[Problem 7.3:] Now define the statement ``$x$ is an interior point of angle $PQR$.''

\item[Problem 7.4:] What would you mean by saying ``$x$ is an interior point of triangle $PQR$''? 

\end{description}

Now we want to find some way of assigning a number to an angle so that we can compare angles. See if you can come up with a definition of measure of an angle without introducing any terms other than those we've already talked about. Try to finish this statement:

\begin{description}

\item[Problem 7.5:] The statement that $\theta$ is the measure of the angle $ABC$ means - - -.

\end{description}

Do not go on further until you have what you consider to be an acceptable definition of angle measure, or you have expended considerable effort trying to do so.

As you have probably seen from your efforts, there are several ways that one could define angle measure. The following is one which we will adopt, but keep your own in mind, and see as we go along if it is better than this one.

\begin{description}

\item[Definition 7.4:] The statement that $\theta$ is the measure of angle $ABC$ means:
  \begin{enumerate}[\hspace{1cm}1)]
    \item $\theta$ is a number, and
    \item If $A_{1}$ is a point of ray $BA$ and $C_{1}$ is a point of ray $BC$, such that $BA_{1} = BC_{1} = 1$, then $\theta = l(\arc{A_{1}C_{1}})$.
  \end{enumerate}

\end{description}

Note that we define angle measure in terms of arc length on a unit circle.

\textbf{Notation}: The measure of angle $ABC$ is denoted $m \angm ABC$. Hence, if $\theta$ is the measure of angle $ABC$, symbolically we state this as $\theta = m \angm ABC$.

% p205 figure added

\begin{center}
%TeXCAD Picture [p205.pic]. Options:
%\grade{\on}
%\emlines{\off}
%\epic{\off}
%\beziermacro{\on}
%\reduce{\on}
%\snapping{\on}
%\quality{8.00}
%\graddiff{0.01}
%\snapasp{1}
%\zoom{1.0000}
\unitlength .25mm % = .71pt
\linethickness{0.4pt}
\ifx\plotpoint\undefined\newsavebox{\plotpoint}\fi % GNUPLOT compatibility
\begin{picture}(250,200)(150,150)
\put(166,197){\line(1,0){308}}
%\emline(166,197)(416,331)
\multiput(166,197)(.2238137869,.1199641898){1117}{\line(1,0){.2238137869}}
%\end
\qbezier(253,292)(387.5,278.5)(348,171)
\put(259,218){\makebox(0,0)[cc]{$\theta$}}
\put(149,189){\makebox(0,0)[cc]{$B$}}
\put(255,168){\makebox(0,0)[cc]{$1$}}
\put(315,293){\makebox(0,0)[cc]{$A_1$}}
\put(377,179){\makebox(0,0)[cc]{$C_1$}}
\put(438,180){\makebox(0,0)[cc]{$C$}}
\put(377,322){\makebox(0,0)[cc]{$A$}}
\put(166,197){\circle*{5.66}}
\put(355,197){\circle*{4.47}}
\put(437,197){\circle*{4.47}}
\put(382,312){\circle*{5.66}}
\put(314,277){\circle*{4.47}}
%\vector[both](169,186)(342,186)
\put(342,186){\vector(1,0){.28}}\put(169,186){\vector(-1,0){.28}}\put(169,186){\line(1,0){173}}
%\end
\end{picture}
\end{center}


In the picture, $\theta = m \angm ABC = l(\arc{A_{1}C_{1}})$.

Note that, by virtue of our definition of angle measure, if $\theta$ is the measure of an angle, then $0 \leq \theta \leq \pi$.

An angle of measure $0$ consists of just one ray.

An angle of measure $\pi$ consists of two rays whose union is a line. It is called a straight angle.

An angle of measure $\frac{\pi}{2}$ also has a special name. It is called a right angle. The picture below illustrates how such an angle is represented in the literature.

\begin{center}
%TeXCAD Options
%\grade{\on}
%\emlines{\off}
%\epic{\off}
%\beziermacro{\on}
%\reduce{\on}
%\snapping{\off}
%\quality{8.00}
%\graddiff{0.01}
%\snapasp{1}
%\zoom{4.0000}
\unitlength 1mm % = 2.85pt
\linethickness{0.4pt}
\ifx\plotpoint\undefined\newsavebox{\plotpoint}\fi % GNUPLOT compatibility
\begin{picture}(59.75,50.75)(0,85)
\put(22,135){\line(0,-1){40.5}} 
\put(22,94.5){\line(1,0){37}}
\put(22.25,99.5){\line(1,0){4.25}}
\put(26.5,99.5){\line(0,-1){4.75}} 
\put(22.25,94.25){\circle{.71}}
\put(58.75,94.5){\circle{.5}} 
\put(22,134.75){\circle{.5}}
\put(18,135.75){\makebox(0,0)[cc]{$A$}}
\put(21,90){\makebox(0,0)[cc]{$B$}}
\put(59.75,91.5){\makebox(0,0)[cc]{$C$}}
\end{picture}
\end{center}

$M \angm ABC = \frac{\pi}{2}$, and the $\square$ indicates this; i.e., $ABC$ is a right angle.

\vspace{0.3cm}

\textbf{Problem 7.6}: Consider the following definition of measure of an angle.

The statement that $\theta$ is the measure of angle $ABC$ means:
\begin{enumerate}[\hspace{1cm}1)]
  \item $\theta $ is a number, and
  \item If $A_{1}$ is a point of ray $BA$, and $C_{1}$ is a point of ray $BC$ such that $BA_{1} = BC_{1} = 1$, then $\theta = A_{1}C_{1}$.
\end{enumerate}

Do you think this notion of angle measure is as useful to us as the one given earlier? If you think it is not as good, what is your objection to it?

It turns out that many of the problems we encounter, for which trigonometry gives us useful tools, involve triangles that are called ``right triangles''; i.e., triangles in which one of the angles is a right angle. Let's see how we can use the cosine and sine functions to ``solve'' right triangles.

Consider the following right triangle; i.e., a triangle with a right angle.

\begin{center}
%TeXCAD Options
%\grade{\on}
%\emlines{\off}
%\epic{\off}
%\beziermacro{\on}
%\reduce{\on}
%\snapping{\off}
%\quality{8.00}
%\graddiff{0.01}
%\snapasp{1}
%\zoom{4.0000}
\unitlength 1mm % = 2.85pt
\linethickness{0.4pt}
\ifx\plotpoint\undefined\newsavebox{\plotpoint}\fi % GNUPLOT compatibility
\begin{picture}(66.25,49.75)(0,70)
\put(62.25,118.5){\line(0,-1){38.25}}
\put(62.25,80.25){\line(-1,0){45.75}}
%\emline(16.5,80.25)(62.25,118.75)
\multiput(16.5,80.25)(.040061296,.0337127846){1142}{\line(1,0){.040061296}}
%\end
\put(58,84.25){\line(0,-1){3.75}}
\put(58.25,84.25){\line(1,0){3.75}}
\put(66.25,119.75){\makebox(0,0)[cc]{$C$}}
\put(65.75,80.75){\makebox(0,0)[cc]{$B$}}
\put(15.5,77){\makebox(0,0)[cc]{$A$}}
\put(36.25,74){\makebox(0,0)[cc]{Figure 7.1}}
\end{picture}
\end{center}

Note that angle $ABC$ is depicted as a right angle.

Problems involving such triangles usually involve finding the lengths of the sides or the measures of the angles. Let's redraw the triangle and label it more conveniently.

\begin{center}
%TeXCAD Options
%\grade{\on}
%\emlines{\off}
%\epic{\off}
%\beziermacro{\on}
%\reduce{\on}
%\snapping{\off}
%\quality{8.00}
%\graddiff{0.01}
%\snapasp{1}
%\zoom{4.0000}
\unitlength 1mm % = 2.85pt
\linethickness{0.4pt}
\ifx\plotpoint\undefined\newsavebox{\plotpoint}\fi % GNUPLOT compatibility
\begin{picture}(66.25,49.75)(0,70)
\put(62.25,118.5){\line(0,-1){38.25}}
\put(62.25,80.25){\line(-1,0){45.75}}
%\emline(16.5,80.25)(62.25,118.75)
\multiput(16.5,80.25)(.040061296,.0337127846){1142}{\line(1,0){.040061296}}
%\end
\put(58,84.25){\line(0,-1){3.75}}
\put(58.25,84.25){\line(1,0){3.75}}
\put(66.25,119.75){\makebox(0,0)[cc]{$C$}}
\put(65.75,80.75){\makebox(0,0)[cc]{$B$}}
\put(15.5,77){\makebox(0,0)[cc]{$A$}}
\put(22.5,82){\makebox(0,0)[cc]{$\theta$}}
\put(66.25,101.5){\makebox(0,0)[cc]{$b$}}
\put(36.5,102){\makebox(0,0)[cc]{$c$}}
\put(40,68.5){\makebox(0,0)[cc]{Figure 7.2}}
\put(40,76.75){\makebox(0,0)[cc]{$a$}}
\end{picture}
\end{center}

The question is: What relationships exist among the lengths of the sides ($a$, $b$, and $c$) and the measure of the angle, $\theta$? $\theta = m \angm CAB$.

Let's impose a coordinate system on our picture, with origin at $A$ and $x$-axis containing the interval $AB$.

% Figures 7.3 and 7.4

\begin{center}
%TeXCAD Options
%\grade{\on}
%\emlines{\off}
%\epic{\off}
%\beziermacro{\on}
%\reduce{\on}
%\snapping{\off}
%\quality{8.00}
%\graddiff{0.01}
%\snapasp{1}
%\zoom{4.0000}
\unitlength 0.9mm % = 2.85pt
\linethickness{0.4pt}
\ifx\plotpoint\undefined\newsavebox{\plotpoint}\fi % GNUPLOT compatibility
\begin{picture}(121.5,140.5)(0,3)
\put(62.25,95.75){\line(-1,0){45.75}}
\put(58,99.75){\line(0,-1){3.75}}
\put(58.25,99.75){\line(1,0){3.75}}
\put(66.25,117){\makebox(0,0)[cc]{$b$}}
\put(40,92.25){\makebox(0,0)[cc]{$a$}}
%\emline(16.5,95.75)(62.75,124.25)
\multiput(16.5,95.75)(.0547337278,.0337278107){845}{\line(1,0){.0547337278}}
%\end
\put(62.75,124.25){\line(0,-1){28.25}}
%\emline(62.75,96)(62.5,95.75)
\multiput(62.75,96)(-.03125,-.03125){8}{\line(0,-1){.03125}}
%\end
\put(62.5,95.75){\line(1,0){.25}}
\put(16.25,148.5){\line(0,-1){66.5}}
\put(62.75,95.75){\line(1,0){58.75}} 
\put(16,95.75){\line(-1,0){11}}
\put(12.5,142.75){\makebox(0,0)[cc]{$Y$}}
\put(16.5,93){\makebox(0,0)[cc]{$(0,0)$}}
\put(63.75,92.5){\makebox(0,0)[cc]{$(a,0)$}}
\put(96.75,93){\makebox(0,0)[cc]{$X$}}
\put(35.5,110.5){\makebox(0,0)[cc]{$c$}}
\put(63.5,126.75){\makebox(0,0)[cc]{$(?,?)$}}
\put(23.5,97.25){\makebox(0,0)[cc]{$\theta$}}
\put(86.25,139.75){\makebox(0,0)[cc]{$A = O:(0,0)$}}
\put(86,130){\makebox(0,0)[cc]{$B:(a,0)$}}
\put(84,123){\makebox(0,0)[cc]{$C:(?,?)$}}
\put(55.25,81.75){\makebox(0,0)[cc]{Figure 7.3}}
\put(60.75,19){\line(-1,0){45.75}} 
\put(56.5,23){\line(0,-1){3.75}}
\put(56.75,23){\line(1,0){3.75}}
\put(64.75,40.25){\makebox(0,0)[cc]{$b$}}
\put(38.5,15.5){\makebox(0,0)[cc]{$a$}}
%\emline(15,19)(61.25,47.5)
\multiput(15,19)(.0547337278,.0337278107){845}{\line(1,0){.0547337278}}
%\end
\put(61.25,47.5){\line(0,-1){28.25}}
%\emline(61.25,19.25)(61,19)
\multiput(61.25,19.25)(-.03125,-.03125){8}{\line(0,-1){.03125}}
%\end
\put(61,19){\line(1,0){.25}} 
\put(14.75,71.75){\line(0,-1){66.5}}
\put(61.25,19){\line(1,0){58.75}} 
\put(14.5,19){\line(-1,0){11}}
\put(11,66){\makebox(0,0)[cc]{$Y$}}
\put(15,16.25){\makebox(0,0)[cc]{$(0,0)$}}
\put(62.25,15.75){\makebox(0,0)[cc]{$(a,0)$}}
\put(95.25,16.25){\makebox(0,0)[cc]{$X$}}
\put(34,33.75){\makebox(0,0)[cc]{$c$}}
\put(22,20.5){\makebox(0,0)[cc]{$\theta$}}
\put(61.25,49.75){\makebox(0,0)[cc]{$(a,b)$}}
\put(80.75,46.25){\makebox(0,0)[cc]{Note that $C:(a,b)$}}
\put(54.75,8.25){\makebox(0,0)[cc]{Figure 7.4}}
\end{picture}
\end{center}

But look back at Corollary 5.3. This tells us that $a = c \cos \theta$ and $b = c \sin \theta$.

\begin{center}
%TeXCAD Options
%\grade{\on}
%\emlines{\off}
%\epic{\off}
%\beziermacro{\on}
%\reduce{\on}
%\snapping{\off}
%\quality{8.00}
%\graddiff{0.01}
%\snapasp{1}
%\zoom{4.0000}
\unitlength 0.9mm % = 2.85pt
\linethickness{0.4pt}
\ifx\plotpoint\undefined\newsavebox{\plotpoint}\fi % GNUPLOT compatibility
\begin{picture}(121.5,75)(0,75)
\put(62.25,95.75){\line(-1,0){45.75}}
\put(58,99.75){\line(0,-1){3.75}}
\put(58.25,99.75){\line(1,0){3.75}}
\put(66.25,117){\makebox(0,0)[cc]{b}}
\put(40,92.25){\makebox(0,0)[cc]{a}}
%\emline(16.5,95.75)(62.75,124.25)
\multiput(16.5,95.75)(.0547337278,.0337278107){845}{\line(1,0){.0547337278}}
%\end
\put(62.75,124.25){\line(0,-1){28.25}}
%\emline(62.75,96)(62.5,95.75)
\multiput(62.75,96)(-.03125,-.03125){8}{\line(0,-1){.03125}}
%\end
\put(62.5,95.75){\line(1,0){.25}}
\put(16.25,148.5){\line(0,-1){66.5}}
\put(62.75,95.75){\line(1,0){58.75}} 
\put(16,95.75){\line(-1,0){11}}
\put(12.5,142.75){\makebox(0,0)[cc]{$Y$}}
\put(16.5,93){\makebox(0,0)[cc]{$(0,0)$}}
\put(63.75,92.5){\makebox(0,0)[cc]{$(a,0)$}}
\put(96.75,93){\makebox(0,0)[cc]{$X$}}
\put(35.5,110.5){\makebox(0,0)[cc]{$c$}}
\put(23.5,97.25){\makebox(0,0)[cc]{$\theta$}}
\put(62.25,128){\makebox(0,0)[cc]{$(c \cos \theta , c \sin \theta)$}}
\put(59.5,83.75){\makebox(0,0)[cc]{Figure 7.5}}
\end{picture}
\end{center}

Extracting the picture from the coordinate system, we have:

\begin{center}
%TeXCAD Options
%\grade{\on}
%\emlines{\off}
%\epic{\off}
%\beziermacro{\on}
%\reduce{\on}
%\snapping{\off}
%\quality{8.00}
%\graddiff{0.01}
%\snapasp{1}
%\zoom{4.0000}
\unitlength 1mm % = 2.85pt
\linethickness{0.4pt}
\ifx\plotpoint\undefined\newsavebox{\plotpoint}\fi % GNUPLOT compatibility
\begin{picture}(67.75,48.75)(0,68)
\put(62.25,118.5){\line(0,-1){38.25}}
\put(62.25,80.25){\line(-1,0){45.75}}
%\emline(16.5,80.25)(62.25,118.75)
\multiput(16.5,80.25)(.040061296,.0337127846){1142}{\line(1,0){.040061296}}
%\end
\put(58,84.25){\line(0,-1){3.75}}
\put(58.25,84.25){\line(1,0){3.75}}
\put(73.75,100.75){\makebox(0,0)[cc]{$b = c \sin \theta$}}
\put(38,77){\makebox(0,0)[cc]{$a = c \cos \theta$}}
\put(33,100.5){\makebox(0,0)[cc]{$c$}}
\put(24.5,83.25){\makebox(0,0)[cc]{$\theta$}}
\put(38.25,69.5){\makebox(0,0)[cc]{Figure 7.6}}
\end{picture}
\end{center}

Hence, $\cos \theta = \frac{a}{c}$ and $\sin \theta = \frac{b}{c}$.

If you look in most any trigonometry text, this is where they start. Isn't this great? After only 6 chapters, we finally got up to page 1 of a standard textbook! It always seems to me that, in the standard textbooks, it was tacitly (look this one up) assumed that the student already had a large amount of knowledge of things: For example, a right triangle would be given with sides $1$, $2$, and $\sqrt{3}$. Is it obvious to you that one of the angles in that triangle has measure $\frac{\pi}{6}$? Experience can make some things obvious that were utterly obscure when you were a beginner. So if you are one of the unfortunate few in the world not born with a lot of experience with trigonometry, this textbook is for you.

In a right triangle, we have three sides and three angles, one of which is a right angle. Consider the following: The sides have lengths $a$, $b$, and $c$, and the angles have measures $\theta$, $\phi$, and $\frac{\pi}{2}$.

\begin{center}
%TeXCAD Options
%\grade{\on}
%\emlines{\off}
%\epic{\off}
%\beziermacro{\on}
%\reduce{\on}
%\snapping{\off}
%\quality{8.00}
%\graddiff{0.01}
%\snapasp{1}
%\zoom{4.0000}
\unitlength 1mm % = 2.85pt
\linethickness{0.4pt}
\ifx\plotpoint\undefined\newsavebox{\plotpoint}\fi % GNUPLOT compatibility
\begin{picture}(65.25,53.75)(0,65)
\put(62.25,118.5){\line(0,-1){38.25}}
\put(62.25,80.25){\line(-1,0){45.75}}
%\emline(16.5,80.25)(62.25,118.75)
\multiput(16.5,80.25)(.040061296,.0337127846){1142}{\line(1,0){.040061296}}
%\end
\put(58,84.25){\line(0,-1){3.75}}
\put(58.25,84.25){\line(1,0){3.75}}
\put(33,100.5){\makebox(0,0)[cc]{$c$}}
\put(23,82.5){\makebox(0,0)[cc]{$\theta$}}
\put(59,111.5){\makebox(0,0)[cc]{$\phi$}}
\put(65.25,98){\makebox(0,0)[cc]{$b$}}
\put(38.5,77.5){\makebox(0,0)[cc]{$a$}}
\put(38.5,69.75){\makebox(0,0)[cc]{Figure 7.7}}
\end{picture}
\end{center}

\begin{description}

\item[Problem 7.7:]
  \begin{enumerate}[\hspace{1cm}a)]
    \item Suppose you know $a$ and $b$. Can you find $c$, $\theta$, and $\phi$?
    \item Suppose you know $a$ and $c$. Can you find $b$, $\theta$, and $\phi$?
    \item Suppose you know $c$ and $\theta$. Can you find $a$, $b$, and $\phi$?
    \item Suppose you know $c$ and $\phi$. Can you find $a$, $b$, and $\theta$?
    \item Suppose you know $a$ and $\theta$. Can you find $b$, $c$, and $\phi$?
    \item Suppose you know $b$ and $\theta$. Can you find $a$, $c$, and $\phi$?
  \end{enumerate}

\end{description}

By now, you should be getting some idea of what might be meant by solving a right triangle. Let's relate all this to the other functions that we defined in terms of the sine and cosine functions.

\begin{description}

\item[Problem 7.8:] From Figure 7.8, find each of the following in terms of the lengths of the sides of the triangle.

\end{description}

\begin{center}
%TeXCAD Options
%\grade{\on}
%\emlines{\off}
%\epic{\off}
%\beziermacro{\on}
%\reduce{\on}
%\snapping{\off}
%\quality{8.00}
%\graddiff{0.01}
%\snapasp{1}
%\zoom{4.0000}
\unitlength 1mm % = 2.85pt
\linethickness{0.4pt}
\ifx\plotpoint\undefined\newsavebox{\plotpoint}\fi % GNUPLOT compatibility
\begin{picture}(65,48.75)(0,70)
\put(62.25,118.5){\line(0,-1){38.25}}
\put(62.25,80.25){\line(-1,0){45.75}}
%\emline(16.5,80.25)(62.25,118.75)
\multiput(16.5,80.25)(.040061296,.0337127846){1142}{\line(1,0){.040061296}}
%\end
\put(23,82.5){\makebox(0,0)[cc]{$\theta$}}
\put(59,111.5){\makebox(0,0)[cc]{$\phi$}}
\put(33.75,100.25){\makebox(0,0)[cc]{$r$}}
\put(40.25,77){\makebox(0,0)[cc]{$x$}}
\put(65,96.75){\makebox(0,0)[cc]{$y$}}
\put(39.75,71.25){\makebox(0,0)[cc]{Figure 7.8}}
\end{picture}
\end{center}

\begin{align*}
  \sin \theta & =   & \csc \theta & =   & \sin \phi & =   & \csc \phi & = \\
  \cos \theta & =   & \sec \theta & =   & \cos \phi & =   & \sec \phi & = \\
  \tan \theta & =   & \cot \theta & =   & \tan \phi & =   & \cot \phi & = 
\end{align*}

\textbf{Problem 7.9}: Now, what are we going to do about solving triangles that are \emph{not} right triangles? Consider the following.

\begin{center}
%TeXCAD Options
%\grade{\on}
%\emlines{\off}
%\epic{\off}
%\beziermacro{\on}
%\reduce{\on}
%\snapping{\off}
%\quality{8.00}
%\graddiff{0.01}
%\snapasp{1}
%\zoom{4.0000}
\unitlength 1mm % = 2.85pt
\linethickness{0.4pt}
\ifx\plotpoint\undefined\newsavebox{\plotpoint}\fi % GNUPLOT compatibility
\begin{picture}(69.25,49.5)(0,81)
\put(26.5,91.75){\line(1,0){24.25}}
%\emline(50.75,91.75)(69.25,130.5)
\multiput(50.75,91.75)(.033697632,.070582878){549}{\line(0,1){.070582878}}
%\end
%\emline(69.25,130.5)(26.75,92)
\multiput(69.25,130.5)(-.0372154116,-.0337127846){1142}{\line(-1,0){.0372154116}}
%\end
\put(32.25,94){\makebox(0,0)[cc]{$\theta$}}
\put(41,108.5){\makebox(0,0)[cc]{$c$}}
\put(60,104.25){\makebox(0,0)[cc]{$b$}}
\put(39.25,88.25){\makebox(0,0)[cc]{$a$}}
\put(38.25,82.75){\makebox(0,0)[cc]{Figure 7.9}}
\end{picture}
\end{center}


We cannot say that $\sin \theta = \frac{b}{c}$ and $\cos \theta = \frac{a}{c}$, but surely we can figure something out. Give it a serious try before going on.

Once before, when attacking a problem, we imposed a coordinate system on our picture and it worked to our advantage. Let's try that technique again.

\begin{description}

\item[Problem 7.10:] $\ $

\end{description}

%TeXCAD Picture [fig7.10.pic]. Options:
%\grade{\on}
%\emlines{\off}
%\epic{\off}
%\beziermacro{\on}
%\reduce{\on}
%\snapping{\off}
%\quality{8.00}
%\graddiff{0.01}
%\snapasp{1}
%\zoom{4.0000}
\unitlength 0.9mm % = 2.85pt
\linethickness{0.4pt}
\ifx\plotpoint\undefined\newsavebox{\plotpoint}\fi % GNUPLOT compatibility
\begin{picture}(103.38,68.5)(0,73)
\put(24,83.75){\line(1,0){33.25}}
%\emline(57.25,83.75)(85,135.25)
\multiput(57.25,83.75)(.0337181045,.0625759417){823}{\line(0,1){.0625759417}}
%\end
%\emline(85,135.25)(24,84)
\multiput(85,135.25)(-.0401579987,-.0337393022){1519}{\line(-1,0){.0401579987}}
%\end
\put(45.75,106.5){\makebox(0,0)[cc]{$c$}}
\put(69.5,99.75){\makebox(0,0)[cc]{$b$}}
\put(42,80.75){\makebox(0,0)[cc]{$a$}}
\put(34.75,86.5){\makebox(0,0)[cc]{$\theta$}}
\put(57.5,81.25){\makebox(0,0)[cc]{$(a,0)$}}
\put(22,81.25){\makebox(0,0)[cc]{$O:(0,0)$}}
%\qbezvec(103,127.75)(103.38,141.5)(87.25,136.25)
\put(87.25,136.25){\vector(-3,-1){.07}}\qbezier(103,127.75)(103.38,141.5)(87.25,136.25)
%\end
\put(103.25,124){\makebox(0,0)[cc]{what are the coordinates}}
\put(103,120.25){\makebox(0,0)[cc]{of this point?}}
\put(41.25,74.75){\makebox(0,0)[cc]{Figure 7.10}}
\end{picture}

Find the coordinates, and if you can, show that $b^{2} =a^{2} + c^{2} - 2ac \cos \theta$.

\begin{description}

\item[Problem 7.11:] Consider the following picture. Show that:
  \begin{enumerate}[\hspace{1cm}1)]
    \item $a^{2} = b^{2} + c^{2} - 2bc \cos A$
    \item $b^{2} = a^{2} + c^{2} - 2ac \cos B$
    \item $c^{2} = a^{2} + b^{2} - 2ab \cos C$
  \end{enumerate}

\end{description}

\begin{center}
%TeXCAD Options
%\grade{\on}
%\emlines{\off}
%\epic{\off}
%\beziermacro{\on}
%\reduce{\on}
%\snapping{\off}
%\quality{8.00}
%\graddiff{0.01}
%\snapasp{1}
%\zoom{4.0000}
\unitlength 1mm % = 2.85pt
\linethickness{0.4pt}
\ifx\plotpoint\undefined\newsavebox{\plotpoint}\fi % GNUPLOT compatibility
\begin{picture}(66.5,49.75)(0,69)
\put(62.25,118.5){\line(0,-1){38.25}}
\put(62.25,80.25){\line(-1,0){45.75}}
%\emline(16.5,80.25)(62.25,118.75)
\multiput(16.5,80.25)(.040061296,.0337127846){1142}{\line(1,0){.040061296}}
%\end
\put(60,113.75){\makebox(0,0)[cc]{$B$}}
\put(22.25,82.25){\makebox(0,0)[cc]{$A$}}
\put(59.25,82.25){\makebox(0,0)[cc]{$C$}}
\put(39.5,77){\makebox(0,0)[cc]{$b$}}
\put(66.5,97.25){\makebox(0,0)[cc]{$a$}}
\put(35.5,101.25){\makebox(0,0)[cc]{$c$}}
\put(39.75,70){\makebox(0,0)[cc]{Figure 7.11}}
\end{picture}
\end{center}

Note that this looks somewhat like the Pythagorean Theorem, with the last term as some adjustment for the fact that the triangle is not a right triangle.

Are the formulas above valid if the triangle \emph{is} a right triangle?

\begin{description}

\item[Problem 7.12:] The formula you just derived involves the cosine function and not the sine function. In fact, those formulas go by the name ``The Law of Cosines''. I know what you're thinking: Is there a Law of Sines? You weren't thinking that? Well, believe it or not, there is one, and you are asked to show that it is valid. 

  Use the following picture as an aid, and show that $\frac{a}{\sin A} = \frac{b}{\sin B} = \frac{c}{\sin C}$. This relationship is known as the Law of Sines.

\end{description}

%TeXCAD Picture [fig7.12.pic]. Options:
%\grade{\on}
%\emlines{\off}
%\epic{\off}
%\beziermacro{\on}
%\reduce{\on}
%\snapping{\off}
%\quality{8.00}
%\graddiff{0.01}
%\snapasp{1}
%\zoom{4.0000}
\unitlength 1mm % = 2.85pt
\linethickness{0.4pt}
\ifx\plotpoint\undefined\newsavebox{\plotpoint}\fi % GNUPLOT compatibility
\begin{picture}(72.25,42.75)(0,86)
\put(18.75,94.75){\line(1,0){50}}
\put(68.75,94.75){\line(0,1){34}}
%\emline(68.75,128.75)(19,95.25)
\multiput(68.75,128.75)(-.0501007049,-.0337361531){993}{\line(-1,0){.0501007049}}
%\end
%\emline(19,95.25)(18.25,94.75)
\multiput(19,95.25)(-.05,-.033333){15}{\line(-1,0){.05}}
%\end
\put(48.5,95){\line(3,5){20.25}}
\put(65,94.75){\line(0,1){3.75}}
\put(65,98.5){\line(1,0){3.75}}
\put(36.75,93.25){\vector(1,0){10}}
\put(29.25,92.75){\vector(-1,0){11.5}}
\put(32.75,92.5){\makebox(0,0)[cc]{$b$}}
\put(55.25,97){\makebox(0,0)[cc]{$\pi - C$}}
\put(45.5,97.75){\makebox(0,0)[cc]{$C$}}
\put(62,121.25){\makebox(0,0)[cc]{$B$}}
\put(25,96.75){\makebox(0,0)[cc]{$A$}}
\put(38,112.5){\makebox(0,0)[cc]{$c$}}
\put(72.25,110.5){\makebox(0,0)[cc]{$h$}}
\put(61,110.5){\makebox(0,0)[cc]{$a$}}
\put(44.5,87.5){\makebox(0,0)[cc]{Figure 7.12}}
\end{picture}

Let's recall our original problem, and see how far we can go with it with our present arsenal of tools. We are given a circle of radius $r>0$ and two points $P$ and $Q$ of the circle. We want to find a relationship between the chord length $PQ$ and the arc length $l(\arc{PQ})$. We lose no generality by imposing a coordinate system with origin at the center of our circle.

\begin{center}
%TeXCAD Options
%\grade{\on}
%\emlines{\off}
%\epic{\off}
%\beziermacro{\on}
%\reduce{\on}
%\snapping{\off}
%\quality{8.00}
%\graddiff{0.01}
%\snapasp{1}
%\zoom{4.0000}
\unitlength 0.9mm % = 2.85pt
\linethickness{0.4pt}
\ifx\plotpoint\undefined\newsavebox{\plotpoint}\fi % GNUPLOT compatibility
\begin{picture}(110,102.75)(0,34)
%\circle(58,90){67.54}
\put(91.77,90){\line(0,1){1.297}}
\put(91.74,91.3){\line(0,1){1.295}}
\multiput(91.67,92.59)(-.0311,.3228){4}{\line(0,1){.3228}}
\multiput(91.54,93.88)(-.02899,.21427){6}{\line(0,1){.21427}}
\multiput(91.37,95.17)(-.03188,.18257){7}{\line(0,1){.18257}}
\multiput(91.15,96.45)(-.03023,.14094){9}{\line(0,1){.14094}}
\multiput(90.88,97.72)(-.03206,.12571){10}{\line(0,1){.12571}}
\multiput(90.55,98.97)(-.03352,.11308){11}{\line(0,1){.11308}}
\multiput(90.19,100.22)(-.032014,.09452){13}{\line(0,1){.09452}}
\multiput(89.77,101.45)(-.033077,.086562){14}{\line(0,1){.086562}}
\multiput(89.31,102.66)(-.03183,.074574){16}{\line(0,1){.074574}}
\multiput(88.8,103.85)(-.032632,.068985){17}{\line(0,1){.068985}}
\multiput(88.24,105.02)(-.033299,.06392){18}{\line(0,1){.06392}}
\multiput(87.64,106.17)(-.032156,.056335){20}{\line(0,1){.056335}}
\multiput(87,107.3)(-.032663,.052436){21}{\line(0,1){.052436}}
\multiput(86.31,108.4)(-.033078,.048818){22}{\line(0,1){.048818}}
\multiput(85.59,109.48)(-.03341,.045446){23}{\line(0,1){.045446}}
\multiput(84.82,110.52)(-.033667,.04229){24}{\line(0,1){.04229}}
\multiput(84.01,111.54)(-.032554,.037815){26}{\line(0,1){.037815}}
\multiput(83.16,112.52)(-.032724,.035183){27}{\line(0,1){.035183}}
\multiput(82.28,113.47)(-.032835,.032689){28}{\line(-1,0){.032835}}
\multiput(81.36,114.38)(-.035328,.032567){27}{\line(-1,0){.035328}}
\multiput(80.41,115.26)(-.039477,.033681){25}{\line(-1,0){.039477}}
\multiput(79.42,116.11)(-.04244,.033479){24}{\line(-1,0){.04244}}
\multiput(78.4,116.91)(-.045594,.033208){23}{\line(-1,0){.045594}}
\multiput(77.35,117.67)(-.048965,.032861){22}{\line(-1,0){.048965}}
\multiput(76.28,118.4)(-.052581,.03243){21}{\line(-1,0){.052581}}
\multiput(75.17,119.08)(-.05945,.033585){19}{\line(-1,0){.05945}}
\multiput(74.04,119.71)(-.064068,.033014){18}{\line(-1,0){.064068}}
\multiput(72.89,120.31)(-.069129,.032324){17}{\line(-1,0){.069129}}
\multiput(71.71,120.86)(-.079696,.033598){15}{\line(-1,0){.079696}}
\multiput(70.52,121.36)(-.086708,.032691){14}{\line(-1,0){.086708}}
\multiput(69.3,121.82)(-.094661,.031593){13}{\line(-1,0){.094661}}
\multiput(68.07,122.23)(-.11322,.03301){11}{\line(-1,0){.11322}}
\multiput(66.83,122.59)(-.12585,.0315){10}{\line(-1,0){.12585}}
\multiput(65.57,122.91)(-.15871,.03331){8}{\line(-1,0){.15871}}
\multiput(64.3,123.18)(-.18271,.03107){7}{\line(-1,0){.18271}}
\multiput(63.02,123.39)(-.25727,.03364){5}{\line(-1,0){.25727}}
\put(61.73,123.56){\line(-1,0){1.292}}
\put(60.44,123.68){\line(-1,0){1.295}}
\put(59.15,123.75){\line(-1,0){1.297}}
\put(57.85,123.77){\line(-1,0){1.297}}
\put(56.55,123.74){\line(-1,0){1.295}}
\multiput(55.26,123.66)(-.3227,-.0325){4}{\line(-1,0){.3227}}
\multiput(53.97,123.53)(-.21414,-.02994){6}{\line(-1,0){.21414}}
\multiput(52.68,123.35)(-.18243,-.03269){7}{\line(-1,0){.18243}}
\multiput(51.41,123.12)(-.14081,-.03086){9}{\line(-1,0){.14081}}
\multiput(50.14,122.84)(-.12556,-.03262){10}{\line(-1,0){.12556}}
\multiput(48.88,122.51)(-.10352,-.03118){12}{\line(-1,0){.10352}}
\multiput(47.64,122.14)(-.094377,-.032434){13}{\line(-1,0){.094377}}
\multiput(46.41,121.72)(-.086414,-.033461){14}{\line(-1,0){.086414}}
\multiput(45.2,121.25)(-.074432,-.032161){16}{\line(-1,0){.074432}}
\multiput(44.01,120.74)(-.068839,-.032938){17}{\line(-1,0){.068839}}
\multiput(42.84,120.18)(-.063772,-.033582){18}{\line(-1,0){.063772}}
\multiput(41.69,119.57)(-.056191,-.032406){20}{\line(-1,0){.056191}}
\multiput(40.57,118.92)(-.05229,-.032896){21}{\line(-1,0){.05229}}
\multiput(39.47,118.23)(-.048671,-.033295){22}{\line(-1,0){.048671}}
\multiput(38.4,117.5)(-.045297,-.033612){23}{\line(-1,0){.045297}}
\multiput(37.36,116.73)(-.040455,-.032501){25}{\line(-1,0){.040455}}
\multiput(36.35,115.91)(-.03767,-.032722){26}{\line(-1,0){.03767}}
\multiput(35.37,115.06)(-.035037,-.03288){27}{\line(-1,0){.035037}}
\multiput(34.42,114.18)(-.032543,-.03298){28}{\line(0,-1){.03298}}
\multiput(33.51,113.25)(-.033656,-.036837){26}{\line(0,-1){.036837}}
\multiput(32.64,112.29)(-.033505,-.039627){25}{\line(0,-1){.039627}}
\multiput(31.8,111.3)(-.03329,-.042588){24}{\line(0,-1){.042588}}
\multiput(31,110.28)(-.033005,-.045741){23}{\line(0,-1){.045741}}
\multiput(30.24,109.23)(-.032643,-.04911){22}{\line(0,-1){.04911}}
\multiput(29.52,108.15)(-.032196,-.052725){21}{\line(0,-1){.052725}}
\multiput(28.85,107.04)(-.03332,-.059598){19}{\line(0,-1){.059598}}
\multiput(28.21,105.91)(-.032729,-.064214){18}{\line(0,-1){.064214}}
\multiput(27.62,104.75)(-.032017,-.069272){17}{\line(0,-1){.069272}}
\multiput(27.08,103.58)(-.033243,-.079844){15}{\line(0,-1){.079844}}
\multiput(26.58,102.38)(-.032306,-.086853){14}{\line(0,-1){.086853}}
\multiput(26.13,101.16)(-.031172,-.094801){13}{\line(0,-1){.094801}}
\multiput(25.72,99.93)(-.03251,-.11337){11}{\line(0,-1){.11337}}
\multiput(25.37,98.68)(-.03094,-.12599){10}{\line(0,-1){.12599}}
\multiput(25.06,97.42)(-.0326,-.15885){8}{\line(0,-1){.15885}}
\multiput(24.8,96.15)(-.03026,-.18285){7}{\line(0,-1){.18285}}
\multiput(24.58,94.87)(-.0325,-.25742){5}{\line(0,-1){.25742}}
\put(24.42,93.59){\line(0,-1){1.292}}
\put(24.31,92.29){\line(0,-1){1.296}}
\put(24.25,91){\line(0,-1){1.297}}
\put(24.23,89.7){\line(0,-1){1.297}}
\put(24.27,88.4){\line(0,-1){1.294}}
\multiput(24.36,87.11)(.02718,-.25804){5}{\line(0,-1){.25804}}
\multiput(24.49,85.82)(.03089,-.214){6}{\line(0,-1){.214}}
\multiput(24.68,84.53)(.03351,-.18228){7}{\line(0,-1){.18228}}
\multiput(24.91,83.26)(.03149,-.14067){9}{\line(0,-1){.14067}}
\multiput(25.19,81.99)(.03318,-.12542){10}{\line(0,-1){.12542}}
\multiput(25.53,80.74)(.03164,-.10338){12}{\line(0,-1){.10338}}
\multiput(25.91,79.5)(.032853,-.094231){13}{\line(0,-1){.094231}}
\multiput(26.33,78.27)(.031589,-.080513){15}{\line(0,-1){.080513}}
\multiput(26.81,77.06)(.032492,-.074288){16}{\line(0,-1){.074288}}
\multiput(27.33,75.88)(.033244,-.068692){17}{\line(0,-1){.068692}}
\multiput(27.89,74.71)(.032083,-.060273){19}{\line(0,-1){.060273}}
\multiput(28.5,73.56)(.032656,-.056047){20}{\line(0,-1){.056047}}
\multiput(29.15,72.44)(.033128,-.052144){21}{\line(0,-1){.052144}}
\multiput(29.85,71.35)(.033511,-.048522){22}{\line(0,-1){.048522}}
\multiput(30.59,70.28)(.032404,-.043266){24}{\line(0,-1){.043266}}
\multiput(31.37,69.24)(.032681,-.04031){25}{\line(0,-1){.04031}}
\multiput(32.18,68.23)(.032889,-.037524){26}{\line(0,-1){.037524}}
\multiput(33.04,67.26)(.033036,-.034891){27}{\line(0,-1){.034891}}
\multiput(33.93,66.32)(.034351,-.033596){27}{\line(1,0){.034351}}
\multiput(34.86,65.41)(.036986,-.033492){26}{\line(1,0){.036986}}
\multiput(35.82,64.54)(.039775,-.033329){25}{\line(1,0){.039775}}
\multiput(36.81,63.7)(.042736,-.0331){24}{\line(1,0){.042736}}
\multiput(37.84,62.91)(.045888,-.032801){23}{\line(1,0){.045888}}
\multiput(38.89,62.16)(.049255,-.032424){22}{\line(1,0){.049255}}
\multiput(39.98,61.44)(.055511,-.033559){20}{\line(1,0){.055511}}
\multiput(41.09,60.77)(.059746,-.033055){19}{\line(1,0){.059746}}
\multiput(42.22,60.14)(.064359,-.032443){18}{\line(1,0){.064359}}
\multiput(43.38,59.56)(.073752,-.03369){16}{\line(1,0){.073752}}
\multiput(44.56,59.02)(.079991,-.032888){15}{\line(1,0){.079991}}
\multiput(45.76,58.53)(.086995,-.031919){14}{\line(1,0){.086995}}
\multiput(46.98,58.08)(.10285,-.03331){12}{\line(1,0){.10285}}
\multiput(48.21,57.68)(.11351,-.032){11}{\line(1,0){.11351}}
\multiput(49.46,57.33)(.12612,-.03038){10}{\line(1,0){.12612}}
\multiput(50.72,57.02)(.159,-.0319){8}{\line(1,0){.159}}
\multiput(52,56.77)(.18298,-.02944){7}{\line(1,0){.18298}}
\multiput(53.28,56.56)(.25756,-.03135){5}{\line(1,0){.25756}}
\put(54.56,56.41){\line(1,0){1.293}}
\put(55.86,56.3){\line(1,0){1.296}}
\put(57.15,56.24){\line(1,0){1.297}}
\put(58.45,56.23){\line(1,0){1.297}}
\put(59.75,56.28){\line(1,0){1.294}}
\multiput(61.04,56.37)(.25791,.02833){5}{\line(1,0){.25791}}
\multiput(62.33,56.51)(.21386,.03184){6}{\line(1,0){.21386}}
\multiput(63.61,56.7)(.15936,.03003){8}{\line(1,0){.15936}}
\multiput(64.89,56.94)(.14053,.03211){9}{\line(1,0){.14053}}
\multiput(66.15,57.23)(.12527,.03374){10}{\line(1,0){.12527}}
\multiput(67.41,57.57)(.10323,.0321){12}{\line(1,0){.10323}}
\multiput(68.65,57.95)(.094084,.033272){13}{\line(1,0){.094084}}
\multiput(69.87,58.39)(.080372,.031947){15}{\line(1,0){.080372}}
\multiput(71.07,58.86)(.074143,.032822){16}{\line(1,0){.074143}}
\multiput(72.26,59.39)(.068543,.033549){17}{\line(1,0){.068543}}
\multiput(73.43,59.96)(.06013,.032351){19}{\line(1,0){.06013}}
\multiput(74.57,60.58)(.055901,.032905){20}{\line(1,0){.055901}}
\multiput(75.69,61.23)(.051996,.03336){21}{\line(1,0){.051996}}
\multiput(76.78,61.93)(.048373,.033726){22}{\line(1,0){.048373}}
\multiput(77.84,62.68)(.043121,.032596){24}{\line(1,0){.043121}}
\multiput(78.88,63.46)(.040164,.032859){25}{\line(1,0){.040164}}
\multiput(79.88,64.28)(.037377,.033056){26}{\line(1,0){.037377}}
\multiput(80.85,65.14)(.034743,.03319){27}{\line(1,0){.034743}}
\multiput(81.79,66.04)(.033443,.034501){27}{\line(0,1){.034501}}
\multiput(82.69,66.97)(.033327,.037135){26}{\line(0,1){.037135}}
\multiput(83.56,67.93)(.033151,.039923){25}{\line(0,1){.039923}}
\multiput(84.39,68.93)(.03291,.042882){24}{\line(0,1){.042882}}
\multiput(85.18,69.96)(.032597,.046033){23}{\line(0,1){.046033}}
\multiput(85.93,71.02)(.033738,.051751){21}{\line(0,1){.051751}}
\multiput(86.64,72.1)(.033312,.055659){20}{\line(0,1){.055659}}
\multiput(87.3,73.22)(.032789,.059892){19}{\line(0,1){.059892}}
\multiput(87.93,74.36)(.032156,.064503){18}{\line(0,1){.064503}}
\multiput(88.51,75.52)(.033362,.073901){16}{\line(0,1){.073901}}
\multiput(89.04,76.7)(.032532,.080137){15}{\line(0,1){.080137}}
\multiput(89.53,77.9)(.031532,.087136){14}{\line(0,1){.087136}}
\multiput(89.97,79.12)(.03286,.103){12}{\line(0,1){.103}}
\multiput(90.36,80.36)(.0315,.11365){11}{\line(0,1){.11365}}
\multiput(90.71,81.61)(.03314,.14029){9}{\line(0,1){.14029}}
\multiput(91.01,82.87)(.03119,.15914){8}{\line(0,1){.15914}}
\multiput(91.26,84.14)(.0334,.21363){6}{\line(0,1){.21363}}
\multiput(91.46,85.43)(.03021,.2577){5}{\line(0,1){.2577}}
\put(91.61,86.71){\line(0,1){1.293}}
\put(91.71,88.01){\line(0,1){1.993}}
%\end
\put(9.5,89.75){\line(1,0){97.5}}
\put(57.75,42.5){\line(0,1){94.25}}
%\vector(58,89.75)(109,116)
\put(109,116){\vector(2,1){.07}}
\multiput(58,89.75)(.0654685494,.0336970475){779}{\line(1,0){.0654685494}}
%\end
%\vector(57.75,90)(81.25,136.75)
\put(81.25,136.75){\vector(1,2){.07}}
\multiput(57.75,90)(.0337159254,.0670731707){697}{\line(0,1){.0670731707}}
%\end
\put(55,87.5){\makebox(0,0)[cc]{$O$}}
\put(72,124.25){\makebox(0,0)[cc]{$P$}}
\put(92.25,103.75){\makebox(0,0)[cc]{$Q$}}
\put(58,35.25){\makebox(0,0)[cc]{Figure 7.13}}
\put(63,95.5){\makebox(0,0)[cc]{$\theta$}}
\end{picture}
\end{center}

Let's suppose that $\theta$ is the measure of angle $POQ$. We know how to find $\theta$ by our definition of angle measure. Since the radius of the circle is $r>0$, we have the following picture extracted from the one above. Can you calculate $PQ$?

% Figure 7.14 added

%TeXCAD Picture [fig7.14.pic]. Options:
%\grade{\on}
%\emlines{\off}
%\epic{\off}
%\beziermacro{\on}
%\reduce{\on}
%\snapping{\off}
%\quality{8.00}
%\graddiff{0.01}
%\snapasp{1}
%\zoom{4.0000}
\unitlength 1mm % = 2.85pt
\linethickness{0.4pt}
\ifx\plotpoint\undefined\newsavebox{\plotpoint}\fi % GNUPLOT compatibility
\begin{picture}(70.25,44.75)(0,0)
%\emline(58,40.75)(66.75,15.25)
\multiput(58,40.75)(.033653846,-.098076923){260}{\line(0,-1){.098076923}}
%\end
\put(66.75,15.25){\line(-1,0){45}}
%\emline(21.75,15.25)(57.75,41)
\multiput(21.75,15.25)(.0471204188,.0337041885){764}{\line(1,0){.0471204188}}
%\end
\put(18.75,12.75){\makebox(0,0)[cc]{$O$}}
\put(70.25,12.75){\makebox(0,0)[cc]{$Q$}}
\put(59.5,43.75){\makebox(0,0)[cc]{$P$}}
\put(65,29.25){\makebox(0,0)[cc]{$PQ$}}
\put(35.25,29.25){\makebox(0,0)[cc]{$r$}}
\put(45.5,11.75){\makebox(0,0)[cc]{$r$}}
\put(45.25,4.75){\makebox(0,0)[cc]{Figure 7.14}}
\end{picture}

From the Law of Cosines, we have: $(PQ)^{2} = r^{2} + r^{2} - 2rr \cos \theta$ $= 2r^{2} - 2r^{2} \cos \theta$ $= 2r^{2} (1 - \cos \theta)$. Hence, $PQ = r \sqrt{2} \sqrt{(1 - \cos \theta)}$.

But how do we calculate $l(\arc{PQ})$? At this point, we really have no way of getting at the length of an arc of a circle that is not the unit circle. But consider the following picture.

\begin{center}
%TeXCAD Options
%\grade{\on}
%\emlines{\off}
%\epic{\off}
%\beziermacro{\on}
%\reduce{\on}
%\snapping{\off}
%\quality{8.00}
%\graddiff{0.01}
%\snapasp{1}
%\zoom{4.0000}
\unitlength 0.9mm % = 2.85pt
\linethickness{0.4pt}
\ifx\plotpoint\undefined\newsavebox{\plotpoint}\fi % GNUPLOT compatibility
\begin{picture}(124.25,86.88)(10,60)
\put(12,75){\line(1,0){112.25}}
%\emline(41.75,75)(121.5,137.5)
\multiput(41.75,75)(.04303831624,.03372908797){1853}{\line(1,0){.04303831624}}
%\end
\qbezier(43.75,114)(80.38,117)(81.5,75)
\qbezier(111.25,75)(104,146.88)(42.75,141.25)
\put(52.25,78){\makebox(0,0)[cc]{$\theta$}}
\put(40.25,72.25){\makebox(0,0)[cc]{$O:(0,0)$}}
\put(81.25,71){\makebox(0,0)[cc]{$I:(1,0)$}}
\put(110.5,71.5){\makebox(0,0)[cc]{$R:(r,0)$}}
\put(89,101){\makebox(0,0)[cc]{$P:(\cos \theta, \sin \theta)$}}
\put(111.25,117.75){\makebox(0,0)[cc]{$Q:(r \cos \theta, \sin \theta)$}}
\put(58.75,62.25){\makebox(0,0)[cc]{Figure 7.15}}
\end{picture}
\end{center}

We have two circles with center at the origin, one with radius $1$ and the other with radius $r>0$. Note that angle $QOR$ = angle $POI$, and that $m \angm POI = \theta$.

We have been toiling over a relationship between arc length and chord length on circles, and we have pretty well solved the problem on the unit circle; i.e., a circle with radius $1$. How do we extend our results to circles with radii different from $1$? Do you think that the relationship between $l(\arc{PI})$ and $PI$ on the unit circle should be the same as the relationship between $l(\arc{QR})$ and $QR$ on the circle with radius $r>0$? Shouldn't it be the case that the ratio of $l(\arc{PI})$ and $PI$ should be the same as the ratio of $l(\arc{QR})$ and $QR$; i.e., shouldn't $\frac{l(\arc{PI})}{PI} = \frac{l(\arc{QR})}{QR}$? If this looks compelling to you, let's see what it leads to.

Suppose, just for the sake of argument, that $\frac{l(\arc{PI})}{PI} = \frac{l(\arc{QR})}{QR}$. Note that the only one of these that we don't know is $l(\arc{QR})$. From this equation, $l(\arc{QR}) = (\frac{QR}{PI}) l(\arc{PI})$. 

Now $QR = \sqrt{ (r \cos \theta - r)^{2} + (r \sin \theta - 0)^{2} }$ $= r \sqrt{2} \sqrt{ (1 -\cos \theta) }$, 

and $PI = \sqrt{ (\cos \theta - 1)^{2} + (\sin \theta - 0)^{2} }$ $= \sqrt{2} \sqrt{ (1 - \cos \theta) }$. 

$l(\arc{PI}) = \theta$. Do you see where this came from? If not, look back at the definition of angle measure.

Then $l(\arc{QR}) = (\frac{QR}{PI}) l(\arc{PI}) = r \theta$.

Now, that's a real nice formula, if it's true. But how can we find out whether or not it is true? Unfortunately, we'll have to wait until the next course, where we'll get a definition of arc length. Then, on the basis of that definition, we can see whether our conjecture that $l(\arc{QR}) = r \theta$ is valid. 

There is something we can do in the meantime, however. We can see if this result of ours leads us to something that contradicts some other thing that we might know. For example, suppose the point $Q$ is such that $QOR$; i.e., $Q$, $O$, and $R$ are colinear. Then, the measure we are talking about, the arc length, is the length of a semi-circle of the circle with radius $r$. Our formula tells us that the length of the semi-circle should be $r \theta$. But what is $\theta$? Well, if $QOR$, then the point $P$ of the unit circle must be $J$. This means $\theta = \pi$. Do you see that? Then the length of a semi-circle of the circle with radius $r>0$ is $r \pi$. The whole length around the circle, the circumference of the circle, is $2(r \pi)$, or $2 \pi r$. Do you recognize that? This is encouraging, since our conjecture led to a result that we were told was true back a long time ago.

Let's accept our conjecture as an axiom, at least temporarily, and see if we can solve our main problem.

Does it matter that we used $I:(1,0)$ and $R:(r,0)$ in our analysis? What if the picture looked like the following?

\begin{center}
%TeXCAD Picture [fig7.16.pic]. Options:
%\grade{\on}
%\emlines{\off}
%\epic{\off}
%\beziermacro{\on}
%\reduce{\on}
%\snapping{\off}
%\quality{8.00}
%\graddiff{0.01}
%\snapasp{1}
%\zoom{4.0000}
\unitlength 1mm % = 2.85pt
\linethickness{0.4pt}
\ifx\plotpoint\undefined\newsavebox{\plotpoint}\fi % GNUPLOT compatibility
\begin{picture}(107.5,59)(10,84)
\put(19.5,94.25){\line(1,0){96.5}}
%\emline(13.5,124.25)(63.5,94.5)
\multiput(13.5,124.25)(.0566893424,-.0337301587){882}{\line(1,0){.0566893424}}
%\end
%\emline(63.5,94.5)(117.5,143)
\multiput(63.5,94.5)(.0375521558,.0337273992){1438}{\line(1,0){.0375521558}}
%\end
%\emline(44.75,105.5)(85,114)
\multiput(44.75,105.5)(.159722222,.033730159){252}{\line(1,0){.159722222}}
%\end
%\emline(106.25,133)(27.25,116)
\multiput(106.25,133)(-.156746032,-.033730159){504}{\line(-1,0){.156746032}}
%\end
\put(106.25,133){\circle{.5}}
\put(85,113.75){\circle{.71}}
\put(45,105.5){\circle{.5}}
\put(27.25,116){\circle{.71}}
\put(87.5,133.5){\makebox(0,0)[cc]{$T:(r \cos \phi,r \sin \phi)$}}
\put(27.75,120.5){\makebox(0,0)[cc]{$r \sin \theta)$}}
\put(28,124.75){\makebox(0,0)[cc]{$S:(r \cos \theta,$}}
\put(32.5,104){\makebox(0,0)[cc]{$P:(\cos \theta,\sin \theta)$}}
\put(101.25,112.75){\makebox(0,0)[cc]{$Q:(\cos \phi,\sin \phi)$}}
\put(63.75,90.5){\makebox(0,0)[cc]{$O:(0,0)$}}
\put(63.75,85){\makebox(0,0)[cc]{Figure 7.16}}
\put(88.5,91){\makebox(0,0)[cc]{$I:(1,0)$}}
\put(108,91.5){\makebox(0,0)[cc]{$R:(r,0)$}}
\put(88.25,94.25){\circle{.5}}
\put(108,94.25){\circle{.5}}
\end{picture}
\end{center}

Our conjecture, generalized, is $\frac{ l(\arc{ST}) }{ST} = \frac{ l(\arc{PQ}) }{PQ}$. The only thing we don't know is $l(\arc{ST})$. 

$l(\arc{ST}) = \frac{ST}{PQ} l(\arc{PQ})$.

Now, find $ST$, $PQ$, and $l(\arc{PQ})$ before you go on.

$ST = \sqrt{ (r \cos \theta - r \cos \phi)^{2} + (r \sin \theta - r \sin \phi)^{2} }$ $= r \sqrt{2} \sqrt{ 1 - \cos (\theta - \phi) }$

$PQ = \sqrt{2} \sqrt{ 1 - \cos (\theta - \phi) }$

$l(\arc{PQ}) = \theta - \phi$.

Hence $l(\arc{ST}) = r(\theta - \phi)$, which is certainly consistent with our previous result.

Can you see how this solves our main problem? That problem, stated in terms of the points $S$ and $T$ of the circle with radius $r>0$, is the following:

\begin{enumerate}[\hspace{1cm}1)]
  \item If you know $ST$, can you calculate $l(\arc{ST})$?
  \item If you know $l(\arc{ST})$, can you calculate $ST$?
\end{enumerate}

Solve this now before you go on.

\begin{enumerate}[\hspace{1cm}1)]
  \item Suppose you know $ST$.

        We have $ST = r \sqrt{2} \sqrt{ 1 - \cos (\theta-\phi) }$. Our problem is to calculate $l(\arc{ST})$ which, by our formula, is $r(\theta-\phi)$. Hence we must extract $\theta - \phi$ from $r \sqrt{2} \sqrt{ 1 - \cos (\theta-\phi) }$.

        $(ST)^{2} = 2r^{2} (1 - \cos (\theta-\phi) )$.

        Hence, $\cos (\theta-\phi) = \frac{ 1-(ST)^{2} }{ 2r^{2} }$

        so that $\theta - \phi = \cos^{-1} ( \frac{ 1-(ST)^{2} }{ 2r^{2} } )$ 

        and $l(\arc{ST}) = r(\theta-\phi) = r \cos^{-1} ( \frac{ 1-(ST)^{2} }{ 2r^{2} } )$.

        This looks very much like our solution on the unit circle, and it certainly seems reasonable that it should.

  \item Now, suppose we know $l(\arc{ST})$. Can we calculate $ST$?

        $l(\arc{ST}) = r (\theta-\phi)$ so $\theta - \phi = \frac{ l(\arc{ST}) }{ r }$.

        Hence, $ST = r \sqrt{2} \sqrt{1 - \cos (\theta-\phi)} = r \sqrt{2} \sqrt{ 1 -\cos ( \frac{ l(\arc{ST}) }{ r } ) }$.

        Eureka! Hallelujah! Everybody gets an A!
\end{enumerate}

\begin{description}

\item[Problem 7.13:] Refer to Figure 7.16. Suppose that, instead of assuming that

  $\frac{ l(\arc{ST}) }{ST} = \frac{ l(\arc{PQ}) }{PQ}$, we assume that the circumference of a circle with radius $r>0$ is $2 \pi r$. Can you solve our main problem on the basis of this assumption?

\item[Problem 7.14:] Refer to Figure 7.16. Do you remember from some course you took in the distant past that the area of a circle with radius $r>0$ is $\pi r^{2}$? Well, let's assume it, anyway, for the sake of argument. Can you see how to find the area of the piece of the interior of the circle with radius $r$ bounded by (a) the circle, (b) the ray $OS$, and (c) the ray $OT$? This region is known as a ``sector'' of a circle.

\end{description}

If your curiosity has driven you to look at a standard textbook on trigonometry, then you have seen a vastly different approach to the subject we have taken here. When it comes to solving triangles, we draw a picture like the following:

\begin{center}
%TeXCAD Options
%\grade{\on}
%\emlines{\off}
%\epic{\off}
%\beziermacro{\on}
%\reduce{\on}
%\snapping{\off}
%\quality{8.00}
%\graddiff{0.01}
%\snapasp{1}
%\zoom{4.0000}
\unitlength 1mm % = 2.85pt
\linethickness{0.4pt}
\ifx\plotpoint\undefined\newsavebox{\plotpoint}\fi % GNUPLOT compatibility
\begin{picture}(80,22.5)(0,103)
%\emline(16,106.25)(54,125.5)
\multiput(16,106.25)(.066549912,.033712785){571}{\line(1,0){.066549912}}
%\end
%\emline(54,125.5)(80,106.75)
\multiput(54,125.5)(.04676259,-.033723022){556}{\line(1,0){.04676259}}
%\end
\put(80,106.75){\line(-1,0){64}}
\put(68.75,117.5){\makebox(0,0)[cc]{$a$}}
\put(48,104.25){\makebox(0,0)[cc]{$c$}}
\put(32.5,118.25){\makebox(0,0)[cc]{$b$}}
\put(26,108.5){\makebox(0,0)[cc]{$\alpha$}}
\put(53,121.25){\makebox(0,0)[cc]{$\gamma$}}
\put(70,109){\makebox(0,0)[cc]{$\beta$}}
\end{picture}
\end{center}

$\alpha$, $\beta$, and $\gamma$ are measures of the angles, and, as a consequence, each of these numbers is in the interval $[ 0, \pi ]$. But in the standard texts, the picture may look like the following:

\begin{center}
%TeXCAD Options
%\grade{\on}
%\emlines{\off}
%\epic{\off}
%\beziermacro{\on}
%\reduce{\on}
%\snapping{\off}
%\quality{8.00}
%\graddiff{0.01}
%\snapasp{1}
%\zoom{4.0000}
\unitlength 1mm % = 2.85pt
\linethickness{0.4pt}
\ifx\plotpoint\undefined\newsavebox{\plotpoint}\fi % GNUPLOT compatibility
\begin{picture}(80,22.5)(0,103)
%\emline(16,106.25)(54,125.5)
\multiput(16,106.25)(.066549912,.033712785){571}{\line(1,0){.066549912}}
%\end
%\emline(54,125.5)(80,106.75)
\multiput(54,125.5)(.04676259,-.033723022){556}{\line(1,0){.04676259}}
%\end
\put(80,106.75){\line(-1,0){64}}
\put(68.75,117.5){\makebox(0,0)[cc]{$a$}}
\put(48,104.25){\makebox(0,0)[cc]{$c$}}
\put(32.5,118.25){\makebox(0,0)[cc]{$b$}}
\put(25.25,108.5){\makebox(0,0)[cc]{$\alpha \degrees$}}
\put(53,121.5){\makebox(0,0)[cc]{$\gamma \degrees$}}
\put(69.5,109){\makebox(0,0)[cc]{$\beta \degrees$}}
\end{picture}
\end{center}

What are $\alpha \degrees$, $\beta \degrees$, and $\gamma \degrees$?

These are read ``alpha degrees'', ``beta degrees'', and ``gamma degrees'', respectively. 

What is this ``degree'' business? I'm sure you've heard of degrees in terms of angles, but when we've talked of angle measure in this course, we've never mentioned anything about degrees. Here is how it works. Degree measure is just a different way of measuring an angle. We measure an angle by relating the angle to an arc length on a unit circle. The measure that we get from this method is a number in the interval $[0,\pi]$ and is called the ``radian measure'' of the angle. 

Standard textbooks usually measure an angle using the following method. Subdivide our upper semi-circle $\arc{JiI}$ into $180$ non-overlapping arcs filling up the whole semi-circle. Each of these little arcs corresponds to one degree ($1 \degrees$) of angle measure. Here is how it works. Suppose we have angle $PQR$, $P_{1}$ is a point on ray $QP$, $R_{1}$ is a point of ray $QR$, and $QP_{1} = QR_{1} = 1$. Then, by our method, $m \angm PQR = l(\arc{P_{1}R_{1}})$. 

By the degree measure of angle PQR is meant the number of non-overlapping arcs in the subdivision of $\arc{JiI}$ that fill up the arc $P_{1}R_{1}$. The number of such $1 \degrees$ arcs filling up arc $\arc{P_{1}R_{1}}$ is the degree measure of $PQR$. If there are five such $1 \degrees$ arcs filling up arc $\arc{P_{1}R_{1}}$, then the degree measure of angle $PQR$ is $5$; i.e., the measure of angle $PQR$ is $5$ degrees, written $5 \degrees$. 

I can hear you protesting. What if seven of the arcs are not enough to fill up the arc $\arc{P_{1}R_{1}}$, but eight are too many? Our clairvoyant mathematical forefathers anticipated your question. But before I tell you what they did, ponder your question and see if you can come up with a solution yourself.

Did you figure out what to do if the $1 \degrees$ arcs don't come out even in filling up arc $\arc{P_{1}R_{1}}$? What our clever ancestors did was simply to divide up each $1 \degrees$ arc into $60$ sub-arcs. They called each of these sub-arcs one minute, denoted $1'$. Then, guess what they did. They divided each minute up into $60$ sub-arcs and called each one of these a second, denoted $1''$. Hence $1 \degrees = 60'$ and $1' = 60''$. So $1 \degrees = 3600''$. And I guess they figured that was good enough for government work. Today, we cut it even finer, but we simply use fractions of a second in terms of decimals; e.g., $.7258''$. Or, depending on what the circumstances are, the degree may be broken down using decimals instead of minutes and seconds. 

This may be what you conjectured earlier as a solution to the uneven fit of degrees in an arc. But, when you come right down to it, most people don't use degree measure anymore at all. Our measure, radian measure, is the measure du jour and is preferred by doctors 2 to 1 over degree measure. Does it really make a difference? Not really. It's just a matter of which is the most convenient to use in your particular circumstance. But, what if your problem is stated in radians and you, for some strange reason, want your answer in degrees? I'll bet you can find a formula to convert back and forth between degrees and radians. Try to do so before going on.

We note that $0 \degrees = 0$ radians, $180 \degrees = \pi$ radians, $90 \degrees = \frac{\pi}{2}$ radians, etc. The relationship looks linear. Draw a graph. Put radians along the $x$-axis and degrees along the $y$-axis. Or the other way around. It doesn't matter.

\begin{center}
%TeXCAD Options
%\grade{\on}
%\emlines{\off}
%\epic{\off}
%\beziermacro{\on}
%\reduce{\on}
%\snapping{\off}
%\quality{8.00}
%\graddiff{0.01}
%\snapasp{1}
%\zoom{4.0000}
\unitlength 1mm % = 2.85pt
\linethickness{0.4pt}
\ifx\plotpoint\undefined\newsavebox{\plotpoint}\fi % GNUPLOT compatibility
\begin{picture}(67.5,40)(0,96)
\put(21,136){\line(0,-1){39.75}} \put(12.5,105.5){\line(1,0){55}}
%\emline(21,105.5)(66.25,135.25)
\multiput(21,105.5)(.0513038549,.0337301587){882}{\line(1,0){.0513038549}}
%\end
%\dashline{1}(66,134.75)(66,105.75)
\put(65.93,134.68){\line(0,-1){.967}}
\put(65.93,132.75){\line(0,-1){.967}}
\put(65.93,130.81){\line(0,-1){.967}}
\put(65.93,128.88){\line(0,-1){.967}}
\put(65.93,126.95){\line(0,-1){.967}}
\put(65.93,125.01){\line(0,-1){.967}}
\put(65.93,123.08){\line(0,-1){.967}}
\put(65.93,121.15){\line(0,-1){.967}}
\put(65.93,119.21){\line(0,-1){.967}}
\put(65.93,117.28){\line(0,-1){.967}}
\put(65.93,115.35){\line(0,-1){.967}}
\put(65.93,113.41){\line(0,-1){.967}}
\put(65.93,111.48){\line(0,-1){.967}}
\put(65.93,109.55){\line(0,-1){.967}}
\put(65.93,107.61){\line(0,-1){.967}}
%\end
%\dashline{1}(66,135)(21,134)
\put(65.93,134.93){\line(-1,0){.978}}
\put(63.97,134.89){\line(-1,0){.978}}
\put(62.02,134.84){\line(-1,0){.978}}
\put(60.06,134.8){\line(-1,0){.978}}
\put(58.1,134.76){\line(-1,0){.978}}
\put(56.15,134.71){\line(-1,0){.978}}
\put(54.19,134.67){\line(-1,0){.978}}
\put(52.23,134.63){\line(-1,0){.978}}
\put(50.28,134.58){\line(-1,0){.978}}
\put(48.32,134.54){\line(-1,0){.978}}
\put(46.36,134.49){\line(-1,0){.978}}
\put(44.41,134.45){\line(-1,0){.978}}
\put(42.45,134.41){\line(-1,0){.978}}
\put(40.49,134.36){\line(-1,0){.978}}
\put(38.54,134.32){\line(-1,0){.978}}
\put(36.58,134.28){\line(-1,0){.978}}
\put(34.63,134.23){\line(-1,0){.978}}
\put(32.67,134.19){\line(-1,0){.978}}
\put(30.71,134.15){\line(-1,0){.978}}
\put(28.76,134.1){\line(-1,0){.978}}
\put(26.8,134.06){\line(-1,0){.978}}
\put(24.84,134.02){\line(-1,0){.978}}
\put(22.89,133.97){\line(-1,0){.978}}
%\end
%\dashline{1}(21.25,121)(45,121.5)
\put(21.18,120.93){\line(1,0){.99}}
\put(23.16,120.97){\line(1,0){.99}}
\put(25.14,121.01){\line(1,0){.99}}
\put(27.12,121.05){\line(1,0){.99}}
\put(29.1,121.1){\line(1,0){.99}}
\put(31.08,121.14){\line(1,0){.99}}
\put(33.05,121.18){\line(1,0){.99}}
\put(35.03,121.22){\line(1,0){.99}}
\put(37.01,121.26){\line(1,0){.99}}
\put(38.99,121.3){\line(1,0){.99}}
\put(40.97,121.35){\line(1,0){.99}}
\put(42.95,121.39){\line(1,0){.99}}
%\end
%\dashline{1}(45,121.5)(45,106)
\put(44.93,121.43){\line(0,-1){.969}}
\put(44.93,119.49){\line(0,-1){.969}}
\put(44.93,117.55){\line(0,-1){.969}}
\put(44.93,115.62){\line(0,-1){.969}}
\put(44.93,113.68){\line(0,-1){.969}}
\put(44.93,111.74){\line(0,-1){.969}}
\put(44.93,109.8){\line(0,-1){.969}}
\put(44.93,107.87){\line(0,-1){.969}}
%\end
\put(23.5,103){\makebox(0,0)[cc]{$0$}}
\put(17,108.25){\makebox(0,0)[cc]{$0 \degrees$}}
\put(15.25,120.25){\makebox(0,0)[cc]{$90 \degrees$}}
\put(14.25,133.75){\makebox(0,0)[cc]{$180 \degrees$}}
\put(44.75,102.5){\makebox(0,0)[cc]{$\frac{\pi}{2}$}}
\put(65.25,103){\makebox(0,0)[cc]{$\pi$}}
\end{picture}
\end{center}

You have two points of a line; $(0,0)$ and $(\pi,180)$. All you have to do is find the equation of the line containing $(0,0)$ and $(\pi,180)$. This equation should be degrees $= (180 \degrees,\pi)$ radians. 

Hence, degree measure $= (180 \degrees,\pi)$ (radian measure),

and radian measure $= (\frac{\pi}{180 \degrees})$ (degree measure).

Certain of these crop up so often as to suggest noting them.

$\frac{\pi}{2}$ radians $= 90 \degrees$, $\frac{\pi}{3}$ rad $= 60 \degrees$, $\frac{\pi}{6}$ rad $= 30 \degrees$, $\frac{\pi}{4}$ rad $= 45 \degrees$. 

We can even write $\cos 30 \degrees = \cos \frac{\pi}{6} = \frac{\pi 3}{2}$, and $\sin 30 \degrees = \sin \frac{\pi}{6} = \frac{1}{2}$.

Everything will be fine until you see something like this: $\sin 240 \degrees$. There isn't any such thing as an angle of measure $240 \degrees$. We know this because our angle measure is a number in the interval $[0,\pi]$. Hence, degree measure of an angle must be in the interval $[0,180]$. $240$ is not in that interval. But, just as we can talk about $\sin x$ where $x$ is not in $[0,\pi]$, in fact $x$ can be any number whatsoever, then perhaps we can use radians $= (\frac{\pi}{180} \degrees)$ (degree measure) to make some sense of $\sin 240 \degrees$.

$(\frac{\pi}{180} \degrees) (240 \degrees) = (\frac{4}{3}) \pi$, so $\sin 240 \degrees = \sin ((\frac{4}{3}) \pi)$, which we can find. 

Similarly, $\cos (-759 \degrees) = \cos (-(\frac{\pi}{180} \degrees) 759 \degrees) = \cos ((\frac{759}{180})\pi)$, which we can find -- with some effort and hair pulling, but we can find it.

There are lots of other questions about degree measure, but I'll bet you can figure them out yourself now that you can convert from degrees to radians.

The following is a set of problems that are typical of those found in normal geeky textbooks. I have included these to give you a taste of how you can use the trigonometric functions (cosine, sine, and those defined in terms of them -- namely tangent, cotangent, secant, and cosecant) to solve what some people call ``practical problems''. Have fun!

\begin{description}

\item[Problem 7.15:] In each of the following, find everything required to four decimal places, calculating angles in terms of both radians and degrees.

\end{description}

\begin{center}
%TeXCAD Picture [prob7.15_01-05.pic]. Options:
%\grade{\on}
%\emlines{\off}
%\epic{\off}
%\beziermacro{\on}
%\reduce{\on}
%\snapping{\off}
%\quality{8.00}
%\graddiff{0.01}
%\snapasp{1}
%\zoom{4.0000}
\unitlength 1mm % = 2.85pt
\linethickness{0.4pt}
\ifx\plotpoint\undefined\newsavebox{\plotpoint}\fi % GNUPLOT compatibility
\begin{picture}(75.75,110.25)(0,25)
\put(16.75,118.75){\line(1,0){47}}
\put(63.75,118.75){\line(0,1){16.5}}
%\emline(63.75,135.25)(16.75,119)
\multiput(63.75,135.25)(-.097510373,-.033713693){482}{\line(-1,0){.097510373}}
%\end
\put(60.75,120.75){\line(0,-1){1.5}}
\put(60.75,121.5){\line(0,-1){2.75}}
\put(60.75,121.5){\line(1,0){3}}
\put(16,96.5){\line(1,0){47}}
\put(63,96.5){\line(0,1){16.5}}
%\emline(63,113)(16,96.75)
\multiput(63,113)(-.097510373,-.033713693){482}{\line(-1,0){.097510373}}
%\end
\put(60,98.5){\line(0,-1){1.5}}
\put(60,99.25){\line(0,-1){2.75}}
\put(60,99.25){\line(1,0){3}}
\put(15.75,73.25){\line(1,0){47}}
\put(62.75,73.25){\line(0,1){16.5}}
%\emline(62.75,89.75)(15.75,73.5)
\multiput(62.75,89.75)(-.097510373,-.033713693){482}{\line(-1,0){.097510373}}
%\end
\put(59.75,75.25){\line(0,-1){1.5}}
\put(59.75,76){\line(0,-1){2.75}}
\put(59.75,76){\line(1,0){3}}
\put(16,52){\line(1,0){47}}
\put(63,52){\line(0,1){16.5}}
%\emline(63,68.5)(16,52.25)
\multiput(63,68.5)(-.097510373,-.033713693){482}{\line(-1,0){.097510373}}
%\end
\put(60,54){\line(0,-1){1.5}}
\put(60,54.75){\line(0,-1){2.75}}
\put(60,54.75){\line(1,0){3}}
\put(15.75,30){\line(1,0){47}}
\put(62.75,30){\line(0,1){16.5}}
%\emline(62.75,46.5)(15.75,30.25)
\multiput(62.75,46.5)(-.097510373,-.033713693){482}{\line(-1,0){.097510373}}
%\end
\put(59.75,32){\line(0,-1){1.5}}
\put(59.75,32.75){\line(0,-1){2.75}}
\put(59.75,32.75){\line(1,0){3}}
\put(41,131){\makebox(0,0)[cc]{$50$}}
\put(28.5,120.25){\makebox(0,0)[cc]{$20 \degrees$}}
\put(60.5,130.75){\makebox(0,0)[cc]{$\theta$}}
\put(67,127.25){\makebox(0,0)[cc]{$x$}}
\put(75.75,133){\makebox(0,0)[lc]{$x = $}}
\put(75.5,122.5){\makebox(0,0)[lc]{$\theta = $}}
\put(75.25,110){\makebox(0,0)[lc]{$x = $}}
\put(75.25,100.25){\makebox(0,0)[lc]{$\theta = $}}
\put(27.75,98.5){\makebox(0,0)[cc]{$\theta$}}
\put(40.5,108.75){\makebox(0,0)[cc]{$25$}}
\put(66.25,104.25){\makebox(0,0)[cc]{$12$}}
\put(40,94.25){\makebox(0,0)[cc]{$x$}}
\put(14.5,135){\makebox(0,0)[cc]{1)}}
\put(14.75,108){\makebox(0,0)[cc]{2)}}
\put(14,87){\makebox(0,0)[cc]{3)}}
\put(15.25,65.75){\makebox(0,0)[cc]{4)}}
\put(15.5,43.5){\makebox(0,0)[cc]{5)}}
\put(27,32.25){\makebox(0,0)[cc]{$\theta$}}
\put(39,41.5){\makebox(0,0)[cc]{$r$}}
\put(39.25,26.75){\makebox(0,0)[cc]{$x$}}
\put(66.25,39.75){\makebox(0,0)[cc]{$y$}}
\put(59.75,40.75){\makebox(0,0)[cc]{$\phi$}}
\put(26,74.75){\makebox(0,0)[cc]{$\alpha$}}
\put(60.5,63.5){\makebox(0,0)[cc]{$\beta$}}
\put(26,54){\makebox(0,0)[cc]{$\alpha$}}
\put(39,63){\makebox(0,0)[cc]{$65$}}
\put(66.5,59.75){\makebox(0,0)[cc]{$50$}}
\put(75.25,86.5){\makebox(0,0)[lc]{$x = $}}
\put(75.25,76.5){\makebox(0,0)[lc]{$\alpha = $}}
\put(75.25,33.75){\makebox(0,0)[lc]{$\cos \phi = $}}
\put(75.25,43.5){\makebox(0,0)[lc]{$\sin \theta = $}}
\put(75,65.75){\makebox(0,0)[lc]{$\alpha = $}}
\put(75,55){\makebox(0,0)[lc]{$\beta = $}}
\put(39.75,84.5){\makebox(0,0)[cc]{$x$}}
\put(66.5,81.75){\makebox(0,0)[cc]{$35$}}
\put(58.25,85.25){\makebox(0,0)[cc]{$75 \degrees$}}
\end{picture}
\end{center}

Are your results here consistent with our expansion formulas?

\vspace{0.3cm}

\begin{center}
%TeXCAD Picture [prob7.15_06.pic]. Options:
%\grade{\on}
%\emlines{\off}
%\epic{\off}
%\beziermacro{\on}
%\reduce{\on}
%\snapping{\off}
%\quality{8.00}
%\graddiff{0.01}
%\snapasp{1}
%\zoom{4.0000}
\unitlength 1mm % = 2.85pt
\linethickness{0.4pt}
\ifx\plotpoint\undefined\newsavebox{\plotpoint}\fi % GNUPLOT compatibility
\begin{picture}(69.25,24.25)(0,77)
%\emline(18,83.75)(37.25,98)
\multiput(18,83.75)(.045508274,.033687943){423}{\line(1,0){.045508274}}
%\end
%\emline(37.25,98)(53,85)
\multiput(37.25,98)(.040803109,-.033678756){386}{\line(1,0){.040803109}}
%\end
\put(18,83.5){\line(1,0){36.25}}
\put(54.25,83.5){\line(1,0){.75}}
%\emline(55,83.5)(53,85)
\multiput(55,83.5)(-.0444444,.0333333){45}{\line(-1,0){.0444444}}
%\end
\put(10.25,101.25){\makebox(0,0)[cc]{6)}}
\put(24.75,92.25){\makebox(0,0)[cc]{$40$}}
\put(50.25,91.75){\makebox(0,0)[cc]{$40$}}
\put(36.5,80.25){\makebox(0,0)[cc]{$x$}}
\put(37,94.75){\makebox(0,0)[cc]{$\theta$}}
\put(24.25,85.25){\makebox(0,0)[cc]{$35 \degrees$}}
\put(69.25,100){\makebox(0,0)[cc]{$x = $}}
\put(69.25,89.5){\makebox(0,0)[cc]{$\theta = $}}
\end{picture}

\vspace{0.3cm}

%TeXCAD Picture [prob7.15_07.pic]. Options:
%\grade{\on}
%\emlines{\off}
%\epic{\off}
%\beziermacro{\on}
%\reduce{\on}
%\snapping{\off}
%\quality{8.00}
%\graddiff{0.01}
%\snapasp{1}
%\zoom{4.0000}
\unitlength 1mm % = 2.85pt
\linethickness{0.4pt}
\ifx\plotpoint\undefined\newsavebox{\plotpoint}\fi % GNUPLOT compatibility
\begin{picture}(98.75,44.25)(0,77)
%\emline(16.25,87.25)(77.25,121.25)
\multiput(16.25,87.25)(.060515873,.0337301587){1008}{\line(1,0){.060515873}}
%\end
%\emline(77.25,121.25)(97.25,87.75)
\multiput(77.25,121.25)(.033726813,-.056492411){593}{\line(0,-1){.056492411}}
%\end
%\emline(97.25,87.75)(98.25,86)
\multiput(97.25,87.75)(.033333,-.058333){30}{\line(0,-1){.058333}}
%\end
\put(98.25,86){\line(-1,0){84}}
%\emline(14.25,86)(16.5,87.5)
\multiput(14.25,86)(.05,.0333333){45}{\line(1,0){.05}}
%\end
%\emline(77.25,121)(55.75,86.5)
\multiput(77.25,121)(-.03369906,-.054075235){638}{\line(0,-1){.054075235}}
%\end
%\emline(56,87)(55.75,86.5)
\multiput(56,87)(-.03125,-.0625){8}{\line(0,-1){.0625}}
%\end
\put(9.25,119.25){\makebox(0,0)[cc]{7)}}
\put(24.75,88.25){\makebox(0,0)[cc]{$\theta$}}
\put(51.25,88.5){\makebox(0,0)[cc]{$\gamma$}}
\put(60.75,88.25){\makebox(0,0)[cc]{$\delta$}}
\put(77.5,114){\makebox(0,0)[cc]{$\beta$}}
\put(68.5,113){\makebox(0,0)[cc]{$\alpha$}}
\put(91,89){\makebox(0,0)[cc]{$\phi$}}
\put(42.5,105.75){\makebox(0,0)[cc]{$4$}}
\put(92.5,101.75){\makebox(0,0)[cc]{$3$}}
\put(36.75,83.5){\makebox(0,0)[cc]{$x$}}
\put(78.5,82.75){\makebox(0,0)[cc]{$y$}}
\put(61.75,78.75){\makebox(0,0)[cc]{$5$}}
\put(33.25,84.5){\vector(-1,0){19}}
\put(38.75,84.5){\vector(1,0){14.75}}
\put(75,84){\vector(-1,0){18.75}}
\put(81.25,84){\vector(1,0){16.75}}
\put(66,79.75){\vector(1,0){32.75}}
\put(57,80){\vector(-1,0){43}}
\end{picture}
\begin{align*}
       x & =   &      y & =   &  \theta & =   &    \phi & =  \\
  \alpha & =   &  \beta & =   &  \gamma & =   &  \delta & =  \\
\end{align*}

\end{center}


Can you write each of the angles in terms of $\theta$ only?


\begin{center}
%TeXCAD Options
%\grade{\on}
%\emlines{\off}
%\epic{\off}
%\beziermacro{\on}
%\reduce{\on}
%\snapping{\off}
%\quality{8.00}
%\graddiff{0.01}
%\snapasp{1}
%\zoom{4.0000}
\unitlength 1mm % = 2.85pt
\linethickness{0.4pt}
\ifx\plotpoint\undefined\newsavebox{\plotpoint}\fi % GNUPLOT compatibility
\begin{picture}(80.5,45.75)(0,90)
%\emline(41.75,135.75)(12.5,97.25)
\multiput(41.75,135.75)(-.0337370242,-.0444059977){867}{\line(0,-1){.0444059977}}
%\end
\put(12.5,97.25){\line(1,0){60.5}}
%\emline(73,97.25)(42,135.75)
\multiput(73,97.25)(-.0337323177,.0418933624){919}{\line(0,1){.0418933624}}
%\end
\put(21.5,117.5){\makebox(0,0)[cc]{$x$}}
\put(60.75,117.5){\makebox(0,0)[cc]{$x$}}
\put(42.25,92.75){\makebox(0,0)[cc]{$225$}}
\put(41.75,130.25){\makebox(0,0)[cc]{$\theta$}}
\put(18.25,100.25){\makebox(0,0)[cc]{$65 \degrees$}}
\put(80,132){\makebox(0,0)[cc]{$x = $}}
\put(80,112.25){\makebox(0,0)[cc]{$\theta = $}}
\put(13,135){\makebox(0,0)[cc]{8)}}
\end{picture}
\end{center}


Use the following figure for questions 9 through 12.

\begin{center}
%TeXCAD Options
%\grade{\on}
%\emlines{\off}
%\epic{\off}
%\beziermacro{\on}
%\reduce{\on}
%\snapping{\off}
%\quality{8.00}
%\graddiff{0.01}
%\snapasp{1}
%\zoom{4.0000}
\unitlength 1mm % = 2.85pt
\linethickness{0.4pt}
\ifx\plotpoint\undefined\newsavebox{\plotpoint}\fi % GNUPLOT compatibility
\begin{picture}(84,43.25)(0,96)
%\emline(55.75,132)(35.25,101)
\multiput(55.75,132)(-.033717105,-.050986842){608}{\line(0,-1){.050986842}}
%\end
\put(35.25,101){\line(1,0){48.75}}
%\emline(84,101)(59.25,136.5)
\multiput(84,101)(-.033719346,.0483651226){734}{\line(0,1){.0483651226}}
%\end
%\emline(59.25,136.5)(55.75,132)
\multiput(59.25,136.5)(-.03365385,-.04326923){104}{\line(0,-1){.04326923}}
%\end
\put(59.25,136.25){\line(0,-1){35}}
\put(55.75,104.75){\line(0,-1){3.5}} 
\put(56,104.75){\line(1,0){3}}
\put(35.25,101){\line(-1,0){19.75}}
%\emline(15.5,101.25)(59,125.5)
\multiput(15.5,101.25)(.0605006954,.0337273992){719}{\line(1,0){.0605006954}}
%\end
\put(60.25,139.25){\makebox(0,0)[cc]{U}}
\put(61.75,125.75){\makebox(0,0)[cc]{T}}
\put(14.75,97.75){\makebox(0,0)[cc]{P}}
\put(34.5,97.75){\makebox(0,0)[cc]{Q}}
\put(60,97.75){\makebox(0,0)[cc]{R}}
\put(82.75,97.75){\makebox(0,0)[cc]{S}}
\put(22.25,103){\makebox(0,0)[cc]{$\phi$}}
\put(39.25,103.25){\makebox(0,0)[cc]{$\theta$}}
\end{picture}
\end{center}

\begin{enumerate}
  \item[ 9)] If $QR=40$ and $\theta = 50 \degrees$, find $QU$.
  \item[10)] If $PT=QU$ and $\phi = 35 \degrees$, what is $PR$? (assuming $QR=40$)
  \item[11)] If $PT=QU$ and $PQ=UT$, show $QR=RT$.
  \item[12)] If $PT=QU$ and $QR=RT$, show $PQ=UT$.
\end{enumerate}

\vspace{0.3cm}

\begin{enumerate}
  \item[13)] In the following diagram, suppose $AB$ is an arc of a circle with center $O$ and radius $OA$.
\end{enumerate}

\begin{center}
%TeXCAD Options
%\grade{\on}
%\emlines{\off}
%\epic{\off}
%\beziermacro{\on}
%\reduce{\on}
%\snapping{\off}
%\quality{8.00}
%\graddiff{0.01}
%\snapasp{1}
%\zoom{4.0000}
\unitlength 1mm % = 2.85pt
\linethickness{0.4pt}
\ifx\plotpoint\undefined\newsavebox{\plotpoint}\fi % GNUPLOT compatibility
\begin{picture}(79.5,50.25)(0,70)
\put(23.5,114.75){\line(1,0){52.5}}
%\emline(76,114.75)(51.75,78.75)
\multiput(76,114.75)(-.0337273992,-.050069541){719}{\line(0,-1){.050069541}}
%\end
\put(51.75,78.75){\line(-3,-4){3}}
%\emline(48.75,74.75)(23.5,114.75)
\multiput(48.75,74.75)(-.0337116155,.0534045394){749}{\line(0,1){.0534045394}}
%\end
\qbezier(23.5,114.75)(51.75,140.25)(76,114.75)
\qbezier(45.5,80)(49.13,82.25)(52.25,79.5)
\put(48.75,78.5){\makebox(0,0)[cc]{$\theta$}}
\put(20.75,114.75){\makebox(0,0)[cc]{A}}
\put(79.5,114.75){\makebox(0,0)[cc]{B}}
\put(49.5,72){\makebox(0,0)[cc]{O}}
\end{picture}
\end{center}

The first known trigonometric table was constructed by Hipparchus of Nicaea (you knew that, didn't you?), and in his table he gave, essentially, values of the ratio $AB/AO$ for various value of $\theta$. Can you relate this ratio to the sine function?

\vspace{0.3cm}

For problems 14 through 17, find $x$ in each triangle, and leave your answers in radical form; i.e., no conservative decimal approximations.

\begin{center}
%TeXCAD Picture [prob7.15_14-17.pic]. Options:
%\grade{\on}
%\emlines{\off}
%\epic{\off}
%\beziermacro{\on}
%\reduce{\on}
%\snapping{\off}
%\quality{8.00}
%\graddiff{0.01}
%\snapasp{1}
%\zoom{4.0000}
\unitlength 1mm % = 2.85pt
\linethickness{0.4pt}
\ifx\plotpoint\undefined\newsavebox{\plotpoint}\fi % GNUPLOT compatibility
\begin{picture}(106.5,131.25)(0,0)
\put(56.5,123){\line(-1,-1){23.75}}
\put(32.75,99.25){\line(1,0){31.25}}
\put(64,99.25){\line(0,1){32}}
%\emline(64,131.25)(56.75,123.25)
\multiput(64,131.25)(-.03372093,-.0372093){215}{\line(0,-1){.0372093}}
%\end
\put(60.25,103){\line(0,-1){3.5}}
\put(60.5,103){\line(1,0){3.5}}
\put(9,4){\line(1,0){71}}
\put(8.75,69){\line(1,0){71}}
\put(8.75,37.75){\line(1,0){71}}
\put(80,4){\line(0,1){24.75}}
\put(79.75,69){\line(0,1){24.75}}
\put(79.75,37.75){\line(0,1){24.75}}
%\emline(80,28.75)(9,4.25)
\multiput(80,28.75)(-.0976616231,-.0337001376){727}{\line(-1,0){.0976616231}}
%\end
%\emline(79.75,93.75)(8.75,69.25)
\multiput(79.75,93.75)(-.0976616231,-.0337001376){727}{\line(-1,0){.0976616231}}
%\end
%\emline(79.75,62.5)(8.75,38)
\multiput(79.75,62.5)(-.0976616231,-.0337001376){727}{\line(-1,0){.0976616231}}
%\end
\put(76,8.25){\line(0,-1){4}}
\put(75.75,73.25){\line(0,-1){4}}
\put(75.75,42){\line(0,-1){4}}
\put(76.25,8.25){\line(1,0){3.75}}
\put(76,73.25){\line(1,0){3.75}}
\put(76,42){\line(1,0){3.75}}
\put(10.5,128.5){\makebox(0,0)[cc]{14)}}
\put(9.75,91.25){\makebox(0,0)[cc]{15)}}
\put(8.75,59.75){\makebox(0,0)[cc]{16)}}
\put(8,26){\makebox(0,0)[cc]{17)}}
\put(75.5,23){\makebox(0,0)[cc]{$60 \degrees$}}
\put(75.25,88){\makebox(0,0)[cc]{$60 \degrees$}}
\put(75.25,56.75){\makebox(0,0)[cc]{$60 \degrees$}}
\put(43,102.25){\makebox(0,0)[cc]{$45 \degrees$}}
\put(30.75,72.5){\makebox(0,0)[cc]{$30 \degrees$}}
\put(44.75,34){\makebox(0,0)[cc]{$5 \sqrt{3}$}}
\put(43.5,114){\makebox(0,0)[cc]{$24$}}
\put(83.5,14.5){\makebox(0,0)[cc]{$11$}}
\put(83.25,79.5){\makebox(0,0)[cc]{$20$}}
\put(83.25,48.25){\makebox(0,0)[cc]{$x$}}
\put(67,113.75){\makebox(0,0)[cc]{$x$}}
\put(44.5,19.5){\makebox(0,0)[cc]{$x$}}
\put(44.25,84.5){\makebox(0,0)[cc]{$x$}}
\put(106,113.25){\makebox(0,0)[cc]{$x = $}}
\put(106,83.75){\makebox(0,0)[cc]{$x = $}}
\put(106,48.5){\makebox(0,0)[cc]{$x = $}}
\put(106,15.5){\makebox(0,0)[cc]{$x = $}}
\end{picture}
\end{center}

For questions 18 through 23, find the sine and the cosine of the angle given. First do this with your calculator, and then try to find a way of doing it without the calculator.

\begin{enumerate}
  \item[18)] $\sin 15 \degrees = $ 

           $\cos 15 \degrees = $
  \item[19)] $\sin 22 \degrees 30' = $

           $\cos 22 \degrees 30' = $
  \item[20)] $\sin 75 \degrees = $

           $\cos 75 \degrees = $
  \item[21)] $\sin 105 \degrees = $

           $\cos 105 \degrees = $ 
  \item[22)] $\sin 37 \degrees 30' = $ 

           $\cos 37 \degrees 30' = $
  \item[23)] $\sin 52 \degrees 30' = $ 

           $\cos 52 \degrees 30' = $
\end{enumerate}

For questions 24 through 29, find $\theta$ to four decimal places.

\begin{enumerate}
  \item[24)] $\sin \theta = 0.514$
  \item[25)] $\sin \theta = \frac{1}{4}$
  \item[26)] $\cos \theta = 0.859$
  \item[27)] $\cos \theta = \frac{2}{3}$
  \item[28)] $\sin \theta = 0.148$
  \item[29)] $\cos \theta = \sqrt{ \frac{5}{9} }$
\end{enumerate}


\begin{enumerate}
  \item[30)] Given $m \angm PQR = 90 \degrees$, $m \angm P = 55 \degrees 30'$, and $PR=100$, find these using the picture below.
  \begin{align*}
             SR &=  &           PS &=  &  PQ &=  &  QR &= \\
    m \angm PQS &=  &  m \angm RQS &=  &  QS &=
  \end{align*}
\end{enumerate}

\begin{center}
%TeXCAD Options
%\grade{\on}
%\emlines{\off}
%\epic{\off}
%\beziermacro{\on}
%\reduce{\on}
%\snapping{\off}
%\quality{8.00}
%\graddiff{0.01}
%\snapasp{1}
%\zoom{4.0000}
\unitlength 1mm % = 2.85pt
\linethickness{0.4pt}
\ifx\plotpoint\undefined\newsavebox{\plotpoint}\fi % GNUPLOT compatibility
\begin{picture}(51.5,28.25)(0,112)
%\emline(37.75,137)(13,117.5)
\multiput(37.75,137)(-.042820069,-.033737024){578}{\line(-1,0){.042820069}}
%\end
\put(13,117.5){\line(1,0){38.25}}
%\emline(51.25,117.5)(38,137)
\multiput(51.25,117.5)(-.033715013,.049618321){393}{\line(0,1){.049618321}}
%\end
\put(38,137){\line(0,-1){19.25}} 
\put(35.5,120.5){\line(0,-1){2.5}}
\put(35.5,120.25){\line(1,0){2.5}}
%\emline(35.5,134.75)(37,132.75)
\multiput(35.5,134.75)(.0333333,-.0444444){45}{\line(0,-1){.0444444}}
%\end
%\emline(37,132.75)(39.5,134.75)
\multiput(37,132.75)(.0416667,.0333333){60}{\line(1,0){.0416667}}
%\end
\put(37.5,140.25){\makebox(0,0)[cc]{$Q$}}
\put(13,114){\makebox(0,0)[cc]{$P$}}
\put(37.75,114){\makebox(0,0)[cc]{$S$}}
\put(51.5,114){\makebox(0,0)[cc]{$R$}}
\end{picture}
\end{center}


\begin{enumerate}
  \item[31)] Given $PQ=1$, $SQ=1$, $RS=1$, $SP=x$, and $RQ=1+x$, find these using the picture below.
  \begin{align*}
    \delta &=  &  \sigma &=  &  \theta &=  &  x &= \\
    \alpha &=  &  \beta  &=  &  \gamma &=
  \end{align*}
\end{enumerate}

\begin{center}
%TeXCAD Options
%\grade{\on}
%\emlines{\off}
%\epic{\off}
%\beziermacro{\on}
%\reduce{\on}
%\snapping{\off}
%\quality{8.00}
%\graddiff{0.01}
%\snapasp{1}
%\zoom{4.0000}
\unitlength 1mm % = 2.85pt
\linethickness{0.4pt}
\ifx\plotpoint\undefined\newsavebox{\plotpoint}\fi % GNUPLOT compatibility
\begin{picture}(84,54.25)(0,93)
%\emline(12.5,101)(56.25,141.75)
\multiput(12.5,101)(.0362168874,.0337334437){1208}{\line(1,0){.0362168874}}
%\end
%\emline(56.25,141.75)(81,99)
\multiput(56.25,141.75)(.033719346,-.0582425068){734}{\line(0,-1){.0582425068}}
%\end
\put(81,99){\line(-1,0){70.75}}
%\emline(10.25,99)(12.75,101)
\multiput(10.25,99)(.0416667,.0333333){60}{\line(1,0){.0416667}}
%\end
%\emline(81,99.25)(38.75,125.5)
\multiput(81,99.25)(-.0542362003,.0336970475){779}{\line(-1,0){.0542362003}}
%\end
\put(25.75,116.5){\vector(1,1){9.75}}
%\vector(20.5,112.75)(8.5,101)
\put(8.5,101){\vector(-1,-1){.07}}\multiput(20.5,112.75)(-.034383954,-.033667622){349}{\line(-1,0){.034383954}}
%\end
\put(22.75,114.5){\makebox(0,0)[cc]{$x$}}
\put(34.75,128.5){\makebox(0,0)[cc]{$S$}}
\put(43,133.5){\makebox(0,0)[cc]{$1$}}
\put(56.5,144.75){\makebox(0,0)[cc]{$R$}}
\put(72.5,123.25){\makebox(0,0)[cc]{$1+x$}}
\put(57.75,116.25){\makebox(0,0)[cc]{$1$}}
\put(74,106){\makebox(0,0)[cc]{$\beta$}}
\put(70.5,101.75){\makebox(0,0)[cc]{$\alpha$}}
\put(38.75,120.75){\makebox(0,0)[cc]{$\delta$}}
\put(43.75,125.5){\makebox(0,0)[cc]{$\sigma$}}
\put(16.75,101.25){\makebox(0,0)[cc]{$\theta$}}
\put(9.25,95){\makebox(0,0)[cc]{$P$}}
\put(47.75,95){\makebox(0,0)[cc]{$1$}}
\put(84,95){\makebox(0,0)[cc]{$Q$}}
\end{picture}
\end{center}

\begin{enumerate}
  \item[32)] A physics problem.

             Light changes velocity (which involves both speed and direction) when it goes from one medium (like air) into another medium (like glass). The following diagram and formula illustrate something called ``Snell's Law of Optics'' (physics departments are called upon to enforce such laws), which tells us how the direction and speed of light is changed when it goes from one medium into another.
\end{enumerate}

\begin{center}
Snell's Law says:

%TeXCAD Picture [prob7.15_32.pic]. Options:
%\grade{\on}
%\emlines{\off}
%\epic{\off}
%\beziermacro{\on}
%\reduce{\on}
%\snapping{\off}
%\quality{8.00}
%\graddiff{0.01}
%\snapasp{1}
%\zoom{4.0000}
\unitlength 1mm % = 2.85pt
\linethickness{0.4pt}
\ifx\plotpoint\undefined\newsavebox{\plotpoint}\fi % GNUPLOT compatibility
\begin{picture}(101.5,63.25)(0,0)
\put(14.75,30.75){\line(1,0){86.75}}
%\dashline{1}(55.75,63)(54.25,3.25)
\put(55.68,62.93){\line(0,-1){.996}}
\put(55.63,60.94){\line(0,-1){.996}}
\put(55.58,58.95){\line(0,-1){.996}}
\put(55.53,56.95){\line(0,-1){.996}}
\put(55.48,54.96){\line(0,-1){.996}}
\put(55.43,52.97){\line(0,-1){.996}}
\put(55.38,50.98){\line(0,-1){.996}}
\put(55.33,48.99){\line(0,-1){.996}}
\put(55.28,47){\line(0,-1){.996}}
\put(55.23,45){\line(0,-1){.996}}
\put(55.18,43.01){\line(0,-1){.996}}
\put(55.13,41.02){\line(0,-1){.996}}
\put(55.08,39.03){\line(0,-1){.996}}
\put(55.03,37.04){\line(0,-1){.996}}
\put(54.98,35.05){\line(0,-1){.996}}
\put(54.93,33.05){\line(0,-1){.996}}
\put(54.88,31.06){\line(0,-1){.996}}
\put(54.83,29.07){\line(0,-1){.996}}
\put(54.78,27.08){\line(0,-1){.996}}
\put(54.73,25.09){\line(0,-1){.996}}
\put(54.68,23.1){\line(0,-1){.996}}
\put(54.63,21.1){\line(0,-1){.996}}
\put(54.58,19.11){\line(0,-1){.996}}
\put(54.53,17.12){\line(0,-1){.996}}
\put(54.48,15.13){\line(0,-1){.996}}
\put(54.43,13.14){\line(0,-1){.996}}
\put(54.38,11.15){\line(0,-1){.996}}
\put(54.33,9.15){\line(0,-1){.996}}
\put(54.28,7.16){\line(0,-1){.996}}
\put(54.23,5.17){\line(0,-1){.996}}
%\end
%\vector(55.25,30.5)(62.75,4.75)
\put(62.75,4.75){\vector(1,-3){.07}}\multiput(55.25,30.5)(.03363229,-.11547085){223}{\line(0,-1){.11547085}}
%\end
%\vector[middle](14.5,56.75)(54.75,31)
\put(34.63,43.88){\vector(3,-2){.07}}\multiput(14.5,56.75)(.0526832461,-.0337041885){764}{\line(1,0){.0526832461}}
%\end
\qbezier(49.75,34.25)(51.75,36.75)(54.75,36.25)
%\qbezvec(68,62)(60.38,63.25)(57.25,59.5)
\put(57.25,59.5){\vector(-1,-1){.07}}\qbezier(68,62)(60.38,63.25)(57.25,59.5)
%\end
\put(85,58.5){\makebox(0,0)[cc]{a line perpendicular}}
\put(85,53.5){\makebox(0,0)[cc]{to the media interface,}}
\put(85,48.75){\makebox(0,0)[cc]{called the normal}}
\put(81.5,33.75){\makebox(0,0)[cc]{medium 1}}
\put(81.5,27.5){\makebox(0,0)[cc]{medium 2}}
\put(21.5,18.25){\makebox(0,0)[cc]{$\frac{\sin \alpha}{\sin \beta} = \frac{s_{1}}{s_{2}}$}}
\put(50.5,37.5){\makebox(0,0)[cc]{$\alpha$}}
\put(57,19){\makebox(0,0)[cc]{$\beta$}}
%\qbezvec(54.75,22.75)(56.88,21.63)(57.5,23)
\put(57.5,23){\vector(1,2){.07}}\qbezier(54.75,22.75)(56.88,21.63)(57.5,23)
%\end
\end{picture}

where $s_1 =$ speed of light in medium 1

and $s_2 =$ speed of light in medium 2
\end{center}


Suppose you have a pane of glass with air on both sides of the glass, and suppose that light passes from the air through the glass and into the air again on the other side of the glass, as in the picture.

\begin{center}
%TeXCAD Options
%\grade{\on}
%\emlines{\off}
%\epic{\off}
%\beziermacro{\on}
%\reduce{\on}
%\snapping{\off}
%\quality{8.00}
%\graddiff{0.01}
%\snapasp{1}
%\zoom{4.0000}
\unitlength 1mm % = 2.85pt
\linethickness{0.4pt}
\ifx\plotpoint\undefined\newsavebox{\plotpoint}\fi % GNUPLOT compatibility
\begin{picture}(81.75,42.25)(0,96)
\put(17.25,124.25){\line(1,0){64.25}}
\put(81.75,113.5){\line(-1,0){65.25}}
%\dashline{1}(69.25,138.25)(69.25,99.25)
\put(69.18,138.18){\line(0,-1){.975}}
\put(69.18,136.23){\line(0,-1){.975}}
\put(69.18,134.28){\line(0,-1){.975}}
\put(69.18,132.33){\line(0,-1){.975}}
\put(69.18,130.38){\line(0,-1){.975}}
\put(69.18,128.43){\line(0,-1){.975}}
\put(69.18,126.48){\line(0,-1){.975}}
\put(69.18,124.53){\line(0,-1){.975}}
\put(69.18,122.58){\line(0,-1){.975}}
\put(69.18,120.63){\line(0,-1){.975}}
\put(69.18,118.68){\line(0,-1){.975}}
\put(69.18,116.73){\line(0,-1){.975}}
\put(69.18,114.78){\line(0,-1){.975}}
\put(69.18,112.83){\line(0,-1){.975}}
\put(69.18,110.88){\line(0,-1){.975}}
\put(69.18,108.93){\line(0,-1){.975}}
\put(69.18,106.98){\line(0,-1){.975}}
\put(69.18,105.03){\line(0,-1){.975}}
\put(69.18,103.08){\line(0,-1){.975}}
\put(69.18,101.13){\line(0,-1){.975}}
%\end
%\vector[middle](52,136.75)(69,124.5)
\put(60.5,130.63){\vector(4,-3){.07}}\multiput(52,136.75)(.046703297,-.033653846){364}{\line(1,0){.046703297}}
%\end
\put(41.25,126.5){\makebox(0,0)[cc]{air}}
\put(41.25,117.75){\makebox(0,0)[cc]{glass}}
\put(41.5,109.5){\makebox(0,0)[cc]{air}}
\put(67,128.75){\makebox(0,0)[cc]{$\alpha$}}
\end{picture}
\end{center}

If the light entering the glass (the incident light) makes an angle $\alpha$ with the normal (as in the picture), what is the angle that the emerging light makes with the normal?

Also, what if your glass was very, very special, so that it took light 6 months to travel through $\frac{1}{4}$'' of your glass? What could you do with such glass? Let your imagination soar!

\vspace{0.3cm}

%TeXCAD Picture [prob7.15_33.pic]. Options:
%\grade{\on}
%\emlines{\off}
%\epic{\off}
%\beziermacro{\on}
%\reduce{\on}
%\snapping{\off}
%\quality{8.00}
%\graddiff{0.01}
%\snapasp{1}
%\zoom{4.0000}
\unitlength 1mm % = 2.85pt
\linethickness{0.4pt}
\ifx\plotpoint\undefined\newsavebox{\plotpoint}\fi % GNUPLOT compatibility
\begin{picture}(98.25,42)(0,0)
%\emline(33.5,33.75)(17,5.5)
\multiput(33.5,33.75)(-.033673469,-.057653061){490}{\line(0,-1){.057653061}}
%\end
\put(17,5.5){\line(1,0){15.75}}
\put(42.75,5.25){\line(1,0){14.75}}
%\emline(57.5,5.25)(36.75,38.75)
\multiput(57.5,5.25)(-.033739837,.054471545){615}{\line(0,1){.054471545}}
%\end
%\emline(36.75,38.75)(33.25,33.25)
\multiput(36.75,38.75)(-.03365385,-.05288462){104}{\line(0,-1){.05288462}}
%\end
\qbezier(6.75,20.75)(4.38,14.88)(15.5,21.5)
\qbezier(15.75,21.5)(21.88,27.38)(23.5,24.75)
\qbezier(23.5,24.75)(31.5,25.25)(37.5,31.75)
\qbezier(37.75,31.25)(46.13,35.13)(51,30.5)
\qbezier(51,30.5)(58.75,31.63)(63.5,35.25)
\qbezier(63.5,35)(72.75,37.5)(84,36)
\qbezier(4.25,6.75)(7.5,11)(11.75,10.25)
\qbezier(12,10.25)(17.75,12.75)(23.5,13.25)
\qbezier(23.75,13)(28.38,12.38)(35.5,14.25)
\qbezier(35.75,14.25)(51,20.63)(50.25,18.5)
\qbezier(50.5,18.5)(56.13,19.13)(65.25,17.25)
\qbezier(65.75,17.5)(70.88,14.5)(79.5,19.5)
\qbezier(79.5,19)(80.88,20.88)(86.75,24.25)
\put(38,5.25){\makebox(0,0)[cc]{$100$}}
\put(37,42){\makebox(0,0)[cc]{$A$}}
\put(13.75,3.75){\makebox(0,0)[cc]{$C$}}
\put(60.75,4){\makebox(0,0)[cc]{$B$}}
\put(22.25,8){\makebox(0,0)[cc]{$80 \degrees$}}
\put(50.75,8){\makebox(0,0)[cc]{$75 \degrees$}}
\put(69.75,25.5){\makebox(0,0)[cc]{river}}
\put(98.25,27.25){\makebox(0,0)[cc]{$AB = $}}
\put(7,35){\makebox(0,0)[cc]{33)}}
\end{picture}

%TeXCAD Picture [prob7.15_34.pic]. Options:
%\grade{\on}
%\emlines{\off}
%\epic{\off}
%\beziermacro{\on}
%\reduce{\on}
%\snapping{\off}
%\quality{8.00}
%\graddiff{0.01}
%\snapasp{1}
%\zoom{3.3636}
\unitlength 1mm % = 2.85pt
\linethickness{0.4pt}
\ifx\plotpoint\undefined\newsavebox{\plotpoint}\fi % GNUPLOT compatibility
\begin{picture}(83.75,33.98)(0,0)
\put(8.5,33.98){\makebox(0,0)[cc]{34)}}
\put(17.25,4.98){\framebox(50.5,25.25)[]{}}
%\emline(17.25,29.98)(67.75,17.98)
\multiput(17.25,29.98)(.141853933,-.033707865){356}{\line(1,0){.141853933}}
%\end
\put(64.5,29.98){\line(0,-1){2.75}}
\put(64.5,27.23){\line(1,0){3.25}}
\put(64.25,8.23){\line(0,-1){3}}
\put(64.25,8.23){\line(1,0){3.5}}
\put(17.25,8.73){\line(1,0){3.75}}
\put(21,8.73){\line(0,-1){3.5}}
\put(34,28.23){\makebox(0,0)[cc]{$20 \degrees$}}
\put(13,16.73){\makebox(0,0)[cc]{$150$}}
\put(39.75,1.23){\makebox(0,0)[cc]{$300$}}
\put(83.75,19.98){\makebox(0,0)[cc]{$h = $}}
\put(69.75,11.23){\makebox(0,0)[cc]{$h$}}
\put(69.75,9.73){\vector(0,-1){4.5}}
\put(69.87,14.57){\vector(0,1){4.76}}
\end{picture}

%TeXCAD Picture [prob7.15_35.pic]. Options:
%\grade{\on}
%\emlines{\off}
%\epic{\off}
%\beziermacro{\on}
%\reduce{\on}
%\snapping{\off}
%\quality{8.00}
%\graddiff{0.01}
%\snapasp{1}
%\zoom{2.3784}
\unitlength 1mm % = 2.85pt
\linethickness{0.4pt}
\ifx\plotpoint\undefined\newsavebox{\plotpoint}\fi % GNUPLOT compatibility
\begin{picture}(84.5,46.16)(0,0)
\put(24,7.66){\line(1,0){31.75}}
%\emline(55.75,7.66)(71.5,46.16)
\multiput(55.75,7.66)(.03372591,.082441113){467}{\line(0,1){.082441113}}
%\end
%\emline(71.5,46.16)(23.75,7.66)
\multiput(71.5,46.16)(-.0418126095,-.0337127846){1142}{\line(-1,0){.0418126095}}
%\end
\put(83.99,39.08){\makebox(0,0)[cc]{$\beta = $}}
\put(84.5,30.16){\makebox(0,0)[cc]{$\alpha = $}}
\put(84.5,21.66){\makebox(0,0)[cc]{$x = $}}
\put(69.5,12.66){\makebox(0,0)[cc]{$\alpha > 90 \degrees$}}
\put(30.5,10.16){\makebox(0,0)[cc]{$36 \degrees$}}
\put(52,10.16){\makebox(0,0)[cc]{$\alpha$}}
\put(65.25,37.91){\makebox(0,0)[cc]{$\beta$}}
\put(65,24.66){\makebox(0,0)[cc]{$6$}}
\put(43.03,28.16){\makebox(0,0)[cc]{$x$}}
\put(10.25,43.16){\makebox(0,0)[cc]{35)}}
\put(38.75,5.66){\makebox(0,0)[cc]{$8$}}
\end{picture}

%TeXCAD Picture [prob7.15_36.pic]. Options:
%\grade{\on}
%\emlines{\off}
%\epic{\off}
%\beziermacro{\on}
%\reduce{\on}
%\snapping{\off}
%\quality{8.00}
%\graddiff{0.01}
%\snapasp{1}
%\zoom{4.0000}
\unitlength 1mm % = 2.85pt
\linethickness{0.4pt}
\ifx\plotpoint\undefined\newsavebox{\plotpoint}\fi % GNUPLOT compatibility
\begin{picture}(71.25,26.75)(0,0)
%\emline(37.5,26.75)(21.75,4)
\multiput(37.5,26.75)(-.03372591,-.048715203){467}{\line(0,-1){.048715203}}
%\end
\put(21.75,4){\line(1,0){34.5}}
%\emline(56.25,4)(37.75,26.5)
\multiput(56.25,4)(-.033697632,.040983607){549}{\line(0,1){.040983607}}
%\end
\put(37.5,26.5){\line(0,-1){22.25}}
\put(35,6.75){\line(0,-1){2.75}}
\put(35.25,6.75){\line(1,0){2.25}}
\put(11.25,24.5){\makebox(0,0)[cc]{36)}}
\put(25.5,20.5){\makebox(0,0)[cc]{$17 \degrees$}}
\put(47.5,23.25){\makebox(0,0)[cc]{$17 \degrees$}}
\put(52.25,12.25){\makebox(0,0)[cc]{$x$}}
\put(25.25,12.75){\makebox(0,0)[cc]{$x$}}
\put(38.5,1.75){\makebox(0,0)[cc]{$13$}}
\put(71.25,20.25){\makebox(0,0)[cc]{$x = $}}
%\qbezvec(27.25,18.75)(34.38,17.38)(36,21.5)
\put(36,21.5){\vector(1,2){.07}}\qbezier(27.25,18.75)(34.38,17.38)(36,21.5)
%\end
%\qbezvec(48.5,21.75)(41.88,15.75)(38.75,21.75)
\put(38.75,21.75){\vector(-2,3){.07}}\qbezier(48.5,21.75)(41.88,15.75)(38.75,21.75)
%\end
\end{picture}

% Problem 7.15, Question 37
%TeXCAD Options
%\grade{\on}
%\emlines{\off}
%\epic{\off}
%\beziermacro{\on}
%\reduce{\on}
%\snapping{\off}
%\quality{8.00}
%\graddiff{0.01}
%\snapasp{1}
%\zoom{4.0000}
\unitlength 1mm % = 2.85pt
\linethickness{0.4pt}
\ifx\plotpoint\undefined\newsavebox{\plotpoint}\fi % GNUPLOT compatibility
\begin{picture}(89.5,26.75)(0,119)
%\emline(16.25,123.5)(47.25,145.25)
\multiput(16.25,123.5)(.048062016,.03372093){645}{\line(1,0){.048062016}}
%\end
%\emline(47.25,145.25)(63,125.25)
\multiput(47.25,145.25)(.03372591,-.042826552){467}{\line(0,-1){.042826552}}
%\end
%\emline(63,125.25)(64,124)
\multiput(63,125.25)(.033333,-.041667){30}{\line(0,-1){.041667}}
%\end
%\emline(64,124)(64.25,123.75)
\multiput(64,124)(.03125,-.03125){8}{\line(0,-1){.03125}}
%\end
\put(64.25,123.75){\line(-1,0){47.5}}
\put(47.5,145.25){\line(0,-1){21.5}}
\put(45,126.5){\line(0,-1){2.5}} 
\put(45.25,126.5){\line(1,0){2}}
%\emline(45.25,143.25)(46.75,141.75)
\multiput(45.25,143.25)(.0333333,-.0333333){45}{\line(0,-1){.0333333}}
%\end
%\emline(46.75,141.75)(48.75,143.25)
\multiput(46.75,141.75)(.0444444,.0333333){45}{\line(1,0){.0444444}}
%\end
\put(44.75,134.5){\makebox(0,0)[cc]{h}}
\put(58,134.5){\makebox(0,0)[cc]{z}}
\put(34.75,120.75){\makebox(0,0)[cc]{x}}
\put(55.5,120.75){\makebox(0,0)[cc]{y}}
\put(73,142){\makebox(0,0)[cc]{x = }}
\put(89,142){\makebox(0,0)[cc]{y = }}
\put(73,132){\makebox(0,0)[cc]{z = }}
\put(89,132){\makebox(0,0)[cc]{h = }}
\put(7,145.75){\makebox(0,0)[cc]{37)}}
\put(23.75,126){\makebox(0,0)[cc]{$35 \degrees$}}
\put(57.25,126){\makebox(0,0)[cc]{$55 \degrees$}}
\put(27.5,137){\makebox(0,0)[cc]{100}}
\end{picture}

%TeXCAD Picture [prob7.15_38.pic]. Options:
%\grade{\on}
%\emlines{\off}
%\epic{\off}
%\beziermacro{\on}
%\reduce{\on}
%\snapping{\off}
%\quality{8.00}
%\graddiff{0.01}
%\snapasp{1}
%\zoom{4.0000}
\unitlength 1mm % = 2.85pt
\linethickness{0.4pt}
\ifx\plotpoint\undefined\newsavebox{\plotpoint}\fi % GNUPLOT compatibility
\begin{picture}(88,26.5)(0,0)
%\emline(15,4.25)(38.25,24.75)
\multiput(15,4.25)(.038240132,.033717105){608}{\line(1,0){.038240132}}
%\end
\put(38.25,24.75){\line(0,-1){20.25}}
\put(38.25,4.5){\line(0,-1){1.5}}
\put(38.25,3){\line(-1,0){25.25}}
%\emline(13,3)(15,4.5)
\multiput(13,3)(.0444444,.0333333){45}{\line(1,0){.0444444}}
%\end
\put(38.5,3){\line(1,0){40.5}}
%\emline(79,3)(38,12)
\multiput(79,3)(-.153558052,.033707865){267}{\line(-1,0){.153558052}}
%\end
\put(35,5.75){\line(0,-1){2.25}}
\put(35,3){\line(0,1){.75}}
\put(35.25,5.75){\line(1,0){3.25}}
\put(5.75,26.5){\makebox(0,0)[cc]{38)}}
\put(40,18.75){\makebox(0,0)[cc]{x}}
\put(21.25,14){\makebox(0,0)[cc]{40}}
\put(55.5,11.75){\makebox(0,0)[cc]{40}}
\put(60,5){\makebox(0,0)[cc]{$30 \degrees$}}
\put(24.75,5){\makebox(0,0)[cc]{$48.590378 \degrees$}}
\put(40.5,6.75){\makebox(0,0)[cc]{20}}
\put(88,21){\makebox(0,0)[cc]{x = }}
\put(40,21.5){\vector(0,1){3.75}}
\put(40,16.75){\vector(0,-1){4.75}}
\end{picture}

%TeXCAD Picture [prob7.15_39.pic]. Options:
%\grade{\on}
%\emlines{\off}
%\epic{\off}
%\beziermacro{\on}
%\reduce{\on}
%\snapping{\off}
%\quality{8.00}
%\graddiff{0.01}
%\snapasp{1}
%\zoom{4.0000}
\unitlength 1mm % = 2.85pt
\linethickness{0.4pt}
\ifx\plotpoint\undefined\newsavebox{\plotpoint}\fi % GNUPLOT compatibility
\begin{picture}(87.25,35)(0,0)
%\emline(14.75,9.75)(63,35)
\multiput(14.75,9.75)(.0644192256,.0337116155){749}{\line(1,0){.0644192256}}
%\end
\put(63,35){\line(0,-1){25}}
\put(63,10){\line(0,-1){3.25}}
\put(63,6.75){\line(-1,0){55.25}}
%\emline(7.75,6.75)(15.25,10)
\multiput(7.75,6.75)(.07731959,.03350515){97}{\line(1,0){.07731959}}
%\end
%\emline(62.75,35)(40.75,7)
\multiput(62.75,35)(-.033690658,-.04287902){653}{\line(0,-1){.04287902}}
%\end
\put(59.25,10.25){\line(0,-1){3.25}}
\put(59.5,10.25){\line(1,0){3.5}}
\put(28.25,4.25){\vector(1,0){11.5}}
\put(20.25,4.5){\vector(-1,0){12.75}}
\put(5.75,31.25){\makebox(0,0)[cc]{39)}}
\put(48.75,9.25){\makebox(0,0)[cc]{$62.3 \degrees$}}
\put(24.25,4.25){\makebox(0,0)[cc]{60}}
\put(87.25,21.5){\makebox(0,0)[cc]{x = }}
\put(22.75,9.25){\makebox(0,0)[cc]{$37.1 \degrees$}}
\put(46.5,18.75){\makebox(0,0)[cc]{x}}
\end{picture}

%TeXCAD Picture [prob7.15_40.pic]. Options:
%\grade{\on}
%\emlines{\off}
%\epic{\off}
%\beziermacro{\on}
%\reduce{\on}
%\snapping{\off}
%\quality{8.00}
%\graddiff{0.01}
%\snapasp{1}
%\zoom{4.0000}
\unitlength 1mm % = 2.85pt
\linethickness{0.4pt}
\ifx\plotpoint\undefined\newsavebox{\plotpoint}\fi % GNUPLOT compatibility
\begin{picture}(88.5,29)(0,0)
%\emline(22.25,28.5)(33.25,.5)
\multiput(22.25,28.5)(.033639144,-.085626911){327}{\line(0,-1){.085626911}}
%\end
%\emline(33.25,.5)(88.5,15)
\multiput(33.25,.5)(.128488372,.03372093){430}{\line(1,0){.128488372}}
%\end
%\emline(88.5,15)(22.25,28.75)
\multiput(88.5,15)(-.162377451,.03370098){408}{\line(-1,0){.162377451}}
%\end
\put(6,29){\makebox(0,0)[cc]{40)}}
\put(24.5,14.5){\makebox(0,0)[cc]{x}}
\put(55.25,24.75){\makebox(0,0)[cc]{140}}
\put(54,2.25){\makebox(0,0)[cc]{100}}
\put(74.25,15){\makebox(0,0)[cc]{$40 \degrees$}}
\put(87.75,26.25){\makebox(0,0)[cc]{x = }}
\end{picture}

%TeXCAD Picture [prob7.15_41.pic]. Options:
%\grade{\on}
%\emlines{\off}
%\epic{\off}
%\beziermacro{\on}
%\reduce{\on}
%\snapping{\off}
%\quality{8.00}
%\graddiff{0.01}
%\snapasp{1}
%\zoom{4.0000}
\unitlength 1mm % = 2.85pt
\linethickness{0.4pt}
\ifx\plotpoint\undefined\newsavebox{\plotpoint}\fi % GNUPLOT compatibility
\begin{picture}(88.25,29.5)(0,0)
\put(22.75,2.75){\line(1,0){39.25}}
\put(62,2.75){\line(0,1){25.75}}
\put(62,28.5){\line(-4,-3){24}}
%\emline(22.75,3)(62,15.75)
\multiput(22.75,3)(.103835979,.033730159){378}{\line(1,0){.103835979}}
%\end
%\emline(38,10.5)(32.5,6.25)
\multiput(38,10.5)(-.04365079,-.03373016){126}{\line(-1,0){.04365079}}
%\end
\put(6.75,29.5){\makebox(0,0)[cc]{41)}}
\put(64.5,21.25){\makebox(0,0)[cc]{30}}
\put(88.25,16.25){\makebox(0,0)[cc]{x = }}
\put(45.5,19.75){\makebox(0,0)[cc]{x}}
\put(35.5,4.5){\makebox(0,0)[cc]{$26 \degrees$}}
%\vector(50.5,10.5)(40.75,7.5)
\put(40.75,7.5){\vector(-3,-1){.07}}\multiput(50.5,10.5)(-.1095506,-.0337079){89}{\line(-1,0){.1095506}}
%\end
\put(56,12.25){\vector(3,1){5.25}}
\put(53,10.5){\makebox(0,0)[cc]{40}}
\put(58.75,5.75){\line(0,-1){3}}
\put(59,5.75){\line(1,0){3}}
\put(64.5,24.25){\vector(0,1){4.5}}
\put(64.5,19.75){\vector(0,-1){4}}
\end{picture}

%TeXCAD Picture [prob7.15_42.pic]. Options:
%\grade{\on}
%\emlines{\off}
%\epic{\off}
%\beziermacro{\on}
%\reduce{\on}
%\snapping{\off}
%\quality{8.00}
%\graddiff{0.01}
%\snapasp{1}
%\zoom{4.0000}
\unitlength 1mm % = 2.85pt
\linethickness{0.4pt}
\ifx\plotpoint\undefined\newsavebox{\plotpoint}\fi % GNUPLOT compatibility
\begin{picture}(71,35.57)(0,0)
%\emline(13.5,12.82)(50,34.82)
\multiput(13.5,12.82)(.055895865,.033690658){653}{\line(1,0){.055895865}}
%\end
\put(50,34.82){\line(0,-1){22.25}}
\put(50,12.57){\line(0,-1){1.5}}
\put(50,11.07){\line(-1,0){40}}
%\emline(10,11.07)(14,12.82)
\multiput(10,11.07)(.0769231,.0336538){52}{\line(1,0){.0769231}}
%\end
\put(47,13.82){\line(0,-1){2.75}}
\put(47,13.57){\line(1,0){3.25}}
%\emline(10.75,11.32)(50.25,20.82)
\multiput(10.75,11.32)(.140070922,.033687943){282}{\line(1,0){.140070922}}
%\end
\put(52.25,26.07){\vector(0,-1){5}}
\put(52.5,30.57){\vector(0,1){4.5}}
\qbezier(24,19.07)(29.38,16.07)(26.25,11.07)
\qbezier(23.5,14.07)(25.75,13.57)(24,11.07)
\put(8.25,35.57){\makebox(0,0)[cc]{42)}}
\put(52.25,27.57){\makebox(0,0)[cc]{h}}
\put(71,22.07){\makebox(0,0)[cc]{h = }}
\put(31.25,8.32){\makebox(0,0)[cc]{250}}
\put(22.75,15.82){\makebox(0,0)[cc]{$60 \degrees$}}
\put(20.5,12.07){\makebox(0,0)[cc]{$30 \degrees$}}
\end{picture}

%TeXCAD Picture [prob7.15_43.pic]. Options:
%\grade{\on}
%\emlines{\off}
%\epic{\off}
%\beziermacro{\on}
%\reduce{\on}
%\snapping{\off}
%\quality{8.00}
%\graddiff{0.01}
%\snapasp{1}
%\zoom{4.0000}
\unitlength 1mm % = 2.85pt
\linethickness{0.4pt}
\ifx\plotpoint\undefined\newsavebox{\plotpoint}\fi % GNUPLOT compatibility
\begin{picture}(98.67,32.26)(0,0)
\put(20.17,9.51){\line(1,0){52.25}}
\put(72.42,9.51){\line(0,1){22}}
%\emline(72.42,31.51)(20.17,9.76)
\multiput(72.42,31.51)(-.081007752,-.03372093){645}{\line(-1,0){.081007752}}
%\end
\put(72.67,31.51){\line(1,0){9.25}}
\put(72.42,31.51){\line(-1,0){48.75}}
\put(75.17,31.26){\line(0,-1){2.75}}
\put(75.17,28.51){\line(-1,0){2.75}}
\put(68.92,12.51){\line(0,-1){2.75}}
\put(69.17,12.51){\line(1,0){3.25}}
\put(60.42,29.26){\makebox(0,0)[cc]{$20 \degrees$}}
\put(75.92,18.51){\makebox(0,0)[cc]{100}}
\put(98.67,21.01){\makebox(0,0)[cc]{x = }}
\put(49.17,7.01){\makebox(0,0)[cc]{x}}
\put(8.92,32.26){\makebox(0,0)[cc]{43)}}
\end{picture}

%TeXCAD Picture [prob7.15_44.pic]. Options:
%\grade{\on}
%\emlines{\off}
%\epic{\off}
%\beziermacro{\on}
%\reduce{\on}
%\snapping{\off}
%\quality{8.00}
%\graddiff{0.01}
%\snapasp{1}
%\zoom{4.0000}
\unitlength 1mm % = 2.85pt
\linethickness{0.4pt}
\ifx\plotpoint\undefined\newsavebox{\plotpoint}\fi % GNUPLOT compatibility
\begin{picture}(97.5,25.8)(0,0)
\put(18.25,5.3){\line(1,0){46}}
\put(64.25,5.3){\line(0,1){20.25}}
%\emline(64.25,25.55)(16.25,5.3)
\multiput(64.25,25.55)(-.079866889,-.033693844){601}{\line(-1,0){.079866889}}
%\end
\put(16.25,5.3){\line(1,0){3}}
\put(60.75,8.55){\line(0,-1){3}}
\put(61,8.3){\line(1,0){2.75}}
\put(9.75,25.8){\makebox(0,0)[cc]{44)}}
\put(41.5,2.8){\makebox(0,0)[cc]{x}}
\put(31,7.55){\makebox(0,0)[cc]{$38 \degrees 30'$}}
\put(97.5,14.55){\makebox(0,0)[cc]{x = }}
\put(67.75,14.55){\makebox(0,0)[cc]{300}}
\end{picture}

%TeXCAD Picture [prob7.15_45.pic]. Options:
%\grade{\on}
%\emlines{\off}
%\epic{\off}
%\beziermacro{\on}
%\reduce{\on}
%\snapping{\off}
%\quality{8.00}
%\graddiff{0.01}
%\snapasp{1}
%\zoom{4.0000}
\unitlength 1mm % = 2.85pt
\linethickness{0.4pt}
\ifx\plotpoint\undefined\newsavebox{\plotpoint}\fi % GNUPLOT compatibility
\begin{picture}(76.75,24.06)(0,0)
\put(16.5,4.81){\line(1,0){39}}
\put(55.5,4.81){\line(0,1){16}}
%\emline(55.5,20.81)(16.75,5.06)
\multiput(55.5,20.81)(-.082976445,-.03372591){467}{\line(-1,0){.082976445}}
%\end
\put(52.25,8.06){\line(0,-1){3}}
\put(52.25,8.06){\line(1,0){2.75}}
\put(9.5,24.06){\makebox(0,0)[cc]{45)}}
\put(59,12.56){\makebox(0,0)[cc]{24}}
\put(76.75,14.81){\makebox(0,0)[cc]{$\alpha = $}}
\put(28.5,7.06){\makebox(0,0)[cc]{$\alpha$}}
\put(34.25,15.56){\makebox(0,0)[cc]{30}}
\end{picture}

%TeXCAD Picture [prob7.15_46.pic]. Options:
%\grade{\on}
%\emlines{\off}
%\epic{\off}
%\beziermacro{\on}
%\reduce{\on}
%\snapping{\off}
%\quality{8.00}
%\graddiff{0.01}
%\snapasp{1}
%\zoom{4.0000}
\unitlength 1mm % = 2.85pt
\linethickness{0.4pt}
\ifx\plotpoint\undefined\newsavebox{\plotpoint}\fi % GNUPLOT compatibility
\begin{picture}(115,30.5)(7,0)
\put(17.5,7.5){\line(1,0){71.75}}
\put(89.25,7.5){\line(0,1){23}}
\put(89.25,30.5){\line(-1,0){29.5}}
\put(59.75,30.5){\line(0,-1){22.75}}
%\emline(59.5,30.5)(14.75,7.5)
\multiput(59.5,30.5)(-.0656158358,-.0337243402){682}{\line(-1,0){.0656158358}}
%\end
\put(14.75,7.5){\line(1,0){2.5}}
\put(56.25,10.75){\line(0,-1){3}}
\put(56.5,10.75){\line(1,0){3.5}}
\put(41,5.25){\vector(1,0){18.5}}
\put(29,5.25){\vector(-1,0){14.25}}
\put(9.75,29){\makebox(0,0)[cc]{46)}}
\put(25,9.5){\makebox(0,0)[cc]{$35 \degrees$}}
\put(34.5,4.75){\makebox(0,0)[cc]{300}}
\put(56.75,20.25){\makebox(0,0)[cc]{$h$}}
\put(72.25,20){\makebox(0,0)[cc]{building}}
\put(107,29.5){\makebox(0,0)[cc]{How tall is}}
\put(106.5,22.75){\makebox(0,0)[cc]{the building?}}
\put(106.75,15.75){\makebox(0,0)[cc]{$h = $}}
\end{picture}

%TeXCAD Picture [prob7.15_47.pic]. Options:
%\grade{\on}
%\emlines{\off}
%\epic{\off}
%\beziermacro{\on}
%\reduce{\on}
%\snapping{\off}
%\quality{8.00}
%\graddiff{0.01}
%\snapasp{1}
%\zoom{4.0000}
\unitlength 1mm % = 2.85pt
\linethickness{0.4pt}
\ifx\plotpoint\undefined\newsavebox{\plotpoint}\fi % GNUPLOT compatibility
\begin{picture}(78,53)(0,0)
%\circle(37,26.25){49}
\put(61.5,26.25){\line(0,1){1.061}}
\put(61.48,27.31){\line(0,1){1.059}}
\put(61.41,28.37){\line(0,1){1.055}}
\multiput(61.29,29.43)(-.03208,.20987){5}{\line(0,1){.20987}}
\multiput(61.13,30.48)(-.02938,.14878){7}{\line(0,1){.14878}}
\multiput(60.93,31.52)(-.03133,.12894){8}{\line(0,1){.12894}}
\multiput(60.68,32.55)(-.03278,.1133){9}{\line(0,1){.1133}}
\multiput(60.38,33.57)(-.03081,.09145){11}{\line(0,1){.09145}}
\multiput(60.04,34.57)(-.03185,.08253){12}{\line(0,1){.08253}}
\multiput(59.66,35.56)(-.032674,.074837){13}{\line(0,1){.074837}}
\multiput(59.24,36.54)(-.033322,.068112){14}{\line(0,1){.068112}}
\multiput(58.77,37.49)(-.031711,.058279){16}{\line(0,1){.058279}}
\multiput(58.26,38.42)(-.032194,.053506){17}{\line(0,1){.053506}}
\multiput(57.71,39.33)(-.032566,.049169){18}{\line(0,1){.049169}}
\multiput(57.13,40.22)(-.03284,.045201){19}{\line(0,1){.045201}}
\multiput(56.5,41.08)(-.033029,.041549){20}{\line(0,1){.041549}}
\multiput(55.84,41.91)(-.033141,.038171){21}{\line(0,1){.038171}}
\multiput(55.15,42.71)(-.033183,.035031){22}{\line(0,1){.035031}}
\multiput(54.42,43.48)(-.03467,.033561){22}{\line(-1,0){.03467}}
\multiput(53.65,44.22)(-.03781,.033553){21}{\line(-1,0){.03781}}
\multiput(52.86,44.92)(-.041189,.033477){20}{\line(-1,0){.041189}}
\multiput(52.04,45.59)(-.044843,.033328){19}{\line(-1,0){.044843}}
\multiput(51.19,46.23)(-.048813,.033096){18}{\line(-1,0){.048813}}
\multiput(50.31,46.82)(-.053154,.032771){17}{\line(-1,0){.053154}}
\multiput(49.4,47.38)(-.057932,.03234){16}{\line(-1,0){.057932}}
\multiput(48.48,47.9)(-.06323,.031787){15}{\line(-1,0){.06323}}
\multiput(47.53,48.37)(-.074478,.033483){13}{\line(-1,0){.074478}}
\multiput(46.56,48.81)(-.08218,.03274){12}{\line(-1,0){.08218}}
\multiput(45.57,49.2)(-.09111,.0318){11}{\line(-1,0){.09111}}
\multiput(44.57,49.55)(-.10165,.03061){10}{\line(-1,0){.10165}}
\multiput(43.55,49.86)(-.1286,.03272){8}{\line(-1,0){.1286}}
\multiput(42.53,50.12)(-.14845,.03099){7}{\line(-1,0){.14845}}
\multiput(41.49,50.34)(-.1746,.02862){6}{\line(-1,0){.1746}}
\multiput(40.44,50.51)(-.2635,.0316){4}{\line(-1,0){.2635}}
\put(39.38,50.63){\line(-1,0){1.059}}
\put(38.33,50.71){\line(-1,0){1.061}}
\put(37.27,50.75){\line(-1,0){1.061}}
\put(36.2,50.74){\line(-1,0){1.06}}
\put(35.14,50.68){\line(-1,0){1.057}}
\multiput(34.09,50.58)(-.21021,-.0298){5}{\line(-1,0){.21021}}
\multiput(33.04,50.43)(-.17393,-.0324){6}{\line(-1,0){.17393}}
\multiput(31.99,50.23)(-.12928,-.02993){8}{\line(-1,0){.12928}}
\multiput(30.96,49.99)(-.11365,-.03156){9}{\line(-1,0){.11365}}
\multiput(29.94,49.71)(-.10096,-.0328){10}{\line(-1,0){.10096}}
\multiput(28.93,49.38)(-.08287,-.03096){12}{\line(-1,0){.08287}}
\multiput(27.93,49.01)(-.075186,-.031861){13}{\line(-1,0){.075186}}
\multiput(26.95,48.6)(-.068469,-.032582){14}{\line(-1,0){.068469}}
\multiput(26,48.14)(-.062527,-.03315){15}{\line(-1,0){.062527}}
\multiput(25.06,47.64)(-.057217,-.033588){16}{\line(-1,0){.057217}}
\multiput(24.14,47.1)(-.049519,-.032031){18}{\line(-1,0){.049519}}
\multiput(23.25,46.53)(-.045554,-.032349){19}{\line(-1,0){.045554}}
\multiput(22.39,45.91)(-.041905,-.032577){20}{\line(-1,0){.041905}}
\multiput(21.55,45.26)(-.038528,-.032726){21}{\line(-1,0){.038528}}
\multiput(20.74,44.57)(-.035389,-.032802){22}{\line(-1,0){.035389}}
\multiput(19.96,43.85)(-.032459,-.032813){23}{\line(0,-1){.032813}}
\multiput(19.21,43.1)(-.032417,-.035742){22}{\line(0,-1){.035742}}
\multiput(18.5,42.31)(-.032306,-.03888){21}{\line(0,-1){.03888}}
\multiput(17.82,41.5)(-.032121,-.042255){20}{\line(0,-1){.042255}}
\multiput(17.18,40.65)(-.033623,-.048452){18}{\line(0,-1){.048452}}
\multiput(16.57,39.78)(-.033345,-.052796){17}{\line(0,-1){.052796}}
\multiput(16.01,38.88)(-.032966,-.057578){16}{\line(0,-1){.057578}}
\multiput(15.48,37.96)(-.03247,-.062882){15}{\line(0,-1){.062882}}
\multiput(14.99,37.02)(-.031838,-.068817){14}{\line(0,-1){.068817}}
\multiput(14.55,36.05)(-.03363,-.08182){12}{\line(0,-1){.08182}}
\multiput(14.14,35.07)(-.03279,-.09076){11}{\line(0,-1){.09076}}
\multiput(13.78,34.07)(-.03171,-.10131){10}{\line(0,-1){.10131}}
\multiput(13.47,33.06)(-.03032,-.11399){9}{\line(0,-1){.11399}}
\multiput(13.19,32.03)(-.0326,-.14811){7}{\line(0,-1){.14811}}
\multiput(12.96,31)(-.03051,-.17428){6}{\line(0,-1){.17428}}
\multiput(12.78,29.95)(-.02752,-.21052){5}{\line(0,-1){.21052}}
\put(12.64,28.9){\line(0,-1){1.058}}
\put(12.55,27.84){\line(0,-1){4.24}}
\multiput(12.64,23.6)(.02752,-.21052){5}{\line(0,-1){.21052}}
\multiput(12.78,22.55)(.03051,-.17428){6}{\line(0,-1){.17428}}
\multiput(12.96,21.5)(.0326,-.14811){7}{\line(0,-1){.14811}}
\multiput(13.19,20.47)(.03032,-.11399){9}{\line(0,-1){.11399}}
\multiput(13.47,19.44)(.03171,-.10131){10}{\line(0,-1){.10131}}
\multiput(13.78,18.43)(.03279,-.09076){11}{\line(0,-1){.09076}}
\multiput(14.14,17.43)(.03363,-.08182){12}{\line(0,-1){.08182}}
\multiput(14.55,16.45)(.031838,-.068817){14}{\line(0,-1){.068817}}
\multiput(14.99,15.48)(.03247,-.062882){15}{\line(0,-1){.062882}}
\multiput(15.48,14.54)(.032966,-.057578){16}{\line(0,-1){.057578}}
\multiput(16.01,13.62)(.033345,-.052796){17}{\line(0,-1){.052796}}
\multiput(16.57,12.72)(.033623,-.048452){18}{\line(0,-1){.048452}}
\multiput(17.18,11.85)(.032121,-.042255){20}{\line(0,-1){.042255}}
\multiput(17.82,11)(.032306,-.03888){21}{\line(0,-1){.03888}}
\multiput(18.5,10.19)(.032417,-.035742){22}{\line(0,-1){.035742}}
\multiput(19.21,9.4)(.032459,-.032813){23}{\line(0,-1){.032813}}
\multiput(19.96,8.65)(.035389,-.032802){22}{\line(1,0){.035389}}
\multiput(20.74,7.93)(.038528,-.032726){21}{\line(1,0){.038528}}
\multiput(21.55,7.24)(.041905,-.032577){20}{\line(1,0){.041905}}
\multiput(22.39,6.59)(.045554,-.032349){19}{\line(1,0){.045554}}
\multiput(23.25,5.97)(.049519,-.032031){18}{\line(1,0){.049519}}
\multiput(24.14,5.4)(.057217,-.033588){16}{\line(1,0){.057217}}
\multiput(25.06,4.86)(.062527,-.03315){15}{\line(1,0){.062527}}
\multiput(26,4.36)(.068469,-.032582){14}{\line(1,0){.068469}}
\multiput(26.95,3.9)(.075186,-.031861){13}{\line(1,0){.075186}}
\multiput(27.93,3.49)(.08287,-.03096){12}{\line(1,0){.08287}}
\multiput(28.93,3.12)(.10096,-.0328){10}{\line(1,0){.10096}}
\multiput(29.94,2.79)(.11365,-.03156){9}{\line(1,0){.11365}}
\multiput(30.96,2.51)(.12928,-.02993){8}{\line(1,0){.12928}}
\multiput(31.99,2.27)(.17393,-.0324){6}{\line(1,0){.17393}}
\multiput(33.04,2.07)(.21021,-.0298){5}{\line(1,0){.21021}}
\put(34.09,1.92){\line(1,0){1.057}}
\put(35.14,1.82){\line(1,0){1.06}}
\put(36.2,1.76){\line(1,0){1.061}}
\put(37.27,1.75){\line(1,0){1.061}}
\put(38.33,1.79){\line(1,0){1.059}}
\multiput(39.38,1.87)(.2635,.0316){4}{\line(1,0){.2635}}
\multiput(40.44,1.99)(.1746,.02862){6}{\line(1,0){.1746}}
\multiput(41.49,2.16)(.14845,.03099){7}{\line(1,0){.14845}}
\multiput(42.53,2.38)(.1286,.03272){8}{\line(1,0){.1286}}
\multiput(43.55,2.64)(.10165,.03061){10}{\line(1,0){.10165}}
\multiput(44.57,2.95)(.09111,.0318){11}{\line(1,0){.09111}}
\multiput(45.57,3.3)(.08218,.03274){12}{\line(1,0){.08218}}
\multiput(46.56,3.69)(.074478,.033483){13}{\line(1,0){.074478}}
\multiput(47.53,4.13)(.06323,.031787){15}{\line(1,0){.06323}}
\multiput(48.48,4.6)(.057932,.03234){16}{\line(1,0){.057932}}
\multiput(49.4,5.12)(.053154,.032771){17}{\line(1,0){.053154}}
\multiput(50.31,5.68)(.048813,.033096){18}{\line(1,0){.048813}}
\multiput(51.19,6.27)(.044843,.033328){19}{\line(1,0){.044843}}
\multiput(52.04,6.91)(.041189,.033477){20}{\line(1,0){.041189}}
\multiput(52.86,7.58)(.03781,.033553){21}{\line(1,0){.03781}}
\multiput(53.65,8.28)(.03467,.033561){22}{\line(1,0){.03467}}
\multiput(54.42,9.02)(.033183,.035031){22}{\line(0,1){.035031}}
\multiput(55.15,9.79)(.033141,.038171){21}{\line(0,1){.038171}}
\multiput(55.84,10.59)(.033029,.041549){20}{\line(0,1){.041549}}
\multiput(56.5,11.42)(.03284,.045201){19}{\line(0,1){.045201}}
\multiput(57.13,12.28)(.032566,.049169){18}{\line(0,1){.049169}}
\multiput(57.71,13.17)(.032194,.053506){17}{\line(0,1){.053506}}
\multiput(58.26,14.08)(.031711,.058279){16}{\line(0,1){.058279}}
\multiput(58.77,15.01)(.033322,.068112){14}{\line(0,1){.068112}}
\multiput(59.24,15.96)(.032674,.074837){13}{\line(0,1){.074837}}
\multiput(59.66,16.94)(.03185,.08253){12}{\line(0,1){.08253}}
\multiput(60.04,17.93)(.03081,.09145){11}{\line(0,1){.09145}}
\multiput(60.38,18.93)(.03278,.1133){9}{\line(0,1){.1133}}
\multiput(60.68,19.95)(.03133,.12894){8}{\line(0,1){.12894}}
\multiput(60.93,20.98)(.02938,.14878){7}{\line(0,1){.14878}}
\multiput(61.13,22.02)(.03208,.20987){5}{\line(0,1){.20987}}
\put(61.29,23.07){\line(0,1){1.055}}
\put(61.41,24.13){\line(0,1){1.059}}
\put(61.48,25.19){\line(0,1){1.061}}
%\end
%\emline(18.25,42.25)(49.5,47.25)
\multiput(18.25,42.25)(.20973154,.03355705){149}{\line(1,0){.20973154}}
%\end
%\emline(49.5,47.25)(61.25,26)
\multiput(49.5,47.25)(.033667622,-.060888252){349}{\line(0,-1){.060888252}}
%\end
%\emline(61.25,26)(48.25,5)
\multiput(61.25,26)(-.033678756,-.054404145){386}{\line(0,-1){.054404145}}
%\end
%\emline(48.25,5)(17.5,11.25)
\multiput(48.25,5)(-.16532258,.03360215){186}{\line(-1,0){.16532258}}
%\end
%\emline(17.5,11.25)(18.25,42.25)
\multiput(17.5,11.25)(.032609,1.347826){23}{\line(0,1){1.347826}}
%\end
\put(61,26){\line(-1,0){24.25}}
\put(37.25,26){\circle{.71}}
\put(47.75,29){\makebox(0,0)[cc]{3}}
\put(78,29){\makebox(0,0)[cc]{$x = $}}
\put(8.75,53){\makebox(0,0)[cc]{47)}}
\put(21,26){\makebox(0,0)[cc]{$x$}}
\put(35.25,11.25){\makebox(0,0)[cc]{$x$}}
\put(51.75,17.25){\makebox(0,0)[cc]{$x$}}
\put(34.25,42){\makebox(0,0)[cc]{$x$}}
\put(52.25,37.5){\makebox(0,0)[cc]{$x$}}
\end{picture}

%TeXCAD Picture [prob7.15_48.pic]. Options:
%\grade{\on}
%\emlines{\off}
%\epic{\off}
%\beziermacro{\on}
%\reduce{\on}
%\snapping{\off}
%\quality{8.00}
%\graddiff{0.01}
%\snapasp{1}
%\zoom{4.0000}
\unitlength 1mm % = 2.85pt
\linethickness{0.4pt}
\ifx\plotpoint\undefined\newsavebox{\plotpoint}\fi % GNUPLOT compatibility
\begin{picture}(115,39.31)(10,0)
\put(16.5,6.81){\line(1,0){60}}
\put(76.5,6.81){\line(0,1){32.25}}
%\emline(76.5,39.06)(16.25,7.06)
\multiput(76.5,39.06)(-.063487882,-.033719705){949}{\line(-1,0){.063487882}}
%\end
%\emline(76.25,38.81)(40.25,7.06)
\multiput(76.25,38.81)(-.0382165605,-.0337048832){942}{\line(-1,0){.0382165605}}
%\end
%\dashline{1}(76.75,39.06)(98.5,39.31)
\put(76.68,38.99){\line(1,0){.989}}
\put(78.66,39.01){\line(1,0){.989}}
\put(80.63,39.03){\line(1,0){.989}}
\put(82.61,39.06){\line(1,0){.989}}
\put(84.59,39.08){\line(1,0){.989}}
\put(86.57,39.1){\line(1,0){.989}}
\put(88.54,39.12){\line(1,0){.989}}
\put(90.52,39.15){\line(1,0){.989}}
\put(92.5,39.17){\line(1,0){.989}}
\put(94.48,39.19){\line(1,0){.989}}
\put(96.45,39.21){\line(1,0){.989}}
%\end
%\dashline{1}(98.5,39.31)(98.5,7.31)
\put(98.43,39.24){\line(0,-1){.97}}
\put(98.43,37.3){\line(0,-1){.97}}
\put(98.43,35.36){\line(0,-1){.97}}
\put(98.43,33.42){\line(0,-1){.97}}
\put(98.43,31.48){\line(0,-1){.97}}
\put(98.43,29.54){\line(0,-1){.97}}
\put(98.43,27.6){\line(0,-1){.97}}
\put(98.43,25.66){\line(0,-1){.97}}
\put(98.43,23.72){\line(0,-1){.97}}
\put(98.43,21.78){\line(0,-1){.97}}
\put(98.43,19.84){\line(0,-1){.97}}
\put(98.43,17.9){\line(0,-1){.97}}
\put(98.43,15.96){\line(0,-1){.97}}
\put(98.43,14.03){\line(0,-1){.97}}
\put(98.43,12.09){\line(0,-1){.97}}
\put(98.43,10.15){\line(0,-1){.97}}
\put(98.43,8.21){\line(0,-1){.97}}
%\end
%\dashline{1}(98.5,7.31)(76.75,6.81)
\put(98.43,7.24){\line(-1,0){.989}}
\put(96.45,7.19){\line(-1,0){.989}}
\put(94.48,7.15){\line(-1,0){.989}}
\put(92.5,7.1){\line(-1,0){.989}}
\put(90.52,7.06){\line(-1,0){.989}}
\put(88.54,7.01){\line(-1,0){.989}}
\put(86.57,6.96){\line(-1,0){.989}}
\put(84.59,6.92){\line(-1,0){.989}}
\put(82.61,6.87){\line(-1,0){.989}}
\put(80.63,6.83){\line(-1,0){.989}}
\put(78.66,6.78){\line(-1,0){.989}}
%\end
\put(74,20.81){\makebox(0,0)[cc]{$h$}}
\put(29.25,3.56){\makebox(0,0)[cc]{200}}
\put(87.75,20.81){\makebox(0,0)[cc]{building}}
\put(114.61,20.13){\makebox(0,0)[cc]{$h = $}}
\put(114.61,31.69){\makebox(0,0)[cc]{the building?}}
\put(114.82,36.11){\makebox(0,0)[cc]{How tall is}}
\put(13.71,38.21){\makebox(0,0)[cc]{48)}}
\put(47,9.06){\makebox(0,0)[cc]{$45 \degrees$}}
\put(25.5,9.31){\makebox(0,0)[cc]{$30 \degrees$}}
\end{picture}

%TeXCAD Picture [prob7.15_49.pic]. Options:
%\grade{\on}
%\emlines{\off}
%\epic{\off}
%\beziermacro{\on}
%\reduce{\on}
%\snapping{\off}
%\quality{8.00}
%\graddiff{0.01}
%\snapasp{1}
%\zoom{4.0000}
\unitlength 1mm % = 2.85pt
\linethickness{0.4pt}
\ifx\plotpoint\undefined\newsavebox{\plotpoint}\fi % GNUPLOT compatibility
\begin{picture}(91,41.5)(0,0)
\put(17.75,14.5){\line(1,0){52.25}}
\put(43.25,14.5){\line(0,1){26.75}}
%\emline(43.25,41.25)(60.75,29)
\multiput(43.25,41.25)(.048076923,-.033653846){364}{\line(1,0){.048076923}}
%\end
\put(60.75,29){\line(0,-1){14.25}}
\put(57.25,18.25){\line(0,-1){3.5}}
\put(57.25,18.25){\line(1,0){3.25}}
\put(39.75,18){\line(0,-1){3.25}}
\put(40,18){\line(1,0){3.25}}
\qbezier(19.25,24.75)(46.25,-1.38)(67.25,41)
\put(90.75,30){\makebox(0,0)[cc]{Show that}}
\put(91,24.25){\makebox(0,0)[cc]{$h = L (1 - \cos \alpha)$}}
\put(63,21){\makebox(0,0)[cc]{$h$}}
\put(11.25,41.5){\makebox(0,0)[cc]{49)}}
\put(54.75,37){\makebox(0,0)[cc]{$L$}}
\put(45,37.75){\makebox(0,0)[cc]{$\alpha$}}
\qbezier(43.5,35.5)(46.88,35)(48.75,37.5)
\end{picture}

% Problem 7.15, Question 50
%TeXCAD Options
%\grade{\on}
%\emlines{\off}
%\epic{\off}
%\beziermacro{\on}
%\reduce{\on}
%\snapping{\off}
%\quality{8.00}
%\graddiff{0.01}
%\snapasp{1}
%\zoom{4.0000}
\unitlength 1mm % = 2.85pt
\linethickness{0.4pt}
\ifx\plotpoint\undefined\newsavebox{\plotpoint}\fi % GNUPLOT compatibility
\begin{picture}(101,34.5)(0,111)
\put(11.75,115.5){\line(1,0){86.25}}
\put(19.25,115.5){\line(2,1){42}} 
\put(61.25,136.5){\line(1,-1){21}}
\put(61.25,136.25){\line(0,-1){20.5}}
\put(54.5,113){\vector(1,0){27.25}}
\put(46,112.75){\vector(-1,0){26.75}}
\put(61,140.25){\makebox(0,0)[cc]{$R$}}
\put(73.75,128.25){\makebox(0,0)[cc]{$p$}}
\put(84.75,112.5){\makebox(0,0)[cc]{$Q$}}
\put(16.5,113){\makebox(0,0)[cc]{$P$}}
\put(50.25,112.5){\makebox(0,0)[cc]{$r$}}
\put(31.75,126){\makebox(0,0)[cc]{$q$}}
\put(58.5,125.75){\makebox(0,0)[cc]{$h$}}
\put(28.75,117.5){\makebox(0,0)[cc]{$\phi$}}
\put(75.5,118){\makebox(0,0)[cc]{$\theta$}}
\put(101,135){\makebox(0,0)[cc]{Show that}}
\put(101,130){\makebox(0,0)[cc]{$h = \frac{r}{(\cot \phi + \cot \theta)}$}}
\put(10.75,145.5){\makebox(0,0)[cc]{50)}}
\end{picture}


Some of these have been interesting; i.e., from mildly challenging to maddeningly difficult. I hope you found some ways of working some of the problems that were different from those presented in the hints section. It's much more satisfying to find your own way to do things. 

In any event, you probably see how useful these new functions (the trigonometric functions) can be. Physicists use them a lot to describe real-world goings-on. You should definitely have a real working familiarity with these functions before you take any physics courses. And please, regardless of what the college catalog says about prerequisites for college physics, take at least one year of calculus before you tackle a college-level physics course. It's not that the material is so difficult. It is simply that without some of the tools of calculus, you will spend too much time trying to follow the mathematics and you will miss the beauty and the point of the physics. It would be a real shame if were to be turned off of physics simply because you took it prematurely.

\vspace{0.3cm}

Here is another irrelevant, hair-pulling, brain-cramping puzzle for you.

Suppose you have 12 identical-looking coins, but you know that exactly one of them is defective. The defective one is either lighter or heavier than it should be. The only tool you have to discover the defective coin is a balance scale. It won't tell you how much something weighs, but it will tell you whether one thing weighs more than, less than, or exactly the same as, another thing. 

Now, here is the problem. With no more than three weighings, determine which coin is defective, and whether it is lighter or heavier than it should be.


\newpage
\begin{center}
\section*{Chapter VII Hints and Solutions}
\end{center}

$\ $

\textbf{Problem 7.1}: The statement that ``$T$ is the triangle $ABC$'' means:
\begin{enumerate}[\hspace{1cm}1)]
  \item $T$ is a pointset,
  \item $A$, $B$, and $C$ are non-colinear points, and
  \item a point $P$ belongs to $T$ if and only if
  \begin{enumerate}[\hspace{1cm}a)]
    \item $P$ is a point of interval $[A,B]$,
    \item $P$ is a point of interval $[B,C]$, or
    \item $P$ is a point of interval $[C,A]$.
  \end{enumerate}
\end{enumerate}
\hspace{1cm} Draw a picture.

\vspace{0.3cm}

\textbf{Problem 7.2}: The statement that A is the angle PQR means
\begin{enumerate}[\hspace{1cm}1)]
  \item $A$ is a pointset,
  \item $P$, $Q$, and $R$ are distinct points, and
  \item a point $x$ belongs to $A$ if and only if
  \begin{enumerate}[\hspace{1cm}a)]
    \item $x$ is a point of ray $QP$, or
    \item $x$ is a point of ray $QR$.
  \end{enumerate}
\end{enumerate}
\hspace{1cm} Draw a picture.

\vspace{0.3cm}

\textbf{Problem 7.3}:
\begin{enumerate}[\hspace{1cm}1)]
  \item Use the notion of side of a line.
  \item The statement that $x$ is an interior point of angle $PQR$ means
  \begin{enumerate}[\hspace{1cm}a)]
    \item $x$ is on the $p$-side of line $QR$, and
    \item $x$ is on the $R$-side of line $QP$.
  \end{enumerate}
\end{enumerate}
\hspace{1cm} Draw a picture.

\vspace{0.3cm}

\textbf{Problem 7.4}: Use the notion of side of a line.

\vspace{0.3cm}

\textbf{Problem 7.5}: No hints.

\vspace{0.3cm}

\textbf{Problem 7.6}: Consider the following picture.

\begin{center}
%TeXCAD Options
%\grade{\on}
%\emlines{\off}
%\epic{\off}
%\beziermacro{\on}
%\reduce{\on}
%\snapping{\off}
%\quality{8.00}
%\graddiff{0.01}
%\snapasp{1}
%\zoom{4.0000}
\unitlength 1mm % = 2.85pt
\linethickness{0.4pt}
\ifx\plotpoint\undefined\newsavebox{\plotpoint}\fi % GNUPLOT compatibility
\begin{picture}(93,31.5)(0,111)
\put(12,115.5){\line(1,0){80.75}}
%\emline(22.25,115.5)(50.5,141.5)
\multiput(22.25,115.5)(.0366407263,.0337224384){771}{\line(1,0){.0366407263}}
%\end
%\emline(22.75,115.75)(93,142.5)
\multiput(22.75,115.75)(.0885876419,.0337326608){793}{\line(1,0){.0885876419}}
%\end
\put(39.75,131.5){\circle{.5}} 
\put(45.75,137.25){\circle{.5}}
\put(45.5,124.5){\circle{.5}} 
\put(76.25,136.25){\circle{.5}}
\put(47.5,115.5){\circle{.5}} 
\put(80,115.5){\circle{.5}}
\put(42.25,130.5){\makebox(0,0)[cc]{$P_{1}$}}
\put(48.5,136.75){\makebox(0,0)[cc]{$P$}}
\put(45.75,121.75){\makebox(0,0)[cc]{$Q_{1}$}}
\put(47.25,112){\makebox(0,0)[cc]{$R_{1}$}}
\put(79.25,112){\makebox(0,0)[cc]{$R$}}
\put(76.75,132.5){\makebox(0,0)[cc]{$Q$}}
\put(22.75,112){\makebox(0,0)[cc]{$O$}}
\end{picture}
\end{center}

Suppose $OP_{1} = OQ_{1} = OR_{1} = 1$.

Can you see that, on the basis of this definition,

$m \angm POQ + m \angm QOR \neq m \angm POR$?

If we want angle measure to have this property, then we must eschew (I always wanted to use that word) this definition.

\vspace{0.3cm}

\textbf{Problem 7.7}:
\begin{enumerate}[\hspace{1cm}a)]
  \item $c=a^{2}+b^{2}$, $\theta = \tan^{-1}(\frac{b}{a})$, $\phi = \tan^{-1}(\frac{a}{b})$
  \item $b=c^{2}-a^{2}$, $\theta = \cos^{-1}(\frac{a}{c})$, $\phi = \sin^{-1}(\frac{a}{c})$
  \item $a = c \cos \theta$, $b = c \sin \theta$, $\phi = \frac{\pi}{2} - \theta$
  \item $a = c \sin \phi$, $b = c \cos \phi$, $\theta = \frac{\pi}{2} - \phi$
  \item $b = a \tan \theta$, $c = \frac{a}{\cos \theta}$, $\phi = \frac{\pi}{2} - \theta$
  \item $a = b \tan \phi$, $c = \frac{b}{\cos \phi}$, $\theta = \frac{\pi}{2} - \phi$
\end{enumerate}

\textbf{Problem 7.8}:
\begin{enumerate}[\hspace{1cm}a)]
  \item $\sin \theta = \frac{y}{r}$
  \item $\cos \theta = \frac{x}{r}$ 
  \item $\tan \theta = \frac{y}{x}$
  \item $\csc \theta = \frac{r}{y}$
  \item $\sec \theta = \frac{r}{x}$
  \item $\cot \theta = \frac{x}{y}$
\end{enumerate}

\textbf{Problem 7.9}: No hints.

\vspace{0.3cm}

\textbf{Problem 7.10}:
\begin{enumerate}[\hspace{1cm}1)]
  \item Draw the following, and find the coordinates of $P$.
\end{enumerate}

\begin{center}
%TeXCAD Options
%\grade{\on}
%\emlines{\off}
%\epic{\off}
%\beziermacro{\on}
%\reduce{\on}
%\snapping{\off}
%\quality{8.00}
%\graddiff{0.01}
%\snapasp{1}
%\zoom{4.0000}
\unitlength 1mm % = 2.85pt
\linethickness{0.4pt}
\ifx\plotpoint\undefined\newsavebox{\plotpoint}\fi % GNUPLOT compatibility
\begin{picture}(88.75,39.25)(0,105)
\put(16.25,109.75){\line(1,0){42}}
%\emline(58.25,109.75)(84.5,143.25)
\multiput(58.25,109.75)(.0336970475,.0430038511){779}{\line(0,1){.0430038511}}
%\end
%\emline(84.5,143.25)(16,110)
\multiput(84.5,143.25)(-.0694726166,-.0337221095){986}{\line(-1,0){.0694726166}}
%\end
\qbezier(26.75,115)(29.5,113.63)(27.25,109.75)
\put(16.25,106.5){\makebox(0,0)[cc]{$O:(0,0)$}}
\put(58,106.5){\makebox(0,0)[cc]{$(a,0)$}}
\put(25,111.75){\makebox(0,0)[cc]{$\theta$}}
\put(88.75,144.25){\makebox(0,0)[cc]{$P:(?,?)$}}
\put(43,127.5){\makebox(0,0)[cc]{$c$}}
\put(42.75,112.25){\makebox(0,0)[cc]{$a$}}
\put(63.25,120.75){\makebox(0,0)[cc]{$b$}}
\end{picture}
\end{center}

\begin{enumerate}
  \item[\hspace{1cm}2)] $P: (c \cos \theta, c \sin \theta)$, from Corollary 5.3.
  \item[\hspace{1cm}3)] Add the following to your drawing:
\end{enumerate}

\begin{center}
%TeXCAD Options
%\grade{\on}
%\emlines{\off}
%\epic{\off}
%\beziermacro{\on}
%\reduce{\on}
%\snapping{\off}
%\quality{8.00}
%\graddiff{0.01}
%\snapasp{1}
%\zoom{4.0000}
\unitlength 1mm % = 2.85pt
\linethickness{0.4pt}
\ifx\plotpoint\undefined\newsavebox{\plotpoint}\fi % GNUPLOT compatibility
\begin{picture}(88.25,26.25)(0,110)
\put(18.5,115){\line(1,0){40}} \put(58.5,115){\line(1,1){21}}
%\emline(79.5,136)(18,115.25)
\multiput(79.5,136)(-.1,-.033739837){615}{\line(-1,0){.1}}
%\end
%\dashline{1}(79.75,136)(79.75,114.75)
\put(79.68,135.93){\line(0,-1){.966}}
\put(79.68,134){\line(0,-1){.966}}
\put(79.68,132.07){\line(0,-1){.966}}
\put(79.68,130.13){\line(0,-1){.966}}
\put(79.68,128.2){\line(0,-1){.966}}
\put(79.68,126.27){\line(0,-1){.966}}
\put(79.68,124.34){\line(0,-1){.966}}
\put(79.68,122.41){\line(0,-1){.966}}
\put(79.68,120.48){\line(0,-1){.966}}
\put(79.68,118.54){\line(0,-1){.966}}
\put(79.68,116.61){\line(0,-1){.966}}
%\end
%\dashline{1}(79.75,114.75)(58.75,115.25)
\put(79.68,114.68){\line(-1,0){.955}}
\put(77.77,114.73){\line(-1,0){.955}}
\put(75.86,114.77){\line(-1,0){.955}}
\put(73.95,114.82){\line(-1,0){.955}}
\put(72.04,114.86){\line(-1,0){.955}}
\put(70.13,114.91){\line(-1,0){.955}}
\put(68.23,114.95){\line(-1,0){.955}}
\put(66.32,115){\line(-1,0){.955}}
\put(64.41,115.04){\line(-1,0){.955}}
\put(62.5,115.09){\line(-1,0){.955}}
\put(60.59,115.13){\line(-1,0){.955}}
%\end
\put(76.25,117.75){\line(0,-1){3}} 
\put(76.5,118){\line(1,0){3.5}}
\qbezier(31,119.75)(34,117.63)(31,115)
\put(92,136.25){\makebox(0,0)[cc]{$(c \cos \theta, c \sin \theta$)}}
\put(79.5,111.5){\makebox(0,0)[cc]{$(c \cos \theta, \theta)$}}
\put(58.25,111.5){\makebox(0,0)[cc]{$(a,0)$}}
\put(17.5,111.5){\makebox(0,0)[cc]{$O:(0,0)$}}
\put(46.75,117){\makebox(0,0)[cc]{$a$}}
\put(64.25,124){\makebox(0,0)[cc]{$b$}}
\put(46.75,129.25){\makebox(0,0)[cc]{$c$}}
\put(28.75,117){\makebox(0,0)[cc]{$\theta$}}
\end{picture}
\end{center}

\begin{enumerate}
  \item[\hspace{1cm}4)] Consider the right triangle.
\end{enumerate}

\begin{center}
%TeXCAD Options
%\grade{\on}
%\emlines{\off}
%\epic{\off}
%\beziermacro{\on}
%\reduce{\on}
%\snapping{\off}
%\quality{8.00}
%\graddiff{0.01}
%\snapasp{1}
%\zoom{4.0000}
\unitlength 1mm % = 2.85pt
\linethickness{0.4pt}
\ifx\plotpoint\undefined\newsavebox{\plotpoint}\fi % GNUPLOT compatibility
\begin{picture}(60.25,37)(0,145.75)
%\emline(18.75,154)(57.25,182.75)
\multiput(18.75,154)(.0451348183,.0337045721){853}{\line(1,0){.0451348183}}
%\end
\put(57.25,182.75){\line(0,-1){31.75}}
\put(57.25,151){\line(-1,0){42.5}}
%\emline(14.75,151)(19,154.25)
\multiput(14.75,151)(.04381443,.03350515){97}{\line(1,0){.04381443}}
%\end
\put(53,155.25){\line(0,-1){4}} 
\put(53,155.25){\line(1,0){4.25}}
\put(31.75,166.75){\makebox(0,0)[cc]{$b$}}
\put(63,164.75){\makebox(0,0)[cc]{$c \sin \theta$}}
\put(36.5,147){\makebox(0,0)[cc]{$c \cos \theta - a$}}
\end{picture}
\end{center}

\begin{enumerate}
  \item[\hspace{1cm}5)] Use the Pythagorean Theorem on the triangle in Step 4. You should get the result that you want.
\end{enumerate}


\textbf{Problem 7.12}:
\begin{enumerate}[\hspace{1cm}1)]
  \item Note $\sin A = \frac{h}{c}$, and $\sin (\pi - c) = \frac{h}{a}$ (From Fig. 7.12).
  \item Eliminate $h$ from the equations in Step 1.
  \item Expand $\sin (\pi - c)$.
  \item You should have $\frac{\sin A}{a} = \frac{\sin C}{c}$. And the problem now is to get $\frac{\sin B}{b}$ into the picture.
  \item Redraw the picture from a different perspective.
  \item Consider the following:
\end{enumerate}

\begin{center}
%TeXCAD Options
%\grade{\on}
%\emlines{\off}
%\epic{\off}
%\beziermacro{\on}
%\reduce{\on}
%\snapping{\off}
%\quality{8.00}
%\graddiff{0.01}
%\snapasp{1}
%\zoom{4.0000}
\unitlength 1mm % = 2.85pt
\linethickness{0.4pt}
\ifx\plotpoint\undefined\newsavebox{\plotpoint}\fi % GNUPLOT compatibility
\begin{picture}(89.25,26.75)(0,111)
\put(15.25,114.5){\line(1,0){74}}
%\emline(24.25,114.5)(51.25,137.75)
\multiput(24.25,114.5)(.0391304348,.0336956522){690}{\line(1,0){.0391304348}}
%\end
%\emline(51.25,137.75)(80.5,114.5)
\multiput(51.25,137.75)(.0423913043,-.0336956522){690}{\line(1,0){.0423913043}}
%\end
\put(51.25,137.75){\line(0,-1){23}} 
\put(47.5,118){\line(0,-1){3.5}}
\put(47.75,118){\line(1,0){3.5}}
\put(30,116.5){\makebox(0,0)[cc]{$B$}}
\put(73.5,116.5){\makebox(0,0)[cc]{$A$}}
\put(34,125.25){\makebox(0,0)[cc]{$a$}}
\put(66.25,128.75){\makebox(0,0)[cc]{$b$}}
\put(53,124.75){\makebox(0,0)[cc]{$h$}}
\put(51.75,112){\makebox(0,0)[cc]{$c$}}
\end{picture}
\end{center}

\begin{enumerate}
  \item[\hspace{1cm}7)] Note $\sin B = \frac{h}{a}$ and $\sin A = \frac{h}{b}$. This should get you the result you want.
\end{enumerate}


\textbf{Problem 7.13}: Suppose the circumference of a circle with radius $r>0$ is $2 \pi r$.
\begin{enumerate}[\hspace{1cm}1)]
  \item Find each of the following:
  \begin{enumerate}[\hspace{1cm}a)]
    \item The length of the whole circle; i.e., its circumference.
    \item The length of a semi-circle.
    \item The length of a quarter of the circle.
    \item The length of the arc from $R$ to $R$; i.e. $l(\arc{RR})$.
  \end{enumerate}
  \item You should have gotten the following in Step 1:
  \begin{enumerate}[\hspace{1cm}a)]
    \item $2 \pi r$ 
    \item $\pi r$
    \item $(\frac{\pi}{2}) r$
    \item $0$
  \end{enumerate}
  \item Tabulate your results from Step 1, and look for a pattern.
  \item Your table could look like the following:
\end{enumerate}

\begin{center}
\begin{table}[htbp]
\begin{center}
\begin{tabular}{|l|l|}
  \hline 
  Arc            &  Arc Length \\
  \hline 
  whole circle   & $2 \pi r$ \\
  semi-circle    & $\pi r$ \\
  quarter-circle & $(\frac{\pi}{2})r$ \\
  RR             & $0$ \\
\hline
\end{tabular}
\label{tbl7.13hint1}
\end{center}
\end{table}
\end{center}

\begin{enumerate}
  \item[\hspace{1cm}5)] What is the relationship between the left column and the right column?
  \item[\hspace{1cm}6)] Graph the entries in your table, letting whole circle = $1$, semi-circle = $\frac{l}{2}$, etc.
  \item[\hspace{1cm}7)] The relationship is linear, is it not?
  \item[\hspace{1cm}8)] Fill in the following table:
\end{enumerate}

\begin{center}
\begin{table}[htbp]
\begin{center}
\begin{tabular}{|l|l|}
  \hline 
  Arc Length         &  Angle Measure \\
  \hline 
  $\pi r$            &  \\
  $(\frac{\pi}{2})r$ &  \\ 
  $(\frac{\pi}{4})r$ &  \\
  $0$                &  \\
  \hline
\end{tabular}
\label{tbl7.13hint2}
\end{center}
\end{table}
\end{center}

\begin{enumerate}
  \item[\hspace{1cm}9)] If you made another entry in the table above under Angle Measure of $\theta$, what would be the corresponding entry under Arc Length?
  \item[\hspace{1cm}10)] Since this relationship is linear, the entry in Step 9 should be $\theta r$. Looks familiar, doesn't it?
  \item[\hspace{1cm}11)] Using the results in Step 10, you can solve our main problem.
\end{enumerate}


\textbf{Problem 7.14}:
\begin{enumerate}
  \item[\hspace{1cm}1)] Complete the following table:
\end{enumerate}

\begin{center}
\begin{table}[htbp]
\begin{center}
\begin{tabular}{|l|l|}
  \hline 
  Sector         & Area   \\
  \hline 
  whole circle   & \qquad \\
  half circle    &        \\
  quarter circle &        \\
  eighth circle  &        \\
  0              &        \\
  \hline
\end{tabular}
\label{tbl7.14hint}
\end{center}
\end{table}
\end{center}

\begin{enumerate}
  \item[\hspace{1cm}2)] Note that the relationship between the sector as it relates to a whole circle and its area is linear.
  \item[\hspace{1cm}3)] State the linear relationship that you see in Step 2 as a formula.
  \item[\hspace{1cm}4)] The relationship you should have gotten in Step 3 is ``area of a sector is $(\frac{1}{2})r^{2} \theta$''. Here is a picture.
\end{enumerate}

\begin{center}
%TeXCAD Picture [prob7.14hint.pic]. Options:
%\grade{\on}
%\emlines{\off}
%\epic{\off}
%\beziermacro{\on}
%\reduce{\on}
%\snapping{\off}
%\quality{8.00}
%\graddiff{0.01}
%\snapasp{1}
%\zoom{4.0000}
\unitlength 1mm % = 2.85pt
\linethickness{0.4pt}
\ifx\plotpoint\undefined\newsavebox{\plotpoint}\fi % GNUPLOT compatibility
\begin{picture}(55.73,41.73)(0,0)
%\circle(36.1,22.1){39.25}
\put(55.73,22.1){\line(0,1){.911}}
\put(55.71,23.01){\line(0,1){.909}}
\put(55.64,23.92){\line(0,1){.905}}
\multiput(55.54,24.83)(-.02951,.17992){5}{\line(0,1){.17992}}
\multiput(55.39,25.73)(-.03153,.14863){6}{\line(0,1){.14863}}
\multiput(55.2,26.62)(-.03291,.126){7}{\line(0,1){.126}}
\multiput(54.97,27.5)(-.03012,.09671){9}{\line(0,1){.09671}}
\multiput(54.7,28.37)(-.03112,.08568){10}{\line(0,1){.08568}}
\multiput(54.39,29.23)(-.03188,.0765){11}{\line(0,1){.0765}}
\multiput(54.04,30.07)(-.03245,.06869){12}{\line(0,1){.06869}}
\multiput(53.65,30.89)(-.032866,.061946){13}{\line(0,1){.061946}}
\multiput(53.22,31.7)(-.033156,.056042){14}{\line(0,1){.056042}}
\multiput(52.76,32.48)(-.033342,.050812){15}{\line(0,1){.050812}}
\multiput(52.26,33.25)(-.033436,.046134){16}{\line(0,1){.046134}}
\multiput(51.72,33.98)(-.033452,.041912){17}{\line(0,1){.041912}}
\multiput(51.15,34.7)(-.033397,.038073){18}{\line(0,1){.038073}}
\multiput(50.55,35.38)(-.033281,.034561){19}{\line(0,1){.034561}}
\multiput(49.92,36.04)(-.03485,.032978){19}{\line(-1,0){.03485}}
\multiput(49.26,36.67)(-.038363,.033065){18}{\line(-1,0){.038363}}
\multiput(48.57,37.26)(-.042201,.033085){17}{\line(-1,0){.042201}}
\multiput(47.85,37.82)(-.046423,.033033){16}{\line(-1,0){.046423}}
\multiput(47.11,38.35)(-.051101,.032898){15}{\line(-1,0){.051101}}
\multiput(46.34,38.84)(-.056329,.032667){14}{\line(-1,0){.056329}}
\multiput(45.55,39.3)(-.06223,.032325){13}{\line(-1,0){.06223}}
\multiput(44.74,39.72)(-.06897,.03185){12}{\line(-1,0){.06897}}
\multiput(43.92,40.1)(-.07677,.03121){11}{\line(-1,0){.07677}}
\multiput(43.07,40.45)(-.08595,.03038){10}{\line(-1,0){.08595}}
\multiput(42.21,40.75)(-.10909,.03294){8}{\line(-1,0){.10909}}
\multiput(41.34,41.02)(-.12629,.03182){7}{\line(-1,0){.12629}}
\multiput(40.45,41.24)(-.1489,.03024){6}{\line(-1,0){.1489}}
\multiput(39.56,41.42)(-.18017,.02795){5}{\line(-1,0){.18017}}
\put(38.66,41.56){\line(-1,0){.906}}
\put(37.75,41.66){\line(-1,0){.91}}
\put(36.84,41.71){\line(-1,0){.912}}
\put(35.93,41.73){\line(-1,0){.911}}
\put(35.02,41.7){\line(-1,0){.909}}
\put(34.11,41.62){\line(-1,0){.905}}
\multiput(33.21,41.51)(-.17966,-.03108){5}{\line(-1,0){.17966}}
\multiput(32.31,41.36)(-.14835,-.03283){6}{\line(-1,0){.14835}}
\multiput(31.42,41.16)(-.11,-.02976){8}{\line(-1,0){.11}}
\multiput(30.54,40.92)(-.09644,-.03096){9}{\line(-1,0){.09644}}
\multiput(29.67,40.64)(-.08541,-.03187){10}{\line(-1,0){.08541}}
\multiput(28.82,40.32)(-.07622,-.03255){11}{\line(-1,0){.07622}}
\multiput(27.98,39.97)(-.0684,-.03305){12}{\line(-1,0){.0684}}
\multiput(27.16,39.57)(-.061657,-.033404){13}{\line(-1,0){.061657}}
\multiput(26.36,39.13)(-.055751,-.033643){14}{\line(-1,0){.055751}}
\multiput(25.58,38.66)(-.047362,-.031671){16}{\line(-1,0){.047362}}
\multiput(24.82,38.16)(-.043144,-.031846){17}{\line(-1,0){.043144}}
\multiput(24.09,37.62)(-.039306,-.031937){18}{\line(-1,0){.039306}}
\multiput(23.38,37.04)(-.037781,-.033728){18}{\line(-1,0){.037781}}
\multiput(22.7,36.43)(-.03427,-.03358){19}{\line(-1,0){.03427}}
\multiput(22.05,35.8)(-.032674,-.035136){19}{\line(0,-1){.035136}}
\multiput(21.43,35.13)(-.032729,-.038649){18}{\line(0,-1){.038649}}
\multiput(20.84,34.43)(-.032717,-.042488){17}{\line(0,-1){.042488}}
\multiput(20.28,33.71)(-.032627,-.046709){16}{\line(0,-1){.046709}}
\multiput(19.76,32.96)(-.032451,-.051385){15}{\line(0,-1){.051385}}
\multiput(19.27,32.19)(-.032175,-.056611){14}{\line(0,-1){.056611}}
\multiput(18.82,31.4)(-.031781,-.062509){13}{\line(0,-1){.062509}}
\multiput(18.41,30.59)(-.03125,-.06924){12}{\line(0,-1){.06924}}
\multiput(18.03,29.76)(-.0336,-.08474){10}{\line(0,-1){.08474}}
\multiput(17.7,28.91)(-.03292,-.09579){9}{\line(0,-1){.09579}}
\multiput(17.4,28.05)(-.03199,-.10937){8}{\line(0,-1){.10937}}
\multiput(17.14,27.17)(-.03071,-.12656){7}{\line(0,-1){.12656}}
\multiput(16.93,26.29)(-.02894,-.14916){6}{\line(0,-1){.14916}}
\multiput(16.76,25.39)(-.033,-.2255){4}{\line(0,-1){.2255}}
\put(16.62,24.49){\line(0,-1){.907}}
\put(16.53,23.58){\line(0,-1){1.822}}
\put(16.48,21.76){\line(0,-1){.911}}
\put(16.52,20.85){\line(0,-1){.908}}
\multiput(16.6,19.94)(.0303,-.2259){4}{\line(0,-1){.2259}}
\multiput(16.72,19.04)(.03264,-.17938){5}{\line(0,-1){.17938}}
\multiput(16.88,18.14)(.02924,-.12691){7}{\line(0,-1){.12691}}
\multiput(17.09,17.25)(.03072,-.10973){8}{\line(0,-1){.10973}}
\multiput(17.33,16.37)(.0318,-.09617){9}{\line(0,-1){.09617}}
\multiput(17.62,15.51)(.03261,-.08513){10}{\line(0,-1){.08513}}
\multiput(17.95,14.66)(.03321,-.07593){11}{\line(0,-1){.07593}}
\multiput(18.31,13.82)(.03364,-.06811){12}{\line(0,-1){.06811}}
\multiput(18.71,13)(.031515,-.056981){14}{\line(0,-1){.056981}}
\multiput(19.16,12.21)(.031852,-.051759){15}{\line(0,-1){.051759}}
\multiput(19.63,11.43)(.032083,-.047085){16}{\line(0,-1){.047085}}
\multiput(20.15,10.68)(.032221,-.042865){17}{\line(0,-1){.042865}}
\multiput(20.69,9.95)(.032278,-.039027){18}{\line(0,-1){.039027}}
\multiput(21.28,9.25)(.032263,-.035513){19}{\line(0,-1){.035513}}
\multiput(21.89,8.57)(.032184,-.032277){20}{\line(0,-1){.032277}}
\multiput(22.53,7.93)(.035419,-.032366){19}{\line(1,0){.035419}}
\multiput(23.2,7.31)(.038933,-.032391){18}{\line(1,0){.038933}}
\multiput(23.91,6.73)(.042771,-.032345){17}{\line(1,0){.042771}}
\multiput(24.63,6.18)(.046991,-.032219){16}{\line(1,0){.046991}}
\multiput(25.38,5.66)(.051666,-.032003){15}{\line(1,0){.051666}}
\multiput(26.16,5.18)(.056889,-.031681){14}{\line(1,0){.056889}}
\multiput(26.96,4.74)(.062784,-.031236){13}{\line(1,0){.062784}}
\multiput(27.77,4.33)(.07583,-.03343){11}{\line(1,0){.07583}}
\multiput(28.61,3.96)(.08503,-.03286){10}{\line(1,0){.08503}}
\multiput(29.46,3.64)(.09608,-.03208){9}{\line(1,0){.09608}}
\multiput(30.32,3.35)(.10964,-.03103){8}{\line(1,0){.10964}}
\multiput(31.2,3.1)(.12682,-.02961){7}{\line(1,0){.12682}}
\multiput(32.09,2.89)(.17928,-.03317){5}{\line(1,0){.17928}}
\multiput(32.98,2.73)(.2258,-.031){4}{\line(1,0){.2258}}
\put(33.89,2.6){\line(1,0){.908}}
\put(34.79,2.52){\line(1,0){.911}}
\put(35.7,2.48){\line(1,0){.912}}
\put(36.62,2.48){\line(1,0){.911}}
\put(37.53,2.53){\line(1,0){.907}}
\multiput(38.43,2.62)(.2256,.0323){4}{\line(1,0){.2256}}
\multiput(39.34,2.74)(.14924,.0285){6}{\line(1,0){.14924}}
\multiput(40.23,2.92)(.12665,.03035){7}{\line(1,0){.12665}}
\multiput(41.12,3.13)(.10946,.03167){8}{\line(1,0){.10946}}
\multiput(41.99,3.38)(.09589,.03264){9}{\line(1,0){.09589}}
\multiput(42.86,3.67)(.08484,.03335){10}{\line(1,0){.08484}}
\multiput(43.71,4.01)(.06933,.03105){12}{\line(1,0){.06933}}
\multiput(44.54,4.38)(.062601,.0316){13}{\line(1,0){.062601}}
\multiput(45.35,4.79)(.056704,.032011){14}{\line(1,0){.056704}}
\multiput(46.15,5.24)(.051479,.032302){15}{\line(1,0){.051479}}
\multiput(46.92,5.72)(.046804,.032492){16}{\line(1,0){.046804}}
\multiput(47.67,6.24)(.042583,.032593){17}{\line(1,0){.042583}}
\multiput(48.39,6.8)(.038744,.032617){18}{\line(1,0){.038744}}
\multiput(49.09,7.39)(.03523,.032571){19}{\line(1,0){.03523}}
\multiput(49.76,8)(.03368,.034172){19}{\line(0,1){.034172}}
\multiput(50.4,8.65)(.032056,.0357){19}{\line(0,1){.0357}}
\multiput(51.01,9.33)(.032051,.039214){18}{\line(0,1){.039214}}
\multiput(51.58,10.04)(.031972,.043051){17}{\line(0,1){.043051}}
\multiput(52.13,10.77)(.031809,.04727){16}{\line(0,1){.04727}}
\multiput(52.64,11.53)(.031551,.051943){15}{\line(0,1){.051943}}
\multiput(53.11,12.31)(.033583,.06156){13}{\line(0,1){.06156}}
\multiput(53.55,13.11)(.03324,.06831){12}{\line(0,1){.06831}}
\multiput(53.94,13.93)(.03277,.07612){11}{\line(0,1){.07612}}
\multiput(54.3,14.76)(.03212,.08532){10}{\line(0,1){.08532}}
\multiput(54.63,15.62)(.03124,.09635){9}{\line(0,1){.09635}}
\multiput(54.91,16.48)(.03008,.10991){8}{\line(0,1){.10991}}
\multiput(55.15,17.36)(.03326,.14825){6}{\line(0,1){.14825}}
\multiput(55.35,18.25)(.0316,.17956){5}{\line(0,1){.17956}}
\put(55.51,19.15){\line(0,1){.904}}
\put(55.62,20.05){\line(0,1){.909}}
\put(55.7,20.96){\line(0,1){1.139}}
%\end
%\emline(25.85,38.85)(35.85,22.35)
\multiput(25.85,38.85)(.033670034,-.055555556){297}{\line(0,-1){.055555556}}
%\end
\put(55.6,22.35){\line(-1,0){19.75}}
%\emline(41.6,23.1)(54.1,29.6)
\multiput(41.6,23.1)(.06476684,.03367876){193}{\line(1,0){.06476684}}
%\end
%\emline(36.1,23.85)(52.1,32.6)
\multiput(36.1,23.85)(.061538462,.033653846){260}{\line(1,0){.061538462}}
%\end
%\emline(33.6,26.6)(50.1,35.6)
\multiput(33.6,26.6)(.061797753,.033707865){267}{\line(1,0){.061797753}}
%\end
%\emline(32.35,30.35)(46.85,38.1)
\multiput(32.35,30.35)(.06304348,.03369565){230}{\line(1,0){.06304348}}
%\end
%\emline(30.1,33.35)(42.35,40.35)
\multiput(30.1,33.35)(.05889423,.03365385){208}{\line(1,0){.05889423}}
%\end
%\emline(28.35,36.35)(38.1,41.6)
\multiput(28.35,36.35)(.0625,.03365385){156}{\line(1,0){.0625}}
%\end
%\emline(49.35,23.35)(54.1,25.6)
\multiput(49.35,23.35)(.0708955,.0335821){67}{\line(1,0){.0708955}}
%\end
\put(39.35,27.6){\makebox(0,0)[cc]{$\theta$}}
\put(45.1,19.35){\makebox(0,0)[cc]{r}}
\end{picture}
\end{center}

The shaded area is the sector, and its area is $(\frac{1}{2}) r^{2} \theta$.

\vspace{0.3cm}

\textbf{Problem 7.15}:
\begin{enumerate}[\hspace{1cm}1)]
  \item $x = 50 \sin 20 \degrees = 17.1010$

      $\phi = 90 \degrees - 20 \degrees = 70 \degrees = 1.2217$ rad
  \item $x = \sqrt{ (25)^{2} - (12)^{2} } = 21.9317$

      $\theta = \sin^{-1} (\frac{12}{25}) = 28.6854 \degrees = .5007$ rad
  \item $x = \frac{35}{\cos 75 \degrees} = 135.2296$

      $\alpha = 90 \degrees - 75 \degrees = 15 \degrees = .2618$ rad
  \item $\alpha = sin^{-1} (\frac{50}{60}) = 56.4427 \degrees = .9851$ rad

      $\beta = \cos (\frac{50}{60}) = 33.5573 \degrees = .5857$ rad
  \item $\sin \theta = \frac{y}{r}$

      $\cos \phi = \frac{y}{r}$

      Note that $\phi = \frac{\pi}{2} - \theta$.

      Hence $\cos \phi = \cos (\frac{\pi}{2} - \theta) = \sin \theta$.
  \item $\theta = 180 \degrees - 2 (35 \degrees) = 110 \degrees$

      $x^{2} = (40)^{2} + (40)^{2} - 2 (40) (40) \cos \theta$ (Law of Cosines)

      $x = 65.5322$
  \item \begin{enumerate}[\hspace{1cm}a)]
          \item Note that $3^{2} + 4^{2} = 5^{2}$, so that the outer triangle is a right triangle. Hence $\alpha + \beta = 90 \degrees$.
          \item Now note that $\theta + \phi = 90 \degrees$.
          \item The right-hand inner triangle is isosceles. Hence, $\phi = \delta$ and $\beta = 180 \degrees - 2 \phi$ $= 180 \degrees - 2 (90 \degrees - \theta)$ $= 2 \theta$.
          \item $\gamma = 180 \degrees - \delta = 180 \degrees - \phi = 180 \degrees - (90 \degrees - \theta) = 90 \degrees + \theta$
          \item $\theta + \gamma + \alpha = 180 \degrees$.

              $\theta + (90 \degrees - \theta) + \alpha = 180 \degrees$

              $2 \theta + \alpha = 90 \degrees$ implies that $\alpha = 90 \degrees - 2 \theta$. Hence, we have $\phi = 90 \degrees - \theta$, $\delta = 90 \degrees - \theta$, $\beta = 2 \theta$, $\gamma = 90 \degrees + \theta$, and $d = 90 \degrees - 2 \theta$.
          \item Now note that $x+y=5$, so that $y=5-x$.
          \item Using the Law of Cosines, $x^{2} = 3^{2} + 4^{2} - 2(3)(4) \cos \alpha$, or $x^{2} = 25 - 24 \cos (90 \degrees - 2 \theta)$.
          \item Note that the outer triangle is a right triangle. Hence, $\cos \theta = \frac{4}{5}$ and $\sin \theta = \frac{3}{5}$.
          \item Hence, $x = 1.4$ and $y = 3.6$.
        \end{enumerate}
  \item $\theta = 180 \degrees - 2(65) = 50 \degrees$

      $(225)^{2} = x^{2} + x^{2} - 2x^{2} \cos \theta$ (Law of Cosines)

      $x = 266.1977$
  \item $QU = \frac{QR}{\cos 50 \degrees} = 62.2290$
  \item $\frac{PR}{\cos 35 \degrees}=\frac{QR}{\cos 50 \degrees}$ implies that $PR=50.9750$.
  \item \begin{enumerate}[\hspace{1cm}a)]
          \item $(PT)^{2} = (PR)^{2} + (RT)^{2}$ and $(QU)^{2} = (QR)^{2} + (RU)^{2}$
          \item $PR = PQ+QR$ and $RU = UT+TR$
        \end{enumerate}
  \item Use the same technique and substitutions as in 11.
  \item \begin{enumerate}[\hspace{1cm}a)]
          \item Consider the following picture.
        \end{enumerate}
\end{enumerate}

\begin{center}
%TeXCAD Options
%\grade{\on}
%\emlines{\off}
%\epic{\off}
%\beziermacro{\on}
%\reduce{\on}
%\snapping{\off}
%\quality{8.00}
%\graddiff{0.01}
%\snapasp{1}
%\zoom{4.0000}
\unitlength 0.9mm % = 2.85pt
\linethickness{0.4pt}
\ifx\plotpoint\undefined\newsavebox{\plotpoint}\fi % GNUPLOT compatibility
\begin{picture}(63,56.25)(0,94)
%\emline(14.5,133)(34.5,97.25)
\multiput(14.5,133)(.033726813,-.060286678){593}{\line(0,-1){.060286678}}
%\end
%\emline(34.5,97.25)(59,132.25)
\multiput(34.5,97.25)(.0337001376,.0481430536){727}{\line(0,1){.0481430536}}
%\end
\qbezier(14.5,132.75)(38.88,150.25)(58.75,131.75)
%\emline(14.75,132.75)(58.25,131.75)
\multiput(14.75,132.75)(1.45,-.033333){30}{\line(1,0){1.45}}
%\end
\put(34.5,97.75){\line(0,1){34.5}}
\put(34.75,129.5){\line(1,0){2.5}}
\put(37.25,132.25){\line(0,-1){2.5}}
\put(34.75,95){\makebox(0,0)[cc]{$O$}}
\put(32.5,106){\makebox(0,0)[cc]{$\theta$}}
\put(37.5,106){\makebox(0,0)[cc]{$\theta$}}
\put(10.25,132.75){\makebox(0,0)[cc]{$A$}}
\put(61.75,132.25){\makebox(0,0)[cc]{$B$}}
\end{picture}
\end{center}

\begin{enumerate}
  \item[\hspace{2cm}b)] Note $\sin \theta = (\frac{1}{2}) \frac{AB}{AO}$.
\end{enumerate}

\begin{enumerate}
  \item[\hspace{1cm}14)] $x = 24 \sin 45 \degrees = 24(\frac{\sqrt{2}}{2}) = 12 \sqrt{2}$
  \item[\hspace{1cm}15)] $x = \frac{20}{\sin 30 \degrees} = \frac{20}{(1/2)} = 40$

       The $45 \degrees$, $45 \degrees$, $90 \degrees$ triangle and the $30 \degrees$, $60 \degrees$, $90 \degrees$ triangle are always cropping up, and they are easy to remember.
\end{enumerate}

\begin{center}
%TeXCAD Options
%\grade{\on}
%\emlines{\off}
%\epic{\off}
%\beziermacro{\on}
%\reduce{\on}
%\snapping{\off}
%\quality{8.00}
%\graddiff{0.01}
%\snapasp{1}
%\zoom{4.0000}
\unitlength 0.9mm % = 2.85pt
\linethickness{0.4pt}
\ifx\plotpoint\undefined\newsavebox{\plotpoint}\fi % GNUPLOT compatibility
\begin{picture}(127,31.25)(10,111)
%\emline(15.25,115.25)(43.5,138.75)
\multiput(15.25,115.25)(.0405308465,.0337159254){697}{\line(1,0){.0405308465}}
%\end
%\emline(43.5,138.75)(47.25,142.25)
\multiput(43.5,138.75)(.03605769,.03365385){104}{\line(1,0){.03605769}}
%\end
\put(47.25,142.25){\line(0,-1){26.5}}
\put(47.25,115.75){\line(0,-1){1.25}}
\put(47.25,114.5){\line(-1,0){33.25}}
%\emline(14,114.5)(15.5,115.75)
\multiput(14,114.5)(.0394737,.0328947){38}{\line(1,0){.0394737}}
%\end
\put(43.75,118){\line(0,-1){3}} 
\put(43.75,118){\line(1,0){3.25}}
\put(68.25,115){\line(1,0){62.5}}
\put(130.75,115){\line(0,1){25.25}}
\put(130.75,140.25){\line(-5,-2){62.5}}
\put(127.25,118.5){\line(0,-1){3.25}}
\put(127.25,118.5){\line(1,0){3.25}}
\put(21.25,117){\makebox(0,0)[cc]{$45 \degrees$}}
\put(28.75,130.5){\makebox(0,0)[cc]{$2$}}
\put(50.5,127){\makebox(0,0)[cc]{$1$}}
\put(33.25,112){\makebox(0,0)[cc]{$1$}}
\put(102,112){\makebox(0,0)[cc]{$3$}}
\put(94.5,130.5){\makebox(0,0)[cc]{$2$}}
\put(135,127){\makebox(0,0)[cc]{$1$}}
\put(80.5,117.5){\makebox(0,0)[cc]{$30 \degrees$}}
\end{picture}
\end{center}

Any triangle similar to one of these has the same ratios of lengths and sides. For example,

\begin{center}
%TeXCAD Options
%\grade{\on}
%\emlines{\off}
%\epic{\off}
%\beziermacro{\on}
%\reduce{\on}
%\snapping{\off}
%\quality{8.00}
%\graddiff{0.01}
%\snapasp{1}
%\zoom{4.0000}
\unitlength 0.9mm % = 2.85pt
\linethickness{0.4pt}
\ifx\plotpoint\undefined\newsavebox{\plotpoint}\fi % GNUPLOT compatibility
\begin{picture}(81.5,86.25)(0,60)
%\emline(18.25,119)(49,146.25)
\multiput(18.25,119)(.0380569307,.0337252475){808}{\line(1,0){.0380569307}}
%\end
\put(49,146.25){\line(0,-1){28.5}}
\put(49,117.75){\line(-1,0){32.25}}
%\emline(17,117.75)(18.25,119.25)
\multiput(17,117.75)(.0328947,.0394737){38}{\line(0,1){.0394737}}
%\end
\put(45.5,120.5){\line(0,-1){3}} 
\put(45.75,120.75){\line(1,0){3}}
\put(19.75,64.25){\line(1,0){48.25}} 
\put(68,64.25){\line(0,1){29}}
%\emline(68,93.25)(16.25,64.25)
\multiput(68,93.25)(-.0601744186,-.0337209302){860}{\line(-1,0){.0601744186}}
%\end
\put(16.25,64.25){\line(1,0){4.75}}
\put(64.75,67.25){\line(0,-1){2.75}}
\put(64.75,67.5){\line(1,0){3.5}}
\put(30.75,133.5){\makebox(0,0)[cc]{$y$}}
\put(51.5,130.75){\makebox(0,0)[cc]{$7$}}
\put(75,130.75){\makebox(0,0)[cc]{$x=7$ and $y=7\sqrt{2}$}}
\put(34,115){\makebox(0,0)[cc]{$x$}}
\put(24.25,120.25){\makebox(0,0)[cc]{$45 \degrees$}}
\put(15.25,98.5){\makebox(0,0)[cc]{also}}
\put(33,77.25){\makebox(0,0)[cc]{$y$}}
\put(64,87.75){\makebox(0,0)[cc]{$60 \degrees$}}
\put(70.75,78.5){\makebox(0,0)[cc]{$7$}}
\put(39,61.75){\makebox(0,0)[cc]{$x$}}
\put(85,86.5){\makebox(0,0)[cc]{$x = 7 \sqrt{3}$}}
\put(95,73){\makebox(0,0)[cc]{and $y = (2)(7) = 14$}}
\end{picture}
\end{center}

\begin{enumerate}
  \item[\hspace{1cm}16)] From the discussion above, $x = 5$.
  \item[\hspace{1cm}17)] $x = 22$
  \item[\hspace{1cm}18)] $\sin 15 \degrees = .2588$

       $\cos 15 \degrees = .9659$

       Now, note that $15 = (1/2) 30$. Hence,

       $\sin 15 \degrees = \sin (\frac{30 \degrees}{2}) = \sqrt{\frac{(1 - \cos 30 \degrees)}{2}} = \sqrt{ \frac{(1-\sqrt{3} }{2}/2 }$.

       Similarly with $\cos 15 \degrees$. You might also note that 15 = 60-45.
  \item[\hspace{1cm}19)]  $\sin 22 \degrees 30' = .3827$

       $\cos 22 \degrees 30' = .9239$

       Note $22 \degrees 30' = (\frac{1}{2}) 45 \degrees$.
  \item[\hspace{1cm}20)] Note $75 = 45+30$
  \item[\hspace{1cm}21)] $105 = 60+45$
  \item[\hspace{1cm}22)] $37 \degrees 30' = \frac{75 \degrees}{2}$
  \item[\hspace{1cm}23)] $52 \degrees 30' = \frac{105 \degrees}{2}$
  \item[\hspace{1cm}24)] $\theta = 30.9306 \degrees = .5398$ rad
  \item[\hspace{1cm}25)] $\theta = 14.4775 \degrees = .2527$ rad
  \item[\hspace{1cm}26)] $\theta = 30.7955 \degrees = .5375$ rad
  \item[\hspace{1cm}27)] $\theta = 48.1897 \degrees = .8411$ rad
  \item[\hspace{1cm}28)] $\theta = 8.5110 \degrees = .1485$ rad
  \item[\hspace{1cm}29)] $\theta = 75.6141 \degrees = 1.3197$ rad
  \item[\hspace{1cm}30)] \begin{enumerate}[\hspace{1cm}a)] 
                           \item $QR = PR \sin (m \angm P) = 82.4126$
                           \item $PQ = PR \cos (m \angm P) = 56.6406$
                           \item $PS = PQ \cos (m \angm P) = PR \cos^{2} (m \angm P) = 32.0816$
                           \item $QS = PQ \sin (m \angm P) = PR \cos (m \angm P) \sin (m \angm P) = 46.6790$
                           \item $m \angm PQS = 90 \degrees - (m \angm P) = 34.5 \degrees$
                           \item $m \angm RQS = 90 \degrees - (m \angm R) = 90 \degrees - (90 \degrees - (m \angm P) ) = m \angm P$
                           \item $SR = QR \cos (m \angm R)$

                               $= PR \sin (m \angm P) \cos (90 \degrees - (m \angm P) ) = 67.9184$
                         \end{enumerate}
  \item[\hspace{1cm}31)] \begin{enumerate}[\hspace{1cm}a)]
                           \item $\delta = \theta$ , $\delta = 180 \degrees - \theta$, $\alpha = 180 \degrees - 2 \theta$
                           \item $\gamma = 180 \degrees - 2 \theta$, $\beta = \gamma = 180 \degrees - 2 \theta$
                           \item $\alpha + \beta = \theta$ implies that $\beta = \theta - \alpha = \theta -(180 \degrees - 2 \theta) = 3 \theta - 180 \degrees$
                           \item $\beta = 180 \degrees - 2 \theta$ and $\beta = 3 \theta - 180 \degrees$ implies that $\theta = 72 \degrees$
                           \item Use Law of Cosines to get $1^{2} = x^{2} + 1^{2} -2(x)(1) \cos \theta$, which implies that $x = .6180$
                           \item Find a way to show that $x = \frac{(-1 + \sqrt{5})}{2}$.
                           \item Then show that $\cos 72 \degrees = \frac{ (-1 + \sqrt{5}) }{4}$ so that 

                               $\cos 36 \degrees = (\frac{1}{2}) \sqrt{ (3+\sqrt{5})/2 }$.
                         \end{enumerate}
  \item[\hspace{1cm}32)] \begin{enumerate}[\hspace{1cm}a)]
                           \item Draw the picture large enough so that you have room to maneuver.
                           \item Let $\beta$ be the angle that the light path makes with the normal as it travels through the glass. Snell's Law says $\frac{\sin \alpha}{\sin \beta} = \frac{ s_{1} }{ s_{2} }$.
                           \item Now let $\gamma$ be the angle that the light path makes with the normal as it emerges from the glass into the air again. Snell's Law says $\frac{\sin \beta}{\sin \gamma} = \frac{ s_{2} }{ s_{3} }$.
                           \item What is $s_{3}$?
                           \item $s_{3}$ is the speed of light in air, so it must be $s_{1}$. It won't travel at a different speed above the glass than below it.
                           \item Hence, $\frac{\sin \beta}{\sin \gamma} = \frac{ s_{2} }{ s_{1} }$. Solve for $\sin \gamma$.
                           \item From your first application of Snell's Law, solve for $\sin \alpha$.
                           \item You should have that $\sin \gamma = \sin \alpha$.
                           \item Does it follow that $\gamma = \alpha$? This is certainly the case if the angles involved must be between $0 \degrees$ and $90 \degrees$. Are they?
                         \end{enumerate}
  \item[\hspace{1cm}33)] \begin{enumerate}[\hspace{1cm}a)]
                           \item Use either the Law of Sines or the Law of Cosines.
                           \item Note that $m \angm CAB = 180 \degrees - (80 \degrees + 75 \degrees)$.
                           \item $\frac{\sin (m \angm A)}{100} = \frac{\sin 80 \degrees}{AB}$.
                           \item $AB = 233.0254$.
                         \end{enumerate}
  \item[\hspace{1cm}34)] \begin{enumerate}[\hspace{1cm}a)]
                           \item $\tan 20 \degrees = \frac{(150-h)}{300}$.
                           \item $h = 40.8089$.
                         \end{enumerate}
  \item[\hspace{1cm}35)] \begin{enumerate}[\hspace{1cm}a)]
                           \item Use Law of Cosines to find $x$.
                           \item Which solution for $x$ must you use in your resulting quadratic in $x$?
                           \item $x = 10.1989$. Why not choose $x = 2.7454$?
                           \item Use Law of Sines to get $\alpha$.
                           \item $\frac{\sin \alpha}{x} = \frac{\sin 36 \degrees}{6}$.
                           \item $\alpha = 92.3981$ or $\alpha = 92.3941$, depending on whether you used $x$ rounded to 4 decimal places or did not round $x$.
                           \item Did you get $\alpha = 87.6019 \degrees$? Why is this not correct?
                           \item Note, or recall, that $\sin^{-1}$ is a function when we restrict ourselves to the interval $[-\frac{\pi}{2},\frac{\pi}{2}]$, or, in degrees, $[-90 \degrees , 90 \degrees]$. We are looking for a solution to $\sin \alpha = \frac{(x \sin 36 \degrees)}{6}$ \emph{outside} of that interval, because we are given that $\alpha > 90 \degrees$.
                           \item If you had used the Law of Cosines instead of the Law of Sines, would this curiosity have been avoided?
                         \end{enumerate}
  \item[\hspace{1cm}36)] $x = 22.2320$
  \item[\hspace{1cm}37)] $h = 100 \sin 35 \degrees = 57.3576$

                       $x = 100 \cos 35 \degrees = 81.9152$

                       $z = 100 \tan 35 \degrees = 70.0208$

                       $y = z \cos 55 \degrees = 40.1623$
  \item[\hspace{1cm}38)] $x + 20 = 40 \sin 48.590378 \degrees$
  \item[\hspace{1cm}39)] \begin{enumerate}[\hspace{1cm}a)]
                           \item Consider the figure
                         \end{enumerate}
\end{enumerate}

\begin{center}
%TeXCAD Picture [prob7.15hint_39.pic]. Options:
%\grade{\on}
%\emlines{\off}
%\epic{\off}
%\beziermacro{\on}
%\reduce{\on}
%\snapping{\off}
%\quality{8.00}
%\graddiff{0.01}
%\snapasp{1}
%\zoom{4.0000}
\unitlength 1mm % = 2.85pt
\linethickness{0.4pt}
\ifx\plotpoint\undefined\newsavebox{\plotpoint}\fi % GNUPLOT compatibility
\begin{picture}(64,32.94)(0,0)
%\emline(17.5,8.69)(64,32.94)
\multiput(17.5,8.69)(.0646731572,.0337273992){719}{\line(1,0){.0646731572}}
%\end
%\emline(14.5,6.94)(17.5,8.69)
\multiput(14.5,6.94)(.0576923,.0336538){52}{\line(1,0){.0576923}}
%\end
\put(57.5,26.94){\makebox(0,0)[cc]{$\beta$}}
\put(60.75,17.44){\makebox(0,0)[cc]{x}}
\put(51,9.69){\makebox(0,0)[cc]{$\alpha$}}
\put(26.5,9.44){\makebox(0,0)[cc]{$37.1 \degrees$}}
\put(37.5,3.44){\makebox(0,0)[cc]{60}}
%\emline(64,32.5)(53.75,7)
\multiput(64,32.5)(-.033717105,-.083881579){304}{\line(0,-1){.083881579}}
%\end
\put(58.25,27.75){\line(0,1){.25}}
\put(14.57,6.84){\line(1,0){39.1}}
\end{picture}
\end{center}

    \begin{enumerate}
      \item[\hspace{2cm}b)] Determine $\alpha$ and then $\beta$.
      \item[\hspace{2cm}c)] Use the Law of Sines to calculate $x$.
      \item[\hspace{2cm}d)] $\alpha = 117.7 \degrees$ implies that $\beta = 25.2 \degrees$
      \item[\hspace{2cm}e)] $x = 60 (\frac{\sin 37.1 \degrees}{\sin 25.2 \degrees}) = 85.0029$
    \end{enumerate}

\begin{enumerate}
  \item[\hspace{1cm}40)] Use Law of Cosines. $x = 90.2815$
  \item[\hspace{1cm}41)] \begin{enumerate}[\hspace{1cm}a)]
                           \item Consider the figure
                         \end{enumerate}
\end{enumerate}


\begin{center}
%TeXCAD Options
%\grade{\on}
%\emlines{\off}
%\epic{\off}
%\beziermacro{\on}
%\reduce{\on}
%\snapping{\off}
%\quality{8.00}
%\graddiff{0.01}
%\snapasp{1}
%\zoom{4.0000}
\unitlength 1mm % = 2.85pt
\linethickness{0.4pt}
\ifx\plotpoint\undefined\newsavebox{\plotpoint}\fi % GNUPLOT compatibility
\begin{picture}(72.5,36.5)(0,115)
\put(15.75,115.75){\line(1,0){53.25}}
\put(69,115.75){\line(0,1){35.25}} 
\put(69,151){\line(-5,-4){32.5}}
%\emline(15.5,115.75)(69,132.5)
\multiput(15.5,115.75)(.107645875,.033702213){497}{\line(1,0){.107645875}}
%\end
%\emline(36.75,125)(31,120.75)
\multiput(36.75,125)(-.04563492,-.03373016){126}{\line(-1,0){.04563492}}
%\end
\put(65.5,118.75){\line(0,-1){2.5}} 
\put(66,119){\line(1,0){2.75}}
\put(65.5,116.5){\line(0,-1){.5}}
\put(45.75,136.75){\makebox(0,0)[cc]{$x$}}
\put(71.75,141){\makebox(0,0)[cc]{$30$}}
\put(56.5,126.5){\makebox(0,0)[cc]{$40$}}
\put(28.25,117.25){\makebox(0,0)[cc]{$20 \degrees$}}
\put(72.5,126.5){\makebox(0,0)[cc]{$\alpha$}}
\put(66.5,134.25){\makebox(0,0)[cc]{$\beta$}}
%\qbezvec(72,123.75)(66,120.88)(66,129.5)
\put(66,129.5){\vector(0,1){.07}}\qbezier(72,123.75)(66,120.88)(66,129.5)
%\end
\put(60,128){\vector(3,1){7.5}}
%\vector(53.25,126.25)(33.75,120.25)
\put(33.75,120.25){\vector(-3,-1){.07}}\multiput(53.25,126.25)(-.10955056,-.03370787){178}{\line(-1,0){.10955056}}
%\end
\put(71.75,140){\vector(0,-1){6.75}}
\put(71.75,144.75){\vector(0,1){6.75}}
\end{picture}
\end{center}

    \begin{enumerate}
      \item[\hspace{2cm}b)] Determine $\alpha$, and consequently $\beta$.
      \item[\hspace{2cm}c)] Use the Law of Cosines.
      \item[\hspace{2cm}d)] $x = 59.5994$
    \end{enumerate}

\begin{enumerate}
  \item[\hspace{1cm}42)] \begin{enumerate}[\hspace{1cm}a)]
                           \item Consider the figure
                         \end{enumerate}
\end{enumerate}


\begin{center}
%TeXCAD Options
%\grade{\on}
%\emlines{\off}
%\epic{\off}
%\beziermacro{\on}
%\reduce{\on}
%\snapping{\off}
%\quality{8.00}
%\graddiff{0.01}
%\snapasp{1}
%\zoom{4.0000}
\unitlength 1mm % = 2.85pt
\linethickness{0.4pt}
\ifx\plotpoint\undefined\newsavebox{\plotpoint}\fi % GNUPLOT compatibility
\begin{picture}(58.5,31.5)(0,116)
\put(15.75,120.5){\line(1,0){39.25}} 
\put(55,120.5){\line(0,1){27}}
%\emline(55,147.5)(15.75,120.75)
\multiput(55,147.5)(-.0494955864,-.0337326608){793}{\line(-1,0){.0494955864}}
%\end
%\emline(15.75,120.75)(55,130.25)
\multiput(15.75,120.75)(.139184397,.033687943){282}{\line(1,0){.139184397}}
%\end
\put(52.25,123){\line(0,-1){2.25}} 
\put(52.5,123){\line(1,0){2.75}}
\qbezier(23.75,122.5)(24.75,121.5)(23.75,120.5)
\qbezier(29.5,120.5)(32.63,125.5)(27.25,128.5)
\put(39.25,117){\makebox(0,0)[cc]{250}}
\put(58.5,125.5){\makebox(0,0)[cc]{$x$}}
\put(57.5,140){\makebox(0,0)[cc]{$h$}}
\put(26.25,125.5){\makebox(0,0)[cc]{$60 \degrees$}}
\put(27.25,121.75){\makebox(0,0)[cc]{$30 \degrees$}}
\end{picture}
\\(Note that $\tan 60 \degrees = \sqrt{3}$)
\end{center}

    \begin{enumerate}
      \item[\hspace{2cm}b)] Calculate $x$
      \item[\hspace{2cm}c)] $x = \frac{250}{\sqrt{3}}$
      \item[\hspace{2cm}d)] $\tan 60 \degrees = \frac{(h+x)}{250}$
      \item[\hspace{2cm}e)] $h = 288.6751$
    \end{enumerate}

\begin{enumerate}
  \item[\hspace{1cm}43)] $x = 100 \tan 70 \degrees = 274.7477$
  \item[\hspace{1cm}44)] $x = \frac{300}{\tan 38.5} \degrees = 377.1517$
  \item[\hspace{1cm}45)] $= sin^{-1} (\frac{24}{30}) = 53.1301 \degrees = .9273$ rad
  \item[\hspace{1cm}46)] $h = 300 \tan 35 \degrees = 210.0623$
  \item[\hspace{1cm}47)] \begin{enumerate}[\hspace{1cm}a)]
                           \item The whole circle has $360 \degrees$. $2 \pi r = 360 \degrees$.
                           \item Consider the figure
                         \end{enumerate}
\end{enumerate}


\begin{center}
%TeXCAD Picture [prob7.15hint_47.pic]. Options:
%\grade{\on}
%\emlines{\off}
%\epic{\off}
%\beziermacro{\on}
%\reduce{\on}
%\snapping{\off}
%\quality{8.00}
%\graddiff{0.01}
%\snapasp{1}
%\zoom{4.0000}
\unitlength 1mm % = 2.85pt
\linethickness{0.4pt}
\ifx\plotpoint\undefined\newsavebox{\plotpoint}\fi % GNUPLOT compatibility
\begin{picture}(58.09,54.59)(0,0)
%\circle(31.75,28.25){52.69}
\put(58.09,28.25){\line(0,1){1.113}}
\put(58.07,29.36){\line(0,1){1.111}}
\put(58,30.47){\line(0,1){1.107}}
\multiput(57.88,31.58)(-.03281,.2202){5}{\line(0,1){.2202}}
\multiput(57.72,32.68)(-.03006,.15616){7}{\line(0,1){.15616}}
\multiput(57.51,33.77)(-.03205,.13541){8}{\line(0,1){.13541}}
\multiput(57.25,34.86)(-.03355,.11905){9}{\line(0,1){.11905}}
\multiput(56.95,35.93)(-.03154,.09616){11}{\line(0,1){.09616}}
\multiput(56.6,36.99)(-.03261,.08684){12}{\line(0,1){.08684}}
\multiput(56.21,38.03)(-.033459,.078821){13}{\line(0,1){.078821}}
\multiput(55.78,39.05)(-.031858,.067026){15}{\line(0,1){.067026}}
\multiput(55.3,40.06)(-.032495,.061519){16}{\line(0,1){.061519}}
\multiput(54.78,41.04)(-.033002,.056557){17}{\line(0,1){.056557}}
\multiput(54.22,42.01)(-.033397,.05205){18}{\line(0,1){.05205}}
\multiput(53.62,42.94)(-.033694,.04793){19}{\line(0,1){.04793}}
\multiput(52.98,43.85)(-.03229,.042039){21}{\line(0,1){.042039}}
\multiput(52.3,44.74)(-.03249,.03879){22}{\line(0,1){.03879}}
\multiput(51.58,45.59)(-.032617,.035758){23}{\line(0,1){.035758}}
\multiput(50.83,46.41)(-.032678,.032917){24}{\line(0,1){.032917}}
\multiput(50.05,47.2)(-.035519,.032877){23}{\line(-1,0){.035519}}
\multiput(49.23,47.96)(-.038552,.032772){22}{\line(-1,0){.038552}}
\multiput(48.39,48.68)(-.041802,.032595){21}{\line(-1,0){.041802}}
\multiput(47.51,49.36)(-.045299,.03234){20}{\line(-1,0){.045299}}
\multiput(46.6,50.01)(-.049079,.031998){19}{\line(-1,0){.049079}}
\multiput(45.67,50.62)(-.056314,.033413){17}{\line(-1,0){.056314}}
\multiput(44.71,51.19)(-.06128,.032942){16}{\line(-1,0){.06128}}
\multiput(43.73,51.71)(-.066792,.032346){15}{\line(-1,0){.066792}}
\multiput(42.73,52.2)(-.072963,.031602){14}{\line(-1,0){.072963}}
\multiput(41.71,52.64)(-.0866,.03324){12}{\line(-1,0){.0866}}
\multiput(40.67,53.04)(-.09593,.03224){11}{\line(-1,0){.09593}}
\multiput(39.61,53.39)(-.10692,.03097){10}{\line(-1,0){.10692}}
\multiput(38.54,53.7)(-.13517,.03304){8}{\line(-1,0){.13517}}
\multiput(37.46,53.97)(-.15594,.0312){7}{\line(-1,0){.15594}}
\multiput(36.37,54.19)(-.1833,.02868){6}{\line(-1,0){.1833}}
\multiput(35.27,54.36)(-.2765,.0314){4}{\line(-1,0){.2765}}
\put(34.17,54.48){\line(-1,0){1.11}}
\put(33.05,54.56){\line(-1,0){1.113}}
\put(31.94,54.59){\line(-1,0){1.113}}
\put(30.83,54.58){\line(-1,0){1.111}}
\put(29.72,54.52){\line(-1,0){1.108}}
\multiput(28.61,54.41)(-.22044,-.0312){5}{\line(-1,0){.22044}}
\multiput(27.51,54.25)(-.18244,-.03374){6}{\line(-1,0){.18244}}
\multiput(26.41,54.05)(-.13564,-.03106){8}{\line(-1,0){.13564}}
\multiput(25.33,53.8)(-.11929,-.03268){9}{\line(-1,0){.11929}}
\multiput(24.25,53.51)(-.09639,-.03084){11}{\line(-1,0){.09639}}
\multiput(23.19,53.17)(-.08708,-.03197){12}{\line(-1,0){.08708}}
\multiput(22.15,52.78)(-.079063,-.032884){13}{\line(-1,0){.079063}}
\multiput(21.12,52.36)(-.07206,-.033609){14}{\line(-1,0){.07206}}
\multiput(20.11,51.89)(-.061754,-.032045){16}{\line(-1,0){.061754}}
\multiput(19.12,51.37)(-.056796,-.032589){17}{\line(-1,0){.056796}}
\multiput(18.16,50.82)(-.052292,-.033017){18}{\line(-1,0){.052292}}
\multiput(17.22,50.22)(-.048174,-.033344){19}{\line(-1,0){.048174}}
\multiput(16.3,49.59)(-.044387,-.033582){20}{\line(-1,0){.044387}}
\multiput(15.41,48.92)(-.040884,-.03374){21}{\line(-1,0){.040884}}
\multiput(14.56,48.21)(-.035994,-.032355){23}{\line(-1,0){.035994}}
\multiput(13.73,47.47)(-.033154,-.032437){24}{\line(-1,0){.033154}}
\multiput(12.93,46.69)(-.033135,-.035278){23}{\line(0,-1){.035278}}
\multiput(12.17,45.88)(-.033052,-.038312){22}{\line(0,-1){.038312}}
\multiput(11.44,45.03)(-.032899,-.041564){21}{\line(0,-1){.041564}}
\multiput(10.75,44.16)(-.03267,-.045062){20}{\line(0,-1){.045062}}
\multiput(10.1,43.26)(-.032355,-.048844){19}{\line(0,-1){.048844}}
\multiput(9.48,42.33)(-.031944,-.052954){18}{\line(0,-1){.052954}}
\multiput(8.91,41.38)(-.033388,-.061039){16}{\line(0,-1){.061039}}
\multiput(8.37,40.4)(-.032832,-.066554){15}{\line(0,-1){.066554}}
\multiput(7.88,39.4)(-.032133,-.072731){14}{\line(0,-1){.072731}}
\multiput(7.43,38.39)(-.031265,-.079717){13}{\line(0,-1){.079717}}
\multiput(7.03,37.35)(-.03294,-.09569){11}{\line(0,-1){.09569}}
\multiput(6.66,36.3)(-.03175,-.10669){10}{\line(0,-1){.10669}}
\multiput(6.35,35.23)(-.03024,-.11993){9}{\line(0,-1){.11993}}
\multiput(6.07,34.15)(-.03233,-.1557){7}{\line(0,-1){.1557}}
\multiput(5.85,33.06)(-.03001,-.18309){6}{\line(0,-1){.18309}}
\multiput(5.67,31.96)(-.0334,-.2763){4}{\line(0,-1){.2763}}
\put(5.53,30.86){\line(0,-1){1.11}}
\put(5.45,29.75){\line(0,-1){2.226}}
\put(5.42,27.52){\line(0,-1){1.112}}
\put(5.47,26.41){\line(0,-1){1.109}}
\multiput(5.57,25.3)(.02959,-.22066){5}{\line(0,-1){.22066}}
\multiput(5.72,24.2)(.03241,-.18268){6}{\line(0,-1){.18268}}
\multiput(5.91,23.1)(.03007,-.13586){8}{\line(0,-1){.13586}}
\multiput(6.15,22.01)(.03181,-.11953){9}{\line(0,-1){.11953}}
\multiput(6.44,20.94)(.03315,-.10627){10}{\line(0,-1){.10627}}
\multiput(6.77,19.88)(.03134,-.08731){12}{\line(0,-1){.08731}}
\multiput(7.15,18.83)(.032306,-.079301){13}{\line(0,-1){.079301}}
\multiput(7.57,17.8)(.033083,-.072303){14}{\line(0,-1){.072303}}
\multiput(8.03,16.78)(.033701,-.066119){15}{\line(0,-1){.066119}}
\multiput(8.54,15.79)(.032174,-.057032){17}{\line(0,-1){.057032}}
\multiput(9.08,14.82)(.032635,-.052531){18}{\line(0,-1){.052531}}
\multiput(9.67,13.88)(.032992,-.048416){19}{\line(0,-1){.048416}}
\multiput(10.3,12.96)(.033257,-.04463){20}{\line(0,-1){.04463}}
\multiput(10.96,12.07)(.033441,-.041129){21}{\line(0,-1){.041129}}
\multiput(11.66,11.2)(.033551,-.037876){22}{\line(0,-1){.037876}}
\multiput(12.4,10.37)(.033594,-.034841){23}{\line(0,-1){.034841}}
\multiput(13.18,9.57)(.035036,-.033391){23}{\line(1,0){.035036}}
\multiput(13.98,8.8)(.03807,-.03333){22}{\line(1,0){.03807}}
\multiput(14.82,8.07)(.041323,-.033201){21}{\line(1,0){.041323}}
\multiput(15.69,7.37)(.044823,-.032997){20}{\line(1,0){.044823}}
\multiput(16.58,6.71)(.048607,-.03271){19}{\line(1,0){.048607}}
\multiput(17.51,6.09)(.05272,-.032329){18}{\line(1,0){.05272}}
\multiput(18.46,5.51)(.057218,-.031842){17}{\line(1,0){.057218}}
\multiput(19.43,4.96)(.066313,-.033316){15}{\line(1,0){.066313}}
\multiput(20.42,4.46)(.072494,-.032662){14}{\line(1,0){.072494}}
\multiput(21.44,4.01)(.079487,-.031845){13}{\line(1,0){.079487}}
\multiput(22.47,3.59)(.09545,-.03363){11}{\line(1,0){.09545}}
\multiput(23.52,3.22)(.10646,-.03253){10}{\line(1,0){.10646}}
\multiput(24.59,2.9)(.11971,-.03111){9}{\line(1,0){.11971}}
\multiput(25.66,2.62)(.15546,-.03347){7}{\line(1,0){.15546}}
\multiput(26.75,2.38)(.18286,-.03135){6}{\line(1,0){.18286}}
\multiput(27.85,2.2)(.22083,-.02831){5}{\line(1,0){.22083}}
\put(28.95,2.05){\line(1,0){1.109}}
\put(30.06,1.96){\line(1,0){1.112}}
\put(31.17,1.91){\line(1,0){1.113}}
\put(32.29,1.91){\line(1,0){1.112}}
\put(33.4,1.96){\line(1,0){1.109}}
\multiput(34.51,2.05)(.22087,.02798){5}{\line(1,0){.22087}}
\multiput(35.61,2.19)(.18291,.03108){6}{\line(1,0){.18291}}
\multiput(36.71,2.38)(.15551,.03324){7}{\line(1,0){.15551}}
\multiput(37.8,2.61)(.11975,.03094){9}{\line(1,0){.11975}}
\multiput(38.88,2.89)(.10651,.03237){10}{\line(1,0){.10651}}
\multiput(39.94,3.21)(.09549,.03349){11}{\line(1,0){.09549}}
\multiput(40.99,3.58)(.079534,.031728){13}{\line(1,0){.079534}}
\multiput(42.03,3.99)(.072543,.032555){14}{\line(1,0){.072543}}
\multiput(43.04,4.45)(.066362,.033218){15}{\line(1,0){.066362}}
\multiput(44.04,4.95)(.057265,.031757){17}{\line(1,0){.057265}}
\multiput(45.01,5.49)(.052768,.032251){18}{\line(1,0){.052768}}
\multiput(45.96,6.07)(.048655,.032638){19}{\line(1,0){.048655}}
\multiput(46.89,6.69)(.044872,.032931){20}{\line(1,0){.044872}}
\multiput(47.78,7.35)(.041372,.03314){21}{\line(1,0){.041372}}
\multiput(48.65,8.04)(.03812,.033274){22}{\line(1,0){.03812}}
\multiput(49.49,8.77)(.035085,.033339){23}{\line(1,0){.035085}}
\multiput(50.3,9.54)(.033645,.034791){23}{\line(0,1){.034791}}
\multiput(51.07,10.34)(.033607,.037826){22}{\line(0,1){.037826}}
\multiput(51.81,11.17)(.033502,.04108){21}{\line(0,1){.04108}}
\multiput(52.51,12.03)(.033323,.044581){20}{\line(0,1){.044581}}
\multiput(53.18,12.93)(.033063,.048367){19}{\line(0,1){.048367}}
\multiput(53.81,13.85)(.032712,.052483){18}{\line(0,1){.052483}}
\multiput(54.4,14.79)(.032258,.056984){17}{\line(0,1){.056984}}
\multiput(54.95,15.76)(.031686,.061939){16}{\line(0,1){.061939}}
\multiput(55.45,16.75)(.03319,.072254){14}{\line(0,1){.072254}}
\multiput(55.92,17.76)(.032424,.079253){13}{\line(0,1){.079253}}
\multiput(56.34,18.79)(.03147,.08726){12}{\line(0,1){.08726}}
\multiput(56.72,19.84)(.0333,.10622){10}{\line(0,1){.10622}}
\multiput(57.05,20.9)(.03199,.11948){9}{\line(0,1){.11948}}
\multiput(57.34,21.98)(.03027,.13581){8}{\line(0,1){.13581}}
\multiput(57.58,23.06)(.03268,.18263){6}{\line(0,1){.18263}}
\multiput(57.78,24.16)(.02992,.22062){5}{\line(0,1){.22062}}
\put(57.92,25.26){\line(0,1){1.108}}
\put(58.03,26.37){\line(0,1){1.88}}
%\end
%\emline(14.5,48.25)(9,15)
\multiput(14.5,48.25)(-.03353659,-.2027439){164}{\line(0,-1){.2027439}}
%\end
%\emline(9,15)(40,3.75)
\multiput(9,15)(.092814371,-.033682635){334}{\line(1,0){.092814371}}
%\end
%\emline(40,3.75)(58,29.75)
\multiput(40,3.75)(.033707865,.048689139){534}{\line(0,1){.048689139}}
%\end
%\emline(58,29.75)(41.25,53)
\multiput(58,29.75)(-.033702213,.046780684){497}{\line(0,1){.046780684}}
%\end
%\emline(41.25,53)(14.75,48.25)
\multiput(41.25,53)(-.18794326,-.03368794){141}{\line(-1,0){.18794326}}
%\end
%\emline(31.75,28.25)(40,3.75)
\multiput(31.75,28.25)(.033673469,-.1){245}{\line(0,-1){.1}}
%\end
%\emline(32,28.25)(57.75,29.75)
\multiput(32,28.25)(.5722222,.0333333){45}{\line(1,0){.5722222}}
%\end
\qbezier(33.75,23.25)(38,23.63)(37.25,28.5)
\put(43,31){\makebox(0,0)[cc]{3}}
\put(35,26){\makebox(0,0)[cc]{$\alpha$}}
\put(50.5,14){\makebox(0,0)[cc]{$x$}}
\put(23.25,7.25){\makebox(0,0)[cc]{$x$}}
\put(8,31){\makebox(0,0)[cc]{$x$}}
\put(26.25,52){\makebox(0,0)[cc]{$x$}}
\put(51.75,42.5){\makebox(0,0)[cc]{$x$}}
\end{picture}
\end{center}

    \begin{enumerate}
      \item[\hspace{2cm}c)] Calculate $\alpha$.
      \item[\hspace{2cm}d)] Use the Law of Cosines.
      \item[\hspace{2cm}e)] $\alpha = 72$, so $x^{2} = 3^{2} + 3^{2} - 2(3)(3) \cos 72 \degrees$, so that $x = 3.5267$
    \end{enumerate}

\begin{enumerate}
  \item[\hspace{1cm}48)] \begin{enumerate}[\hspace{1cm}a)]
                           \item Use the large triangle and show that $200 + h = h \sqrt{3}$.
                           \item $h = 273.2051$
                         \end{enumerate}
  \item[\hspace{1cm}49)] \begin{enumerate}[\hspace{1cm}a)]
                           \item Consider the figure
                         \end{enumerate}
\end{enumerate}


\begin{center}
%TeXCAD Options
%\grade{\on}
%\emlines{\off}
%\epic{\off}
%\beziermacro{\on}
%\reduce{\on}
%\snapping{\off}
%\quality{8.00}
%\graddiff{0.01}
%\snapasp{1}
%\zoom{4.0000}
\unitlength 1mm % = 2.85pt
\linethickness{0.4pt}
\ifx\plotpoint\undefined\newsavebox{\plotpoint}\fi % GNUPLOT compatibility
\begin{picture}(64.75,32.63)(0,110)
\put(12.5,115.88){\line(1,0){52.25}}
\put(38,115.88){\line(0,1){26.75}}
%\emline(38,142.63)(55.5,130.38)
\multiput(38,142.63)(.048076923,-.033653846){364}{\line(1,0){.048076923}}
%\end
\put(55.5,130.38){\line(0,-1){14.25}}
\put(52,119.63){\line(0,-1){3.5}} 
\put(52,119.63){\line(1,0){3.25}}
\put(34.5,119.38){\line(0,-1){3.25}}
\put(34.75,119.38){\line(1,0){3.25}}
\qbezier(14,126.13)(41,100)(62,142.38)
\put(57.75,122.38){\makebox(0,0)[cc]{$h$}}
\put(49.5,138.38){\makebox(0,0)[cc]{$L$}}
\put(39.75,139.13){\makebox(0,0)[cc]{$\alpha$}}
\qbezier(38.25,136.88)(41.63,136.38)(43.5,138.88)
\put(35,130){\makebox(0,0)[cc]{$L$}}
\put(32.5,111){\makebox(0,0)[cc]{$\beta$}}
\put(40.25,120.5){\makebox(0,0)[cc]{$\theta$}}
\put(50.75,129.5){\makebox(0,0)[cc]{$\theta$}}
\put(45.75,124.25){\makebox(0,0)[cc]{x}}
%\emline(38.25,116)(55.25,130)
\multiput(38.25,116)(.040963855,.03373494){415}{\line(1,0){.040963855}}
%\end
\qbezier(42.5,119.25)(45.38,119.5)(42.75,115.75)
%\qbezvec(37,112)(42.38,112.13)(42.25,116.75)
\put(42.25,116.75){\vector(0,1){.07}}\qbezier(37,112)(42.38,112.13)(42.25,116.75)
%\end
\end{picture}
\end{center}

    \begin{enumerate}
      \item[\hspace{2cm}b)] Calculate $\theta$ in terms of $\alpha$
      \item[\hspace{2cm}c)] Calculate $\beta$ in terms of $\alpha$
      \item[\hspace{2cm}d)] Show $h = x \sin \beta$
      \item[\hspace{2cm}e)] Calculate $x$ using the Law of Cosines. You can get the result you want by using the Law of Sines, but it's harder.
      \item[\hspace{2cm}f)] You should get $x = L \sqrt{2} \sqrt{ 1 - \cos \alpha }$.
      \item[\hspace{2cm}g)] You should also have gotten $\beta = \frac{\alpha}{2}$.
      \item[\hspace{2cm}h)] Recall that $\frac{\sin \alpha}{2} = \sqrt{ \frac{(1 - \cos \alpha)}{2} }$.
    \end{enumerate}

\begin{enumerate}
  \item[\hspace{1cm}50)] \begin{enumerate}[\hspace{1cm}a)]
                           \item Redraw the picture
                         \end{enumerate}
\end{enumerate}


\begin{center}
%TeXCAD Picture [prob7.15hint_50.pic]. Options:
%\grade{\on}
%\emlines{\off}
%\epic{\off}
%\beziermacro{\on}
%\reduce{\on}
%\snapping{\off}
%\quality{8.00}
%\graddiff{0.01}
%\snapasp{1}
%\zoom{6.7272}
\unitlength 1mm % = 2.85pt
\linethickness{0.4pt}
\ifx\plotpoint\undefined\newsavebox{\plotpoint}\fi % GNUPLOT compatibility
\begin{picture}(71.5,37.5)(0,107.5)
%\emline(14.25,112.25)(52,145)
\multiput(14.25,112.25)(.0388774459,.0337281153){971}{\line(1,0){.0388774459}}
%\end
%\emline(52,145)(70.5,113)
\multiput(52,145)(.033697632,-.058287796){549}{\line(0,-1){.058287796}}
%\end
%\emline(70.5,113)(71.5,111.25)
\multiput(70.5,113)(.033333,-.058333){30}{\line(0,-1){.058333}}
%\end
\put(71.5,111.25){\line(-1,0){58.75}}
%\emline(12.75,111.25)(14.5,112.25)
\multiput(12.75,111.25)(.058333,.033333){30}{\line(1,0){.058333}}
%\end
\put(52,144.75){\line(0,-1){33.25}}
\put(49,114.25){\line(0,-1){3}}
\put(49,114.25){\line(1,0){3}}
\put(37.5,112.75){\vector(1,0){10.75}}
\put(29.75,112.75){\vector(-1,0){7.75}}
\put(58,113){\vector(-1,0){4.75}}
\put(27.75,128.25){\makebox(0,0)[cc]{$q$}}
\put(49.25,124.5){\makebox(0,0)[cc]{$h$}}
\put(66.5,127.75){\makebox(0,0)[cc]{$p$}}
\put(45,108.5){\makebox(0,0)[cc]{$r$}}
\put(33.5,112.75){\makebox(0,0)[cc]{$x$}}
\put(59.85,112.91){\makebox(0,0)[cc]{$y$}}
\put(18.5,113.25){\makebox(0,0)[cc]{$\theta$}}
\put(68.25,113){\makebox(0,0)[cc]{$\theta$}}
%\vector(61.69,112.85)(65.85,113)
\put(65.85,113){\vector(1,0){.07}}\multiput(61.69,112.85)(.83244,.02973){5}{\line(1,0){.83244}}
%\end
\end{picture}
\end{center}

    \begin{enumerate}
      \item[\hspace{2cm}b)] You want to show that $h = \frac{r}{(\cot \phi + \cot \theta)}$. Start with $\cot \phi + \cot \theta$ and see what that is.
      \item[\hspace{2cm}c)] $\cot \phi + \cot \theta = (\frac{x}{h}) + (\frac{y}{h})$.
      \item[\hspace{2cm}d)] Note that $x+y = r$.
    \end{enumerate}

If you spent a considerable amount of time trying to get this formula directly, you will really appreciate this approach to the problem. The direct approach is almost impossible. This approach is so simple it almost seems like cheating.

\vspace{2cm}

\begin{center}
Does your doctor measure up?: Hospitals are sued by 7 foot doctors.
\vspace{0.5cm}

Just Figures: Possible amnesty for tax delinquents on table.
\end{center}



\chapter*{Chapter VIII}
\addcontentsline{toc}{chapter}{\numberline{}Chapter VIII}
\markboth{Chapter VIII}{Chapter VIII}


There is another topic that we should cover in order to have done justice to the subject of trigonometry. The development of the sine and cosine functions is by far the most significant aspect of the study of trigonometry, but there is something else that we need to
delve into. 

It has been my observation in teaching high school mathematics over the last century that, after a student has taken algebra four times, that he's is ready for calculus. Now, before you give up and pronounce me brain-dead, let me explain. You take Algebra I. That's the first time through it. Then, you take Geometry. You should have gotten a pretty good review of Algebra 1 while you were taking Geometry. Not that it's really part of Geometry, but you should have gotten the review in some homework at the very least. That's twice. Then you take Algebra II. Here you developed some pretty sophisticated techniques for manipulating with the symbols, and Algebra I looks easy in retrospect. That's thrice! Then you take Trigonometry. You've had some algebraic workout in this course, but not all that you need. 

Here's the deal. Assuming that you're going to take Calculus, you need certain preparation. I don't mean that you need to be familiar with a bunch of esoteric mumbo-jumbo. The preparation that you need is to be able to do algebraic manipulation with ease and confidence. We are entering the stage of this course where that ease and confidence will be built. This is where, at least in my experience, the student develops his algebraic prowess to an extent where it is a useful tool to him, not a stumbling-block. 

You see, almost all of the development of the tools of calculus depend upon a lot of algebra. I can't tell you how many times I have seen an otherwise very bright student stumble and fall in calculus not because he is unable to understand the calculus but because he cannot follow the algebra. He has to spend so much of his time and effort trying to follow the algebra and fill in the gaps that he misses the point of the discussion. To miss the beauty and power of a theorem in calculus because you are not proficient in algebra is unfortunate to say the least. 

I think a lot of people have understood this, and perhaps this is why almost all trigonometry texts have a section devoted to \emph{identities}. What on earth are identities? Regardless of what they are, their purpose here is to give you a thorough workout with your algebra so that you'll be prepared for the most exciting and satisfying mathematics course you may ever take.

So, what are identities? Our five main properties of sine and cosine are examples of identities; i.e., $\sin^{2} x + \cos^{2} x = 1$ and the four expansion formulas.

Here are some that I'll bet you can show without developing a brain cramp.
\begin{enumerate}[\hspace{1cm}a)]
  \item $\tan^{2} x  + 1 = \sec^{2} x $ if $\cos x  \ne 0$.
  \item $1 + \cot^{2} x = \csc^{2} x $ if $\sin x  \ne 0$.
\end{enumerate}

Do you remember all those formulas we got like $\sin (2x) = 2 \sin x \cos x$? Those are all identities. For our purposes, they are equations that are valid everywhere the terms are defined. 

Note this one: $\sin^{2} x + \cos^{2} x = 1$. This statement is true regardless of what number $x$ is, since $\sin x$ and $\cos x$ are defined everywhere. But look at $\tan^{2} x  + 1 = \sec^{2} x $. Neither $\tan x$ nor $\sec x$ is defined if $\cos x = 0$. Hence, this is an identity as long as $\cos x \ne 0.$ Where is $\cos x = 0$? If $x$ is $\frac{\pi}{2}$, $-\frac{\pi}{2}$, $\frac{3 \pi}{2}$, $-\frac{3 \pi}{2}, \cdots$ . In other words, if $x$ is an odd multiple of $\frac{\pi}{2}$ -- i.e., $x = (2n + 1) \frac{\pi}{2}$ where $n$ is an integer -- then $\cos x = 0$. 

See if you can show that each of the following is an identity. 

Here are a couple of hints.

Work from one side only. When you are asked to show that an equation is an identity, you are being asked to show not only that the equation is valid (i.e., there is \emph{some} number for which it is true), but that it is valid for all numbers for which the functions
involved are defined. So don't start off assuming that the equation is valid. 

Here is a pitiful example. Show that $5 = 6$ is an identity. Well, obviously it's not one, but I want to illustrate what happens if you start off assuming what you're trying to prove. So, to show $5 = 6$ is an identity, let's start with $5 = 6$. Then $6 = 5$. Okay? Then $11 = 11$. I simply added. And since we know that $11 = 11$ is a valid statement, it must follow that $5 = 6$.

Clearly, something is wrong here. Of course, the problem is that we started with something that was false (even if we didn't know it) and ended with something that we know to be true. What did we prove? If $5 = 6$, then $11 = 11$. Startling, isn't it? What if we had started with some assumption (not knowing whether it is true or false), and, as a result of that assumption, we arrive at something we know to be false? What would that mean? Consider the previous example. 

\begin{quote}
  Conjecture: $5 = 6$

  then $5 + (-5) = 6 + (-5)$ (we added -5 to both sides)

  $0 = 1$ (do the arithmetic)

  But we know that $0 \ne 1$. It's an axiom from Algebra. Remember? Hence, if $5 = 6$, $0 = 1$. Since $0 \ne 1$, it follows that $5 \ne 6$.
\end{quote}

All this has to do with your handling of the identities. You will be tempted to write the equation down (assuming that it's true) and work with both sides. Don't! If you get stuck, reduce everything to sines and cosines. You'll have fun doing some of them. You'll find the hints section helpful with some of these. It will introduce you to some algebraic wizardry which you can use to amaze and astound you friends. Some of these are really difficult; but then others are really, really difficult. 

Please don't lose sight of what these are for. You will never need any of these relationships. So, their only purpose is to develop your manipulative prowess in algebra. You will learn some neat tricks for handling things that you will find very satisfying (after a great deal of frustration, of course). Onward!



\newpage
\begin{center}
\section*{Identities}
\end{center}

Show that each of the following statements is either:
\begin{enumerate}[\hspace{1cm}1)]
    \item true at each number $x$, or
    \item if it is not true at some number $x$, then stipulate the number or set of numbers on which it is not true, and show that it is true everywhere else.
\end{enumerate}

\renewcommand\theenumi{\Roman{enumi}}
\renewcommand\theenumii{\arabic{enumii}}

\begin{enumerate}

  \item \begin{enumerate}  % I
    \item $\sin x \cot x = \cos x $
    \item $\frac{\sin x}{\tan x} = \cos(-x) $
    \item $\sin(-x) \frac{\sec x}{\tan(-x) } = 1$
    \item $\tan^{2} x \cos^{2} x = \sin^{2} x $
    \item $\cos^{2} x  (1 + \tan^{2} x) = 1$
    \item $(\csc^{2} x - 1) \sin^{2} x = \cos^{2} x $
    \item $\frac{(1 + \sin x) (1 - \sin x)}{\cos x} = \cos x $
    \item $\frac{(1 + \tan^{2} x)}{\csc^{2} x} = \tan^{2} x $
    \item $\frac{\sin^{2} x}{\tan^{2} x - \sin^{2} x} = \cot^{2} x $
    \item $\frac{1}{\csc x - \sin x} = \frac{\tan x}{\cos x}$

          \vspace{0.3cm}

          Pretty easy so far, aren't they?
  \end{enumerate}

  \item \begin{enumerate}  % II
    \item $\cos x \tan x \csc x = 1$
    \item $\sec x - \cos x = \tan x \sin x $
    \item $\cos^{2} x - \sin^{2} x = \cos(2x) $
    \item $\cos^{2} x - \sin^{2} x = 2 \cos^{2} x - 1$
    \item $\cos^{2}(2x) = 1 - 2 \sin^{2} x $
    \item $(\sec^{2} x - 1) (\csc^{2} x - 1) = 1$
    \item $\tan^{2} x + \sec^{2} x = \frac{2}{\cos^{2} x} - 1$
    \item $(\tan x + \cot x) \frac{\sin x}{\cos x} = 1$
    \item $\frac{\cos x + \sin x}{\sin x} = 1 + \cot x $
    \item $\frac{1 - \tan^{2} x}{\tan x} = \cot x - \tan x$
    \item $\tan^{2} x - \sin^{2} x = \tan^{2} x \sin^{2} x$
    \item $\frac{1}{1 + \sin x} + \frac{1}{1 - \sin x} = \frac{2}{\cos^{2} x}$
    \item $\tan x + \cot x = \sec x \csc x$
    \item $\frac{1 + \tan^{2} x}{\tan^{2} x} = \csc^{2} x$
  \end{enumerate}

  \item \begin{enumerate}  % III
    \item $\cos^{4} x - \sin^{4} x = \cos(2x)$
    \item $\frac{1 + \sec x}{\sec x} = \frac{\sin^{2} x}{1 - \cos x}$
    \item $\frac{1 - \cos x}{\sin x} = \frac{\sin x}{1 + \cos x}$
    \item $\frac{1 + \sin x}{\cos x} = \frac{\cos x}{1 - \sin x}$
    \item $\csc x = \frac{\sec x + \csc x}{1 + \tan x}$
    \item $\frac{1}{\sec x - \tan x} = \sec x + \tan x$
    \item $\frac{\sec x - 1}{\sec x + 1} = \frac{1 - \cos x}{1 + \cos x}$
    \item $\frac{\cos^{2} x}{1 - \sin x} = \frac{\cos x}{\sec x - \tan x}$
    \item $\sec^{2} x + \csc^{2} x = \sec^{2} x   \csc^{2} x$
    \item $\tan^{4} x + \tan^{2} x = \sec^{4} x - \sec^{2} x$
    \item $\cos(-x) \tan(-x) \csc x = -1$
    \item $\frac{ \cos(-x) - \sin(-x) }{\cos x} = 1 + \tan x$

          \vspace{0.3cm}
          
          Are they getting a little more interesting? Not yet? Well, here are some of a slightly different nature.
  \end{enumerate}

  \item \begin{enumerate}  % IV
    \item For this problem, find each of the six trigonometric functions in terms of sine only. For example, $\sin x = \sin x $. That was easy. What about $\cos x$ in terms of $\sin x$? Well, $\sin^{2} x + \cos^{2} x = 1$ so that $\cos^{2} x = 1 - \sin^{2} x$, which leads to $\cos x = \pm \sqrt{1 - \sin^{2} x}$. Get the idea? Now find $\tan x$, $\cot x$, $\sec x$, and $\csc x$ in terms of $\sin x$ only.
    \item Find all six trig functions in terms of $\cos x$ only.
    \item Find all six trig functions in terms of $\tan x$ only.
    \item Find all six trig functions in terms of $\cot x$ only.
    \item Find all six trig functions in terms of $\sec x$ only.
    \item Find all six trig functions in terms of $\csc x$ only.
  \end{enumerate}

  \item Now, back to the more mundane identities. \begin{enumerate}  % V
    \item $\sin x \cos x \sec x \csc x = 1$
    \item $\cos x + \tan x \sin x = \sec x $
    \item $\frac{\tan x - \cot x}{\tan x + \cot x} = 2 \sin^{2} x - 1$
    \item $\frac{\sin x}{1 + \cos x} + \frac{1 + \cos x}{\sin x} = 2 \csc x$
    \item $\cos^{2} x - \sin^{2} x = \frac{ 1 - \tan^{2} x }{ 1 + \tan^{2} x }$
    \item $\cot x + \tan x = \cot x \sec^{2} x$
    \item $(\csc x - \cot x)^2 = \frac{1 - \cos x}{1 + \cos x}$
    \item $(\sec x - \tan x)^2 = \frac{1 - \sin x}{1 + \sin x}$
    \item $(\cos x - \sin x)^2 + 2  \sin x \cos x = 1$
    \item $\frac{ \cot^{2} x - 1}{1 + \cot^{2} x} = 2 \cos^{2} x - 1$
    \item $\frac{\tan x}{1 - \cot x} + \frac{\cot x}{1 - \tan x} = 1 + \tan x + \cot x$
    \item $\frac{2 \sin^{2} x - 1}{\sin x \cos x} = \tan x - \cot x$
    \item $\sec x \csc x - 2  \cos x \csc x = \tan x - \cot x$
    \item $\frac{\cos x - \sin x}{\cos x + \sin x} = \frac{\cot x - 1}{\cot x  + 1}$
    \item $\frac{\sin x}{\csc x - \cot x} = 1 + \cos x$
    \item $\frac{ (\cos^{2} x - \sin^{2} x{^2} }{ \cos^{4} x - \sin^{4} x } =1-2 \sin^{2} x$
    \item $\frac{\sec^{2} x}{1 + \sin x} = \frac{ \sec^{2} x - \sec x \tan x }{\cos^{2} x}$
    \item $1 + \cot x = \frac{(1 - \cot^{2} x) \sin x}{\sin x - \cos x}$
    \item $\frac{\sec x + \tan x}{\cos x - \tan x - \sec x} = - \csc x$
    \item $\frac{\sin x + \cos x + \sin x}{\cot x} =\frac{\sec x + \csc x - \cos x}{\tan x}$
  \end{enumerate}

  \item I think you will find these more challenging. \begin{enumerate}  % VI 
    \item $(\sin x + \cos x)^2 = 1 + \sin (2x)$
    \item $\cos^{4} x - \sin^{4} x = \cos (2x)$
    \item $2 \cos^{2} x = 1 + \cos (2x)$
    \item $\frac{\cos x + \sin x}{\cos x - \sin x} = \tan (2x) + \sec (2x)$
    \item $2 \cos^{2}(\frac{x}{2} ) - \cos x = 1$
    \item $2 \csc (2x) \cot x = 1 + \cot^{2} x$
    \item $2 \sin x + \sin (2x) = \frac{2  \sin^{3} x}{1 - \cos x}$
    \item $\tan (3x) \cos x - \sin x = \sin (2x) \sec (3x)$
    \item $2 \cot (2x) = \cot x - \tan x$
    \item $\tan (2x) - \tan x = \tan x \sec (2x)$
    \item $\tan (\frac{x}{2}) + \cot(\frac{x}{2}) = 2 \csc x$
    \item $\frac{\sin x - \cos(2x)}{\cos^{2} x} = \frac{2 \sin x - 1}{1 - \sin x}$
    \item $\sin (\frac{x}{2}) = \frac{ 1 - \cos x }{ 2  \sin (\frac{x}{2}) }$
    \item $\cot (\frac{x}{2}) - \tan(\frac{x}{2}) = 2 \cot x$
    \item $\sin (4x) ( \cos^{2} (2x) - \sin^{2} (2x) ) = \sin (\frac{8x}{2})$
    \item $\frac{2}{1 + \cos(2x) } = \sec^{2} x$
    \item $\frac{\sin (3x) }{\sin x} + \frac{\cos (3x) }{\cos x} = 4 \cos (2x)$
    \item $\frac{\sin (3x) }{\sin x} - \frac{\cos (3x) }{\cos x} = 2$
    \item $\frac{3 \cos x + \cos (3x) }{3 \sin x - \sin (3x) } = \cot^{3} x$
    \item $\cos^{6} x - \sin^{6} x = (1 - \sin^{2} x \cos^{2} x) \cos (2x)$
    \item $\sec x + \sec (2x) \sec x + \sec (3x) \sec (2x) = \csc x \tan (3x)$

          \vspace{0.3cm}
          
          Have they gotten interesting yet?
  \end{enumerate}

  \item Some of these involve both an $x$ and a $y$. I state that so you won't think the $y$s are misprints.
  \begin{enumerate}  % VII
    \item $\sin (2x) \cos (3x) + \cos (2x) \sin (3x) = \sin (5x)$
    \item $\cos x \cos (2x) - \sin x \sin (2x) = \cos (3x)$
    \item $\sin (x+y) \cos y - \cos (x+y) \sin y = \sin x$
    \item $\cos (x-y) \cos y - \sin (x-y) \sin y = \cos x$
    \item $\frac{ \tan (x+y) - \tan y }{ 1 + \tan (x+y) \tan y } = \tan x$
    \item $\csc x (\tan (3x) - \tan (2x) ) = \sec (2x) \sec (3x)$
    \item $\frac{\cos x -\cos y}{\cos x +\cos y} =-\tan (\frac{x-y}{2}) \tan(\frac{x+y}{2})$
    \item $\frac{\cos (3x) - \cos (4x)}{\sin (4x) + \sin (3x)} = \tan (\frac{x}{2})$
    \item Suppose $n$ is an integer. $\frac{\cos ((n-2)x) - \cos (nx)}{\sin ((n-2)x) + \sin (nx)} = \tan x$
    \item $\sin (x+y) \sin (x-y) = \sin^{2} x - \sin^{2} y$
    \item $\cos (x+y) \cos (x-y) = \cos^{2} x - \sin^{2} y$
    \item $\sin (5x) - \sin (\frac{\pi}{2}-x) = 2  \cos (2x+\frac{\pi}{4}) \sin (3x-\frac{\pi}{4})$
    \item $\cos (3x) + \cos(\frac{\pi}{2}-x) = 2  \cos (x+\frac{\pi}{4}) \cos (2x-\frac{\pi}{4})$
    \item $\frac{\sin (2x-y) + \sin y}{\cos (2x-y) + \cos y} = \tan x$
    \item $\frac{\sin (2x) + \sin (2y)}{\cos (2x) + \cos (2y)} = \tan (x+y)$
    \item $\cos (3x-2y) + \cos (3x+2y) = 2 \cos (3x) \cos (2y)$
    \item $\sin x (\sin x + \sin (3x) ) = \cos x (\cos x - \cos (3x) )$
    \item $\sin x - \cos (2x) - \sin (3x) = - \cos (2x) (2 \sin x + 1)$
    \item $\sin (5x) = 2 \cos (4x) \sin x + \sin (3x)$
    \item $\sin (3x) = 2 \sin (2x) \cos x - \sin x$
    \item $\frac{\sin (5x) + \sin (3x)}{\sin (2x) \cos (2x)} = 4 \cos x$
    \item $\frac{\sin (6x) + \sin (4x)}{2 \sin x} = \cos (5x)$
    \item $\frac{\cos (3x) + \cos x}{2 \cos^{3} x - 2 \cos x} = 2 - \csc^{2} x$
    \item $\sin (4x) = 2 \sin x \cos (3x) + \sin (2x)$
    \item $\sin x + \sin (3x) + \sin (5x) = \frac{\sin^{2} (3x)}{\sin x}$
    \item $\frac{ \sin (3x) - \sin (2x) }{ 2 \cos (2x) +4 \sin^{2} (x/2) - 1 } = \sin x$
    \item $4 \cos (\frac{x}{3}) \cos (\frac{2\pi+x}{3}) \cos (\frac{4 \pi+x}{3}) = \cos x$
    \item If $\tan x + \tan y + \tan z = \tan x \tan y \tan z$, \\
          show that $\tan(2x) + \tan(2y) + \tan(2z) = \tan(2x) \tan(2y) \tan(2z)$
  \end{enumerate}

\end{enumerate}


\newpage
\begin{center}
\section*{Chapter VIII Hints and Solutions}
\end{center}

\begin{enumerate}

  \item \begin{enumerate}  % I
    \item Use the definition of $\cot x$.
    \item Use the definition of $\tan x$ and note that $\cos(-x) = \cos x$.
    \item Writing all the entries on the left in terms of sines and cosines will help.
    \item Same as 3 above.
    \item Same as 3 above.
    \item Same as 3 above.
    \item Recall $(1-a)(1+a) = 1 - a^2$.
    \item Write everything on the left in terms of sines and cosines.
    \item Same as 8 above.
    \item Same as 8 above.
    \end{enumerate}

  \item \begin{enumerate}  % II
    \item No hints.
    \item No hints.
    \item No hints.
    \item No hints.
    \item No hints.
    \item No hints.
    \item \begin{enumerate}[\hspace{1cm}a)]
            \item Write everything on the left in terms of sines and cosines.
            \item Note that $\sin^{2} x = 1 - \cos^{2} x$.
          \end{enumerate}
    \item No hints.
    \item No hints.
    \item No hints.
    \item No hints.
    \item Add the fractions on the left together, and recall that $(1+a)(1-a)= 1-a^2$.
    \item \begin{enumerate}[\hspace{1cm}a)]
            \item Write the left side in terms of sine and cosine, and add the resulting fractions together.
            \item Note that $\sin^{2} x + \cos^{2} x = 1$.
            \item Recall the definition of $\sec x$ and $\csc x$.
          \end{enumerate}
    \item No hints.

          \vspace{0.3cm}
          
          You may want to go back over some of these problems to determine if there are some restrictions on where the equations are valid. For example, can you see that in 12), $\sin x  \ne \pm 1$? Also, $\cos x \ne 0$. $\cos x = 0$ if $\sin x = 1 or -1$, so we need to find where $\sin x = -1 or 1$. $\sin x = 1$ if $x = \frac{\pi}{2}, 2\pi + \frac{\pi}{2}, -2\pi + \frac{\pi}{2}, 2(2\pi) + \frac{\pi}{2}, -2(2\pi) + \frac{\pi}{2} \dots$ . In general, $\sin x = 1$ if $x = (2n)\pi + \frac{\pi}{2}$ where $n$ is an integer. Where is $\sin x = -1$? See if you can show that $\sin x = -1$ if $x=(2n +1)\pi + \frac{\pi}{2}$ where $n$ is an integer. What all this means is that $\frac{1}{1 + \sin x} + \frac{1}{1 - \sin x} = \frac{2}{\cos^{2} x}$ as long as $\sin x \ne -1$, $\sin x \ne 1$, and $\cos x \ne 0$. These conditions occur if $x = (2n)\pi + \frac{\pi}{2}$ or $x = (2n+1)\pi + \frac{\pi}{2}$ where $n$ is an integer.

          \vspace{0.3cm}
          
          Can you see that the combination of these two conditions $x = (2n)\pi + \frac{\pi}{2}$ and $x = (2n + 1)\pi + \frac{\pi}{2}$ is equivalent to saying $x = (2n+1) \frac{\pi}{2}$, so $x$ is an odd multiple of $\frac{\pi}{2}$?

          \vspace{0.3cm}
          
          Hence, we write our conclusion as follows: $\frac{1}{1 + \sin x} + \frac{1}{1 - \sin x} = \frac{2}{\cos^{2} x}$ if $x \ne (2n+1) \frac{\pi}{2}$ where $n$ is an integer. The part ``if $x \ne (2n+1)\frac{\pi}{2}$ where $n$ is an integer'' states the conditions under which the equation is valid. Ignoring the conditions under which equations are valid can be quite serious, especially if the engineer that built the airplane you're riding in did so.
    \end{enumerate}


  \item \begin{enumerate}  % III
    \item \begin{enumerate}[\hspace{1cm}a)]
            \item Factor $\cos^{4} x - \sin^{4} x$.
            \item Recall $a^4 - b^4 = (a^2 + b^2)(a^2 - b^2)$.
            \item $sin^{2} x + \cos^{2} x = 1$.
            \item Note that $\cos^{2} x - \sin^{2} x = \cos(x+x)$.
          \end{enumerate}
    \item \begin{enumerate}[\hspace{1cm}a)]
            \item Write the entries on the left in terms of $\cos x$.
            \item Multiply $\frac{1+ \cos x}{1}$ by $1$, but write $1$ as $\frac{1- \cos x}{1- \cos x}$. Tricky, but effective.

                  \vspace{0.3cm}
                  
                  It turns out that two of the most difficult operations in mathematics are multiplying by $1$ and adding $0$. This sounds silly, but the trick is writing $1$ or $0$ in some disguised form that gives you a result that you want. It's certainly not obvious that you want to write $1$ as $\frac{1- \cos x}{1- \cos x}$. But in this problem, you got $\frac{1+ \sec x}{\sec x}$ reduced to $1+ \cos x$, and you could see that what you wanted was $\frac{\sin x}{1- \cos x}$. Multiplying $1+ \cos x$ by $\frac{1- \cos x}{1- \cos x}$ will give you what you want.

                  \vspace{0.3cm}
                  
                  I can imagine what you're thinking. What blue sky rack did \emph{that} come from? Experience is probably the best answer. And these problems will give you some experience in such witchcraft so that you will look for such things in your future attempts to untie algebraic knots.
          \end{enumerate}
    \item \begin{enumerate}[\hspace{1cm}a)]
            \item Multiply the left side by $1$ (in disguise, of course).
            \item Let $1 = \frac{1 + \cos x}{1 + \cos x}$ in a) above.
          \end{enumerate}
    \item Multiply by $1$.
    \item Work with the right side here and convert everything on the right to sines and cosines.
    \item \begin{enumerate}[\hspace{1cm}a)]
            \item Reduce the left side to $\frac{\cos x}{1- \sin x}$.
            \item Multiply by $\frac{1+ \sin x}{1+ \sin x}$
            \item Write $\frac{1+ \sin x}{\cos x}$ as two fractions.
          \end{enumerate}
    \item No hints.
    \item Work with the right side.
    \item Convert to sines and cosines and add to make one fraction.
    \item \begin{enumerate}[\hspace{1cm}a)]
            \item Factor $\tan^{2} x$ out.
            \item Recall $\tan^{2} x + 1 = \sec^{2} x$.
            \item Write $\tan^{2} x = \sec^{2} x - 1$.
          \end{enumerate}
    \item No hints.
    \item No hints.
  \end{enumerate}

  \item \begin{enumerate}  % IV
    \item \begin{enumerate}[\hspace{1cm}a)]
            \item $\tan x = \frac{\sin x}{\cos x}=\frac{\sin x}{\pm\sqrt{1-\sin^{2} x}}$
            \item $\cot x = \frac{\cos x}{\sin x}=\frac{\pm\sqrt{1-\sin^{2} x}}{\sin x}$
            \item $\sec x = \frac{1}{\cos x} = \frac{1}{\pm \sqrt{1-\sin^{2} x}}$
            \item $\csc x = \frac{1}{\sin x}$
          \end{enumerate}
    \item \begin{enumerate}[\hspace{1cm}a)]
            \item $\sin x = \pm \sqrt{1 - \cos^{2} x}$
            \item $\cos x = \cos x$
            \item $\tan x = \frac{\sin x}{\cos x}=\frac{\pm\sqrt{1-\cos^{2} x}}{\cos x}$
            \item $\cot x = \frac{\cos x}{\sin x}=\frac{\cos x}{\pm\sqrt{1-\cos^{2} x}}$
            \item $\sec x = \frac{1}{\cos x}$
            \item $\csc x = \frac{1}{\sin x} = \frac{1}{\pm \sqrt{1 - \cos^{2} x}}$
          \end{enumerate}
    \item First, note that $tan^{2} x + 1 = \sec^{2} x$ so that $\frac{1}{cos^{2} x} = \tan^{2} x + 1$ and $\cos x = \frac{1}{\pm \sqrt{\tan^{2} x + 1}}$.
          \begin{enumerate}[\hspace{1cm}a)]
            \item $\sin x = \pm\sqrt{1-\cos^{2} x} = \pm\sqrt{1-\frac{1}{\tan^{2} x+1}} = \pm\sqrt{\frac{\tan^{2} x}{\tan^{2} x+1}} = \frac{\tan x}{\pm\sqrt{\tan^{2} x+1}}$
            \item $\cos x = \frac{1}{\pm \sqrt{\tan^{2} x + 1} }$
            \item $\tan x = \tan x$
            \item $\cot x = \frac{1}{\tan x}$
            \item $\sec x = \frac{1}{\cos x} = \pm \sqrt{\tan^{2} x + 1}$
            \item $\csc x = \frac{1}{\sin x} = \pm \sqrt{ \frac{\tan^{2} x + 1}{\tan^{2} x} } = \frac{\pm \sqrt{\tan^{2} x + 1}}{\tan x}$
          \end{enumerate}
    \item Note $1 + \cot^{2} x = \csc^{2} x = \frac{1}{\sin^{2} x}$. Hence,
          \begin{enumerate}[\hspace{1cm}a)]
            \item $\sin x = \frac{1}{\pm \sqrt{1+ \cot^{2} x}}$
            \item $\cos x = \pm\sqrt{1-\sin^{2}x} = \pm\sqrt{1-\frac{1}{1+\cot^{2} x}} = \pm\sqrt{\frac{\cot^{2} x}{1+ \cot^{2}x}} = \frac{\cot x}{\pm\sqrt{1+ \cot^{2} x}}$
            \item $\tan x = \frac{1}{\cot x}$
            \item $\sec x = \frac{1}{\cos x} = \pm \sqrt{\frac{1+ \cot^{2} x}{ \cot^{2} x}} = \frac{\pm \sqrt{1+ \cot^{2} x}}{\cot x}$
            \item $\cot x = \cot x$
            \item $\csc x = \pm \sqrt{1+ \cot^{2} x}$
          \end{enumerate}
    \item \begin{enumerate}[\hspace{1cm}a)]
            \item $\sin x = \pm \sqrt{1 - \cos^{2} x} = \pm \sqrt{1 - \frac{1}{\sec^{2} x}} = \pm \sqrt{\frac{\sec^{2} x -1}{\sec^{2} x}} = \frac{\pm \sqrt{\sec^{2} x -1}}{\sec x}$
            \item $\cos x = \frac{1}{\sec x}$
            \item $\tan x = \frac{\sin x}{\cos x} = \pm \sqrt{\frac{\sec^{2} x -1}{\sec^{2} x}}  \sec x = \pm \sqrt{\sec^{2} x -1}$
            \item $\cot x = \frac{1}{\tan x} = \frac{1}{\pm \sqrt{\sec^{2} x -1}}$
            \item $\sec x = \sec x $
            \item $\csc x = \frac{1}{\sin x} = \pm \sqrt{\frac{\sec^{2} x}{\sec^{2} x -1}} = \frac{\sec x}{\pm \sqrt{\sec^{2} x -1}}$
          \end{enumerate}
    \item \begin{enumerate}[\hspace{1cm}a)]
            \item $\sin x = \frac{1}{\csc x}$
            \item $\cos x = \pm \sqrt{1- \sin^{2} x} = \pm \sqrt{1- \frac{1}{\csc^{2} x}} = \pm \sqrt{\frac{\csc^{2} x -1}{\csc^{2} x}} = \frac{\pm \sqrt{\csc^{2} x -1}}{\csc x}$
            \item $\tan x = \frac{\sin x}{\cos x} = \frac{\frac{1}{\csc x}}{\pm \sqrt{\frac{\csc^{2} x -1}{\csc^{2} x}}} = \frac{1}{\pm \sqrt{\csc^{2} x -1}}$
            \item $\cot x = \frac{1}{\tan x} = \pm \sqrt{\csc^{2} x -1}$
            \item $\sec x = \frac{1}{\cos x} = \pm \sqrt{\frac{\csc^{2} x}{\csc^{2} x -1}} = \frac{\csc x}{\pm \sqrt{\csc^{2} x -1}}$
            \item $\csc x = \csc x$

            \vspace{0.3cm}
            
            I'm sure you'll use all this material every day of your life. But, just in case you don't, you now have a handy reference in the unlikely event that some of this should fade from your memory. 
          \end{enumerate}
    \end{enumerate}

    \item \begin{enumerate}  % V
        \item No hints.
        \item Write the left side as two fractions and add them together.
        \item Write the left side in terms of sines and cosines and add the fractions.
        \item Remember that $\sin^{2} x + \cos^{2} x = 1$.
        \item Work with the right side.
        \item Write the left side in terms of sines and cosines and add the fractions.
        \item \begin{enumerate}[\hspace{1cm}a)]
                \item Write the left side in terms of sines and cosines
                \item Write $\sin^{2} x = 1 - \cos^{2} x$.
                \item Divide out $1 - \cos x $ from both the numerator and the denominator.
              \end{enumerate}
        \item \begin{enumerate}[\hspace{1cm}a)]
                \item Convert the left side to sines and cosines.
                \item Add the fractions and square the result.
                \item Write the denominator as $1- \sin^{2} x$.
              \end{enumerate}
        \item No hints.
        \item \begin{enumerate}[\hspace{1cm}a)]
                \item Convert the left side to sines and cosines.
                \item Add the fractions.
              \end{enumerate}
        \item \begin{enumerate}[\hspace{1cm}a)]
                \item Subtract the right side from the left side. If you can show that this difference is $0$, then you have verified that the original equation is an identity.
                \item Convert all terms to sines and cosines. I'm sure you've seen that this is a rather standard step. It's really not necessary, but some people are more comfortable dealing with sines and cosines.
                \item Reduce the expression to $\frac{\sin^{2} x}{\cos x (\sin x - \cos x)} - \frac{\cos^{2} x}{\sin x (\sin x - \cos x)} - (1+ \frac{1}{\cos x  \sin x})$
                \item Reduce the above to $(\frac{1}{\sin x - \cos x}) (\frac{1}{\cos x  \sin x}) (\sin^{3} x - \cos^{3} x) - (1+ \frac{1}{\cos x  \sin x})$
                \item Factor $\sin x - \cos x $ out of $\sin^{3} x - \cos^{3} x$
                \item Recall that $a^3 - b^3 = (a-b) (a^2 +ab + b^2)$.
              \end{enumerate}
        \item Work with the right side.
        \item \begin{enumerate}[\hspace{1cm}a)]
                \item Subtract the right side from  the left side.
                \item Convert all terms to sines and cosines.
                \item Add the four functions together.
              \end{enumerate}
        \item Work with the right side.
        \item No hints.
        \item Factor the denominator on the left.
        \item \begin{enumerate}[\hspace{1cm}a)]
                \item Work with the right side converting to sines and cosines.
                \item Write the right side with $\cos^{4} x$ as the denominator.
                \item Write $\cos^{4} x = \cos^{2} x (1 - \sin^{2} x)$.
                \item Divide out $1 - \sin x$.
              \end{enumerate}
        \item \begin{enumerate}[\hspace{1cm}a)]
                \item Work with the right side.
                \item Factor $1 - \cot^{2} x$.
                \item Multiply $\sin x$ times one of the factors in b).
              \end{enumerate}
        \item Convert the left side to sines and cosines.
        \item \begin{enumerate}[\hspace{1cm}a)]
                \item Subtract the right side from the left side and convert to sines and cosines. Again, the object is to show that this difference is $0$.
                \item Add all of your terms together, with a common denominator of $\sin x \cos x$.
                \item Find a way of grouping the terms in the numerator in order to factor out the term $\sin x + \cos x$.
              \end{enumerate}
    \end{enumerate}

    \item \begin{enumerate} % VI
        \item No hints.
        \item Factor the left side.
        \item No hints.
        \item \begin{enumerate}[\hspace{1cm}a)]
                \item Work with the right side converting to sines and cosines. Don't overlook the $2$ in $\tan(2x)$ and $\sec(2x)$.
                \item Write $\sin(2x)$ and $\cos(2x)$ in terms of $\sin x$ and $\cos x$.
                \item Here's a neat trick. In your numerator, write $1$ as $\sin^{2} x + \cos^{2} x$.
                \item Your numerator should be a perfect square and you should be able to write it as $(\sin x + \cos x)^2$.
                \item Divide out what you can from both numerator and denominator.
              \end{enumerate}
        \item No hints.
        \item Recall that $\sin(2x) = 2 \sin x \cos x$ and that $1 + \cot^{2}x = \csc^{2}x$.
        \item After you get a little ways into this one, multiply by $1$ disguised as $\frac{1 - \cos x}{1 - \cos x}$.
        \item \begin{enumerate}[\hspace{1cm}a)]
                \item Convert the left side to sines and cosines.
                \item Add the two terms on the left together getting $\cos(3x)$ in the denominator.
                \item You should have $\sin(3x) \cos x - \sin x \cos(3x)$ in the numerator. Do you recognize this as fitting one of our expansion formulas?
                \item $\sin(3x) \cos x - \sin x \cos(3x) = \sin(3x-x) = \sin(2x)$
              \end{enumerate}
        \item \begin{enumerate}[\hspace{1cm}a)]
                \item Write $\cot(2x)$ in terms of $\cos(2x)$ and $\sin(2x)$.
                \item Expand $\cos(2x)$ and $\sin(2x)$.
                \item Write the result as the difference of two fractions.
              \end{enumerate}
        \item \begin{enumerate}[\hspace{1cm}a)]
                \item Write the left side in terms of sines and cosines.
                \item Add the fractions together and recognize the numerator as fitting one of our expansion formulas.
              \end{enumerate}
        \item \begin{enumerate}[\hspace{1cm}a)]
                \item Convert the left side to sines and cosines and add the fractions together.
                \item You should have $\frac{1}{\sin(\frac{x}{2})\cos(\frac{x}{2})}$.
                \item Apply $\sin(2\theta) = 2 \sin(\theta) \cos(\theta)$ to what you have in b).
              \end{enumerate}
        \item \begin{enumerate}[\hspace{1cm}a)]
                \item Convert the numerator on the left to all sines and factor.
                \item Convert the denominator to sines.
              \end{enumerate}
        \item Note that $\frac{1 - \cos x}{2} = \sin^{2}(\frac{x}{2})$.
        \item \begin{enumerate}[\hspace{1cm}a)]
                \item Convert the left side to sines and cosines and add the fractions.
                \item Recognize that $\sin(\frac{x}{2} ) \cos(\frac{x}{2} ) = \frac{1}{2} ( 2 \sin(\frac{x}{2}) \cos(\frac{x}{2}) )$
                   \\ $= \frac{1}{2} \sin (\frac{x}{2} + \frac{x}{2}) = \frac{1}{2} \sin x$.
              \end{enumerate}
        \item \begin{enumerate}[\hspace{1cm}a)]
                \item Note that $\cos^{2}(2x) - \sin^{2}(2x) = \cos(2(2x)) = \cos(4x)$.
                \item Now note that $\sin(4x)\cos(4x)=\frac{1}{2}(2 \sin(4x) \cos(4x))$.
              \end{enumerate}
        \item No hints.
        \item \begin{enumerate}[\hspace{1cm}a)]
                \item Add the left side together, and note that it fits one of our expansion formulas.
                \item Write $\sin(4x)$ as $2 \sin(2x) \cos(2x)$.
              \end{enumerate}
        \item See 17 a).
        \item You'll probably have to grind a little on this one, expanding $\cos(3x)$ and $\sin(3x)$ until you've got them in terms of $\sin x$ and $\cos x$.
        \item \begin{enumerate}[\hspace{1cm}a)]
                \item Write $\cos^{6} x - \sin^{6} x = (\cos^{2} x)^3 - (\sin^{2} x)^3$.
                \item Factor the expression in a) and note that one of the factors can be written as $\cos(2x)$.
                \item Write $\cos^{4} x = \cos^{2} x (1 - \sin^{2} x)$ and $\sin^{4} x = \sin^{2} x (1 - \cos^{2} x)$.
                \item Everything should reduce to $1 - \cos^{2} x \sin^{2} x$ times the factor $\cos(2x)$.
              \end{enumerate}
        \item \begin{enumerate}[\hspace{1cm}a)]
                \item The trick here, as in a lot of these, is to maneuver so that you can take advantage of some of our basic expansion formulas.
                \item Subtract the right side from the left side and convert to sines and cosines.
                \item Factor $\frac{1}{\cos x}$ out of two terms, and $\frac{1}{\cos(3x)}$ out of the other two.
                \item Say ``Abra Cadabra'' three times, and reduce what you have to 
                   \\ $(\frac{1}{\cos(2x)}) (2 \cos x + \frac{\sin x - \cos(2x) \sin(3x)}{\sin x  \cos(3x)})$. 
                \item By skillful manipulation, you should be able to contract some things inside the right parentheses by the use of one of your expansion formulas.
                \item Lo and behold, the entry in the right parentheses is $0$!
              \end{enumerate}
        \end{enumerate}

    \item \begin{enumerate} % VII
        \item Direct application of one of your expansion formulas.
        \item No hints.
        \item Direct application of one of your expansion formulas.
        \item No hints.
        \item No hints.
        \item \begin{enumerate}[\hspace{1cm}a)]
                \item Convert the left side to sines and cosines.
                \item Add the fractions inside the parentheses and look for an application of one of your expansion formulas.
              \end{enumerate}
        \item \begin{enumerate}[\hspace{1cm}a)]
                \item Work with the right side and write everything in terms of sines and cosines.
                \item Use you expansion formulas, multiply numerators, multiply denominators, and collect terms.
                \item Write all the entries in the numerator in terms of sines and all those in the denominator in terms of cosines.
                \item Recall that $\sin(\frac{x}{2}) = \pm \sqrt{\frac{1 - \cos x}{2}}$ and $\cos(\frac{x}{2}) = \pm \sqrt{\frac{1+ \cos x}{2}}$.
              \end{enumerate}
        \item \begin{enumerate}[\hspace{1cm}a)]
                \item This is a particularly frustrating one if you don't see a neat little trick. Note that $\cos(a-b) - \cos(a+b) = 2 \sin a \sin b$ and that $\sin(a-b) + \sin(a+b) = 2 \sin a \cos b$. I can hear you. Don't think those words!
                \item Let $a-b=3x$ and $a+b=4x$.
                \item Then $a=\frac{7x}{2}$ and $b=\frac{x}{2}$.
                \item Write $\frac{\cos(3x) - \cos(4x)}{\sin(4x) + \sin(3x)} = \frac{\cos(a-b) - \cos(a+b)}{\sin(a+b) + \sin(a-b)}$.
                \item Cool, huh?
              \end{enumerate}
        \item Same procedure as in 8).
        \item No hints.
        \item No hints.
        \item \begin{enumerate}[\hspace{1cm}a)]
                \item Work with the right side and expand both expressions by your expansion formulas.
                \item The right side should reduce to
                   \\ $(\cos(2x) - \sin(2x) )(\sin(3x) - \cos(3x) )$.
                \item Multiply the expressions in b) and look for applications of your expansion formulas in the result.
                \item You should have the right side reduced to $\sin(5x) - \cos x$.
                \item Is $\cos x = \sin(\frac{\pi}{2}-x)$?
              \end{enumerate}
        \item Follow the same strategy as in 11).
        \item \begin{enumerate}[\hspace{1cm}a)]
                \item Work with the left side and use the trick you used in problem 8).
                \item Let $2x-y=a-b$ and $y=a+b$ in this application. Hence, $a=x$ and $b=y-x$.
              \end{enumerate}
        \item \begin{enumerate}[\hspace{1cm}a)]
                \item Try the same trick here by letting $a+b=2x$ and $a-b=2y$.
                \item Isn't this a neat little device? You'll run across this idea in your next course, where you will be interested in writing $\cos x - \cos y$ in a more convenient form. Here, you will let $x=a+b$ and $y=a-b$. Can you show that $\cos x - \cos y = -( \sin(\frac{x+y}{2}) \sin(\frac{x-y}{2}) )$?
              \end{enumerate}
        \item No hints.
        \item \begin{enumerate}[\hspace{1cm}a)]
                \item Subtract the right side from the left side.
                \item Look for an application of your expansion formulas.
              \end{enumerate}
        \item \begin{enumerate}[\hspace{1cm}a)]
                \item Subtract the right side from the left side.
                \item Reduce what is left to $\sin x - \sin(3x) +2 \sin x  \cos(2x)$.
                \item Use the trick from question 14 to write $\sin x - \sin(3x)$ as
                   \\ $- 2 \sin x \cos(2x)$.
                \item Let $x=a-b$ and $3x=a+b$ to get c).
                \item Then $a=2x$ and $b=x$.
                \item $\sin x - \sin(3x) = \sin(a-b) - \sin(a-b)$
                \item Hence, $\sin x - \sin(3x)=-2 \sin b \cos a = -2 \sin x  \cos(2x)$.
              \end{enumerate}
        \item \begin{enumerate}[\hspace{1cm}a)]
                \item Subtract the right side from the left side.
                \item Expand $\sin(5x)$ as $\sin(4x+x)$.
                \item Collect terms and look for an application of your expansion formulas.
                \item You should end up with $\sin(4x-x) - \sin(3x)$ which is $0$.
              \end{enumerate}
        \item \begin{enumerate}[\hspace{1cm}a)]
                \item Use the same technique as in 18).
                \item Expand $\sin(3x)$ as $\sin(2x+x)$.
              \end{enumerate}
        \item \begin{enumerate}[\hspace{1cm}a)]
                \item Use the cool trick on $\sin(5x) + \sin(3x)$ to write it as $2 \sin(4x)$ $\cos x$.
                \item $\frac{2\sin(4x)\cos x}{\sin(2x)\cos(2x)}$ reduces to $4\cos x$.
              \end{enumerate}
        \item \begin{enumerate}[\hspace{1cm}a)]
                \item The trick works on this one as well. Let $a+b=6x$ and $a-b=4x$.
                \item You can also just write $\sin(6x) - \sin(4x) -2 \sin x  \cos(5x)$ and show this is $0$ by expanding $\sin(6x)$ as $\sin(5x+x)$ and collecting terms.

                   \vspace{0.3cm}
                
                   I'm sure you are aware that the hints I've given you on these problems do \emph{not} in any way suggest that these are the only ways to solve them. In fact, in many cases my way may be much more complicated and obscure compared to the way you have done it. It has been my experience that, if I find an easy and elegant way to work a problem, it is only after having found the most difficult way possible first.
              \end{enumerate}
        \item \begin{enumerate}[\hspace{1cm}a)]
                \item Use the trick. Let $a+b=3x$ and $a-b=x$.
                \item The left side will then reduce to $\frac{\cos(2x)}{\cos^{2}x-1}$.
                \item Expand $\cos(2x)$ and write $\cos^{2} x - 1 = -sin^{2} x$.
                \item Recall that $1 + \cot^{2} x = \csc^{2} x$.
              \end{enumerate}
        \item \begin{enumerate}[\hspace{1cm}a)]
                \item Subtract the right side from the left side and expand $\sin(4x)$ as $\sin(3x+x)$.
                \item Collect terms and look for an application of your expansion formulas.
              \end{enumerate}
        \item \begin{enumerate}[\hspace{1cm}a)]
                \item Consider $\sin x (\sin x + \sin(3x) + \sin(5x) ) - \sin(2x)$ and show that this is $0$.
                \item Write $\sin x+\sin(3x)+\sin(5x)=\sin(3x-2x)+\sin(3x+2x)+\sin(3x)$.
                \item Hence, $\sin x+\sin(3x)+\sin(5x) = 2\sin(3x)\cos(2x)+\sin(3x)$.
                \item Factor out $\sin(3x)$, and work with the remainder 

                   $\sin x (1+2 \cos(2x) ) - \sin(3x)$. You want to show this is $0$.
                \item Expand $\sin(3x)$ and factor out $\sin x$ so that what you have left is $1+2 \cos(2x) -2 \cos^{2} x - \cos(2x)$.
                \item The expression in e) above reduces to $0$.
              \end{enumerate}
        \item \begin{enumerate}[\hspace{1cm}a)]
                \item Consider $\sin(3x) - \sin(2x) - \sin x (2 \cos(2x) + 4 \sin^{2} (\frac{x}{2} ) - 1)$ and show this is $0$.
                \item Expand $\sin(3x)$ and $\sin(2x)$ and then factor $\sin x$ out and work with the remainder.
                \item Convert everything to cosines.
              \end{enumerate}
        \item \begin{enumerate}[\hspace{1cm}a)]
                \item Expand $\cos(\frac{2+x}{3})$ and $\cos(\frac{4+x}{3})$ to get rid of the $\frac{2}{3}$ and $\frac{4}{3}$.
                \item Multiply the terms together. You should be able to get 

                   $\cos(\frac{x}{3}) (4 \cos^{2}(\frac{x}{3}) -3)$.
                \item Dredge up from your memory or show that 

                   $\cos(3\theta) = 4 \cos^{3} (\theta) - 3 \cos(\theta)$, and apply it to what you got in b).
              \end{enumerate}
        \item This is probably the nastiest of the whole bunch of identities. Here are some clues that may help if you get stuck.
              \begin{enumerate}[\hspace{1cm}a)]
                \item Show that $\frac{1}{\tan x \tan y} + \frac{1}{\tan x \tan z} + \frac{1}{\tan y \tan x} = 1$.
                \item Show that $\tan x \tan y + \tan x \tan z + \tan y \tan z = \frac{\tan x}{\tan z} + \frac{\tan y}{\tan z} + 1 + \frac{\tan x}{\tan y } + 1 + \frac{\tan z}{\tan y} + 1 + \frac{\tan y}{\tan x} + \frac{\tan z}{\tan x}$.

                      \vspace{0.3cm}
                      
                      If you haven't put in much effort on this problem, you won't see much relevance that these so-called clues have to the problem.
                \item Now consider $\frac{\tan(2x)+\tan(2y)+\tan(2z)}{\tan(2x) \tan(2y) \tan(2z)}$ and try to show that this is $1$.
                \item Recall (or show, if you don't remember) that
                      \begin{equation*}
                      \tan(2\theta) = \frac{2 \tan(\theta)}{1- \tan^{2}(\theta)}
                      \end{equation*}
                \item This will take some algebraic skill on your part, but with patience (with yourself), you'll be able to use the first two clues to your advantage.
                \item What conditions must be imposed on $x$, $y$, and $z$ in order for this equation to be valid?
              \end{enumerate}

    \end{enumerate}

\end{enumerate}

\vspace{2cm}

\begin{center}
I hope that you will ponder the following: It is not from the answers to our questions that we gain meaning; it is from the pursuit of those answers.
\end{center}


\end{document}
